\documentclass[11pt]{article}
\usepackage[T1]{fontenc}
\usepackage[utf8]{inputenc}
\DeclareUnicodeCharacter{2082}{\ensuremath{_2}}
\DeclareUnicodeCharacter{2022}{\ensuremath{\bullet}}
\DeclareUnicodeCharacter{2227}{\ensuremath{\wedge}}
\DeclareUnicodeCharacter{22A4}{\ensuremath{\top}}
\usepackage{lmodern,amsmath,amssymb,amsthm,mathtools}
\usepackage[a4paper,margin=27mm]{geometry}
\usepackage{booktabs,longtable,array,enumitem,microtype,xurl}
\usepackage[colorlinks=true,linkcolor=blue,citecolor=blue,urlcolor=blue,hypertexnames=false]{hyperref}
\usepackage{fancyhdr}
\pagestyle{fancy}\fancyhf{}\fancyhead[L]{THETA3: three signatures}\fancyhead[R]{21 November 2026}\fancyfoot[C]{\thepage}
\setlength{\headheight}{14pt}
\newcommand{\Q}{\mathbb Q}\newcommand{\C}{\mathbb C}\newcommand{\Z}{\mathbb Z}\newcommand{\R}{\mathbb R}\newcommand{\A}{\mathbb A}\newcommand{\OO}{\mathcal O}
\newcommand{\F}{\mathbb F}
\providecommand{\mathscr}[1]{\mathcal{#1}}
\newcommand{\Hdg}{\operatorname{Hdg}}\newcommand{\CH}{\operatorname{CH}}\newcommand{\Hom}{\operatorname{Hom}}\newcommand{\Tr}{\operatorname{Tr}}\newcommand{\SU}{\operatorname{SU}}\newcommand{\U}{\operatorname{U}}
\newcommand{\im}{\operatorname{im}}\newcommand{\vol}{\operatorname{vol}}\newcommand{\Res}{\operatorname{Res}}\newcommand{\Nm}{\operatorname{N}}\newcommand{\cl}{\operatorname{cl}}\newcommand{\Span}{\operatorname{span}}
\newcommand{\framework}[1]{\par\smallskip\noindent\textit{Framework and scope.} #1\par}
\newtheorem{theorem}{Theorem}[section]\newtheorem{proposition}[theorem]{Proposition}\newtheorem{lemma}[theorem]{Lemma}\newtheorem{corollary}[theorem]{Corollary}
\theoremstyle{definition}\newtheorem{definition}[theorem]{Definition}\newtheorem{remark}[theorem]{Remark}
\title{THETA3\\Theta kernels, definite periods and integral models\\\large Exact comparisons and quantitative certificates}
\author{Run THETA3-20261121\\GPT-6 Astra Pro, with independent agent proof checks}
\date{21 November 2026}
\begin{document}\maketitle
\begin{abstract}
We compare the CM Hodge construction with the theta and period mechanisms in Alp\"oge's archive, retaining the Hermitian spaces, splitting characters, coefficient types, measures and integral structures. Exact zero-orbit and restriction-seesaw identities give finite quantitative certificates and a positive finite-frame formula for the ball pairing. A direct construction on an $n$-ball realizes balanced $2n$-factor CM tensors by codimension-$n$ cycles; Liu's componentwise theorem supplies their precise CM sources in all ranks. A finite-dimensional kernel argument realizes the wedge as a genuine single theta lift; its true adjoint supplies a common double-lift and last-step Bessel identity. The completed Taylor tensor agrees with the matched canonical harmonic jet. For non-CM sources, finite \'etale covers of one arithmetic ball surface provide infinitely many pairwise nonisogenous elliptic Albanese factors, explicit new $j$-invariants, a simple genus-two Jacobian with endomorphism ring $\Z$, balanced non-CM Weil sources in every even dimension, and a geometrically simple exceptional Weil fourfold with endomorphism ring $\Z[i]$ on a ball cover of degree at most $43264$. A paired-branch construction gives explicit integral pure Weil cycles in every even dimension; its exceptional fourfold cycle has self-intersection $73728$, with a positive multiple realized on the fourth power of the actual ball cover. Exact Albanese decompositions explain why the corresponding map-only constructions vanish. Actual finite geometric hosts and standard theta inputs give constructive derivative interpolation, with a certified inverse. Finite CM theta modules give a common period normalization for fixed-type coefficients. Literal admissible arithmetic character families, primary coefficient-prime assignments and same-ambient-datum jets are constructed on the existing full archive host. The finite definite-period inverse preserves the actual theta image. Its distinct integral-kernel inverse is projectively algebraic in the totally definite case and in one irreducible source Hecke sector; the actual mixed source has multiplicity-free sectors with orthogonal target images, giving an exact scalar-weighted sum of algebraic sector inverses and an explicit inverse-norm bound from finite CM evaluations. The integral part constructs the actual horizontal cycle and its balanced Poincar\'e kernels, proves full isotropic canonical-frame cancellation, and identifies the labelled higher-CM graphs with actual horizontal interior special-cycle closures. A relative PEL boundary correspondence recovers the exact chosen compact-ball cycle, retaining the cover degree and projector coefficients. Derived vertical intersections remain separate. Checked implications and cited inputs are distinguished from the unresolved instantiation of the complete character-independent iterated functional. No Masullo claim is used, and no general non-CM Hodge-conjecture conclusion is asserted.
\end{abstract}
\noindent\textbf{Working status.} This is the evolving research paper for the requested run. Checked mathematical statements below have the scope stated in their hypotheses; an unchecked global identification is never used as a theorem. The final status and current gates are recorded in \texttt{QUESTIONS.tsv} and \texttt{LEDGER.md}.

\tableofcontents\clearpage
\part{Scope, sources and the geometric dictionary}
\section{What is being compared}
The requested sketch is a map of proposed connections, not a collection of background theorems. In particular, sketch~1.7 does not specify the test vectors, splittings, inner forms or normalizing constants of its ``one seesaw pairing.'' We therefore distinguish exact identities, exact nonvanishing implications, and a still-uninstantiated identification with the character-independent iterated functional in Alp\"oge's step~12. This is a mathematical distinction: a local Bessel integral, a global period on a definite Shimura set, and a cup-product integral on a surface have different domains.

Throughout, $E/F$ is a CM extension with $F$ totally real. A skew-Hermitian multiplicity space is denoted $W$; writing $W=\delta P$ for a Hermitian space $P$ requires a fixed $\delta\in E^\times$ with $\bar\delta=-\delta$. The associated symplectic form is the trace of the tensor-product form. A change of $\delta$, an additive character, a Weil splitting or a determinant-line trivialization can change a numerical period. These choices are not suppressed.

In the surface construction, $V$ is Hermitian of dimension~3, with signature $(2,1)$ at one real place $\mu$ and $(3,0)$ at the others. There is at least one definite place. Thus $V$ is anisotropic over $F$: an isotropic rational vector would remain isotropic at a definite real place. A sufficiently small level gives a compact projective ball surface. We distinguish $\SU(V)$, used in the original geometric construction, from $\U(V)$, used for adelic theta identities. Passing between their appropriate components and finite covers retains determinant characters and component measures.

\subsection{Background accepted and what was checked}
The supplied 032 theorem for CM abelian varieties, the supplied human proof of van Hoften--Sempliner, and the stated Alp\"oge RH and Dirichlet GRH theorems are carried as cited background, as requested. This run does not purport to referee those global theorems. It reads and checks the particular hypotheses and formulas used. No Masullo claim enters any deduction.

The supplied ZIP contains two Lean entry files and a source map, not the full tree. Named modules were subsequently obtained from a public mirror at commit
\begin{center}\small\nolinkurl{7cbf7a0c32fca1badf05931879a37c3f8a2b03b4}.\end{center}
The mirror's \texttt{Main.lean} and \texttt{Statements.lean} match the supplied files byte for byte. This does not authenticate every mirror module against the original multi-part archive. No Lean compilation was performed. A Lean \texttt{def ... : Prop} is a definition of a proposition; it is not itself a proof of that proposition. Source readings and elementary consequences of those definitions are recorded accordingly.

A mathematical finding here means a deduction formulated and proved for this comparison; it does not mean a claim of research priority. Standard consequences included to prevent an incorrect inference are marked as such. The independent checks are whole-proof readings or reconstructions, not independent formal certifications. The model was GPT-6 Astra Pro; no numerical reasoning-effort setting was exposed. Actual check events, including background uses, appear in the TSV register.

\subsection{Realizations and correspondences}
A nonzero mixed period is a nonzero functional on a Betti realization. It produces an algebraic cycle only after an algebraic morphism or correspondence has been constructed and an extraction criterion has been verified. In the 032 setting rational maps from CM Hodge structures into $H^1(X,\Q)$ come from maps to abelian varieties after multiplying by an integer. Pushing forward $[X]$ then gives the surface cycle. An analytic inverse, an evaluation of a differential at a CM point, or a positive height is not by itself such an algebraic correspondence.

We use motives modulo homological equivalence when discussing projectors whose idempotence is proved on cohomology. We do not silently assert idempotence in the Chow ring. A homologically trivial cycle is numerically trivial, although it can carry nonzero Abel--Jacobi or height data. These distinctions are essential in comparing the ordinary and arithmetic theta correspondences.

\subsection{Grothendieck anchors}
The archival anchors in this paper locate the framework of a finding, rather than claim that its theta formula appears in Grothendieck's text. The original scans of cote~44, scans~7--8, were read: the letter of 12 August 1964 discusses the intermediate Jacobians $J^i$, their duality and the correspondence pairing used to reduce positivity to a polarization question. Cote~16, scans~43--46, the letter to Coates dated 4 January 1967, discusses the standard conjectures, equality of dimensions of complementary algebraic classes, and algebraic inverses associated with hyperplane sections. In particular the date and the first dimension statement are on scan~43, not scan~44. The published \emph{Standard conjectures on algebraic cycles}, pp.~196--198, supplies the complementary realization/polarization language. Later sections give verified SGA/EGA locators for specialization, cycle classes and spreading out.

No separate ``text of record'' transcription of these archival pages was supplied. Consequently a typographical problem in extracted modern-paper text is not entered as a correction to Grothendieck's page. The transcription register records only actual comparisons for which both readings exist.


\section{Oscillator and arithmetic source conventions}

% THETA3 contribution by /root/theta_032, 21 November 2026.
% GPT-6 Astra Pro; inherited reasoning-effort label not separately exposed.
% External theorems are cited inputs. A1--A5 deductions were checked in full.

\subsection{The exact theta dictionary and an applicable Rallis formula}

The attached 032 paper fixes a Galois CM field $E$, its real subfield
$F=E^+$ with $[F:\Q]\geq2$, and a Hermitian $E$-space $V$ of rank three,
of signature $(2,1)$ at the distinguished embedding $\mu$ and $(3,0)$
elsewhere. Its group is $\SU(V)$, and its domain $D$ parametrizes negative
lines $L\subset V_\mu$. Put $Q=V_\mu/L$; thus
\[
 T_D=\Hom(L,Q),\qquad \Omega_D^1=L\otimes Q^*.
\]
For a multiplicity space $P$ of rank $k$, the active positive Lagrangian
has pieces $L\otimes P_\mu$ and $Q^*\otimes P_\mu^*$. The half-determinant
line is $L^k$ times a constant line. These are the precise identifications
in 032, Section 4.1 and Lemma 4.1, printed pp.13--16. The character signs
are $d(\tau)=v(\tau^{-1})$, not $v(\tau)$. The finite Schwartz input, the
rational symplectic space, the embedding of $\SU(V)$, and $V$ itself remain
fixed in Proposition 4.6. Only multiplicity subspaces at compact places
move. This operation does not change the signature of $V$.

The first Taylor tensor projects to a holomorphic one-form. Proposition
4.5, pp.21--22, proves that finitely many actual finite inputs span any
specified cotangent fibre; the common level is chosen after selecting
those inputs. For $P=P_1\perp P_2$, double alternation of the completed
second Taylor tensor gives a two-form. Proposition 4.6, pp.22--24, proves
the product identity. In its divided-power convention the degree-two
coefficient of a product is
\[
 T_2^{(1)}T_0^{(2)}+T_1^{(1)}T_1^{(2)}+T_0^{(1)}T_2^{(2)}.
\]
Double alternation kills the two unmixed terms and sends the mixed term
to the wedge of the two first tensors with coefficient one. The general
completed second tensor need not be holomorphic. Intermediate compact
parameters are asserted only to give smooth forms; the rational
decomposition endpoints are holomorphic. Theorem 5.1, p.26, proved in
Section 6.3, and Corollary 6.7, p.45, then extract the rational CM Hodge
maps on a common principal cover. These arithmetic extraction results
are distinct from the analytic product identity.

The companion paper of van Hoften--Sempliner makes the global character
data explicit. In Sections 3.1 and 4.1, the $\mu_i$ are unitary
conjugate-symplectic Hecke characters of weight one, with restriction
$\epsilon_{E/F}$; their CM types satisfy
$\Phi_{\mu_i}=\Phi_i^{-1}$. The pair consists of $\U(V)$ and a
skew-Hermitian line $W_i$, with splitting characters $(\mu_i,\nu_5)$.
The character $\pi_i$ on $[\U(W_i)]$ is a source automorphic character,
and is not the splitting character $\mu_i$. Section 4.2 imposes the
global equality $\mu_1\mu_2=\mu_3\mu_4$. Lemma 4.2.1 supplies
$W_1\perp W_2\simeq W_3\perp W_4$. Doubling gives a split rank-four
skew-Hermitian space and the model $\mathcal S(V(\A_E)^2)$.
Remark 4.2.5 explains the cancellation of the diagonal $\U(V)$ character.
With that cancellation its action is $\Phi(x)\mapsto\Phi(g^{-1}x)$.
The doubling map has normalization
\[
 \delta(f\otimes\bar f)(0)=\|f\|_2^2.
\]
Proposition 4.2.6 and Section 4.2.8 express the fourfold cup period as a
specified geometric scalar times a doubled theta integral. Lemmas
4.2.7, 4.2.9, and 4.2.10 provide, respectively, nonvanishing, the
archimedean zero evaluation, and finite torus-character extraction.

For the rank-one to rank-three lift, an applicable explicit Rallis
formula is Haan, \emph{The Bessel period of $\U(3)$ and $\U(2)$ involving
a non-tempered representation}, arXiv:1403.5061v7, Section 3.3,
Theorem 3.4, printed p.8. Its proof cites J.-S. Li, \emph{Non-vanishing
theorems for the cohomology of certain arithmetic quotients}, Theorem
2.1, and the doubling factorization. This is a convergent-range
formula; no appeal to a second-term-range theorem with incompatible
dimensions is needed.

Here is its normalization. Let $V_1$ be Hermitian of rank one and $W_3$
skew-Hermitian of rank three, and put $T=\U(V_1)$ and $G=\U(W_3)$.
Haan uses the symplectic form
$\frac12\Tr_{E/F}(h_{V_1}\otimes h_{W_3})$, a fixed nontrivial additive
character $\psi$, and a unitary character $\gamma$ of
$E^\times\backslash\A_E^\times$ with
$\gamma|_{\A_F^\times}=\epsilon_{E/F}$. The splitting pair is
$(\gamma,\gamma^3)$. The theta lift is
\[
 \theta(\bar f,\varphi)(g)
 =\int_{[T]}\theta(t,g;\varphi)\overline{f(t)}\,dt,
\]
with Tamagawa measure. If $\tau$ is a unitary automorphic character of
$T$, $f_i=\bigotimes_v f_{i,v}\in\tau$, and
$\varphi_i=\bigotimes_v\varphi_{i,v}$, then
\begin{equation}\label{A-Rallis}
 \left\langle\theta(\bar f_1,\varphi_1),
                    \theta(\bar f_2,\varphi_2)\right\rangle_G
 =\frac{L_E(3/2,\operatorname{BC}(\tau)\otimes\gamma^3)}
        {L_F(3,\epsilon_{E/F})}
       \prod_v Z_v^{\#}(1,f_{1,v},f_{2,v},\Phi_{1,v}).
\end{equation}
The target Petersson measure is also Tamagawa. Local pairings are
chosen so that their product is the global Petersson pairing on
$\tau$, and local Haar measures multiply to the relevant global
Tamagawa measure. The section $\Phi_{s,v}$ is the holomorphic
Siegel--Weil section obtained from
$\rho_{1,3}(\varphi_{1,v}\otimes\bar\varphi_{2,v})$. With these
conventions,
\[
\begin{split}
 Z_v(1)&=\int_{T(F_v)}\Phi_{1,v}(\iota(t,1))
               \langle\tau_v(t)f_{1,v},f_{2,v}\rangle_v\,dt,\\
 Z_v^{\#}(1)&=
 \frac{L_v(3,\epsilon_{E_v/F_v})}
 {L_{E_v}(3/2,\operatorname{BC}(\tau_v)\otimes\gamma_v^3)}Z_v(1),\\
 \Phi_{1,v}(\iota(t,1))&=
       \langle\omega_v(t)\varphi_{1,v},\varphi_{2,v}\rangle_v.
\end{split}
\]
The normalized local factors equal one for normalized unramified data
outside a finite set. These conventions, including the last matrix
coefficient identity, are in Haan Sections 2.1, 2.4, and 3.1--3.2,
pp.3--8. The $L$-values at $3/2$ and $3$ are nonzero for the unitary
characters occurring here. Thus choosing nonzero local factors gives a
nonzero global lift. For the anisotropic rank-three group in the present
application, square integrability and cuspidality cause no additional
global convergence issue. Local occurrence and local test vectors still
must be checked; the formula does not assert that every prescribed
Schwartz vector has nonzero norm.

This displayed formula fixes one splitting convention completely.
To transport it to arbitrary splitting pairs, first compute the actual
unitary characters $\lambda(g)$ and $\kappa(t)$ in an identity
$\omega'(g,t)=\lambda(g)\kappa(t)\omega(g,t)$. Then the new lift with
input $\bar\tau$ is the old lift with input
$\overline{\tau\kappa^{-1}}$, multiplied by $\lambda$; its norm is
unchanged by the latter unitary twist. This gives the correct character
substitution without silently choosing a conjugation convention.
The factor $1/2$ in Haan's symplectic form must likewise be matched by
the form or additive-character normalization. Replacing both quotient
measures by probability measures multiplies the displayed target norm
by $\operatorname{vol}([G])^{-1}\operatorname{vol}([T])^{-2}$ when the
underlying input function is kept fixed. One must also account for
passage from $\U(V)$ and all adelic components to a connected
$\SU(V)$ quotient. No such factor is implicitly set equal to one here.



\subsection{The arithmetic dictionary in 032, Sections 5--6}

The geometric target of the degree-one theta tensor is the Albanese
realization. The arithmetic extraction uses numerical characters of a
commuting algebra of algebraic Hecke endomorphisms; these are not the
global splitting characters $\mu_i$ or the torus inputs $\pi_i$.
Here are the exact formulas, with the conventions that determine them.

At a suitable split good prime with residue cardinality $q$, let
$\Phi$ denote the $q$-power Frobenius endomorphism of the reduced
Albanese. In 032, Proposition~5.5, p.~32, the integral endomorphisms
$D_1,D_2,D_3$ satisfy
\[
 \Phi^3-D_1\Phi^2+qD_2\Phi-q^3D_3=0.
\]
The Frobenius on crystalline degree-one cohomology is its pullback
$\varphi=\Phi^*$, with the same polynomial in $D_i^*$.
Replacing this by reciprocal Frobenius would invert the slopes and
would not be this formula. The valuation is normalized by
$\nu_\iota(q)=1$ in Section~5.2, p.~30.

Fix the finite theta input away from the active split place $w$.
Lemma~6.2, pp.~37--38, constructs $r\in E^\times$, $M>0$, and
$\omega\in\overline\Q^\times$ such that
\[
 r\bar r=1,\quad \operatorname{val}_w(r)=M,
 \quad\operatorname{val}_{\bar w}(r)=-M,
 \quad D_q^M=\omega,
 \qquad
 \omega^2=\zeta\prod_{\tau\in\Sigma}\tau(r)^{-3d(\tau)}.
\]
Here $D_q$ is the normalized radial dilation, the relation holds on
the image in theta classes, and $\zeta$ is a root of unity. The
square records the metaplectic ambiguity. The exponent $-3$ is
the active rank, not the dimension of the multiplicity line.

Proposition~6.3, pp.~39--40, gives the half-modulus projective
spherical parameter
\[
 \{q^{1/2},q^{-1/2},e\},\qquad
 e=e_0^{-1},\quad e_0^M=\omega,
 \qquad
 \nu_\iota(\sigma e)=-\frac32v(\sigma|_E).
\]
Every complex conjugate of $e$ has modulus one. A scalar-conjugate
active character has third parameter $e'=\pm\sigma e$.
For a full numerical character $\chi$, write its full parameter as
$b(q^{1/2},q^{-1/2},e')$. The central relation in Lemma~6.4
implies that $b^3e'$ is a root of unity. In equation~(6.9), p.~42,
\[
 \chi(D_1)=q e_1(\alpha),\qquad
 \chi(D_2)=q e_2(\alpha),\qquad
 \chi(D_3)=e_3(\alpha),
 \quad\alpha=b(q^{1/2},q^{-1/2},e').
\]
The $e_i$ are elementary symmetric polynomials. These are numerical
eigenvalues; the Hecke operators on their generalized spaces can
still have nilpotent parts.

Proposition~6.6, pp.~43--44, combines this polynomial, Weil weight,
filtered comparison, and weak admissibility. It gives the pure
Hodge index
\[
 t_\eta=\frac12-\frac13\nu_\iota(e')\in\{0,1\}.
\]
The generalized space is holomorphic for $t_\eta=1$ and
antiholomorphic for $t_\eta=0$. The proof uses polynomial
idempotents in extending algebraic Hecke endomorphisms, so the
projectors preserve the Hodge filtration. For each scalar
conjugate the filtration ranks are calculated over the actual
compositum of the definition field and the projector coefficient
field. Thus it does not assume independence of those fields or
semisimplicity of all ambient Hecke operators.

Theorem~5.1, p.~26, proved in Section~6.3, pp.~44--45, then applies
scalar descent to the whole selected generalized space. Formula
(6.13) reads
\[
 t_{\sigma\eta}=\frac{1+v(\sigma|_E)}2.
\]
It supplies the exact rational-Hodge maps from $U(v)$ into
$H^1(X,\Q)$; Corollary~6.7, p.~45, places a finite collection
on a common principal cover. Only then does the Albanese convert
the maps into algebraic maps to the prescribed CM abelian varieties.
This completes the original rank-three dictionary. A6 uses Liu's
all-rank component theorem to obtain the corresponding identification
in higher ball dimensions, without claiming that the displayed
cubic Frobenius calculation itself has already been generalized.

The relation to the definite machine is therefore at two levels:
the oscillator kernels and splitting rules are shared, whereas the
selected numerical Hecke eigenspaces, their integral projectors,
and the definite class-function normalization have to be matched.
Q9 makes the projector denominators explicit when their integral
matrices and annihilating polynomials have been supplied.


\part{Checked findings and proved conditional implications}

\section{Common kernels, ball tensors and exact pairing recovery}

\subsection{THETA3-A1: the anisotropic Siegel zero-orbit identity}\label{A1}

\paragraph{Statement.}
Let $E/F$ be CM and let $V$ be an anisotropic Hermitian $E$-space.
Let $r\geq1$, and let $W^{\square}$ be split skew-Hermitian of rank
$2r$, with chosen complementary maximal isotropic subspaces. Use its
Schr\"odinger model $\mathcal S(V(\A_E)^r)$ and compatible splittings
for which
\[
 (\omega(g)\Phi)(x)=\Phi(g^{-1}x),\qquad
 (\omega(n(b))\Phi)(x)=\psi(\langle b,Q_V(x)\rangle)\Phi(x).
\]
Here $N$ is the Siegel unipotent, identified with the additive
$F$-space of Hermitian $r$ by $r$ matrices, $Q_V(x)$ is the Gram
matrix, and the pairing is nondegenerate. A conventional factor
$1/2$ may be absorbed in this identification. Give $[G]$ and $[N]$
probability Haar measures, where $G=\U(V)$, and define
\[
 T_\Phi(h)=\int_{[G]}\sum_{x\in V(E)^r}(\omega(g,h)\Phi)(x)\,dg.
\]
Then
\begin{equation}\label{A1-eq}
                    \int_{[N]}T_\Phi(n)\,dn=\Phi(0).
\end{equation}

\paragraph{Proof.}
Anisotropy makes $[G]$ compact; the centre is an anisotropic norm-one
torus as well. The quotient $[N]$ of an adelic vector group is compact.
On compact sets of $g$, the finite support condition confines rational
points to finitely many lattice cosets, and uniform archimedean
Schwartz decay dominates their polynomial counting function. Thus the
theta series converges absolutely and uniformly there. The $N$-action
adds only phases of modulus one. Summation and both integrals may
therefore be interchanged.

Since $Q_V(g^{-1}x)=Q_V(x)$, the inner unipotent integral is
\[
 \int_{[N]}\psi(\langle b,Q_V(x)\rangle)\,db
 =\begin{cases}1&Q_V(x)=0,\\0&Q_V(x)\ne0.\end{cases}
\]
For a nontrivial character the vanishing follows by translation by an
element on which its value differs from one. Nondegeneracy of the
pairing ensures that the character is trivial precisely for zero Gram
matrix. If the Gram matrix vanishes, every $x_i$ has zero norm, so
anisotropy forces $x=0$. The surviving term is fixed by $G$ and is
$\Phi(0)$, proving the formula. $\square$

The identity applies to both the ball Hermitian space and a totally
definite Hermitian space. It identifies the normalized zero-orbit
functional, not their nonzero Fourier coefficients, unaveraged periods,
or motives. A residual diagonal character changes its right-hand side
to $\Phi(0)\int_{[G]}\chi(g)\,dg$; isotropic $V$ instead admits
nonzero zero-Gram tuples. These hypotheses cannot be omitted.

\subsection{THETA3-A2: finite translation and quantitative nonvanishing}\label{A2}

\paragraph{Statement.}
In A1 write $\Phi=\Phi_\infty\otimes\Phi_f$, assume
$a=\Phi_\infty(0)\ne0$ and $\Phi_f(0)=1$, and keep the infinite input
fixed. For every $0<\varepsilon<1$ there is $h_f\in N(\A_{F,f})$
such that
\[
 \left|T_{\Phi_\infty\otimes\omega_f(h_f)\Phi_f}(1)\right|
                   \geq(1-\varepsilon)|a|.
\]
The finite Schwartz norm is unchanged. Every compact subgroup of
$G(\A_{F,f})$ fixing the total doubled input still fixes its translate.

\paragraph{Proof.}
The normalized mean in A1 is $a$. Compactness and continuity give
$n=(n_\infty,n_f)$ with $|T_\Phi(n)|\geq|a|$. Additive weak
approximation gives $\gamma_j\in N(F)$ with
$\gamma_{j,\infty}\longrightarrow n_\infty$. Automorphy and
continuity imply
\[
 T_\Phi(\gamma_{j,\infty},n_f)
 =T_\Phi(1,\gamma_{j,f}^{-1}n_f)
 \longrightarrow T_\Phi(n).
\]
Choose $j$ making the error at most $\varepsilon|a|$. Commutation of
the dual-pair actions identifies the value with the translated-input
integral. The finite Weil action is unitary. The same commutation shows
that every old $G$-stabilizer fixes the translated total input.
$\square$

This assertion changes finite theta data. Its finite tensor factors
may require a smaller common level after decomposition. It supplies
no uniform bound for their number or coefficients. The explicit
additive grid and Gaussian derivative certificate in THETA3-Q7
strengthen this existence argument under named effective data.

\subsection{THETA3-A3: the exact finite extraction cost}\label{A3}

\paragraph{Statement.}
Suppose an inverse-doubling decomposition and the finite torus
expansions of van Hoften--Sempliner Lemma 4.2.10 give
\[
 P=\sum_{a=1}^{R}c_a
       \sum_{\chi_i\in\mathcal X_{i,a},\ 1\leq i\leq4}
                       P_{a,\chi_1,\chi_2,\chi_3,\chi_4}.
\]
The inner terms are actual periods of integrated character lifts, with
conjugation on the last two factors and one fixed geometric
normalization. If $P\ne0$, at least one satisfies
\[
 |P_{a,\chi_1,\chi_2,\chi_3,\chi_4}|
 \geq \frac{|P|}{\sum_a|c_a|\prod_{i=1}^4|\mathcal X_{i,a}|}.
\]

\paragraph{Proof.}
Let $M$ be the maximum modulus of the finitely many periods in the
sum. The triangle inequality gives
$|P|\leq M\sum_a|c_a|\prod_i|\mathcal X_{i,a}|$; divide by the
positive denominator. $\square$

For an effective certificate use finite permissible character sets or
computable upper bounds for their sizes; determining the exact nonzero
support may itself require zero testing. Further rational Hodge-map or
connected-component extraction has its own finite coefficient cost.
The estimate is the same finite-sum principle developed in THETA3-Q5.

\subsection{THETA3-A4: higher-rank ball tensors and a surface section}\label{A4}

\paragraph{Statement.}
Replace the rank-three $V$ in the analytic part of 032 by a Hermitian
space of rank $m\geq3$, of signature $(m-1,1)$ at one real place and
positive definite at the others. Write $d=m-1$. For a multiplicity line
the first projected Taylor tensor again takes values in $\Omega_D^1$;
for a multiplicity plane, double alternation of the second tensor takes
values in $\Omega_D^2$. Finite actual inputs span every cotangent
fibre, and decomposable inputs satisfy the same wedge identity as in
032. Balanced compact multiplicity signs therefore give a nonzero
mixed pairing
\begin{equation}\label{A4-pair}
 \int_X\xi\wedge\bar\eta\wedge c_1(A)^{d-2}\ne0
\end{equation}
on a sufficiently deep compact principal ball quotient $X$, where
$A$ is ample and $\xi,\eta$ are the rational-decomposition endpoint
two-forms. A smooth complete-intersection surface $S\subset X$ of
divisors in $|A^{n_j}|$, $1\leq j\leq d-2$, satisfies
\[
 \int_S\xi|_S\wedge\overline{\eta|_S}
 =\left(\prod_jn_j\right)
      \int_X\xi\wedge\bar\eta\wedge c_1(A)^{d-2}\ne0.
\]
For $d=2$ take $S=X$ and the empty product one.

\paragraph{Proof.}
Let $L$ be the negative tautological line and $Q=V_\mu/L$, of rank
$m-1$. For multiplicity rank $k$ the active Lagrangian is
$L\otimes P\oplus Q^*\otimes P^*$. The determinant identity is
\[
 \det\mathcal F_P
 =L^{2k}\otimes(\det V_\mu)^{-k}\otimes(\det P)^{2-m}.
\]
The last two factors, and all compact-place factors, are constant on
$D$. On $\SU(V)$ the first of those constant representations is
trivial. After choosing the same constant-line trivializations as in
032, $(\det\mathcal F_P)^{1/2}=L^k$ times a constant line. Thus
\[
 (\det\mathcal F_{P_1})^{1/2}\otimes Q^*\otimes P_1^*
       \simeq L\otimes Q^*=\Omega_D^1
\]
for a line, and
\[
 (\det\mathcal F_P)^{1/2}\otimes
             \bigwedge^2Q^*\otimes\bigwedge^2P^*
       \simeq L^2\otimes\bigwedge^2Q^*=\Omega_D^2
\]
for a plane. The latter is the canonical bundle only when $m=3$.

The Gaussian-isolation argument of 032 Proposition 4.5 uses a position
space of dimension $m$ instead of three. In a fixed graph frame,
indicators of $x_0+k\widehat{\OO}_F^{\,m}$ give first tensors tending
to a nonzero Gaussian scalar times the Minkowski vector $x_0$.
Uniform Gaussian decay proves the convergence. These lattice vectors
span the complexified position space; the image of the span of actual
inputs is a finite-dimensional, hence closed, subspace. Projection
onto $L\otimes Q^*$ is surjective. Consequently actual theta values
span every cotangent fibre, and a finite selected list has a common
principal level. The covariance argument uses the same unitary Weil
splittings and triviality of determinant twists on $\SU(V)$, and is
independent of the value three.

For a multiplicity decomposition the completed theta series multiply.
The degree-two coefficient and double-alternation calculation stated
above is valid for any rank of $Q$. Select two first tensors with
independent values at one point; their wedge $\xi$ is nonzero.
Compactness of $X$ follows from a definite real place, and projectivity
from the arithmetic ball-quotient construction. Choose a positive
curvature representative $\omega$ of $c_1(A)$. Pointwise
$\xi\wedge\bar\xi\wedge\omega^{d-2}$ is a positive constant
times $\|\xi\|^2\omega^d$, so its integral is positive.

Keep the rational total multiplicity plane and its finite input fixed.
On a compact parameter neighborhood, the same Gaussian domination
with polynomial derivatives proves $C^\infty$ continuity of the
projected second tensor, exactly as in 032 Proposition 4.6(iii).
Therefore its mixed pairing with $\xi$, using $\omega^{d-2}$, remains
nonzero near the original decomposition. Balanced signs permit a
rational orthogonal decomposition with the requested endpoint signs
arbitrarily close to this compact-place datum, by the same weak
approximation argument. Only multiplicity subspaces move; $V$ is
unchanged. At that rational endpoint the tensor is a finite sum of
wedges of holomorphic first tensors, so is holomorphic and closed.
Pass to a common principal cover for the selected finite families.
The integral scales by the positive covering degree. This proves
\eqref{A4-pair}.

For sufficiently large $n_j$, Bertini gives a smooth projective
complete-intersection surface. Its cohomology class is
$(\prod_jn_j)c_1(A)^{d-2}$. Both endpoint forms are closed because
holomorphic forms on a projective manifold are closed. The
cohomological restriction formula gives the displayed surface
identity. This also shows directly that every such smooth complete
intersection retains a nonzero holomorphic two-form when $\eta=\xi$.
$\square$

This extends the analytic and geometric wedge mechanism. It does not
by itself extend 032's precise $U(v)$ arithmetic extraction theorem
with its rank-three Frobenius calculation. The exact all-rank source
identification is supplied separately in A6 using Liu's componentwise
Theorems 4.15 and 4.18 and the reflex-type calculation. Liu's
Corollary 4.20 also gives the CM Albanese conclusion in
all ranks $m\geq3$, so simply raising this ball rank does not create
non-CM $H^1$. Nor does adding positive directions turn this family
into a hyperbolic Witt tower or a chain of totally definite theta
lifts. A theorem imposing signatures at every real place still
requires those signatures to match: one active ball and positive
definite compact places do not satisfy such a hypothesis merely by
increasing $m$.

\subsection{THETA3-A5: an exact finite sum of definite Bessel periods}\label{A5}

\paragraph{Setup.}
Return to rank three and let $G=\U(V)$, with the ball signature at one
real place and positive definite signature elsewhere. Choose
$z\in V(E)$ negative at the active place, using weak approximation;
then $H=\U(z^\perp)$ is totally definite of rank two. View $H$ as a
subgroup of $G$ fixing $z$. Let $K_f\subset G(\A_{F,f})$ be compact
open and let $K_\infty$ be the stabilizer of the negative line $Ez$ at
the active place, times the compact groups at the other real places.
Fix a unitary scalar character $\kappa$ of $K_\infty$.

Let $\mathcal E$ be a nonzero finite-dimensional space of smooth
automorphic functions on $[G]$, right $K_f$-invariant, satisfying
$f(gk)=\kappa(k)^{-1}f(g)$ for $k\in K_\infty$. Give it the
Petersson inner product inherited from probability measure on $[G]$.
The application is a fixed-level span of holomorphic canonical theta
forms in their equivariant adelic realization. On a ball surface the
canonical fibre is one-dimensional, which gives precisely a scalar
$K_\infty$ type. Fix the splitting and central type and use the full
adelic realization, as in van Hoften--Sempliner, before extracting
connected components. This is not an assertion that an arbitrary
unrealized form on a connected $\SU(V)$ quotient has already been
extended to $G$.

For $a\in G(\A_{F,f})$ put
\[
 R_af(h)=f(ha),\qquad
 K_H(a)=H(\A_{F,f})\cap aK_fa^{-1}.
\]
Give $[H]$ probability measure, and let $\mathcal M_a$ be the space of
automorphic functions on $H$ right $K_H(a)$-invariant and of right
$H(F\otimes_\Q\R)$-type $\kappa^{-1}|_H$.

\paragraph{Statement.}
Writing $D=\dim_\C\mathcal E$, there are $s\leq D$ finite adelic
elements $a_1,\ldots,a_s$ for which
\[
 \mathcal Q=\sum_{j=1}^sR_{a_j}^*R_{a_j}
\]
is positive definite on $\mathcal E$. For all $u,v\in\mathcal E$,
\begin{equation}\label{A5-frame}
 \langle u,v\rangle_G
  =\sum_{j=1}^s
          \langle R_{a_j}\mathcal Q^{-1}u,R_{a_j}v\rangle_H.
\end{equation}
Choose an orthonormal basis $e_{j,\alpha}$ of $\mathcal M_{a_j}$ and
put
\[
 \mathcal P_{j,\alpha}(f)=
    \int_{[H]}f(ha_j)\overline{e_{j,\alpha}(h)}\,dh.
\]
With inner products conjugate-linear in the first variable, the
identity becomes the finite sum
\begin{equation}\label{A5-periods}
 \langle u,v\rangle_G
  =\sum_{j,\alpha}
       \overline{\mathcal P_{j,\alpha}(\mathcal Q^{-1}u)}
                   \mathcal P_{j,\alpha}(v).
\end{equation}
These are genuine global adjacent-unitary-group periods of $\U(3)$
against totally definite $\U(2)$. Moreover
\begin{equation}\label{A5-bounds}
 \lambda_{\min}(\mathcal Q)\|u\|_G^2
 \leq\sum_{j,\alpha}|\mathcal P_{j,\alpha}(u)|^2
 \leq\lambda_{\max}(\mathcal Q)\|u\|_G^2.
\end{equation}

\paragraph{Proof: well-defined finite restriction spaces.}
The left $H(F)$-invariance is inherited from $G(F)$-invariance. If
$k\in K_H(a)$, then $a^{-1}ka\in K_f$, so
$f(hka)=f(ha)$. At infinity, finite $a$ commutes with $H(F\otimes\R)$,
which is contained in $K_\infty$; this proves the asserted right type.
Thus $R_a$ maps into $\mathcal M_a$.

Because $H$ is totally definite,
$H(F)\backslash H(\A_{F,f})/K_H(a)$ is finite. After choosing finite
representatives, a scalar-type function is determined by its values
at those representatives. A finite stabilizer can force a value to
vanish when its character is nontrivial, but introduces no infinite
dimensional space. Thus $\mathcal M_a$ is finite-dimensional. Its
norm is a finite weighted sum of squared values, with positive masses
proportional to the reciprocal orders of the finite stabilizers. In
particular, a smooth nonzero restriction has positive norm.

\paragraph{Proof: separation and at most $D$ translates.}
The group $G$ satisfies weak approximation at the archimedean places
needed here. A direct proof uses the Cayley transform
\[
 X\longmapsto(1-X)(1+X)^{-1},\qquad X^*=-X.
\]
The skew-adjoint matrices form an $F$-vector space; its $F$-points
are dense in the product of its real completions. The Cayley transform
is an isomorphism of open sets, with inverse
$(1-g)(1+g)^{-1}$. Its real image, excluding the proper locus with
eigenvalue $-1$, is dense. Hence $G(F)$ is dense in
$G(F\otimes_\Q\R)$. This uses $\U(V)$, not an unproved
approximation assertion for a chosen $\SU(V)$ presentation.

If $f\in\mathcal E$ is nonzero, choose $(g_\infty,g_f)$ where its
value is nonzero. Continuity and weak approximation provide
$\gamma\in G(F)$ such that $f(\gamma_\infty,g_f)\ne0$.
Automorphy gives
\[
 f(\gamma_\infty,g_f)=f(1,\gamma_f^{-1}g_f).
\]
For $a=\gamma_f^{-1}g_f$, therefore $R_af(1)\ne0$, and
$R_af$ has positive norm. The family of restrictions separates
$\mathcal E$.

Start with its full kernel space $\mathcal E_0=\mathcal E$.
If $\mathcal E_i\ne0$, choose nonzero $f_i\in\mathcal E_i$ and
one restriction not annihilating it, and intersect $\mathcal E_i$
with that restriction's kernel. Dimension drops strictly. After at
most $D$ steps the common kernel is zero. Thus
$\langle u,\mathcal Q u\rangle_G=\sum_j\|R_{a_j}u\|_H^2$
is positive for every nonzero $u$.

\paragraph{Proof: identity, Parseval, and bounds.}
The operator $\mathcal Q$ is positive and self-adjoint, hence
invertible. By its definition,
\[
 \sum_j\langle R_{a_j}\mathcal Q^{-1}u,R_{a_j}v\rangle_H
 =\langle\mathcal Q^{-1}u,\mathcal Q v\rangle_G
 =\langle u,v\rangle_G.
\]
Finite-dimensional Parseval in each $\mathcal M_{a_j}$ gives
\eqref{A5-periods}. The spectral bounds for a positive self-adjoint
operator give \eqref{A5-bounds}. $\square$

\paragraph{Exact matrices and conditional quantitative bounds.}
In a basis of $\mathcal E$, let $G_{\rm ball}$ be its Petersson Gram
matrix. Let $R_j$ be the finite evaluation matrix and $D_j$ the
positive diagonal mass matrix for $\mathcal M_{a_j}$. Then
\[
 H_{\rm fin}=\sum_jR_j^*D_jR_j,\qquad
                  \mathcal Q=G_{\rm ball}^{-1}H_{\rm fin}.
\]
This formula keeps the analytic Gram matrix separate from definite
finite masses. If $D>1$, $G_{\rm ball}\leq M I$, and a square
$D$ by $D$ minor $B$ of the stacked matrix with blocks
$D_j^{1/2}R_j$ has determinant of modulus at least $\delta>0$, then
Cauchy--Binet and the eigenvalue product give the valid lower bound
\[
 \lambda_{\min}(\mathcal Q)
 \geq\frac{\delta^2}{M\bigl(\operatorname{tr}H_{\rm fin}\bigr)^{D-1}}.
\]
The case $D=1$ is immediate. A height or denominator argument can
supply $\delta$ only after algebraicity of appropriately normalized
entries has actually been proved. No such algebraicity, and no
rationality of $\mathcal Q^{-1}$, follows from finite dimensionality.

\paragraph{Meaning of the comparison.}
Formula \eqref{A5-periods} is an exact finite-sum refinement of the
desired comparison with definite periods, with an explicitly defined
frame correction. A fixed restriction alone computes its own norm,
not the ball norm. The evaluation of canonical forms at these CM
points is a fibrewise operation, not a pullback of a degree-two
cohomology class to a zero-dimensional variety. The analytic operator
$\mathcal Q^{-1}$ is not automatically an algebraic correspondence.
An algebraic, rational, or uniform-integral realization is an
additional theorem required for a motivic transfer.

\subsection{Direct sums, genuine seesaws, and iterated lifts}

For compatible splittings and $W=W_1\perp W_2$, the common oscillator
representation gives
\[
 \theta_W(g,h_1\oplus h_2;\varphi_1\otimes\varphi_2)
 =\theta_{W_1}(g,h_1;\varphi_1)
             \theta_{W_2}(g,h_2;\varphi_2).
\]
The differential version is the double-alternating wedge identity.
It belongs to the actual seesaw whose pairs are
$\U(V)\times\U(V)$ against $\U(W_1)\times\U(W_2)$, and diagonal
$\U(V)$ against $\U(W)$. Integrating against characters or test
automorphic forms gives a seesaw identity by Fubini on the compact
quotients in this application. A second orthogonal decomposition of
the same $W$ supplies two other tori in the same group. It does not
change $V$ to a definite Hermitian space.

This construction is bilinear in two independent inputs and combines
two kernels by a direct sum. An iterated lift is instead a composition
of linear integral operators, with the output automorphic
representation of one step serving as the input to the next. One
cannot identify the two merely from their number of theta factors.
The literal step 8 of Alp\"oge's note, p.4, describes a chain of
totally definite groups and changes of automorphic parameters. Its
recovered mirror contains actual successive operators \texttt{up}
and \texttt{dn}, rather than the above wedge operation. A comparison
therefore needs intertwining maps and matching representations,
splittings, and vectors. The finite-frame identity A5 and the
restriction seesaw for $V=z^\perp\perp Ez$ are concrete alternative
comparisons; neither asserts that the two independent degree-one
lifts already form that iterated chain.


% THETA3-A6 contribution by /root/theta_032, 21 November 2026.
% GPT-6 Astra Pro; inherited reasoning-effort label not separately exposed.
% Liu/BMM/CM classification are cited inputs, not independently re-proved.

\subsection{THETA3-A6: a direct construction for balanced tensors of arbitrary even length}\label{A6}

\paragraph{Statement and scope.}
Let $n\geq2$, let $E$ be a Galois CM field with real subfield $F$
and $[F:\Q]\geq2$, and let $c\in\operatorname{Gal}(E/\Q)$ be complex
conjugation. Suppose that the odd sign functions
$v_1,\ldots,v_{2n}:\operatorname{Gal}(E/\Q)\to\{1,-1\}$ satisfy
\begin{equation}\label{A6-balance}
 v_i(1)=1,\qquad \sum_{i=1}^n v_i=\sum_{i=n+1}^{2n}v_i.
\end{equation}
Put $\lambda=(1,\ldots,1,c,\ldots,c)$, with $n$ entries of each kind.
The theta construction on a compact ball quotient of dimension $n$
gives a direct cycle construction for the line
\begin{equation}\label{A6-line}
 \bigotimes_{i=1}^{2n}U(v_i)_{\lambda_i}
 \ \subset\ H^{2n}\!\left(\prod_{i=1}^{2n}B(v_i),\overline\Q\right).
\end{equation}
The arithmetic input is Liu's exact oscillator-component theorem in
arbitrary rank, together with the reflex-type argument of
van Hoften--Sempliner. The proof below does not infer the prescribed
$U(v_i)$ sources from CM-ness of the Albanese alone. It does not use
the reduction to four-factor relations.

All auxiliary abelian varieties in this construction remain CM. The algebraicity
conclusion is already covered by the supplied CM Hodge-conjecture
results; the contribution here is the direct $2n$-factor construction,
its exact multiplicity and differential degree, and the comparison
with the stated theta inputs. No claim of research priority is made.

\paragraph{The exact arithmetic source in rank $m\geq3$.}
Let $V$ be Hermitian of rank $m$, with signature $(m-1,1)$ at the
distinguished real place and $(m,0)$ elsewhere. It is anisotropic
over $E$, because there is a definite real place. Replacing its active
real component by a positive definite component while retaining the
finite components gives a totally definite incoherent Hermitian
space. The original $V$ is its active nearby coherent space.
Consequently Liu's setup applies to this precise Shimura tower:
Definitions C.3--C.4 and Proposition C.5, printed pp.108--109, and
Section 4.2, pp.45--47. The compact case holds because $[F:\Q]>1$.
There is no all-place indefinite-signature hypothesis in this use
of Liu's theorem.

For an odd $v$ with $v(1)=1$, write
$\Phi(v)=\{g:v(g)=1\}$ and $d(\tau)=v(\tau^{-1})$. Choose a
unitary conjugate-symplectic Hecke character $\mu$ of weight one with
\[
 \mu|_{\A_F^\times}=\epsilon_{E/F},\qquad
 \Phi_\mu=\Phi(v)^{-1}.
\]
Existence is van Hoften--Sempliner, Lemma 3.1.5, pp.9--10.
Liu's admissibility condition is the existence of a skew element
$e\in E^\times$, $e^c=-e$, such that
$\operatorname{Im}\tau(e)<0$ for $\tau\in\Phi_\mu$, with the
specified finite norm classes; see Definition 4.12, p.47. This is
exactly the 032 sign prescription: for $e=\delta h_P$, it says
\[
 \operatorname{sgn}(h_{P,\tau})
   =-d(\tau)\operatorname{sgn}(\operatorname{Im}\tau(\delta)).
\]
The skew-Hermitian source line therefore has active signature
$(0,1)$.

Liu, Definition 4.3, p.41, puts
$\mu_{\rm alg}=\mu|\cdot|_E^{-1/2}$ and defines a coefficient
field $M_\mu$ containing the reflex field $M'_\mu$ of
$(E,\Phi_\mu)$, with reflex type $\Psi_\mu$.
Definition 4.5, p.42, specifies a CM abelian variety $A_\mu$ and its
distinguished $M_\mu$-inclusion eigenline, whose Galois character is
$\mu_{\rm alg}$; Proposition 4.6, pp.42--43, supplies such data.
Theorem 4.15, pp.50--51, identifies the Galois character of the
particular admissible oscillator component with this exact character
for every rank $m\geq3$. Theorem 4.18, pp.52--54, then gives
\begin{equation}\label{A6-Liu}
 \Hom_E(A_\infty,A_\mu)_\Q\otimes_{M_\mu}\C
       \simeq\bigoplus_{\varepsilon,\chi}
                         \omega(\mu,\varepsilon,\chi).
\end{equation}
The sum uses $\mu$-admissible $\varepsilon$. The map in its proof
is pullback of a basis of the distinguished inclusion eigenline;
its finite-level assertion identifies $K$-invariants with
$\Hom_E(A_K,A_\mu)_\Q$. Here $\chi$ is the finite central
automorphic character of Liu, Definition 4.11. It is neither the
splitting character $\mu$ nor the source torus character.
Corollary 4.20, pp.54--55, describes the whole CM Albanese, but the
stronger componentwise statement \eqref{A6-Liu} is what is needed.

Let $(E_0,\Phi_0)$ be the primitive CM type inducing
$(E,\Phi(v))$. Then $B(v)\sim B_0^{[E:E_0]}$. Since $E$ is
Galois and $\Phi_\mu=\Phi(v)^{-1}$, the reflex pair
$(M'_\mu,\Psi_\mu)$ is $(E_0,\Phi_0)$, and
$A_{\mu,\C}\sim B_0^{[M_\mu:E_0]}$. These are precisely the
CM-type computations in van Hoften--Sempliner, Corollary 3.1.13,
p.12, and Proposition 4.1.8, pp.15--16. Their proof with the
primitive factor and the selected inclusion line gives
\begin{equation}\label{A6-extraction}
 \omega(\mu,\varepsilon,\chi)
 \subseteq
 \operatorname{span}_{\C}\{h(U(v)_1):
       h\in\Hom_{\Hdg}(U(v),H^1(X,\Q))\},
\end{equation}
on a sufficiently deep common cover. This argument is independent
of $m$. The inverse in $\Phi(v)^{-1}$ is essential.

\paragraph{The first theta vector and the scalar conventions.}
The local comparison needed to use \eqref{A6-extraction} is explicit.
In a graph frame at a base point, 032's first Taylor tensor is the
ordinary gradient: a nonzero scalar times a linear position function
multiplied by the positive Gaussian. The Bargmann transform sends
linear position times the vacuum to a nonzero scalar times the
corresponding linear Fock coordinate. After projection to the active
cotangent summand, this is the contraction
\[
 \psi_{1,0}=\sum_{\alpha=1}^{m-1}
                 z''_\alpha\otimes\xi'_{\alpha,m}.
\]
It is the $q=1$ case of BMM, published Lemma 5.5, pp.43--44;
equation (5.11), p.54, defines its Schwartz realization by the
inverse Bargmann transform. One may check the comparison on this
fibre and transport it by $\SU(V)$, since both sides have the same
covariance. The Gaussian creation-operator formulas in BMM,
Appendix C, Lemmas C.2--C.3, pp.115--116, also verify directly that
the coefficient is nonzero.

The numerical scalar depends on the Bargmann, Gaussian and
constant-line normalizations. It need not be algebraic: only a
complex span and nonvanishing are used in this argument. It must be
retained if a numerical period identity is later asserted.
BMM, Lemma 5.18, p.49, fixes the twist
$\det^{(a-b)/2}$ required for the special cochain on $\U(V)$.
For the holomorphic degree-one cochain $a=0,b=1$, it is
$\det^{-1/2}$ in their oscillator convention. The reciprocal
half-determinant in a covariant Taylor-value convention is a dual
representation convention, not a change that may be silently made
on one side of an identity. On $\SU(V)$ the determinant twists
are trivial. As in 032, Lemma 4.4, pp.20--21, the circle and double
lifts agree there: their ratio is a continuous character of a
connected semisimple real group and hence is trivial.

At the compact places the appropriately split Gaussian is the
trivial representation; the exact possibilities are Liu,
Lemma D.2(1), pp.127--128. At the active place the holomorphic
component is $\omega_{m-1,1}^{-1,-,0}$ by Lemma D.2(2).
Fix a source splitting $\nu_m$ with
$\nu_m|_{\A_F^\times}=\epsilon_{E/F}^m$. Changing $\nu_m$
twists the source $\U(1)$ character, denoted by $\pi$; it does not
change the indicated $\mu$-component after that twist. In the
van Hoften--Sempliner convention the lift integrates
$\theta(g,t)\pi(t)$. Writing it instead as
$\theta(g,t)\eta(t)^{-1}$ requires $\eta=\pi^{-1}$.

At fixed finite level and prescribed infinity type the source
torus quotient has only finitely many character components.
Thus the raw first theta gradient is a finite sum of the corresponding
integrated lifts, all of whose nonzero components satisfy
\eqref{A6-extraction}. The degree-one argument of
van Hoften--Sempliner, Proposition 4.1.8, applies in every rank:
BMM, Section 9.1 and Theorem 9.1, pp.72--75, require
$3(a+b)+|a-b|<2m$, which here is $4<2m$. This theorem is being
used in degree one, not for the subsequent top-degree wedge.

For a finite list of forms, construct the full $\U(V)$ lifts with
these splittings and restrict to a connected component. A sufficiently
deep principal $\SU(V)$ level is contained in their common finite
stabilizer. Equivalently, at a lattice-preserving unitary level the
determinants of rational arithmetic elements are norm-one algebraic
units, hence roots of unity; a deep principal level kills them.
This gives the required component cover and its rational Hodge
pullbacks. No extension of an arbitrary unrealized special-unitary
form is being assumed.

There is also a split-local check on the rank parameter. Liu,
Lemma D.1, proof p.126, gives, for
$\mu=\nu\boxtimes\nu^{-1}$,
\[
 \omega(\mu,\varepsilon,\chi)
 =\operatorname{Ind}_{\mathrm{GL}_{m-1}\times\mathrm{GL}_1}
                       ^{\mathrm{GL}_m}
       ((\nu\circ\det)\boxtimes\chi\nu^{1-m}),
\]
with normalized induction. Its projective unramified Satake tuple is
\[
 \{q^{(m-2)/2},q^{(m-4)/2},\ldots,q^{-(m-2)/2},
                   \chi(\varpi)\nu(\varpi)^{-m}\}.
\]
Thus the exponent is $-m$, agreeing with $-3$ in 032. An inverse
geometric Hecke convention inverts the tuple. This is a consistency
check; Theorems 4.15 and 4.18 supply the actual CM source.

\paragraph{Degree $n$ alternation on the $n$-ball.}
Now set $m=n+1$. Choose multiplicity lines $P_1,\ldots,P_n$ with
the first $n$ prescribed signs, and put $P=\perp_{i=1}^nP_i$.
At the active place all line signs agree, and
\[
 \mathcal F_P=(L\otimes P)\oplus(Q^*\otimes P^*),
 \qquad \operatorname{rank}Q=n.
\]
The determinant computation in A4 gives
\[
 \det\mathcal F_P=L^{2n}(\det V_\mu)^{-n}(\det P)^{2-m}
\]
times compact constant factors. Its square root is $L^n$ times a
constant line. Define
\[
 \operatorname{Alt}_n:\operatorname{Sym}^n(Q^*\otimes P^*)
       \longrightarrow\bigwedge^nQ^*\otimes\bigwedge^nP^*
\]
by the product of the two exterior wedges. Swapping two factors
changes both wedge signs, so it factors through the symmetric power.
After the fixed constant-line trivializations, its target with the
half-determinant is
$L^n\otimes\bigwedge^nQ^*=\Omega_D^n$.

For a decomposable finite input the completed theta series is the
product of its $n$ line theta series. Its degree-$n$ coefficient,
in the divided-power convention, is the sum indexed by
$k_1+\cdots+k_n=n$. A summand with some $k_i\geq2$ contains
two factors in the same multiplicity line and is killed by
$\bigwedge^nP^*$. The sole surviving index is $(1,\ldots,1)$.
Consequently the projected degree-$n$ tensor is exactly the wedge
of the $n$ first tensors, with coefficient one.

For an arbitrary finite input the same assertion holds as a finite
sum of such wedges. Indeed a compactly supported locally constant
function on a finite-adelic direct sum is a finite sum of rectangle
indicators after choosing a sufficiently fine open lattice, so its
Schwartz space is the algebraic tensor product of the factor spaces.
A change to a rational adapted polarization uses the rational
metaplectic covariance of 032, Lemma 4.2 and Remark 4.3, in the
corresponding dimension. Its nonzero lift scalar is retained in the
constant-line normalization; it is not discarded outside $\SU(V)$.

Here is the spanning argument in the new dimension. Fix a point
$x_*\in D$ and, for one prescribed multiplicity line, choose rational
Schr\"odinger positions $F^m$. Their Minkowski image is a real vector
space of dimension $D_0=m[F:\Q]$, with full lattice
$\Lambda=\OO_F^m$. In the graph frame of its period point
$Z=(Z_\tau)_\tau$, put
\[
 P_i:\C^{D_0}\longrightarrow
       (L\otimes Q^*)_{x_*}.
\]
This is the complex-linear projection of the first Taylor fibre
onto its $c\mu$ summand, followed by the fixed nonzero
half-determinant and multiplicity-line trivializations. It is
surjective. The complex span of $\Lambda$ is $\C^{D_0}$, so choose
$x_{i,1},\ldots,x_{i,n}\in\Lambda$ such that
$P_i x_{i,1},\ldots,P_i x_{i,n}$ is a basis of the cotangent fibre.
For each $x_0=x_{i,j}$ use the actual finite input
$\phi_{x_0,k}=1_{x_0+k\widehat{\OO}_F^{\,m}}$. Its rational
support is exactly $x_0+k\OO_F^m$, and its projected first tensor is
\[
 2\pi i\sum_{x\in x_0+k\Lambda}
       P_i x\,
       \exp\left(\pi i\sum_\tau x_\tau^tZ_\tau x_\tau\right).
\]
Choose $\lambda_0>0$ below the least eigenvalue of every
$\operatorname{Im}Z_\tau$, and let $\rho>0$ be a lower bound for
the norm of a nonzero lattice vector. For
$k\geq2\|x_0\|/\rho$ and $0\ne y\in\Lambda$,
$\|x_0+ky\|\geq k\|y\|/2$. The tail is therefore bounded by
\[
 2\|P_i\|\sqrt{\frac{\pi}{e\lambda_0}}
   \sum_{r\geq1}(2r+3)^{D_0}
      \exp\left(-\frac{\pi\lambda_0 k^2\rho^2 r^2}{8}\right).
\]
Indeed $u e^{-\pi\lambda_0u^2/2}\leq
(\pi e\lambda_0)^{-1/2}$ absorbs the linear factor, and disjoint
balls of radius $\rho/2$ give at most $(2r+3)^{D_0}$ lattice
points in the shell $r\rho\leq\|y\|<(r+1)\rho$.
The displayed series tends to zero by dominated convergence;
THETA3-Q2 also gives a finite certified tail bound. Thus the
projected tensors tend to
\[
 2\pi i\,P_i x_{i,j}\,
 \exp\left(\pi i\sum_\tau
              x_{i,j,\tau}^tZ_\tau x_{i,j,\tau}\right).
\]
Each exponential is nonzero. The determinant of these $n$ limit
columns is nonzero, so continuity of the determinant gives $n$
actual finite inputs whose values form a basis. This construction
works separately for every prescribed type, with its own positive
Gaussian and projection $P_i$.

Now choose one input from each of the first $n$ prescribed families
in order. The already chosen values span a proper subspace until
the last step, while the next family's selected finite basis spans
the full fibre. One of its values lies outside that subspace.
The resulting $n$ values are independent. Their first theta forms
are holomorphic because the first completed tensor is the ordinary
holomorphic gradient and the active embedding projection is
holomorphic. Their wedge $\xi$ is nonzero. Only after this finite
selection do we choose a common principal level, using smoothness
of the finite Weil action. On the resulting smooth projective
compact ball quotient $X$,
\begin{equation}\label{A6-positive}
 i^{n^2}\int_X\xi\wedge\bar\xi>0.
\end{equation}
The phase is necessary when $n$ is odd: in complex coordinates,
$i^{n^2}dz_1\wedge\cdots\wedge dz_n\wedge
d\bar z_1\wedge\cdots\wedge d\bar z_n=2^n\,d\operatorname{vol}$.

\paragraph{A nearby rational orthogonal decomposition.}
Keep $V$, $P$, its rational symplectic realization, the finite total
input, and the total special-unitary level fixed. At a compact real
place, \eqref{A6-balance} says that the first and last lists contain
the same number of each sign. Relabel an initial real orthogonal
basis to have the desired last-list labels. This leaves its aggregate
positive and negative multiplicity subspaces unchanged. At the
active place every line has the same sign, so the total active
positive Lagrangian is unchanged as well.

Approximate these ordered orthogonal bases at all real places by
vectors in $P(E)$, using weak approximation. Apply Hermitian
Gram--Schmidt. The denominators are nonzero norms; being nonzero
with the required signs is an open condition. The resulting
rational orthogonal basis can therefore be chosen arbitrarily close
and with all the prescribed signs. It gives
$P=P_{n+1}\perp\cdots\perp P_{2n}$ over $E$.

Positive Gaussian majorants, with polynomial factors for derivatives,
give $C^\infty$ continuity of the projected total tensor on a compact
parameter neighborhood, exactly as in 032, Proposition 4.6(iii),
pp.22--24. Choose its constant-line trivializations continuously,
agreeing with the original ones at the reference point. The mixed
pairing with $\xi$ remains nonzero near that point by
\eqref{A6-positive}. Intermediate parameters need only give smooth
forms. At the chosen rational endpoint, finite tensor decomposition
gives a finite sum of wedges of holomorphic first tensors of the last
$n$ types. Select a summand whose mixed period is nonzero. After
a common principal cover for the selected inputs, one has
\begin{equation}\label{A6-period}
 \int_X\theta_1\wedge\cdots\wedge\theta_n
       \wedge\bar\theta_{n+1}\wedge\cdots\wedge\bar\theta_{2n}
        \ne0.
\end{equation}
The cover multiplies the integral by its positive degree.

The full unitary splitting characters may also be matched globally.
Choose $\mu_1,\ldots,\mu_{2n-1}$ with their prescribed types and set
\[
 \mu_{2n}=\frac{\prod_{i=1}^n\mu_i}
                    {\prod_{i=n+1}^{2n-1}\mu_i}.
\]
Its restriction to $\A_F^\times$ is
$\epsilon_{E/F}^{n-(n-1)}=\epsilon_{E/F}$, and
\eqref{A6-balance} gives its required weight-one infinity type.
Thus $\prod_{i\leq n}\mu_i=\prod_{i>n}\mu_i$ globally, with the
same $\nu_m$ retained on the source side. The special-unitary
argument has already removed these determinant twists, but they
are available for a compatible full-unitary tensor rule.

\paragraph{The direct cycle extraction, including its sign.}
By \eqref{A6-extraction}, expand the finitely many classes in
\eqref{A6-period} in complex spans of rational Hodge images.
Conjugate the last $n$ expansions. Nonvanishing selects rational
Hodge maps
$h_i:U(v_i)\to H^1(X,\Q)$ and nonzero
$\alpha_i\in U(v_i)_{\lambda_i}$ with algebraic Betti
coefficients, whose integrated product is nonzero. The arbitrary
analytic scalars are coefficients of the expansion; each selected
eigenvector itself is chosen over $\overline\Q$.

By the Albanese and the rational Riemann theorem, realize suitable
positive integer multiples of the $h_i$ by morphisms
$f_i:X\to B(v_i)$, as in 032, Lemma 2.2, p.5. Put
$A=\prod_iB(v_i)$, $g=\dim A$, $f=(f_i)$, and
$Z=\operatorname{cl}(f_*[X])\in H^{2g-2n}(A,\Q)$.
Writing $t_\alpha=\bigotimes_i\alpha_i$ in the distinct
$H^1$ slots, the projection formula gives
\[
 \int_A Z\smile t_\alpha=k\in\overline\Q^\times.
\]
For each $i$ take a copied factor $B(v_i)'$ and a nonzero
algebraic-Betti vector $\beta_i$ in its line labelled by
$c\lambda_i$. The pure class
$D_i=\alpha_i\otimes\beta_i\in H^2(B_i\times B_i',\overline\Q)$
is algebraic. Indeed every scalar conjugate puts its two factors
in opposite Hodge types, by centrality of $c$; scalar descent and
the Lefschetz $(1,1)$ theorem give its exact class as an algebraic
linear combination of divisor classes. No projection of an
unspecified divisor is involved.

Let $A'=\prod_i B_i'$, let $p,p'$ be the two projections from
$A\times A'$, and let $q_i$ select $B_i\times B_i'$. With the old
factors ordered before the new ones,
\[
 K:=\prod_iq_i^*D_i=(-1)^n t_\alpha\otimes t_\beta,
 \qquad t_\beta=\bigotimes_i\beta_i.
\]
The number of odd-degree transpositions is
$(2n)(2n-1)/2=n(2n-1)$, of parity $n$. Hence
\begin{equation}\label{A6-cycle}
 p'_*(p^*Z\smile K)=(-1)^n k\,t_\beta.
\end{equation}
Its codimension is $(g-n)+2n-g=n$. All operations on the left
preserve algebraicity, and $k$ is nonzero algebraic. Divide by $k$
and apply scalar complex conjugation, which preserves the
$\overline\Q$ span of rational cycle classes and changes the labels
back from $(c^n,1^n)$ to $(1^n,c^n)$. Identifying the copied factors
proves algebraicity of \eqref{A6-line}. This is a cycle on the line,
with the codimension and coefficient given by \eqref{A6-cycle}.
\hfill$\square$

\paragraph{The Grothendieck framework.}
This is an explicit composition of divisor correspondences,
pullback and proper pushforward in homological motives. The
relevant framework locators in the supplied Grothendieck scans
are \emph{cote 44}, scans 7--8, for intermediate Jacobians and
polarizations through correspondences, and \emph{cote 16},
scans 43--46, for the dimensions of Lefschetz operators and
algebraic inverse correspondences. Here the actual correspondence
$K$ is displayed and its output is computed in
\eqref{A6-cycle}. No general algebraic inverse Lefschetz
operator or Chow projector is assumed. Its action on the
specified tensor line follows directly from the divisor classes
and the projection formula.

\paragraph{An explicit family not generated by divisor products.}
For every $n\geq2$, choose a Galois CM field with at least $n+1$
real pairs of embeddings; a sufficiently large cyclotomic field
does so. Choose compact-place representatives
$\tau_1,\ldots,\tau_n$. Let $v_0$ be $+1$ on a chosen set of
representatives, including the distinguished one, and extend it
oddly. For $1\leq i\leq n$, let $v_i$ differ from $v_0$ only at
$\tau_i$ and its conjugate. Let $v_{n+1}$ differ at all those
$n$ pairs and put $v_{n+2}=\cdots=v_{2n}=v_0$. Then
\[
 \sum_{i=1}^n v_i=v_{n+1}+(n-1)v_0.
\]
No first-list function equals a second-list function.

The selected line \eqref{A6-line} is not in the subalgebra generated
by divisor classes. To prove this using the full Galois criterion,
project a product of $n$ divisor classes to that pure
one-$H^1$-per-factor line. A potentially nonzero contribution must
pair the $2n$ factors. At the distinguished embedding, each pair
must join a first-list factor to a last-list factor. For its selected
two eigenlines to have type $(1,1)$ under every scalar conjugation,
the two complete type functions must be equal. There are no such
pairs. Every divisor product therefore has zero projection to the
line. This argument uses the fixed eigenline labels, not a claim
that the underlying abelian varieties are pairwise nonisogenous.
Galois translates or isogenies involving other labels do not change
the conclusion. Formula \eqref{A6-cycle} nevertheless constructs
the line by a codimension-$n$ algebraic cycle, for all $n$, including
$n\geq4$.

\paragraph{Transfer to actual Weil pieces.}
Let a CM abelian variety $A_0$ of dimension $2n$ have multiplication
by an imaginary quadratic field $K$ with signature $(n,n)$.
Each of its two Weil lines over $\overline\Q$ is a wedge of
$2n$ CM embedding lines and has exactly $n$ of each Hodge type
under every scalar conjugation. Choose a Galois CM field containing
the CM fields involved, and let $s_i(g)$ be the complete type
functions of those lines. Reorder by their type at the distinguished
embedding and put $v_i=s_i$ for the first $n$, $v_i=-s_i$ for the
last $n$. Then \eqref{A6-balance} holds, and the selected original
line has the same complete type as $U(v_i)_{\lambda_i}$.

Extend a map between these matching lines by zero on the other
source embedding lines. Every scalar conjugate is a Hodge map,
so scalar descent expresses it as a $\overline\Q$ linear
combination of rational Hodge maps; this is 032, Lemma 2.4, p.8.
Lemma 2.2 realizes those maps by pullbacks of abelian morphisms.
Apply their tensor product to \eqref{A6-line} and take cup product.
After expansion this is a linear combination of pullbacks along
product morphisms from $A_0$ to the auxiliary product, and hence
preserves algebraicity. It proves the prescribed Weil line
algebraic, without an assumed algebraic projector.

For an explicit family in every $n$, choose a totally real Galois
field $F$ of degree $2n$, put $E=F(i)$, and select a CM type
containing $n$ embeddings above each embedding of $\Q(i)$.
The CM abelian variety has dimension $2n$ and the required Weil
signature. One field construction takes a prime
$\ell\equiv1\pmod{4n}$, the unique degree-$2n$ subfield of
$\Q(\zeta_\ell)$, which is real, and then adjoins $i$.
This does not assert simplicity or that every such Weil class is
outside the divisor algebra.

The same argument applies to the embedding-line monomials of any
rational CM Hodge class: rationality makes every occurring monomial
balanced under all scalar conjugations. Thus the direct
$2n$-factor construction can replace the four-factor combinatorial
reduction, with Liu's all-rank theorem retained as a deep cited input.

The remaining requirements for the proposed iterated Bessel bridge
are different. They include an actual tuple satisfying the recovered
\texttt{LawsX} and link conditions, identification of these
cohomological vectors and their normalizers, and the relevant
integral theta frames. Increasing $m$ does not satisfy the
every-place signature hypotheses of a different tower theorem,
and an independent $n$-fold product is not thereby an iteration
of theta operators.


% THETA3-A7 contribution by /root/theta_032, 21 November 2026.
% GPT-6 Astra Pro; inherited reasoning-effort label not separately exposed.

\subsection{THETA3-A7: an independent wedge is a genuine single theta lift}\label{A7}

\paragraph{Setup and statement.}
Let $V$ have rank $n+1$, signature $(n,1)$ at one real place and
positive definite signature elsewhere, with $[F:\Q]\geq2$.
Let $P=P_1\perp\cdots\perp P_n$ be the skew-Hermitian multiplicity
space of the first $n$ types in A6. Its active component is definite.
Put
\[
 G=\U(V),\qquad S=\U(P),\qquad T=\prod_{i=1}^n\U(P_i).
\]
The quotient $[S]=S(F)\backslash S(\A_F)$ is compact: a nonzero
global isotropic vector would remain isotropic in the definite
active component. This assertion does not require $S$ to be compact
at every real place.

Fix compatible theta splittings, a finite Schwartz input, and the
degree-$n$ holomorphic Fock cocycle. At the active place this is
the determinant coefficient of the wedge of $n$ degree-one
cocycles; at the other places take the specified Gaussian
coefficients. Let
\[
 \mathcal K_h=\theta_V(\,\cdot\,,h;\Phi)\in\mathcal E,\qquad
 \mathcal E=H^0(X,\Omega_X^n),
\]
in the equivariant adelic realization on a fixed sufficiently small
compact Shimura level, including its finitely many components.
The notation includes the indicated differential-form projection.
Fix a product of unitary source characters
$\pi(t)=\prod_i\pi_i(t_i)$ on $[T]$, and use the
van Hoften--Sempliner convention
\[
 F=\int_{[T]}\mathcal K_t\,\pi(t)\,dt.
\]
For a decomposable input the tensor rule and degree-$n$ alternation
identify $F$ with the wedge of the corresponding $n$ integrated
degree-one lifts. For a finite sum of inputs it is the corresponding
finite sum of wedges.

There is a smooth automorphic function $v_\pi$ on $[S]$, of finite
level and finite maximal-compact type, such that the exact identity
\begin{equation}\label{A7-lift}
 F=\int_{[S]}\theta_V(\,\cdot\,,h;\Phi)\,
                                  \overline{v_\pi(h)}\,dh
\end{equation}
holds with the displayed measures. Thus the independent wedge is
an actual single theta lift from $\U(P)$. The representative
$v_\pi$ is allowed to depend on the source characters and on the
kernel. No character independence is asserted.

\paragraph{Why the kernel has finite rank.}
First we verify that $\mathcal K_h$ belongs to the same
finite-dimensional holomorphic space for every $h\in S(\A_F)$.
At infinity the source Weil action commutes with the target
$G$-action, its maximal compact action, and its relative Lie
algebra differential. Applying the source action to a closed
cochain of Hodge type $(n,0)$ therefore gives another closed
cochain of that type. The corresponding differential form is
$\bar\partial$-closed and hence holomorphic. The original cochain
is the wedge of the degree-one holomorphic cocycles; its
closedness follows from the derivation rule for the cochain
differential. The explicit source is BMM, Lemma 5.4 and
Proposition 5.17(1),(4), pp.43 and 48. No stable-range assertion
about generation in degree $n$ is used.

At finite places the source and target Weil actions commute.
Every element of a compact open target group fixing the original
finite input consequently fixes its source translate. The target
level is thus independent of $h$. The chosen target splitting
and canonical compact type are retained. Compactness of $X$
implies that $\mathcal E$ is finite-dimensional.

The theta series and its derivatives converge uniformly on compact
sets of representatives, by Schwartz decay and the positive
Gaussian estimates. This gives a smooth $\mathcal E$-valued
function of $h$. It is left $S(F)$-invariant by rational theta
automorphy. A sufficiently small finite source subgroup fixes
$\Phi_f$, giving a common finite level. At infinity a polynomial
times a Gaussian is finite under a source maximal compact
subgroup, so the resulting coefficient functions have finite
maximal-compact type. These statements use compact representative
sets for the anisotropic adelic quotients, not a finite-dimensional
assertion about all functions with one compact type on a
noncompact real group.

Choose a Petersson orthonormal basis
$f_1,\ldots,f_D$ of $\mathcal E$, with inner products conjugate
linear in the first variable. Then
\begin{equation}\label{A7-kernel}
 \mathcal K_h=\sum_{j=1}^D f_j\,e_j(h),
 \qquad e_j(h)=\langle f_j,\mathcal K_h\rangle_G .
\end{equation}
The functions $e_j$ are smooth automorphic functions on $[S]$
with the level and compact-type properties just proved.
Let $\mathcal C=\operatorname{span}_{\C}\{e_1,\ldots,e_D\}$.
It is finite-dimensional. Its $L^2([S])$ inner product is positive
definite: a nonzero smooth function is nonzero on an open set
of positive Haar measure. In particular the passage from smooth
functions to their $L^2$ classes is injective here.

\paragraph{The exact Riesz representative.}
Define a complex-linear functional on $\mathcal C$ by
\[
 \Lambda_\pi(e)=\int_{[T]}e(t)\pi(t)\,dt.
\]
The inclusion $T\hookrightarrow S$ gives a well-defined continuous
map $[T]\to[S]$. The torus quotient is compact, so this integral
exists for every $e\in\mathcal C$. Its boundedness as a functional
on $\mathcal C$ follows from finite dimensionality; no assertion
about boundedness of pointwise restriction on an unrestricted
$L^2$ space is needed.

By finite-dimensional Riesz representation there is a unique
$v_\pi\in\mathcal C$ such that
\[
 \Lambda_\pi(e)=\langle v_\pi,e\rangle_S
               =\int_{[S]}\overline{v_\pi(h)}e(h)\,dh
 \quad(e\in\mathcal C).
\]
Substitute \eqref{A7-kernel} and interchange the finite sum and
the integrals:
\[
 F=\sum_jf_j\Lambda_\pi(e_j)
   =\int_{[S]}\left(\sum_jf_je_j(h)\right)
                                  \overline{v_\pi(h)}\,dh.
\]
This is \eqref{A7-lift}. If $F\ne0$, then
$\Lambda_\pi$ is nonzero on $\mathcal C$, so $v_\pi\ne0$.
For an orthonormal basis $\eta_1,\ldots,\eta_r$ of $\mathcal C$,
the representative and its norm are explicitly
\[
 v_\pi=\sum_{a=1}^r
              \overline{\Lambda_\pi(\eta_a)}\,\eta_a,
 \qquad
 \|v_\pi\|_S^2=\sum_{a=1}^r|\Lambda_\pi(\eta_a)|^2.
\]
These are analytic finite-dimensional quantities; their values
or normalizations have not been proved algebraic.
\hfill$\square$

The assertion is an equality of theta integrals with a smooth
automorphic input. It does not select a prescribed irreducible
source representation. When $P$ is totally definite, the entire
archimedean source group is compact; the stated finite source
level and finite archimedean type then place $v_\pi$ in a
finite-dimensional automorphic module, which has a finite
spectral decomposition. For the more general anisotropic source,
finite rank of the particular coefficient space $\mathcal C$
is what proves the identity.

\paragraph{The signature condition for a totally definite middle source.}
At a compact real place $\tau$, the sign of the Hermitian form on
$P_i$ is
$-d_i(\tau)\operatorname{sgn}(\operatorname{Im}\tau(\delta))$.
Consequently $P$ is definite there if and only if all
$d_i(\tau)$ agree. It is already definite at the active place.
Thus the actual multiplicity space is totally definite if and
only if
\[
 v_1=v_2=\cdots=v_n
\]
as complete type functions. For $n=2$, a place where $v_1$
and $v_2$ differ gives signature $(1,1)$ on $P$. Multiplying
the form by a scalar can reverse the signs but cannot make that
component definite.

There are always nonzero wedges from one repeated type: by A6's
fibre-spanning proof, one theta family supplies $n$ forms with
independent values. Such a wedge has a totally definite middle
source in \eqref{A7-lift}. But for the particular balanced tensor
of A6, equality of the first $n$ types forces equality of all
$2n$ types. Indeed at each $g$, the first sum is either $n$ or
$-n$; a sum of $n$ numbers in $\{1,-1\}$ attains that value
only when each summand has the same sign. Hence every last-list
$v_i$ equals the common first-list function. The selected
balanced tensor line is then a product of divisor lines.

This separates an arbitrary nonzero canonical wedge from the
particular wedge needed to reach a selected non-Lefschetz
Weil component. Changing to isogenous CM factors does not
remove this condition: the signs follow the complete types of
the selected eigenlines, including their embedding labels.

In the direct $m=2$ instantiation of the recovered
\texttt{MainX}, the middle group is required to be totally
definite. Formula \eqref{A7-lift} therefore provides the actual
middle source only in the just-described definite case.
For a general exceptional balanced input, it supplies a genuine
single lift from an anisotropic group with compact quotient,
but not the required totally definite middle space. It does
not rule out another source realization of the same canonical
automorphic representation. A second reverse lift, its chosen
vector, and the complete \texttt{LawsX} and arithmetic link data
remain separate requirements for that iterated comparison.


\subsection{THETA3-A8: the completed Taylor tensor is the canonical harmonic jet}\label{A8}

This comparison concerns a fixed polarized abelian variety and its
rigidified theta line bundle.  When the polarized torus and its CM
action in 032 are realized by an archive datum, it identifies the
actual analytic section and its derivatives.  It does not identify
an arbitrary representation in a theta packet with that section.
The distinction between a divided Taylor coefficient and a
differential operator is retained throughout.

\paragraph{The completed theta function is the Appell--Humbert reading.}
First use principal symplectic coordinates.  Write $Z=X+iY$ with
$Y>0$, put $\Lambda=\Z^g+Z\Z^g$, and set
\[
 H(z,w)=z^tY^{-1}\overline w,\qquad
 \widehat\vartheta(Z,z)
       =\exp\left(\frac\pi2z^tY^{-1}z\right)\vartheta(Z,z).
 \tag{A8a}
\]
For $\lambda=a+Zb$ with $a,b\in\Z^g$, the usual theta multiplier
and a direct calculation give
\[
 \widehat\vartheta(Z,z+\lambda)
  =(-1)^{a^tb}
    \exp\left(\pi H(z,\lambda)+\frac\pi2H(\lambda,\lambda)\right)
       \widehat\vartheta(Z,z).
 \tag{A8b}
\]
Indeed, the ordinary multiplier is
$\exp(-\pi i b^tZb-2\pi i b^tz)$.  The new linear term is
$\pi z^t(Y^{-1}\lambda-2ib)=\pi z^tY^{-1}\overline\lambda$.
For the constant term, use
$\lambda-\overline\lambda=2iYb$ to obtain
\[
 \frac\pi2\lambda^tY^{-1}\lambda-\pi i b^tZb
  =\frac\pi2\lambda^tY^{-1}\overline\lambda+\pi i a^tb.
\]
This proves (A8b), including its sign.

For a finite Schwartz input, choose a common lattice level $N$ so
that the positions in its support lie in $N^{-1}\Z^g$ and its
coefficients are invariant under $N\Z^g$.  The same reindexing
argument gives (A8b) on $N\Lambda$.  Equivalently,
$w\mapsto\widehat\vartheta_\varphi(Z,Nw)$ is a section of the
$N^2$-th power of the line bundle on $\C^g/\Lambda$ with the
displayed multiplier.  A nonprincipal polarization is handled in
an integral symplectic lattice with its actual Riemann form; the
same calculation uses its corresponding Hermitian form and
semicharacter.  In the restriction-of-scalars coordinates of 032,
the exponent and Hermitian form are the sums over real embeddings.

The archive's \texttt{AnFr.IsTh}, in
\path{Definitions/ThetaValues/ThConst.lean}, is exactly the law
\[
 \theta(z+\lambda)
   =\chi(\lambda)^k
       \exp\left(k\pi H(z,\lambda)
                +\frac{k\pi}{2}H(\lambda,\lambda)\right)\theta(z).
\]
Thus (A8a), with the stated level scaling, is its analytic reading.
Here the polarized lattice, semicharacter, and rigidification must
be matched, and the two readings must represent the same actual
section.  For that section, two nowhere-vanishing trivialization
identifications with the same multiplier have a holomorphic
periodic frame ratio.  This ratio descends to a compact complex
torus, hence is constant, and the rigidification fixes that
constant.  This does not claim that the multiplier determines an
arbitrary section of a higher tensor power.  The $E$-adjointness of
$H$ is the Rosati adjointness of
the polarization and supplies \texttt{AnPt.adj}.  The tangent
frame supplies \texttt{AnPt.der}.

There is a conjugation in the embedding labels that must be kept.
The torus with period $Z$ is naturally
$A=\mathcal F_P^\vee/\iota(\Lambda_{\mathrm{rat}})$, with
$\iota(u)(v)=\psi_P(v,u)$ for $v\in\mathcal F_P$.
In the graph frame this is $\iota(p,q)=Zq-p$, which gives the
displayed period lattice.  Rosati adjointness gives
$\iota(au)(v)=\iota(u)(\overline a v)$.
Consequently the tangent variable dual to the tensor block
$(\mathcal F_P)_\tau$ has the archive tangent label $c\tau$.
The $c\mu$ Taylor-tensor block of 032 therefore corresponds to
the $\mu$ tangent variables in this convention.

\paragraph{What the archive's canonical derivative means.}
The exact definitions are as follows.  In
\path{Formalism/LocalToGlobal/ProjectiveLattices/IsThRule.lean},
\texttt{NabV} assigns to a framed point $y$, a section $F$ of
$L^k$, and a monomial $\alpha$ in the $E$-eigenvectors of the
tangent space a scalar $\operatorname{nab}_y^\alpha(F)$.
\texttt{ThGerm.lean}, \texttt{ChartAt} and \texttt{IsNab}, require
that, for $k\geq0$, these scalars be computed from a formal-group
chart, a compatible biextension frame, a normalized primitive,
and a canonical gauge.  They are not arbitrary ordinary
derivatives in an unspecified local generator.

The precise correction is in
\texttt{ExpTiePf.lean}, theorem \texttt{tie\_holds}: if $f$ is the
section's formal coefficient in such a chart, $\phi$ is the formal
exponential with the prescribed tangent frame, and $h_L$ is the
chart's canonical potential, then the corrected exponential germ
is
\[
       g=\bigl(f\exp(kh_L)\bigr)\circ\phi,
 \qquad
       \operatorname{nab}_y^\alpha(F)=\alpha!\,[z^\alpha]g.
 \tag{A8c}
\]
This is the actual sign and factorial convention.  The definition
\texttt{expVal} in \texttt{IsExpChart.lean} includes $\alpha!$.
The analytic chart constructed in
\nolinkurl{ExistsAnalyticBiextensionChartComplexPoint.lean} has
additive formal group and $h=0$, and its germ is the Taylor series
of the Appell--Humbert reading.  The two proof files
\nolinkurl{HarmonicValueCanonicalConnectionEqPolyDeriv.lean} and
\nolinkurl{ThetaValueAnalyticFrameEqPolyDerivFock.lean} use exactly
these data, rather than identifying a raw theta series with an
integral local coefficient.

For completeness, the residual ambiguity has a precise and much
smaller form.  The normalized exponential-frame comparison
\texttt{GaugeUniqSt}, proved in \texttt{GaugeUniqPf.lean}, gives
\[
       g'_b=g_b\exp(-kq),\qquad
 q(z)=\sum_{\tau,i,j}q_{\tau ij}z_{\tau i}z_{c\tau,j},
 \tag{A8d}
\]
where one representative of each conjugate pair is used.
In particular $q$ is homogeneous of degree two; there are no
linear or higher terms.  These restrictions come from the
rigidification and the CM adjointness cocycle in
\texttt{IsExpGauge}, not from a choice to disregard other terms.
The formal exponential is unique for its tangent frame.
For the chart used by \texttt{IsNab}, compatibility under jet
truncation follows from existence of germs for every section of
every nonnegative tensor power: this is
\nolinkurl{GlobalTheory.NablaPinCompare.compat_of_germs} in
\nolinkurl{ExistsIsNabNotIsNablaPin.lean}, using properness and
ampleness.  It is therefore legitimate to apply the exponential
and gauge comparison to the actual \texttt{IsNab} data.  All the
archive statements in this paragraph are source inputs read at
the pinned mirror commit; no Lean build or transitive dependency
certificate is claimed here.

\paragraph{A stronger harmonic comparison, with proof.}
Let $z_{\tau i b}$ denote the coordinate in tangent embedding
$\tau$, row $i$, and section column $b$.  Let
\[
 \Delta_{\tau ij}
     =\sum_b
       \frac{\partial^2}
       {\partial X_{\tau i b}\partial X_{c\tau,jb}}.
\]
For any polynomial $P$ satisfying
$\Delta_{\tau ij}P=0$ for every conjugate pair and all rows, and
sections $F_b$ of the same nonnegative power $L^k$, one has
\[
 \sum_\alpha P_\alpha
      \prod_b\operatorname{nab}_y^{\alpha_b}(F_b)
   =P(\partial)\left(\prod_b\theta_b(z_b)\right)\bigg|_{z=0},
 \tag{A8e}
\]
where $\theta_b$ are their Appell--Humbert readings in the matched
tangent frame.  The scalar coefficients can be complex; the
arithmetic version uses the indicated coefficient field.

Here is a direct proof, which also explains why highest-weight
conditions are unnecessary for this particular equality.  Pair a
polynomial and a formal series by
$[P,F]=P(\partial)F(0)$.  On monomials,
\[
       [P,z_aF]=[\partial_{X_a}P,F].
\]
Indeed both sides are zero unless the relevant exponents agree,
and in the remaining case both are the same product of
factorials.  Iteration proves
$[P,Q(z)F]=[Q(\partial_X)P,F]$ for every polynomial $Q$.
Only finitely many terms of an exponential contribute to this
pairing.  Therefore, writing $Q(z)=\sum_b q(z_b)$, (A8d) gives
\[
 [P,\exp(-kQ)F]
      =[\exp(-kQ(\partial_X))P,F]=[P,F],
\]
because $Q(\partial_X)$ is a linear combination of the displayed
$\Delta_{\tau ij}$ and annihilates $P$.  Apply (A8c) in each
chart, use uniqueness of the exponential coordinates, and compare
with the analytic chart.  This proves (A8e).  In particular every
polynomial supported in just one embedding is covered, regardless
of its row weights.

The archive's \texttt{HHa} is more restrictive: by
\texttt{RowDeg.lean}, it is the joint harmonic space together with
lower-row-unipotent invariance and specified degree in every row.
Those extra conditions select highest-weight vectors for the
automorphic construction.  They are not needed for the
coefficient identity just proved.  Equation (A8e) thus extends the
archive's stated harmonic comparison at this precise point; it
does not assert an extension of its automorphic generation or
integrality theorems to arbitrary weights.

\paragraph{The determinant for A6 and all normalization factors.}
If $n$ theta sections live on the same torus and the chosen
embedding has at least $n$ tangent rows, then
\[
       P=\det(X_{\sigma i b})_{1\leq i,b\leq n}
 \tag{A8f}
\]
belongs to \texttt{HHa}.  It is harmonic because it uses no
$c\sigma$ variables.  Its row degrees are one in the first $n$
rows of $\sigma$ and zero elsewhere.  Replacing a selected lower
row by that row plus a multiple of an earlier row preserves its
determinant; operations on an unselected row leave it unchanged.
This proves the required row invariance, including all rows beyond
$n$.  Its coefficients are integers, so it also satisfies the
coefficient condition in \texttt{harmA}.

In the A6 construction the prescribed line types can differ at
compact places.  Their honest common torus is then the product of
the individual polarized tori.  At the active embedding $c\mu$,
the tensor block has dual tangent label $\mu$ by the preceding
calculation.  Write those tangent coordinates as $z_{ij}$, where $i$ is a
cotangent direction and $j$ is the multiplicity-line factor.  The
product theta section is
$\widehat\Theta(z)=\prod_{j=1}^n\widehat\theta_j(z_j)$, and the
required operator is
\[
 \mathcal D=\det(\partial_{z_{ij}}),\qquad
 \mathcal D\widehat\Theta(0)
       =\det\bigl(\partial_{z_{ij}}
                         \widehat\theta_j(0)\bigr)_{i,j}.
 \tag{A8g}
\]
Expand the determinant.  In each summand precisely one derivative
acts on each separate factor, so (A8g) follows from the product
rule with coefficient one.  This polynomial need not satisfy the
row-degree or row-invariance conditions for the full
$n^2$-dimensional embedding block of the product torus.  It does
use only $\mu$ tangent coordinates, so the proved extension (A8e)
applies.  This closes the analytic comparison for the actual
different-type product used in A6, whenever the indicated
polarized product is placed in the archive datum.

To check the factorial separately, write
$\widehat\Theta(z)=\sum_d\langle T_d,z^{\otimes d}\rangle$
as in 032, (4.10).  For homogeneous $P$ of degree $d$, with the
natural symmetric-tensor pairing,
\[
          P(\partial)\widehat\Theta(0)
                 =d!\,\langle P,T_d\rangle.
 \tag{A8h}
\]
For a multiindex $\alpha$ of degree $d$, both sides equal
$\alpha!$ times the coefficient of $z^\alpha$; this proves the
formula.  The double alternation in A6 maps a symmetric product
of $n$ degree-one tensors to the two exterior products with
coefficient one.  Its determinant functional is $n!$ times the
natural symmetric-tensor pairing in (A8h).  Consequently (A8g),
the A6 wedge, and the archive's ordinary differential convention
have precisely the same determinant coefficient.  There is no
extra $n!$ in the wedge itself.

For the archive's translated theta values, its second comparison
theorem retains the additional scalar $M^{-\deg P}$ and uses
\[
 \mathcal F_{q,a}\theta(z)
   =\exp\left(-q\pi H(z,a)-\frac{q\pi}{2}H(a,a)\right)
       \theta(z+a),\qquad q=M^2.
 \tag{A8i}
\]
The translation is $a=M^{-2}\operatorname{mk}(x_2)$ and the
untranslated function is $w\mapsto\theta(Mw)$.  The oddness,
ideal-divisibility and representative conditions on $M,x,x_2$ in
that theorem remain part of its hypothesis.  They are not
consequences of an arbitrary finite Schwartz input.  Likewise,
the final automorphic convention must retain the archive's
complex conjugations and its negative standard additive
character; the analytic equality (A8e) does not silently change
032's positive additive character.

\paragraph{Specialization to the CM elliptic frame in C18.}
For a CM elliptic curve with the invariant coordinate $z$ and
Weierstrass sigma function, the reduced theta function has the
form
\[
 \theta_{\mathrm{red}}(z)
       =\exp(-e_2^*z^2/2)\sigma(z).
 \tag{A8j}
\]
This is the standard reduced-theta normalization used in
Bannai--Kobayashi, Example 1.9, pp.~10--11 of
\url{https://arxiv.org/abs/math/0610163v4}.
If $H(z,a)=h z\overline a$ with $h>0$ and
$\eta(a)=e_2^*a+\pi h\overline a$, direct substitution in (A8i)
gives
\[
 \exp\left(-\pi H(z,a)-\frac\pi2H(a,a)\right)
          \theta_{\mathrm{red}}(z+a)
       =\exp(-e_2^*z^2/2)\,
             e^{-\eta(a)(z+a/2)}\sigma(z+a).
 \tag{A8k}
\]
Thus the common factor for the canonical analytic reading of
C18's tame translated sections is
\[
      q_{\mathrm{can}}(t)
       =\exp\bigl(-e_2^*\lambda(t)^2/2\bigr)
                         \frac{\sigma(\lambda(t))}{t}.
 \tag{A8l}
\]
Equal-degree ratios cancel this entire factor.  For tensor degree
$d$ the exact normalized operator on its integral coefficient is
$q_{\mathrm{can}}^{-d}D(q_{\mathrm{can}}^d\,\cdot)$.
On an elliptic curve with a CM field acting through a single
tangent embedding, the opposite tangent multiplicity is zero.
The quadratic ambiguity (A8d) is therefore zero, and (A8c)
identifies every nonnegative-order canonical jet with the actual
Taylor derivative of the reduced-theta reading.  Harmonicity
does not discard the nonquadratic factor in (A8l): it is the
specified conversion from the integral local generator to the
canonical analytic generator.

\paragraph{Exact archive source locations.}
All these locators refer to the recovered mirror at commit
\nolinkurl{7cbf7a0c32fca1badf05931879a37c3f8a2b03b4}.
The common prefix for \texttt{IsThRule.lean},
\texttt{ThGerm.lean}, \texttt{IsExpChart.lean}, and
\texttt{RowDeg.lean} is
\path{Formalism/LocalToGlobal/ProjectiveLattices/}.
The prefix for \texttt{ExpTiePf.lean},
\texttt{GaugeUniq.lean}, and \texttt{GaugeUniqPf.lean} is
\path{Formalism/SharpRestriction/Divisibility/GeneratorPowers/}.
The prefix for the three files beginning
\nolinkurl{ExistsAnalyticBiextensionChart},
\nolinkurl{HarmonicValueCanonicalConnection}, and
\nolinkurl{ThetaValueAnalyticFrame}, whose full basenames are given
above, is
\path{Development/AbelianVarietyTheory/Proofs/CuspFormEvaluations/}.
The file \nolinkurl{ExistsIsNabNotIsNablaPin.lean} is under
\path{Development/AbelianVarietyTheory/Proofs/UnitaryExistence/}.
Finally \texttt{ThConst.lean} is under
\path{Development/AbelianVarietyTheory/Definitions/ThetaValues/}.

The arithmetic assertion requires that the chosen imaginary
quadratic CM order, its action, the rigidified polarization and
the invariant differential be the ones in the archive datum.
Neither a general CM field nor an unspecified analytic torus
provides those data automatically.  For the weak tensor-degree
estimate of C18, passage from its raw sigma normalization to
(A8l) additionally requires control of $e_2^*$; equal-degree
ratios do not require it.  If $e_2^*$ is integral, the same Bell
proof extends by allowing the quadratic term: its new weight
inequality is $1/(p-1)\leq2p/(p^2-1)$.  This last observation is
an elementary implication, not an unproved integrality assertion
about $e_2^*$.

\paragraph{The resulting inert-prime specialization.}
Under C19's hypotheses, the required integrality of $e_2^*$ is
proved.  Suppose in addition that the algebraic differential,
rigidification and tame torsion lifts are matched as above, and
that any change of the integral tangent frame has unit determinant.
For a local tensor expression of degree $d$, its canonical
analytic reading is $Q^dv$ with $Q=q_{\mathrm{can}}$ and
$v\in R[[t]]$.  Its order-$r$ canonical local jet is therefore
$D^r(Q^dv)(0)$.  C19 gives the explicit bound
\[
 v_p\bigl(D^r(Q^dv)(0)\bigr)
 \geq\max\left\{0,\frac{r}{p^2-1}
                         -\frac{s_p(r)}{p-1}\right\}.
\]
For equal-degree quotients the common factor cancels and the
bound improves to
\[
 v_p\bigl(D^rv(0)\bigr)
 \geq\max\left\{0,
          \left\lfloor\frac{rp}{p^2-1}\right\rfloor-1\right\}.
\]
For actual archive sections of nonnegative power these are the
jets in (A8c).  For local ratios and negative tensor degree, they
are the induced local connection expressions, obtained by the
product and quotient rules.  They make no assertion about the
unconstrained negative-power values of an arbitrary
\texttt{NabV} satisfying only \texttt{IsNab}.

The level normalization can also be checked exactly.  For the
degree-one reduced theta function and integer $M$, linearity of
$H$ gives
\[
 \mathcal F_{M^2,a}\bigl(z\mapsto\theta_{\mathrm{red}}(Mz)\bigr)(z)
      =\bigl(\mathcal F_{1,Ma}\theta_{\mathrm{red}}\bigr)(Mz).
\]
An order-$r$ derivative on the right contributes $M^r$, which
cancels the archive's $M^{-r}$ in (A8i).  Thus, when $M$ and the
relevant torsion order are prime to $p$, the displayed C19 bounds
apply to this normalized value without an additional level
denominator.  This is an elliptic specialization with specified
normalizations, not an assertion that the higher-dimensional
032 torus has such an elliptic realization.

\framework{The comparison fixes a rigidified line bundle, its
formal completion, and the correspondence between an algebraic
section and its analytic reading.  It is compatible with the
correspondence and polarization viewpoint in Grothendieck's
\emph{cote44}, scans 7--8, and with the dimension and Lefschetz
discussion in \emph{cote16}, scans 43--46.  No algebraic inverse
Lefschetz operator or Chow projector is inferred from equality of
analytic jets.}


\subsection{THETA3-A9: the exact factor-depth invariant and isotropic theta-frame cancellation}\label{A9}

There is a positive comparison behind the archive's divisibility
statements.  It uses more than complex harmonicity.  A specified
matrix of columns is Hermitian-isotropic in the completed CM
algebra, and has a unit minor.  The canonical theta-frame
correction restricts to an integral unit along those columns.
Its tensor power consequently does not change the derivative
image ideal.  This section proves that mechanism, identifies
the exact recovered statements, and distinguishes the
determinant exponent from the degree of the theta line bundle.

\paragraph{The three different integers.}
Let $i$ be the number of tangent rows in one determinant block,
let $a\geq0$ be its exponent, and let $k\geq0$ be the degree of
the theta line bundle.  On $i$ copies of a formal group $G$,
with an integral tangent frame $u_1,\ldots,u_i$, put
\[
 P_{i,a}(X)=\det(X_{rb})_{1\leq r,b\leq i}^{\,a}.
 \tag{A9a}
\]
The differential order is $ia$.  The number $k$ does not occur
in this polynomial.  For the normalized formal exponential
$\operatorname{Exp}_u$, define
\[
 \nu_{\det}(G,u,a)
   =\inf_{f\in R[[G^i]]}
       v_p\!\left(
         P_{i,a}(\partial)
            (f\circ\operatorname{Exp}_u)(0)\right),
 \qquad v_p(p)=1 .
 \tag{A9b}
\]
Here $R[[G^i]]$ means integral formal functions in an integral
coordinate system, and the exponential is evaluated separately
in each copy.  The notation is invariant under an integral
change of coordinates.  The operator is an actual derivative:
for $P=\sum_\alpha P_\alpha X^\alpha$ its value is
$\sum_\alpha P_\alpha\alpha![z^\alpha]f$.
It is not the divided Taylor coefficient.  These are exactly
\texttt{FGDepth.derivAt0}, \texttt{depth}, \texttt{detPow}, and
\texttt{detDepth} in
\path{Development/AbelianVarietyTheory/Definitions/Foundations/IsSatFrame.lean}.

In the archive, $\nu_{\mathrm{Fac}}$ is (A9b), expressed in an
integral chart and then made independent of the chart and of a
closed valued extension.  The canonical local depth
$\nu_{\mathrm{Loc}}(P,k)$ is instead the infimum of valuations
of
\[
 P(D)\left(g\exp\!\left(k\sum_{b=1}^n h_L(t_b)\right)\right)(0)
 \quad\text{over }g\in R[[G^n]],
 \tag{A9c}
\]
where $D$ denotes invariant differentiation and $h_L$ is the
canonical potential from A8.  Thus the exponential correction
is present in this definition.  Removing it without a further
argument would change the invariant.

\paragraph{A formal integral-unit lemma.}
Let $R$ be the valuation ring of a complete characteristic-zero
valued field $K$, with $v_p(p)=1$, and let $G$ be a commutative
formal group over $R$.  Let $\mathcal O$ be a commutative order, finite free
over $\Z$, acting on $G$, with involution $c\mapsto\bar c$.
Write $\ell$ for the normalized formal logarithm and
$\operatorname{Exp}$ for its inverse.  Suppose a formal series
$h(x,y)\in K[[G\times G]]$ vanishes on both axes and has the
following integral biextension cocycles:
\[
\begin{split}
 &\exp\bigl(h(x,z)+h(y,z)-h(x+_G y,z)\bigr),\\
 &\exp\bigl(h(z,x)+h(z,y)-h(z,x+_G y)\bigr),\\
 &\exp\bigl(h([c]x,y)-h(x,[\bar c]y)\bigr)
       \qquad(c\in\mathcal O)
\end{split}
 \tag{A9d}
\]
all belong to $1+(x,y,z)R[[x,y,z]]$, in the applicable variables.
No integrality of $h$ or $\exp h$ is assumed.
Fix $w_0\in\mathcal O$ and put
\[
 q(x)=\exp h(x,[w_0]x),\qquad
 \mathcal Q_n(x_1,\ldots,x_n)=\prod_{j=1}^n q(x_j).
 \tag{A9e}
\]
For the actual canonical potential, $w_0+\bar w_0=1$ and
$q=\exp h_L$.

Suppose matrices $a_N=(a_{N,bj})\in M_{r,n}(\mathcal O)$
satisfy
\[
 \sum_j\overline{a_{N,bj}}a_{N,b'j}\in p^N\mathcal O
       \quad\text{for every }b,b',
 \tag{A9f}
\]
and their formal group maps
\[
 \Phi_N:G^r\longrightarrow G^n,\qquad
 (\Phi_N)_j(x)=\sum_b^G[a_{N,bj}]x_b
\]
converge coefficientwise to a formal map $\Phi$.
Then $\Phi$ is integral and
\[
                 U:=\Phi^*\mathcal Q_n
                     \in 1+(x)R[[G^r]] .
 \tag{A9g}
\]
It suffices for convergence of $\Phi_N$ that the matrices
$d[a_{N,bj}]$ converge: the identity
$\ell([c]x)=d[c]\ell(x)$ expresses every finite jet of
$\Phi_N$ as a polynomial over $K$ in these matrix entries.

\textbf{Proof.}
Every $\Phi_N$ is integral.  Coefficientwise limits remain
integral because the valuation ring is closed.
For a fixed set of formal variables, define the additive
subgroup
\[
 \mathcal L
  =\{s\in(x)K[[x]]:\exp s\in1+(x)R[[x]]\}.
\]
It is a subgroup: multiplication of integral units corresponds
to addition, and the inverse of a series with integral
coefficients and constant term one is integral, by the
coefficient recursion for its inverse.
For integral formal points $x,y$, let $B(x,y)$ be the class of
$h(x,y)$ modulo $\mathcal L$.  Substitution of integral
zero-constant-term points into (A9d) shows that $B$ is
biadditive for the formal group law and satisfies
\[
             B([c]x,y)=B(x,[\bar c]y).
 \tag{A9h}
\]
In particular it is additive in the endomorphism applied to
either argument, including negative integer multiples.

Put $S_N=\sum_j h((\Phi_N)_j,[w_0](\Phi_N)_j)$.
Repeated use of biadditivity and (A9h) gives, modulo
$\mathcal L$,
\[
 S_N\equiv
  \sum_{b,b'}h\left(x_b,
       \left[w_0\sum_j\overline{a_{N,bj}}a_{N,b'j}\right]x_{b'}
            \right).
 \tag{A9i}
\]
Choose a fixed $\Z$-basis $e_1,\ldots,e_s$ of $\mathcal O$.
By (A9f), the endomorphism in brackets has the expression
\[
 \sum_\lambda p^N m_{N,bb'\lambda}e_\lambda,
 \qquad m_{N,bb'\lambda}\in\Z.
\]
Hence
\[
 S_N-R_N\in\mathcal L,\qquad
 R_N=\sum_{b,b',\lambda}
       p^N m_{N,bb'\lambda}h(x_b,[e_\lambda]x_{b'}).
 \tag{A9j}
\]
Every coefficient of $R_N$ tends to zero.  Indeed there are
only finitely many fixed series in (A9j), and integers have
nonnegative $p$-adic valuation.  All these series have zero
constant term.  Each coefficient of $\exp(-R_N)$ is therefore
a finite polynomial in coefficients of $R_N$, and
$\exp(-R_N)\to1$ coefficientwise.

By construction
$\exp(S_N-R_N)=\Phi_N^*\mathcal Q_n\exp(-R_N)$ is an integral
unit with constant term one.  Its coefficientwise limit is
$\Phi^*\mathcal Q_n$, because composition of zero-constant-term
formal maps uses only finitely many coefficients in each
degree.  Closedness of $R$ proves (A9g), and its constant term
is one.  This proves the lemma. \hfill$\square$

The proof uses the entire canonical frame, including its
higher terms.  It is distinct from the quadratic
mixed-Laplacian cancellation in A8.  Its additional input is
the Hermitian congruence (A9f).

\paragraph{Equality of derivative image ideals.}
Here is the precise consequence for arbitrary differential
polynomials.  Let $P$ be a polynomial in variables describing
chosen tangent directions on the $r$ source copies, and write
$E_u=\operatorname{Exp}_u$ for the corresponding formal
exponential map.  Suppose there is an integral zero-constant-term
map $\rho:G^n\to G^r$ with
\[
                   \rho\circ\Phi\circ E_u=E_u.
 \tag{A9k}
\]
Only this identity along the chosen exponential directions is
needed; a retraction on the entire formal group is sufficient
but is not required.  For any integer $k$, define
\[
 J_k(g)
   =P(\partial)\bigl(((\Phi^*g)\,U^k)
                         \circ E_u\bigr)(0),
          \qquad g\in R[[G^n]].
 \tag{A9l}
\]
Then
\[
 \{J_k(g):g\in R[[G^n]]\}
   =\{P(\partial)(f\circ E_u)(0):f\in R[[G^r]]\}.
 \tag{A9m}
\]
Both are $R$-submodules of $K$.  In particular their
valuation infima agree and are independent of $k$.

\textbf{Proof.}
For any integral $g$, the series
$(\Phi^*g)U^k$ is integral, because $U$ and its inverse are
integral.  This proves one inclusion.
Conversely take an integral $f$ and set
$g=\rho^*(fU^{-k})$.  This is integral.  Equation (A9k) implies
\[
 ((\Phi^*g)U^k)\circ E_u
   =((fU^{-k})\circ E_u)(U^k\circ E_u)
   =f\circ E_u.
 \]
Applying $P(\partial)$ proves the other inclusion.
The proof also establishes the assertion for negative $k$
as a local statement.  The global \texttt{IsNab} axioms in A8
still constrain actual sections only for nonnegative degree.
\hfill$\square$

For a product of determinant polynomials in disjoint blocks,
the ordinary image ideal is the product of the block image
ideals.  To see this without a tensor-completion assumption,
the operator has finite order and hence depends on a finite
jet of its argument.  Every monomial in that finite jet is a
product of block monomials, on which the functional factors.
This gives containment in the product ideal.  Conversely,
products of functions realizing block values realize the
products of those values.  The image ideals are finitely
generated by the finitely many jet monomials.  Their valuation
infima are minima, and valuation of a product adds.  Therefore
the total depth is the sum of the factor depths.

\paragraph{The exact recovered column condition.}
The preceding hypotheses are concrete in the recovered source.
In
\path{Development/AbelianVarietyTheory/Definitions/Divisors/NonemptyValExtOfOverLa.lean},
lines 327--344, \texttt{polyP} and \texttt{polyT} are products
of determinant powers
\[
 P_B(X)=\prod_f
   \det\left(\sum_{j=1}^n B_{f,bj}X_{\tau_f,rj}
                    \right)_{1\leq b,r\leq i_f}^{a_f}.
 \tag{A9n}
\]
For each $f$ with $a_f\ne0$, \texttt{IsCols} requires a
disjoint block of $2i_f$ section columns and
\texttt{IsIsoCols} requires the following three properties:
the matrix is supported in that block; an $i_f$ by $i_f$
minor is a unit; and for every $N$ it is congruent modulo
$p^N$, at the specified CM embedding, to a matrix
$x_N$ with entries in $\OO_E$ satisfying
$x_N\overline{x_N}^{\,t}\equiv0\pmod{p^N\OO_E}$.
The condition is stronger than belonging to \texttt{HHa}.
In particular a single nonzero column on an inert CM
elliptic factor does not satisfy it.

The source's prime-idempotent approximation constructs
limits of the $\OO_E$-actions whose tangent actions are the
prescribed entries $B_{f,bj}$ on the selected embedding
directions.  Applied blockwise, (A9f) persists because the
blocks are disjoint.  Its integral formal map is the
\texttt{InclSt} map $\psi$.
The unit minor constructs the \texttt{RetrSt} map $\rho$:
choose that minor, invert it in the local integral ring,
and use the limiting CM actions of its inverse entries.
The inverse entries remain in the same completed coefficient
ring by the adjugate formula and unit determinant.  On the
selected tangent directions the resulting matrices multiply
to the identity; uniqueness of the formal exponential then
gives (A9k).  Other blocks contribute zero in the selected
columns, by disjoint support.  This is precisely the
calculation in
\nolinkurl{ColumnInclusionRetractionOfLimitSteps.lean}.
It does not require an integral projector onto one complex
embedding inside an inert local field.

Here are exact source locators for this mechanism.  In the
following list the common full prefix is
\path{Development/AbelianVarietyTheory/Proofs/DivisorProofs/}.
The definitions \texttt{IsIntFB}, \texttt{LogUnit},
\texttt{BiaddSt}, \texttt{LimFrameSt}, \texttt{InclSt}, and
\texttt{RetrSt} are in \texttt{MvIsIntMapMap.lean}.
The full proofs of the integral-unit mechanism are
\nolinkurl{CanonGaugeBiadditiveOfLogSumActionLinear.lean},
\nolinkurl{GaugeHermitianCombinationOfSteps.lean},
\nolinkurl{LimitIsotropicFrameIntegralOfSteps.lean}, and
\nolinkurl{ColumnInclusionIntegralOfLimitSteps.lean}.
The exact operator substitution is
\nolinkurl{NablaOperatorFrameSeriesOfSteps.lean},
the inverse-unit construction is in
\nolinkurl{ExistsIntegralGermOnCopiesSameValue.lean}, and
the depth equality is theorem \texttt{nuLocOf\_eq\_nuP}
in \nolinkurl{ColumnPolynomialLocalDepthEqPlaceDepth.lean}.
It states, for every theta degree $k$ and every
\texttt{IsCols} matrix,
\[
 \nu_{\mathrm{Loc}}(P_B,k)
       =\sum_f\nu_{\mathrm{Fac}}(f,a_f),
 \tag{A9o}
\]
with the local valuation normalized by dividing the DVR
valuation by $v(p)$.  These proof bodies were read in this
run.  The source is the pinned secondary mirror commit
\nolinkurl{7cbf7a0c32fca1badf05931879a37c3f8a2b03b4};
this is a source reconstruction, not a Lean build or a
transitive certification of all its dependencies.

\paragraph{The exact inert-prime bound and its scope.}
Write $j=\mathrm{Dw.q}$ and $P_0=\mathrm{Dw.p}$ to avoid
confusing the datum with the rational prime $p$.
The host contains the imaginary quadratic field generated
by its specified square root of $-d$, and
\texttt{InPsi} selects one of that field's embeddings.
For an odd prime unramified in $E$ and inert in this
quadratic field, the recovered lower bound is
\[
 \nu_{\mathrm{Fac}}(f,a)
 \geq
   \frac{(i-c_r)ap}{p^2-1}
       -i^2(\log_p a+1),\qquad
 0\leq c_r\leq i,\quad
 c_r\leq j\left(\frac32+\frac12\log_p a\right),
 \quad a\geq1,
 \tag{A9p}
\]
for a saturated rank-$i$ frame in an \texttt{InPsi} direction.
The precise formal-group theorem is
\texttt{detDepth\_lower} in
\nolinkurl{DetDepthLowerBoundSaturatedEvenType.lean};
its application to the PEL datum is
\texttt{facBetween\_fine} in
\nolinkurl{FactorDepthBoundedOddInertPrimeCmType.lean}.
Both full proof bodies were read, with their imported
formal-group radius and typed-basis inputs kept as source
inputs.  At an algebraic point the companion upper bound is
\[
 \nu_{\mathrm{Fac}}(f,a)
       \leq a\left(\frac{j}{p-1}
                    +\frac{(P_0-j)p}{p^2-1}\right);
 \tag{A9q}
\]
its exact statement is \texttt{nuFac\_le\_odd\_unram} in
\nolinkurl{FactorDepthUpperBoundOddUnramifiedPrime.lean}.

For $j=0$ and $i=1$, (A9p) becomes
\[
 \nu_{\mathrm{Fac}}(f,a)
       \geq \frac{ap}{p^2-1}-\log_p a-1
       =\frac ap+\frac a{p^3}+\frac a{p^5}+\cdots
                                      -\log_p a-1.
 \tag{A9r}
\]
This is an exact recovered occurrence of the odd-power
divisibility slope in step 2 of the supplied timeline.
It is a statement about (A9b); applying it to a theta
section uses (A9o), its isotropic column data, and the
actual canonical normalization.
The timeline's period symbol must still be identified with
the chosen algebraic invariant differential before its
display can be substituted literally.

The source does transfer this depth to actual theta
generators.  Theorem
\nolinkurl{harmonic_value_ge_local_depth} in
\nolinkurl{HarmonicValueGeLocalDepthPrimeToP.lean}
uses a good integral chart, \texttt{IsNab}, an arithmetic
harmonic polynomial, a lower bound uniform in $k$, and
odd tame square-section data with level a unit at $p$.
It concludes the corresponding lower bound for the
canonical harmonic value.  Its proof reads reached
sections as integral local germs, uses the actual positive
exponential correction, and then applies the defining
infimum.  Thus the frame issue is resolved for these
specified operators and integral theta inputs.  It is
not resolved for every raw one-copy Taylor coefficient.

\paragraph{An actual elliptic instance and the comparison still required.}
C26 proves the integral-unit mechanism without relying on
the recovered archive.  For a good CM elliptic curve at an
inert prime, choose a unit $b$ in the completed CM order with
$b\bar b=-1$.  If $q(t)=\theta_{\rm red}(\lambda(t))/t$, then
\[
                 q(t)q([b]t)\in1+tR[[t]].
 \tag{A9s}
\]
The two columns $(1,b)$ have a unit minor and are
Hermitian-isotropic.  Actual normalized tame translated
theta sections on the two copies therefore restrict to an
integral function.  C19 gives the stronger explicit bound
\[
 v_p\!\left(D^a
    \bigl(s_{T_1}(t)s_{T_2}([b]t)\bigr)\big|_{t=0}\right)
   \geq
   \max\left\{0,\left\lfloor\frac{ap}{p^2-1}\right\rfloor-1\right\}.
 \tag{A9t}
\]
The differential is exactly $(D_1+bD_2)^a$ on the product.
This proves a positive-degree, actual-theta realization of
the strong slope on a specified formal isotropic direction.
The completed CM element need not define a map of complex
abelian varieties.

\paragraph{An explicit improvement in the product-formal case.}
There is a stronger quantitative statement under an additional
integral splitting hypothesis.  Suppose the selected $i$ tangent
directions on each of the $i$ source copies come from independent
one-dimensional strict height-two CM formal modules over a complete
DVR $R$, as in C19 and C21.  Normalize each invariant derivative
by its integral invariant differential.  If a CM isogeny is used
to obtain this product, require its degree to be prime to $p$;
it then gives an isomorphism of the formal completions.  Extra
formal coordinates may be retained as coefficient variables.
For $m\geq0$ put
\[
 \ell_p(m)=\max\left\{0,
          \left\lfloor\frac{mp}{p^2-1}\right\rfloor-1\right\},
 \qquad
 B_{i,p}(a)=
 \min_{\substack{\alpha\in\Z_{\geq0}^{i\times i}\!:\\
               \sum_b\alpha_{rb}=a\ (\forall r),\,
               \sum_r\alpha_{rb}=a\ (\forall b)}}
                   \sum_{r,b}\ell_p(\alpha_{rb}).
 \tag{A9u}
\]
The minimum is over a nonempty finite set, including when
$a=0$.  Then
\[
 \nu_{\det}(G,u,a)\geq B_{i,p}(a)
       \geq\max\left\{0,
                  \frac{iap}{p^2-1}-2i^2\right\}.
 \tag{A9v}
\]
This gives an explicit finite combinatorial certificate and a
constant-loss estimate in this product-formal regime.

\textbf{Proof.}
C19 proves operator divisibility, before evaluation at the
origin: the $m$th power of the invariant derivative on a strict
factor carries integral formal functions into
$p^{\ell_p(m)}R[[t]]$.  Apply this coefficientwise to all the
other variables.  The derivatives on the independent factors
commute, so the monomial $\prod_{r,b}D_{rb}^{\alpha_{rb}}$
carries integral functions into
$p^{\sum_{r,b}\ell_p(\alpha_{rb})}R[[G^i]]$.
Every monomial of $\det(D_{rb})^a$ has all row sums and all
column sums equal to $a$: each of its $a$ permutation factors
uses exactly one entry in every row and column.  Its coefficient
is an integer.  Consequently the determinant expansion introduces
no denominator, and every summand is divisible by
$p^{B_{i,p}(a)}$.  Evaluation and the defining infimum give
the first inequality in (A9v).  Finally
$\ell_p(m)\geq mp/(p^2-1)-2$ and $\ell_p(m)\geq0$.
Summing over the $i^2$ entries, whose total is $ia$, proves
the second inequality. \hfill$\square$

Under the isotropic-column and retraction hypotheses above,
(A9m) transfers the same bound to every theta degree $k$.
This improvement does not assert that a general typed formal
group in the recovered archive has such an integral product
decomposition.  For an actual 032 CM evaluation point, the
isogeny and differential normalizations must be checked at the
chosen prime before using it.

For the 032 top wedge, A8 identifies its actual complex
canonical derivative once the torus, section, frame, and
level are matched.  Its determinant in the original
section columns does not automatically satisfy
\texttt{IsCols}.  Replacing it by an isotropic doubled-column
operator changes the input unless a separate comparison
identifies the resulting automorphic coefficient.
This is now a concrete comparison of column matrices,
tensor degree, and derivative normalization.  C22 does
not contradict the recovered isotropic statement.
Finally, (A9s)--(A9t) are identities of formal germs and
their derivatives, not algebraic correspondences.
Within the requested Grothendieck framework, a passage
to homological motives must still specify the
correspondence and its polarization (cote44, scans 7--8)
and any Lefschetz inverse used for a degree change
(cote16, scans 43--46).


\subsection{THETA3-A10: stabilization to one finite CM theta module}
\label{A10}

The different CM theta tori can be put into one actual product
without identifying their CM types.  At a fixed CM point this
produces a common finite theta module and one comparison scalar.
The price is an explicit theta-null multiplier on each original
input.  This construction removes a finite-dimensional matching
problem for the stabilized values; it retains that multiplier
when one returns to the original automorphic coefficients.

\paragraph{The product theta module.}
Let the finite list of tori and theta lines be
$(A_i,M_i)$, $1\leq i\leq h$, with the origin rigidifications,
finite theta structures and normalizations of B19.  Write
\[
 J_i:W_i\longrightarrow H^0(A_i,M_i),\qquad
 \mathcal T_i=c_i(J_i\otimes\C),\quad c_i\ne0.
 \tag{A10a}
\]
Here $J_i$ is algebraic over $\overline\Q$, and
$\mathcal T_i$ is the actual completed analytic theta map.
Assume that an algebraic even section $R_i=J_i\rho_i$ has
been chosen with value one at zero in the specified
rigidification.  B20 gives such sections integrally at an
odd good prime under its explicit evaluation-surjectivity
hypothesis.  Then
\[
 \mathcal T_i(\rho_i)(0)=c_i,
       \qquad d\mathcal T_i(\rho_i)(0)=0.
 \tag{A10b}
\]
The derivative statement uses the symmetric origin frame;
the canonical correction has no linear term.

Set
\[
 A_* =\prod_{i=1}^h A_i,\quad
 M_* =\boxtimes_{i=1}^h M_i,\quad
 W_* =\bigotimes_{i=1}^h W_i,\quad
 J_* =\bigotimes_{i=1}^hJ_i,
 \qquad c_* =\prod_{i=1}^h c_i.
 \tag{A10c}
\]
The map $J_*$ is an algebraic isomorphism onto
$H^0(A_*,M_*)$.  The product finite theta structure is
specified, and the completed theta map of the direct-sum
symplectic datum is exactly
\[
                    \mathcal T_*=c_*(J_*\otimes\C).
 \tag{A10d}
\]
Indeed the Gaussian completion is a product, and the theta
series of a pure tensor is the product of the factor series.
Absolute convergence justifies the separation of the sums.
Pure tensors span $W_*$, so this proves (A10d) on the whole
finite module.  Equivalently, the product finite-Heisenberg
representation supplies the same identity.  If
$\dim W_i=N_i^{2\dim A_i}$, its dimension is explicitly
$\prod_iN_i^{2\dim A_i}$.

For $\varphi\in W_i$ put
\[
 \iota_i\varphi=
       \rho_1\otimes\cdots\otimes\rho_{i-1}
            \otimes\varphi\otimes
           \rho_{i+1}\otimes\cdots\otimes\rho_h.
 \tag{A10e}
\]
This is an actual finite theta input.  If $u$ is a tangent
vector in the $i$th factor, then
\[
 D_u\mathcal T_*(\iota_i\varphi)(0)
       =\frac{c_*}{c_i}\,D_u\mathcal T_i(\varphi)(0).
 \tag{A10f}
\]
This follows from the product rule and (A10b).  The factor
$c_*/c_i$ is exactly the product of the missing reference
theta-null values.  It has not been declared algebraic.

\paragraph{One algebraic tangent frame and the determinant values.}
Use the algebraic first-jet frames of B19.  Absorb its
algebraic matrices into the tangent vectors and write
$w_{ir}$ for the resulting vectors, $r=1,2$.  Put
\[
 w_{*r}=(w_{1r},\ldots,w_{hr})\in\operatorname{Lie}(A_*),
 \qquad
 j_{ir}(\varphi)=D_{w_{ir}}(J_i\varphi)(0).
 \tag{A10g}
\]
The product rule gives
\[
 D_{w_{*r}}J_*(\iota_i\varphi)(0)=j_{ir}(\varphi),
 \qquad
 D_{w_{*r}}\mathcal T_*(\iota_i\varphi)(0)
                                     =c_*j_{ir}(\varphi).
 \tag{A10h}
\]
All terms differentiating a reference vanish.  The same
algebraic rows therefore work for every input family.
For two inputs, the determinant of their stabilized rows is
$c_*^2$ times the determinant of their algebraic jet rows.
There is no derivative factorial.

To compare with the original analytic ball directions, retain
B19's exact normalization
\[
  D_{u^z_{ir}}\widehat\vartheta_{i,\varphi}(0)
          =\frac{c_ip_i}{N_i}j_{ir}(\varphi).
 \tag{A10i}
\]
The periods $p_i$ are external scalars in this equation;
they are not coefficients of the integral row frame in
(A10g).  Define the stabilized value by multiplying the
$i$th original row by $c_*/c_i$.  For the pair $i,j$ its
first-jet determinant is consequently
\[
 \widetilde\Delta_{ij}
   =\frac{c_*^2}{c_ic_j}\Delta_{ij}
   =\frac{c_*^2p_ip_j}{N_iN_j}
       \det\begin{pmatrix}j_{i1}&j_{i2}\\j_{j1}&j_{j2}\end{pmatrix}.
 \tag{A10j}
\]
Any prescribed constant-line factor $\kappa_{ij}$ multiplies
both sides.  With compatible algebraic Betti frames as in
B21, ratios of these factors are algebraic after the common
oscillator normalization has been separated.

The construction works on the finite list of transformed
inputs in B19's source class sums and definite Bessel sums.
Choose the references once.  Equations (A10f)--(A10j) then
hold term by term; all Haar volumes, stabilizer masses and
finite character values remain those of B19.  Thus the
stabilized finite functionals have the one common theta
factor $c_*^2$.

\paragraph{The balanced tangent-period relation.}
Here is the exact positive CM argument for the remaining
$p_i$.  Assume the accepted CM Hodge conjecture for products
of the auxiliary CM abelian varieties.  For sign functions
$d_1+d_2=d_3+d_4$ and the common selected label $\mu=1$,
the tensor lines
\[
 U(d_1)_\mu\otimes U(d_2)_\mu,
 \qquad
 U(d_3)_\mu\otimes U(d_4)_\mu
 \tag{A10k}
\]
have the same Hodge type after every scalar conjugation.
In fact their holomorphic degrees after $g$ are
$1+(d_1(g)+d_2(g))/2$ and the corresponding expression for
$3,4$, which are equal.  Choose algebraic-Betti basis
vectors and extend a nonzero map of the selected lines by
zero on the other tensor embedding lines.  Every scalar
conjugate is a Hodge map.  The scalar-descent proof of 032,
Lemma~2.1, p.5, in the rational Hom space therefore expresses
this map as an algebraic-coefficient combination of rational
Hodge maps; the label convention is 032, equation~(2.3), p.7.

By the stated CM Hodge input, these maps are induced by
algebraic correspondences.  The degree-one K\"unneth factors
are cut out by polynomials in multiplication-by-an-integer,
and the CM embedding factors by the algebraic $E$ action.
Polarization correspondences implement the duals.  The cycle
classes can be represented over $\overline\Q$: spread a
cycle in the fixed projective product to its Hilbert-scheme
component and specialize to a $\overline\Q$ point, preserving
its locally constant Betti class.  Thus the resulting maps
also have algebraic de Rham matrices.

Write the comparison on the two selected tensor lines as
multiplication by their period products.  Commutation with
a nonzero correspondence gives
$P_{12}a_{\rm dR}=P_{34}a_{\rm B}$ with
$a_{\rm dR},a_{\rm B}\in\overline\Q^\times$.
The quotient of the periods is algebraic.  Dualizing gives
exactly the same assertion for B19's tangent convention:
\[
                         \frac{p_1p_2}{p_3p_4}
                                  \in\overline\Q^\times.
 \tag{A10l}
\]
This proof uses the selected embedding lines.  It does not
replace an ordinary top Hodge determinant by an $E$-linear
determinant; B21 shows why that replacement would be wrong.

\paragraph{Integral columns and the precise remaining comparisons.}
At a chosen good prime, impose the integral reference,
frame and saturation hypotheses of B20 on $A_*$ and $M_*$.
The diagonal $E$ action has the same eigenvalue on every
selected $\mu$ row.  B20 therefore supplies the completed-CM
isotropic columns and identifies the exponent-one determinant
value with a doubled-column value on copies of this same
$A_*$.  A9 removes the entire canonical frame along those
columns.  This is a common actual finite theta module and a
common actual first-jet functional, not a formal definition
chosen to mimic an analytic value.

It is consequently sufficient to check an asserted arithmetic
comparison on one basis of $W_*$ and the two indicated
first-jet rows.  Such a check still requires a genuine
archive datum realizing this polarized CM torus, its integral
model and its tangent and theta frames.  The product need
not have the uniform CM multiplicities or the real signatures
of a specified recovered datum.  Those hypotheses are not
provided by taking a product.

Returning to the raw pair values in (A10j) requires division
by the explicit factors $c_*^2/(c_ic_j)$.  Stabilization does
not prove their cross-type ratios algebraic.  Nor is it an
equality of the original global automorphic forms: the extra
reference theta-null functions vary with the ball point and
carry their own line-bundle weights.  The fixed-point
construction has removed the finite-module and first-jet
matching issue for the stabilized functional; the global
operator, signature and period comparisons remain separate.
Within the requested Grothendieck framework, (A10l) is a
statement about algebraic correspondences in homological
motives and their polarizations (cote44, scans 7--8), whereas
(A10f)--(A10j) are equalities of explicitly normalized
analytic and algebraic first jets.\hfill$\square$


\subsection{THETA3-A11: an explicit host extension and actual finite-jet realization}
\label{A11}

The finite CM host condition can be achieved constructively.  The
construction below retains the actual CM constituents, produces a
principal polarization and the archive's primitive symmetric line,
and then realizes a prescribed finite family of first jets.  Its
output is a generic \texttt{Host}/\texttt{Datum}/\texttt{Obj}.  An
identification with a particular \texttt{StdPair}, oscillator kernel,
or global iterated theta chain is additional data.

\paragraph{The exact source conditions.}
Put
\[
 \Psi=\{\tau:E\longrightarrow\C:\tau(\mathrm{rt})=\mathrm{rt}_0\}.
\]
The recovered source
\path{ReciprocityTheory/ReductiveGroups/ArithmeticGeometry/}
\path{BirationalGeometry/ToroidalCompactification/}
\nolinkurl{TwoWeightEqOneAddRoot.lean},
\texttt{Host}, \texttt{Datum}, \texttt{Level}, \texttt{SAB},
and \texttt{Obj}, requires the following objects.  The field host
contains an integral square root of $-d$, where $d$ is a prime
congruent to $7$ modulo $16$, and $w_0=(1+\mathrm{rt})/2$ is
integral.  The extension over the maximal real subfield is
unramified at every finite prime.  A datum is an actual finite
projective $\OO_E$-lattice with a unimodular alternating integral
form, with adjoint involution complex conjugation, and a compatible
positive complex structure.  Its Lie multiplicities are the same
integer $p$ at every embedding in $\Psi$ and the same integer $q$
at the conjugate embeddings.  There is no condition $p\geq q$ or
parity condition in this definition.  The exact multiplicity
function is also written in \texttt{FlagBasis.lean},
\texttt{Datum.mult}.  The additional line-bundle identity is
\[
       L\simeq(\mathrm{id},[w_0])^*P_L,\qquad
       P_L=m^*L\otimes\operatorname{pr}_1^*L^{-1}
                    \otimes\operatorname{pr}_2^*L^{-1},
 \tag{A11a}
\]
with the origin rigidification understood.  This sign is the
literal two-variable \texttt{cubeLB} in
\path{ReciprocityTheory/ReductiveGroups/ArithmeticGeometry/}
\path{SemiabelianSchemes/CubicalStructures/}
\nolinkurl{CubeSymmetricSwapFirst.lean}.

\paragraph{A universal field host of degree at most four.}
Let $E$ be any CM field.  Define
\[
 K=\Q(\sqrt{-7}),\qquad F_H=\Q(\sqrt{42}),\qquad
 H=K F_H,\qquad E'=EH,\qquad F'=(E')^+.
 \tag{A11b}
\]
Then
\[
 [E':E]\leq4,\qquad
 E'=F'K,\qquad
 \OO_{E'}\text{ is finite \'etale over }\OO_{F'}.
 \tag{A11c}
\]
Here $\mathrm{rt}=\sqrt{-7}$ and $w_0=(1+\sqrt{-7})/2$ give
all the field-generator conditions with $d=7$.

We prove the finite-prime assertion directly.  The quadratic field
$K$ has discriminant $-7$, so its integer ring is finite
\'etale away from $7$.  Base change to $F_H$ therefore gives
an unramified quadratic algebra away from $7$.  At $7$, the element
$-6=42/(-7)$ is a square in $\Q_7$: apply Hensel's lemma to
$X^2+6$ at $X=1$, whose derivative is a unit modulo $7$.
Consequently
\[
 \Q_7(\sqrt{42})=\Q_7(\sqrt{-7}),
\]
and $H/F_H$ is split at the prime above $7$.  It is thus unramified
at every finite prime.  The discriminant check is
$(-7)(-24)(168)=168^2$ for the three quadratic subfields, in
agreement with this local proof.

A compositum of CM fields is CM, so $E'$ is CM.  Its maximal real
subfield contains $F_H$, and $K$ is not contained in that real
field.  Since $[E':F']=2$, one has $E'=F'K$.  The finite
\'etale $\OO_{F'}$-algebra
$\OO_H\otimes_{\OO_{F_H}}\OO_{F'}$ has generic algebra $E'$.
It is normal, because finite \'etale algebras over a Dedekind
domain are normal.  It therefore equals the integral closure
$\OO_{E'}$.  This proves (A11c), including the integral assertion
required by \texttt{Host.unram}.  Finite \'etale base change
is Stacks, Lemma 29.37.4, tag \texttt{02GO}.  The CM compositum
and Galois-closure facts are Milne, \emph{Complex Multiplication}
(2020), Proposition 1.4 and Corollary 1.5, printed p.~10.

If $E/\Q$ is Galois, as in the original 032 setup, $E'/\Q$ is
already Galois.  In general one may replace $E'$ by its CM Galois
closure.  This closure still contains $H$, so precisely the same
base-change argument proves finite unramifiedness over its own
maximal real subfield.  The bound $[E':E]\leq4$ refers to (A11b),
before this optional Galois closure.

The same construction is available for any prescribed prime
$d\equiv7\pmod {16}$.  Replace $H$ by
$H_d=\Q(\sqrt{-d},\sqrt{d(d-1)})$ and its real subfield by
$\Q(\sqrt{d(d-1)})$.  Away from $d$ the quadratic field of
discriminant $-d$ is unramified, and at $d$ the ratio
$d(d-1)/(-d)=1-d$ is a square in $\Q_d$ by Hensel's lemma
at $1$.  The same local and integral base-change proof applies.
This variant permits a prescribed size condition on $d$
once such a prime has been chosen; it does not itself supply
the additional arithmetic parameters of the global archive.

\paragraph{Scalar extension of the actual CM abelian variety.}
Let the positive-dimensional $A/\overline\Q$ have an action of $R=\OO_E$ and an
$E$-compatible polarization.  Let
$s=[E':E]$, $S=\OO_{E'}$, and $S^*=\Hom_R(S,R)$.  Define
\[
       A'=\Hom_R(S^*,A),
       \qquad H_1(A',\Q)=H_1(A,\Q)\otimes_E E'.
 \tag{A11d}
\]
The first expression denotes a represented group functor.
Since $S$ is finite projective over the Dedekind domain $R$,
$S^*$ is a direct summand of a finite free module.  Applying
$\Hom_R(-,A)$ to a splitting realizes $A'$ as an actual direct
factor of a power of $A$.  It is therefore an abelian variety.
Multiplication on $S$ gives its $S$-action.  The same idempotent
construction on homology proves the second equality, and gives
\[
       \dim A'=s\dim A,\qquad A'\sim A^s.
 \tag{A11e}
\]
Here the isogeny is $E$-linear; it is not an assertion that $A$
already carried the larger field action.  This is the projective
Serre construction in Milne, \emph{Complex Multiplication},
Proposition 7.33, Remark 7.34, and Proposition 7.35, printed
pp.~62--63.  The proof just given also explains its compatibility
with abelian schemes and base change.

There is even an $R$-linear direct-factor inclusion of $A$ in
$A'$.  Indeed, $S/R$ is torsion free: if $0\ne a\in R$ and
$ax\in R$ for $x\in S$, then $x\in E$ and its integrality gives
$x\in R$.  Thus $S/R$ is projective and $R\cdot1\subset S$
splits.  Tensoring this splitting with $A$ gives the inclusion.
It is generally not $E'$-linear.  For an embedding $\tau'$ of
$E'$ extending $\tau$ of $E$, the induced Lie decomposition is
\[
       \operatorname{Lie}(A')_{\tau'}
                       =\operatorname{Lie}(A)_\tau.
 \tag{A11f}
\]
More precisely, project the preceding inclusion to the
$\tau'$-eigenspace; its restriction to the $\tau$-space is the
displayed isomorphism.  All these maps and eigenspaces are
algebraic over a finite number field.  Integrality of an
eigenspace projection is only asserted after its denominators
are inverted, or at an unramified split coefficient base where
the corresponding idempotent is integral.

For completeness, the extended action has a compatible
polarization.  If $\psi$ is the original rational alternating
Riemann form, trace nondegeneracy gives a unique $E$-valued form
$h$ with
\[
       \Tr_{E/\Q}(a h(x,y))=\psi(ax,y),\qquad a\in E.
\]
The adjoint identity for $\psi$ implies that $h$ is linear in
the first variable, conjugate linear in the second, and
$h(y,x)=-\overline{h(x,y)}$.  Extend $h$ to $E'$ and use
$\Tr_{E'/\Q}h$.  At every complex embedding this is the same
positive Riemann form as the original one, on each of the
additional summands.  It satisfies the Riemann relations.
Multiplication by a positive integer makes it integral on
$H_1(A',\Z)$ and supplies the asserted polarization.

\paragraph{Minimal uniform multiplicity padding.}
For this paragraph write $E$ for the enlarged field and
$r=\dim_E H_1(A',\Q)$.  Put
$m_\tau=\dim_{\overline\Q}\operatorname{Lie}(A')_\tau$.
The Hodge decomposition and conjugation give
$m_\tau+m_{c\tau}=r$.  Define
\[
 P=\max_{\tau\in\Psi}m_\tau,\qquad
 m=\min_{\tau\in\Psi}m_\tau,\qquad
 Q=r-m,\qquad R=P-m.
 \tag{A11g}
\]
There are $R$ rank-one CM abelian varieties $B_1,\ldots,B_R$
such that
\[
        B=A'\times\prod_{j=1}^R B_j
 \tag{A11h}
\]
has multiplicities $(P,Q)$ at every pair $(\tau,c\tau)$ with
$\tau\in\Psi$.  To construct $B_j$, choose its CM type to
contain $\tau\in\Psi$ exactly when $j\leq P-m_\tau$, and
otherwise to contain $c\tau$.  This selects exactly one
embedding from every conjugate pair.  The added multiplicities
at that pair are $P-m_\tau$ and $m_\tau-m$, respectively,
which proves the assertion.

Every such type is realized by a polarized rank-one CM torus:
use an $\OO_E$-lattice in $E$, choose a purely imaginary
element with the signs prescribed by the type, and take the
trace alternating form.  Weak approximation in the maximal
real subfield supplies the signs, and a positive integer clears
its denominators.  The Riemann relations give an abelian
variety, and CM descent gives a model over $\overline\Q$
(Milne, Corollary 7.10, printed p.~55).  Weak approximation is
Milne, \emph{Algebraic Number Theory} (2020), Theorem 7.20,
printed p.~114.  These are actual added constituents.  No
isogeny between different full CM types is inferred from their
local multiplicity counts.

The number $R$ is minimal for padding by rank-one CM factors.
Indeed, any uniform target $(P',Q')$ must satisfy
$P'\geq P$ and $Q'\geq r-m$.  The number of added factors is
$P'+Q'-r$, hence is at least $P-m=R$.  In particular $R\leq r$.

\paragraph{An explicit compatible principal polarization.}
Let $\lambda:B\to B^\vee$ be the product compatible
polarization.  Choose $N\geq1$ annihilating $\ker\lambda$;
then $\lambda'=N\lambda^{-1}:B^\vee\to B$ is an actual
homomorphism.  Choose integers $a,b,c,d$ whose squares sum to
$N-1$, and form
\[
 T=\begin{pmatrix}
 a&b&c&d\\-b&a&-d&c\\-c&d&a&-b\\-d&-c&b&a
 \end{pmatrix},\qquad
 TT^t=T^tT=(N-1)I_4.
\]
Existence of the four integers is Lagrange's four-square
theorem.  Give $B^\vee$ the conjugate-dual $E$-action
$u\mapsto\iota_B(\bar u)^\vee$, and put
\[
 X=B^4\times(B^\vee)^4,\qquad
 \Lambda=
 \begin{pmatrix}\lambda I_4&T^t\\T&\lambda'I_4\end{pmatrix}
       :X\longrightarrow X^\vee.
 \tag{A11i}
\]
This is a principal polarization with Rosati involution $c$
on $E$.  We verify every assertion.  The matrix is symmetric,
and its inverse is the actual homomorphism
\[
 \Lambda^{-1}=
 \begin{pmatrix}\lambda'I_4&-T^t\\-T&\lambda I_4\end{pmatrix},
 \tag{A11j}
\]
as multiplication gives $N-(N-1)=1$ on the diagonal and zero
off the diagonal.  Positivity follows from the positive first
block and its Schur complement
\[
       \lambda'I_4-T(\lambda^{-1}I_4)T^t
                         =\lambda^{-1}I_4>0.
\]
Thus the symmetric map is a polarization, and (A11j) makes it
principal.  The equation
$\iota_B(u)^\vee\lambda=\lambda\iota_B(\bar u)$,
the analogous equation for $\lambda'$, and the integral scalar
entries of $T$ verify block by block that
\[
       \iota_X(u)^\vee\Lambda=\Lambda\iota_X(\bar u).
\]
Finally, $\lambda$ is $E$-linear from $B$ to the
conjugate-dual $B^\vee$.  The latter consequently has the same
full $E$-Hodge constituents and the same Lie multiplicities.
The uniform multiplicities of $X$ are $(8P,8Q)$.

Let $P_\Lambda$ be the normalized Poincar\'e biextension on
$X\times X$ classified by $\Lambda$, and define
\[
                 L=(\mathrm{id},[w_0])^*P_\Lambda.
 \tag{A11k}
\]
The polarization of this rigidified line is
\[
  \phi_L=\Lambda[w_0]+[w_0]^\vee\Lambda
          =\Lambda[w_0+\bar w_0]=\Lambda.
\]
Hence $L$ is ample and principal.  Pullback by $[-1]$ on both
arguments of a biextension leaves it unchanged, so $L$ is
symmetric.  Its normalized biextension $P_L$ is $P_\Lambda$,
which proves (A11a) as an isomorphism of actual rigidified
lines.  The theorem of the cube supplies its cubical structure,
and the adjoint relation follows from the displayed Rosati
identity.  These are the Poincar\'e/line-bundle constructions of
SGA~7 I, VII, Corollary 2.9.4 and Example 2.9.5, printed
pp.~174--175, and VIII, 3.1--3.2, printed pp.~254--255.

Take $\mathcal L=H_1(X,\Z)$ and its principal alternating
Riemann form, with the convention $EF(Jx,x)>0$ used in the
source.  The lattice is finite and torsion free over the
Dedekind domain $\OO_E$, hence projective.  Principalness gives
unimodularity.  The Rosati equation gives the required adjoint
identity.  Multiplication by $i$ gives $J$.  If
\[
 t=\sum_{\tau\in\Psi}\bigl(8P\tau(u)+8Q\overline{\tau(u)}\bigr),
\]
then the real traces are $2\operatorname{Re}t$ and
$-2\operatorname{Im}t$ for $u$ and $Ju$, exactly as in
\texttt{Datum.dom}.  This constructs the actual datum.

The dimension cost is explicit.  If $g=\dim A$ before the
field extension, then
\[
       \dim X=8s\frac{r+R}{r}\,g\leq64g,
 \tag{A11l}
\]
using $s\leq4$ and $R\leq r$.  An optional Galois closure of
a non-Galois original field replaces the bound on $s$ by its
actual extension degree.  The generic unitary moduli dimension
for the constructed datum is
$[E:\Q](8P)(8Q)/2$; this is the exact \texttt{Datum.dim}
in \texttt{PushInvertible.lean}.

\paragraph{Levels and an actual complex/algebraic object.}
There are arbitrarily large admissible odd levels avoiding any
one prescribed prime $p$.  Choose $\ell=11$ or $23$, with
$\ell\ne p$, and a prime $\mathfrak l$ of $K$ above $\ell$.
Both primes split in $K$: $X^2-X+2$ has roots $5,7$ modulo
$11$ and $10,14$ modulo $23$.  For $t\geq1$, put
\[
       M=\ell^t,\qquad
       \mathfrak M=\mathfrak l^t\OO_E,\qquad A_{\rm aux}=1.
 \tag{A11m}
\]
Then $\mathfrak M\overline{\mathfrak M}=(M)$ and
$\mathfrak M+\overline{\mathfrak M}=\OO_E$.  Also
$\mathfrak M$ is generated by its elements in
$\Z[w_0]=\OO_K$, proving precisely \texttt{Level.fromK}.
All integers are invertible over the base $\overline\Q$,
and $\rho=\sqrt{-7}$ and $\zeta=1$ supply \texttt{Base}.

For clarity, the remaining level data can be constructed
without a moduli existence assertion.  In
$X(\C)=\operatorname{Lie}(X)/\mathcal L$, set
$\mathrm{ev}_n(v)=v/n\pmod{\mathcal L}$, with a compatible
choice of roots of unity matching the Riemann-form convention.
These maps give \texttt{TateIso}: their kernel is
$n\mathcal L$, their image is all $n$-torsion, they commute
with the $\OO_E$-action, and they have the required tower and
Weil-pairing identities.  The Chinese remainder decomposition
\[
 \mathcal L/M\mathcal L
       =(\mathfrak M\mathcal L/M\mathcal L)
        \oplus(\overline{\mathfrak M}\mathcal L/M\mathcal L)
\]
defines $\eta$ by projecting to the second summand and then
applying $\mathrm{ev}_M$.  It kills
$\mathfrak M\mathcal L$ and equals $\mathrm{ev}_M$ on
$\overline{\mathfrak M}\mathcal L$, as required.  Take
$\eta_A=0$ when $A_{\rm aux}=1$.  The finite level points are
algebraic, so the object and its level structure are defined
over a finite extension of a number field.  The auxiliary
compatible torsion identifications for all $n$ are required
only over algebraically closed geometric fibres; their
simultaneous descent to that number field is not asserted.
Here is explicitly why the determinant condition survives
every coefficient embedding.  Enlarge the number field of
definition to contain $\rho=\sqrt{-7}$ and all the embedding
values of $E$.  For a field embedding $\beta$ into an
algebraically closed field of characteristic zero, extend
$\beta$ to an embedding of $\overline\Q$.  A Lie eigenspace
labelled by $\tau$ then becomes the eigenspace labelled by
$\beta\circ\tau$.  Since
\[
 (\beta\circ\tau)(\mathrm{rt})=\beta(\rho)
                  \quad\Longleftrightarrow\quad
                  \tau(\mathrm{rt})=\rho,
\]
its multiplicity is still $8P$ in this case and $8Q$
otherwise.  This is exactly \texttt{DetCond} with the
transported root $\beta(\rho)$; the root is not held fixed
while conjugating the object.  Transport every auxiliary
$\mathrm{ev}_n$ and its compatible root of unity through
the same extension of $\beta$.  The integral Riemann-form
entries, tower relations, and algebraic Weil pairings are
preserved.  The finite level maps already descend to the
number field, so the transported identifications agree with
their prescribed base changes.  This proves the entire
\texttt{geom} condition.  Principalness gives
\texttt{Out.KTriv}: a point whose translation preserves $L$
lies in the trivial kernel of $\phi_L$.  Thus the constructed
abelian variety, line, action, and level form an actual
\texttt{Obj}, over $\overline\Q$ and after finite descent.

\paragraph{Preserving a finite family of first-jet values.}
Changing the polarization in (A11i) changes its theta line.
The following precise interpolation argument accounts for
that change.  Let $U$ be the finite-dimensional algebraic input
space containing the finite class-sum inputs under consideration.
Transport each prescribed algebraic tangent direction through
(A11f), the product inclusion in (A11h), and one of the direct
factors of $X$.  This gives the selected tangent subspace
$V_0\subset\operatorname{Lie}(X)$.  Retain the algebraic
values and first derivatives of the old sections on $V_0$.
Choose an algebraic complement of $V_0$ and extend the desired
linear functionals by zero on that complement.

Let $Z_1$ be the first infinitesimal neighborhood of zero,
defined by the square of its ideal.  A normalized symmetric
origin frame identifies
\[
       H^0(Z_1,L^{M^2}|_{Z_1})
                    =\overline\Q\oplus\operatorname{Lie}(X)^*.
\]
The prescription is therefore one linear map from $U$ to this
finite-dimensional jet space.  For all sufficiently large
$M$ in the sequence (A11m), the map
\[
 H^0(X,L^{M^2})\longrightarrow H^0(Z_1,L^{M^2}|_{Z_1})
 \tag{A11n}
\]
is surjective.  Indeed, tensor the ideal sequence of $Z_1$
with $L^{M^2}$ and use Serre vanishing for its coherent ideal
(Stacks, Lemma 30.16.2, tag \texttt{02O1}).  Choose one linear
right inverse to (A11n) and compose with the prescribed jet
map.  This produces a single linear family of actual algebraic
theta sections realizing every prescribed value and first
derivative.  Consequently it preserves every specified
determinant of first derivatives, including all the finite
linear combinations and class sums in B19 and A10.

The symmetric frame used here has no linear change relative
to the canonical completed frame.  Explicitly, the ratio of
two normalized symmetric frames has value one and is even,
so its differential at zero vanishes.  The quadratic canonical
completion also has no linear term.  Thus the interpolated
first derivatives are the same first derivatives needed for
the canonical finite theta values; higher jets would require
the corresponding higher-order comparison.

After choosing a finite theta structure, let
$J_{\rm new}:W_{\rm new}\to H^0(X,L^{M^2})$ be the algebraic
Heisenberg isomorphism and
\[
        \mathcal T_{\rm new}
               =c_{\rm new}(J_{\rm new}\otimes\C),
        \qquad \dim W_{\rm new}=M^{2\dim X}.
 \tag{A11o}
\]
Applying $J_{\rm new}^{-1}$ to the preceding linear family
gives actual finite theta inputs.  Their analytic first jets
are $c_{\rm new}$ times the prescribed algebraic jets.  The
old analytic jet factors $c_i p_i/N_i$, or the stabilized
factors of A10, remain explicit external comparison scalars.
No assertion that their ratios with $c_{\rm new}$ are
algebraic is needed for this realization.  The right inverse
may have a substantial denominator or norm; it is a recorded
cost and not a uniform quantitative estimate.

There is an integral version at any chosen odd good prime
where the actions, polarization, specified eigenprojections,
and tangent lattice have been extended integrally.  Over a
complete discrete valuation ring, $Z_1$ is finite flat and its
jet module is free.  Relative Serre vanishing makes (A11n)
surjective for sufficiently large $M$ prime to $p$.  A map
onto a finite free module has a linear right inverse, so an
integral prescribed jet family has integral lifts.  Symmetric
references with value one are obtained by averaging with
$[-1]$ and dividing by $2$; their first derivatives vanish.
Prime-to-$p$ level structures can be acquired after a finite
extension.  If an old tangent lattice was not saturated, its
denominator must be retained before applying this assertion.
This argument does not assert an integral identification at a
prime where these hypotheses fail.

\paragraph{What the construction settles.}
This proves an actual generic finite CM host realization,
with a field extension of degree at most four, minimal
rank-one multiplicity padding, an explicit factor-eight
principalization, and a finite theta level realizing the
required first-jet functional.  The added constituents keep
their full CM types.  Equal local multiplicities alone have
not been used to identify two $E$-Hodge structures or two
polarized lattices.  The construction does not identify the
new sections with a prescribed global automorphic family,
preserve its weights or signatures, or supply the separate
definite Gram-one \texttt{StdPair} of the recovered global
chain.  Those comparisons concern more than a finite CM
evaluation at one point.

There is also a precise distinction between two source records
named \texttt{Host}.  A11 constructs
\nolinkurl{ReciprocityTheory.ChartRestriction.G.Geo.Tor.Host}.
The global \texttt{StdPair} is over a \texttt{D7} whose host is
\nolinkurl{Formalism.BasicNotions.Host}, defined in
\path{Formalism/LocalRingCardinality/DataHost.lean}.
The latter extends \texttt{HC.HostAx} and additionally
specifies a Galois/Frobenius frame, auxiliary primes and
size conditions, and a factorization of a relative different.
Neither (A11b) nor multiplicity padding is asserted to
construct those additional records.

The construction uses actual homomorphisms, polarizations,
Poincar\'e biextensions, and sections.  In the requested
Grothendieck framework, the correspondence and polarization
viewpoint is the one recorded in \emph{cote44}, scans 7--8;
the dimension and Lefschetz discussion is \emph{cote16}, scans
43--46.  No unproved Chow projector is used here.  The
compatible principalization and finite-jet arguments above
are proved here in full; the explicitly named representability, CM descent,
cube, and cohomology theorems are cited background inputs.


\subsection{THETA3-A12: a constructive standard finite theta comparison}
\label{A12}

There is a stronger finite comparison than interpolation on an
unspecified ample line.  On an actual positive standard oscillator
torus, explicit finite Heisenberg operators produce a basis of first
jets, with a quantitative inverse.  Each constructed vector has an
actual \texttt{IsRatDatum} witness.  The resulting rational coefficient
array is specified by the construction.  Identifying that array with
the \texttt{xiAt} of a previously fixed global tuple is a separate,
finite equality, described below.  The analytic-to-algebraic scalar
and the tangent period remain in every comparison.
Throughout this section the prescribed first jets mean
their derivative part, of dimension $g$.  The value at
the origin is not prescribed by the inverse matrix below.

\paragraph{Exact source conventions and scope.}
Fix an actual archive \texttt{StdPair} with ranks $a,b$ and
$0<a<b$, both even, and put
\[
                  g=ab[F:\Q].
\]
Its Gram matrices are both the identity.  The definition in
\path{Development/AbelianVarietyTheory/Definitions/}
\path{LatticeFrames/SplitPart.lean}, \texttt{StdPair}, requires
the standard additive character with parameter $-1$,
$\xi_0=i/\sqrt d$ at each selected embedding, and
$\mathrm{sc}|\delta|^2=2/\sqrt d$.  The complex structure
\texttt{cplxO 1 1} is multiplication by $i$.  Its holomorphic
reading is the literal \texttt{xiN} in that file.  We retain a
fixed compatible self-dual standard lattice and a normalized
vacuum $f_0$ satisfying \texttt{IsVacO}.  A rational Schr\"odinger
polarization and its theta reading are fixed once.

For the algebraic conclusion we choose the finite vacuum in the
algebraic finite theta module, up to its one overall normalization
scalar.  This is a choice that can be made, not a property inferred
from the bare assertion \texttt{IsFinComb}.  Indeed, at a level
containing a nonzero vacuum, the lattice-vacuum conditions are
linear equations for finite Heisenberg operators.  Those operators
have algebraic matrix entries, and their action factors through a
finite quotient.  The solution subspace is therefore defined over
$\overline\Q$.  A nonzero complex solution gives a nonzero
algebraic solution by Gaussian elimination.  Its archimedean
Gaussian is unchanged; rescaling its positive $L^2$ norm gives
the required normalization.  The rescaling is kept as an external
scalar.  If a previously specified vacuum is used instead, the
assertion that its finite vector is algebraic up to one scalar
must be supplied for that vacuum.

The finite model conventions are particularly useful here.  Write
$h(u)$ for \texttt{heisO}, $B$ for \texttt{symplO}, and
$B_{1/2}=B/2$ for \texttt{symplOH}.  If $e$ is the sum of
selected split-place idempotents, $X=eW_f$, $Y=\bar eW_f$, and
$L=e\Lambda$, then a normalized place model $T$ obeys
\[
 \begin{split}
 T(\Phi(\,\cdot+x))&=h(x)T\Phi,\qquad x\in X,\\
 T\bigl(\psi(B(\,\cdot,y))\Phi\bigr)&=h(y)T\Phi,
                                      \qquad y\in Y,\\
 T1_L&=f_0,\qquad
 h(x)h(y)=\psi(B_{1/2}(x,y))h(x+y).
 \end{split}
 \tag{A12a}
\]
The first three equations are the literal
\path{IsPlaceModel.trans/modul/vac} in
\path{Formalism/LocalToGlobal/IndefiniteReal/}
\nolinkurl{LocalNewformFunction.lean}.  The last equation follows
directly from its definition of \texttt{heisO}.  The predicate
\texttt{IsThetaAt} means both existence of a normalized model
and equality with its value for every such model.

Existence for the fixed self-dual lattice and vacuum is the
theorem \nolinkurl{GlobalTheory.FunObj.VecEx.exists_isThetaAt_of_vac}
in
\path{Development/AbelianVarietyTheory/Proofs/DivisorProofs/}
\nolinkurl{ExistsIsThetaAtOfVac.lean}.  The proof of the precise
finite-operator evaluation used below is
\nolinkurl{GlobalTheory.HalfFourier.thetaAt_eq_sum_heis_translates}
in
\path{Development/AbelianVarietyTheory/Proofs/}
\path{DecoratedLatticeProofs/}
\nolinkurl{ThetaFourierSumEqSumHeisenbergTranslates.lean}.
These are source locators and a check of the displayed
mathematics, not a claim that the recovered tree was compiled.

\paragraph{An explicit half-space input for every finite operator sum.}
Suppose $\eta_\nu$ are finitely many vectors supported at these
split places, and $c_\nu\in\overline\Q$.  Decompose
\[
 x_\nu=e\eta_\nu,\qquad y_\nu=\bar e\eta_\nu,
 \qquad \eta_\nu=x_\nu+y_\nu.
\]
The last equality uses that $e+\bar e$ is the identity on
the selected support; it need not be the identity at other
finite places.  Define
\[
 \Phi(t)=\sum_\nu c_\nu\psi(-B_{1/2}(x_\nu,y_\nu))
       \psi(B(t+x_\nu,y_\nu))1_L(t+x_\nu),
 \qquad
 f=\sum_\nu c_\nu h(\eta_\nu)f_0.
 \tag{A12b}
\]
Choose an integer $k>0$ clearing all translation denominators
and all character conductors in this finite list.  Then
$\Phi$ is supported in $k^{-1}L$, invariant under $kL$, and
zero off $X$.  Thus it satisfies \texttt{IsSBHalf}.  For
every normalized $T$, apply the modulation rule first and
then the translation rule in (A12a).  Its value on the
$\nu$th summand is
\[
 c_\nu\psi(-B_{1/2}(x_\nu,y_\nu))h(x_\nu)h(y_\nu)f_0
                         =c_\nu h(\eta_\nu)f_0.
\]
Linearity proves $T\Phi=f$ for every model.  The existence
theorem just cited proves the existential part as well.
Consequently this is an actual \texttt{IsThetaAt} witness.
Every finite character value occurring here is a root of
unity: on the finite support quotient, some positive integer
annihilates its argument modulo the character kernel.
Therefore $\Phi$ is algebraic-valued.

Define, on the rational matrix space used by the source,
\[
                         \Xi(\gamma)=\Phi(e\gamma).
 \tag{A12c}
\]
Equations (A12b)--(A12c) are exactly the three conditions in
\texttt{IsRatDatum}, in
\path{Development/AbelianVarietyTheory/Definitions/}
\path{LatticeFrames/IsRatDatum.lean}.  They also give
\texttt{IsFinComb} explicitly.  No irreducibility argument
is needed to leave an unspecified combination or scalar.

\paragraph{A rational Schr\"odinger lattice, not a sum over the full CM lattice.}
Use the fixed rational Lagrangian of the underlying
$2g$-dimensional rational symplectic space.  Its standard
theta reading, after the fixed overall scalar has been
separated, has the form
\[
 \vartheta_\phi(Z,z)=
   \sum_{x\in\Q^g}\phi(x)
          \exp(\pi i x^tZx+2\pi i x^tz),
 \qquad Z=X_0+iY_0,\quad Y_0>0.
 \tag{A12d}
\]
Here $\phi$ is a finite Schwartz function on $\A_f^g$.
Rational theta characteristics and finite bad-place factors
are included in $\phi$.  The sum is over a rational
Schr\"odinger Lagrangian of real dimension $g$, not over
the underlying $2g$-dimensional CM lattice.  This distinction
prevents a spurious dimension or Gaussian normalization.
The fixed canonical completion has value one and first
derivative zero at the origin, so it leaves these first
jets unchanged.  The choice of rational symplectic chart
fixes the sign in (A12d) and its positive $Y_0$; its
comparison with the archive reading is retained below.

A rational symplectic basis need not preserve the original
integral polarization lattice.  If necessary, use its
commensurable unimodular lattice $\Gamma_0$ for the
principally polarized torus underlying (A12d).  This is
an actual isogeny comparison, with explicit indices:
if $\Gamma$ is the original lattice and
$\Gamma_1=\Gamma\cap\Gamma_0$, the torus with lattice
$\Gamma_1$ maps to the two tori with degrees
$[\Gamma:\Gamma_1]$ and $[\Gamma_0:\Gamma_1]$.
The positive complex structure and rational symplectic
space are the same.  Hence both are polarized CM tori,
their rational $E$-Hodge actions agree, and these isogenies
are algebraic over $\overline\Q$.  The new unimodular
lattice is not asserted to be $\OO_E$-stable.  Its finite
theta structure and the isogeny tangent matrices are
recorded, with denominators included in the levels below.
The source vacuum, source lattice, and operators in
(A12a)--(A12c) remain the original ones; only their
Schr\"odinger coordinates and algebraic theta comparison
use this isogenous principal model.  This does not
identify Hermitian self-duality with unimodularity of
the rational trace pairing without checking its different.

Let $\phi_0$ be the nonzero finite coefficient of the chosen
vacuum, with its overall scalar removed.  Write this scalar
as $u_0\ne0$: the actual source reading of $f_0$ is
$u_0\vartheta_{\phi_0}$ in the fixed chart, with the
canonical completion understood.  It follows by equivariance
that the source vector constructed from $\phi_j$ below has
reading $u_0\vartheta_{\phi_j}$.  We do not change the
normalized vacuum when taking its finite Heisenberg translates.
By compact support
and local constancy, an integer $M\geq1$ can be chosen such
that
\[
 \operatorname{supp}\phi_0\subset M^{-1}\widehat\Z^g,
 \qquad \phi_0(x+M\widehat\Z^g)=\phi_0(x).
 \tag{A12e}
\]
There is a rational $x_0\in M^{-1}\Z^g$ with
$a_0:=\phi_0(x_0)\ne0$.  To see this, the support contains
a coset modulo $M\widehat\Z^g$ with nonzero value, and
the Chinese remainder theorem represents that coset by
a rational vector.  Put $B_0=\max|\phi_0|$, which is the
maximum of a finite list.  The values $a_0,B_0$ and the
support and period data are part of the certificate.

Choose a positive integer $r$ such that the $g$ columns
\[
                    x_j=x_0+Mr e_j\quad(1\leq j\leq g)
 \tag{A12f}
\]
are linearly independent.  Such an $r$ is explicit: their
determinant is
$(Mr)^{g-1}(Mr+\sum_i(x_0)_i)$, so only one positive
value could fail.  Each $x_j$ is nonzero and has coefficient
$\phi_0(x_j)=a_0$.

Fix an odd rational prime $\ell$ that splits completely in
$E$, avoids $M$ and all lattice and character bad primes,
and at which the fixed rational symplectic chart is integral.
Such primes exist by Chebotarev (Milne, \emph{Algebraic
Number Theory}, 2020, Theorem 8.31 and Corollary 8.32,
printed p.~147, applied to a Galois closure); no effective bound for the
first such prime is used.  Once one is specified, all
subsequent conductor and tail bounds are explicit.  Select
one of the two $E/F$ halves at every place above $\ell$;
their sum is the allowed split idempotent $e$.  The local
Schr\"odinger vacuum is $1_{\Z_\ell^g}$, up to the already
removed scalar.  This also follows directly from lattice
invariance: modulation by the dual integral lattice forces
support in $\Z_\ell^g$, and translation by that lattice
forces constancy there.  All these choices are at good
places, so no bad-place characteristic is discarded.

For $t\geq1$, let $k=\ell^t$ and define
\[
 \phi_j(x)=\phi_0(x)
           1_{\,x_j+k\Z_\ell^g}(x_\ell).
 \tag{A12g}
\]
This operation is an exact finite Heisenberg projector.
If $C=\Z_\ell^g/k\Z_\ell^g$ and $\widehat C$ its
character group, character orthogonality gives
\[
 1_{\,x_j+k\Z_\ell^g}(x)
       =\frac1{k^g}\sum_{\chi\in\widehat C}
                         \chi(-x_j)\chi(x),
                         \qquad x\in\Z_\ell^g.
 \tag{A12h}
\]
Every $\chi(x)$ is a modulation by a vector $\eta_\chi$
in the complementary rational Schr\"odinger Lagrangian,
supported at $\ell$.  Thus the vector with coefficient
$\phi_j$ is exactly
$k^{-g}\sum_\chi\chi(-x_j)h(\eta_\chi)f_0$.
There are $k^g$ terms, and their coefficient absolute
values sum to one.  Apply (A12b) to this very list.
This converts the Schr\"odinger projector into the
archive's $e,\bar e$ model, with its precise half-pairing
phase, and gives an actual \texttt{IsRatDatum} for every
$j$.  At our good places $x_\chi,y_\chi\in k^{-1}\Lambda$,
so the same $k$ supplies its support and period bounds.
This avoids an unspecified change-of-polarization scalar.

\paragraph{A terminating quantitative first-jet certificate.}
Fix a certified $\lambda>0$ with $Y_0\geq\lambda I_g$,
put $R=\max_j\|x_j\|$, and let $h=k/M$.
Every rational point contributing to (A12g) differs from
$x_j$ by an element of $h\Z^g$.  Indeed, both points
lie in $M^{-1}\Z^g$, and $M(x-x_j)$ is divisible by
$k$ in $\Z_\ell^g$, hence in $\Z^g$.  In particular,
when $h>2R$, the point $-x_j$ is excluded and $x_j$
is the unique point of this coarse coset with norm at
most $R$.  There is no unrecorded $\pm x_j$ cancellation.

Let $J_t$ be the matrix whose $j$th column is
$\nabla_z\vartheta_{\phi_j}(Z,0)$.  Its leading matrix
$A_0$ has columns
\[
         (A_0)_j=2\pi i a_0x_j e^{\pi i x_j^tZx_j}.
 \tag{A12i}
\]
It is invertible by (A12f).  Choose any certified
$\sigma_0>0$ with $\sigma_{\min}(A_0)\geq\sigma_0$.
For example, a determinant lower bound and column-norm
upper bounds give such a number by the product formula
for singular values.  For $h>2R$ the error in each column
is at most
\[
 E(h)=2\pi B_0\sum_{m\geq1}(2m+3)^g
        ((m+1)h+R)e^{-\pi\lambda(mh-R)^2}.
 \tag{A12j}
\]
Here is the full estimate.  Group the nonzero lattice
differences $u\in h\Z^g$ by
$mh\leq\|u\|<(m+1)h$.  Disjoint balls of radius
$h/2$ around them fit into the ball of radius
$(m+3/2)h$, so there are at most $(2m+3)^g$ such
points.  For $x=x_j+u$ one has
\[
 mh-R\leq\|x\|\leq(m+1)h+R,
 \quad |e^{\pi ix^tZx}|\leq e^{-\pi\lambda\|x\|^2},
 \quad |\phi_0(x)|\leq B_0.
\]
Multiplying these bounds by the derivative factor
$2\pi\|x\|$ proves (A12j).  The argument allows zero
or varying coefficients on the other bad-place residue
classes, and therefore does not replace $\phi_0$ by
a single characteristic function.

The bound tends to zero and has a finite tail certificate.
For $h\geq2R$ put $a=\pi\lambda h^2/4$.  Using
$(2m+3)^g\leq5^gm^g$ and
$((m+1)h+R)\leq(5/2)mh$ for $m\geq1$ gives
\[
 E(h)\leq5^{g+1}\pi B_0h
                   \sum_{m\geq1}m^{g+1}e^{-am^2}.
\]
For any integer $L\geq0$ with
$q_L=2^{g+1}e^{-a(2L+3)}<1$, the part of the last
sum with $m>L$ is at most
\[
          \frac{(L+1)^{g+1}e^{-a(L+1)^2}}{1-q_L}.
\]
Indeed the ratio of consecutive terms for $m\geq L+1$
is at most $q_L$.  This bounds the tail of (A12j)
after any computed finite initial segment by multiplying
this last expression by $5^{g+1}\pi B_0h$.
It also proves $E(h)\longrightarrow0$, for example by
taking $L=0$ once $a$ is large.  Increase $t$ until
\[
                  h>2R,\qquad
                  \sqrt g E(h)\leq\sigma_0/2.
 \tag{A12k}
\]
The procedure terminates.  The Frobenius bound gives
$\|J_t-A_0\|_{\rm op}\leq\sqrt g E(h)$.
For each unit vector $v$,
\[
 \|J_tv\|\geq\|A_0v\|-\|(J_t-A_0)v\|
                       \geq\sigma_0/2.
\]
Consequently
\[
       \sigma_{\min}(J_t)\geq\sigma_0/2,
                   \qquad \|J_t^{-1}\|_{\rm op}\leq2/\sigma_0.
 \tag{A12l}
\]
Alternatively one can keep $t=1$ and choose the good
split prime $\ell$ sufficiently large.  The quantities
$M,R,B_0,A_0,\lambda,\sigma_0$ were fixed before this
choice, whereas $h=\ell/M$ tends to infinity through
the permitted primes.  The same finite certificate
therefore proves (A12k)--(A12l) at residue level
$\ell$, without a high power of $\ell$.  Its additive
projector has coefficients in the cyclotomic field
generated by $\ell$th roots of unity; if one wants
all coefficients in a specified number field, that
field must first contain these roots, or be enlarged
once after $\ell$ is chosen.  This observation by
itself does not replace the archive's multiplicative
tuple characters by additive projector characters.
This is a quantitative spanning theorem for all $g$
first derivatives, with the finite input and its
\texttt{IsRatDatum} supplied explicitly.

\paragraph{Removing the common periods before taking an algebraic inverse.}
Take a common even level $N$ containing the finitely many
inputs just constructed; for example an even multiple
of $Mk$ suffices in the fixed Schr\"odinger coordinates.
Choose the actual algebraic finite theta isomorphism
$\mathcal J_N$ for the standard CM torus, or for the
just specified isogenous principal model.  Its analytic
comparison is one scalar $c_N$ on the whole finite
module, as in B19 and B21.  Here $c_N$ refers to the
normalized series (A12d); the actual source theta map
on these vectors has comparison scalar $u_0c_N$.

There is one common tangent period on this standard
torus.  Its $E$-type consists of all embeddings restricting
to the chosen embedding of $K=\Q(\sqrt{-d})$.  As a
rational $K$-Hodge structure it is therefore a sum of
copies of the same one-dimensional CM type.  A rational
$K$-basis identifies it with the homology of a power
of a CM elliptic curve, and clearing denominators
gives an isogeny.  The isogeny and its tangent maps are
algebraic over $\overline\Q$ by CM descent and descent
of homomorphisms (Milne, \emph{Complex Multiplication},
2020, Proposition 7.9 and Corollary 7.10, printed
pp.~54--55).  Thus, in a chosen algebraic tangent
frame, the coordinate vectors of the analytic rows
are $p_{\rm st}B$, where $p_{\rm st}\ne0$ is one
elliptic tangent period and
$B\in\operatorname{GL}_g(\overline\Q)$ is fixed.
This argument uses the full induced type, not only
equality of numerical Lie multiplicities.

Let $A_t$ be the $g$ by $g$ matrix of algebraic first
jets of $\mathcal J_N\phi_j$ in that frame.  Recall
that the standard finite theta coordinate is $z=Nw$.
The chain rule and the symmetric origin frame give
the exact matrix identity
\[
       J_t=\gamma_N B A_t,
             \qquad \gamma_N=\frac{c_Np_{\rm st}}{N}.
 \tag{A12m}
\]
The actual source first-jet matrix is $u_0J_t$.
All entries of $A_t$ are algebraic.  Its invertibility
follows from (A12l), after the nonzero scalar and
invertible row matrix have been retained.  For any
prescribed algebraic jet matrix $C$ with $g$ rows,
put
\[
         L_t=A_t^{-1}C\in M(\overline\Q),
 \qquad
 \|L_t\|_{\rm op}\leq
       \frac{2|\gamma_N|\,\|B\|_{\rm op}}{\sigma_0}
                         \|C\|_{\rm op}.
 \tag{A12n}
\]
Indeed $A_t^{-1}=\gamma_NJ_t^{-1}B$.  The linear
combinations $\sum_j(L_t)_{jv}\phi_j$ have exactly
the algebraic jet columns $C_v$.  The corresponding
linear combinations of (A12b) give their actual
\texttt{IsRatDatum} witnesses.  A column's finite
Heisenberg coefficient cost is at most
$\sum_j|(L_t)_{jv}|\leq\sqrt g\|(L_t)_v\|_2$,
because each congruence projector had cost one.

The factor $|\gamma_N|$ is not being inferred to be
algebraic.  B22 gives an optional explicit calibration
by a nonzero odd reference first jet, with Gaussian
lower and upper bounds, after a specified algebraic
theta and tangent normalization.  Its change-of-frame
matrix must then be included in $B$.  None of these
analytic bounds bounds the heights or denominators
of $A_t^{-1}$, or proves integral saturation at a
prescribed prime.

In the exact archive polynomial reading, each first
derivative also has the fixed scalar
$s=c_K(-\delta)$, where
$c_K=(\sqrt{2\pi/\sqrt d})^{-1}$; a degree-$r$
determinant has the factor $s^r$.  This is the
explicit witness in
\path{Development/AbelianVarietyTheory/Proofs/}
\path{DecoratedLatticeProofs/PureTensorOZeroEqJetP.lean},
\nolinkurl{pureTensorO_zero_eq_jetP}.
The real reading of finite rational Heisenberg operators
is proved in the same directory's
\nolinkurl{GaussianReadingFiniteRationalHeisenbergTranslate.lean},
\texttt{xiN\_heisO\_finOf}; the classical-theta identification
with the specified rational array is
\nolinkurl{DatumReadingSplitPlacesEqClassicalTheta.lean},
\texttt{reads\_datum\_eq\_thC}.
Thus the scalars in (A12m) and this polynomial scalar
are kept at their different stages rather than
silently set equal to one.
Here the source's $s$ multiplies derivatives in its
literal $\xi_N$ coordinates.  If $C_{\rm arch}$ is the
Jacobian taking those derivative rows to the chosen
rational Schr\"odinger chart, the source operator is
$s^rP(C_{\rm arch}\partial_z)$, not automatically
$s^rP(\partial_z)$.  All row changes caused by the
commensurable-lattice isogeny must be retained as well.
One may absorb this Jacobian into the matrix $B$ only
after establishing the stated algebraic change of
algebraic tangent frames and the common period factor;
otherwise $C_{\rm arch}$ remains a separate, explicit
comparison matrix in the following scalar rules.

There are two distinct scalar rules.  A determinant of
$r$ independent first-jet columns has comparison factor
$(u_0c_Np_{\rm st}/N)^r$, together with the appropriate
minor of $B$; if one first creation operator is used in
each factor, it also has $s^r$.  By contrast, a single
homogeneous differential operator of degree $r$ applied
to one theta section has the one section-comparison
factor $u_0c_N$, and tangent factor
$(p_{\rm st}/N)^r$, with the induced degree-$r$ action
of $B$.  Its polynomial oscillator normalization is
again $s^r$.  The latter higher-jet statement requires
the matching canonical connection or harmonic frame
comparison of A8; it is not inferred from the absence
of a linear term in the origin-frame ratio.

\paragraph{Transfer of a finite family and the exact remaining array test.}
Apply (A12n) to the finite list of algebraic first-jet
columns obtained from B19, A10, or A11, extending their
rows into $g$ rows if necessary.  This application requires
$g$ to be at least the number of prescribed tangent rows;
it does not assert that an arbitrary fixed \texttt{StdPair}
has sufficient dimension.  This gives one linear
map on that entire finite input space and preserves
all its specified determinants, finite sums and
linear relations at the algebraic-jet level.  Original
032 rows still carry their factors $c_ip_i/N_i$;
stabilized rows still carry the exact missing-reference
theta-null factors in A10.  The standard replacement
carries $u_0\gamma_NB$ and, if the polynomial oscillator
coefficient is used, the displayed power of $s$.
Their equality as unnormalized analytic coefficients
has not been inferred from algebraic interpolation.

There is a finite, exact test for a proposed global
tuple to use these very inputs.  Suppose its prescribed
\texttt{xiAt} descends to an algebraic locally constant
function on the same half-space, and a common integer
$k'$ is provided for its support in $(k')^{-1}L$
and invariance under $k'L$.  Enlarge $k'$ to contain
our support and period bounds.  Then equality with
(A12c) is equivalent to equality on
\[
                        (k')^{-1}L/k'L.
 \tag{A12o}
\]
For $k'=\ell^{t'}$ at the $g$-dimensional good
$\ell$-adic half-space this has exactly
$\ell^{2t'g}$ entries.  Rational matrices attain all
its classes by weak approximation, so equality of
the rational values is precisely this finite-array
equality.  If several places or additional periods
are used, the size is the corresponding product of
local indices.  One must supply the descent and
support/period bounds of the prescribed function;
they do not follow merely from its name.

This is the literal distinction in
\texttt{IsRatDatum.lean}: \texttt{ColDat.Is.datum}
uses \texttt{A.xi}, and \texttt{A.xi} is defined
from the already chosen \texttt{xiAt} of the global
tuple.  The construction proves \texttt{IsRatDatum}
for the new explicit array, and reduces an identification
with a specified tuple to (A12o).  It does not replace
that tuple by fiat.  The global Weil characters,
signatures, source levels, and the other \texttt{ColDat}
or \texttt{MainX} conditions remain their stated
conditions.

\paragraph{Why the finite transfer is not a hidden CM isogeny.}
There is also a precise geometric limitation.  Let
$U(\mathfrak d)$ be a rank-one $E$-Hodge structure
whose complete type differs from the induced type
$\Psi$ of the standard torus.  Then
\[
           \Hom_{E,\mathrm{Hdg}}
                  (U(\mathfrak d),U(\Psi)^r)=0.
 \tag{A12p}
\]
An $E$-linear map is a column of elements of $E$.
At an embedding where the two types differ, the Hodge
condition forces every entry of that column to have
zero image.  A field embedding is injective, so every
entry is zero.  This proves (A12p).  Adding reference
factors does not change this calculation for an
original constituent.  Thus matching finite first
jets by (A12n) is a genuine finite linear construction,
but supplies no geometric embedding of arbitrary
032 CM tori into this induced standard torus, and
no equality of global theta operators.

\paragraph{The Grothendieck framework.}
The finite theta comparison here produces actual
sections and algebraic linear maps on their first
jets.  It is compatible with the explicitly constructed
CM correspondences of A6 and A10 when their period
factors are retained.  It does not promote evaluation
at a CM point into a weight-two cohomological restriction.
That distinction is the correspondence-and-polarization
framework in the supplied \emph{cote 44}, scans 7--8,
and the inverse-correspondence discussion in
\emph{cote 16}, scans 43--46.  B23 supplies a different,
cohomological block inverse; neither inverse is silently
substituted for the other.


\subsection{THETA3-A13: the actual finite arrays, their spans, and tame globalization}
\label{A13}

This section gives restrictions and positive constructions for the
\emph{prescribed} column arrays.  It does not replace a prescribed
\texttt{xiAt} by an arbitrary Schwartz function.  The finite-field and
idele arguments below are new proofs.  Literal archive definitions and
the named archive existence theorem are source statements; no Lean
build or transitive dependency certification is asserted.

\paragraph{The defining objects and the coefficient field.}
All paths in this paragraph are relative to the pinned source root.
The literal global definition is
\path{Development/AbelianVarietyTheory/Definitions/AutomorphicMembership/}
\nolinkurl{XiAtPenultimateEqXiCpl1.lean}, lines 146--162.
It uses \texttt{XiEntry} when $i+2<m$, \texttt{XiCpl1}
when $i+2=m$, and \texttt{xiLast} when $i+1=m$.
The last is the displayed \texttt{xiOver}; all three carry
the same \texttt{phiT} factor.
The notation \texttt{TC} is an abbreviation for \texttt{ColIn},
in \texttt{LatAtLattice.lean}, line 40, in that same directory.
The latter structure is in
\nolinkurl{TrueColumnChoicesHeckeCharReferenceEqOne.lean},
lines 165--245: it has a finite injective family of tuples,
the reference tuple $1$, the specified reference characters,
local constants, local double-coset indices and their tying laws.
Thus a \texttt{TC} is more than a list of local characters.

The local definitions are in
\path{Formalism/LocalToGlobal/Columns/FiniteSumNonzeroLargeDepth.lean}:
\texttt{phiT}, lines 646--648;
\texttt{PhiQ}, lines 711--714;
\texttt{PhiG}/\texttt{PhiGS}, lines 780--805;
\texttt{om}, lines 553--556;
\texttt{xx1}/\texttt{E1}, lines 568--579;
\texttt{PhiAt}/\texttt{XiAt}, lines 927--948;
\texttt{XiEntry}/\texttt{XiCpl1}, lines 1116--1138.
The last local factor is \texttt{locTop}, lines 391--393 of
\path{Development/AbelianVarietyTheory/Definitions/GlobalTwisting/}
\nolinkurl{IntegralGLSuccEmbeddingMemParahoricTimesPrincipa.lean}.

Write $K$ for the coefficient field, to distinguish it from the
imaginary quadratic subfield of the CM host.  The literal
\texttt{Fields.K} is a number field
(\nolinkurl{AbelianHostFieldConditions.lean}, lines 19--44).
For the actual base it is $B.\mathrm{LL}$
(\nolinkurl{BaseDataEmbeddingConj.lean}, lines 13--20).
There is an additional, concrete restriction on a fixed $K$.
In the actual canonical case \texttt{IsTw}, the proof
\texttt{root\_of\_ok} in \texttt{IsTwIffCanonEq.lean}
shows that every new twisting prime of residue cardinality $q$
requires a primitive $q^{2m}$th root of unity in $K$;
\texttt{sqrt\_of\_ok} requires a square root of $q$.
The proof uses the stipulated local additive character at depth
$2m$, its tying to the global character, and the fact that the
reference $\xi_0$ is a unit at this prime.  A fixed number field
therefore admits only finitely many such rational primes.
A12's choice $t=1$ removes a large power from its isolating
congruence, but not this separate canonical-depth requirement.

\paragraph{Two unconditional restrictions on the actual arrays.}
For a valid field and prime datum, put $f=[F:\Q]$.
The source has
\[
        |v_T(u)|=f>0,\qquad |v_G(j)|=f>0.
        \tag{A13a}
\]
These are \texttt{Primes.Ok.vT\_card} and
\texttt{vG\_card}, lines 216 and 231 of
\nolinkurl{AbelianHostFieldConditions.lean}; positivity follows
from \texttt{Fields.Ok.card\_Tau} and $[F:\Q]>0$.
Every valid decoration has a nonzero residue functional at
each $T$-prime, by \texttt{Deco.Ok.f\_ne}, line 88 of
\texttt{StdLatUnimodular.lean}.
Consequently
\[
                         \mathrm{xiAt}(0)=0.
                         \tag{A13b}
\]
Indeed, \texttt{phiT} requires $f'_P\circ x=f_P$ on the
localized source lattice.  At $x=0$ its left side is zero:
$f'_P(0)=0$ follows from the displayed scalar-linearity
condition with scalar zero.  Its right side is nonzero on
some lattice vector.  Thus $\mathrm{phiT}(0)=0$ at any one
$T$-prime.  All branches of the global array contain this
factor, or are identically zero.

There is also the stronger rank restriction
\[
 \mathrm{xiAt}_i(x)\ne0
           \quad\Longrightarrow\quad
             \operatorname{rank}_E(x)=N(i),
       \qquad N(i)=k+i(t+1),\quad 0\leq i<m.
 \tag{A13c}
\]
Here the target dimension is $N(i+1)$.
To prove this directly, choose any one $G$-prime using
(A13a).  Set $c_i=m-1-i$ and $r=k+m-1$.
In \texttt{PhiG} and \texttt{PhiGS} at step $i$,
the nonzero reciprocal determinant factor is the determinant
of a square block of size
\[
               r+it=N(i)+c_i.
\]
Its reduction is a unit whenever the local function is nonzero.
The coordinate operation \texttt{entB} is an invertible row
and column permutation followed by nonzero diagonal scalings;
the displayed piecewise insertion \texttt{fIdx} has inverse
\texttt{fPos}.  Hence the underlying enlarged matrix has full
column rank $N(i)+c_i$.  The matrices in \texttt{om} are
invertible and do not change this conclusion.

For $i+2<m$, each summand of \texttt{E1} has matrix
\texttt{xx1}(x,t).  Its first $N(i)$ columns are the columns
of $x$ with zero auxiliary rows, and only $c_i$ columns
remain.  Thus
\[
       \operatorname{rank}(\mathrm{xx1}(x,t))
                     \leq\operatorname{rank}(x)+c_i.
\]
If $x$ is not injective, every summand is zero, so is
\texttt{XiAt}, and so is the global product.
For $i+2=m$, one has $c_i=1$ and
\texttt{joinCol}(x,y) has rank at most
$\operatorname{rank}(x)+1$, giving the same conclusion.
For $i+1=m$, one has $c_i=0$; the last local factor is
\texttt{PhiTop} applied to a row projection of an invertible
coordinate change of $x$.  Its required unit minor has size
$N(i)$, and a row projection cannot increase rank.
This proves (A13c) in all three cases.  In particular, the
actual array span is a proper subspace of all finite arrays;
this statement does not imply that its first-derivative image
is a proper subspace.

\paragraph{The exact incidence span at a $T$-prime.}
Let $\kappa$ have $q$ elements, let
$M=\operatorname{Mat}_{b\times a}(\kappa)$, and fix a nontrivial
additive character $\psi$ of $\kappa$.
For nonzero row vectors $v\in(\kappa^b)^*$ and
$u\in(\kappa^a)^*$, put
\[
                    I_{v,u}(X)=1_{\{vX=u\}}.
                    \tag{A13d}
\]
Their complex span is exactly
\[
 \{F:M\to\C:\ F(0)=0,
       \ \widehat F(A)=0\text{ whenever }\operatorname{rank}(A)>1\}.
 \tag{A13e}
\]
The Fourier pairing is $\psi(\operatorname{Tr}(A^tX))$;
using its inverse only reverses signs and leaves the assertion
unchanged.  Its dimension is
\[
                 \frac{(q^a-1)(q^b-1)}{q-1}.
                 \tag{A13f}
\]

Here is a full proof.  For a fixed projective row $[v]$,
the functions $F(vX)$ with $F(0)=0$ form a space of
dimension $q^a-1$, spanned by its nonzero-fibre indicators.
Their Fourier support is contained in the frequency block
$\{v^tt:t\in\kappa^a\}$: summing an additive character over
the kernel of $X\mapsto vX$ gives zero unless the frequency
annihilates that kernel.  Distinct projective rows have disjoint
nonzero frequency blocks, and these blocks partition all
nonzero rank-one matrices.  The sum of these spaces is direct.
Indeed, in a zero sum each summand has all nonzero Fourier
coefficients zero, so is constant; its value at zero makes
that constant zero.  There are $(q^b-1)/(q-1)$ projective
rows.  The entire Fourier-rank-at-most-one space has one more
dimension, its zero frequency included; evaluation at zero
is a nonzero functional on it.  The direct sum just constructed
has the dimension of its kernel, proving (A13e)--(A13f).

This is a literal local description of \texttt{phiT}, not
merely an analogy.  Keep both lattices standard and their
$G$-frames and $Q$-frames fixed.  Reduction at $P$ identifies
the source and target residue modules with $\kappa^a$ and
$\kappa^b$.  Every chosen nonzero $u,v$ gives a valid
decoration: define its functional by that residue-linear
map on the localized lattice and by zero off it.  Additivity,
scalar-linearity and the off-lattice condition are immediate;
nonvanishing on the global standard lattice follows from
surjectivity of $\OO_E\to\kappa$.  The same construction
works independently at every $T$-prime.  The matrix-integrality
part of \texttt{phiT} is kept as well.  Thus, on integral
residue matrices, the tensor product of the spaces (A13e)
is precisely the span of the local \texttt{phiT} factors
obtained by varying these raw decorations.  The corresponding
global-array span is the image of this space under multiplication
by the other fixed local factors, followed by the stated class
equivariance.  Equality with an unmultiplied space is asserted
only on a CRT slice where those other factors have fixed nonzero
values.
Passing to the chosen class representatives still uses the
archive's equivariance maps; no new list of representatives
or new full \texttt{TC} is being silently declared.

These restricted incidence functions can isolate a finite
window even though they cannot isolate every residue matrix.
Let $\mathcal W\subset M$ be finite, with $X_0\in\mathcal W$
nonzero.  If $q>|\mathcal W|$, there is $v\ne0$ such that
\[
 vX_0\ne0,\qquad
 I_{v,vX_0}(X)=0\quad(X\in\mathcal W,\ X\ne X_0).
 \tag{A13g}
\]
Each forbidden equation $vX_0=0$ or $v(X-X_0)=0$
is a proper linear subspace of the row space, of cardinality
at most $q^{b-1}$.  At most $|\mathcal W|$ such subspaces
are forbidden.  Their union has cardinality less than $q^b$.
Choosing a row outside it proves the assertion and also
ensures that the source decoration $u=vX_0$ is nonzero.
For a finite set of rational lattice matrices, the same test
applies whenever their reductions are defined, distinct from
$X_0$, and $X_0$ has nonzero reduction.  These are exact
finite divisibility tests.  If several $T$-primes are present,
one can impose (A13g) at one of them and merely preserve
$X_0$ at the others, provided its reduction there is nonzero.

\paragraph{What tame determinant characters and matrix translates span.}
There is an equally explicit answer for the first-step
\texttt{PhiQ} seed.  Before its Jacquet-type finite summation,
its source dimension is $r=k+m-1$.  After the displayed row
permutation, its reduction on
$\operatorname{Mat}_{b\times r}(\kappa)$ is
\[
 F_\chi(X)=
 \begin{cases}
   \chi(\det X_0),&X=\binom{X_0}{0},\quad X_0\in\mathrm{GL}_r(\kappa),\\
   0,&\text{otherwise}.
 \end{cases}
 \tag{A13h}
\]
Here $\chi$ is a multiplicative character of $\kappa^\times$.
Let $p(X)$ be the wedge of the $r$ columns.  The span of
all left and right general-linear translates of (A13h),
as $\chi$ ranges over all multiplicative characters, is
exactly the pullback under $p$ of all functions on the
nonzero decomposable $r$-vectors; these functions are extended
by zero on matrices of rank less than $r$.

Indeed $\mathrm{SL}_r(\kappa)$ acts transitively on the
bases of an $r$-plane with any specified exterior product.
Every function (A13h), and every translate of it, is invariant
under this right action.  Conversely, a left translate can
place the image plane at any specified $r$-plane $W$.
Choose a nonzero $w\in\bigwedge^rW$.  On that plane its
functions are exactly $\chi(d)$ for $p(X)=dw$.
Character orthogonality gives
\[
 1_{\{d=d_0\}}
       =\frac1{q-1}\sum_{\chi}\chi(dd_0^{-1}).
 \tag{A13i}
\]
This proves the span claim.  If the trivial character is
excluded, the exact span consists of the functions whose
sum over the $q-1$ nonzero multiples of each $w$ is zero.
In particular
\[
 1_{\{d=d_0\}}-1_{\{d=d_1\}}
   =\frac1{q-1}\sum_{\chi\ne1}
           (\chi(d_0)^{-1}-\chi(d_1)^{-1})\chi(d)
 \tag{A13j}
\]
has coefficient absolute sum at most $2$.
Thus a second residue can cancel the trivial component
exactly, but the result is a delta on an \emph{oriented
image plane}, not on an individual full-rank matrix.

The source's first local \texttt{XiAt} retains a corresponding
restriction even after the finite summation.  For a rational integral matrix
$A\in\mathrm{GL}_k(E)\cap\mathrm{GL}_k(\OO_{E,P})$, its defining formulas give
\[
       \Xi_{0,P}(xA)=\chi(\operatorname{red}_P(\det A))\Xi_{0,P}(x).
       \tag{A13k}
\]
For clarity this identity is in the displayed local coordinates;
the locally constant identity extends to the completed local
matrix group, and integral changes of source frame preserve its
$\mathrm{SL}_k(\OO_{E,P})$-invariance.
At step zero, \texttt{ChIdx.Ok.zero\_h} sets the source
translate to $1$.  The enlarged source has precisely $r$
columns.  Right multiplication by
$\operatorname{diag}(A,1_{m-1})$ preserves the integrality
and the zero lower block of (A13h) and multiplies its unit
minor by $\det A$.  Also
\[
 \mathrm{xx1}(xA,t)=\mathrm{xx1}(x,t)
                    \operatorname{diag}(A,1_{m-1}),
\]
while its auxiliary character \texttt{chArg} is unchanged.
Every summand and its normalization therefore obey (A13k).
When $m=2$, the first step is \texttt{PsiAt} instead;
the identical calculation uses
$\mathrm{joinCol}(xA,y)=\mathrm{joinCol}(x,y)
\operatorname{diag}(A,1)$.
This proves (A13k), including right special-linear invariance,
without assuming that arbitrary local translates are admitted
by the actual finite index set.  The exact seed span (A13i)
does not by itself prove surjectivity after those index and
Jacquet constraints.

\paragraph{Exact tame anticyclotomic globalization.}
Assume $E/F$ is CM quadratic and unramified at every finite
place.  Let $\mathcal S$ be a finite set of split finite
places of $F$, with a selected prime $P_v$ of $E$ above each
$v$.  Write $\bar P_v$ for its conjugate and $\kappa_v$
for the common residue field.  Let
\[
 R_{\mathcal S}=\prod_{v\in\mathcal S}\kappa_v^\times,
 \quad
 U(\mathfrak m)=
 \prod_{P\nmid\mathfrak m}\OO_{E,P}^\times
 \prod_{P\mid\mathfrak m}(1+P\OO_{E,P}),
 \quad
 \mathfrak m=\prod_{v\in\mathcal S}P_v\bar P_v.
\]
The finite group relevant here is
\[
 C_{\mathcal S}=
 \A_{E,f}^\times/(E^\times\A_{F,f}^\times U(\mathfrak m)).
 \tag{A13l}
\]
It is a quotient of the usual finite ray class group.
For this standard background see Milne,
\emph{Class Field Theory}, v4.03 (2020), Theorem V.1.7,
printed pp.~150--151, and Proposition V.4.6,
printed p.~173; both statements and their complete displayed
proofs were read for this comparison.

At a split place use the ratio convention
$(a,b)\mapsto\bar a\bar b^{-1}\in\kappa_v^\times$,
where bars on $a,b$ in this sentence mean residue reduction.
Local unit representatives give a well-defined homomorphism
$R_{\mathcal S}\to C_{\mathcal S}$.  Its kernel is exactly
\[
 \Delta_{\mathcal S}\mu(E)
    =\{(\zeta\bmod P_v)_v:\zeta\in\mu(E)\}.
 \tag{A13m}
\]
Here is a proof that includes the possible global-unit issue.
If a local-unit idele represents zero in (A13l), write it
as $a f w$ with $a\in E^\times$, $f\in\A_{F,f}^\times$,
$w\in U(\mathfrak m)$.  The valuations of $a$ at conjugate
primes are equal, so $a/\bar a$ is a global unit.  Every
complex conjugate of this unit has absolute value one.
It is therefore a root of unity: the powers are algebraic
integers of bounded degree whose conjugates all have
absolute value one; their monic minimal-polynomial
coefficients belong to fixed finite sets of integers,
so the set of powers is finite.  Taking the local ratios
in the chosen order gives precisely its image in (A13m).

Conversely, for every $\zeta\in\mu(E)$ one can choose
$a\in E^\times$ with $a/\bar a=\zeta$.
For $\zeta\ne-1$ take $a=1+\zeta$; for $\zeta=-1$
take any nonzero purely imaginary element of $E$.
Its valuations at conjugate primes agree.  Since every
finite ramification index in $E/F$ is one, choose an
$F$-idele $f$ with valuation $-v_P(a)$ at the underlying
place.  Then $af$ is an $E$-unit idele at every finite
place.  Let $u$ equal $af$ at the places in $\mathcal S$
and $1$ elsewhere.  The quotient $u/(af)$ lies in
$U(\mathfrak m)$, and the ratio of the two components
of $u$ at $v$ reduces to $\zeta\bmod P_v$.
Thus its class is zero, proving (A13m).  The ramification
index condition is used exactly in the choice of $f$;
no principal-ideal or class-number-one assumption was used.

It follows that prescribed characters
$\lambda_v:\kappa_v^\times\to\C^\times$ occur as the tame
inertia restrictions of a finite-order character trivial
on $E^\times\A_{F,f}^\times$ if and only if
\[
              \prod_{v\in\mathcal S}
                   \lambda_v(\zeta\bmod P_v)=1
                         \quad(\zeta\in\mu(E)).
 \tag{A13n}
\]
Necessity follows from (A13m).  For sufficiency, the
product character descends to the image of $R_{\mathcal S}$
in $C_{\mathcal S}$.  Every character of a subgroup of
a finite abelian group extends to the group.  One elementary
proof adjoins one group element at a time: if $x$ has order
$d$ modulo the old subgroup, choose a $d$th root of the
specified value on $x^d$, and set the new character on
$hx^j$ equal to the old value on $h$ times that root to
the power $j$.  This is well-defined and enlarges the
domain.  Iteration terminates.  The final character has
values in roots of unity of order dividing
$|C_{\mathcal S}|$.  Its pullback to finite ideles has
the requested restrictions, is trivial on all wild inertia,
and is unramified outside $\mathcal S$.
The class group influences the choices and possible orders
of the extension, not an additional condition on (A13n).

This gives an actual \texttt{Tup} once all places in
$\mathcal S$ are split degree-one places outside the
forbidden set: compare \texttt{GoodPl}, \texttt{finChar}
and \texttt{tupGrp}, lines 64--115 of
\path{Formalism/LocalToGlobal/Divisibility/Tuples.lean}.
The exact residue character at $(a,b)$ is
$\lambda_v(\bar a\bar b^{-1})$; changing the selected
prime inverts this convention.

\paragraph{One auxiliary prime removes the root-of-unity condition.}
Prescribe any finite collection of local characters at the
main places.  Choose one additional allowed split degree-one
place $v_*$ whose residue characteristic does not divide
$|\mu(E)|$.  Reduction injects $\mu(E)$ into
$\kappa_{v_*}^\times$: roots of unity of order prime to
the residue characteristic cannot lie nontrivially in
$1+P_{v_*}$, as follows either by the binomial expansion
of $\zeta^n=1$ with $n$ a unit or by uniqueness of simple
Hensel lifts.  The inverse of the main product in (A13n)
is a character of that embedded subgroup.  Extend it to
$\kappa_{v_*}^\times$ by the same finite-group argument.
Then (A13n) holds and the prescribed main characters
globalize without alteration.  One auxiliary prime works
for the entire finite collection, with a possibly different
auxiliary character for each member.

For a finite forbidden set such primes are available by
Chebotarev, choosing rational primes split completely in
a Galois closure of $E$ and excluding the finite list.
The exact background locator is Milne, \emph{Algebraic Number
Theory}, v3.08 (2020), Theorem 8.31 and Corollary 8.32,
printed p.~147; these give the positive density of the
identity Frobenius class and hence infinitude.
For a \emph{fixed} actual \texttt{TC}, however, its set
$I.S$ is literally an arbitrary set of rational integers;
\texttt{ColIn.Ok} only requires it to contain the specified
$T$-primes, the two $G$-primes, and the primes dividing
the host conductor.  Thus the assertion for that fixed
datum requires an exhibited spare allowed place; it does
not follow solely from the type of \texttt{TC}.

The auxiliary prime changes its local \texttt{PhiQ}/\texttt{XiAt}
factor.  Its value and support are part of the new input.
In particular, the exact character identity (A13j)
at a main prime is not a global-array identity until
these auxiliary factors are included.  There are only
$|\mu(E)|$ possible restrictions to cancel, so one can
choose at most that many auxiliary residue characters
for this purpose.  This is a finite, explicit local cost,
not a proof that the auxiliary kernel equals $1$.

\paragraph{Coefficient enlargement and the remaining finite comparison.}
For the finite character family, a single number field
$K'$ can contain its values, all required
$\zeta_{q^{2m}}$, and all required $\sqrt q$.
Take the compositum of the old normal coefficient field
with the corresponding finite cyclotomic and quadratic
extensions, in the chosen complex embedding.  Complex
conjugation restricts to the old involution.  The
\texttt{Fields.Ok} conditions transport: the $E/F$ data
are unchanged; every embedding of $E$ into the new field
already has image in the old normal field, so the list
of embeddings and the defining equalizer $K_q\subset E$
are unchanged; total complexity and normality persist.
This is a statement about the coefficient package, not
an alteration of the rigid full CM host field or its
Galois frame.

In fact the entire literal \texttt{Base} survives this
coefficient enlargement.  Its complete definition is
\path{Development/AbelianVarietyTheory/Definitions/Foundations/}
\nolinkurl{RootOfUnityEightQMemCoeffGaloisSubfield.lean},
lines 261--304.  Keep \texttt{Pre} unchanged.  Embed the
old $\mathrm{LL}$ in the chosen normal field $K'$ and
use the images of the \emph{old} intermediate fields
$\mathrm{LL}_0$ and $K_s$ as the new intermediate fields.
They are isomorphic to their originals, rather than
being replaced by their composita with $K'$.
The old $\mathrm{LL}_0$ is still Galois over $\Q$;
the inclusions $E\subset K_s\subset\mathrm{LL}_0$
and every stipulated root of unity transport unchanged.
An isomorphism of number fields identifies their rings
of integers, ideals and principal ideals.  It therefore
transports each generator in \texttt{capit} and
\texttt{kPsi}, with exactly the same principal-ideal
equalities.  Every complex embedding of the image of
$K_s$ is uniquely the old embedding composed with this
isomorphism's inverse, so the norm equalities quantified
over \emph{all} embeddings in \texttt{kPsi} also persist.
The roots required merely in $\mathrm{LL}$ remain in
$K'$.  Conjugation on the chosen normal subfield of
$\C$ gives the required involution, compatible with
the old one and the prescribed complex embedding.
These are all the additional clauses of \texttt{Base}.
In particular the CM host, its Galois frame, the ranks,
the $T$-primes and their fixed windows, and the old
integral coefficient ring $\OO_{K_s}$ are unchanged.

The canonical local indices are supplied by
\texttt{twIdx\_ok\_closed}, lines 87--102 of
\path{Formalism/LocalToGlobal/Columns/}
\nolinkurl{UnfoldingNonvanishingDoubleCosetHolds.lean}.
Its hypotheses explicitly include the existing field,
prime, reference-character and Weil tying conditions,
the new good-prime and unit conditions, and the roots
and square roots just listed.  Its displayed proof was
read; it invokes the archive's nonvanishing and finite-index
theorems.  We do not replace those invocations by a claim
of a fresh independent proof.

\paragraph{Complete realization of a prescribed finite character family.}
Here the source provides a stronger positive conclusion
than merely the existence of the local indices.  Fix the
original \texttt{Pre}, an oriented \texttt{PreX}, the
reference frame $f$, a prime-selection map \texttt{sel}
and a valid ambient theta datum $J$.  Retain the displayed
diagonal-space and $\mathrm{duStar}$ identities of that
geometric input.  Let $S$ contain its $T$-primes,
the two $G$-primes and all primes dividing the host
conductor.  Suppose that the finite family $\Xi$ contains
$1,\delta,\delta'$ and that every $\xi\in\Xi$ is a
finite-order character trivial on $\A_F^\times$, tame
at every finite place, and ramified only at split,
absolutely unramified degree-one places outside $S$.
These are precisely the clauses of \texttt{Chr.IsTwist}
in
\path{Development/AbelianVarietyTheory/Definitions/AutomorphicMembership/}
\nolinkurl{IsRationalUnitaryDeterminantMapRationalUnitaryTo.lean},
lines 111--114.  Thus the characters constructed above
satisfy them whenever the main and auxiliary places
have been selected as stipulated.

Assume the original coefficient field contains the
normalized reference-character values required by
$\mu_{0,j}.\mathrm{ValuesIn}(0)$ and
$\eta.\mathrm{ValuesIn}(1/2)$, and the square roots
of the two fixed $G$-primes.  These are explicit reference
conditions, satisfied by the original fully tied
column construction.  Enlarge only $\mathrm{LL}$ as
above to contain the values of all $\xi\in\Xi$ and
$\sqrt q,\zeta_{q^{2m}}$ for the finite set of their
ramified residue cardinalities.  This is finite: if
the character orders divide $M$ and $\mathcal Q$ is
that set of residue cardinalities, one possible
compositum has the coarse relative degree bound
\[
 [K':K]\ \leq\ \varphi(M)
                   \prod_{q\in\mathcal Q}2\varphi(q^{2m}).
 \tag{A13o}
\]
All factors on the right bound the degrees of the
explicit cyclotomic or quadratic extensions being
adjoined.  Intersections can only decrease this bound.
For the characters extended from a single
$C_{\mathcal S}$ above one may use $M=|C_{\mathcal S}|$,
also including the orders of the fixed reference
characters if necessary.

There then exists an actual finite \texttt{TC} on this
same geometric \texttt{Pre}, with exactly the character
set $\Xi$, the prescribed \texttt{sel} and the same
ambient datum $J$, satisfying \texttt{Ok},
\texttt{Laws} and \texttt{Law0}.  It may be chosen
canonical in the sense \texttt{IsTw}.  If the original
\texttt{Pre} has \nolinkurl{StagesInhab.HasLetters}, its
\texttt{Later} data exist as well.

\textbf{Proof and exact source scope.}
The full constructor is \texttt{InEx.exists\_colIn},
lines 195--289 of
\path{Development/AbelianVarietyTheory/Proofs/IwasawaDecomposition/}
\nolinkurl{ExistsColumnInputAssembled.lean}.
Its condition \texttt{Holds}, lines 126--128, is exactly
the value-field, square-root and depth-$2m$ root-of-unity
condition just arranged.  Its proof first constructs
the reference finite characters, reads each member of
$\Xi$ as a tuple, applies the local-index existence
theorem to each member, and uses $\Xi$ itself as its
finite injective index set.  It returns \texttt{Ok},
the asserted equalities of character range, $S$,
\texttt{sel} and $J$, and the exact constants
$M=A=B=m$, $D_Q=2m$.
In particular its \texttt{Deep} condition holds.
The complete proof of \texttt{canon\_ok} and
\texttt{canon\_isTw} in \texttt{IsTwIffCanonEq.lean}
then replaces only its constants and indices by their
canonical versions.  The ambient theta datum,
characters and selected primes are preserved.

For this actual valid input, the statements needed next
are proved theorems, not additional Prop declarations.
In
\path{Development/AbelianVarietyTheory/Definitions/AutomorphicMembership/}
\texttt{ExistsOkPins.lean},
\texttt{TrueColsTest.laws}, lines 86--103, proves
$I.\mathrm{Ok}\Rightarrow I.\mathrm{Laws}$;
\texttt{law0}, lines 152--159, proves
$I.\mathrm{Ok}\Rightarrow I.\mathrm{Law0}(l)$ for
every $l:I.\mathrm{Laws}$; and \texttt{later},
lines 194--204, constructs all five later stages
from $l$ and \texttt{HasLetters}.
The first assembles the chain, level, period and lattice
theorems.  The second pulls functions back through the
step-zero class bijection, whose representatives have
the same underlying lattice, and hence preserves the
integral lattices as well as their complex embeddings.
The third constructs the finite-dimensional row model
and its divisor, level, local and jet stages.
All three complete displayed proofs were read.
\texttt{HasLetters} depends only on the unchanged
\texttt{Pre}; in the archive's contextual construction
\nolinkurl{JetEx.hasLetters_of_over}, lines 16--26 of
\nolinkurl{ExistsColInJetsIn.lean}, supplies it from
the existing \texttt{Ctx.Over} hypothesis.
This proves the composite realization statement using
the named source theorems.\hfill$\square$

The stronger specialized \texttt{ColW} carries another
restriction.  Its literal definition, lines 23--36 of
\texttt{PositionOfIntro.lean}, fixes its character range
to $T.\mathrm{set}(f)$ for one \texttt{TwDat}.
The latter definition, lines 79--97 of
\nolinkurl{CmQuadraticFormEqAlternatingForm.lean}, requires
the auxiliary character indexed by a coefficient prime
$p$ to have $p$-power order, with its \texttt{IsPsi},
conjugation and avoidance conditions.  Its range consists
of $1,\delta,\delta'$ and the specified products of
one of $\delta,\delta'$ with these auxiliary characters.
An arbitrary globally admissible finite family is not
automatically of this special form.
For this specialized witness the forbidden-prime set is,
however, finite.  The host's \texttt{levelPrimes} is
literally a finite set (\path{Formalism/LocalRingCardinality/}
\texttt{DataHost.lean}, line 69), and the complete proof
\texttt{finite\_offSet}, lines 208--216 of
\nolinkurl{CmQuadraticFormEqAlternatingForm.lean}, places
all excluded primes in a finite divisor set together
with those level primes and one distinguished prime.
The \texttt{ColW.S\_eq} clause therefore proves that
its $S$ is finite.  Hence an auxiliary completely split
prime outside it, the main support, and the divisors of
$|\mu(E)|$ is available by the cited Chebotarev theorem.
The unresolved specialized-family condition is its
\texttt{TwDat} shape, not the existence of a spare prime.

There is also a separate positive arithmetic-jet source
theorem: \texttt{exists\_bound\_jetsIn}, lines 34--40
of \nolinkurl{ExistsColInJetsIn.lean}, gives a finite list
of algebraic coefficient generators after which a valid
$I'$ with \texttt{JetsIn} exists and has the same
reference frame, forbidden set and character range.
Its conclusion does not retain the ambient datum $J$.
The following refinement retains $J$ by allowing that
finite generator list to depend on it.

\paragraph{A fixed-ambient-datum arithmetic-jet refinement.}
Under the preceding realization hypotheses and
\texttt{HasLetters}, after a further finite coefficient
extension the complete canonical \texttt{TC} can be
chosen with the \emph{same} $J$ and with
\texttt{JetsIn}.  One may simultaneously require the
same field to contain all normalized jets of any
prescribed finite list of total degrees, including
degree one, on the admitted lattices and torsion
arguments of these finitely many steps.  This is
an arithmetic-field assertion; it asserts neither
nonzero jets nor surjectivity of their span.

Here is the full refinement argument.  The crucial
source theorem is \texttt{JetEx.finAlg\_step},
lines 193--226 of
\path{Development/AbelianVarietyTheory/Proofs/IwasawaDecomposition/}
\texttt{DegreeEqTotDeg.lean}.  It applies to \emph{any}
valid ambient datum $J$ in an $I.\mathrm{Ok}$ input,
and to every total degree $n$.  It states that the
corresponding normalized jet set is contained in a
number field.  Its proof invokes the normalized-vacuum
algebraicity and finite-jet-set theorems, with the
displayed self-duality, parity, bad-prime and ambient
normalization hypotheses.  This source result is used
with those hypotheses, not inferred from arbitrary
analytic theta normalization.

Let $\mathcal B$ be the finite set of primes of $E$
above the fixed $T$- and $G$-primes and the ramification
of the finitely many characters in $\Xi$.
For a word $u$, a step $i<m$, and subsets
$S_0,R_0\subset\mathcal B$, put
\[
 M(u,R_0)=T(u)
             \left(\prod_{Q\in R_0}\mathrm NQ\right)^{2(m+1)}.
 \tag{A13p}
\]
The exponent is precise: the canonical
\texttt{twCst} has $A=B=m+1$, $M=m$,
as displayed in \path{Formalism/LocalToGlobal/Columns/}
\texttt{ChainTw.lean}, lines 217--226.
For each such finite index and each selected degree,
form the source's normalized \texttt{jetSet} using
$J$, $S_0$, modulus (A13p), and the fixed ideal
character $\chi$.  These sets are independent of the
coefficient field: \texttt{hostAmb} and
\texttt{jetSet\_amb}, lines 43--66 of
\texttt{DegreeEqTotDeg.lean}, prove this exactly.

Keep the members of this finite list which are
contained in a number field, and take their union.
It is contained in the compositum of finitely many
number fields.  Adjoin a finite set of algebraic
generators for that compositum to the character
coefficient field, using the same \texttt{Base}
enlargement construction.  This finite-choice argument
is the explicit mathematical content of the source's
\texttt{FinAlg} union step; it does not require a
decision procedure for testing the members of the list.

Now construct and canonicalize the actual input in
this enlarged field, retaining $J$.  For every one of
its words and steps, its actual sets
$\mathrm{Sb}(a,u)$ and $\mathrm{rowSG}(a)$ are subsets
of $\mathcal B$.  The complete proofs of
\texttt{Sb\_subset} and \texttt{rowSG\_subset},
lines 126--156 of the same file, establish precisely
this containment from $I.\mathrm{Ok}$ and the character
range.  Its actual modulus is (A13p), by the definitions
of \texttt{cAof}, \texttt{cBof} and \texttt{twCst}.
Apply \texttt{finAlg\_step} to this actual valid input.
Its jet set is consequently one of the retained sets
in the finite union, and all its values lie in the
chosen field.  Taking, for each $u,i$, the total degree
of the required embedding multidegree
$\nu\mathrm{Deg}(u,i,\tau)$ proves
\texttt{JetsIn}; the identity between these total
degrees is proved in \texttt{degree\_eq\_totDeg},
lines 75--84.  Adding any finitely many other degrees
to the initial list proves the additional assertion.
The universal laws and later-stage constructors apply
again without a change of \texttt{Pre}.\hfill$\square$

This proof improves the wrapper's preservation statement,
but only for a datum already satisfying the literal
\texttt{AmbIn.Ok}.  The ambient structure and that validity definition fix its
discriminant-normalized vacuum and differential generators;
it is not a statement that an arbitrary 032 normalization
is such a datum.  The comparison with a particular 032
section must still retain the actual section, tangent
and period factors, as in A8 and A12.  Nor does (A13o)
bound the degree of this further arithmetic-jet field:
the latter is supplied by the stated finite algebraic
generation theorem and its finite compositum.

Finally, (A13g) is an actual \emph{array}-window selector.
The precise source passage from this array to a holomorphic
theta section is \texttt{reads\_datum\_eq\_thC}, in
\path{Development/AbelianVarietyTheory/Proofs/DecoratedLatticeProofs/}
\nolinkurl{DatumReadingSplitPlacesEqClassicalTheta.lean}.
Its complete proof expresses the output as a finite
Fourier sum of Heisenberg translates of the chosen vacuum;
the half-pairing phases, semicharacters, and lattice
representatives all occur in that sum.  Thus a finite
window in the raw $E$-matrix array is not automatically
the same window in the rational Schr\"odinger Gaussian
series of A12.  For any given full datum, use that explicit
finite transform and then the derivative matrix: apply it
to the incidence basis (A13d), or to the oriented-plane
character combinations (A13i)--(A13j), and compute its rank
and a nonzero minor with all common section and tangent
periods removed as in A12.  This is a smaller, exactly
specified finite comparison.  The new results above close
the global tame-character compatibility condition and
give literal array selectors and invariant subspaces;
they do not assert the unproved surjectivity of this
final transformed derivative matrix.

In the requested Grothendieck framework this last distinction
is essential: the finite theta-module comparison belongs
to the period and correspondence comparison discussed in
\emph{cote44}, scans 7--8.  It is not an algebraic Lefschetz
inverse or a cohomological projector of the kind discussed
in \emph{cote16}, scans 43--46.  A later inverse matrix is
algebraic only after its stated period factors have been
removed, and an actual correspondence still requires the
cohomological identification.


\subsection{THETA3-A14. A prescribed tame primary family in one actual \texttt{TwDat}}
\label{A14}

This is a construction inside the fixed full host, rather than a change
of its CM field, ranks, or signatures.  Write $E/F$ for the host CM
extension, and $c$ for complex conjugation.  We retain the archive
hypotheses used in A13: $E/\Q$ is Galois and $E/F$ is unramified at every
finite place.  We also retain the entire \texttt{Pre}, the reference
character frame $f$, and the complex images of the two distinguished
coefficient subfields $K_s$ and $L_0$.  The extension of the larger
coefficient field $L$ below is exactly an \texttt{Enl} in the sense of
\nolinkurl{CmQuadraticFormEqAlternatingForm.lean}, lines 71--74.

Here and below a primary character means a finite-order character of
prime-power order; it does not mean a character of residue-characteristic
power conductor.  In this construction the conductor is tame and is
prime to the primary prime.

\paragraph{Statement and the literal extra local condition.}
There are infinitely many odd rational primes $p$, avoiding any
prescribed finite set and $|\mu(E)|$, with the following property.
Let $\mathcal S$ be a finite collection of split places of $F$, lying
over rational primes that split completely in $E$ and avoid
\texttt{offSet} and $p$.  At each $v\in\mathcal S$ choose an ordering of
its two places in $E$.  Given finitely many systems
\[
 \lambda_j=\bigotimes_{v\in\mathcal S}\lambda_{j,v},
 \qquad 1\leq j\leq t,
 \tag{A14a}
\]
where $\lambda_{j,v}$ is any $p$-primary character of
$\kappa_v^\times$, there exist finite-order anticyclotomic Hecke
characters $\psi_j$ whose restriction at this split pair is
$\lambda_{j,v}(a/b)$.  They are unramified outside the indicated split
pairs, are tame everywhere, have $p$-power order, and are
\emph{trivial on all of $E_w^\times$} at every $w\mid p$.

After a finite normal coefficient-field enlargement, these characters
occur together in a single literal \texttt{TwDat}, with
\[
 \texttt{pl}=\{p\},\qquad
 \texttt{TwDat.set}(f)=
 \{1,\delta,\delta'\}\ \cup\ \{\delta\psi_j:1\leq j\leq t\}.
 \tag{A14b}
\]
Repeated characters are of course discarded in this set.  We choose
$p\ne p_{\delta},p_{\delta'}$.  If there are no prescribed characters,
the untwisted datum suffices.

The distinction between unramified and locally trivial is essential.
The actual definition is
\path{Formalism/LocalToGlobal/UnitaryNorms/CharHost.lean},
\texttt{IsPsi}, lines 77--84.  Besides order, anticyclotomicity,
tameness, exact ramification set, and split ramification, its
\texttt{atl} clause requires $\psi_w(x)=1$ for \emph{every}
$x\in E_w^\times$ when $w\mid p$.  The construction proves this clause,
not merely triviality on local units.

\paragraph{Proof of the character construction.}
Complex conjugation is central in $\operatorname{Gal}(E/\Q)$.
Choose $p$ unramified in $E$ with Frobenius $c$ and outside the indicated
finite set.  Chebotarev supplies infinitely many such primes; the precise
primary reference used here is Milne, \emph{Algebraic Number Theory},
version 3.08 (2020), Theorem 8.31 and Corollary 8.32, printed page 147.
Its restriction to $F$ is trivial, so $p$ splits completely in $F$;
at each resulting place, $E_w/F_v$ is the unramified quadratic
extension.  The finite exclusions include $2$, $|\mu(E)|$, the two
reference primes, and any already fixed bad primes.

Use the ratio convention at a split pair exactly as in A13.  Put
\[
 R_{\mathcal S}=\prod_{v\in\mathcal S}\kappa_v^\times,
 \qquad
 C_{\mathcal S}=
 \A_{E,f}^\times/
       \bigl(E^\times\A_{F,f}^\times U(\mathcal S)\bigr).
\]
Here $U(\mathcal S)$ consists of the principal local units at the two
places above each $v\in\mathcal S$ and all units elsewhere.  This is a
finite quotient of a ray class group.  A13 proves, including the
unramified valuation cancellation, that the kernel of
$R_{\mathcal S}\longrightarrow C_{\mathcal S}$ is the image of
$\mu(E)$ under the split ratios.  Each $\lambda_j$ has $p$-power
order.  Because $p\nmid |\mu(E)|$, its restriction to this kernel is
trivial.  Thus it defines a character on the image of
$R_{\mathcal S}$, and the elementary finite-abelian character extension
argument in A13 extends it to $C_{\mathcal S}$.

One can make this extension $p$-primary without changing the prescribed
local characters.  Write
$|C_{\mathcal S}|=p^a M$ with $(p,M)=1$, and choose an integer $e$ with
\[
 e\equiv1\pmod{p^a},\qquad e\equiv0\pmod M.
\]
Raising any extension to the $e$th power kills its prime-to-$p$ part
and acts identically on its prescribed $p$-primary restriction.  The
result is $\psi_j$.  Extend it trivially at infinity.  The defining
quotient gives triviality on $E^\times$ and on the finite ideles of
$F$; with trivial infinity type it is anticyclotomic on all ideles.
The subgroup $U(\mathcal S)$ gives tameness and the asserted conductor
support.  Let $Q_j$ be its exact finite ramification set, possibly
smaller than the prescribed collection of pairs.

At a place $w\mid p$, the character is unramified.  Since
$E_w/F_v$ is unramified, a uniformizer of $F_v$ is also one of $E_w$,
and therefore
\[
 E_w^\times=F_v^\times\,\OO_{E_w}^\times.
 \tag{A14c}
\]
Anticyclotomicity makes the local character trivial on the first
factor; unramifiedness makes it trivial on the second.  This proves
\texttt{atl}.  Every place in $Q_j$ is split, has prime absolute norm,
and lies above a rational prime outside \texttt{offSet}.  Such a prime
splits completely in $E$, so its ramification index is one.  These are
exactly the further clauses of \texttt{Chr.IsTwist}.  The defining
source is
\path{Development/AbelianVarietyTheory/Definitions/AutomorphicMembership/}
\nolinkurl{IsRationalUnitaryDeterminantMapRationalUnitaryTo.lean},
lines 111--114.  We have proved all clauses of
\texttt{IsPsi} and \texttt{IsTwist}, with $Q_j$ the exact ramification
set.  In particular $Q_j$ contains no place above $p$.

There is no shortage of possible conductor primes: those splitting
completely in $E(\zeta_p)$ avoid any fixed finite set and satisfy
$q\equiv1\pmod p$.  Hence nontrivial tame $p$-primary characters occur
at infinitely many allowed primes.  This assertion does not prescribe
the power-residue values of a previously fixed finite list of matrix
entries.

\paragraph{Enough coefficient-place orbits, with a degree bound.}
The index $w$ of \texttt{TwDat} is a finite place of the coefficient
field, not a place of $E$.  First enlarge the normal coefficient field
$L$ to contain all finitely many character values and any already
specified algebraic constants.  For $t\geq1$ choose $r\geq0$ with
$2^r\geq t$.  There is a totally real multiquadratic extension
\[
 M=\Q(\sqrt{q_1},\ldots,\sqrt{q_r}),
 \qquad [M:\Q]=2^r,
 \tag{A14d}
\]
which is linearly disjoint from $L$, and in which $p$ splits
completely.  To construct it, choose distinct rational primes
$q_i\equiv1\pmod{4p}$, none dividing the discriminant of $L$.
The split-prime case of the same Chebotarev theorem applied to
$\Q(\zeta_{4p})$ supplies them.  Independence of their rational
squareclasses gives the stated degree.  Every nontrivial quadratic
subfield of $M$ is generated by the square root of a nonempty product
of the $q_i$ and is ramified at at least one of these primes.
It cannot be contained in $L$.  A nontrivial subfield of a
multiquadratic extension contains a quadratic subfield, so
$M\cap L=\Q$.  Also $q_i\equiv1\pmod p$, and Hensel's lemma at the
odd prime $p$ makes $q_i$ a square in $\Q_p$.  Thus $p$ splits
completely in $M$.

Replace $L$ by $LM$.  Each prime of $L$ above $p$ has $2^r$ primes
above it in $LM$: locally this is
$L_w\otimes_{\Q_p}(M\otimes_\Q\Q_p)
 \simeq L_w^{\,2^r}$.
Complex conjugation fixes $M$ pointwise, so it fixes each prime of
$M$ above $p$ and acts only on the $L$ component of these labels.
Consequently every conjugation-orbit of primes of $L$ above $p$
produces $2^r$ distinct conjugation-orbits in $LM$.  There are at
least $t$ such orbits.  The extra degree for this orbit-separation
step is
\[
 [LM:L]=2^{\lceil\log_2 t\rceil}<2t
 \quad(t>1),
 \qquad [LM:L]=1\quad(t=1).
 \tag{A14e}
\]
This bound concerns only the orbit-separation step.  It excludes
the earlier adjoining of character values and the later normalized
jet constants.

Choose $t$ distinct conjugation-orbits of coefficient primes above
$p$.  On the $j$th orbit set
$\psi(w)=\psi_j$, $k(w)=k_j$, and $Q(w)=Q_j$, where
$\psi_j^{p^{k_j}}=1$.  On every remaining coefficient prime put
$\psi(w)=1$, $k(w)=0$, and $Q(w)=\varnothing$.
Set \texttt{pl}$=\{p\}$.  The full definition
\nolinkurl{CmQuadraticFormEqAlternatingForm.lean}, lines 79--89,
now checks clause by clause.  The three character conditions and
\texttt{Q\_l} were just proved; \texttt{off} holds by definition;
\texttt{conj} holds because the assignment is constant on each orbit.
At non-$p$ places the trivial character satisfies \texttt{IsPsi} for
their own residue characteristic.  At $p$ the prescribed characters
satisfy \texttt{IsPsi} by (A14c).  The formula for
\texttt{xiTw}, lines 91--95 of that source, gives (A14b), because
$p\ne p_\delta$ and all other assigned characters are trivial.

This datum survives every subsequent finite normal coefficient-field
enlargement.  Assign to a new coefficient prime the data of its
contraction to the old field.  Contraction preserves residue
characteristic and commutes with complex conjugation; lying over
ensures that every old character still occurs.  Thus all
\texttt{TwDat} clauses and its exact character range persist.
The A13 construction of \texttt{Base.Enl} keeps the two distinguished
subfields as their old images, so it also preserves the literal
\texttt{Base}, not just the list of fields.

\paragraph{An actual correct-range TC and, from an existing witness, ColW.}
For the TC assertion retain A13's diagonal-space and
\texttt{duStar} identities, the required reference-character values,
and its oriented \texttt{PreX}, valid selection, and valid
\texttt{AmbIn}.  A retained $K.\texttt{Over}\,\texttt{Pre}$
supplies the first two geometric identities and
\texttt{HasLetters}; otherwise state those conditions separately.
Apply A13 to the finite set (A14b), adjoining its finite
\texttt{Holds} constants at depth $2m$ and, if requested, the finite
normalized-jet generators for a fixed valid \texttt{AmbIn}.
The preceding pullback of \texttt{TwDat} preserves its range through
these extensions.  The exact source
\path{Development/AbelianVarietyTheory/Proofs/IwasawaDecomposition/}
\nolinkurl{ExistsColumnInputAssembled.lean},
\texttt{isTwist\_of\_mem\_set}, lines 301--315, proves that every
member of this range is an allowed twist for
\[
 S=\{q:q\in\texttt{offSet},\ q\ne p_\delta,p_{\delta'}\}.
\]
Its proof includes both reference characters and their products with
the new characters.  The same file's
\texttt{exists\_colIn}, lines 195--289, then gives the correct-range
TC with the retained orientation, selection and valid ambient datum;
A13 supplies its canonical constants, \texttt{Laws} and
\texttt{Law0}; it supplies \texttt{Later} when
\texttt{HasLetters} holds.  Thus the new family is not merely
a collection of abstract Hecke characters.

There is also a full \texttt{ColW} statement with a concrete starting
hypothesis.  Suppose a witness \texttt{ColW} for this same
\texttt{Pre} and frame is already given, and the retained context
satisfies $K.\texttt{Over}\,\texttt{Pre}$.  Then the preceding construction
produces a new \texttt{ColW} with character range (A14b).
Here is the verification of the additional fields, which is needed
beyond the TC construction.  Orient its \texttt{PreX} using
\texttt{exists\_oriented}, lines 91--103 of
\nolinkurl{ExistsColumnInputAssembled.lean}.  That explicit constructor
changes only the two \texttt{piG} choices and leaves
\texttt{piT} unchanged.  The two clauses of \texttt{TC.Pins}
in \path{Definitions/AutomorphicMembership/LatAtLattice.lean},
lines 125--130, involve only \texttt{Pre} and \texttt{piT}; hence they
are retained.  The equality \texttt{KsExact} is an equality of the
fixed complex field $K_s$ with the field generated by the unchanged
reference data, so it is retained as well.  The original
\texttt{ValLL0} asserts that those reference values lie in the fixed
complex image of $L_0$; it is likewise unchanged.

Finally \texttt{FieldOK} consists of the presence of the fixed
elliptic $j$-invariant and the principality of specified extended
fractional ideals.  This is its literal definition in
\path{Definitions/AutomorphicMembership/IsHalfPtIff.lean},
lines 31--41.  A generator of a principal ideal remains a generator
after coefficient extension: extending the equality
$I=x\OO_L$ gives $I\OO_{L'}=x\OO_{L'}$.  The defining spanning
generators of \texttt{ellLat} are the same elements of $E$, so its
new ideal is precisely that extension.  The $j$-invariant remains
in the larger field.  Thus \texttt{FieldOK} is preserved.
Combine these data with the new TC, its proven \texttt{Laws} and
\texttt{Law0}.  They are exactly all fields of \texttt{ColW},
\path{Definitions/AutomorphicMembership/PositionOfIntro.lean},
lines 23--36.  The retained $K.\texttt{Over}$ gives
\texttt{HasLetters}, and hence \texttt{Later}, as in A13.
This argument does not manufacture the original full host or
the starting \texttt{ColW} from a numerically matching generic CM torus.

\paragraph{Exactly which local character span has been obtained.}
At a chosen conductor prime write
$|\kappa_v^\times|=p^a h$, with $(p,h)=1$, and let
$H_v\subset\kappa_v^\times$ be its unique subgroup of order $h$.
The $p$-primary characters are precisely the character group of
$G_v=\kappa_v^\times/H_v$.  By finite Fourier inversion they span
all functions on $G_v$; the nontrivial ones span its mean-zero
subspace.  Explicitly, for distinct $z_0,z_1\in G_v$,
\[
 \mathbf1_{z=z_0}-\mathbf1_{z=z_1}
  =\frac1{|G_v|}\sum_{\chi\ne1}
       \bigl(\overline{\chi(z_0)}-\overline{\chi(z_1)}\bigr)\chi(z),
 \qquad
 \sum_{\chi\ne1}
   \left|\frac{\overline{\chi(z_0)}-\overline{\chi(z_1)}}{|G_v|}\right|
   \leq2.
 \tag{A14f}
\]
Every finite list of characters used in this identity can occur
together in the single datum just constructed.  On the literal
first-column \texttt{PhiQ} seed, this is a function of the
nonzero oriented wedge modulo $H_v$, retaining the full-rank support
and the right-$\mathrm{SL}$ covariance proved in A13.  It is not a
delta function on all matrices or even on all nonzero determinants.
At different primes the prescribed systems can be chosen independently,
so the corresponding tensor spans also occur in a common datum.
Actual global arrays still include their other local factors and
their class equivariance.  Their theta sections are obtained through
the literal finite Heisenberg transform, not by replacing those
factors with arbitrary Schwartz functions.

The new conclusion is the exact specialization of a prescribed finite
primary local family to one \texttt{TwDat}, and, with the stated
starting witness, to one \texttt{ColW}.  The remaining derivative
question is now a concrete finite matrix on this specified quotient
character space together with the incidence space of A13.  Arithmetic
inversion of such a matrix still requires a nonzero minor and the
section and tangent normalizations of A8 and A12.  In the requested
Grothendieck framework, this is a finite period comparison of the
kind isolated in \emph{cote44}, scans 7--8; it does not by itself
supply the cohomological correspondence or algebraic Lefschetz
inverse discussed in \emph{cote16}, scans 43--46.


\subsection{THETA3-A15. The literal finite transform, its derivative matrix, and rank-one sampling}
\label{A15}

We first compute the actual finite transform for the whole allowed
array.  We then prove a positive first-derivative theorem for its
single-incidence Fourier model.  The final paragraph specifies the
additional tensor-product test required when several $T$-primes occur.

\paragraph{The exact operator and first derivative.}
Use the actual data of
\path{Development/AbelianVarietyTheory/Proofs/DecoratedLatticeProofs/}
\nolinkurl{DatumReadingSplitPlacesEqClassicalTheta.lean},
\texttt{reads\_datum\_eq\_thC}.  In particular the source has a valid
\texttt{StdPair}, the stipulated half idempotent, self-dual lattices,
vacuum and \texttt{IsThetaAt} identity, and the required support and
periodicity of the actual array $\Xi$.  Set
\[
 H=\mathrm{HV}(\mathrm{sc}),\quad
 \alpha=\mathrm{alV}(d,\mathrm{sc}),\quad
 F=\mathrm{xiN}(f_0).
\]
Here $H$ is linear in its first argument and conjugate-linear in its
second.  The exact definition, including its sign, is
\path{Formalism/LocalToGlobal/TorsionIntegrality/QuadraticSemicharacters/}
\texttt{IsSemiChar.lean}, lines 9--24:
\[
 U_xF(z)=F(z+x)
       \exp\bigl(-\pi H(z,x)-\tfrac\pi2H(x,x)\bigr).
 \tag{A15a}
\]
Write $S_X,S_A,S_B$ for the three finite representative sets in the
source proof, and put
\[
 b(y,s)=\exp(-2\pi i\,\operatorname{Im}H(s,y)),
 \qquad
 K_\Xi(x,y)=\frac1{|S_A|}\sum_{s\in S_A}
          \Phi(s-x)b(y,s)^{-1},
 \tag{A15b}
\]
where $\Phi=\mathrm{onV}(\Xi)$.  Then the literal conclusion of that
proof is
\[
 \Theta_\Xi(z)=
  \sum_{x\in S_X,\ y\in S_B}
      \alpha(x)\alpha(y)K_\Xi(x,y)U_xU_yF(z).
 \tag{A15c}
\]
The definitions are
\path{Formalism/LocalToGlobal/KernelLattices/XiIndCClass.lean},
\texttt{brk}, \texttt{hatS}, \texttt{xiL}, lines 18--28,
and \texttt{thC}, line 239.  The actual \texttt{thC} uses volume
$1$.  If instead one uses the separately defined \texttt{xiC},
its displayed positive volume contributes the additional factor
$\sqrt{\mathrm{vol}}$; it is not absorbed in (A15c).

For a tangent vector $u$, direct differentiation of (A15a) gives
\[
 \begin{split}
 d(U_xU_yF)_0(u)
  ={}&\exp\bigl(-\pi H(x,y)
              -\tfrac\pi2H(x,x)-\tfrac\pi2H(y,y)\bigr)\,V_u(x+y),\\
 V_u(w)={}&dF_w(u)-\pi H(u,w)F(w).
 \end{split}
 \tag{A15d}
\]
Indeed $H(z,x)+H(z,y)=H(z,x+y)$, and the constant term in
the two exponential factors is precisely the exponent displayed.
This proves (A15d) without a convention-dependent appeal to a
Heisenberg commutation rule.  Equivalently,
$U_xU_y=\exp(-\pi i\operatorname{Im}H(x,y))U_{x+y}$.

Fix tangent vectors $u_1,\ldots,u_g$.  The complete first-derivative
matrix for any finite list of actual arrays $\Xi_j$ is therefore
\[
 J_{rj}=\frac1{|S_A|}
  \sum_{x,y,s}\alpha(x)\alpha(y)b(y,s)^{-1}
       \Phi_j(s-x)
       e^{-\pi H(x,y)-\pi H(x,x)/2-\pi H(y,y)/2}
       V_{u_r}(x+y).
 \tag{A15e}
\]
Every index and scalar in this formula comes from the stated finite
transform.  In particular $\Xi_j(0)=0$ does not set $F(0)$ or the
translated vacuum values to zero.  A prescribed full-rank support
condition or right-$\mathrm{SL}$ covariance is retained inside
$\Phi_j$ throughout this formula.

There is a useful rigorous error certificate for (A15e).  Suppose
computed values of $F(w)$ and $dF_w(u_r)$ have respective error bounds
$\epsilon_0(w)$ and $\epsilon_r(w)$.  Since the semicharacters and
commutator factors have modulus one, the resulting error in $J_{rj}$
is at most
\[
 \frac1{|S_A|}\sum_{x,y,s}|\Phi_j(s-x)|
 e^{-\pi H(x+y,x+y)/2}
 \bigl(\epsilon_r(x+y)+\pi|H(u_r,x+y)|\epsilon_0(x+y)\bigr).
 \tag{A15f}
\]
The real part of the exponent in (A15e) is
$-\pi H(x+y,x+y)/2$, which proves this estimate.  The theta and
derivative tail bounds of Q2 and A12 can supply the errors here.
Thus a numerical nonzero-minor certificate need not omit any of the
finite half-pairing phases.

\paragraph{The exact restricted Fourier rank test.}
Fix all non-$T$ data, including an actually permitted primary tuple
from A14 and its allowed local indices.  At a $T$-prime $v$, let
$M_v=\operatorname{Mat}_{b\times a}(\kappa_v)$ in the source's residue
frames.  Integrality is included in the remaining factor, so on the
finite window in (A15e) the actual array has the form
\[
 C_j(t)\prod_{v\in\mathcal T}I_{v_v,u_v}(\operatorname{red}_v t).
 \tag{A15g}
\]
Here $C_j$ is the \emph{actual} product of other factors, including
the indicated finite auxiliary sums.  It is not replaced by a
constant.  Collecting all terms with a fixed tuple of reductions
defines a completely explicit weighted kernel
$\mathcal K_{rj}$ on $\prod_v M_v$, so that (A15e) becomes
\[
 J_{r,(j,I)}=\sum_{X\in\prod_v M_v}
                    \mathcal K_{rj}(X)\prod_v I_v(X_v).
 \tag{A15h}
\]
The collection operation simply sums the coefficient of
$C_j(t)$ in (A15e) over all triples $(x,y,s)$ with $t=s-x$ and
the specified reductions.  Terms off the required integral domain
are zero.  Its definition therefore involves only the literal
finite window, not an infinite summation or a changed support.

Choose nontrivial additive residue characters.  For a matrix $A$
put $e_A(X)=\psi(\operatorname{Tr}(A^tX))$.  A13 implies that the
single-prime incidence space is also exactly
\[
 \operatorname{span}\{e_A-1:\operatorname{rank}A=1\}.
 \tag{A15i}
\]
For completeness, if $A=v^tt$, then
\[
 e_A(X)-1
       =\sum_{u\ne0}\bigl(\psi(\textstyle\sum_j u_jt_j)-1\bigr)I_{v,u}(X).
 \tag{A15j}
\]
The value at $vX=0$ is zero, so the missing $u=0$ term causes no
error.  Conversely the Fourier support of each incidence function
lies in $\{v^tt:t\in\kappa^a\}$ and its value at zero is zero,
as proved in A13.  Fourier independence then proves (A15i).
The coefficient sum in (A15j) is at most $2(q^a-1)$.

Use the positive-sign finite transform
\[
 \widehat{\mathcal K}_{rj}(A_1,\ldots,A_f)
   =\sum_X\mathcal K_{rj}(X)\prod_v e_{A_v}(X_v).
\]
For a derivative covector $\ell$, put
$\mathcal K_{\ell j}=\sum_r\ell_r\mathcal K_{rj}$.
The actual incidence derivative matrix has full row rank if and
only if there is no nonzero $\ell$ for which
\[
 \sum_{S\subset\{1,\ldots,f\}}(-1)^{f-|S|}
       \widehat{\mathcal K}_{\ell j}(A^S)=0
 \quad\text{for all }j\text{ and all rank-one }A_v.
 \tag{A15k}
\]
Here $A^S_v=A_v$ for $v\in S$ and $A^S_v=0$ otherwise.
Indeed the left side is exactly the pairing in (A15h) with
$\prod_v(e_{A_v}-1)$, and these products span the tensor incidence
space.  This proves both directions of the criterion.  In the
single-prime case it says that
$\widehat{\mathcal K}_{\ell j}$ is constant on all matrices of rank
at most one.  For several primes it is the full mixed difference
(A15k), rather than the separate one-variable condition.
Equivalently, on this rank-one grid the transform is the alternating
sum of its proper coordinate faces.  The number of local basis
columns before multiplying the other factors is
\[
       \prod_{v\in\mathcal T}
       \frac{(q_v^a-1)(q_v^b-1)}{q_v-1}.
 \tag{A15l}
\]

When varying the primary twisting characters, the index $j$ runs
only over the quotient-character family actually constructed in
A14.  In particular the finite auxiliary sum in a \texttt{XiEntry}
is not silently replaced by one determinant character.  When passing
from raw decorated lattices to the archive's chosen class
representatives, include its stated equivariance scalar and its row
change of coordinates in (A15e).  The criterion then tests the
resulting actual matrix.  No enlarged list of class representatives
is assumed.

\paragraph{Normalization and a finite algebraic certificate.}
The literal differential convention is especially simple at degree
one.  In
\path{Formalism/LocalToGlobal/KernelLattices/ThDistC.lean},
lines 54--62, \texttt{jetC} is the undivided derivative multiplied by
$c_D(\tau)$ in each differentiated coordinate.  For
$r=(\tau,i,j)$ use
$u_r=c_D(\tau)\,\mathrm{coordX}(r)$ in (A15e).
Then the result is exactly its degree-one \texttt{jetC}; there is
no factorial.  A different rational Schr\"odinger chart requires
the actual row Jacobian retained in A12.

For fixed source and target lattices, divide the entries by the
common factor $\chi_F(N,N')$ specified in A13.  The fixed-$J$
arithmetic-jet construction there, with degree one included, puts
all entries for the finite actual family in one number field after
a further \texttt{Enl}.  This invokes the literal normalized-jet
theorem at its stated \texttt{AmbIn.Ok} hypotheses; it does not
deduce algebraicity of unnormalized analytic derivatives.
Changing one theta section into a linear combination still uses
its section comparison factor once.  A determinant of independent
first-derivative columns uses the corresponding product of factors,
as in A12.

Thus a nonzero minor of this \emph{normalized} matrix gives an
algebraic inverse on its image.  If a computed square submatrix
$J_0$ has least singular value at least $s>0$, then a target
derivative vector $w$ can be obtained with coefficient sum at most
$\sqrt g\,\|w\|/s$.  This follows from
$\|J_0^{-1}w\|_1\leq\sqrt g\|J_0^{-1}w\|_2$.
For an approximate matrix replace $s$ by
$\sigma_{\min}(J_{0,\mathrm{approx}})-\|\mathrm{error}\|$, provided
this is positive; (A15f) bounds the error.
These are concrete finite certificates.  An algebraic inverse of
this derivative matrix is still not, without a cohomological
identification, an algebraic correspondence on the ball variety.

\paragraph{Faithfulness and modulation in the precise finite block.}
There is a useful positive nonvanishing criterion before any derivative
is taken.  Suppose that $\Lambda,A,B$ are commensurable full lattices,
that all three are self-dual for $\operatorname{Im}H$, and that
$\alpha$ is a semicharacter on each.  Retain the actual
\texttt{GoodLat} conditions: in particular
\[
 \Lambda=(\Lambda\cap A)+(\Lambda\cap B),\qquad
 a\in A,\ b\in B,\ a+b\in\Lambda\ \Longrightarrow\ a\in\Lambda.
 \tag{A15u}
\]
These are the first two lattice conditions in
\path{Formalism/LocalToGlobal/PairingIndexBounds/ModelLatticeMoves/}
\texttt{GoodLat.lean}, lines 25--29; its last two conditions assert
the corresponding nondegeneracy of the cross pairing.
Let $F_0\ne0$ satisfy $\alpha(\lambda)U_\lambda F_0=F_0$.
Then (A15c) is injective on the finite space of functions
\[
                 \Phi:A/(A\cap B)\longrightarrow\C.
 \tag{A15v}
\]
In this statement $S_X$ represents $A/(\Lambda\cap A)$,
$S_A$ represents $(\Lambda\cap A)/(A\cap B)$, and $S_B$
represents $B/(\Lambda\cap B)$, precisely as in the source.
The inclusion $A\cap B\subset\Lambda$ follows from (A15u): an
element of $A\cap B$ pairs integrally with both summands of
$\Lambda$, hence belongs to its self-dual lattice.

Here is a full proof of injectivity.  For every $\lambda\in\Lambda$,
the nonzero entire function $U_xU_yF_0$ is an eigenvector of
$\alpha(\lambda)U_\lambda$, with character
\[
             e^{-2\pi i\operatorname{Im}H(\lambda,x+y)}.
\]
These characters are distinct as $(x,y)$ runs over $S_X\times S_B$.
Equality would imply $x+y-x'-y'\in\Lambda$ by self-duality;
(A15u) then gives $x-x'\in\Lambda$, forcing $x=x'$ and $y=y'$
by the representative conditions.  Eigenvectors with distinct
characters of an abelian group are linearly independent.  For
example, given a finite relation and a chosen character, successive
operators of the form $\rho(\lambda)-\chi(\lambda)$ eliminate
each of the other characters while leaving that chosen term
nonzero.  Thus $\Theta_\Phi=0$ forces $K_\Phi(x,y)=0$ for every
pair.

The pairing $b(y,s)$ between
$(\Lambda\cap A)/(A\cap B)$ and $B/(\Lambda\cap B)$ is perfect.
Its right kernel is exactly $\Lambda\cap B$ by the last
\texttt{GoodLat} condition.  Its left kernel is $A\cap B$, since
an element of $\Lambda\cap A$ pairing integrally with all of $B$
belongs to $B$ by self-duality.  Both groups are finite, so
nondegeneracy on both sides makes their orders equal and gives a
perfect finite-character pairing.  Finite Fourier inversion in
(A15b) now gives $\Phi(s-x)=0$ for every $s,x$.
These elements represent all of $A/(A\cap B)$: choose $x$
representing $-t$ modulo $\Lambda\cap A$, then $s$ representing
$t+x$ modulo $A\cap B$.  This proves (A15v).

There is also an exact operator identity on this whole block,
including an already present array of other local factors:
\[
       \Theta_{\,b(y,\cdot)\Phi}
                =\alpha(y)U_y\Theta_\Phi,
                  \qquad y\in B.
 \tag{A15w}
\]
To verify it without dropping a phase, multiplication of $\Phi(s-x)$
by $b(y,s-x)$ changes $K_\Phi(x,z)$ to
$b(y,-x)K_\Phi(x,z-y)$.
The coefficient $K_\Phi(x,z)$ is periodic in
$z\in B$ modulo $\Lambda\cap B$, because $s\in\Lambda\cap A$.
On $B$ the operators $W_y=\alpha(y)U_y$ form a representation:
the semicharacter identity and
$U_xU_y=e^{-\pi i\operatorname{Im}H(x,y)}U_{x+y}$ give
$W_xW_y=W_{x+y}$, since $\operatorname{Im}H(B,B)\subset\Z$.
Also $W_\lambda F_0=F_0$ for $\lambda\in\Lambda\cap B$.
Reindexing $z-y$ in $S_B$ therefore moves $W_y$ to the left of
$W_zF_0$.  Commuting it past $U_x$ contributes $b(y,x)$,
which cancels $b(y,-x)$.  This gives exactly (A15w).
It is the statement called \texttt{XiModulate} in
\texttt{XiIndCClass.lean}, lines 98--102; the argument here proves
the finite identity directly, rather than treating that
\texttt{Prop} declaration as a proof.

Consequently, for any collection $y_1,\ldots,y_f\in B$,
\[
 \Theta_{\,\Phi\prod_i(b(y_i,\cdot)-1)}
      =(W_{y_f}-1)\cdots(W_{y_1}-1)\Theta_\Phi.
 \tag{A15x}
\]
This proves the ordered character-operator prerequisite in C37
whenever the literal input has been placed in this finite block.
The representatives need not be bounded; C37 retains the lattice
characters introduced by reducing them modulo the baseline period
lattice.  The positive-sign finite character $e_A$ in (A15i) may
label $-y$ rather than $y$ under the source's negative-sign
$b(y,t)=\exp(-2\pi i\operatorname{Im}H(t,y))$.
Use that actual perfect-pairing label.  Negation preserves the
rank-one sample and does not change any spanning assertion.

In particular a nonzero remaining array gives a nonzero baseline
\emph{when its membership in (A15v) and these self-duality and
vacuum hypotheses have been checked}.  This is more specific than
an assertion about an arbitrary theta summation map.  To use one
fixed baseline while a new split prime varies, the added local
factor must be the unramified integral vacuum at that prime.
In the finite Fourier formula its integral support sets the new
$x$-index to zero, and its constant value on the new $S_A$ makes
the new Fourier index zero.  Thus this factor acts as the identity.
Equivalently, the orthogonal prime components of the finite
Heisenberg group factor by the Chinese remainder theorem, and
the new factor is its vacuum vector.  This statement keeps the
remaining lattices, section normalization and local factors fixed;
it does not assert that an arbitrary changing archive datum has a
fixed baseline.  For fixed lattices, the new sufficiently large
primes avoid their finite nonstandard set, so the added local
factor of the volume is one.  Any pre-existing positive volume
scalar and every later class-representative or row-frame factor
remain explicit, as in (A15e) and the normalization paragraph.

\paragraph{The fixed covering lattice and the new primary samples.}
The baseline in the preceding paragraph need not be a section on
the original principally polarized torus.  Let $M$ be its original
$\OO_E$-period module, with complex image $\Lambda$, and choose a
positive integer $N_0$ such that
\[
 \Lambda_0=N_0\Lambda\subset A_{\rm old}\cap B_{\rm old}.
 \tag{A15y}
\]
Such an integer exists because the old lattices are commensurable.
For $\lambda\in\Lambda_0$, the character $b(\lambda,t)$ is one
on $A_{\rm old}$, by self-duality of $A_{\rm old}$.
Equation (A15w) therefore proves
$\alpha(\lambda)U_\lambda F_{\rm base}=F_{\rm base}$.
Thus the actual fixed positive theta line for C37 is its line on
$V/\Lambda_0$.  No descent of $F_{\rm base}$ to $V/\Lambda$
is being assumed.

Here is the precise new-prime compatibility for the half-lattices
used above.  Exclude primes dividing $N_0$, the old levels, the
different, and the finitely many nonstandard lattice indices.
At a new split rational prime $q$, the source definitions
\texttt{halfAO}, \texttt{halfBO} in
\texttt{XiIndCClass.lean}, lines 181--191, give the local pairs
$q^{-1}M,qM$ at each selected prime $\mathfrak p$ and the
opposite pair at $\overline{\mathfrak p}$.
The baseline has the integral vacuum at these places.
The $\overline{\mathfrak p}$-primary torsion group
\[
 \overline{\mathfrak p}^{-1}(N_0M)/(N_0M)
\]
on $V/\Lambda_0$ maps isomorphically to the corresponding
primary group on $V/\Lambda$.
Indeed, the kernel of $V/\Lambda_0\to V/\Lambda$ has order
prime to $q$, so multiplication by $q$ is invertible on it.
This proves injectivity on $q$-torsion, and a lift of a
$q$-torsion point can be adjusted by this kernel to become
$q$-torsion.  The map is $\OO_E$-linear, so it preserves the
individual prime components.

One can use these very primary points in (A15w).
Represent a point by $y\in\overline{\mathfrak p}^{-1}(N_0M)$.
By the Chinese remainder theorem for the finite module
$N_0M/qN_0M$, add an element of $N_0M$ so that $y$ belongs
to $qM$ at $\mathfrak p$ and at every other prime above $q$
except $\overline{\mathfrak p}$.  At the old primes its
component remains in $N_0M\subset A_{\rm old}\cap B_{\rm old}$.
The resulting representative belongs to $B_{\rm new}$.
On the baseline's integral new-prime support its character
$b(y,t)$ is precisely the residue trace character at
$\mathfrak p$, with the negative sign and scalar different
factor already recorded in (A15p); all other prime components
are trivial.  Perfect local trace duality proves this assertion
coordinate by coordinate.  With the stipulated source and target
frames the labels are matrices, and replacing $M$ by $N_0M$
only multiplies their local labels by a nonzero scalar.  Thus the
rank-one labels remain rank one.  The same construction applies
independently to every selected prime above $q$.

Consequently the torsion samples needed by C37 are samples on
the fixed torus $V/\Lambda_0$, with its fixed baseline section.
Their trace-dual Fourier test is (A15p), now for the dual of
$\Lambda_0$; a fixed nonzero dual element reduces nontrivially
away from finitely many primes.  The representatives in C37 may
be unbounded, but their changes of lift are elements of
$\Lambda_0$, as required there.  A change merely by an element
of the larger lattice $\Lambda$ is not declared harmless.
This argument proves compatibility under the stated literal
half-lattice, integral-vacuum and tensor-frame hypotheses; it
does not replace those hypotheses for an unrelated finite
transform.

\paragraph{A positive theorem for one incidence factor.}
Here is a nonvanishing result for the restricted rank-one frequencies
themselves.  Let $V/\Lambda$ be a complex torus with a positive theta
line, Hermitian form $H$, and semicharacter $\alpha$.  Let $F\ne0$
be one of its holomorphic theta sections in the entire-function
frame used in (A15a).  Set
\[
       v(y)=d_z[U_yF](0)\in V^*,\qquad a_0=v(0).
 \tag{A15m}
\]
Although $v$ depends on a unitary frame, its norm and the norm of
each constant covector pairing are well defined on the torus.
More precisely, $\alpha(\lambda)U_\lambda F=F$ implies
\[
 U_{y+\lambda}F
    =\alpha(\lambda)^{-1}
       e^{\pi i\operatorname{Im}H(y,\lambda)}U_yF.
 \tag{A15n}
\]
The same scalar multiplies every component of $v$, so those
components are continuous sections of the same nontrivial theta
line in a unitary frame.

This parameter-space statement does not make every $U_y$ an
automorphism of the original theta line on $V/\Lambda$.
For a point with $Ny\in\Lambda$, $U_yF$ and $F$ have the same
automorphy law on $N\Lambda$, so they are sections together on
$V/N\Lambda$ of the pullback of the original theta line under
the identity-linear covering
$\pi_N:V/N\Lambda\longrightarrow V/\Lambda$.
Indeed the extra multiplier on $N\lambda$ is
$e^{2\pi i\operatorname{Im}H(y,N\lambda)}=1$.
In this covering's $z$ coordinate the derivative of $\pi_N$ is
the identity, and (A15m) is its literal derivative.  If the cover
is instead identified with the original torus by $z=Nw$, its
line is $[N]^*L$ and its first derivative is multiplied by $N$.
Only the explicitly normalized factor $N^{-1}$ cancels that
change of coordinate.  The fixed positive line in (A15o) is
the parameter line of $v(y)$; the output columns for growing
torsion levels live on these stated pullback lines.

For any fixed Hermitian norm on $V^*$, put
\[
 \delta_F=\min_{\|\ell\|=1}
     \int_{V/\Lambda}
       \bigl(|\ell v(y)|-|\ell a_0|\bigr)^2\,dy,
 \tag{A15o}
\]
where Haar measure is a probability measure.  Then
$\delta_F>0$.  To prove this, if the integrand has zero integral
it vanishes everywhere by continuity.  If its constant norm is
positive, the corresponding nowhere-zero section trivializes the
theta line topologically, contradicting its nonzero positive
first Chern class.  This contradiction can be read directly from
curvature: a global unit frame would make its curvature exact,
whereas its positive curvature has a positive integral against
the $(g-1)$st power of a K\"ahler form.  If the constant is zero,
the directional form of (A15d) gives, for the nonzero direction
$u$ corresponding to $\ell$,
\[
             dF_y(u)=\pi H(u,y)F(y)\quad\hbox{for all }y.
\]
Apply the antiholomorphic derivative in direction $u$.  Holomorphy
of $F$ gives $0=\pi H(u,u)F(y)$, which contradicts $F\ne0$ and
positive definiteness.  Thus the integral in (A15o) is positive
for every unit $\ell$.  It is continuous in $\ell$, and the unit
sphere is compact, proving the positive minimum.

Suppose the finite torsion groups at a sequence of allowed CM split
primes are labelled, in tensor-compatible residue frames, by
$\operatorname{Mat}_{b\times a}(\kappa)$, with $q=|\kappa|\to\infty$.
Write $y_A$ for the point labelled by $A$.  The compatibility required
here is the actual trace-pairing compatibility: for every torus
character $\gamma$ there is a residue matrix $X_\gamma$ with
\[
        \gamma(y_A)=\psi(\operatorname{Tr}(A^tX_\gamma)),
 \tag{A15p}
\]
and every fixed nonzero $\gamma$ has $X_\gamma\ne0$ at all but
finitely many of the primes.  This is the usual CM lattice
construction: the rational trace dual is an $E$-lattice, and a
nonzero element of it has nonzero reduction away from finitely many
primes.  At good split primes the trace pairing is perfect.
Left and right changes of the stipulated source and target frames,
and the scalar different factor, preserve rank one.  An arbitrary
unrecorded linear relabelling of matrix coordinates is not allowed
in (A15p).

The probability measure on the $y_A$ with $\operatorname{rank}A=1$
converges to Haar measure.  Indeed every such matrix is $uv^t$ with
$u,v\ne0$, in exactly $q-1$ ways.  If $r=\operatorname{rank}X$, direct
summation gives
\[
 \frac{\sum_{u\ne0,\ v\ne0}\psi(u^tXv)}
            {(q^b-1)(q^a-1)}
  =\frac{q^{a+b-r}-q^a-q^b+1}{(q^b-1)(q^a-1)}.
 \tag{A15q}
\]
For fixed $u$ the inner sum is $q^a-1$ if $X^tu=0$ and $-1$
otherwise; there are $q^{b-r}-1$ nonzero $u$ of the first kind.
This proves the formula.  For $r\ge1$ its absolute value is at
most $1/(q-1)$: the largest positive numerator occurs at $r=1$,
and the negative magnitude is at most
$\min(q^a,q^b)-1$, both bounded by the indicated denominator
divided by $q-1$.  Equations (A15p)--(A15q) give convergence on
each torus character, hence on trigonometric polynomials, and then
on continuous functions by uniform approximation.

For \emph{any} choices of phases $|\eta_A|=1$, the columns
\[
                 \eta_Av(y_A)-a_0,
                 \qquad \operatorname{rank}A=1,
 \tag{A15r}
\]
span $V^*$ for every sufficiently large such prime.  More precisely,
if $R_q=(q^a-1)(q^b-1)/(q-1)$ is their number, their matrix $J_q$
satisfies
\[
             J_qJ_q^*\ \ge\ R_q\delta_F/2\;\mathrm{Id}.
 \tag{A15s}
\]
The triangle inequality gives
$|\eta_A\ell v(y_A)-\ell a_0|
 \ge\bigl||\ell v(y_A)|-|\ell a_0|\bigr|$.
Equidistribution is uniform on the compact family of continuous
functions in (A15o): use a finite uniform approximation net in
$\ell$ and then the proved convergence for that finite net.
Their discrete integrals are therefore at least $\delta_F/2$,
uniformly in unit $\ell$ and in every phase.  This proves (A15s).
The minimum-norm right inverse gives coefficient sum at most
$\sqrt{2/\delta_F}\,\|w\|$ for a desired derivative vector $w$.
To express each difference by the literal incidence indicators,
multiply this bound by at most $2(q^a-1)$ using (A15j).

There is also a finite quantitative threshold.  Choose coordinates
on the real torus of dimension $d=2g$, and let $P_M$ be the product
Fej\'er kernel of frequency range $[-M,M]^d$.  Choose a uniform
bound $B\ge\sup_y\|v(y)\|$ and a bound $L$ for its covariant first
derivative.  The functions in (A15o) have bound $4B^2$ and
Lipschitz constant at most $4BL$.  Put
$\epsilon_M=\int\operatorname{dist}(t,0)P_M(t)\,dt$.
After excluding the finitely many primes where one of the nonzero
frequencies in this range reduces to zero, (A15q) shows that the
discrete average in (A15o) is at least
\[
 \delta_F-4BL\epsilon_M
       -\frac{4B^2((2M+1)^d-1)}{q-1}.
 \tag{A15t}
\]
Convolution by $P_M$ preserves the Haar average, approximates every
function in the family within $4BL\epsilon_M$, and has the
indicated finite Fourier support; each Fourier coefficient has
absolute value at most $4B^2$.  This proves (A15t), including its
constants.  Taking the last two terms at most $\delta_F/4$ each
is an explicit sufficient test for (A15s).  The constants concern
this fixed theta section; they do not bound an arbitrary varying
section or an unnormalized ball period.

\paragraph{Application to several factors.}
In the single-incidence Fourier model the exact finite transform
sends $e_A$ to a unit scalar times $U_{y_A}F$, and sends $1$ to
$F$.  Thus (A15r) is exactly its first-derivative column.  The
theorem proves that low Fourier rank alone is not an obstruction
to first-derivative spanning, while retaining every unit phase.
For the literal full \texttt{xiAt}, however, all the $T$-factors,
the other local factors, and the allowed primary tuple list remain
those in (A15g).  Several incidence factors give the alternating
sum (A15k).  C37 proves a simultaneous theorem for these ordered
differences, at one growing rational prime, even with unbounded
half-lattice representatives.  The exact identity (A15x) supplies
its operator comparison in the stated finite block.  The
injectivity proof supplies a nonzero baseline from a nonzero
array in that block, and the unramified-factor argument explains
when this is one fixed baseline as the new prime grows.
Together these results give full first-derivative rank whenever
those finite-block, fixed-data, and tensor-compatible residue
conditions are met for the sequence of actual inputs.

For a particular archive family one must still verify these
conditions for its selected non-$T$ local functions and permitted
prime choices; they are not consequences merely of a nonzero raw
array or a numerical signature match.  The exact finite
nonzero-minor test, normalized arithmetic field, and error bounds
remain (A15e)--(A15l) and (A15f).  In particular the literal
$v_T(u)$ has several places over the same rational $T(u)$;
C37 addresses that simultaneous situation without assuming
independent successive prime choices.  All these comparisons
retain the actual allowed array space of A13--A14.

In the requested Grothendieck framework these are comparisons of
finite theta values and their normalized derivatives, as in
\emph{cote44}, scans 7--8.  Turning their inverse into a map of
homological motives requires the separate cohomological
identification; it is not supplied by finite-dimensionality or by
the algebraic Lefschetz inverses discussed in \emph{cote16},
scans 43--46.


\subsection{THETA3-A16: the literal full host has an abelian-field restriction}\label{A16}

The distinction at the end of A11 is structural. An arbitrary
CM field cannot be enlarged into the field of the literal
\nolinkurl{Formalism.BasicNotions.Host}. This rules out one
overly broad matching procedure; it does not contradict the
general oscillator identities or the conditional constructions
A13--A15 on an existing full host.

\paragraph{The exact field restriction in the source.}
The pinned archive defines \texttt{HostRaw} in
\path{Formalism/SharpRestriction/RamificationFrobenius/}
\path{TotientGaloisClasses/ConjEMemHostFIff.lean},
lines 12--47. Its fields include
\[
 \operatorname{Gal}(E/\Q)\simeq\phi.\mathrm{GG},
 \qquad
 E\hookrightarrow\Q(\zeta_{\mathrm{QS}\,q\,d}),
 \tag{A16a}
\]
as well as the explicit assertion that $E/\Q$ is Galois.
The group $\phi.\mathrm{GG}$ is
$\{\pm1\}^2\times\phi.\mathrm{MA}$, where
$\phi.\mathrm{MA}$ is the multiplicative version of a finite
abelian additive group. The exact definitions are
\nolinkurl{ZsmulGeneratorEqZeroOfProjection.lean},
lines 15--23, 29--33, 87--89, in the same directory;
\path{Formalism/FrobeniusOrbits/Statement.lean},
line 27, defines the two sign factors;
\texttt{FrameFacts.lean}, lines 30--35, also expands the
two sign factors directly. Hence $E/\Q$ is abelian.
This conclusion follows directly from the displayed group
isomorphism; it does not require any converse to the
Kronecker--Weber theorem.

The record \texttt{HostAx} extends \texttt{HostRaw} in
\texttt{ConjEMemHostFIff.lean}, line 174. The stronger
\texttt{Host} extends \texttt{HostAx} in
\path{Formalism/LocalRingCardinality/DataHost.lean},
line 39. Thus every full host used by the stated
\texttt{D7}/\texttt{StdPair} construction has the restriction
(A16a). Enlarging the separate coefficient field
\texttt{LL}, as in A13--A14, does not change this host field.

\paragraph{A concrete CM field excluded by every such host.}
Put
\[
 F=\Q(\sqrt2),\qquad
 K=F(\alpha),\qquad
 \alpha^2=-(3+\sqrt2).
 \tag{A16b}
\]
Then $K$ is a quartic CM field which is not Galois over
$\Q$. Indeed, both real embeddings of $F$ send
$3+\sqrt2$ to a positive number. Its negative is not
a square in $F$, since it is negative at either real
embedding. Consequently $[K:F]=2$, and $K/F$ is a
totally imaginary quadratic extension of the totally
real field $F$. Also $\sqrt2=-\alpha^2-3$, so
$K=\Q(\alpha)$ and its minimal polynomial is
\[
                     X^4+6X^2+7.
 \tag{A16c}
\]

Fix the embedding with $\alpha=i\sqrt{3+\sqrt2}$.
Every element of $K$ is uniquely $a+b\alpha$ with
$a,b\in F\subset\R$, and it is real exactly when
$b=0$. Thus $K\cap\R=F$. If $K/\Q$ were normal,
it would contain another root $\beta$ of (A16c) with
$\beta^2=-(3-\sqrt2)$. Then
$(\alpha\beta)^2=7$, so $\sqrt7\in K\cap\R=F$.
This is impossible: if $\sqrt7=a+b\sqrt2$ with
$a,b\in\Q$, squaring gives $2ab=0$, and then either
$a^2=7$ or $b^2=7/2$, neither of which has a rational
solution. This proves the claimed nonnormality.

No $\Q$-embedding of $K$ into a full-host field exists.
Every subfield of a finite abelian Galois extension is
Galois over $\Q$: its corresponding subgroup is normal,
and the Galois correspondence applies. An embedded copy
of $K$ would therefore be Galois, contrary to the proof
above. More strongly, for every finite extension
$K'/K$ there is no $\Q$-embedding $K'\hookrightarrow E$
into a full-host field, since its restriction would
embed $K$.

The same obstruction applies if one starts with a
Galois CM field as in 032. For completeness, the normal
closure here is
\[
       \widetilde K=K(\sqrt7),\qquad
       \widetilde F=\Q(\sqrt2,\sqrt7).
 \tag{A16d}
\]
The field $\widetilde F$ is totally real of degree four;
$\widetilde K=\widetilde F(\alpha)$ is its totally
imaginary quadratic extension. The four roots
$\pm\alpha,\ \pm\sqrt7/\alpha$ of (A16c) all lie
in it, so it is the normal closure and is a Galois
CM field of degree eight. It is nonabelian, since it
contains the nonnormal subfield $K$. It cannot embed
in a full-host field, nor can any finite extension
containing it. The four-factor theorem in 032,
Theorem 2.5 and its reduction Proposition 2.6, printed p.8,
is stated for every Galois CM field; it is not
restricted to the abelian fields of (A16a).

\paragraph{What is refuted, and the remaining comparison.}
The following proposed universal step is refuted:
start with an arbitrary 032 CM field, retain its field
action by a field embedding, enlarge it finitely, and
obtain the literal full-host record of the recovered
arithmetic archive. The explicit counterexample is
(A16b), or the Galois example (A16d). This is stronger
than an unresolved search for auxiliary primes or
for a factorization of a different.

It does not refute A11's generic
\texttt{Geo.Tor.Host}, which has different fields and
hypotheses. It does not refute A13--A15, which start
with an existing full host, or B25's finite definite
period identity on the actual oscillator spaces.
It does not assert the impossibility of an unrelated
correspondence, a different field action, or a
generalized arithmetic theorem. The viable general
route is to use the field-independent proved
oscillator and finite-transform identities with
their own explicit hypotheses, and prove the needed
arithmetic estimates for the given CM field. The
literal archive theorem remains usable on its actual
abelian-field domain.

In the requested Grothendieck framework, an
$E$-linear correspondence must retain its specified
endomorphism-field action on the realization. A
change of coefficient field is distinct from a
change of that action. This is the realization and
correspondence setting of \emph{cote44}, scans 7--8;
the algebraic-projector discussion of \emph{cote16},
scans 43--46, supplies no embedding of a nonnormal
CM field into an abelian extension.


\subsection{THETA3-C37. Simultaneous rank-one differences of a positive theta section}
\label{C37}

This result supplies a positive analytic rank theorem for several
incidence factors at the same rational prime.  It allows the actual
torsion representatives to be unbounded.  Its application to a particular
finite transform still requires that transform to have the ordered
operator form stated below and to have one fixed nonzero baseline theta
section.  The graph-subgroup argument was developed jointly with the
independent review of A15; none of the phase factors is discarded.

\paragraph{The analytic and sampling hypotheses.}
Let $T=V/\Lambda$ be a complex torus of positive dimension, with a
positive definite Hermitian form $H$, linear in its first argument.
Write $E=\operatorname{Im}H$ and assume $E(\Lambda,\Lambda)\subset\Z$.
Let $L$ be the positive Appell--Humbert line with semicharacter
$\alpha$, and let $F\ne0$ be an entire theta function with the
normalization
\[
 \alpha(\lambda)U_\lambda F=F,
 \qquad
 U_yF(z)=F(z+y)
       \exp\bigl(-\pi H(z,y)-\tfrac\pi2H(y,y)\bigr).
 \tag{C37a}
\]
Fix an integer $f\ge1$.  For each sufficiently large index $q$, and
each $i=1,\ldots,f$, let $S_{i,q}$ be a finite nonempty sample set,
with a map $s\mapsto t_{i,q}(s)\in T$ and uniform probability measure
$\mu_{i,q}$ on these points, counting multiplicities.  Suppose that
there are numbers $\epsilon_q\to0$ with the following two properties.
For every character $\gamma\in\Lambda^*=\Hom(\Lambda,\Z)$,
\[
 \widehat\mu_{i,q}(\gamma)=1
 \quad\hbox{or}\quad
 |\widehat\mu_{i,q}(\gamma)|\le\epsilon_q;
 \tag{C37b}
\]
the first alternative occurs precisely when the character is trivial
on the underlying finite torsion group.  For each fixed nonzero
$\gamma$, that first alternative does not occur for all sufficiently
large $q$.  The latter condition implies
$\mu_{i,q}\to m_T$, the probability Haar measure, by approximation
with trigonometric polynomials.

Choose an arbitrary lift $Y_{i,q}(s)\in V$ of every sample point and
an arbitrary scalar $\eta_{i,q}(s)$ of modulus one.  The lifts need
not lie in a bounded set.  Consider the first-derivative columns
\[
 w_q(s_1,\ldots,s_f)=
 d_z\left[
   (\eta_{f,q}(s_f)U_{Y_{f,q}(s_f)}-1)\cdots
   (\eta_{1,q}(s_1)U_{Y_{1,q}(s_1)}-1)F
       \right]_{z=0}\in V^*.
 \tag{C37c}
\]
Each phase here depends on its own factor.  Bilinear phases produced
by reducing the lifts will be retained in the proof.

\paragraph{The simultaneous spanning theorem.}
For every sufficiently large $q$, the columns (C37c) span $V^*$.
More precisely, choose points $t_2,\ldots,t_f\in V$ such that
\[
 \sum_{i\in J}t_i-\sum_{i\in K}t_i\notin\Lambda\otimes\Q
 \quad\text{whenever }J\ne K\subset\{2,\ldots,f\}.
 \tag{C37d}
\]
Such choices exist: the prohibited set is a countable union of
proper real affine subspaces in $V^{f-1}$.  Equidistribution permits
choices $s_{i,q}\in S_{i,q}$, for $i\ge2$, whose images converge to
$t_i$ modulo $\Lambda$.  Then there is a constant $c>0$ such that,
for all sufficiently large $q$, the matrix $J_q$ formed from the
columns $w_q(s_1,s_{2,q},\ldots,s_{f,q})$ satisfies
\[
 J_qJ_q^*\ \ge\ c\,|S_{1,q}|\,\operatorname{Id}_{V^*}.
 \tag{C37e}
\]
Here the Hermitian norm on $V^*$ is fixed.  The conclusion is
uniform in all choices of lifts and phases in (C37c).  The proof
gives an existence threshold and a positive constant, rather than
an explicit numerical threshold for these several-factor data.
For $f=1$, A15 gives the explicit constant $\delta_F/2$ and its
Fej\'er threshold.

We prove (C37e) by contradiction, recording the two geometric
ingredients in full.

\paragraph{Reduction of the lifts and its exact extra characters.}
For $\lambda\in\Lambda$, direct use of (C37a) gives
\[
 U_{y+\lambda}F
  =\alpha(\lambda)^{-1}e^{\pi iE(y,\lambda)}U_yF.
 \tag{C37f}
\]
The multiplier is independent of $z$.  Thus for every fixed linear
operator $P$ on these entire functions and tangent vector $u$,
$d(PU_yF)_0(u)$ is a smooth section of the same theta line in the
parameter $y$; its norm is well defined on $T$.  In particular,
it is not being treated as a holomorphic scalar function of $y$.

For a translation on the left the exact identity is
\[
 U_{t+\lambda}U_yF
   =\alpha(\lambda)^{-1}e^{\pi iE(t,\lambda)}
       e^{-2\pi iE(\lambda,y)}U_tU_yF.
 \tag{C37g}
\]
Indeed $U_xU_y=e^{-\pi iE(x,y)}U_{x+y}$, and (C37f) applied
to $x+y$ proves (C37g).  The last factor is a character of $T$,
because $E(\lambda,\Lambda)\subset\Z$.  This character, which
can vary with $q$, is essential.

For the chosen samples in factors $i\ge2$, choose bounded lifts
$t_{i,q}\to t_i$ and write their actual lifts as
$Y_{i,q}=t_{i,q}+\lambda_{i,q}$, $\lambda_{i,q}\in\Lambda$.
Put $t_{J,q}=\sum_{i\in J}t_{i,q}$.  Expanding the ordered product
of the last $f-1$ differences and applying (C37g) gives, for every
tangent vector $u$,
\[
 d\left[
  (\eta_fU_{Y_f}-1)\cdots(\eta_2U_{Y_2}-1)U_yF
       \right]_0(u)
   =\sum_{J\subset\{2,\ldots,f\}}
         c_{J,q}\prod_{i\in J}\chi_{i,q}(y)\,
               d(U_{t_{J,q}}U_yF)_0(u),
 \tag{C37h}
\]
where $|c_{J,q}|=1$ and
\[
       \chi_{i,q}(y)=e^{-2\pi iE(\lambda_{i,q},y)}.
 \tag{C37i}
\]
All signs, the factors $\eta_i$, the Heisenberg cocycles, and the
constant multipliers in (C37g) are included in $c_{J,q}$.
There is exactly one summand per subset, so none of these
coefficients vanishes.

For a first-factor sample, (C37f) removes its lattice displacement
before the fixed left operator is applied.  Consequently its
pairing with $u$ has the form
\[
       \zeta_{s,q}A_q(t_{1,q}(s),\chi_{2,q}(t_{1,q}(s)),\ldots)
                       -A_q(0,1,\ldots,1),
       \qquad |\zeta_{s,q}|=1,
 \tag{C37j}
\]
where $A_q(y,z_2,\ldots,z_f)$ is the right side of (C37h) with
$\chi_{i,q}(y)$ replaced by $z_i$.  This is a section of the
pullback of $L$ to $T\times(S^1)^{f-1}$.  Taking absolute values
in (C37j) eliminates only the first scalar $\zeta_{s,q}$; it does
not eliminate the characters (C37i).

\paragraph{Haar limits on graph subgroups.}
Push $\mu_{1,q}$ forward under
\[
       y\longmapsto(y,\chi_{2,q}(y),\ldots,\chi_{f,q}(y))
          \in T\times(S^1)^{f-1},
\]
and denote the resulting probability measure by $\nu_q$.
Every subsequential weak limit of these measures is Haar measure
on a closed subgroup $G\subset T\times(S^1)^{f-1}$ whose
projection to $T$ is surjective.  Here is an elementary proof.
For a character $(\gamma,k_2,\ldots,k_f)$ of this product, its
integral against $\nu_q$ is a Fourier coefficient of $\mu_{1,q}$
at the possibly moving character
$\gamma+\sum_i k_i\chi_{i,q}$.  By (C37b), every limit of these
coefficients is either $0$ or $1$.  If the integral of a
unit-modulus character against a probability measure equals $1$,
that character equals $1$ almost everywhere.  The characters with
integral $1$ therefore form a subgroup $D$ of the character group.
The limiting measure and Haar measure on the annihilator of $D$
have identical Fourier coefficients, hence coincide, by density
of trigonometric polynomials.  This proves the subgroup assertion.
For characters with $k_i=0$, the fixed-character part of (C37b)
shows that the marginal is $m_T$.  The projection of $G$ is
closed, and therefore is all of $T$.

For later use, any such projection $G\to T$ has a section after
a finite identity-linear covering of $T$.  Indeed $G$ is a
compact abelian Lie group.  Its identity component $G^0$ still
surjects onto $T$: the quotient by the image of $G^0$ is both
connected and finite.  The surjection of the rational fundamental
groups of these tori has a rational linear right inverse.
Clearing its denominators gives a homomorphism
\[
 j:V/N\Lambda\longrightarrow G^0,
       \operatorname{pr}_Tj(y)=y\pmod\Lambda
 \tag{C37k}
\]
for some positive integer $N$.  Each circle coordinate of this
homomorphism has the form $e^{2\pi i\ell_i(y)}$, where
$\ell_i\in\Lambda^*\otimes\Q$.  This argument does not assume
that $G$ is connected or that its projection has a section
without a finite covering.

\paragraph{Nonvanishing on every such subgroup.}
Let $u\ne0$, let all $c_J$ have modulus one, and set
\[
 A(y,z)=\sum_J c_J\prod_{i\in J}z_i\,
                       d(U_{t_J}U_yF)_0(u),
       \qquad t_J=\sum_{i\in J}t_i.
 \tag{C37l}
\]
For every closed subgroup $G$ projecting onto $T$, the norm of
this section is not constant on $G$.

First we prove that its restriction along (C37k) is not the zero
section.  Since the rational alternating form $E$ is
nondegenerate, for every $\ell_i\in\Lambda^*\otimes\Q$ there
is $s_i\in\Lambda\otimes\Q$ with $\ell_i(y)=E(s_i,y)$.
Put $s_J=\sum_{i\in J}s_i$.  Differentiating (C37a) gives
\[
 d(U_tU_yF)_0(u)=
 e^{-\pi H(y,y)/2-\pi H(t,t)/2-\pi H(t,y)}
 \bigl(dF_{y+t}(u)-\pi H(u,y+t)F(y+t)\bigr).
 \tag{C37m}
\]
The equality
$2\pi iE(s,y)=\pi H(s,y)-\pi H(y,s)$ shows that, after
pullback along $j$ and multiplication by $e^{\pi H(y,y)/2}$,
the summands have antiholomorphic exponential slopes
$-\pi H(t_J-s_J,y)$.  These slopes are distinct: equality for
$J\ne K$ would contradict (C37d).

For completeness, the required exponential independence can be
proved by differential operators.  Choose a direction $v\in V$
such that $H(u,v)\ne0$ and all numbers
\[
           \rho_J=-\pi H(t_J-s_J,v)
\]
are distinct; only finitely many proper complex hyperplanes
are excluded.  Write $D=\partial_{\bar v}$.  Each summand is
an antiholomorphic exponential times a polynomial of degree at
most one in the antiholomorphic variables, with holomorphic
coefficients.  Thus $(D-\rho_J)^2$ kills that summand.
Its image under $D-\rho_J$ is the nonzero real-analytic function
\[
 -\pi H(u,v)c_J
 e^{-\pi H(t_J,t_J)/2-\pi H(y,s_J)}
 e^{-\pi H(t_J-s_J,y)}F(y+t_J).
\]
If the sum vanished identically, applying
$(D-\rho_J)\prod_{K\ne J}(D-\rho_K)^2$ would leave this
nonzero function multiplied by
$\prod_{K\ne J}(\rho_J-\rho_K)^2$, a contradiction.
This proves nonvanishing without a genericity assumption on $F$.

If the norm in (C37l) were constant on $G$, its pullback along
$j$ would have constant norm.  A zero constant contradicts the
nonvanishing just proved.  A positive constant supplies a
nowhere-zero smooth section of the pullback of $L$ to
$V/N\Lambda$, trivializing that line topologically.  This is
impossible: its positive curvature has nonzero top intersection,
whereas a topologically trivial line has zero first Chern class.
This proves the assertion about the norm on $G$.

\paragraph{Completion of the uniform rank proof.}
If (C37e) failed for arbitrarily large $q$ with no positive lower
bound, there would be a sequence of unit tangent vectors $u_q$
for which
\[
  \frac1{|S_{1,q}|}\sum_{s\in S_{1,q}}
       |\langle u_q,w_q(s,s_{2,q},\ldots,s_{f,q})\rangle|^2
                  \longrightarrow0.
 \tag{C37n}
\]
Pass to subsequences so that $u_q\to u$ is a unit vector,
every $c_{J,q}\to c_J$, and $\nu_q\to m_G$ as above.
The bounded translations already satisfy $t_{J,q}\to t_J$.
In local unitary frames the corresponding sections $A_q$ converge
uniformly to (C37l); their transition functions are all those of
the same line $\operatorname{pr}_T^*L$, by (C37f).  Thus the
squared norm-difference functions converge uniformly as well.
The inequality $|\zeta a-b|\ge\bigl||a|-|b|\bigr|$ and
(C37j)--(C37n) imply
\[
       \int_G\bigl(|A(y,z)|-|A(0,1,\ldots,1)|\bigr)^2\,dm_G=0.
 \tag{C37o}
\]
Continuity implies that the integrand vanishes everywhere on
$G$, contrary to the preceding paragraph.  This proves (C37e),
including uniformity over arbitrary lifts and phases.  A
minimum-norm right inverse of $J_q$ has coefficient-sum bound
$c^{-1/2}\|w\|$ for a desired first derivative $w$, by
Cauchy--Schwarz and (C37e).

\paragraph{The exact CM rank-one samples.}
Here are sufficient arithmetic hypotheses for (C37b).  Let
$\Lambda$ be a full lattice over an order in a CM field $E_0$,
with maximal order and perfect trace pairing at the primes under
consideration.  For a prime ideal $\mathfrak p$, the ideal-torsion
group is $\mathfrak p^{-1}\Lambda/\Lambda$.  Define the action
on the rational trace dual by $(a\ell)(x)=\ell(ax)$.  Its exact
annihilator in $\Lambda^*$ is
\[
 \{\ell:\ell(\mathfrak p^{-1}\Lambda)\subset\Z\}
                        =\mathfrak p\Lambda^*.
 \tag{C37p}
\]
This follows locally by dualizing the invertible ideal
$\mathfrak p^{-1}$, and away from $\mathfrak p$ both lattices
are unchanged.  Under an identification using a Hermitian
polarization, conjugate ideals may appear instead; (C37p) uses
the stated contragredient trace-dual convention.  A fixed nonzero
dual element lies in $\mathfrak p\Lambda^*$ for only finitely
many primes, as is seen from any nonzero coordinate in a finite
generating set after clearing denominators.

Assume the actual residue frames identify this torsion group with
$\operatorname{Mat}_{b\times a}(\mathbb F_q)$ and its dual
pairing with $\psi(\operatorname{Tr}(A^tX))$.  For the uniform
sample of nonzero rank-one matrices, a direct sum over
$A=uv^t$, $u,v\ne0$, gives
\[
 \widehat\mu(X)=
 \frac{q^{a+b-r}-q^a-q^b+1}{(q^a-1)(q^b-1)},
                  \qquad r=\operatorname{rank}X.
 \tag{C37q}
\]
Every rank-one matrix has $q-1$ such presentations.  If $r=0$
this coefficient is $1$; if $r\ge1$ its absolute value is at
most $1/(q-1)$.  The latter follows by counting the
$q^{b-r}-1$ nonzero vectors in the left kernel and using that
the inner additive-character sum is $q^a-1$ there and $-1$
elsewhere.  Thus (C37b) holds with $\epsilon_q=1/(q-1)$,
including for the moving characters in the graph proof.
Equation (C37p) proves the required fixed-character condition.
Different selected ideals above the same rational prime are
allowed; the proof never chooses the prime for one factor after
fixing a prime for another factor.

\paragraph{Scope for the literal finite transform.}
If the finite Fourier transform sends each factor character to
$T_y=\alpha(y)U_y$ with $y$ in a common half lattice $B$ and
$E(B,B)\subset\Z$, then the semicharacter relation and (C37a)
give $T_xT_y=T_{x+y}$.  Hence a product of character differences
has exactly the ordered form (C37c).  The rank-one incidence
identity (A15j) then converts the columns back to the permitted
incidence span; it does not enlarge that span.  Source and target
residue frames and the different factor must be the actual
tensor-compatible ones, as in (A15p).

This conditional application retains two exact checks.  The
remaining local factors must give a fixed nonzero positive theta
section $F$, and the actual transform must have the asserted
ordered character-operator realization.  A nonzero raw finite
array is not a proof of the first check.  If the output is viewed
on $V/N\Lambda$, the coordinate is still the $z$ of (C37a).
Identifying that cover with $V/\Lambda$ by $z=Nw$ multiplies
first derivatives by $N$, exactly as in A15.  No such factor,
canonical theta period, arithmetic normalization, or later
cohomological comparison is absorbed in the analytic theorem.


\section{Definite theta periods and the exact iterated seesaw}

There are now two precise global comparisons. First, an orthogonal decomposition of the rank-three Hermitian space gives a totally definite rank-two subgroup. Restriction to that subgroup is a genuine global unitary Bessel period, and the theta tensor rule gives an exact seesaw identity; this is proved in THETA3-B11--B12. Second, the recovered Alp\"oge source supplies a common normalized functional for an entire family of iterated lifts. Its exact statement is \texttt{MainX}, and its source witness and elementary scalar consequence are recorded in THETA3-B13. Identifying the particular cohomological inputs with that family, including its finite lattices and normalization, remains an instantiation problem.

The inner-form calculations below describe a different possible route. Their obstructions do not obstruct the definite-subgroup construction. All new proofs have been read in full for the targeted checks recorded in the accompanying event files. Source statements used as background are distinguished from these proofs. No Lean compilation, complete dependency audit, or validation of the cited RH theorem is claimed or attempted.

\subsection{THETA3-B0: what the supplied and recovered sources establish}\label{B0}

The attached four-page Alp\"oge note gives a timeline. Step 4 mentions determinant positivity, step 7 GGP periods and interlacing, step 8 iterated theta lifts on totally definite unitary groups, and step 12 a seesaw producing a Bessel period independent of the Hecke characters. The note does not specify the integral, test vectors, measures, or comparison with differential forms. The attached source map is explicitly a text-parsing report, with no Lean build. Its paths are useful search locators; they are not substitutes for the definitions at those paths. The attached \texttt{Main.lean} and \texttt{Statements.lean} contain the final analytic declarations, not the theta-period definitions. This is an evidence finding, not a refutation of the intended comparison.

During this run, the named definitions and several proof bodies were recovered from the secondary public mirror \path{CocoSgt/riemann-hypothesis-lean}, pinned to commit
\begin{center}\small\nolinkurl{7cbf7a0c32fca1badf05931879a37c3f8a2b03b4}.\end{center}
The commit date is 20 November 2026. The recovered \texttt{Main.lean} and \texttt{Statements.lean} are byte-for-byte equal to the attachments. A manifest records the pinned URL, SHA256, Git blob hash, size, and retrieval time of each fetched file. This is partial identity evidence; it does not assert a complete primary-archive hash match. Consequently the initial missing-definition gate has been reduced to the exact mathematical matching problem in B13.

For reference, the 032 setting is a Galois CM extension $E/F$, with $[F:\Q]\geq2$, and an $E$-Hermitian space $V$ of dimension three. At one real place $\mu$ it has signature $(2,1)$, and elsewhere signature $(3,0)$. The compact ball surface is attached to $\SU(V)$ at sufficiently small principal level. For multiplicity space $P$ of dimension $k=1$ or $2$, the alternating form is
\[
 \psi_P=\Tr_{E/F}\bigl(\delta h_V\otimes h_P\bigr),
 \qquad \overline\delta=-\delta.
\]
At the active place the positive Lagrangian is
\[
 \mathcal F_P=(L\otimes P_\mu)\oplus(Q^*\otimes P_\mu^*),
 \qquad Q=V_\mu/L.
\]
The determinant square root and projected Taylor tensors give $\Omega^1$ and $K_X$; see 032, Lemma 4.1 and equations (4.5)--(4.8), pp.14--16. Proposition 4.6 and Corollary 4.7 vary the multiplicity-space data at compact places while keeping $V$ fixed. This operation does not change the signature of $V$.

\subsection{THETA3-B1: coherent rank-three definite neighbors}\label{B1}

\textbf{Statement.} There is a global Hermitian space $V^-$ with signature $(0,3)$ at $\mu$ and $(3,0)$ elsewhere, locally isometric to $V$ at every finite place. There is no everywhere positive definite $V^+$ with all those finite local isometry classes. For any finite nonsplit place $v_0$, there is an everywhere positive definite $V^{(v_0)}$ locally isometric to $V$ outside $\{\mu,v_0\}$, with the opposite determinant norm class at $v_0$.

\textbf{Proof.} Let $\epsilon_v$ be the local quadratic character of the determinant. At a real place it is $(-1)^{q_v}$, where $q_v$ is the negative index. A coherent collection has $\prod_v\epsilon_v=1$. Local dimension and determinant classify Hermitian spaces at finite places; together with the specified real signatures, the product condition gives global existence by Landherr's theorem. An exact source is B.H. Gross, \emph{Incoherent definite spaces and Shimura varieties}, 25 February 2018, \S3, Theorem 3.1, pp.6--7.

Changing the negative index at $\mu$ from one to three preserves $\epsilon_\mu$, so all finite data can be retained. Changing it from one to zero reverses $\epsilon_\mu$, contradicting the product condition if all other local data are retained. Reversing the determinant class at one nonsplit finite place reverses one more sign and restores the product condition. Landherr's theorem gives the asserted global spaces. \hfill$\square$

This is a statement about Hermitian isometry data and oscillator embeddings. Multiplying an odd-dimensional Hermitian form by a scalar can leave its abstract unitary group unchanged. Such a group identification does not identify the dual-pair embedding, splitting, or test vector.

\subsection{THETA3-B2: the scalar type distinguishes the definite neighbors}\label{B2}

\textbf{Statement.} Keep the active-place $\delta$ and multiplicity-space sign used in 032, and put $r=t^2$ on the double cover of the scalar circle of $\U(P_\mu)$. In the covariant Taylor model of 032, the degree-$k$ cohomological tensor has weight $t^{-3k}$. For the negative definite replacement in B1 the polynomial model
\[
 \bigoplus_{j\geq0}(\det\mathcal F^-)^{1/2}\otimes
                    \operatorname{Sym}^j\mathcal F^-
\]
has only weights $t^{3k+2j}$. The required weight is absent. For the positive definite replacement, its vacuum has weight $t^{-3k}$.

\textbf{Proof.} The scalar $r$ acts by $r$ on the $k$-dimensional summand $L\otimes P_\mu$ and by $r^{-1}$ on the $2k$-dimensional summand $Q^*\otimes P_\mu^*$. Hence $\det\mathcal F_P$ has weight $r^{-k}$, and its chosen square root has weight $t^{-k}$. The projected degree-$k$ tensor contributes $r^{-k}=t^{-2k}$, giving $t^{-3k}$.

Write $\delta_\mu=ia$, with $a\in\R^*$. The prescribed sign of $h_P$ is $-\operatorname{sgn}(a)$, as in 032, Lemma 4.1. When $h_V$ is negative definite, $a h_Vh_P$ is positive definite. The positive Lagrangian is then the nonconjugate summand of dimension $3k$, on which the scalar acts by $r$. Its half determinant contributes $t^{3k}$ and polynomial degree $j$ contributes $t^{2j}$. None gives $t^{-3k}$. When $h_V$ is positive definite, the positive Lagrangian is the conjugate summand, with scalar $r^{-1}$; its vacuum therefore has weight $t^{-3k}$. \hfill$\square$

The actual infinite Schr\"odinger coefficient vectors carry the dual action in 032, equation (4.13). Inverting both displayed spectra preserves both conclusions. A common honest character twist also preserves the mismatch. The same names for two splitting characters do not establish the same correction from the bare metaplectic action. This computation does not assert that an unspecified global theta packet is empty. Positive definite vacuum compatibility is necessary for this particular tensor comparison, not a sufficient identification of periods.

\subsection{THETA3-B3: an explicit local epsilon obstruction}\label{B3}

Fix a nonsplit quadratic local extension $E/F$, an additive character, $\delta$, a skew-Hermitian line $W$, a splitting character $\chi_W$, and a character $\alpha$ of the norm-one torus. If $H_+$ and $H_-$ are the two Hermitian lines of opposite determinant classes, then exactly one of the corresponding lifts of $\alpha$ to $W$ is nonzero.

Indeed, Borade--Franzel--Girsch--Yao--Yu--Zelingher, \emph{On tori periods of Weil representations of unitary groups}, Selecta Math. 31 (2025), article 49, Theorem 3.5, gives the condition
\[
 \epsilon_{E/F}(\chi_W^{-1}\alpha_E,\psi,\delta)
       =\epsilon(H)\epsilon_\delta(W),
 \qquad \alpha_E(x)=\alpha(x/\overline x).
\]
The left side is fixed and the two right sides are opposite signs, so exactly one equality holds. This proves the assertion as an application of that theorem. The source's \S\S3.2.1 and 4.2.1 state the splitting compatibilities before the seesaw identities. The example shows why local signs cannot be suppressed merely because a global parameter or a branching dimension agrees. It does not identify every Alp\"oge Bessel period with this particular torus period.

\subsection{THETA3-B4: normalization is indispensable for a numerical bound}\label{B4}

\textbf{Statement and proof.} Nonvanishing alone gives no positive lower bound independent of scalar normalization. If a nonzero theta input is replaced by $\varepsilon$ times itself, with $\varepsilon>0$, the resulting sesquilinear norm or simultaneous mixed pairing is multiplied by $\varepsilon^2$. It remains nonzero and tends to zero. Likewise, a nonzero constant determinant-line trivialization can be rescaled. Thus every numerical lower bound must specify the input norms and period conventions. \hfill$\square$

\subsection{THETA3-B5: a quantitative compact projection certificate}\label{B5}

\textbf{Statement.} Let $M$ be an $N$-dimensional unitary representation of a compact group $H$, with probability Haar measure. Let
\[
 P=\int_H\rho(h)\,dh
\]
have rank $r>0$. Some $r\times r$ principal minor of $P$ is at least $\binom Nr^{-1}$. In particular, when $r=1$, a normalized invariant functional has squared value at least $1/N$ on some vector of any orthonormal basis.

\textbf{Proof.} Haar invariance gives $P^2=P$; unitarity and invariance under inversion give $P^*=P$. Thus $P$ is the orthogonal projection onto $M^H$. Choose an $N\times r$ matrix $U$ whose columns form an orthonormal basis of that subspace. Then $P=UU^*$ and $U^*U=I_r$. Cauchy--Binet gives
\[
 1=\det(U^*U)=\sum_{|I|=r}|\det U_I|^2.
\]
There are $\binom Nr$ terms. For some $I$, its term is at least the reciprocal of this number, and that term is $\det P_{I,I}$. For $r=1$ these are exactly the squared coefficients of a normalized invariant vector, or equivalently of its normalized invariant functional. \hfill$\square$

For compact branching, take $M=\pi\otimes\sigma^\vee$. A multiplicity-one statement then gives a selected normalized period of squared size at least $1/(\dim\pi\dim\sigma)$. This does not set an unnormalized global automorphic integral equal to one.

\subsection{THETA3-B6: iteration requires an overlap estimate}\label{B6}

\textbf{Statement.} Positivity and nonzero eigenvalues of individual stages do not bound their composite. If instead
\[
 \|T_jw\|\geq b_j\|w\|,\qquad b_j>0,
\]
holds on the image of $T_{j-1}\cdots T_1$, then
\[
 \|T_m\cdots T_1v\|\geq\left(\prod_{j=1}^m b_j\right)\|v\|.
\]

\textbf{Proof.} In $\C^2$, let $P$ project onto $e_1$ and $Q_\theta$ onto $(\cos\theta,\sin\theta)$. Both have rank one and nonzero eigenvalue one, but $\|Q_\theta P\|=|\cos\theta|$, which can be arbitrarily small or zero. Under the stated additional hypothesis, apply the bound for $T_1$, then that for $T_2$ to $T_1v$, and continue. Each argument lies in the specified image, and multiplying the inequalities proves the result. \hfill$\square$

Thus multiplicity gained at each theta-lift stage does not, by itself, give a lower bound for a fixed composite period.

\subsection{THETA3-B7: an arithmetic Gram-matrix lower bound}\label{B7}

\textbf{Statement.} Let $K$ be CM with real subfield $K_0$ of degree $d$. Let $G\in M_r(K)$ be Hermitian and positive definite at every embedding $\sigma:K_0\hookrightarrow\R$. Suppose $DG$ has algebraic-integer entries for a positive integer $D$, and
\[
 \operatorname{tr}(G_\sigma)\leq T_\sigma,\qquad T_\sigma>0.
\]
For a chosen $\sigma_0$ and $r\geq2$,
\[
 \lambda_{\min}(G_{\sigma_0})\geq
 D^{-rd}\prod_{\sigma\ne\sigma_0}\left(\frac r{T_\sigma}\right)^r
       \left(\frac{r-1}{T_{\sigma_0}}\right)^{r-1}.
 \tag{B7}
\]
For $r=1$, omit the last factor. The right side also bounds
$a^*G_{\sigma_0}a/\|a\|_2^2$ from below for nonzero $a$.

\textbf{Proof.} The determinant $\Delta=\det G$ lies in $K_0$. The number $D^r\Delta$ is a totally positive algebraic integer, so its norm to $\Q$ is a positive integer and at least one. Consequently
\[
 \prod_\sigma\det G_\sigma\geq D^{-rd}.
\]
For $\sigma\ne\sigma_0$, the arithmetic--geometric mean inequality gives
$\det G_\sigma\leq(T_\sigma/r)^r$. Divide to obtain the corresponding lower bound for $\det G_{\sigma_0}$. The product of the other $r-1$ eigenvalues at $\sigma_0$ is at most $(T_{\sigma_0}/(r-1))^{r-1}$, so division gives (B7). When $r=1$ the determinant is the single eigenvalue. The final assertion follows from the spectral theorem. \hfill$\square$

The bounds at the other embeddings are essential: the totally positive unit $(3-2\sqrt2)^m$ tends to zero at one real embedding while its conjugate tends to infinity. On a finite definite adelic module, algebraic-valued vectors and rational mass weights can supply such a matrix once the field, denominator, and conjugate bounds are specified. For cohomological periods one must likewise identify an algebraic normalization and the relevant positivity at all embeddings; arbitrary complex theta coordinates do not have these properties automatically.

\subsection{THETA3-B8: a conditional two-place transfer inequality}\label{B8}

Suppose compatible normalized positive period formulas have been established:
\[
 I_{\mathrm{ball}}=C_{\mathrm{ball}}L\prod_{v\in S}Z_{\mathrm{ball},v},
 \qquad
 I_{\mathrm{def}}=C_{\mathrm{def}}L\prod_{v\in S}Z_{\mathrm{def},v},
\]
with the same nonzero $L$-factor and identical local factors outside
$\{\mu,v_0\}$. If
\[
 I_{\mathrm{def}}\geq b,\qquad
 \frac{C_{\mathrm{ball}}}{C_{\mathrm{def}}}\geq c>0,
 \qquad \frac{Z_{\mathrm{ball},v}}{Z_{\mathrm{def},v}}\geq a_v>0
 \quad(v=\mu,v_0),
\]
then
\[
 I_{\mathrm{ball}}\geq bc\,a_\mu a_{v_0}.
\]
The proof is division of the two identities followed by cancellation and multiplication of the displayed lower bounds. This proves an implication; it does not supply the two identities for the particular 032 and Alp\"oge inputs.

An applicable Rallis statement must be checked in its actual direction and range. Gan--Qiu--Takeda, arXiv:1207.4709v3, \S11.3, p.52, displays the doubling seesaw and its determinant-character twist. Section 11.4 assumes $d(n)<m\leq2d(n)$ before Proposition 11.1, including its convergent clause, and Theorem 11.4 repeats that range. For unitary groups $d(n)=n$, so it does not directly cover the naive $\U(1)\to\U(3)$ substitution. Dividing a local zeta integral by its standard local $L$-factor can still leave the normalizing denominator $d_v(s,\chi_V)$; one must not silently declare all such spherical factors to be one.

An applicable positive locator for the $\U(1)\to\U(3)$ direction is J. Haan, arXiv:1403.5061v7, \S3.3, Theorem 3.4, p.8. With Hermitian source dimension one, skew-Hermitian target dimension three, splitting $(\gamma,\gamma^3)$, $\gamma|_{\A_F^*}=\chi_{E/F}$, and the paper's Tamagawa and symplectic conventions, it states
\[
 \langle\theta(\overline f_1,\phi_1),\theta(\overline f_2,\phi_2)\rangle
 =\frac{L_E(3/2,\operatorname{BC}(\tau)\otimes\gamma^3)}
        {L_F(3,\chi_{E/F})}
   \prod_v Z_v^{\#}(1,f_{1,v},f_{2,v},\Phi_{1,v}).
\]
Here $Z_v^{\#}=L_v(3,\chi_{E_v/F_v})Z_v/L_{E_v}(3/2,\operatorname{BC}(\tau_v)\gamma_v^3)$ and the doubled section comes from $\rho_{1,3}(\phi_1\otimes\overline\phi_2)$. The source proves it by the convergent formula of J.-S. Li; it is not an application of the restricted GQT theorem just audited. Transport of the splitting, the conjugated source input, and the measures is still required. The formula controls an individual torus theta lift; it does not automatically give the norm of an arbitrary wedge of two such lifts. Haan's main Bessel theorem has additional special-pair hypotheses.

The recovered archive also retains a quotient-volume factor explicitly. In
\path{Development/AutomorphicArithmetic/ThetaCorrespondence/RallisInnerProduct/GraphRelationExtendsFromFockVectorsToKFinite.lean}, lines 791--802, the proof of \nolinkurl{rallisInnerProductConvergent_of} passes through the identity
\[
 \frac{\langle\Theta(\Phi_1,f_1),\Theta(\Phi_2,f_2)\rangle_B}
      {\operatorname{vol}([B])}
 =\operatorname{dblInt}_{A,\chi_B^{-1}}
       (\operatorname{ThetaMean}_B(\delta(\Phi_1,\Phi_2));f_1,f_2).
\]
Lines 776--789 explicitly convert the integral over $[B]$ to its volume times the normalized theta mean. The proof uses the conjugated source functions in its Fubini step; its lift convention is not interchangeable with \texttt{liftS}. Indeed, \nolinkurl{GlobalTheory.RallisIP.thetaLift}, defined in \path{PublishedTheorems/TwistedTraceFormula/LowerLiftsZeroOfWittIdxEqZero.lean}, lines 22--26, is $h\mapsto\int_{[A]}\theta(g,h;\Phi)\overline{f(g)}\,dg$; lines 41--59 state the resulting range and volume-normalized zeta identity. Its input \nolinkurl{Published.SiegelWeilUnitaryConvergent}, defined in \path{PublishedTheorems/SiegelWeil$_2$/SiegelWeilUnitaryConvergent.lean}, lines 18--35, requires
\[
 \operatorname{rank}A<\operatorname{rank}B\leq2\operatorname{rank}A,
 \qquad \operatorname{rank}A+\operatorname{wittIdx}B<\operatorname{rank}B.
\]
Thus this archive formula does include the anisotropic adjacent-rank alternatives $1\to2$ and $2\to3$ of B15.

A \emph{different} law, \nolinkurl{SphericalVolumeLift.Convergent}, in \path{BlockTheory/TwistedDiscreteSeries/SphericalVolumeLift.lean}, lines 75--91, imposes $\operatorname{rank}B=\operatorname{rank}A+t$ and $\operatorname{rank}A<t$. Its range is $\operatorname{rank}B>2\operatorname{rank}A$; it is this larger-rank law that the recovered definite-positivity argument uses. Neither of these two ranges includes equal rank. There is also a separate \nolinkurl{Published.SiegelWeilAnisotropicUnitary} statement in \path{PublishedTheorems/TwistedTraceFormula/SiegelWeilAnisotropicUnitary.lean}, lines 109--119: with anisotropic $B$ and positive ranks, its Eisenstein value is $(1\text{ if }\operatorname{rank}A<\operatorname{rank}B,\,2\text{ otherwise})$ times the normalized theta mean. This supplies an analytic identity at equal rank, not automatic nonvanishing of a given theta lift. The volume, character, and rank hypotheses must be kept separately for each formula.

\subsection{THETA3-B9: the canonical local Bessel integral equals one}\label{B9}

The recovered \path{Formalism/PeriodIntegrals/BesselPeriodGGP.lean} defines a real local nonvanishing proposition. Its ambient matrix \texttt{Hpk} has signature $(k+r,r+1)$; the compact subgroup \texttt{UWset} fixes the hyperbolic coordinates and a negative vector. Its integral has the form
\[
 \int_{\U(W_z)\times N}
 \langle\sigma(x_0)w,\sigma(hn)\sigma(x)w\rangle
       \overline{\psi_N(n)}\,\pi(h)^{-1}\,dh\,dn.
 \tag{B9a}
\]
It is operator-valued in the finite-dimensional $\pi$-space. The proposition asks for a discrete-series $\sigma$, a lowest-$K$-type unit vector, compact $\pi$, interlacing and half-integral differences, and asserts existence of $x,x_0$ making the integral nonzero. It does not say that arbitrary periods are zero or one. \texttt{GapBig} is in a separate numerical setting, not in the actual proposition's assumptions.

\textbf{Statement.} Put $k=2$, $r=0$, and $z=e_3$, and use probability Haar measures. Suppose a unitary representation of $\U(2,1)$ contains a unit vector $w$ with
\[
 \sigma(\operatorname{diag}(A,d))w=\det(A)d^{-2}w,
 \qquad A\in\U(2),\ d\in\U(1).
\]
Take $\pi(\operatorname{diag}(A,1))=\det A$ on $\C$. Then (B9a), at $x=x_0=1$, equals one.

\textbf{Proof.} The matrix equations defining \texttt{Nset} at $r=0$ force $N=\{1\}$. A unitary matrix fixing $e_3$ preserves its positive orthogonal complement, so \texttt{UWset} is precisely $\operatorname{diag}(\U(2),1)$. In the convention linear in the second inner-product variable,
\[
 \langle w,\sigma(\operatorname{diag}(A,1))w\rangle=\det A.
\]
Multiplication by $\pi(h)^{-1}$ cancels this factor. The integrand is identically one; compactness gives absolute convergence, and probability Haar gives the stated value. \hfill$\square$

This is the correct type for the right-regular coefficient of a canonical two-form. At a negative line $o$, the tangent space is $\Hom(L,Q)$. For $k=\operatorname{diag}(A,d)$, its derivative is $A/d$. If
\[
 f_\xi(g)=(g^*\xi)_o(e_1,e_2),
\]
then
\[
 f_\xi(gk)=\det(A)d^{-2}f_\xi(g).
\]
Left lattice invariance follows from that of $\xi$. On a compact quotient this coefficient has finite nonzero $L^2$ norm when $\xi\ne0$, and can be normalized. The displayed character is trivial on the center of $\U(2,1)$, so it is compatible with the projective special-unitary action. It is a coefficient convention; the action on the canonical-bundle fibre is dual.

The exact numerical choice in the archive is
\[
 \lambda=(1,0,-1),\quad v=(1),\quad b=(0,1),\quad
 u=(3/2,1/2),\quad\epsilon=(+,+,-).
\]
The definitions in \texttt{UpqRep.lean} give
\[
 \rho_c=(1/2,-1/2,0),\quad\rho_n=(1/2,1/2,-1),\quad
 \lambda+\rho_n-\rho_c=(1,1;-2).
\]
For compact $\U(2)$, $u-\rho_c=(1,1)$, so $\pi=\det$. The interlacing is
$3/2>1>1/2>0$; the Harish--Chandra integrality and half-integral differences follow directly from these tuples. The direct evaluation does not require the general discrete-series theorem. Invoking that theorem instead would additionally require the relevant \texttt{IsDSHpk} representation identification. No global period has been set equal to one.

The same direct calculation works for a pluricanonical coefficient of $K_X^q$, $q\geq1$. Its compact character is $(\det A)^qd^{-2q}$, and taking $\pi=\det^q$ again makes the local period one. The corresponding parameters are
\[
 \lambda=(q,q-1,1-2q),\quad v=q,\quad
 b=(1-q,2q-1),\quad u=(q+1/2,q-1/2).
\]
The same half-sum calculation gives lowest weight $(q,q;-2q)$, and the interlacing follows from
$q+1/2>q>q-1/2>q-1$. The node gap is $3q-2$. Thus powers can enlarge the numerical gap, whereas a determinant twist shifts both nodes equally and cannot do so. A section $\xi^q$ is pluricanonical; it is not the original ordinary cohomology class in $H^2(X)$. Passing a quantitative period bound back to that class needs an additional argument.

\subsection{THETA3-B10: a separate high-rank matching theorem}\label{B10}

The recovered \path{Development/AbelianVarietyTheory/BlockTheory/DivisorBase/WeilTensorModels.lean} defines one global matching theorem. It is distinct from the iterated Bessel family in B13. Its two spaces have dimensions $a,a'$, and its common target has dimension $a+a'$. The exact equality \texttt{IsConst}, line 491, is
\[
 I_V(\Phi)=c\,I_{V'}(\Phi')
\]
for matched data. \texttt{OneConst}, line 503, requires uniqueness and $c\ne0$. \texttt{StepAlpha} and \texttt{StepBeta}, lines 571 and 559, contain recursive constants involving partial Hecke $L$-values and local epsilon/Gauss factors. \texttt{LocalMatching} and \texttt{ThetaIntegralsAgree} occur at lines 581 and 585.

Its \texttt{Geo.IsGood}, lines 360--370, requires
\[
 a'<a,\quad a+a'\equiv0\pmod2,\quad
 2(p+q)+1<\left\lfloor\frac{a-a'}2\right\rfloor,\quad
 2(p+q)+1<a',
\]
and signatures
\[
 V:(a-p,p),\quad V':(q,a'-q),\quad
 W:(a+a'-p-q,p+q).
\]
The definitions \texttt{Spc.IsHerm} and \texttt{Spc.IsSkew}, lines 182--200, impose these signatures at every real place. The base also has CM/Galois and integral unramified-extension restrictions.

For $p=1,q=0$, the inequalities force $a'\geq4$ and the even difference $a-a'\geq8$; the smallest pair is $(12,4)$, with common target dimension sixteen. Thus assigning $a=3$ does not instantiate this theorem. Increasing the dimension of a one-active-place ball does not by itself repair its all-place signature or source-dimension requirements. This calculation audits this one matching theorem; it does not exclude the different definite-subgroup route or the abstract family theorem below.

\subsection{THETA3-B11: the exact restriction theta seesaw}\label{B11}

\textbf{Statement.} Let $V=V_H\perp L$ be Hermitian, with $L$ a line. Put $G=\U(V)$ and $H=\U(V_H)$, embedded by fixing $L$ pointwise. Let $W$ be a skew-Hermitian source and $T=\U(W)$. Assume $[H]$ and $[T]$ are compact, or assume absolute convergence of all sums and integrals below. Fix compatible oscillator splittings for the orthogonal tensor rule. In particular the splitting character for $V$ is the product of those for $V_H$ and $L$, and the source splittings agree.

For a finite sum of Schwartz tensors
$\Phi=\sum_a\Phi_{H,a}\otimes\Phi_{L,a}$ and a unitary automorphic character $\eta$ of $T$, set
\[
 F_\eta(g)=\int_{[T]}\theta_V(g,t;\Phi)\eta(t)^{-1}\,dt,
 \qquad
 B_e(F)=\int_{[H]}F(h)\overline{e(h)}\,dh.
\]
For a continuous automorphic test function $e$ on $[H]$, define
\[
 \Theta_{H\to T}(e;\Phi_H)(t)=
 \int_{[H]}\theta_{V_H}(h,t;\Phi_H)\overline{e(h)}\,dh.
\]
Then
\[
 B_e(F_\eta)=\sum_a\int_{[T]}
 \Theta_{H\to T}(e;\Phi_{H,a})(t)
 \theta_L(1,t;\Phi_{L,a})\eta(t)^{-1}\,dt.
 \tag{B11}
\]
All measures are the displayed measures, so no unspecified scalar occurs.

\textbf{Proof.} A polarization compatible with the orthogonal sum identifies the rational indexing lattice with the direct sum of the two indexing lattices. The tensor identity for the oscillator action identifies each summand for a pure tensor with a product. Absolute convergence permits separation of the sums, giving
\[
 \theta_V((h,1),t;\Phi_H\otimes\Phi_L)
       =\theta_{V_H}(h,t;\Phi_H)\theta_L(1,t;\Phi_L).
\]
Sum over $a$, multiply by $\overline{e(h)}\eta(t)^{-1}$, and integrate. Fubini applies by the convergence hypothesis. The line factor and $\eta$ are independent of $h$, and the inner $h$-integral is the displayed reverse theta lift. This proves (B11). In the compact case, uniform convergence on compact representative sets and boundedness of $e$ supply absolute integrability. \hfill$\square$

For a theta differential form, the statement applies to its polynomial Fock coefficient once that coefficient and its orthogonal tensor decomposition have been identified. This is a required identification of vectors, not merely of form degrees.

The actual archive declaration is
\nolinkurl{GlobalTheory.T.thetaTensorRule_orthogonalSum_source}, in
\path{Proofs/DecoratedLatticeProofs/ThetaTensorRuleOrthogonalSumSourceStatement.lean}, with the proof in the matching file without \texttt{Statement}. The property \texttt{TensorSource}, in \path{BlockTheory/TwistedDiscreteSeries/ConjugateAdjointDagger.lean}, lines 27--53, explicitly imposes the splitting product and tensor action. The adjacent target tensor rule is used by \nolinkurl{LastSeeSawThetaIdentity.lean}.

The precise global Bessel orientation is also recovered. In
\path{Definitions/TwistedInnerProducts/WittSmallGroupEmbeddingMemPointwiseStabilizerIso.lean}, lines 176--181, its integrand is
\[
 \phi(h)f(n\,\operatorname{emb}_A(h))\overline{\psi_{\rm Bes}(n)}.
\]
Its \texttt{liftS\_apply}, line 163, integrates $\theta(g,h)G(h)$ without conjugating $G$. Thus (B11) matches by taking $\phi=\overline e$ and $G=\eta^{-1}$. At Bessel rank zero the unipotent is trivial.

The source character $\eta$ remains visible in (B11). The identity therefore supplies a genuine global definite-theta/Bessel comparison, but does not on its own prove the cancellation claimed for the special iterated family. Nor does it equate the full ball norm with its restriction norm.

\subsection{THETA3-B12: every nonzero canonical form has a definite Bessel coefficient}\label{B12}

\textbf{Statement.} Let $X$ be the compact congruence ball surface of B0 and $0\ne\xi\in H^0(X,K_X)$. There is an $E$-rational negative line $L$ such that $H=\U(L^\perp)$ is totally definite of rank two and a finite-level automorphic test function $e$ with $B_e(f_\xi)\ne0$. This evaluates a canonical-bundle coefficient, not the cohomological pullback of $H^2(X)$ to a point.

More precisely, use probability Haar on $[H]$, a compact open finite level $K_H$ fixing $f_\xi$, and the finite-dimensional space $R$ with this level and the same one-dimensional infinity type. Put $N_H=\dim R$, let $m_0>0$ be the mass of the identity double coset, and let $b=|f_\xi(1)|$. For any orthonormal basis $(e_j)$ of $R$,
\[
 \sum_{j=1}^{N_H}|B_{e_j}(f_\xi)|^2
      =\|f_\xi|_{[H]}\|_2^2\geq m_0b^2,
 \qquad
 \max_j|B_{e_j}(f_\xi)|^2\geq\frac{m_0b^2}{N_H}.
 \tag{B12}
\]
In particular, $\xi\ne0$, $\int_X\xi\wedge\overline\xi>0$, and existence of some such nonzero Bessel coefficient are equivalent.

\textbf{Proof.} First make the adelic realization precise. Fix the $\mathcal O_E$-lattice defining the principal congruence quotient. An element of $\U(V)(F)$ preserving this lattice has determinant in $\mathcal O_E^\times$ with $\alpha\overline\alpha=1$. Every complex conjugate of $\alpha$ has absolute value one, so $\alpha$ is a root of unity. There are finitely many such roots in $E$. Choose a principal congruence level at which no nonidentity root is congruent to one. Then
\[
 \Gamma'=\U(V)(F)\cap K_f\subset\SU(V)(F).
\]
Make the level deeper so that $\Gamma'\subset\Gamma$, where $X=\Gamma\backslash D$; this inclusion has finite index. Thus $X'=\Gamma'\backslash D$ is both a finite principal cover of $X$ and a connected component of the adelic $\U(V)$ quotient. Pull $\xi$ to $X'$ and extend it by zero on the other components. This gives the automorphic canonical-bundle coefficient $f_\xi$ used below. Its central archimedean character is trivial, since $\det(tI_2)t^{-2}=1$. If the covering degree is $d$, then
\[
 \int_{X'}p^*\xi\wedge\overline{p^*\xi}
       =d\int_X\xi\wedge\overline\xi,
\]
so nonvanishing and positivity are unchanged. The masses and finite level in (B12) refer to this chosen realization. We retain the notation $\xi$ for its pullback.

The nonzero locus of $\xi$ is a nonempty open subset of the ball. Rational vectors are dense at the active embedding: the image of $E$ in $\C$ contains $\Q+\Q\alpha$ for a nonreal $\alpha$, hence is dense. In coordinates, $V(E)$ is therefore dense in $V_\mu$. Choose $z\in V(E)$ whose negative line lies in this open locus. At every other real place the form on $V$ is positive definite, so $h_\tau(z,z)>0$. Thus $L^\perp$ has signature $(2,0)$ everywhere and $H$ is totally definite.

Base the canonical coefficient at $L$, choosing a nonzero tangent two-frame there. The calculation of B9 gives the right compact character $\det(A)d^{-2}$, and on $H_\infty$, which fixes $L$ pointwise, it is $\det A$ at the active place. The other compact factors have the prescribed canonical coefficient type, trivial for the trivial coefficient system. Its absolute value is therefore constant along each right $H_\infty$ orbit. By construction $f_\xi(1)\ne0$. The adelic realization just constructed gives an open compact $K_H$ fixing the coefficient.

The set $H(F)\backslash H(\A_F)/(H_\infty K_H)$ is finite. On each double coset a function of the specified scalar infinity type is determined by one complex number, subject to its finite stabilizer condition. Thus $R$ is finite-dimensional. Let $m_i>0$ be the masses of these cosets. The identity coset contributes exactly $m_0b^2$ to the norm square, and all other terms are nonnegative. Consistent Haar choices give masses proportional to reciprocal finite stabilizer orders, normalized by $\sum_i m_i=1$. At neat level the stabilizers are trivial and the masses are equal.

Parseval in $R$ proves the equality in (B12); one term is at least their average. The orthonormal basis may be refined within the automorphic spectral summands of the compact quotient, so these are periods against definite automorphic test vectors. This assertion does not require their values or their Hilbert normalizations to be algebraic.

Conversely, a nonzero Bessel coefficient forces $f_\xi$, and hence $\xi$, to be nonzero. In local holomorphic coordinates, if $\xi=a\,dz_1\wedge dz_2$, then
\[
 \xi\wedge\overline\xi
   =4|a|^2\,dx_1\,dy_1\,dx_2\,dy_2.
\]
Its integral in the complex orientation is positive exactly when $\xi\ne0$. This proves the equivalence. \hfill$\square$

There is also a certified point-selection version. If a unit-frame coefficient has size $b_0>0$ at a point, has Lipschitz constant $L_\xi$ nearby, and the Hermitian form has a known negative margin, choose a rational point within the margin-preserving radius and within $b_0/(2L_\xi)$ when $L_\xi>0$. Then $b\geq b_0/2$ and the lower bound is $m_0b_0^2/(4N_H)$. THETA3-Q3 gives an explicit denominator from the coordinate matrix, margin, and radius. The resulting level, mass, and dimension still require bounds or computation.

For two restrictions in the same $R$, Parseval gives the mixed identity
\[
 \int_{[H]}f\overline{f'}=
 \sum_j B_{e_j}(f)\overline{B_{e_j}(f')}.
\]
THETA3-A5 passes from finitely many restrictions to the full ball inner product by a positive inverse frame operator. A naive quotient disintegration does not do this: it unfolds to $H(F)\backslash G(\A)$, rather than $G(F)\backslash G(\A)$. The rational or algebraic nature of the frame correction is a further issue.

Finally, a central lift of the Shimura cocharacter at $L$ may factor through $\U(L^\perp)\times\U(L)$ rather than $H$ alone. The proof uses the finite definite adelic module and canonical-fibre evaluation. It does not assert an unspecified morphism of Shimura data $H\to G$.

\subsection{THETA3-B13: the recovered character-independent functional}\label{B13}

\paragraph{The exact family.}
The source \path{Definitions/LatticeFrames/ValLastX.lean}, under
\path{Development/AbelianVarietyTheory}, contains the relevant definitions. The following locators separate data from conclusions:
\begin{itemize}
\item \texttt{Comm}, lines 28--57, fixes the indefinite spaces $V_j$, discrete-series models, definite spaces $X_j$, Bessel/Jacobi decompositions, measures, final $T$, and $X_m=T\perp L_z$.
\item \texttt{Tup}, lines 79--113, supplies varying Weil representations, finite translation combinations, row functions, and constants. \texttt{RealEq}, lines 182--184, requires their real Weil actions to agree.
\item \texttt{perF}, lines 119--120, is a finite sum of genuine Bessel periods of finite right translates. \texttt{per}, lines 122--123, inserts the row function and automorphic realization. \texttt{perT}, lines 137--138, is the final definite restriction period.
\item \texttt{up}, lines 166--167, iterates two definite theta operators with the intervening $S$ operation; \texttt{dn}, lines 169--170, composes the reverse automorphic lift maps.
\end{itemize}
The normalizing products, lines 172--174, are
\[
 c_{\rm Top}(n)=
 \left(\prod_{i=0}^{n-1}c_V(2i)c_K(2i)c_K(2i+1)\right)c_P,
 \qquad
 c_{\rm Bot}=\prod_{j=0}^{m-1}c_I(j).
\]
\texttt{LawsX}, lines 212--221, includes common-data validity, tuple validity, finite support, common real action, the double-step and last-step integral laws, definite kernel reduction, final period factorization, and row covariance.

The actual source witness is
\nolinkurl{GlobalTheory.SeeSaw.mainX_of_lawsX}, in
\path{Proofs/DecoratedLatticeProofs/SeesawTransportRelationsOfFamilyLaws.lean}, lines 25--37. Its inputs are \texttt{CuspRapidDecay}, \texttt{HarishChandraModules}, and \texttt{LawsX}; its conclusion is \texttt{MainX}. The matching \texttt{Statement.lean}, line 12, exposes the wrapper
\nolinkurl{GlobalTheory.T.seesawTransportRelations_of_familyLaws}.
These declarations and the displayed proof body were read, but their complete transitive dependencies were not audited or compiled.

\paragraph{Mathematical content.}
Fix $m=2n+2$. For an allowed tuple, input column, and terminal automorphic vector, indexed collectively by $s$, put
\begin{align*}
 A_s&=c_{{\rm Top},s}(n+1)
      \operatorname{tail}_s(\operatorname{up}_s(n+1)\psi_s,F_s),\\
 B_s(y,v)&=c_{{\rm Bot},s}
      \operatorname{per}_s\bigl(\psi_s,
       \operatorname{dn}_s(2n+1)\operatorname{ioT}_s(F_s);y,v\bigr).
\end{align*}
\texttt{MainX}, lines 244--251, says that every finite scalar relation
$\sum_s\alpha_s A_s=0$ implies
$\sum_s\alpha_s B_s(y,v)=0$ for every fixed $y,v$.
Consequently there is a single scalar $a(y,v)$ such that
\[
 B_s(y,v)=a(y,v)A_s
 \quad\text{for every allowed tuple, column, and terminal vector.}
 \tag{B13}
\]
The scalar is unique if some $A_s\ne0$. This is the precise normalized independence supplied by the recovered theorem. It is not independence under arbitrary changes of characters and Schwartz data.

\textbf{Proof of the consequence.} If all $A_s$ vanish, one-term relations force all $B_s(y,v)=0$, and $a=0$ works. Otherwise choose $s_0$ with $A_{s_0}\ne0$ and put $a=B_{s_0}(y,v)/A_{s_0}$. For each $s$, the two-term relation
$A_{s_0}A_s-A_sA_{s_0}=0$ gives
$A_{s_0}B_s-A_sB_{s_0}=0$. Division proves (B13), and also uniqueness. This is a complete elementary re-proof of the logical passage from relation transfer to the common scalar. It does not re-prove the analytic inputs of \texttt{MainX}. \hfill$\square$

\paragraph{The exact instantiation and arithmetic gates.}
The abstract family is more flexible than the particular tower used in the construction. In \path{Definitions/LatticeFrames/JointSmooth.lean}, lines 56--62, \texttt{DefSp.Is} requires a compact quotient and a Cartan involution equal to $+1$ or $-1$ at each real place, hence definiteness, without requiring the same sign everywhere. Its \texttt{BesAt}, lines 107--115, consists of a line orthogonal decomposition and a hyperbolic decomposition. Its rank-zero local geometry therefore includes the active $\U(2,1)/\U(2)$ pair of B9.

By contrast, the concrete \texttt{PiLocalization$_2$.Is} in
\path{Definitions/LatticeFrames/ChainParSucc.lean}, lines 56--60, requires the same signature \texttt{sigYD D.R j} at every embedding. In \texttt{ExistsDatOfTower.lean}, lines 151--171, the Bessel step is built with a globally negative line, of signature $(0,1)$ everywhere. Hence that constructed tower does not immediately equal the one-active-place 032 datum. The theorem's abstract shape does not itself supply a witness satisfying all these hypotheses, nor the discrete-series span and gap requirements used to construct the tower.

The gap obstruction can be stated numerically. In \path{BlockTheory/TwistedDiscreteSeries/ThetaSourceAction.lean}, lines 209--211, \texttt{NodeGap} requires a natural number $N\geq\operatorname{rank}X'+1$ such that
\[
 t+2N^3+1\leq b_{i+1}-b_i
\]
for every adjacent pair and every infinite place. Consequently its right-hand gap is at least three. The canonical choice $b=(0,1)$ in B9 cannot satisfy this condition when identified with the two nodes. This requirement is used to construct the tower's spanning data, for example in \texttt{StepVanOfStepX.lean}, lines 71--84; it is not an additional assumption of the direct local integral evaluated in B9, nor of the elementary consequence (B13) once \texttt{LawsX} has already been supplied.

The bridge to the arithmetic columns is explicitly named. In
\path{Definitions/LatticeFrames/RatioOKOfRootOfUnity.lean}, the property
\texttt{FamAll}, lines 60--62, combines \texttt{LawsX} with
\texttt{LinkZero}, \texttt{LinkTail}, \texttt{LinkDown},
\texttt{LinkOnto}, \texttt{LinkComm}, and \texttt{RatioOK}.
The last condition compares $c_{\rm Top}/c_{\rm Bot}$ across tuples.
\texttt{UnitRatio}, in \path{Definitions/LatticeFrames/UnitRatio.lean}, lines 22--23, means
\[
 y\ne0,\quad \exists n>0,\ r\in\Q^*:
 (x/y)^n=r,
 \qquad v_p(r)\ne0\Longrightarrow p\in S.
\]
Thus the source controls radical-of-rational ratios with specified prime support; it does not automatically identify their absolute values with one. To derive a quantitative period bound, one must instantiate these links for the cohomological vectors, retain the normalizers, and bound the relevant common scalar or normalized Gram matrix. B9 and B12 give explicit local and definite-subgroup nonvanishing inputs, while THETA3-A5 supplies a separate exact finite-frame comparison of the global mixed pairing.

\subsection{THETA3-B14: large weights, tensor contractions, and the return to cohomology}\label{B14}

The large-weight device has an exact scope. For $q\geq1$, the canonical-power coefficient has compact type $(\det A)^qd^{-2q}$ and the parameters displayed after B9. Its local Bessel integral is one and its node gap is $3q-2$. Thus, for a fixed numerical threshold $M$, the condition $3q-2\geq M$ is attained by $q\geq\lceil(M+2)/3\rceil$. This does not alone construct the remaining tower data, and changes the input from a two-form to a pluricanonical section.

\paragraph{A cohomological distinction.}
For $q>1$, the coherent multiplication map
\[
 \operatorname{Sym}^q H^0(X,K_X)\longrightarrow H^0(X,K_X^q)
\]
is not the exterior-product map followed by diagonal pullback. If
$0\ne\xi\in H^{2,0}(X)$, its external product lies in
$H^{2q,0}(X^q)$, whereas diagonal pullback lies in
$H^{2q,0}(X)=0$, since $2q>\dim_\C X=2$.
The coherent section $\xi^q$ is nonzero: in a local canonical frame it is the $q$th power of its coefficient, tensor the $q$th power of the frame. These are different operations.

There is no nonzero global holomorphic bundle map $K_X^q\to K_X$ on a compact ball surface. Indeed, $K_X$ is ample, and such a map would be a section of $K_X^{1-q}$. A nonzero section would give an effective divisor $D$ with
\[
 D\cdot K_X=(1-q)K_X^2<0,
\]
contrary to ampleness; the nowhere-zero case would make $K_X^{1-q}$ trivial, also contradicting $K_X^2>0$. This rules out a global fibrewise holomorphic contraction of that specified kind. It does not rule out other correspondences between auxiliary varieties or maps requiring extra data.

Division by $\xi^{q-1}$ is an inverse on the particular image
$\xi^{q-1}H^0(X,K_X)\subset H^0(X,K_X^q)$, because cancellation gives a holomorphic result on that image. On an arbitrary section it can have poles along the zero divisor of $\xi$. Thus using this inverse requires proof that the high-weight input belongs to this image. A nonzero vector produced somewhere in the full high-weight space does not supply that proof.

\paragraph{The precise tensor contraction and its nonvanishing.}
Let $V=H^2(X,\Q)$ and
\[
 Q_X(u,v)=\int_Xu\smile v:\ V\otimes V\longrightarrow\Q(-2).
\]
Put $A=Q_X(\xi,\overline\xi)>0$. Every contraction of two factors in the pure tensor $\xi^{\otimes q}$ by $Q_X$ vanishes, because $Q_X(\xi,\xi)=0$ by Hodge type. By contrast the contraction
\[
 C_q:V^{\otimes q}\otimes V^{\otimes(q-1)}
       \longrightarrow V(-2(q-1)),
\]
\[
 C_q(u_1\otimes\cdots\otimes u_q\otimes v_2\otimes\cdots\otimes v_q)
       =u_1\prod_{i=2}^q Q_X(u_i,v_i)
\]
satisfies
\[
 C_q(\xi^{\otimes q}\otimes\overline\xi^{\otimes(q-1)})
           =A^{q-1}\xi\ne0.
 \tag{B14a}
\]
There is no new nonvanishing hypothesis in this calculation beyond $\xi\ne0$.

The map $C_q$ is induced by algebraic diagonals and push-forward. On $X^{2q-1}$, pair the $i$th and $(q+i-1)$st factors, for $2\leq i\leq q$, by the pullback of the diagonal class of $X\times X$, and push to the first factor. For an external product, the projection formula gives one factor $\int_Xu_i\smile v_i$ from each pair. All factors have even degree, so no permutation sign occurs. The resulting Tate twist is exactly $-2(q-1)$. This proves both the formula and its algebraic-correspondence realization.

This construction needs an actual external-product cohomological input, including the conjugate factors. A coherent section of $K_X^q$ alone is not such an input. Nor does the fact that $C_q$ is algebraic turn a non-Hodge vector such as $\xi^{\otimes q}\otimes\overline\xi^{\otimes(q-1)}$ into an algebraic cycle. In a proposed cycle construction, the ambient Hodge types and the required projected component must first be identified; after that, algebraic contraction preserves algebraicity and (B14a) supplies the appropriate nonzero test value.

A parallel representation-theoretic check applies to a specified rational Weil component. Suppose a polarizing algebraic torus acts on
$W_\C=\ell_\chi\oplus\ell_{\chi^{-1}}$, where $\chi$ has infinite order, and its invariant pairing pairs the two lines. Any torus-equivariant map from $W^{\otimes q}$ to $W$ sends the pure weight $\chi^q$ to zero for $q>1$, since neither $\chi$ nor $\chi^{-1}$ equals $\chi^q$. Pairwise polarization contractions are equivariant and therefore have this obstruction. With $q-1$ opposite-character factors the total weight is $\chi$, and the contraction is nonzero whenever their pairings are nonzero. This is a statement about that explicitly specified torus and equivariant contractions; it is not an impossibility theorem for every algebraic correspondence. It pinpoints which conjugate component a tensor-power argument must retain.

\paragraph{Analytic return with an explicit norm loss.}
Fix a Hermitian metric on $K_X$ and a positive measure on $X$. Write
\[
 A_q=\int_X|\xi|^{2q},\qquad M=\|\xi\|_\infty>0.
\]
Pointwise $|\xi|^{2q}\leq M^{2q-2}|\xi|^2$, hence
\[
 A_1\geq \frac{A_q}{M^{2q-2}}.
 \tag{B14b}
\]
Therefore a certified lower bound for the global norm of the particular pure power $\xi^q$, together with a certified upper bound for $M$, gives a lower bound for the original norm. A local Bessel integral normalized to one is not by itself that global bound.

There is also a nonzero analytic contraction. Use the metric to contract
$\xi^q\otimes\overline{\xi}^{\otimes(q-1)}$ to the smooth canonical section
$|\xi|^{2q-2}\xi$, and take its orthogonal projection to $H^0(X,K_X)$. Its inner product with $\xi$ is $A_q>0$, so the projection is nonzero. This projection depends on the metric and the original form and is not thereby an algebraic correspondence. Equations (B14a) and (B14b) give two valid return mechanisms with their required input data; neither identifies an arbitrary high-weight output with the original Weil component automatically.

\subsection{THETA3-B15: a small-rank local alternative and its exact global gates}\label{B15}

There is an alternative to the large-weight construction: the canonical discrete series itself has a nonzero local theta correspondence with a compact rank-two group. We give the characters and parameters explicitly. This removes a local occurrence obstruction; it does not identify a particular global theta vector with the iterated lift of B13.

\paragraph{The literal one-active-place datum forces the rank-zero step.}
Suppose that $V_0$ in \texttt{Comm} is the rank-three space $V$ of B0 and that $m=2n+2$. Then $n=0$, hence $m=2$. Indeed, \texttt{Comm.bes}, at \texttt{ValLastX.lean}, line 46, requires a \texttt{BesAt} datum of rank $n$. By \texttt{JointSmooth.lean}, lines 107--115, this contains a global hyperbolic decomposition of $V$ with $n$ hyperbolic planes. A nonzero isotropic vector over $E$ remains a nonzero isotropic vector at every completion. Since $V$ is positive definite at another real place, it has no such vector. Its global Witt index is zero, and therefore $n=0$.

At $m=2$ all the intermediate conditions indexed by $j+2(\sigma+1)+2=m$ are empty. The rank-zero \texttt{JacAt} identifies $V_1$ with the definite space $X_1$; write this middle space as $S$. Thus the remaining geometric pattern is
\[
 V=A\perp L,\qquad S=X_1\text{ totally definite},\qquad
 X_2=T\perp L_z\text{ totally definite}.
 \tag{B15a}
\]
Here $A$ is the definite rank-two complement in B12. The abstract \texttt{Comm} does not impose the concrete ranks \texttt{rankY} on its spaces $V_j$. Its field \texttt{par}, line 33, may vary with the real place. The numerical \texttt{Ranks} conditions do allow $m=2$: the choices $k=6,t=13,\operatorname{ord}=1$ satisfy the inequalities in \path{Formalism/LocalRingCardinality/Params.lean}, lines 6--22. This numerical observation does not construct the other fields of \texttt{D7} or a full \texttt{Comm.Is} witness.

The last seesaw imposes a precise line relation, not an arbitrary real signature change. In \texttt{PlaceLaws.lean}, lines 165--166, \texttt{WeilOK.tau} is
\[
 \operatorname{rank}L_z=1,\qquad
 \tau^*G_L\tau=-N_{E/F}(\xi_0)G_{L_z}.
 \tag{B15b}
\]
The norm is positive at every real place. Hence $L_z$ has the opposite sign to $L$: positive at $\mu$, negative elsewhere. This has a coherent global realization: set $L_z=-N_{E/F}(\xi_0)^{-1}L$ and $\tau=1$. Taking $T=L_z$ gives a rank-one top and a definite rank-two $X_2=T\perp L_z$, with positive sign at $\mu$ and negative sign elsewhere. These are actual global forms; no unverified collection of finite determinant classes is being prescribed. Choosing them does not yet prove the needed finite local theta occurrences.

\paragraph{Convention and local theorem.}
Use Ichino's convention in \emph{Theta lifting for tempered representations of real unitary groups}, arXiv:2008.06174v2, \S4.1, p.10. The input space $W$ is skew-Hermitian of dimension $n$ and signature $(p,q)$, the target space is Hermitian of dimension $r+s=m$, and
\begin{equation}
\begin{gathered}
 \chi_V(z)=(z/|z|)^{m_0},\qquad
 \chi_W(z)=(z/|z|)^{n_0},\\
 m_0\equiv m\pmod2,\qquad n_0\equiv n\pmod2,
 \qquad \psi(x)=e^{-2\pi i x}.
\end{gathered}
 \tag{B15c}
\end{equation}
The local lift is characterized by
$\Hom_{\U(W)\times\U(V)}(\omega,\pi\boxtimes\theta(\pi))\ne0$.
Let $\pi_\mu$ have Harish--Chandra parameter $(1,0;-1)$ on $\U(2,1)$, and let $\pi_\tau=1$ on $\U(3)$ for $\tau\ne\mu$. For the pair of dimensions $(n,m)=(3,2)$, take $m_0=-2,n_0=-1$. Then
\[
 \theta_{2,0}(\pi_\mu)=\det,\qquad
 \theta_{2,0}(\pi_\tau)=\det\quad(\tau\ne\mu).
 \tag{B15d}
\]
Here the character $\det$ of compact $\U(2)$ has Harish--Chandra parameter $(3/2,1/2)$.
The representation $\pi_\mu$ is the holomorphic discrete series with lowest compact type $(1,1;-2)$, namely $(A,d)\mapsto\det(A)d^{-2}$. This is the canonical right-regular coefficient convention fixed in B9: $\lambda+\rho_n-\rho_c=(1,1;-2)$. Replacing the right-action convention by its dual would require dualizing this parameter and the subsequent pairings as well.

For the next pair, let $\sigma=\det^{-1}$ on $\U(2)$, with Harish--Chandra parameter $(-1/2,-3/2)$. The following local lifts are nonzero and trivial:
\[
\begin{array}{c|c|c|c|c}
\text{place}&(n,m)&(m_0,n_0)&\text{target signature}&\theta(\sigma)\\ \hline
\mu&(2,1)&(-3,-2)&(1,0)&1\\
\tau\ne\mu&(2,1)&(-1,2)&(0,1)&1.
\end{array}
 \tag{B15e}
\]
Thus there are compatible middle representations for a two-step composition at every real place: the first oscillator has a quotient $\pi_\tau\boxtimes\det$, and the second has a quotient $\det^{-1}\boxtimes1$.

\textbf{Proof.}
In all three calculations the rank drops by one. After subtracting $m_0/2$ from the input Harish--Chandra parameter, the nonvanishing criterion has a single zero chain, of length $k=1$. For (B15d) the shifted parameters are
\[
 (2,1;0)\quad\text{and}\quad(2,1,0),
\]
respectively. In Atobe's notation both have $(k_\lambda,r_\lambda,s_\lambda)=(1,2,0)$. Atobe, \emph{On the non-vanishing of theta liftings of tempered representations of $\U(p,q)$}, Adv. Math. 363 (2020), Definition 1.6 and Theorem 1.7(2), pp.4--5, gives nonvanishing at the base point $(2,0)$. The theorem for discrete series is unconditional, as stated immediately after it. Ichino, \S6.2, explicitly identifies Atobe's additive-character convention with (B15c).

The parameter rule for adjacent ranks deletes the zero, transfers positive entries in the first block and negative entries in the second block to the positive target block, and then adds $n_0/2$. This is the adjacent-rank case of Ichino's Theorem 4.1(ii), pp.11--12, already covered by Paul's adjacent-rank correspondence cited there. In (B15d) it gives $(2,1)-(1/2,1/2)=(3/2,1/2)$. For the first row of (B15e), the shifted parameter is $(1,0)$, giving base point $(1,0)$ and output $1-1=0$. For the second row it is $(0,-1)$, giving base point $(0,1)$ and output $-1+1=0$. The same nonvanishing theorem applies in both cases. Every parity condition in (B15c) holds. Compact rank-one Harish--Chandra parameter zero is the trivial representation. This proves the claims. \hfill$\square$

The abstract groups in this calculation must be identified with the Hermitian/skew-Hermitian spaces of 032 and the archive, including the sign of the imaginary scalar used to pass between the forms. Equations (B15c)--(B15e) specify the convention rather than silently making that identification. Likewise, the archive's \texttt{liftS} integrates the kernel against the input function bilinearly. A quotient $\pi\boxtimes\det$ therefore pairs with an automorphic input of type $\det^{-1}$ to produce $\pi$. This is why the middle input in (B15e) is the dual of the B9 compact test representation. The two pairs are allowed different splitting characters; \texttt{WeilOK.last}, lines 160--164, separately names the second pair's character $\chi'_S$.

\paragraph{The archimedean characters can be made global.}
Let $E/F$ be CM, let $\epsilon$ be its quadratic idele-class character, and fix integers $a_\tau$ all congruent to $d$ modulo two. There is a unitary Hecke character $\chi$ of $E$ satisfying
\[
 \chi|_{\A_F^\times}=\epsilon^d,
 \qquad \chi_\tau(z)=(z/|z|)^{a_\tau}.
 \tag{B15f}
\]
Finite ramification is allowed; prescribed finite local values are not part of this assertion.

\textbf{Proof.}
The natural map $C_F\to C_E$ is a closed embedding. Indeed it is injective, and the idele norm of $C_E$ restricts to the square of the norm on $C_F$; the norm-one idele-class group of $F$ is compact, so the embedding is proper. The character-extension theorem for closed subgroups of locally compact abelian groups supplies a unitary extension $\chi_0$ of $\epsilon^d$. This is precisely the argument in 032, Lemma 4.4, using Folland, \emph{A Course in Abstract Harmonic Analysis}, second edition, \S4.3, Corollary 4.42. Since $\chi_0$ is trivial on each positive real scalar, its archimedean characters are $(z/|z|)^{a^0_\tau}$, with $a^0_\tau\equiv d\pmod2$.

Put $T_1=\operatorname{Res}^{1}_{E/F}\mathbf G_m$. Its real group $T_1(F\otimes\R)$ embeds as a compact subgroup of $T_1(F)\backslash T_1(\A_F)$ by placing one at the finite components: the intersection with the diagonal rational group is trivial. Character extension on this locally compact abelian quotient gives an automorphic character $\rho$ with archimedean weights $(a_\tau-a^0_\tau)/2$. Then
\[
 \chi(z)=\chi_0(z)\rho(z/\overline z)
\]
is trivial on $E^\times$, has the same restriction to $\A_F^\times$ as $\chi_0$, and has the prescribed weights. Its finite components have open kernels on compact unit groups, so it has finite conductor. \hfill$\square$

This proves that the integer choices in (B15d)--(B15e) themselves have no global parity obstruction. It does not force the finite components of those characters to be those in a previously fixed theta construction.

\paragraph{What changing a splitting transports.}
For a fixed dual pair, changing the splitting multiplies the oscillator action by determinant characters of the same two groups. Locally, Atobe's Proposition 3.3 and Proposition 3.5, pp.20--22, give
\begin{equation}
\begin{aligned}
 \omega_{m_0+2a}|_{\U(p,q)}
       &\simeq\omega_{m_0}|_{\U(p,q)}\otimes\det^a,\\
 \theta^{m_0+2a}(\pi)
       &\simeq\theta^{m_0}(\pi\otimes\det^{-a}),
\end{aligned}
 \tag{B15g}
\end{equation}
when the other splitting character is kept fixed. The latter formula includes zero lifts. A simultaneous change of the other character tensors the output by its corresponding determinant character.

The archive exposes its orientation in
\path{Formalism/LocalToGlobal/IndefiniteReal/LocalNewformFunction.lean}, line 51:
\[
 \texttt{splitFacO}(\chi_S,\chi_T;x,y)=\chi_T(x)^{-1}\chi_S(y).
 \tag{B15h}
\]
If $\kappa_S(z)=\rho_S(z/\bar z)$ and $\kappa_T(z)=\rho_T(z/\bar z)$ are character ratios, its split-place scalar changes by
$\rho_T(\det g)^{-1}\rho_S(\det h)$.
The determinants here follow from \texttt{TiedWeil.lean}, lines 71--78, and the actual \nolinkurl{IsWeilDatumO.om_split} condition at lines 72--75 of \nolinkurl{LocalNewformFunction.lean}. Thus an inverse in one convention cannot simply be removed when comparing it with (B15g).

All these determinant factors restrict trivially to the appropriate special unitary group. The canonical $\SU(V)$ geometry can consequently survive a determinant twist of its $\U(V)$ realization. However, the underlying source space is unchanged: a mixed-sign rank-two $P$ in A7 remains that same $P$, and a local group $\U(1,1)$ does not become $\U(2)$ by a character twist. Compensating the target twist as in (B15g) preserves the original local correspondence. To replace the source by a different definite inner form while keeping a fixed target vector, one must establish occurrence in the new dual pair. The original A7 realization does not alone provide that occurrence at every finite place, or an automorphic representation on the new source.

\paragraph{A precise global nonvanishing test for the alternative.}
There is a usable criterion with no unspecified real-place correction. In the conventions of Gan--Qiu--Takeda, arXiv:1207.4709v3, Theorem 11.7, pp.56--57, let $U_n$ be the source of dimension $n=1$ or $2$, and let $V$ be an anisotropic target of dimension $m=n+1$. For an irreducible cuspidal automorphic source representation $\Sigma$, its lift to this particular $V$ is nonzero if and only if
\begin{equation}
\begin{gathered}
 \theta_{V_v}(\Sigma_v)\ne0\quad\text{for every place }v,
 \\ L(1,\Sigma\times\chi_V)\ne0
       \quad\text{with a holomorphic value}.
\end{gathered}
 \tag{B15i}
\end{equation}
Here $L$ is the source's completed standard doubling $L$-function, with the character as named in that paper. The hypotheses are $d(n)=n<m=n+1\leq2n$, target Witt index $r=0$, and hence no lower step in its global Witt tower. The equality $m=d(n)+1$ is exactly the final case in Theorem 11.7(ii); Proposition 11.6 says that it removes the otherwise possible real-sign replacement. No temperedness hypothesis is needed for this criterion.

With Tamagawa measures and the paper's normalized doubling factors, Theorem 11.4(ii), p.54, specializes to
\begin{equation}
\begin{gathered}
 \langle\theta(\phi_1,f_1),\theta(\phi_2,f_2)\rangle
 \\
 =2L(1,\Sigma\times\chi_V)
       Z^*(1/2,\phi_1\otimes\overline{\phi_2},f_1,f_2),
 \\
 Z_v^*(s)=\frac{Z_v(s)}{L(s+1/2,\Sigma_v\times\chi_{V,v})}.
\end{gathered}
 \tag{B15j}
\end{equation}
The global product includes the unramified denominator factors of that normalization; they must not be replaced by one. The archive's bilinear input and Ichino's character names must be converted before substituting a representation in this formula. For comparison, Ichino's \S8.2, p.37, uses an incomplete $L$-function twisted by $\chi_V^{-1}$ and, for dimensions $2\to3$ and tempered source, divides it by $L^S(2,1)L^S(3,\epsilon)$ while retaining the local matrix-coefficient integrals at the excluded places. These are different displayed normalizations of the standard method, not an asserted equality of the named characters.

Thus both arrows in a proposed rank-one to rank-two to rank-three global construction have a specific value-at-one criterion, together with finite local occurrence. This is not a central-value criterion. It still must be shown that the resulting rank-three representation contains the prescribed canonical theta vector, rather than merely some canonical vector. For a fixed rank-three representation, the reverse first-occurrence question is the pole side of the Rallis theory; an upward criterion for an independently chosen $\Sigma$ does not establish that reverse occurrence automatically.

The definite side of the small choice in (B15a) has a different rank pattern: $A_2\to S_2\to X_2$, with $\operatorname{rank}X_2=2$. Its two arrows have equal rank. The adjacent-rank criterion (B15i) does not cover them; their standard Rallis parameter would be the central point, and their nonvanishing requires its own argument. In particular, the strict range of the archive's \nolinkurl{SphericalVolumeLift.Convergent} just recorded in B8 is not satisfied. No claim that every definite-side lift is nonzero is part of the local calculation.

\paragraph{The universal identity and the arithmetic equation left to prove.}
The last-step analytic identity is already a universal source theorem, rather than a large-gap instance. In \path{Proofs/DecoratedLatticeProofs/DoubleLiftPeriodRankZeroFubiniCompactQuotient.lean}, line 136, the source proves \texttt{lastLaw : \ensuremath{\forall} D, LastLaw D}; in \nolinkurl{LastSeeSawThetaIdentity.lean}, line 22, it proves \texttt{seeLast : \ensuremath{\forall} D, SeeLast D}. The proof bodies use compact quotient integration and the exact tensor identity with (B15b). With the requisite Weil data and
\[
 \mu_{V_1}=(\operatorname{emb}_0)_*\mu_{X_1},
 \tag{B15k}
\]
which is \texttt{Comm.VolLast} in \texttt{HzUnit.lean}, lines 23--24, the theorem \nolinkurl{GlobalTheory.LastRed.main}, in \nolinkurl{LastTwoStepsDoubleLiftPeriodEqIntegral.lean}, line 293, gives
\begin{equation}
\begin{gathered}
 B_A\bigl(\varphi,
   \Theta_{V\leftarrow S}(\Theta_{S\leftarrow T}(F;\Phi_2);\Phi_1)\bigr)
 \\
  =\int_{[T]} F(t)\,
       \operatorname{biLift}_{A\to S\to X_2}
       (\varphi;C_{\Phi_1,\Phi_2})(\iota(t,1))\,dt.
\end{gathered}
 \tag{B15l}
\end{equation}
This is the untranslated specialization; the source additionally gives every finite/right translate explicitly. The contracted kernel is its actual \texttt{lastBi}, defined in \texttt{CombAtSingleZero.lean}, lines 29--34: on the decomposition $X_2=T\perp L_z$ it is
\[
 C_{\Phi_1,\Phi_2}(u;v_T,v_z)
       =\Phi_2(v_T)\Phi_1(u,\operatorname{ofAdj}_{L,S}(\tau v_z)).
\]
All characters and measures remain in this identity. There is no arbitrary comparison scalar at this stage. The required global input is the concrete double theta lift appearing on its left.

The additional equation needed to identify these analytic kernels with the arithmetic columns is explicit. In \path{Definitions/LatticeFrames/OnFramesRdApply.lean}, lines 108--113, the condition \texttt{ClsK.Core} is
\begin{equation}
\begin{gathered}
 \operatorname{volLev}(X_j,U_{j,\alpha})\,
 \theta_{K,j}(h_{j,\alpha,a},h_{j+1,\gamma,a'};\Phi_{j,P})
 \\
 =c_{K,j}\,\mathcal A_{j,\alpha,\gamma,a,a'}
       \bigl(\operatorname{mvZ}(1,s_{j,\tau}I)(\operatorname{castP}P)\bigr).
\end{gathered}
 \tag{B15m}
\end{equation}
Here $\mathcal A$ is the specified \texttt{abelian\ensuremath{_2}} functional, not an arbitrary linear functional defined to make the equality true. Its class sums carry the reciprocal stabilizer orders in lines 115--121. The terminal condition \texttt{ClsT.Is.tail}, lines 147--150, similarly includes the top-level covolume and stabilizer masses. Their conjunction with the real branching and level conditions is \texttt{ClsPinD}, lines 152--153, from which the source derives the corresponding \texttt{KerX} and \texttt{PerX} laws. The universal Fubini identity (B15l) does not prove (B15m), supply those arithmetic functionals, or establish that the given 032 vector lies in the indicated double-lift image. These are now concrete representation and kernel/covolume equations, rather than an unspecified appeal to global matching.


% THETA3-B16 contribution by /root/definite_bridge, 21 November 2026.
\subsection{THETA3-B16: a pole criterion that also supplies finite coherence}\label{B16}

The value-at-one criteria in B15 start with an automorphic representation on a chosen smaller group. There is also a useful converse route: a pole on the rank-three group can construct a coherent smaller group. Its character conventions must be retained.

\paragraph{A precise conditional descent statement.}
Use exactly the notation of Jiang--Wu, \emph{On $(\chi,b)$-factors of cuspidal automorphic representations of unitary groups I}, arXiv:1409.0767v3, \S3.1 and Theorem 3.3(1), pp.15--17. Let $X$ be a skew-Hermitian space of rank three, let $\chi$ be a unitary Hecke character of $E$ with $\chi|_{\A_F^\times}=1$, and let $\sigma$ be irreducible and cuspidal on $\U(X)(\A_F)$. Write
\[
 \rho=\sigma\otimes\chi^{-1}
 \tag{B16a}
\]
in the source's notation for its twist on the unitary group. Its global integral is bilinear (\S3.1, p.15), so the corresponding input in the local Howe quotient convention is $\rho^\vee$. Suppose that:
\begin{enumerate}
\item $L^S(s,\sigma\times\chi)$, with the source's convention, has a pole at $s=1$;
\item $\rho$ has a real component which is not one dimensional;
\item at each real place, the only nonzero rank-two local theta lift of the corresponding input, with this splitting and additive character, is on a definite signature.
\end{enumerate}
Then there is a global Hermitian space $S$ of rank two, definite at every real place, for which the global theta lift of $\rho$ to $\U(S)$ is nonzero and cuspidal. In particular its finite local isometry classes are coherent, and its local theta lifts are nonzero at all finite places. No finite local form or lattice has been prescribed in this statement.

\textbf{Proof.}
In Jiang--Wu's Theorem 3.3(1), take their $m=3$, $\epsilon_\chi=0$, and $j=1$. Its distinguished point is
\[
 \tfrac12(m+1-\epsilon_\chi)-j=1,
\]
and it gives $\operatorname{LO}_{\psi,\chi}(\rho)\leq2$. The lowest occurrence minimizes the first-occurrence dimension over global Hermitian Witt towers of the prescribed even parity; see \S3.2, p.16. Therefore the dimension realizing the minimum is zero or two. At dimension zero the theta series is the determinant character $\nu_\chi\circ\det$, as the source states immediately before \S3.2. A nonzero global pairing with it would force every local component of the irreducible input to be one dimensional. Hypothesis 2 excludes this. Thus the lowest occurrence has dimension two, and by its definition its global theta lift is nonzero and cuspidal.

Nonzero global theta occurrence implies the corresponding local occurrences. Hypothesis 3 forces the real signatures of the constructed space to be definite. The theorem constructed a global space, so its finite local forms already satisfy the global coherence relation. This proves the assertion. \hfill$\square$

This result replaces separate finite-place guesses by one global pole test, once its representation and character dictionary has been established. It does not assert that the pole is known for the 032 wedge. In particular, the twist in (B16a) must be resolved against the bilinear integral in Jiang--Wu's \S3.1 and the contragredient in the local Howe quotient before substituting the canonical representation in (B15d). The local computations of B15 give the required unique definite occurrence after that convention matching; they do not by themselves prove hypothesis 1.

\paragraph{The invariant under a compensating splitting change.}
This remaining test is not removed by the independence of the special-unitary action from splitting characters. Let $\tau$ be a global character of $\operatorname{Res}^{1}_{E/F}\mathbf G_m$ and put
\[
 \tau_E(z)=\tau(z/\overline z),\qquad
 \Pi'=\Pi\otimes(\tau\circ\det),\qquad
 \chi'=\chi\tau_E.
\]
In the standard local-parameter convention of (B15c), changing the splitting by $\tau_E$ multiplies the input action by $\tau\circ\det$. Hence the compensating simultaneous twist transports the same local Howe correspondence on the same underlying Hermitian spaces. This is the global version of (B15g). Jiang--Wu, equation (3.1), p.15, gives the same determinant-character rule, with $\nu_\kappa(z/\bar z)=\kappa(z)$.

There is an equally explicit invariant of unramified standard parameters. Whenever standard base change is used to write them, one has
\[
 \operatorname{BC}(\Pi')\otimes(\chi')^{-1}
   =\operatorname{BC}(\Pi)\otimes\chi^{-1}.
 \tag{B16b}
\]
Indeed, twisting a representation by $\tau\circ\det$ tensors its standard base-change parameter by $\tau_E$; the two additional factors then cancel. This also proves equality of the corresponding incomplete standard $L$-functions term by term at every unramified place. Formula (B16b) is stated in this parameter convention; it is not an unexplained replacement of Jiang--Wu's named $L$-function by a differently twisted one.

To change the relative splitting further by a character $\eta$ changes the right side of (B16b) by $\eta^{-1}$. That is a different standard parameter and a different pole question. The original nonzero A7 lift supplies information for its original source $P$ and characters. The compensating transport preserves $P$, including any real signature $(1,1)$; an additional relative change has not inherited that nonvanishing. This identifies the missing representation-theoretic calculation without asserting that another definite realization is impossible.

Finally, a successful application to each of several rank-three representations separately need not produce the same finite local source forms. The \texttt{Comm} in B13 is common to the whole family. Its use still requires common spaces and real action across tuples, the specified finite lattice data, and the kernel/covolume equation (B15m). The pole construction therefore closes a single-representation coherence problem under its hypotheses; it does not supply the entire arithmetic family by itself.


% THETA3-B17 contribution by /root/definite_bridge, 21 November 2026.
\subsection{THETA3-B17: the actual wedge admits an adjoint double lift and last-step seesaw}\label{B17}

The compact-quotient condition in the analytic last-step theorem is
strictly weaker than the totally definite condition in the arithmetic
\texttt{MainX} family. This distinction permits a complete analytic
specialization for the actual multiplicity plane of A7. The second
lift can be its adjoint, with terminal space equal to the original
rank-three space.

\paragraph{Fixed data and the finite-rank operator.}
Use the rank-three case of A7, with its exact scalar automorphic
realization of the canonical coefficient. Thus $V$ has signature
$(2,1)$ at $\mu$ and $(3,0)$ elsewhere, $[F:\Q]\geq2$, and the
multiplicity plane $P=P_1\perp P_2$ is definite at $\mu$.
Write
\[
 G=\U(V),\qquad S=\U(P).
\]
Both adelic quotients $[G]$ and $[S]$ are compact. For $V$, use
the definite component at a place other than $\mu$; for $P$, use
its definite component at $\mu$. A rational isotropic vector would
remain a nonzero isotropic vector in that definite completion.
The quotient compactness is the anisotropic case of the usual
arithmetic compactness theorem already used in A7.

For the nonvanishing assertion about a particular wedge, first
choose its rational negative line $L$ as in B12 and rebase its
canonical coefficient frame at $L$. Make this choice before
fixing the kernel and the operators below. The original
coefficient at an unrelated base point need not be nonzero,
and its maximal compact subgroup need not contain $\U(L^\perp)$.
After this rebasing the following data remain common to all
source characters under consideration.

Fix the additive character, the two splitting characters
$(\chi_V,\chi_P)$, the theta input $\Phi$, finite levels, the
canonical coefficient frame, and positive Haar quotient measures
$dg,dh$. Let
\[
 K(g,h)=\theta_{V,P}(g,h;\Phi).
\]
The coefficient identification here is the one in A7: after
evaluation on the fixed canonical frame, this is an actual scalar
theta kernel, not a separately chosen function having the same
cohomology class. More explicitly, let $\Psi$ be the holomorphic
top-degree relative Lie algebra cochain and choose a unit vector
$\kappa$ in its one-dimensional canonical tangent-frame line.
Its scalar Schwartz coefficient is $\Phi=\Psi(\kappa)$; evaluation
of the theta cochain on the frame translated by $g$ is
$\theta(\omega(g,h)\Phi)$. This fixes the archimedean coefficient
pairing before any adjoint is taken. A fixed conversion between
the Hermitian/skew-Hermitian presentation and the archive's
Gram-matrix presentation is part of the oscillator identification;
the subsequent dagger is applied to this very coefficient.
By A7, $K(\,\cdot\,,h)$ belongs to one
finite-dimensional space $\mathcal E$ of holomorphic canonical
coefficients for every $h$. Its smooth automorphy and the
compactness of the two quotients imply that $K$ is bounded on
$[G]\times[S]$.

Equip $\mathcal E$ with the Petersson inner product inherited from
$L^2([G],dg)$, conjugate linear in the first variable. Define
\begin{equation}
\begin{aligned}
 \mathcal T:L^2([S],dh)&\longrightarrow\mathcal E,\\
 (\mathcal T u)(g)&=\int_{[S]}K(g,h)u(h)\,dh,\\
 \mathcal E_0&=\operatorname{im}\mathcal T.
\end{aligned}
\tag{B17a}
\end{equation}
This is the archive's bilinear \texttt{liftS} convention. Boundedness
follows from Cauchy--Schwarz on the product of the two finite-measure
quotients. Its range is finite dimensional and therefore closed.
Moreover $K(\,\cdot\,,h)\in\mathcal E_0$ for every $h$:
if $f\in\mathcal E$ is orthogonal to the range, then its
continuous coefficient $\langle f,K(\,\cdot\,,h)\rangle_G$
vanishes almost everywhere by testing against all $u$, hence
everywhere. Taking the orthogonal complement proves the claim.
All the A7 character integrals with these fixed data lie in
$\mathcal E_0$: A7 writes the integral for a character $\pi$ as
$f_\pi=\mathcal T(\overline{v_\pi})$.

\paragraph{The adjoint is another genuine theta lift.}
The Hilbert-space adjoint is
\[
 (\mathcal T^*f)(h)
   =\int_{[G]}\overline{K(g,h)}f(g)\,dg.
 \tag{B17b}
\]
Indeed, expanding the inner products and using Fubini gives
$\langle\mathcal T u,f\rangle_G
 =\langle u,\mathcal T^*f\rangle_S$.
In particular $\mathcal T^*f$ is a smooth automorphic function
when $f\in\mathcal E$.

This adjoint kernel is a theta kernel for the exchanged pair with
the same underlying spaces. Since the canonical compact type is
one dimensional, conjugating the chosen scalar coefficient is
exactly the adjoint operation on that coefficient line; there is
no unrecorded sum over compact-type coordinates. Any volume of
the compact fibre is already included in $dg$ and in the
inherited Petersson norm. In the recovered source let
\[
 d_{V,P}(x)=G_V^{-1}x^*G_P,\qquad
 \Phi^\dagger(w)=\overline{\Phi(d_{P,V}(w))},
 \qquad \omega^\dagger=\operatorname{conjAdjoint}(\omega).
\]
Here $G_V,G_P$ are the Gram matrices and $x^*$ is conjugate
transpose. The exact source rule is
\begin{equation}
\begin{gathered}
 \omega\text{ has data }(V,P;\chi_V,\chi_P)
 \quad\Longrightarrow\quad
 \omega^\dagger\text{ has data }(P,V;\chi_P,\chi_V),\\
 \theta_{P,V}^{\dagger}(h,g;\Phi^\dagger)
       =\overline{\theta_{V,P}(g,h;\Phi)}.
\end{gathered}
\tag{B17c}
\end{equation}
There is no inverse on the displayed exchanged splitting
characters in this archive convention.

For precision,
\path{Development/AbelianVarietyTheory/BlockTheory/TwistedDiscreteSeries/ConjugateAdjointDagger.lean}, lines
112--124, defines the rational involution and the conjugated
exchanged representation. Lines 144--156 give its
antisymplectic, theta-function, and Weil-data clauses.
The following proof files have the full directory prefix
\path{Development/AbelianVarietyTheory/Proofs/DecoratedLatticeProofs/}.
The actual proof \nolinkurl{ConjWeilDatumViaAdjointExchangedPair.lean},
theorem \nolinkurl{isWeilDatumO_conjAdjoint}, proves the exchanged
Weil data, including the split-place normalization.
\nolinkurl{HermitianAdjointSymplecticActionThetaHeisenberg.lean}
proves the rational, antisymplectic, and theta-function
transport. The complete rule is proved by
\nolinkurl{WeilConjugationAdjointRuleStandardChar.lean},
theorem \texttt{conjAdjoint\_holds}, and exposed by
\nolinkurl{GlobalTheory.T.weilConjugationAdjointRule_standard_char}
in its statement file.

One can also check the kernel equality directly. The involution
$d$ carries the rational lattice bijectively onto the exchanged
rational lattice, reverses the symplectic form, and exchanges
the two group actions. Complex conjugation reverses the additive
character at the same time, so the standard additive character
is retained on the exchanged pair. The defining formula
$\omega^\dagger(h,g)\Phi^\dagger
 =\overline{\omega(g,h)\Phi}\circ d_{P,V}$
then gives (B17c) at zero, equivalently after summing the
rational Schwartz model. Thus (B17b) is the bilinear theta lift
from $V$ to $P$ with input $f$ and kernel $\Phi^\dagger$.

\paragraph{An exact double lift, common to all source characters.}
On $\mathcal E_0$ put
\[
 Q=\mathcal T\mathcal T^*.
\]
This is positive definite. In fact
$\langle f,Qf\rangle_G=\|\mathcal T^*f\|_S^2$;
if it vanishes, then $f$ is orthogonal to
$\operatorname{im}\mathcal T=\mathcal E_0$, and hence $f=0$.
Consequently, for every $f\in\mathcal E_0$,
\begin{equation}
\begin{gathered}
 F_f=Q^{-1}f,\qquad
 f=\mathcal T\mathcal T^*F_f\\
   =\Theta_{V\leftarrow P}
          \bigl(\Theta_{P\leftarrow V}(F_f;\Phi^\dagger);
                \Phi\bigr).
\end{gathered}
\tag{B17d}
\end{equation}
All functions in this identity are continuous automorphic
functions on compact adelic quotients. It requires neither a new
global nonvanishing theorem nor a rank-one reverse descent:
the second source is $T=V$, of rank three.

The same $Q$ works for every A7 source character with the fixed
spaces, kernel, splittings, levels, coefficient frame, and
measures. Thus (B17d) is a common linear operator identity on the
whole space $\mathcal E_0$, rather than a normalization chosen
separately for each vector or character. Changing the complete
types may change $P$; changing the Schwartz input may change
$Q$. Uniformity across those changes is not asserted.

The exact mixed-pairing identities are
\begin{equation}
\begin{aligned}
 \langle f_1,f_2\rangle_G
    &=\langle\mathcal T^*Q^{-1}f_1,
                        \mathcal T^*f_2\rangle_S\\
    &=\langle\mathcal T^*Q^{-1/2}f_1,
                        \mathcal T^*Q^{-1/2}f_2\rangle_S .
\end{aligned}
\tag{B17e}
\end{equation}
They follow by adjunction and self-adjointness of $Q$.
The positive square root is defined by the finite-dimensional
spectral theorem. In particular, the second identity is an
isometry from $\mathcal E_0$ to a space of actual reverse theta
lifts on the original anisotropic $P$. The conversion from this
Petersson pairing to the geometric canonical cup-product
pairing uses the fixed coefficient frame and quotient measure;
it is the same conversion for every $f_\pi$.

For completeness, when $\mathcal E_0\ne0$, fix an orthonormal basis $f_1,\ldots,f_r$ of
$\mathcal E_0$ and write $K(g,h)=\sum_i f_i(g)e_i(h)$. The matrix
of $Q$ is the positive Gram matrix
\[
 Q_{ij}=\int_{[S]}e_i(h)\overline{e_j(h)}\,dh .
\]
With $t=\operatorname{Tr}Q$ and $\Delta=\det Q>0$, for $r\geq2$
its least eigenvalue satisfies
\[
 \lambda_{\min}(Q)\geq
       \Delta\left(\frac{r-1}{t}\right)^{r-1},
 \qquad
 \|Q^{-1}\|\leq
       \frac{(t/(r-1))^{r-1}}{\Delta}.
 \tag{B17f}
\]
Indeed the product of the other $r-1$ eigenvalues is at most
$(t/(r-1))^{r-1}$ by the arithmetic--geometric mean inequality.
For $r=1$, $\|Q^{-1}\|=\Delta^{-1}$.
Also
$\|\mathcal T^*Q^{-1}f\|_S^2
 =\langle Q^{-1}f,f\rangle_G
 \leq\lambda_{\min}(Q)^{-1}\|f\|_G^2$.
These give an explicit analytic matrix certificate. B7's
arithmetic determinant bound applies only if its arithmetic
normalization hypotheses for this Gram matrix are separately
verified.

\paragraph{Specialization of the actual rank-zero Bessel identity.}
Choose a rational negative line $L\subset V$ at $\mu$ and put
$A=L^\perp$, $H=\U(A)$. Then $A$ is rank two and totally
positive definite. Put
\[
 L_z=-(N_{E/F}\xi_0)^{-1}L,\qquad
 T=V,\qquad X=T\perp L_z .
\]
The identity map between the underlying one-dimensional vector
spaces supplies the exact line relation
$G_L=-N_{E/F}(\xi_0)G_{L_z}$ required by
\texttt{SeeLast}. Use identity rank-zero Witt decompositions
for $V$ and $P$ and the same measures on the identified copies.
The zero-dimensional unipotent measure has mass one.

Choose a splitting $\chi_z$ for $L$ and put
$\chi_A=\chi_V\chi_z^{-1}$. The orthogonal source decomposition
$V=A\perp L$ then has character identity
$\chi_V=\chi_A\chi_z$. The exchanged second pair in (B17c)
has characters $(\chi'_P,\chi_T)=(\chi_P,\chi_V)$, so the
upper pair $(P,X)$ has characters
\[
 (\chi'_P,\chi_T\chi_z^{-1})=(\chi_P,\chi_A),
 \qquad (P,L_z)\text{ has }(\chi_P,\chi_z^{-1}).
\]
These are exactly the character equations in
\path{Development/AbelianVarietyTheory/Definitions/LatticeFrames/IsQuasiFam.lean},
lines 114--117. As in B15f,
such global splitting characters exist; their selection is
part of the fixed input, with no assertion about prescribed
unramified finite values.

Let $e$ be a continuous automorphic test function on $[H]$, put
$\varphi=\overline e$, and use the rank-zero Bessel period
\[
 B_e(f)=\int_{[H]}f(\iota_A(a))\overline{e(a)}\,da .
\]
In the notation of (B15l), let
$C=C_{\Phi,\Phi^\dagger}=\texttt{lastBi}(\Phi,\Phi^\dagger)$
be the exact contracted two-variable theta input. Its formula
on $X=V\perp L_z$ is
\[
 C(u;v_T,v_z)
     =\Phi^\dagger(v_T)\,
       \Phi\bigl((o.\operatorname{res}_S)^{-1}
                       (u,\operatorname{ofAdj}_{L,P}(v_z))\bigr).
\]
The line map here is the identity chosen above; the source
definition incorporates its Gram matrices and $\xi_0$ in
$\operatorname{ofAdj}$. Define $\mathcal J_e(x)$ to be the
archive's upper double lift
$\operatorname{biLift}_{A\to P\to X}(\overline e;C)(x)$,
with the same $da,dh$ as above. Then the exact identity is
\begin{equation}
 B_e(f)=\int_{[G]}(Q^{-1}f)(t)\,
                 \mathcal J_e(\iota_T(t,1))\,dt
       \qquad(f\in\mathcal E_0).
\tag{B17g}
\end{equation}

\textbf{Proof.}
Set
$\beta_e(h)=\int_{[H]}\overline{e(a)}K(\iota_A(a),h)\,da$.
Substitute (B17d) in the left side and use Fubini on
$[H]\times[S]\times[G]$. The result is
\[
 \int_{[G]}F_f(t)\int_{[S]}
       \overline{K(t,h)}\,\beta_e(h)\,dh\,dt.
\]
All factors are continuous on the respective compact quotients,
so their product is absolutely integrable. Formula (B17c)
identifies the first kernel in the inner integral with
$\theta^\dagger_{P,V}(h,t;\Phi^\dagger)$.
The orthogonal tensor identity and the line relation now give
\[
 \mathcal J_e(\iota_T(t,1))
       =\int_{[S]}\theta^\dagger_{P,V}(h,t;\Phi^\dagger)
                         \beta_e(h)\,dh .
\]
This is precisely the integrand equality of
\texttt{SeeLast}, after integration over $[H]$ and $[S]$.
It proves (B17g).

These are the actual hypotheses used by the recovered analytic
theorems. Both definition files in this paragraph have directory
prefix
\path{Development/AbelianVarietyTheory/Definitions/LatticeFrames/}.
\texttt{IsThetaFam.lean}, line 17, defines
\texttt{CoCpt} by a compact set of adelic representatives;
its structures \texttt{BesLev}, \texttt{JacLev}, and
\texttt{Last} ask for this condition, not for definiteness at
every real place. \texttt{IsQuasiFam.lean}, lines 71--75,
states \texttt{LastLaw}, and lines 110--123 state
\texttt{SeeLast}. The complete proofs, in
\path{Development/AbelianVarietyTheory/Proofs/DecoratedLatticeProofs/}, are
\nolinkurl{DoubleLiftPeriodRankZeroFubiniCompactQuotient.lean},
lines 136--159, and \nolinkurl{LastSeeSawThetaIdentity.lean},
lines 22--54. The proof above instantiates their analytic
content with $S=P$ and $T=V$. It does not require the
additional \texttt{Comm.Is} hypotheses used to package
\nolinkurl{GlobalTheory.LastRed.main}. \hfill$\square$

\paragraph{What is now obtained, and what remains arithmetic.}
For a nonzero actual A7 wedge, make the B12 choice of rational
$H$ and rebase the coefficient frame before fixing $K,Q$, as
specified above. B12 then supplies an $e$ for which $B_e(f)\ne0$.
Equation (B17g) is therefore a
nonzero actual last-step Bessel seesaw involving that wedge.
It is uniform in the source characters for the fixed data.
The identity uses a common positive operator $Q^{-1}$, not
an assertion that this operator is a scalar or is algebraic.

Equation (B17g) is an identity for one Bessel coefficient.
The entire geometric mixed period is recovered from the
Petersson pairing in (B17e), with the fixed geometric measure
conversion. If a description solely through definite Bessel
coefficients is wanted, apply the finite Bessel frame of A5 to
$\mathcal E_0$, and then apply (B17g) to each of its coefficients.
For an adelic translate in that frame, translate the first
kernel by the same group element; \texttt{LastLaw} permits the
two kernels to be different. This combines actual last-step
seesaws to recover the whole pairing. A single restriction
period is not being identified with the whole ball integral.

The upper space is completely determined here. Since
$N_{E/F}\xi_0$ is a norm, $L\perp L_z$ is a hyperbolic plane:
if $e$ and $e_z$ have Gram entries $a$ and
$-a/N_{E/F}\xi_0$, the vector $e+\xi_0e_z$ is isotropic.
Thus
\[
 X\simeq A\perp\mathbb H,\qquad
 \operatorname{rank}X=4,\qquad
 \operatorname{sig}_\tau X=(3,1)
        \text{ for every real }\tau .
\]
Its Witt index is one, because $A$ is definite at every real
place. The quotient $[X]$ is not being integrated over in
(B17g); evaluation on the embedded compact $[G]$ is sufficient.
This feature causes no convergence gap in the displayed
identity.

The arithmetic \texttt{MainX} family, by contrast, asks for
totally definite spaces, specified arithmetic functionals,
and the common kernel/covolume and terminal mass equations
(B15m). Neither the mixed $P$ nor this rank-four $X$ is replaced
by a definite inner form in the argument. The result closes
the analytic double-lift and last-step comparison for the
actual A7 kernel while keeping those arithmetic matching
conditions separate. In particular it proves no algebraicity
of $Q^{-1}$ and no new projected algebraic cycle.


% THETA3-B18 contribution by /root/definite_bridge, 21 November 2026.
\subsection{THETA3-B18: the remaining definite alternatives have a finite parameter test}\label{B18}

The analytic identity of B17 uses the actual multiplicity plane and
is unconditional on a new global descent. A different question is
whether a prescribed irreducible rank-three constituent also comes
from a totally definite rank-two group. The following statements
make that question more explicit. The global base-change hypotheses
below are inputs, not calculations of the actual A7 constituent.

\paragraph{The complete real deletion table in the canonical case.}
Use B15's real convention, including $\psi(x)=e^{-2\pi i x}$.
Let the rank-three input have Harish--Chandra parameter
$(1,0;-1)$ at $\mu$, and $(1,0,-1;\varnothing)$ at a compact
place. Write the target rank-two splitting character as
$\chi_2(z)=(z/|z|)^{m_0}$, with $m_0$ even; the other character
has any fixed odd exponent $n_0$. A rank-two lift is nonzero
only when $m_0\in\{-2,0,2\}$. Its unique possible signature is
then given by
\[
\begin{array}{c|c|c|c|c}
 m_0&\lambda_\mu-m_0/2&(r_\mu,s_\mu)
       &\lambda_\tau-m_0/2&(r_\tau,s_\tau)\\ \hline
 -2&(2,1;0)&(2,0)&(2,1,0;\varnothing)&(2,0)\\
 0 &(1,0;-1)&(2,0)&(1,0,-1;\varnothing)&(1,1)\\
 2 &(0,-1;-2)&(1,1)&(0,-1,-2;\varnothing)&(0,2)
\end{array}
\tag{B18a}
\]
All six listed lifts are nonzero. In particular total definiteness
requires $m_{0,\mu}\in\{-2,0\}$ and
$m_{0,\tau}\in\{-2,2\}$ elsewhere. The other splitting character
changes the output by a determinant twist and does not change this
table.

\textbf{Proof.}
Ichino's Theorem 4.1(ii), pp.11--12, with $n=3,m=2$, requires a
zero in the shifted parameter: its condition $p_1+q_1>0$ is
necessary. This proves the asserted restriction on $m_0$.
Delete that zero. Each remaining positive entry in the first
block and negative entry in the second block contributes to the
positive target block; the other entries contribute to the
negative target block. The same theorem gives exactly the six
signatures in (B18a), so no other rank-two signature is possible.
For nonvanishing, apply Atobe's Definition 1.6 and Theorem 1.7(2),
pp.4--5. In five cases the zero chain has length one and the
listed signature is its base point $(r_\lambda,s_\lambda)$.
For the compact input with $m_0=0$, the maximal chain is instead
$(1,0,-1)$ of length three; its base point is $(0,0)$, and the
listed $(1,1)$ is obtained by the theorem's equal increment
$l=1$. This checks that case without mistaking its chain length
for one. The output Harish--Chandra parameter is obtained from
the retained entries by adding $n_0/2$, as in B15.
\hfill$\square$

The conversion between directions is explicit. Start with a
Hermitian rank-three $V$ and skew-Hermitian $W$ of rank two.
For a trace-zero scalar $\delta$, switch to the Hermitian space
$(W,\delta^{-1}h_W)$ and the skew-Hermitian space
$(V,\delta h_V)$. The tensor symplectic form is unchanged.
Ichino, \S2.2, p.5, states that this switches the two factors of
the same oscillator. At a real place take $\delta=i$, or a
positive multiple of it in the fixed complex embedding, so the
signatures in (B18a) are precisely those of the Hermitian
multiplicity space. The splitting characters switch their names;
none is inverted in this factor-switch rule. These identifications
must be made before applying the table to the archive convention.

\paragraph{The exact finite-place parameter and inner-form rule.}
Now label the original pair by $V_3$ Hermitian and $W_2$
skew-Hermitian, with characters $\chi_3|_{\A_F^\times}=\epsilon$
and $\chi_2|_{\A_F^\times}=1$. In the local quotient convention
\[
 \Hom_{\U(V_3)\times\U(W_2)}
       (\omega,\Pi\boxtimes\Sigma)\ne0,
\]
Gan--Ichino, \emph{The Gross--Prasad conjecture and local theta
correspondence}, arXiv:1409.6824v2, Theorem 4.4, pp.14--15,
gives at every nonsplit finite place
\[
 \phi_{\Pi_v}
    =(\phi_{\Sigma_v}\otimes\chi_{3,v}^{-1}\chi_{2,v})
                                      \oplus\chi_{2,v}.
 \tag{B18b}
\]
This theorem covers all local $L$-parameters, not just tempered
ones. For fixed $\Pi_v$ and fixed splitting characters, the
smaller parameter is obtained by deleting one occurrence of
$\chi_{2,v}$. The theorem's packet bijections then determine a
unique smaller representation, including its Hermitian inner
form. If the smaller parameter does not contain $\chi_{3,v}$,
use Theorem 4.4(i)(b) for the fixed target sign; if it does,
use Theorem 4.4(ii)(c). The enhanced packet character is the
restriction prescribed in \S4.6, proved by Theorem 1.3. This
last assertion is a finite-place statement; (B18a) supplies the
real-place assertion separately.

The same parameter formula holds at split finite places, where
there is a unique space of each dimension. One can check the
convention rather than silently omit a contragredient.
M\'inguez, \emph{Correspondance de Howe explicite: paires duales
de type II}, Ann. Sci. ENS 41 (2008), Corollary 6.3, p.734,
gives $\phi^\vee\oplus1$ for his untwisted matrix model in
ranks $2\to3$. His equation (1.4) and footnote on p.721 state
that replacing the first action by transpose inverse replaces
the input by its contragredient. Thus the switched model has
parameter $\phi\oplus1$. Twisting its first action by
$\chi_3\circ\det$ and its second by $\chi_2\circ\det$ gives
(B18b). This is the split normalization in Ichino's \S8.1.3,
Lemma 8.4, p.36; M\'inguez's statement also covers non-tempered
inputs. In particular (B18b) is available at every unramified
place, without a Ramanujan assumption.

\paragraph{A conditional global uniqueness test.}
Suppose that a fixed irreducible automorphic $\Pi$ on $\U(V_3)$
has the real components specified above and occurs in a global
theta pair with an irreducible automorphic $\Sigma$ on $\U(W_2)$.
Here $\Sigma$ labels the local oscillator quotient in the
preceding display. In the archive bilinear integral of B17, the
integrating automorphic input has type $\Sigma^\vee$, or its
chosen unitary conjugate realization. Assume that the standard
weak base changes $\mathcal B_3$ and $\mathcal B_2$ are specified isobaric
sums of unitary cuspidal automorphic representations of total
ranks three and two over $E$, with the local parameters in (B18b)
at almost all places. Then
\[
 \mathcal B_3=
   (\mathcal B_2\otimes\chi_3^{-1}\chi_2)\boxplus\chi_2.
 \tag{B18c}
\]
Indeed (B18b) identifies the unramified parameters at almost all
places, and strong multiplicity one for unitary isobaric
representations gives (B18c). The version being used is the
usual finite-exception case of Rajan, \emph{On strong
multiplicity one for automorphic representations},
arXiv:math/0210235, Theorem 1, p.2, and its isobaric extension
stated after the proof on p.6.

If $\mathcal B_2$ is cuspidal on $\mathrm{GL}_2$, then
$\chi_2$ is the unique one-dimensional constituent of the
isobaric cuspidal support in (B18c). Consequently any other
rank-two global realization of this same $\Pi$, in the same
standard-parameter convention and with base changes of the
same specified isobaric kind, must have the same rank-two character
$\chi'_2=\chi_2$. To prove this, apply (B18c) to the second
realization and compare the unique isobaric cuspidal supports.
Its displayed $\mathrm{GL}_1$ summand must be the unique one
in the first support.

Changing $\chi_3$ while fixing $\chi_2$ only twists the smaller
representation by the corresponding determinant character.
The oscillator action on the smaller group is multiplied by
that character, whereas its action on $\U(V_3)$ is unchanged.
Thus the local space is unchanged. More explicitly, write
$\chi'_3/\chi_3=\eta$, and use Hilbert 90 to define
$\nu_\eta(z/\bar z)=\eta(z)$. The source representation is
transported by $\nu_\eta\circ\det$; its parameter in (B18b)
is multiplied by $\eta$, exactly cancelling $\chi_3'^{-1}$.
The finite-place packet bijection and the real table then
force the original smaller local isometry classes.
Therefore, under these base-change and cuspidality hypotheses,
a mixed real place of the original smaller space obstructs a
totally definite rank-two realization of the same $\Pi$.
This is a conditional obstruction for the specified constituent;
no cuspidal $\mathrm{GL}_2$ transfer has been computed here for
an actual A7 wedge.

A compensating determinant twist of the rank-three realization
does not remove this obstruction. It multiplies every summand
of $\mathcal B_3$ and the distinguished $\chi_2$ by the same
character. At a real place it adds $a$ to every entry of
$\lambda$ and adds $2a$ to $m_0$, leaving
$\lambda-m_0/2$ and hence (B18a) unchanged. This recovers the
invariance proved in B16 at the level of the explicit table.

If instead $\mathcal B_2=\eta_1\boxplus\eta_2$ is a sum of
characters, (B18c) has at most three distinct character
constituents. Any alternate $\chi'_2$ must be one of them,
must have the allowed restriction to $\A_F^\times$, and must
pass (B18a) at every real place. For each candidate, deletion
in (B18b) and its enhanced packet rule determine the finite
local smaller representation. Coherence and global occurrence
remain to be proved for that candidate, for example through
the pole criterion of B16. This is a finite list of specified
character tests, not a claim that every deletion occurs globally.
Nontrivial Arthur $\mathrm{SL}_2$ blocks or base changes outside
the unitary isobaric hypothesis require their own analysis;
they are not covered by the preceding two cases.

\paragraph{An equal-rank central-value implication.}
There is also a precise sufficient condition for constructing a
new definite rank-two group from rank-two input. Use exactly
Jiang--Wu's notation in \S3.1--3.2 and Theorem 3.3(1), p.17.
Let $X$ be skew-Hermitian of rank two, let $\sigma$ be irreducible
cuspidal on $\U(X)$, and let $\chi$ be unitary with
$\chi|_{\A_F^\times}=1$.
Put $\rho=\sigma\otimes\chi^{-1}$ in their unitary-group
character convention. Suppose
\[
 L^S(1/2,\sigma\times\chi)\ne0,
 \tag{B18d}
\]
with a finite value, suppose $\rho$ has at least one local
component which is not one dimensional, and suppose that
at each real place every nonzero rank-two local theta lift
of the corresponding Howe input $\tau_v=\rho_v^\vee$ is
definite. Then some global definite
Hermitian $Y$ of rank two has a nonzero cuspidal theta lift of
$\rho$. Its finite local forms are coherent automatically.

\textbf{Proof.}
In that theorem take $m=2$, $\epsilon_\chi=0$ and $j=1$.
Then $m-\epsilon_\chi$ is even and
$(m+1-\epsilon_\chi)/2-j=1/2$. Hence
$\operatorname{LO}_{\psi,\chi}(\rho)\leq2$.
The dimension is zero or two. The rank-zero theta series is
a determinant character, as in Jiang--Wu's \S3.1; the local
hypothesis excludes its occurrence. The first occurrence is
therefore two and is nonzero cuspidal by the definition of
lowest occurrence. The assumed real local condition makes its
global space definite. The global construction supplies
coherence at the finite places. \hfill$\square$

Jiang--Wu's defining integral in \S3.1, p.15, is bilinear:
it integrates the theta kernel against $f$, with no complex
conjugation. Thus its automorphic integrating input $\rho$
corresponds to the local Howe quotient input
$\tau=\rho^\vee$. In the following local rules it is the
parameter and enhancement of $\tau$ that must be used.

The real condition in this last statement can itself be made
explicit when $\tau$ is discrete series at every real place.
For a mixed $\tau$ on $\U(1,1)$ with parameter $(a;b)$,
$a,b\in\Z+1/2$, $a\ne b$, choose the even exponent $m_0$ of
the target splitting character so that $m_0/2$ lies strictly
between $a$ and $b$. The equal-rank case $k=0$ of Ichino's
Theorem 4.1(ii) then gives the unique signature $(2,0)$ if
$a>b$, and $(0,2)$ if $a<b$. Atobe's base-point criterion
proves nonvanishing. For a compact input, choose $m_0/2$
outside the interval spanned by its two half-integral parameters;
the same rules again give a unique definite target and a
nonzero lift. Suitable integers exist in all these cases.
B15's character-extension lemma globalizes these even infinity
exponents, allowing finite ramification. If the desired input
is fixed as $\rho$, set $\sigma=\rho\otimes\chi$ so that
$\rho=\sigma\otimes\chi^{-1}$ holds in Jiang--Wu's convention.
Condition (B18d) must then be tested for this resulting pair
$(\sigma,\chi)$; the freedom to choose infinity characters
does not prove the central value is nonzero.

At a nonsplit finite place the equal-rank local test is also
explicit. In Gan--Ichino's conventions, their Theorem 4.1(ii),
p.13, says
\[
 \begin{gathered}
 \theta_{Y_v}(\tau_v)\ne0
 \quad\Longleftrightarrow\quad
 \epsilon\bigl(1/2,\phi_{\tau_v}\chi_{Y,v}^{-1},
                      \psi_{2,v}^E\bigr)
               =\epsilon(Y_v)\epsilon(X_v),\\
 \psi_{2,v}^E(x)=\psi_v(\operatorname{Tr}_{E_v/F_v}(\delta x)).
 \end{gathered}
 \tag{B18e}
\]
Its two sides determine one of the two finite inner forms.
The additive character in this epsilon factor has no factor
$1/2$; the $\psi^E$ used for the Whittaker normalization in
that paper does have that factor. These characters must not be
confused. When (B18d) applies, the global lowest-occurrence
theorem constructs a coherent collection of these finite forms;
otherwise guessing their signs separately does not prove
coherence or a nonzero global lift.

This central-value condition is different from B15's upward
value-at-one test and B16's rank-three pole-at-one test. It can
produce an equal-rank definite transfer when its hypotheses are
verified. It does not say that the resulting smaller
representation lifts back to the prescribed $\Pi$, or that it
satisfies the arithmetic kernel/covolume equation of B15.
In particular the conditional uniqueness obstruction above
must still be respected when trying to use such a transfer in
a rank-two-to-rank-three realization of the original wedge.


% THETA3-B19 contribution by /root/definite_bridge, 21 November 2026.
% Static source reading; no Lean build is asserted.
\subsection{THETA3-B19: definite periods are finite CM theta-jet functionals}\label{B19}

The remaining arithmetic comparison can be localized to explicit CM
abelian varieties.  This gives a common period normalization for a
fixed finite family of inputs and source characters.  It does not
assert that this normalization equals the particular constant of an
uninstantiated \texttt{MainX} datum.

\paragraph{The actual CM abelian varieties at the evaluation point.}
Keep the rank-three setting of 032 and B12.  Choose a rational
negative line $L=Ez$ and an orthogonal decomposition
\[
 V=Ee_1\perp Ee_2\perp Ez
\]
over $E$, with the first two lines positive at the active place.
They are positive at the other real places as well.  For the
multiplicity line $P_i$, put
\[
 d_i(\tau)=v_i(\tau^{-1}),\qquad
 d_i^{\mathrm{flip}}(\tau)=
 \begin{cases}-d_i(\tau)&\tau\in\{\mu,c\mu\},\\
 d_i(\tau)&\text{otherwise}.
 \end{cases}
\]
Use the geometric $E$-action on
$A_{x,i}=\mathcal F_i^\vee/\iota(\Lambda_i)$, where
$\iota(u)(v)=\psi_i(v,u)$ and $x=[Ez]$.
Then, as rational $E$-Hodge structures,
\begin{equation}\label{B19-CM}
 H^1(A_{x,i},\Q)
 \simeq U(d_i)^{\oplus2}\oplus U(d_i^{\mathrm{flip}}),
 \qquad
 A_{x,i}\sim_E B(d_i)^2\times B(d_i^{\mathrm{flip}}).
 \tag{B19a}
\end{equation}
Here the last relation is an isogeny, with the chosen polarizations
retained separately.  In particular it is not an assertion that the
polarization is the product principal polarization.

To prove the formula, use 032, \S4.2, equation (4.5), pp.14--15.
For the positive coordinate lines the homology-filtration block
uses $c\tau$ when $d_i(\tau)=1$ and $\tau$ when
$d_i(\tau)=-1$, including the active place.  The negative
coordinate agrees at every compact place and switches the active
pair.  The relation
\[
 \iota(au)(v)=\iota(u)(\overline a v)
\]
conjugates these embedding labels in geometric $H^1$.
This proves the first formula.  The equivalence between rational
weight-one Hodge maps and homomorphisms of abelian varieties up to
isogeny, used in 032, Lemma 2.2, proves the second.
The two-dimensional tangent block used by the first Taylor
coefficient comes from the two $B(d_i)$ factors.  The block has
the archive tangent label $\mu$, whereas the covariant Taylor
tensor has label $c\mu$, as in A8.  The inverse in the definition
of $d_i$ cannot in general be discarded.

The descent input is Milne, \emph{Complex Multiplication},
version 14 July 2020, Proposition 7.9 and Corollary 7.10,
printed pp.54--55: extension from an algebraically closed subfield
of $\C$ is fully faithful for abelian varieties and is essentially
surjective on CM abelian varieties.  Apply it to
$\overline\Q\subset\C$.  The abelian variety in (B19a), together with its polarization,
symmetric theta line, and any specified finite torsion data, has a
model over $\overline\Q$.  This assertion needs no identification
of its polarization with a product polarization.  Start with the
CM factors $B(d)$ over $\overline\Q$.  A rational Hodge
isomorphism gives a complex isogeny; its kernel consists of torsion
points and is defined over a finite extension of their field of
definition.  Taking the quotient produces the required abelian
variety over $\overline\Q$.  Homomorphisms and polarizations of
abelian varieties over an algebraically closed field of
characteristic zero are unchanged by extension to $\C$; thus the
given polarization descends as well.  Its symmetric line-bundle
representatives differ by two-torsion data, which are also
algebraic.  These operations preserve the specified analytic
polarized torus and merely supply its algebraic model.

\paragraph{The finite theta module and its one common comparison scalar.}
Choose rational symplectic coordinates for the actual alternating
form.  Replacing a lattice by a commensurable symplectically
self-dual lattice is an isogeny and allows principal coordinates
without changing the rational CM decomposition.  Record this
isogeny and its effect on the polarization; it is part of the
comparison data.  Write the resulting period matrix as $Z_i$ and
the torus as $\C^g/(\Z^g+Z_i\Z^g)$, where $g=3[F:\Q]$.

For a fixed finite family of finite Schwartz inputs with algebraic
values, including
all the finite translates used below, choose an even integer $N_i$
such that their supports are contained in $N_i^{-1}\Z^g$ and
their restrictions to rational points are invariant under
$N_i\Z^g$.  Enlarge $N_i$ to absorb the chosen isogenies.  They
belong to the finite Schr\"odinger space
\[
 W_i=\operatorname{Fun}
       (N_i^{-1}\Z^g/N_i\Z^g,\overline\Q).
\]
The theta sum followed by $z=N_iw$ gives the section
\[
 \mathcal T_i(\varphi)(w)=
 \exp\!\left(\frac\pi2N_i^2w^t
                       (\operatorname{Im}Z_i)^{-1}w\right)
 \sum_{q\in N_i^{-1}\Z^g}
    \varphi(q)\exp(\pi i q^tZ_iq+2\pi iN_iq^tw)
 \tag{B19b}
\]
of the $N_i^2$-th power $L_i^{N_i^2}$ of the symmetric theta
line.  In restriction-of-scalars coordinates the exponents are
the corresponding sums over embeddings.  Any scalar in the
chosen oscillator or Fourier normalization is included in the
definition of $\mathcal T_i$, not suppressed.

The dimension of $W_i$ is $N_i^{2g}$, equal to
$h^0(A_{x,i},L_i^{N_i^2})$.  The map $\mathcal T_i$ is injective:
its Fourier expansion recovers every value of $\varphi$ on the
finite quotient.  It is therefore an isomorphism after extension
to $\C$.  Translation and multiplication by finite characters
give its theta-group equivariance, with the central character
fixed by the displayed multiplier.  The completion factor is
precisely the one proved in A8.

Choose the matching algebraic theta structure and an equivariant
isomorphism over $\overline\Q$,
\[
 J_i:W_i\longrightarrow
       H^0(A_{x,i},L_i^{N_i^2}).
\]
The matching theta structure is algebraic: its torsion points are
algebraic, and lifts of finite order in the theta group satisfy
algebraic power equations.  A choice of compatible finite lifts
and of the central roots of unity therefore descends after a
finite extension.  It is essential to match this theta structure,
not merely the polarization type.

There is a single $c_i\in\C^\times$ such that
\begin{equation}\label{B19-Heis}
                 \mathcal T_i=c_i(J_i\otimes\C).
                 \tag{B19c}
\end{equation}
For completeness, the finite Heisenberg proof is elementary.
Multiplication characters diagonalize $W_i$ with one-dimensional
eigenspaces spanned by the delta functions.  Translation operators
permute these eigenspaces transitively.  A commuting endomorphism
is diagonal in the delta basis, and commutation with translations
forces all its diagonal entries to agree.  The same representation
is the theta-group representation on the section space.  This is
Mumford, \emph{On the equations defining abelian varieties I},
\S1, Proposition 3 and Theorem 2, printed pp.295--298.
Applying the commuting-endomorphism argument to
$\mathcal T_iJ_i^{-1}$ proves (B19c).

Thus $c_i$ depends on the actual torus, theta structure, oscillator
normalization, and chosen algebraic isomorphism $J_i$.  It does
not depend on $\varphi$.  It may be specified, without an
unmentioned scalar, by choosing a section with nonzero value at
the origin and taking the ratio of its two values in the fixed
rigidification.  Such a section exists because $N_i^2\geq4$.
This ratio is a theta-null period; no assertion that $c_i$ itself
is algebraic is made.

\paragraph{First jets and the CM tangent periods.}
Normalize the symmetric structure on $L_i^{N_i^2}$ to act as the
identity on its fibre at zero, and fix an algebraic rigidification.
A symmetric formal local generator with value one exists by
averaging a generator with its pullback by $[-1]$.  The first jet
of a section in this generator is independent of the choice:
the ratio of two such generators is an even unit with value one,
so its differential at zero vanishes.  First jets of algebraic
sections are therefore algebraic in algebraic tangent frames.
This argument concerns first jets.  A8 supplies the specified
canonical gauge for higher jets.

The Appell--Humbert frame in (B19b) has the same normalized
symmetry.  Its comparison with the symmetric algebraic frame has
value one and zero first derivative at the origin.  Consequently
(B19c) gives, for every input, the exact first-jet comparison.
Let $u^z_{ir}$, $r=1,2$, be the two directions in the original
$z$ variable used by the selected $c\mu$ Taylor tensor.  Let
$u^w_{ir}$ have the same coordinate vectors in the $w$ variable;
thus $d(z=N_iw)(u^w_{ir})=N_i u^z_{ir}$.  It is the
$w$ variable that uniformizes the chosen algebraic torus.  Choose
algebraic tangent vectors $w_{i1},w_{i2}$ through the two CM
factors in (B19a).  There exist $p_i\ne0$ and
$M_i\in\operatorname{GL}_2(\overline\Q)$ such that
\[
 u^w_{ir}=p_i\sum_{s=1}^2 (M_i)_{sr}w_{is}.
 \tag{B19d}
\]
Here $p_i$ is defined by this Betti--de Rham tangent comparison,
so the equation fixes whether a period or its inverse occurs.
To justify the common scalar, use the same algebraic frame on
the one-dimensional $\mu$ tangent eigenspace of $B(d_i)$ in
each of its two copies.  The two copies then have the same
comparison factor.  The rational $E$-linear isogeny and the
chosen rational coordinate lines change their frames by matrices
over $\overline\Q$, giving $M_i$.  This proves (B19d).

If $j_{is}(\varphi)$ denotes the algebraic first jet of
$J_i\varphi$ along $w_{is}$, the result is
\begin{equation}\label{B19-jet}
 D_{u^z_{ir}}\widehat\vartheta_{\varphi}(Z_i,0)
   =\frac{c_ip_i}{N_i}
        \sum_{s=1}^2(M_i)_{sr}j_{is}(\varphi).
 \tag{B19e}
\end{equation}
Indeed $D_{u^w_{ir}}\mathcal T_i(\varphi)(0)
=N_iD_{u^z_{ir}}\widehat\vartheta_{\varphi}(Z_i,0)$.
The factor $1/N_i$ therefore occurs once.  At first order the
quadratic completion contributes zero, and the divided Taylor
coefficient has $1/1!=1$.  If a different rational symplectic
lattice is used, its pullback matrix and its rational dilation
must be included in $M_i$ and $N_i$.

The passage from the two selected Taylor tensors to
$\Omega_x^1$ uses the chosen square-root determinant and
constant-line frames of 032, equation (4.6).  In a fixed rational
coordinate frame of $V$ its remaining scalar for the wedge will
be denoted by $\kappa_x\ne0$.  This scalar is part of the
specified metaplectic and canonical-fibre comparison.  All
algebraic coordinate matrices can be absorbed into $M_1,M_2$.
Define the common CM comparison factor
\[
 C_x=\kappa_x\frac{c_1c_2p_1p_2}{N_1N_2}.
 \tag{B19f}
\]
No Haar volume is included in $C_x$.

\paragraph{The infinity integral and the finite source class sums.}
Let $T_i=\U(P_i)$.  Its entire archimedean group is compact.
Fix a finite open subgroup $K_i$ fixing all source inputs, and
characters $\pi_i$ trivial on $K_i$.  The chosen first-degree
archimedean oscillator coefficient has one character
$\eta_{i,\infty}$ under $T_i(F_\infty)$.  Thus the bilinear
integral of the theta coefficient against $\pi_i$ vanishes unless
\[
                    \eta_{i,\infty}\pi_{i,\infty}=1.
 \tag{B19g}
\]
Indeed the Haar integral of a nontrivial character of a compact
group is zero.  Assume (B19g), as is necessary for a nonzero
degree-one lift in this fixed infinity type.  The integrand is
then invariant under the full archimedean group and under $K_i$.
Choose representatives $t_{ia}$ of the finite quotient
$T_i(F)\backslash T_i(\A_{F,f})/K_i$ and put
\[
 \Gamma_{ia}=T_i(F)\cap t_{ia}K_it_{ia}^{-1},\qquad
 v_i^T=\operatorname{vol}(T_i(F_\infty))\operatorname{vol}(K_i).
\]
The intersection defining $\Gamma_{ia}$ is understood with the
archimedean component unrestricted; it is a finite group of
integral norm-one units.  Quotient Haar integration gives exactly
\[
 \int_{[T_i]}\theta_i(g,t)\pi_i(t)\,dt
  =v_i^T\sum_a\frac{\pi_{i,f}(t_{ia})}{|\Gamma_{ia}|}
                      \theta_i(g,t_{ia}).
 \tag{B19h}
\]
This is a consequence of the verified archimedean character
cancellation, not a replacement of an arbitrary continuous
integral by a finite sum.  Stabilizers and both Haar factors are
retained.  For an unintegrated Taylor form set $v_i^T=1$ and
use its single prescribed input instead of (B19h).

For honest algebraic infinity characters the finite values in
(B19h) are algebraic.  Here is a direct proof.  The finite torus
class group has finite order.  For a representative $t$ some
$r>0$ satisfies $t^r=\gamma k$ with
$\gamma\in T_i(F)$ and $k\in K_i$.  Global triviality on
$T_i(F)$ gives
\[
 \pi_{i,f}(t)^r=\pi_{i,\infty}(\gamma_\infty)^{-1}
                      \in\overline\Q^\times.
\]
The last inclusion follows because an infinity character of the
honest norm-one torus is a product of integral powers of its
algebraic embedding coordinates.  Taking a root proves the claim.
Weil transforms of the rational finite inputs are likewise
$\overline\Q$-valued when the fixed finite splitting characters
are algebraic: translations and character multiplication use
algebraic values, and a finite Fourier transform uses roots of
unity and the specified algebraic self-dual volume factors.
Arbitrary scalar multiples of the input must be carried as
additional coefficients, rather than declared algebraic.

\paragraph{The definite Bessel coefficient is the promised finite functional.}
Let $H=\U(L^\perp)$ and choose a common finite level $K_H$.
Let $e$ be a function in its scalar archimedean type, the same
type as the canonical coefficient $f$ of the wedge.  Choose
representatives $h_b$ of
$H(F)\backslash H(\A_{F,f})/K_H$ and set
\[
 \Gamma_b^H=H(F)\cap h_bK_Hh_b^{-1},\qquad
 v_H=\operatorname{vol}(H(F_\infty))\operatorname{vol}(K_H).
\]
Use these fixed Haar measures in B17 as well.  The product
$f\overline e$ is invariant under $H(F_\infty)K_H$, so
\[
 \mathcal B_e(f)=\int_{[H]}f(h)\overline{e(h)}\,dh
       =v_H\sum_b\frac{f(h_b)\overline{e(h_b)}}{|\Gamma_b^H|}.
 \tag{B19i}
\]
At $h_b$ the finite oscillator input in the $a$-summand is the
actual transform $\varphi_{iab}=\omega_{i,f}(h_b,t_{ia})
\varphi_i$.  The common level $N_i$ above was chosen to include
all these finitely many transforms.  Define the algebraic arrays
\[
 a_{ibr}=\sum_a\frac{\pi_{i,f}(t_{ia})}{|\Gamma_{ia}|}
             \sum_{s=1}^2(M_i)_{sr}j_{is}(\varphi_{iab}).
\]
Equations (B19e)--(B19h), and alternation of the two one-forms,
give the exact identity
\begin{equation}\label{B19-period}
 \mathcal B_e(f)
   =v_Hv_1^Tv_2^T C_x
       \sum_b\frac{\overline{e(h_b)}}{|\Gamma_b^H|}
          (a_{1b1}a_{2b2}-a_{1b2}a_{2b1}).
 \tag{B19j}
\end{equation}
The sign is the order of the chosen coordinate volume form and
of the two input one-forms.  No factorial is missing: this is the
alternation of two first Taylor tensors.

Choose $e$ from the algebraic finite definite module, meaning that
its values in the specified algebraic scalar-type frame belong
to $\overline\Q$.  Then the sum in (B19j) is algebraic.  For
a fixed finite family of inputs, characters and test vectors,
it lies in one number field.  The same $C_x$, $v_H$ and $v_i^T$
work throughout that family, provided the types, theta structures,
levels, splittings, and frames remain fixed.  This proves the
absence of an extra transcendental scalar depending on the
source character in this finite CM-evaluation comparison.

An orthonormal definite basis need not have the same algebraic
normalization.  Its Gram matrix in an algebraic basis is $v_H$
times a matrix with algebraic entries, because the stabilizer
weights are rational.  Gram--Schmidt therefore gives basis values
in $v_H^{-1/2}\overline\Q$.  For that normalization the common
measure factor in (B19j) is $v_H^{1/2}v_1^Tv_2^T C_x$.
This records the change instead of silently imposing Haar
probability or identifying an algebraic vector with a unit vector.

There is also an explicit lower certificate once arithmetic data
for the displayed sum are supplied.  Let its value be $b\ne0$
in a number field $K\subset\C$ of degree $d$ whose chosen
embedding is nonreal (adjoin $i$ if necessary), let $Db$ be integral
for a positive integer $D$, and assume
$|\sigma(b)|\leq M$ with $M>0$ for the $d-2$ embeddings other than the
chosen complex embedding and its conjugate.  Then
\[
 |b|^2\geq D^{-d}M^{-(d-2)},\qquad
 |\mathcal B_e(f)|^2\geq
 |v_Hv_1^Tv_2^TC_x|^2D^{-d}M^{-(d-2)}.
 \tag{B19k}
\]
Indeed $|N_{K/\Q}(Db)|\geq1$, and the two distinguished
factors in the norm product are $D^2|b|^2$.
For a unit test vector replace the Haar factor as just explained.
This certificate does not supply a bound on the CM period factor
or the conjugates without additional data.

\paragraph{What remains to match the actual archive arithmetic functional.}
The recovered source makes this remaining requirement explicit.
Under the directory
\path{Development/AbelianVarietyTheory/Definitions/LatticeFrames/},
\texttt{LevChiEqOne.lean:73--77} defines
\begin{verbatim}
def abelian₂ ... :=
  ((I.chIn u a).rho i c * (I.chIn u a).rho' i c') •
    jetL I.J.cD
      (I.th u a i _ _ ((I.stp u a i).rep c)
                         ((I.stp u a (i + 1)).rep c'))
\end{verbatim}
The omitted arguments are only the typed indices and the declared
domain; the scalar and section in this quotation are the actual
definition.  Thus this functional is tied to a specified
arithmetic theta section and its period frame.
\texttt{OnFramesRdApply.lean:108--113} requires its equality
with the analytic theta kernel after multiplication by
\texttt{volLev}, the prescribed kernel constant \texttt{t.cK},
and the polynomial rescaling \texttt{rho.sc}.

The proved constructor is
\path{Development/AbelianVarietyTheory/Proofs/DecoratedLatticeProofs/KernelPartLawsOfColumnPins.lean},
theorem \texttt{clsKL\_is}, lines 22--52.  Its core step at
lines 29--35 invokes
\nolinkurl{thetaKernelIdentity_at_every_tuple} with
\texttt{PinX} and \texttt{JetS}.
The underlying proved reading is
\nolinkurl{ThetaKernelSplitPlacesEqJetClassicalTheta.lean},
theorem \nolinkurl{thetaKernel_eq_jet_thC}, lines 46--71.
Its \texttt{StdPair} hypothesis is defined in
\texttt{SplitPart.lean:38--49}; it includes
\begin{verbatim}
gS : S.gram = 1
gT : T.gram = 1
rk : 0 < S.rank ∧ S.rank < T.rank ∧ Even S.rank ∧ Even T.rank
cc : D.host.ccK = ⊤
\end{verbatim}
The explicit jet scale is exhibited in
\texttt{PureTensorOZeroEqJetP.lean:342--346} as
$s=c_K(-\delta)$, where lines 267--269 define
$c_K=(\sqrt{2\pi/\sqrt d})^{-1}$.
These are static source readings at the pinned mirror commit
recorded in B0, not a build or an audit of their full dependencies.

Consequently the arithmetic match for (B19j) is now a concrete
comparison of the polarized CM lattices and isogenies, finite
theta structures and sections, the two tangent period frames,
the determinant-line frame, and the displayed covolume/Haar
factors with the corresponding \texttt{I.th}, \texttt{I.J.cD},
\texttt{rho}, \texttt{c0}, and \texttt{t.cK}.  It is sufficient
to check equality on a basis of the finite theta modules and the
relevant first-jet directions; linearity then gives the entire
fixed family and (B19j).  The recovered definite Gram-one,
strictly increasing even-rank instance does not itself make that
comparison for the rank-two/rank-three mixed-signature data.
Defining an abstract \texttt{abelian₂} to equal the analytic side
would satisfy a formal equality but would not establish this
arithmetic comparison.

In particular (B19j) gives an actual definite Bessel functional
of the CM jets associated with the original 032 inputs, and B17
identifies the same analytic Bessel coefficient with the last
theta seesaw.  The remaining scalar comparison is common to
the fixed family, with every measure and frame identified; it is
not a collection of unspecified constants indexed by characters.
\hfill$\square$


% THETA3-B20 contribution by /root/definite_bridge, 21 November 2026.
\subsection{THETA3-B20: isotropic column completion preserves a first determinant value}\label{B20}

For the exponent-one determinant occurring in the original
canonical wedge, the extra section columns required by
\texttt{IsCols} can be supplied without changing the value of
the differential at a fixed CM point.  The reference theta-null
values and the integral formal CM action must be retained.
This is a finite-jet comparison; it does not identify the
resulting theta sections as global automorphic functions.

\paragraph{The elementary determinant identity.}
Let $G$ be a commutative formal group over a characteristic-zero
field, with commuting invariant derivatives $D_1,\ldots,D_i$.
For a section germ use a fixed normalized canonical frame, and
write its scalar coefficient as $F$.
Let $F_1,\ldots,F_i,R_1,\ldots,R_i$ be scalar formal series with
\[
                  R_c(0)=r_c,\qquad D_aR_c(0)=0
                     \quad(1\leq a,c\leq i).
\]
On $2i$ copies write $D_{a,c}$ for $D_a$ on the $c$-th copy.
For any scalar $\beta$ define
\[
 \mathcal D_\beta=
     \det(D_{a,c}+\beta D_{a,c+i})_{1\leq a,c\leq i},
 \qquad
 \mathcal D_0=\det(D_{a,c})_{1\leq a,c\leq i}.
\]
Then
\begin{equation}\label{B20-value}
 \mathcal D_\beta
       \left(\prod_{c=1}^iF_c(t_c)R_c(t_{c+i})\right)\Big|_0
     =\left(\prod_{c=1}^ir_c\right)
          \mathcal D_0\left(\prod_{c=1}^iF_c(t_c)\right)\Big|_0.
 \tag{B20a}
\end{equation}
Expand the determinant by permutations and each of its $i$
factors into its two summands.  Every term selecting an extra
column differentiates its particular $R_c$ exactly once;
that term vanishes at zero.  The only surviving choice selects
all the original columns and leaves the reference factors
undifferentiated.  This proves (B20a), with its sign and without
a factorial or a conjugation of $\beta$.

\paragraph{Actual integral even references.}
Let $R$ be a complete discrete valuation ring of residue
characteristic $p\ne2$, and let $A/R$ be an abelian scheme
with a symmetric relatively ample line bundle $L$, rigidified
at zero.  Normalize the symmetry on $L^k$ to be the identity
on its fibre at zero.  Suppose the evaluation map
\[
                  H^0(A,L^k)\longrightarrow (L^k)_0=R
 \tag{B20b}
\]
is surjective.  This is an explicit condition on the chosen
theta level.  It holds for sufficiently large $k$, by relative
generation and cohomology-and-base-change for ample powers.
It can be imposed by choosing a sufficiently large square
level prime to $p$; any pullback isogeny and its derivative
factor must be recorded when an already fixed input is moved
to that level.

Choose an integral section $s$ with $s(0)=1$ and put
\[
                        r=\frac{s+[-1]^*s}{2},
\]
using the normalized symmetry to regard both summands as
sections of the same line bundle.  Then $r(0)=1$ and its
first derivative in any symmetric origin frame is zero,
because $d[-1]=-1$.  The canonical potential of A8--A9 has
zero constant and linear terms: it is obtained from a
biextension series vanishing on both axes by restriction
to $(x,[w_0]x)$.  The actual canonical-frame first jets of
$r$ therefore vanish as well.  Thus $r$ supplies the reference
germ in (B20a), with $r_c=1$ and with an integral section,
without assuming the integrality of its entire canonical
theta-frame expansion.

If the chosen reference has a unit value $r_c$ rather than one,
division by that unit gives the normalized reference.  If its
value is merely nonzero in the fraction field, the value can
still be divided out in characteristic zero, but that operation
has a valuation cost.  An assertion of integral divisibility
must then retain that cost.  Merely choosing a nonvanishing
complex theta-null value does not prove (B20b).

\paragraph{The actual completed CM columns.}
Assume the formal group has an integral action of $\OO_E$,
all its endomorphisms are defined over $R$, and its coefficientwise
continuous extension to $\OO_E\otimes\Z_p$ is available.
This extension can also be constructed directly: use the formal
logarithm over the fraction field, write an endomorphism as
$\operatorname{Exp}\circ d[a]\circ\log$, and take the
coefficientwise limit of integral global CM actions.  Each finite
jet is a polynomial over the fraction field in the entries of
$d[a]$, and closedness of $R$ preserves integrality.

Fix one CM embedding direction $\tau$ in which the tangent
frame $u_1,\ldots,u_i$ is integral and saturated.  Choose
\[
 b\in(\OO_E\otimes\Z_p)^\times,
                   \qquad b\overline b=-1,
                   \qquad \beta=d[b]|_\tau\in R^\times.
 \tag{B20c}
\]
The bar in (B20c) is the CM involution on the completed algebra;
it is not an assertion about the ordinary complex conjugate
of an arbitrarily chosen complex image of $\beta$.

Such a $b$ exists when $p$ is odd and unramified in $E$.
At a split pair in $E/F$, take the two unit components
$(1,-1)$.  At an unramified quadratic local extension with
residue-field sizes $q^2/q$, the Teichm\"uller units contain an
element with $b^{q+1}=-1$: $2(q+1)$ divides $q^2-1$ because
$q$ is odd.  Its local CM norm is $-1$.
Make these choices at every factor of $E\otimes\Q_p$.
The embedding values required in (B20c) are roots of unity or
$\pm1$, so they can be included in the coefficient field after
a finite extension.

The matrix of section columns is
\[
                         B=[I_i\mid\beta I_i].
 \tag{B20d}
\]
It has an $i\times i$ unit minor.  For every $N$, choose
$a_N\in\OO_E$ congruent to $b$ modulo
$p^N(\OO_E\otimes\Z_p)$.  Then
\[
 X_N=[I_i\mid a_NI_i],\qquad
 X_N\overline{X_N}^{\,t}
                  =(1+a_N\overline{a_N})I_i
                        \equiv0\pmod{p^N\OO_E}.
\]
At the selected embedding $X_N$ converges to $B$.
Thus (B20d) satisfies the literal supported-block, unit-minor
and CM-isotropic-approximation requirements of
\texttt{IsIsoCols}; its single block has precisely $2i$ columns.
For several determinant factors the construction uses disjoint
blocks, as required by \texttt{IsCols}.  The precise definitions
are
\path{Development/AbelianVarietyTheory/Definitions/Divisors/NonemptyValExtOfOverLa.lean:331--342},
read in A9.  The differential polynomial produced by this
matrix is exactly $\mathcal D_\beta$, since the selected
embedding acts on every $u_a$ by the same scalar $\beta$.

Applying A9's integral-unit lemma to these columns makes the
entire canonical theta-frame factor an integral unit after
the corresponding formal substitution.  Applying (B20a) to
the actual section germs then identifies the exponent-one
value with the original determinant value, including the
reference theta-null factors.  Any use of a particular
\texttt{nuFac} lower bound still requires its stated local
type, saturation and valuation-normalization hypotheses.

\paragraph{The two different oscillator tori in B19.}
The construction can be made on one common polarized abelian
variety without identifying $B(d_1)$ with $B(d_2)$.
Let $A_1,A_2$ be the two actual CM theta tori at the fixed
rational negative line, with line bundles $M_1,M_2$ at the
chosen common input levels.  Choose even reference sections
$r_i$ of $M_i$ with $r_i(0)=1$ under (B20b), and put
\[
 A=A_1\times A_2,\qquad M=M_1\boxtimes M_2,
 \qquad S_1=s_1\boxtimes r_2,
 \quad S_2=r_1\boxtimes s_2,
 \quad R_0=r_1\boxtimes r_2.
 \tag{B20e}
\]
All three are actual sections of the same ample line bundle
$M$; $R_0(0)=1$ and its first canonical jets vanish.
For the tangent rows take
\[
             u_a=(u_{1a},u_{2a})\in
                          \operatorname{Lie}(A_1)\oplus
                          \operatorname{Lie}(A_2),\quad a=1,2.
\]
The $u_{ia}$ are the normalized algebraic $\mu$ directions
from B19, with its period factors kept outside the algebraic
jet value.  The diagonal $E$-action has the same eigenvalue
on both summands in this selected embedding direction.
Under the stated integral saturated-frame hypothesis they
therefore admit the columns (B20d).

The product rule and $dr_i(0)=0$ give
\[
       D_{u_a}S_1(0)=D_{u_{1a}}s_1(0),\qquad
       D_{u_a}S_2(0)=D_{u_{2a}}s_2(0).
\]
Consequently the original two-torus first-jet determinant is
the two-column determinant for $S_1,S_2$ on $A^2$.
By (B20a) it is also the four-column isotropic determinant
for $S_1,S_2,R_0,R_0$ on $A^4$.  Nonunit changes of tangent
frame or of level must be kept as explicit factors; the
saturation condition is not supplied by an arbitrary
complex choice of frame.

This works simultaneously for the finite list of transformed
sections in (B19j): they lie in the same finite section spaces
on the fixed $A_i$, and the references can be chosen once.
It therefore completes the exponent-one \emph{value} comparison
for that finite CM-jet functional.  It does not assert that
$A$ has the particular uniform CM Lie multiplicities required
by every recovered \texttt{Datum}, or supply a missing
integral model, tame-generator normalization, or global
\texttt{MainX} witness.

\paragraph{Why the conclusion stops at exponent one.}
For $\mathcal D_\beta^{a}$ with $a>1$, a reference factor
may receive two or more derivatives, and evenness no longer
kills all the added terms.  A sufficient replacement is that
each reference have constant canonical jet through order $a$:
then every extra derivative term still vanishes.  Jet separation
over $\overline\Q$ alone does not give this with integral unit
normalization.  The inverse canonical gauge can have
nonintegral higher coefficients, as C22 demonstrates.
No such high-order integral references are asserted here.

Even at exponent one, the completion is not an equality of
the original automorphic functions on the ball.  Away from
the fixed CM point the extra even theta sections contribute
their varying theta-null factors and additional line-bundle
weights.  Their chosen value at the point is all that was
used in (B20a).  Thus B20 removes a specific finite-jet column
obstruction while leaving the global section and high-exponent
comparisons explicit.\hfill$\square$


% THETA3-B21 contribution by /root/definite_bridge, 21 November 2026.
\subsection{THETA3-B21: the cross-type scalar is a specific determinant period}\label{B21}

B19 gives a common scalar when the complete CM types are fixed.
For different complete types, the remaining comparison can be
specified more precisely.  Balancing the four original CM types
controls the selected tangent periods, but need not control the
ordinary determinant of the full Hodge space of each theta torus.
This distinction does not disprove an equality between particular
period numbers.  It identifies the additional comparison that
the proposed CM-Hodge argument must actually establish.

\paragraph{An unambiguous modular normalization.}
Let $A_i/\overline\Q$ be the actual principally polarized CM
torus of B19, after its recorded isogeny, and set
$g=3[F:\Q]$.  Fix an algebraic basis of regular differentials
and normalized analytic coordinates by
\[
 A_i(\C)=\C^g/(\Z^g+Z_i\Z^g),\qquad
 \boldsymbol\omega_i=\Omega_i\,d\boldsymbol w,
 \qquad D_i=\det\Omega_i\ne0.
 \tag{B21a}
\]
Thus $\Omega_i$ is a period block of algebraic differentials,
not the normalized Siegel matrix $Z_i$.
Write $a\sim_{\overline\Q}b$ for
$a/b\in\overline\Q^\times$.
Use the geometric modular-form convention
\[
 f(A,\omega)=\frac{f(A)}{\omega^{\otimes k}},\qquad
 f(A,\lambda\omega)=\lambda^{-k}f(A,\omega).
\]
The analytic form $f_{\rm an}$ corresponding to the geometric
weight-$k$ form satisfies
\[
 f(A_i,\omega_{i1}\wedge\cdots\wedge\omega_{ig})
       =\frac{(2\pi i)^{gk}}{D_i^k}f_{\rm an}(Z_i).
 \tag{B21b}
\]
This is Lachaud--Ritzenthaler--Zykin,
\emph{Jacobians among abelian threefolds},
\S2.1, equation (1), Proposition 2.2.1 and Proposition 2.4.4
with its proof, printed pp.3--6, arXiv:0802.4017v1.
Their normalized block is called $\Omega_2$.
Although their application is in genus three and their displayed
spaces have full level, the proof of (B21b) is a change of
coordinates on a polarized complex torus and is unchanged after
fixing a finite level structure.

Here is the precise theta input.  Kirschmer--Narbonne--Ritzenthaler--Robert,
\emph{Spanning the isogeny class of a power of an elliptic curve},
arXiv:2004.08315v2, \S4.3, equation (7), Definition 4.3 and
Remark 4.4, printed pp.23--24, prove that pair products of the
finite-level theta-null coordinates are algebraic weight-one
Hodge sections.  They obtain this from the theta transformation
law and the algebraic $q$-expansion principle.  A basis of
algebraic differentials therefore fixes their modular affine
lift up to a common sign.  We use the scalar convention
(B21b) to fix the direction of the period comparison.

The theta coordinates in that source and the coordinates of B19
differ by a finite Fourier transform with algebraic entries.
Indeed, putting $n=N_i^2$ and using $a,j\in(\Z/n\Z)^g$,
\[
 \sum_j e^{-2\pi i a^tj/n}
     \vartheta\left[\begin{smallmatrix}0\\j/n\end{smallmatrix}\right]
                       (w,Z_i/n)
   =n^g\vartheta\left[\begin{smallmatrix}a/n\\0\end{smallmatrix}\right]
                       (nw,nZ_i).
 \tag{B21c}
\]
To prove this identity, expand the left side over $m\in\Z^g$.
The inner character sum is $n^g$ when $m\equiv a\pmod n$
and zero otherwise.  Writing $m=n\ell+a$ gives the right side.
At $w=0$ the completion factor in (B19b) is one.
Thus the pair products of B19's theta-null values have the
same algebraic modular interpretation.  Enlarging the finite
level to kill a multiplier, or taking a common even power,
would give the same conclusion below after taking an algebraic
root; no claim that the eighth power suffices for every chosen
rational characteristic is needed.

This interpretation is on the actual $A_i$, with its recorded
theta level.  In the coordinates of (B21c), the torus with
matrix $Z_i/n$ is the quotient of $A_i$ by the indicated
finite subgroup of its $n$-torsion; the map is induced by
the identity on $\C^g$.  Its algebraic model and this isogeny
are defined over $\overline\Q$, since the subgroup consists
of algebraic torsion points.  Pullback identifies their Hodge
spaces over $\overline\Q$.  Thus one may either use the
level-$n$ Hodge section on $A_i$ directly, as in the cited
theta-moduli construction, or pull it back through this
isogeny.  The latter description changes an algebraic
differential basis by a matrix in
$\operatorname{GL}_g(\overline\Q)$, whose determinant is
the recorded algebraic isogeny factor.  It introduces no
additional transcendental scalar and does not replace $D_i$
by an unspecified period of another abelian variety.

Let $\upsilon_i$ denote any scalar by which the specified
oscillator normalization multiplies the displayed standard
theta sum of (B19b); for exactly that displayed sum
$\upsilon_i=1$.  Define $\widetilde c_i=c_i/\upsilon_i$.
Then
\[
             \widetilde c_i^{\,2}
                    \sim_{\overline\Q}
                         \frac{D_i}{(2\pi i)^g}.
 \tag{B21d}
\]
To prove this, select an algebraic finite input $\varphi$ for
which $J_i\varphi(0)\ne0$.  B19 gives such an input and
identifies $\vartheta_\varphi(0)
=\widetilde c_i J_i\varphi(0)$ in the fixed rigidification.
The square of this theta-null is the value of an algebraic
weight-one modular form in analytic coordinates.  By (B21b),
\[
 \frac{(2\pi i)^g}{D_i}
       \vartheta_\varphi(Z_i,0)^2\in\overline\Q^\times.
\]
Dividing by the nonzero algebraic number
$J_i\varphi(0)^2$ proves (B21d).
If a higher even power was used, the assertion follows because
$\overline\Q$ is algebraically closed.  This proof retains
both the period orientation and the universal power of $2\pi i$.

\paragraph{The exact residual factor for two balanced pairs.}
For pairs $(1,2)$ and $(3,4)$, write B19's constants as
\[
 C_{x,ij}=\kappa_{x,ij}
                    \frac{c_ic_jp_ip_j}{N_iN_j}.
\]
Define
\[
 \mathcal R=
 \frac{C_{x,12}}{C_{x,34}}
 \frac{\kappa_{x,34}}{\kappa_{x,12}}
 \frac{\upsilon_3\upsilon_4}{\upsilon_1\upsilon_2}
 \frac{p_3p_4}{p_1p_2}\frac{N_1N_2}{N_3N_4}.
\]
Then $\mathcal R=\widetilde c_1\widetilde c_2/
(\widetilde c_3\widetilde c_4)$ exactly, and
\[
              \mathcal R^2\sim_{\overline\Q}
                        \frac{D_1D_2}{D_3D_4}.
 \tag{B21e}
\]
The universal powers cancel as
$(2\pi i)^{-g-g+g+g}=1$ because all four theta tori have
the same dimension.  The level factors have been displayed
explicitly; they are nonzero rational numbers.

Suppose $d_1+d_2=d_3+d_4$ as complete type functions,
with the indexing $d_i(\tau)=v_i(\tau^{-1})$ of B19.
The selected tangent eigenspace has a Hodge type which is
linear in $d_i$: at a conjugate embedding it is determined
by $(1+d_i(\tau))/2$, or by its simultaneous dual convention.
Hence the two selected tangent tensor products are isomorphic
as coefficient-field Hodge structures.  The accepted CM Hodge
conjecture makes this tensor comparison algebraic: after the
usual K\"unneth projectors it is a Hodge class on a product of
CM abelian varieties, represented by an algebraic cycle.
Its Betti and de Rham matrices are algebraic, so compatibility
of the comparison isomorphism gives
\[
                         \frac{p_1p_2}{p_3p_4}
                                      \in\overline\Q^\times.
 \tag{B21f}
\]
The factors $p_i$ here have the exact direction fixed by
(B19d); taking all four duals changes neither this conclusion
nor the balancing condition.

There is no hidden inverse $D_i$ in the Betti constant-line
calculation of 032, (4.6).  At the rational point $x=[Ez]$,
the line $L$, its annihilator, and every full compact embedding
summand in (4.5) are defined over $\overline\Q$ in the
complexification of the rational symplectic space.  A rational
symplectic change of coordinates is algebraic, and the graph
matrix $Z_i$ and its graph basis have algebraic entries.
The projection defining this graph is invertible by positivity.
The isomorphism
$\det Q^*=L\otimes\det(V_\mu)^*$ and the tensor determinant
identities use algebraic matrices.  Algebraic square roots can
be chosen in $\overline\Q$.  The ball coordinates obtained
from $E(z+z_1e_1+z_2e_2)$ give an algebraic Betti tangent frame
as well.  Thus choosing these algebraic constant-line frames
makes $\kappa_{x,12}/\kappa_{x,34}$ algebraic after the same
rank- and Taylor-order-dependent oscillator constants are fixed.
The original 032, p.15 immediately after (4.6), permits arbitrary
nonzero complex constant-line vectors and explicitly asserts no
rational trivialization; any such additional rescalings must
instead remain in the displayed $\kappa$ ratio.

Consequently, under the stated algebraic Betti frames, common
oscillator convention, and accepted CM-Hodge comparison,
\[
 \left(\frac{C_{x,12}}{C_{x,34}}\right)^2
       \sim_{\overline\Q}\frac{D_1D_2}{D_3D_4},
 \qquad
 \frac{C_{x,12}}{C_{x,34}}\in\overline\Q^\times
 \Longleftrightarrow
 \frac{D_1D_2}{D_3D_4}\in\overline\Q^\times.
 \tag{B21g}
\]
This is a reduction to one specified period ratio, not an
assertion that the ratio is algebraic.

\paragraph{Why complete-type balance does not settle that ratio.}
Assume for this calculation that $E/\Q$ is Galois, as in the
type-function setting of 032.  Put $h_d=(1+d)/2$ on
$\operatorname{Gal}(E/\Q)$ and let
\[
 A_d(\gamma)=\sum_\tau h_d(\tau)h_d(\gamma\tau).
 \tag{B21h}
\]
The ordinary top Hodge determinant of $B(d)$ is formed by
wedging the eigenlines indexed by $h_d(\tau)=1$.  Its
coefficient-field Betti determinant line has, at the conjugate
embedding $\gamma$, holomorphic degree $A_d(\gamma)$:
the sum simply counts which of the selected eigenlines remain
of type $(1,0)$.  By the actual decomposition (B19a), the
corresponding degree function for $A_{x,i}$ is
\[
                         M_d=2A_d+A_{d^{\rm flip}}.
 \tag{B21i}
\]
This is a quadratic function of the original type, unlike the
selected tangent type in (B21f).

The determinant $D_i$, up to an algebraic factor, is the
period of this selected determinant tensor.  Indeed choose
algebraic CM eigen-differentials.  Each row of their period
matrix is its one CM comparison period times a row with
algebraic entries in a rational Betti basis.  Taking a
nonzero period-block determinant gives the product of those
comparison periods times a nonzero algebraic determinant.
Algebraic changes of differential basis and the recorded
isogenies preserve this statement.

For a direct counterexample to the proposed inference, take
$E=\Q(\zeta_7)$ and order its cyclic Galois group as
$1,\gamma,\ldots,\gamma^5$, with complex conjugation
$\gamma^3$.  Let $\mu=1$ and define the four complete types
\[
 d_{a,b}=(1,a,b,-1,-a,-b),\qquad a,b\in\{1,-1\}.
\]
They satisfy $d_{+,+}+d_{-,-}=d_{+,-}+d_{-,+}$, and all
have the same positive value at $\mu$.
Flipping $\mu,c\mu$ changes entries zero and three.
Direct counting in (B21h) gives
\[
\begin{array}{c|rrrrrr}
 d&M_d(1)&M_d(\gamma)&M_d(\gamma^2)&M_d(\gamma^3)
       &M_d(\gamma^4)&M_d(\gamma^5)\\\hline
 d_{+,+}&9&6&3&0&3&6\\
 d_{-,-}&9&6&3&0&3&6\\
 d_{+,-}&9&4&5&0&5&4\\
 d_{-,+}&9&2&7&0&7&2
\end{array}
\]
For example $A_{+,+}=(3,2,1,0,1,2)$, whereas
$A_{-,+}=(3,0,3,0,3,0)$; flipping interchanges the relevant
entries as prescribed.  Subtraction yields
\[
 M_{+,+}+M_{-,-}-M_{+,-}-M_{-,+}
                         =(0,6,-6,0,-6,6)\ne0.
 \tag{B21j}
\]
Thus the two specified determinant tensors need not be
isomorphic as coefficient-field Hodge structures, even though
the original complete types balance.  Applying the CM Hodge
conjecture to a nonexistent Hodge isomorphism cannot prove
their period ratio algebraic.

If the stronger equality $M_{d_1}+M_{d_2}=M_{d_3}+M_{d_4}$
does hold, the same algebraic-cycle argument used for the
tangent tensors proves the determinant period ratio algebraic.
Without it, (B21j) invalidates that particular deduction; it
does not prove the period ratio transcendental, exclude other
period relations, or refute the original mixed-pairing theorem.
The exact finite-jet functional and the common normalization
for fixed complete types in B19 remain valid.\hfill$\square$


% THETA3-B22 contribution by /root/definite_bridge, 21 November 2026.
\subsection{THETA3-B22: explicit theta-null and tangent normalization bounds}\label{B22}

The scalar in B19 need not be left with an unknown absolute
CM-period size.  An explicitly chosen even reference and an
explicitly nonzero odd first jet give an algebraic normalization
whose analytic factor has a direct lower certificate.  This
does not bound the heights or denominators of the resulting
algebraic coordinates, and does not assert that they are an
integral saturated frame at a specified prime.

\paragraph{The two references and their Gaussian tails.}
Fix $Z=X+iY$ in the Siegel upper half space, a dimension $g$,
and a nonzero specified tangent vector $u\in\C^g$.
Fix a certified number $\lambda>0$ with $Y\geq\lambda I_g$.
At a rational CM point the matrix entries are algebraic, so
one can choose a positive rational $\lambda$ and certify this
inequality by the positive principal minors of $Y-\lambda I_g$.
Choose an index $j$ with $u_j\ne0$ and put $q=e_j$.
Let $N>2$ be an even multiple of a level containing all the
previously prescribed finite inputs.  The two functions
\[
 \varphi_{\rm ev}=1_{N\Z^g},\qquad
 \varphi_{\rm odd}=1_{q+N\Z^g}-1_{-q+N\Z^g}
 \tag{B22a}
\]
belong to $W_N=\operatorname{Fun}(N^{-1}\Z^g/N\Z^g,\Q)$.
Their two odd cosets are distinct.  Use the standard theta
sum and its completed version from (B19b).  At zero the
completion is one, and its first derivative is zero.

For the even theta-null $E_N$ one has
\[
 E_N=\sum_{m\in\Z^g}e^{\pi iN^2m^tZm},\qquad
 |E_N-1|\leq
       \left(\frac{1+r}{1-r}\right)^g-1,
 \quad r=e^{-\pi\lambda N^2}.
 \tag{B22b}
\]
Indeed $m^2\geq |m|$ for an integer, so
$\sum_{m\in\Z}r^{m^2}\leq(1+r)/(1-r)$; taking the
product over the coordinates and deleting zero proves the
bound.  We may increase $N$ until its right side is at most
$1/2$, obtaining
\[
                         \frac12\leq |E_N|\leq\frac32.
 \tag{B22c}
\]

The odd first derivative has the exact expansion
\[
 H_N=D_u\widehat\vartheta_{\varphi_{\rm odd}}(Z,0)
       =4\pi i\sum_{m\in\Z^g}
          ((q+Nm)^tu)e^{\pi i(q+Nm)^tZ(q+Nm)}.
 \tag{B22d}
\]
To see the factor four and the sign, differentiate both
theta series and replace $m$ by $-m$ in the second.  The
linear factor changes sign and the quadratic factor does
not, so the two differentiated series add.
The zero term is
\[
 a_0=4\pi i(q^tu)e^{\pi iq^tZq},\qquad
 |a_0|=4\pi|q^tu|e^{-\pi q^tYq}>0.
\]
Put $a=\pi\lambda N^2/4$ and $t=e^{-a}$.
Since $N\geq2\|q\|$ and $\|m\|\geq1$ for $m\ne0$,
\[
 \|q+Nm\|\geq\frac N2\|m\|,\qquad
 \|q+Nm\|\leq\frac{3N}{2}\|m\|.
\]
Therefore
\[
 |H_N-a_0|\leq
  12\pi gN\|u\|\,
       \frac{t}{(1-t)^2}
       \left(\frac{1+t}{1-t}\right)^{g-1}
       =:T_N.
 \tag{B22e}
\]
For completeness, the initial bound is
$6\pi N\|u\|\sum_{m\ne0}\|m\|e^{-a\|m\|^2}$.
Use $\|m\|\leq\sum_k|m_k|$, then separate the sums.
In the distinguished coordinate,
\[
 \sum_{b\in\Z}|b|t^{b^2}
       \leq2\sum_{b\geq1}bt^b
       =\frac{2t}{(1-t)^2};
\]
each of the other coordinates contributes at most
$(1+t)/(1-t)$.  Summing the $g$ possibilities gives (B22e).

Increase $N$ through even multiples of the prescribed level
until, in addition to (B22c),
\[
                 T_N\leq
                 2\pi|q^tu|e^{-\pi q^tYq}.
 \tag{B22f}
\]
This process terminates: the right side is a fixed positive
number, whereas $T_N$ is a polynomial factor in $N$ times a
Gaussian tending to zero.  It is a finite certification
procedure with explicit input data.  The result is
\[
 |H_N|\geq B(Z,u,q):=
               2\pi|q^tu|e^{-\pi q^tYq}>0.
 \tag{B22g}
\]
No theta nonvanishing theorem or CM-period estimate is used
in this calculation.

\paragraph{The corresponding algebraic theta and tangent frames.}
Now use the actual CM torus of B19 and its algebraic finite
theta isomorphism $J_N$.
First normalize it algebraically so that
\[
                         J_N\varphi_{\rm ev}(0)=1.
 \tag{B22h}
\]
This is possible because its value is algebraic and nonzero:
the analytic value is $E_N\ne0$, and the two theta maps differ
by a nonzero scalar.  Dividing $J_N$ by its previous nonzero
value is an allowed algebraic change of frame.  With the
standard oscillator normalization, the new comparison scalar
is exactly $c_N=E_N$.  For a prescribed overall oscillator
scalar $\upsilon$, it is $c_N=\upsilon E_N$ instead; that
known scalar is retained.

Fix the rational $E$-linear CM decomposition from (B19a).
Choose one nonzero Betti eigenline direction in its first
$B(d)$ copy and let $u=u^w$ be its image in the $w$
coordinate of this actual torus.  The same coordinate vector
in the original theta variable is denoted $u^z$.
The selected complex tangent line is the complexification of
an algebraic CM eigenspace.  Its first-jet functional
\[
       \ell(v)=D_v(J_N\varphi_{\rm odd})(0)
\]
is algebraic on that eigenspace, because the section, line,
and symmetric origin first-jet frame are algebraic.
The exact substitution $z=Nw$ gives
\[
                          c_N\ell(u^w)=N\upsilon H_N.
 \tag{B22i}
\]
The right side is nonzero by (B22g).  Choose an algebraic
generator $v$ of the line with $\ell(v)=1$, by dividing any
algebraic generator by its nonzero algebraic $\ell$-value.
Writing $u^w=p_Nv$ now gives the exact identity
\[
                    \frac{c_Np_N}{N}=\upsilon H_N.
 \tag{B22j}
\]
Thus the tangent period of B19, in this specified algebraic
frame, has not been assumed bounded: its product with the
theta-null normalization has been identified with a directly
bounded analytic first derivative.

Use the same generator of the CM factor in its second
repeated copy, transported by the fixed algebraic isogeny.
The corresponding Betti directions are obtained from the
same Betti generator in the two copies.  They have the same
comparison scalar $p_N$.  Any original pair of selected
ball directions is obtained from these two by an invertible
matrix over $\overline\Q$: both are bases in the same
two-dimensional Betti eigenspace, with the rational
$E$-decomposition and the rational negative-line coordinates
fixed.  This is precisely the matrix $M_i$ in (B19d).
Consequently (B22j) can be imposed for each input type while
preserving B19's common scalar for both repeated directions.
The reference depends on the torus and the chosen direction,
not on the source character or the finite original input.

\paragraph{The resulting period certificate and its arithmetic cost.}
Apply this normalization to the two factors in (B19j).
The common constant becomes exactly
\[
                  C_x=\kappa_x\upsilon_1\upsilon_2
                                      H_{N_1}H_{N_2}.
\]
It follows that
\[
 |C_x|^2\geq
   |\kappa_x\upsilon_1\upsilon_2|^2
       \prod_{i=1}^2
        \left(2\pi|q_i^tu_i|
                      e^{-\pi q_i^tY_iq_i}\right)^2.
 \tag{B22k}
\]
Here $\kappa_x$ is in the fully specified frame of B19;
under the algebraic Betti choices of B21 it is computed by
the indicated nonzero determinant and square-root factors,
with the common universal oscillator factors retained.
Arbitrary original complex rescalings cannot be omitted.

For a nonzero algebraic finite sum $b$ in (B19j), suppose
$b\in K$, $[K:\Q]=d$, $Db$ is integral, and the other
$d-2$ conjugates are bounded by $M>0$ in B19's nonreal
embedding convention.  Combining (B19k) and (B22k) gives
\[
\begin{split}
 |\mathcal B_e(f)|^2\ \geq\ &
 |v_Hv_1^Tv_2^T\kappa_x\upsilon_1\upsilon_2|^2
       D^{-d}M^{-(d-2)}\\
 &{}\times\prod_{i=1}^2
        \left(2\pi|q_i^tu_i|
                      e^{-\pi q_i^tY_iq_i}\right)^2.
\end{split}
 \tag{B22l}
\]
For an orthonormal test vector use the $v_H^{1/2}$ factor
specified in B19 instead.  In combination with Q11 this
also gives a lower certificate for the geometric canonical
norm when its stated chart and evaluation data are supplied.

The constants $D,M$ refer to the algebraic array in these
new frames.  Dividing by the algebraic reference value in
(B22h), and by $\ell(v_0)$ when choosing $v$, can introduce
denominators and large conjugates.  Those factors must be
included in $D,M$ and in any prime-by-prime valuation bound.
The normalized tangent frame need not be integral or
saturated, so B20 or the archive's local divisibility
theorems require their integral hypotheses to be checked
again or their explicit change-of-frame costs to be applied.
No uniform arithmetic height bound follows from the analytic
Gaussian estimate.  The improvement is the removal of an
unspecified absolute CM-period factor from the normalized
finite-period certificate, with its arithmetic costs assigned
to actual algebraic data.\hfill$\square$


% THETA3-B23 contribution by /root/definite_bridge, 21 November 2026.
\subsection{THETA3-B23: an algebraic inverse on a full CM Hodge block}\label{B23}

The finite evaluation inverse in A5 has a motivic replacement on
a full CM Hodge block.  The replacement is an algebraic
correspondence multiplied by one explicitly identified positive
comparison scalar.  It is common to every vector in that block.
This statement does not require the Hodge conjecture for the ball
surface, and does not assert that the image of one fixed B17
kernel is a rational Hodge block.

\paragraph{The CM source and the geometric normalization.}
Let $X/\overline\Q$ be the smooth projective ball surface, or the
finite disjoint union of projective PEL components at the level
used in A5.  A further finite principal cover is allowed, with
its degree retained in pullback and trace.  Let $A/\overline\Q$
be a CM abelian variety and let
\[
 T\subset H^2(A,\Q),\qquad
 \alpha:T\longrightarrow H^2(X,\Q)
 \tag{B23a}
\]
be an algebraically projected rational Hodge summand and an
algebraic correspondence.  Assume that $T$ is a sum of non-Tate
simple Hodge structures.  For the prescribed finite family,
032, Corollary 6.7, printed p.45, gives finitely many rational
Hodge maps on degree-one cohomology from the specified CM sources.
Realize them by morphisms
$X\to B_i$ and combine these into $f:X\to A=\prod_iB_i$.
Then $\alpha=f^*$ on the non-Tate part of $H^2(A)$ contains
all the prescribed theta wedges in its $(2,0)$ image.
This is a full-block enlargement inside their actual CM cup
source; it does not require that the entire Albanese be CM.
Equivalently one may use the cup image of the CM part of
$H^1(\operatorname{Alb}X,\Q)$.  The degree-one realization
is 032, Lemma 2.2, printed p.5, and its use for the CM source
maps is in Section 6.  We use the accepted CM
Hodge conjecture for products of CM abelian varieties.

Put $M=\operatorname{im}\alpha$.  On holomorphic two-forms use
the geometric Hermitian pairing
\[
                  (f,g)_{\rm geom}=\int_X\overline f\wedge g.
 \tag{B23b}
\]
It is positive with the complex orientation.  The Haar
normalization in a particular automorphic realization is
recorded by a fixed positive constant $\nu_G$:
$\langle f,g\rangle_G=\nu_G(f,g)_{\rm geom}$.
One may choose geometric normalization, when $\nu_G=1$.
Finite covering degrees or fixed rational component weights
can instead be included in (B23b).  No unknown factor is
absorbed into a character-dependent normalization.

We retain A5's full adelic $\U(V)$ realization, its fixed
splitting, central transformation law, finite level, and scalar
archimedean canonical type.  If starting with a form on one
principal $\SU(V)$ component, first use the deeper principal
level and component realization of B12 and 032, Lemma 5.2,
printed pp.27--28; extend by zero on the other components.
All CM subspaces and projectors below refer to this realized
finite collection of algebraic components.  Thus the restriction
operators of A5 apply to $\mathcal E_\chi$; an unspecified
extension of an arbitrary $\SU(V)$ form is not being assumed.

\paragraph{An algebraic retract of the CM source.}
Every non-Tate simple summand of $M$ is primitive.  Indeed,
cup product with an ample class is a Hodge map
$M(-1)\to H^4(X,\Q)\simeq\Q(-2)$ on each component,
or equivalently $M\to\Q(-1)$ after untwisting.  A
non-Tate simple Hodge structure has no such quotient.
Consequently $q_X(x,y)=-\int_Xx\cup y$ is a polarization
on $M$.  Choose a rational polarization $q_T$ on $T$.
Define the polarization adjoint $\alpha^\dagger$ using $q_T$
and the nondegenerate middle cup pairing on $H^2(X)$.
This adjoint is algebraic: the transpose of the given
correspondence is algebraic, and the inverse polarization
tensor on the CM summand $T$ is a Hodge tensor on a product
of $A$, hence algebraic by the accepted CM Hodge conjecture.
Only this CM-source inverse is needed; there is no inverse
Lefschetz or Hodge-conjecture assumption on $X\times X$.

For clarity, all cycles needed here may be defined over
$\overline\Q$.  Spread a cycle over $\C$ representing a
specified rational Hodge tensor over a finite-type
$\overline\Q$ parameter scheme, restrict to a smooth connected
open on which its components are flat, and specialize at a
$\overline\Q$ point.  On the constant ambient abelian variety
the cycle class is locally constant in that connected family.
The specialization therefore has the same class.  The usual
algebraic degree projectors of an abelian variety are also
available, for example by polynomial interpolation in
multiplication by an integer greater than one.

Set $S=\alpha^\dagger\alpha$.  The positive Hermitian forms
obtained from the two polarizations make $S$ self-adjoint and
positive semidefinite, with
$\ker S=\ker\alpha$.  Thus it is semisimple.  If $r(t)\in\Q[t]$
is the product of the nonzero distinct factors of its minimal
polynomial, then $r(0)\ne0$.  Choose $a,b\in\Q[t]$ satisfying
$ta(t)+r(t)b(t)=1$ and put
\[
 d(t)=a(t)-\frac{a(0)}{r(0)}r(t),\qquad
 \beta=d(S)\alpha^\dagger,
 \qquad e_M=\alpha d(S)\alpha^\dagger.
 \tag{B23c}
\]
The zero image case is harmless and is discarded.  The
polynomial $d$ vanishes at zero and is $t^{-1}$ on the
nonzero spectrum of $S$.  Hence $e_M$ is the orthogonal
cohomological projector onto $M$, and $\beta$ is a section
of $\alpha$ on $M$.  Every operator in (B23c) is algebraic.
This is a statement about the action of correspondences on
cohomology; it makes no assertion of an idempotent identity
modulo rational equivalence in a Chow ring.

Moreover, every rational Hodge endomorphism $u$ of $M$ is
algebraic.  The map $\widetilde u=\beta u\alpha$ is a Hodge
endomorphism of the CM summand $T$, so is algebraic by the
same CM Hodge assumption.  The correspondence
$\alpha\widetilde u\beta$ acts as $u$ on $M$ and as zero
on its orthogonal complement.  This argument also works
with algebraic coefficients.  In particular, all the
isotypic projectors and matrix units used below are
realized by algebraic correspondences on $X$.

\paragraph{One simple type with its full multiplicity.}
Fix a non-Tate simple CM Hodge structure $S_0$ occurring in
$M$, with multiplicity $m$, and write its isotypic component
as $S_0^m$.  Its endomorphism field $K$ is a CM field and
$S_0$ has rank one over $K$.  Here is the elementary
representation-theoretic reason.  A simple rational
representation of the commutative Mumford--Tate group has
commutative division endomorphism algebra, and is
one-dimensional over that field.  The positive involution
from a polarization is complex conjugation.  If it were
trivial, the Hodge circle would have image in a compact
subgroup of a product of real multiplicative groups and
would be trivial.  A weight-two simple structure would
then be Tate.  This excludes the totally real case here.

Choose an embedding $\chi:K\hookrightarrow\overline\Q$
whose eigenspace has Hodge type $(2,0)$, and put
\[
 \mathcal E_\chi=(S_0^m\otimes_\Q\C)_\chi
                  \subset H^0(X,K_X),
 \qquad \dim_\C\mathcal E_\chi=m.
 \tag{B23d}
\]
The projector $e_\chi$ is allowed algebraic coefficients.
It is algebraic by (B23c), since
$\operatorname{End}_{\rm Hdg}(S_0^m)=M_m(K)$ and
$K\otimes_\Q\overline\Q$ splits into its embedding factors.
Distinct isotypic components and distinct embedding
components are orthogonal for the Hermitian polarization.
For the latter assertion, multiplication by $a\in K$ has
adjoint multiplication by $\overline a$; choosing $a$ whose
two embedding values differ forces the cross pairing to
vanish.

Choose compatible Betti and algebraic de Rham bases using
one copy of $S_0$ and the algebraic isomorphisms to its $m$
copies.  If $\beta_1,\ldots,\beta_m$ are the holomorphic
representatives of the resulting algebraic-coefficient
Betti classes and $\omega_1,\ldots,\omega_m$ the de Rham
basis, compatibility with comparison gives one scalar
$P_\chi\in\C^\times$ such that
\[
 \omega_j^{\rm an}=P_\chi\beta_j\quad(1\leq j\leq m),
 \qquad
 G_\omega=\nu_G|P_\chi|^2G_B,
 \quad
 (G_B)_{ij}=\int_X\overline{\beta_i}\wedge\beta_j.
 \tag{B23e}
\]
The matrix $G_B$ is positive definite with entries in
$\overline\Q$: cup product is rational on Betti cohomology,
and the classes and their conjugates have algebraic
coefficients.  The $\omega_j$ are actual algebraic
canonical forms.  On a smooth projective surface the
$F^2$ piece of algebraic de Rham cohomology is
$H^0(X,K_X)$, and the algebraic projector respects this
filtration.  The common period in (B23e) is a consequence
of the copy isomorphisms, not of independent rescaling of
the individual forms.

\paragraph{The finite evaluation rows and their actual fibre frames.}
Fix a rational negative line $L\subset V$ and its rational
orthogonal basis as in B12.  Use A5 for the full block
$\mathcal E_\chi$.  List all finite definite double cosets
occurring in its finitely many chosen translates, retaining
repetitions when they are different rows.  Their evaluation
points have the form $x_\ell=(L,b_\ell)$ with finite
adelic $b_\ell$.  These points and the finite Hecke branches
joining them are defined over $\overline\Q$.

There are direct proofs in this particular PEL model.
At the rational line, the three rational orthogonal
projectors preserve the Hodge filtration.  Adding integral
multiples of their graphs to the finite-type moduli
problem makes the chosen point isolated: a negative line
preserved by the projectors must be the specified negative
coordinate line.  An isolated complex point of a
finite-type scheme over a number field is algebraic.
This is precisely the argument in 032, Lemma 5.3,
printed pp.29--30.  For a finite adelic modification of
the lattice, clear denominators by an integer.  The
corresponding isogeny has a finite, $\mathcal O_E$-stable
torsion kernel.  Such kernels and finite level structures
are parametrized by finite etale schemes in characteristic
zero; their quotients and polarization homomorphisms are
algebraic.  The target lattice has the prescribed
polarization type because the adelic modification is
unitary.  Both projections of the resulting Hecke
correspondence become finite etale after passing to a
common level.  The quotient and polarization descent are
the same constructions written explicitly for the split
spherical correspondences in 032, Section 5.3,
printed pp.30--32.  This argument allows the finitely many
adelic translates used by A5, not just one split prime.

Let $\kappa_\ell$ denote the analytic canonical-fibre frame
used to define the automorphic scalar at $x_\ell$.  Choose
one algebraic frame $\eta_0$ at the base point and transport
it through the algebraic differentials of these finite
Hecke branches.  Since the two maps on the marked ball
chart have the same archimedean coordinate, their
differentials also identify the frames $\kappa_\ell$.
Consequently the transported algebraic frames satisfy
\[
                  \eta_\ell=t_x\kappa_\ell
                  \quad\hbox{with one }t_x\in\C^\times.
 \tag{B23f}
\]
If independent algebraic choices of the $\eta_\ell$ are
preferred, their nonzero algebraic ratios are retained
in the row matrix below.  Formula (B23f) does not assert
that a Betti ball-coordinate frame is algebraic.

With probability Haar measure on each definite quotient,
let $D$ be the diagonal matrix of the finite double-coset
masses.  Its positive entries are rational: they are
reciprocal finite stabilizer orders divided by their
finite sum.  Rows forced to vanish by a stabilizer
character can be removed.  Define
\[
 A_{\ell j}=\frac{\omega_j(x_\ell)}{\eta_\ell}
       \in\overline\Q,
 \qquad R=t_x A,
 \qquad H=A^*DA\in M_m(\overline\Q).
 \tag{B23g}
\]
The finite family in A5 separates $\mathcal E_\chi$, so
$A$ has full column rank and $H$ is positive definite.

\paragraph{The canonical-fibre period, including the Tate factors.}
The number $t_x$ has a specified CM-period meaning.  Use
the actual reference PEL abelian scheme of 032, Section
5.1.  Its active Hodge bundles have ranks two and one;
write them as $\mathcal P=\omega_\mu$ and
$\mathcal L=\omega_{c\mu}$.  Gauss--Manin and the
polarization give algebraic isomorphisms
\[
 \operatorname{KS}:\mathcal P\otimes\mathcal L
                         \simeq\Omega_X^1,
 \qquad
 \det\mathcal P\otimes\mathcal L^{\otimes2}
                         \simeq K_X.
 \tag{B23h}
\]
The other embedding summands are constant along the ball
and contribute no tangent directions.  To check that
the active map is an isomorphism, differentiate the
negative line and its annihilator in a marked period
chart.  The two independent coordinates of the line
give the two independent components of
$\mathcal P\otimes\mathcal L$ under the nondegenerate
polarization pairing.  Thus its determinant is nonzero.
This also proves (B23h) for the higher-degree CM field,
without importing a dimension-three assertion globally.

At $x$, choose algebraic bases of $\mathcal P_x$ and
$\mathcal L_x$ and write their comparisons with the
rational graph bases as a matrix $\Omega_\mu$ and a
scalar $P_{c\mu}$, respectively.  The graph bases and
the derivative of the rational line have algebraic
entries.  The de Rham polarization on $H^1$ compares
with its Betti alternating pairing by $(2\pi i)^{-1}$.
There are two such pairings in the determinant of
Kodaira--Spencer.  Hence, for the frame obtained from
these algebraic bases,
\[
 t_x=\zeta_x\,
       \frac{\det(\Omega_\mu)P_{c\mu}^{\,2}}
                         {(2\pi i)^2},
                 \qquad\zeta_x\in\overline\Q^\times.
 \tag{B23i}
\]
Here $\zeta_x$ is the explicitly determined determinant
of the algebraic graph, polarization, and tangent-frame
matrices.  It can be included in $A$ instead if desired.
The factor $(2\pi i)^{-1}$ can also be checked by pairing
two de Rham one-classes with the $(g-1)$st power of an
algebraic polarization class: its comparison contributes
$(2\pi i)^{g-1}$, while the top-degree de Rham trace
contributes $(2\pi i)^{-g}$.

An exact primary normalization check is de Shalit--Goren,
\emph{A theta operator on Picard modular forms modulo
an inert prime}, Section 1.8, Proposition 1.22 and
Section 1.11.2, Lemma 1.27 and Proposition 1.28:
equation (1.141) is
$\operatorname{KS}(d\zeta_1\otimes d\zeta_3)
=-\delta(2\pi i)^{-1}du$ and
$\operatorname{KS}(d\zeta_2\otimes d\zeta_3)
=(2\pi i)^{-1}dz$.
The complete displayed computation and its algebraic
interpretation in Section 1.11.3 give precisely the
power in (B23i).  The version on the authors' site has
these statements on printed pp.26--27 and 31--33;
the published article uses the same numbered statements.
No identification $\det\mathcal P\simeq\mathcal L$
with an unspecified scalar is used here.

\paragraph{Algebraic matrix and common positive scalar.}
In the coherent basis $\omega$, the operator of A5 and
its inverse are exactly
\[
 \begin{split}
 Q_\chi&=
     \frac{|t_x|^2}{\nu_G|P_\chi|^2}\,G_B^{-1}H,\\
 Q_\chi^{-1}&=s_\chi K_\chi,
 \qquad
 s_\chi=\frac{\nu_G|P_\chi|^2}{|t_x|^2}>0,
 \qquad K_\chi=H^{-1}G_B\in M_m(\overline\Q).
 \end{split}
 \tag{B23j}
\]
Indeed $R^*DR=|t_x|^2H$ and (B23e) is the domain
Gram matrix; the definition of the adjoint gives the
first equality and inversion gives the second.
Because this is the full multiplicity of a single
embedding component, the matrix $K_\chi$ extends by
zero to an element of
$\operatorname{End}_{\rm Hdg}(M)\otimes_\Q\overline\Q$.
The CM-source retract above therefore realizes it as
an algebraic correspondence on $X$.  On the selected
complex block it is positive and self-adjoint for
the geometric pairing, since it is a positive real
multiple of $Q_\chi^{-1}$.  Positivity of all algebraic
conjugates is neither asserted nor needed.

For every $f,g\in\mathcal E_\chi$, A5 now reads
\[
 \langle f,g\rangle_G
   =s_\chi\sum_{j,\alpha}
       \overline{\mathcal P_{j,\alpha}(K_\chi f)}
                       \mathcal P_{j,\alpha}(g).
 \tag{B23k}
\]
The periods are the actual totally definite $\U(2)$
Bessel periods with A5's fixed translates and probability
Haar measures.  The same $K_\chi$ and $s_\chi$ work for
every source character whose image lies in this block.
When two selected CM tensor lines have the same complete
Hodge type, the algebraic CM comparison of A10 places
their nonzero images in this same isotypic block and
embedding component.  This is the appropriate meaning
of a common normalization for balanced pairs.

There is also an exact one-reference calibration.
Choose $0\ne f_0\in\mathcal E_\chi$ once and put
$B(f,g)=\sum_{j,\alpha}
\overline{\mathcal P_{j,\alpha}(K_\chi f)}
\mathcal P_{j,\alpha}(g)$.  Then $B(f_0,f_0)>0$ and
\[
 \langle f,g\rangle_G
       =\frac{\langle f_0,f_0\rangle_G}{B(f_0,f_0)}
                           B(f,g).
 \tag{B23l}
\]
If $f_0$ is specified by an algebraic-coefficient Betti
class its geometric squared norm is algebraic.  The
denominator in (B23l) still contains the common period;
this formula does not claim that $s_\chi$ is algebraic.
Formula (B23i) identifies the exact remaining
Kodaira--Spencer CM tensor, rather than hiding it in
an arbitrary evaluation normalization.

\paragraph{What block projection repairs, and what it leaves.}
Evaluation at a point is not a morphism of weight-two
Hodge structures.  Before projection, the finite
evaluation Gram matrix can have nonzero cross terms
between distinct CM blocks, although the geometric
cup pairing has none.  Replace each restriction by
$R_ae_\chi$ and apply (B23j) separately on its full
block.  The algebraic projectors remove these cross
terms exactly.  Averaging over a CM action alone
would only force a commuting operator; even in
multiplicity one its positive scalar remains an
analytic norm ratio.  No algebraicity of that scalar
follows from the commutation relation.

The assertion proved here concerns the full CM block
with its algebraic global forms and all selected
finite evaluation rows.  A particular image
$\mathcal E_0$ of the B17 integral operator need not
be a rational Hodge summand, and its adjoint involves
an integral over the original source group.  Neither
algebraicity of that integral Gram matrix nor equality
of its inverse with $K_\chi$ has been established by
the finite-frame argument.  Thus (B23k) closes the
algebraic-frame correction for A5 on a full CM block,
up to the single specified positive factor, while
the literal identification with one fixed arithmetic
\texttt{MainX}/\texttt{WeilTensorModels} package remains
a separate kernel and normalization comparison.
In particular, the coordinates of an original unnormalized
theta vector in the coherent basis $\omega$ can contain
the type-dependent constants retained in B19--B21.
The algebraicity of $K_\chi$ does not make those
coordinates algebraic.

\paragraph{The Grothendieck framework.}
The supplied manuscript locators are
Grothendieck's \emph{cote 44}, scans 7--8, in the letter
to Serre dated 12 August 1964: the polarization is
constructed through a transposed cycle correspondence,
and its kernel is removed before passing to the abelian
quotient.  The inverse-correspondence and critical-degree
discussion is \emph{cote 16}, scans 43--46, especially
the equivalences (iv) and (iv bis) on scan 44 and
their transposition argument on scan 45.  Here (B23c)
constructs the particular inverse on a CM source by a
rational polynomial and the accepted CM Hodge theorem.
It therefore proves the needed homological retract;
it does not invoke the general standard conjectures
or replace the distinction between a cohomological
projector and a Chow idempotent.


% THETA3-B24 contribution by /root/definite_bridge, 21 November 2026.
\subsection{THETA3-B24: the actual fixed-kernel inverse for a definite middle source}\label{B24}

In the totally definite natural-source sector, the integral
Gram matrix of B17 becomes a finite algebraic Gram matrix
times one common comparison scalar.  This identifies the
actual fixed-kernel inverse projectively with an algebraic
correspondence.  The same-type condition in A7 is retained:
this sector by itself gives divisor-generated balanced
tensors, rather than the general exceptional balanced type.

\paragraph{Fixed data and the scalar source type.}
Use the rank-three target $V$ and rank-two multiplicity
space $P=W_1\perp W_2$ of A7, with $d_1=d_2=d$ as
complete sign functions.  Then $P$ is totally definite.
Fix the same additive character, Weil splittings, target
and source levels, coefficient frames, and polynomial
Gaussian as in B17.  Require algebraic finite Schwartz
coefficients and algebraic values of the finite splitting
characters on the finitely many translates used below.
These conditions hold for the usual finite Weil model
with algebraic splitting characters.  Give
$[S]=[\U(P)]$ probability Haar measure and retain the
target conversion constant $\nu_G$ of B23.

Let $K(g,h)$ be the actual canonical scalar coefficient
kernel, including the determinant cocycle.  There is one
unitary character $\kappa_S$ of the full compact group
$S_\infty$ for which
\[
                 K(g,hk)=\kappa_S(k)K(g,h)
                    \quad(k\in S_\infty).
 \tag{B24a}
\]
This is stronger than just finite maximal-compact type.
Indeed the source action preserves the positive
Lagrangian used by the Gaussian at every real place.
At the active place its degree-two exterior factor is
$\bigwedge^2P_\mu^*$, a line; the half-determinant
factor is also a line.  At every other place, equality
of the two signs makes the chosen whole embedding
summand source-stable, and the Gaussian is a vacuum
line.  Their product transforms by a character.
These are precisely the determinant and exterior
factors in 032, Lemma 4.1 and equations (4.6)--(4.8),
printed pp.15--16.
The specified Weil splitting transports the bare
half-determinant characters to this honest character
of $S_\infty$.  This proof does not equate bare and
circle-valued splitting conventions.  Compactness
makes the character unitary.  It is independent of
the torus characters subsequently integrated in A7.

Choose finite representatives $s_b$ for
$S(F)\backslash S(\A_{F,f})/K_{S,f}$.  Let $D_S$ be
the diagonal probability-mass matrix.  Its entries
are positive rational numbers proportional to the
reciprocals of the finite stabilizer orders.
An incompatible stabilizer character forces the
corresponding scalar row to vanish; remove that row.
The surviving masses are not renormalized after
this removal.
The effective integrating source module has right
infinity type $\kappa_S^{-1}$, since B17 uses the
bilinear integral $Tu=\int K(g,h)u(h)\,dh$.
It is determined by its values at the $s_b$, and
its Hermitian norm is $z^*D_Sz$.
The operator on unrestricted $L^2([S])$ factors
through this finite module: averaging over
$S_\infty$ and $K_{S,f}$ annihilates the other
isotypic components.  Thus this reduction does not
alter B17's actual operator $TT^*$.

\paragraph{Algebraicity of the finitely many translated inputs.}
At a finite place, the Schr\"odinger Weil action is
generated by linear substitutions, additive-character
multiplications, and Fourier transforms.  Applied to
a locally constant compactly supported function,
each Fourier transform is a finite sum of roots of
unity times lattice volumes; the determinant factors
are powers of a rational residue cardinality, with
specified square roots, and the finite Weil indices
are roots of unity.  Rational or self-dual Haar
normalizations give algebraic lattice volumes.
Thus all these numbers are algebraic.
For explicit generator formulas with the splitting
character retained, see Chao Li, \emph{Geometric and
arithmetic theta correspondences}, Section 2.1,
printed pp.2--3: the Levi factor, unipotent factor,
and Fourier transform have exactly these forms.
The same argument applies after restricting the
ambient symplectic Weil model to the fixed unitary
splittings; their additional finite character
values are the hypothesis already stated.
The splitting-character values have been retained
as an explicit hypothesis above.  It follows that
the finitely many source and target finite translates
of the chosen input are again Schwartz functions
with algebraic values.  A common level and lattice
contain all of them.

Each such finite function is a finite algebraic sum
of tensor products for $W_1\perp W_2$: choose product
cosets fine enough to support and trivialize the
function.  The canonical determinant coefficient
therefore is a finite algebraic sum of wedges of
the two same-type theta one-forms.  Corollary 6.7
of 032 applies to this finite collection.  After
one common principal cover, all its kernels lie
in the cup image of a finite product of the actual
CM sources, as in B23.

They lie in one simple CM type with multiplicity,
and in one of its $(2,0)$ embedding components.
Here is a precise way to see this assertion even
if the original CM type is imprimitive.  In the
two copies of $U(v)$, project their tensor product
to equal $E$-embedding labels.  The rational
idempotent is supplied by the separable algebra
$E\otimes_\Q E$; its image is a rank-one
$E$-Hodge structure whose type at each embedding
is $(2,0)$ or $(0,2)$ according to $v$.  Its
Mumford--Tate characters form a single Galois
orbit, possibly with repetitions, because the
embeddings of the field $E$ form a single orbit.
It is consequently isotypic as a rational Hodge
structure.  The chosen embedding line belongs
to one embedding component of that simple type.
Every rational pair of CM source maps followed
by cup product preserves this assertion.
The complex linear spans in Corollary 6.7 do
not change the component.  Thus the finite
kernels belong to a full block
$\mathcal E_\chi$ of the kind in B23; their
image can be a proper subspace of that block.

\paragraph{One common CM evaluation scalar determines the source matrix.}
Choose the rational negative line and frame before
fixing the kernel, as in B17.  Select the finite
target translates of A5 separating the full block
$\mathcal E_\chi$, and use its coherent basis
$\omega=P_\chi\beta$.  Write $A$ for B23's
algebraic full-column-rank evaluation matrix, so
the actual target evaluation matrix is $t_xA$.

B19 applies at this one fixed CM period point
to all the finitely many translated inputs.
Choose one common finite theta module and its
algebraic Heisenberg isomorphism for each line.
Its determinant and tangent factors give one
nonzero constant $C_x$ for all these values.
With the oscillator constants, level factors,
and tangent matrices kept exactly as in B19,
there is a matrix $B\in M_{r_0\times b_0}(\overline\Q)$
such that
\[
       K(x_\ell,s_b)=C_x B_{\ell b}.
 \tag{B24b}
\]
Here $r_0$ is the number of target evaluation rows
and $b_0$ the number of surviving source cosets.
No torus integral is replaced by a finite sum:
(B24b) is a theta-kernel value.  The only
infinity integral used to reduce the source
has already been evaluated by (B24a).

Expand $K(\,\cdot\,,s_b)=\sum_a\omega_a E_{ab}$.
Choose a left inverse $L_A$ of $A$ over
$\overline\Q$, by row elimination.  Comparing
the two evaluations gives
\[
          E=\frac{C_x}{t_x}E_{\rm alg},
                  \qquad E_{\rm alg}=L_AB.
 \tag{B24c}
\]
The equality is exact because the evaluation
family separates the block.  This proves that
the image $\mathcal E_0=\operatorname{im}T$
is defined over $\overline\Q$ in the coherent
de Rham coordinates: it is the column span of
$E_{\rm alg}$.  It is not an assertion obtained
merely by declaring fibre evaluation to be a
Hodge operation.

Assume $\mathcal E_0\ne0$.  Choose a full-rank
algebraic column matrix $F$ spanning it and write
$E_{\rm alg}=F E_0$.  In the basis $\omega F$ put
\[
       G_0=F^*G_BF,
       \qquad W_0=E_0D_SE_0^*.
 \tag{B24d}
\]
Both are positive definite algebraic Hermitian
matrices.  Positivity of $W_0$ follows because
$E_0$ has full row rank.  On this basis the
bilinear integral $T$ is the matrix
$(C_x/t_x)E_0D_S$.  Its source Gram matrix is
$D_S$, and its target Gram matrix is
$\nu_G|P_\chi|^2G_0$.  Direct calculation of
the adjoint therefore proves
\[
 \begin{split}
 Q_T=TT^*&=
      \nu_G|P_\chi|^2\left|\frac{C_x}{t_x}\right|^2
                                      W_0G_0,\\
 Q_T^{-1}&=
      \frac{|t_x|^2}{\nu_G|P_\chi|^2|C_x|^2}
                             G_0^{-1}W_0^{-1}.
 \end{split}
 \tag{B24e}
\]
In particular this is the actual B17 inverse,
not the different finite restriction inverse
of A5.

The last matrix is realized by an algebraic
correspondence on the CM cup image.  Indeed,
the $G_B$-orthogonal projector onto the
column space of $F$ is the algebraic matrix
$F(F^*G_BF)^{-1}F^*G_B$.  Extending the
operator on that space by zero gives an
algebraic matrix on the full multiplicity
block.  B23's CM-source retract realizes
every such matrix by a correspondence.
The use of the transpose and the inverse source
polarization here is precisely B23's construction,
with the supplied Grothendieck \emph{cote 44},
scans 7--8, and \emph{cote 16}, scans 43--46,
as the framework locators.  It does not invoke
an inverse-polarization assertion for all of $X$.
It is positive self-adjoint on the selected
complex image, because it is a positive
multiple of $Q_T^{-1}$.

Consequently B17's exact adjoint double lift
and its literal rank-zero \texttt{LastLaw}
and \texttt{SeeLast} specialization use one
algebraic correspondence and the one positive
scalar displayed in (B24e), uniformly for
all torus characters integrated against this
fixed kernel.  The original scalar $C_x$
retains the exact B19 normalization; B22 can
normalize and bound its analytic reference
factors.  Algebraic denominators, conjugate
sizes, and the scalar $P_\chi/t_x$ are not
discarded.  The finite positive matrix $W_0$
admits the lattice/determinant certificates
of B7 whenever their stated arithmetic data
are provided.

\paragraph{The extension to other finite compact types.}
For a totally definite source and a finite
sum of compact infinity types, the same
argument applies whenever all coefficient
values share the common comparison scalar
used in (B24b).  One first expands in
matrix coefficients of the compact
representations, in compatible algebraic
bases.  An algebraic positive invariant Hermitian
Gram matrix can be chosen: infinitesimal invariance
for an algebraic basis of the compact real Lie
algebra, together with any component representatives,
is a finite real linear system over the real
algebraic numbers.  For the connected unitary groups
here the Lie algebra equations suffice.  Averaging
any positive real Gram matrix gives a positive real
solution.  Algebraic points of the solution space
are dense, and positivity is open, so a positive
algebraic solution exists.  Algebraic Cholesky
factorization then gives an orthonormal algebraic
change of basis.  Haar probability integration is
algebraic in those bases: averaging an
operator between two irreducibles gives
zero for distinct types and a scalar
identity for one type; taking the trace
gives the scalar $1/\dim\rho$ in the
usual matrix-coefficient orthogonality
formula.  Finite stabilizer projections
are averages over the stabilizers.
Their orders, the dimensions, and the
chosen algebraic representation matrices
therefore give an algebraic source Gram.
This proves the extension under the stated
common-scalar hypothesis.  That hypothesis
is proved above for the actual scalar
determinant cocycle, not asserted for
arbitrary unrelated infinity types.

\paragraph{The remaining indefinite-source matrix.}
For general anisotropic $P$ with a
noncompact real factor, write a finite-rank
kernel as $K(g,h)=\sum_a\omega_a(g)e_a(h)$.
Its source Gram is
\[
 W_{ab}=\int_{[S]}e_a(h)\overline{e_b(h)}\,dh,
 \qquad Q_T=\nu_G|P_\chi|^2W G_B
 \tag{B24f}
\]
when its image is the specified full block.
There is no compact Schur reduction to
the finitely many values in (B24b).
Projective algebraicity of this positive
matrix $W$ is the precise remaining
condition for a projectively algebraic
inverse in these coordinates.  It cannot
be forced by an invertible algebraic
change of target basis: if $W$ is changed
to $CWC^*$ with $C$ algebraic, one such
matrix is a positive scalar times an
algebraic matrix if and only if the
other is.

Finally, A7 shows that this natural
totally definite rank-two source has
equal complete types.  A balanced second
pair then has the same equal types, so
the selected four-factor line is a
product of divisor lines.  Formula
(B24e) does not settle the indefinite
source in a general exceptional switch.
Nor does a finite sum of identities
for different spaces $P$ constitute
one \texttt{MainX} tuple: the recovered
\texttt{Comm} fixes its Hermitian spaces
and requires its upper spaces to be
\texttt{DefSp}.  The exact clauses are
\path{Development/AbelianVarietyTheory/Definitions/}
\path{LatticeFrames/ValLastX.lean},
\texttt{Comm}, lines 30--69, and
\texttt{LawsX}, lines 212--221.


% THETA3-B25 contribution by /root/definite_bridge, 21 November 2026.
\subsection{THETA3-B25: the actual theta image has an algebraic finite-period inverse}\label{B25}

The finite-source argument of B24 has a part that does not
require the source to be totally definite.  Finite adelic
source translates already span the actual finite theta
image.  Their common CM evaluation scalar proves that
this image is algebraic in the coherent coordinates of
B23.  Consequently the algebraic finite-period correction
can be chosen to preserve the original theta image.
The integral source Gram matrix of B17 is not replaced
by this finite restriction matrix.

\paragraph{Fixed kernel and the finite-adelic spanning lemma.}
Retain all the data of B17, including the full adelic
target realization, its scalar canonical type, the
coefficient pairing, additive character, splittings,
finite levels and Haar measures.  In particular
\[
 K(g,h)=\theta_{V,P}(g,h;\Phi),\qquad
 Tu=\int_{[S]}K(\,\cdot\,,h)u(h)\,dh,
 \qquad \mathcal E_0=\operatorname{im}T,
 \quad S=\U(P).
 \tag{B25a}
\]
The source $P$ is the actual rank-two multiplicity
plane, definite at the active place but possibly
indefinite elsewhere.  Both quotients are compact,
and B17 proves that the image is finite dimensional
and is the span of all $K(\,\cdot\,,h)$.
Assume, as in B24, that the fixed finite input has
algebraic coefficients and that the finite splitting
values used in the finitely many translations below
are algebraic.  All specified original complex
normalization scalars are retained.

Then
\[
 \mathcal E_0=
  \operatorname{span}_{\C}
       \{K(\,\cdot\,,(1,b)):b\in S(\A_{F,f})\}.
 \tag{B25b}
\]
In particular, if $r=\dim\mathcal E_0>0$, one can
choose $r$ such finite adelic elements whose kernel
vectors are a basis.

\textbf{Proof.}
The Cayley transform identifies open sets of $S$
and its $F$-vector space of skew-adjoint matrices:
$X\mapsto(1-X)(1+X)^{-1}$, with the same formula
for its inverse.  Weak approximation for the
coefficients of that vector space shows that
$S(F)$ is dense in $S(F\otimes_\Q\R)$; the omitted
real locus having eigenvalue $-1$ is a proper
algebraic subset and has empty interior.  This is
an assertion for the full unitary group, not an
assumption about a particular special-unitary
arithmetic subgroup.

Let $h=(h_\infty,h_f)$.  Choose
$\gamma_j\in S(F)$ with
$\gamma_{j,\infty}\to h_\infty$.  Rational theta
automorphy gives the exact vector equality
\[
 K(\,\cdot\,,(1,\gamma_{j,f}^{-1}h_f))
       =K(\,\cdot\,,(\gamma_{j,\infty},h_f)).
 \tag{B25c}
\]
The right side converges to $K(\,\cdot\,,h)$ in
the finite-dimensional holomorphic target space,
by the smooth coefficient dependence proved in A7.
The span on the right of (B25b) is a linear
subspace of a finite-dimensional vector space,
so it is closed.  It therefore contains every
kernel vector.  The reverse inclusion is immediate.
Selecting a basis proves the last assertion.
No uniform-continuity assertion over varying finite
representatives is needed in this argument: $h_f$
is fixed on the right of (B25c). \hfill$\square$

\paragraph{The one CM type and the common finite-value scalar.}
Fix the original rational decomposition
$P=P_1\perp P_2$, with complete types $v_1,v_2$.
Each of the finitely many vectors in (B25b) has the
same archimedean coefficient as the original kernel.
Its finite input is the corresponding finite Weil
translate of $\Phi_f$.  The finite generator
calculation in B24 proves that its coefficients
are algebraic.  Every such input is a finite
algebraic sum of product cosets for $P_1\perp P_2$.
The exact double alternation in 032,
Proposition 4.6(i), printed pp.22--23,
then expresses the kernel vector as a finite sum
of wedges of the two prescribed types.

Apply 032, Corollary 6.7, printed p.45, to this
finite collection of actual inputs.  After one
common principal cover, their degree-one classes
come from finitely many rational maps from the
specified CM modules.  Use B23's algebraic CM
cup source on that cover.  This does not require
that all of $H^1$ be generated by theta forms.

The relevant source tensor is the equal-label
summand $U(v_1)\otimes_E U(v_2)$, obtained by the
diagonal idempotent in $E\otimes_\Q E$.
It has rank one over $E$.  At an embedding its
Hodge type is determined by $v_1+v_2$ and is one
of $(2,0),(1,1),(0,2)$.  Its Mumford--Tate
characters form one Galois orbit, possibly with
repetitions, since the embeddings of $E$ form
one orbit.  It is therefore isotypic as a
rational CM Hodge structure.  The active
selected line has type $(2,0)$, so its simple
type is non-Tate.  Taking rational source maps
and then cup product preserves this assertion.
Consequently all the selected kernel vectors,
and hence all of $\mathcal E_0$, lie in one
full multiplicity block $\mathcal E_\chi$ of B23.

In the coherent basis $\omega=P_\chi\beta$ of
that block, choose A5's separating finite
target evaluation rows.  Their matrix is
$t_xA$, with $A$ algebraic and of full column
rank.  Evaluate the finitely many source
basis vectors at these rows.  The archimedean
positive Lagrangian and rational CM base point
are fixed; only finitely many algebraic finite
Weil inputs have changed.  B19 therefore gives
one nonzero comparison scalar $C_x$, with its
same theta-module, tangent, level and constant-line
factors, such that their evaluation matrix is
$C_xB$, $B\in M(\overline\Q)$.

Let $E$ be the coefficient matrix of these
kernel vectors in the basis $\omega$.  For
any algebraic left inverse $L_A$ of $A$,
\[
       E=\frac{C_x}{t_x}L_AB.
 \tag{B25d}
\]
Thus $\mathcal E_0$ is a subspace defined over
$\overline\Q$ in coherent coordinates.  Choose
an algebraic full-rank column matrix $F$ whose
columns are a basis of that subspace.
This proof uses actual CM-point theta values;
it does not treat evaluation as a morphism of
weight-two Hodge structures.  Nor does it
assert that the original, unnormalized theta
basis vectors themselves have algebraic
coherent coordinates: their common factor
$C_x/t_x$ remains in (B25d).

\paragraph{A correction preserving the original image.}
Write
\[
 G_0=F^*G_BF,\qquad A_0=AF,\qquad
 H_0=A_0^*DA_0,
 \tag{B25e}
\]
where $D$ is the finite positive diagonal
probability-mass matrix for the definite
restrictions.  It has rational entries, with
the original masses retained if stabilizer
conditions force some rows to vanish.
Both $G_0$ and $H_0$ are positive definite
algebraic Hermitian matrices.
Let $Q_{\rm fin}$ be A5's actual restriction
operator on $\mathcal E_0$, with its
inherited Petersson pairing.  In the basis
$\omega F$ it satisfies
\[
 \begin{split}
 Q_{\rm fin}&=
  \frac{|t_x|^2}{\nu_G|P_\chi|^2}G_0^{-1}H_0,\\
 Q_{\rm fin}^{-1}&=s_\chi K_0,
 \qquad s_\chi=\frac{\nu_G|P_\chi|^2}{|t_x|^2},
 \qquad K_0=H_0^{-1}G_0.
 \end{split}
 \tag{B25f}
\]
These identities follow by taking the adjoint
with domain Gram matrix $\nu_G|P_\chi|^2G_0$
and range Gram matrix $D$.

The algebraic matrix $K_0$ extends by zero
on the $G_B$-orthogonal complement of the
column space of $F$.  That projection has
matrix $F(F^*G_BF)^{-1}F^*G_B$.
B23's CM-source retract realizes the extended
matrix by an algebraic correspondence with
$\overline\Q$ coefficients on the target.
In particular $K_0$ preserves this very
$\mathcal E_0$, not merely its full containing
CM block.

Choose the definite test functions and
periods $\mathcal P_{j,\alpha}$ exactly as
in A5.  Then, for all $f,g\in\mathcal E_0$,
\[
 \langle f,g\rangle_G
   =s_\chi\sum_{j,\alpha}
       \overline{\mathcal P_{j,\alpha}(K_0f)}
                        \mathcal P_{j,\alpha}(g).
 \tag{B25g}
\]
The tests are actual totally definite rank-two
subgroup periods with the stated Haar measures.
Every corrected form $K_0f$ still lies in the
image of the original fixed theta operator.
Thus B17's actual adjoint double lift and
literal last-step seesaw apply to it using
the same $P,V,\Phi$ and splittings.
This closes the image-preservation condition
in this finite-period comparison.  The
single $K_0$ and $s_\chi$ work for all A7
source characters with these fixed data.

\paragraph{The actual fixed-data variation in 032.}
The parameter space in 032, Section 4.6,
printed pp.22--24, varies only the positive
and negative lines of the same multiplicity
plane at the compact target places.  Its
active positive Lagrangian, rational total
space, finite representation and total finite
input are fixed.  Each such line space is
an orbit of the appropriate $\U(P_\tau)$.
Choose $s_t\in S(F\otimes_\Q\R)$ carrying
the initial lines to those at $t$, and choose
its active component to be the identity.
In a parameter neighborhood it can be
chosen smoothly.

Heisenberg covariance transports the
annihilator equations for the Gaussian.
Consequently the specified Weil action
$\omega(s_t)$ carries the initial nonzero
Gaussian vacuum line to the new Gaussian
vacuum line.  Their chosen vectors differ
by a nonzero scalar.  The active determinant
cochain is unchanged because $s_{t,\mu}=1$.
The selected determinant half-lines and
their square-root frames add only their
recorded nonzero scalar.  These factors
are independent of the target ball variable:
only the compact-place data have changed,
and the active data are identical.  Equivalently
the two coefficient identifications have an
invariant ratio under the transitive target
special-unitary action, hence a constant
ratio, exactly as in the proof of 032,
Proposition 4.6(i).
Thus, with the original splitting retained,
the literal scalar coefficient satisfies
\[
       \Phi_t=c_t\omega(s_t)\Phi_0,
       \qquad
       \Xi(\Phi_f;t)=c_tK(\,\cdot\,,s_t),
       \qquad c_t\ne0.
 \tag{B25h}
\]
Here $c_t$ is not asserted algebraic or
silently set equal to one.  This statement
uses A7's common closed-cochain realization;
the continuity assertion of 032 alone did
not claim holomorphy at intermediate $t$.
It follows from (B25h) and A7 that every
such form belongs to the same $\mathcal E_0$.

There is also an independent check of the
intermediate holomorphy using only 032's
displayed conclusions.  Rational nonisotropic
vectors of $P(E)$ can approximate prescribed
nonisotropic vectors at all real places
simultaneously.  Their lines and orthogonal
complements give rational decompositions
dense in this parameter product; the active
definite place imposes no additional sign
condition.  Proposition 4.6(i) makes the
corresponding forms holomorphic, and (iii)
gives convergence in $C^\infty$ on the fixed
quotient.  Passing to the limit in
$\bar\partial\Xi_t$ proves it is zero.
An individual wedge expansion can be checked
on its deeper principal cover; vanishing of
$\bar\partial$ then descends to the fixed
quotient.  This explains why no holomorphy
assumption at intermediate parameters is
being added.

In particular the rational balanced
endpoints of 032, Corollary 4.7, printed
pp.24--25, use the same preserved-image
correction (B25f).  Their source tori may
be different rational subgroups of this
same $S$; integrating any of the resulting
torus functionals still takes values in
the same finite image.
Balance also identifies their complete
rank-one $E$ tensor Hodge types, since it
is exactly equality of the functions
$v_1+v_2=v_3+v_4$.

\paragraph{The remaining integral Gram.}
It is essential to distinguish (B25f) from
the inverse of $Q_T=TT^*$ in B17.
The latter still involves the actual matrix
$W_{ab}=\int_{[S]}e_a\overline{e_b}\,dh$.
The present result makes its coefficient
space algebraic, but does not prove that
this integral matrix is projectively algebraic.
In a last-step application to $K_0f$, the
reverse input $Q_T^{-1}K_0f$ is still the
analytic inverse specified in B17.
Nor is a finite sum of these period identities
one prescribed arithmetic \texttt{MainX}
tuple.  Its exact kernel, lattice and
normalization clauses remain separate.

\paragraph{Framework and source scope.}
The CM correspondence realization is exactly
B23's transpose and inverse-CM-polarization
construction, with the supplied Grothendieck
\emph{cote 44}, scans 7--8, and
\emph{cote 16}, scans 43--46.
The actual analytic source clauses are
\path{Development/AbelianVarietyTheory/Definitions/}
\path{LatticeFrames/IsQuasiFam.lean},
\texttt{LastLaw}, lines 71--75, and
\texttt{SeeLast}, lines 110--123, as
instantiated with the coefficient pairing
and dagger in B17.  The compact-source
finite sum of B24 is recovered when
$P$ is totally definite; no such finite
source quotient is assumed in (B25b).


% THETA3-B26 contribution by /root/definite_bridge, 21 November 2026.
\subsection{THETA3-B26: the mixed source Gram reduces to algebraic Hecke blocks and their multiplicity periods}\label{B26}

Keep the actual fixed kernel of B17 and B25.  There is a further
reduction of its source Gram matrix which does not replace the
kernel by a finite frame.  First its coefficient functions are
holomorphic sections of a specified algebraic automorphic line,
and their coherent coordinates have one common comparison scalar.
Second the actual Petersson form is constrained by an algebraic
Hecke algebra with adjoints.  On one irreducible Hecke module this
already makes the actual Gram projectively algebraic.  In general
the remaining quantities are positive Hermitian forms on the
Hecke multiplicity spaces.  The cohomological calculation below
identifies the periods represented by those remaining forms.

\paragraph{The actual source line and its central factor.}
Let
\[
 I=\{\tau\ne\mu:d_1(\tau)\ne d_2(\tau)\},\qquad s=|I|.
 \tag{B26a}
\]
The source $S=\U(P)$ has signature $(1,1)$ at the places in $I$
and is definite at the other real places.  Its adelic quotient
is compact because it is definite at $\mu$.  We discuss $s>0$;
the case $s=0$ is B24.  Choose the complex orientation of each
disk to be the one in which the actual Gaussian is holomorphic.
The opposite oscillator convention conjugates that disk and
every occurrence of its holomorphic line; it is not silently
identified with the chosen orientation.

At a place in $I$ the target $V_\tau$ is positive definite of
rank three.  On the source derived group the Gaussian has
lowest compact weight three.  Here is the direct Fock check.
For the compact torus of $\SU(1,1)$, the two source lines have
weights $e^{i\theta}$ and $e^{-i\theta}$.  In the positive
Lagrangian specified in 032, Section 4.6, printed p.22, the
second line is dualized.  Thus all six complex coordinates have
weight $e^{i\theta}$, and the half-determinant vacuum line has
weight $e^{3i\theta}$.  A unitary splitting changes only the
determinant character and is trivial on this derived torus.
The Gaussian is killed by the lowering root.  Equivalently,
the compact rank-three input with shifted parameter $(1,0,-1)$
has source parameter $(1;-1)+n_0/2$ in B18; its difference is
two and its holomorphic discrete-series weight is three.
The actual coefficient-vector convention may dualize both
weights, which is exactly the orientation convention just fixed.

At a definite source place different from $\mu$ the Gaussian
transforms by a determinant character.  At $\mu$ the degree-two
cochain uses the exterior square of the multiplicity factor;
this is again a determinant character of $\U(2)$.  Consequently
the complete source compact type is scalar.  Let $Y$ be its
compact unitary Shimura variety at a sufficiently small level,
with all the adelic components and the fixed central character
retained.  One may use a PEL similitude realization and then its
unitary components.  On a deeper principal component the
arithmetic determinant is one: the determinants of integral
unitary elements are norm-one algebraic units, hence roots of
unity, and a sufficiently small principal level kills them.
This does not remove their Hodge-theoretic determinant lines.

Let $\lambda_\tau$ denote the algebraic Hodge line whose
automorphy factor on the oriented upper half-plane is
$j(\gamma,z)=cz+d$.  The source line is
\[
 \mathscr L=
       \bigotimes_{\tau\in I}\lambda_\tau^{\otimes3}
                  \otimes\mathscr C.
 \tag{B26b}
\]
This equation defines $\mathscr C$ with the actual splitting
and compact characters: it is the quotient of the actual
one-dimensional compact type by the displayed weight-three
characters.  At each real place its remaining character is
an integral determinant power.  Indeed the full unitary
weights are integral after the specified unitary splitting,
and their difference at a switch is three.  Thus the quotient
is an algebraic character of the determinant torus.  Its
finite character part is retained as well.

The bundle $\mathscr C$ is the corresponding fixed CM
determinant line, with a possible finite character local
system.  On the finite component cover just described its
flat monodromy is trivial.  Its Hodge filtration and comparison
period need not be trivial.  Choose a unit flat analytic frame
$u_0$ and an algebraic de Rham frame $u_{\rm dR}$ of this line,
and record
\[
                 u_{\rm dR}=\Omega_{\mathscr C}u_0,
             \qquad \|u_{\rm dR}\|^2=|\Omega_{\mathscr C}|^2.
 \tag{B26c}
\]
If several labelled components are used, transport the frames
by their specified finite Hecke isogenies; their algebraic
transport factors are retained in the component coordinates.
Nothing sets $\Omega_{\mathscr C}$ to one or declares the line
to have Hodge weight zero.

The algebraicity of (B26b), its finite Hecke action and its
comparison at CM fibres can be read directly from the PEL
Hodge bundles.  In general canonical-bundle language the
precise source is Milne, \emph{Canonical models of mixed
Shimura varieties and automorphic vector bundles}, corrected
2018 version, Chapter III, Remark 2.3 and Example 2.4,
pp.55--56; Theorem 5.1 with proof and Remark 5.2, p.61;
Section 6, pp.61--62; and Section 7, pp.62--63.
These give algebraic automorphic bundles and their
$S(\A_{F,f})$-equivariant structures, not rational Betti
coordinates of arbitrary holomorphic sections.

Write the actual kernel in B25's algebraic coherent target
basis as
\[
                  K(g,h)=\sum_{a=1}^r f_a(g)e_a(h).
 \tag{B26d}
\]
The lowering-root equation for the source Gaussian survives
application of each target coefficient functional.  Therefore
each $e_a$ is the scalar automorphic realization of a
holomorphic section of $\mathscr L$.  This is a finite-dimensional
space because $Y$ is projective.  No finite-dimensionality
claim for arbitrary fixed-compact-type smooth functions is used.

\paragraph{One common coherent scalar on the actual source.}
Choose the source CM point $y$ defined by the original rational
decomposition $P=P_1\perp P_2$.  Its Shimura cocharacter lies
in the product of the two line tori.  At a switch it selects
the indicated positive or negative line; at a definite place
it has the prescribed central lift.  This establishes the CM
point for the full unitary datum, not merely a map from its
derived subgroup.

The finite adelic translates of $y$ separate
$H^0(Y,\mathscr L)$.  To see this, represent a section by its
scalar automorphic function $a(h_\infty,h_f)$.  If this is
nonzero at some point, choose $\gamma\in S(F)$ whose real
component is sufficiently close to $h_\infty$.  The Cayley
weak approximation argument of B25 gives such $\gamma$, and
automorphy gives
\[
       a(1,\gamma_f^{-1}h_f)=a(\gamma_\infty,h_f)\ne0.
 \tag{B26e}
\]
Because the section space is finite-dimensional, finitely
many such evaluations separate it.

Fix an algebraic basis $\eta_1,\ldots,\eta_D$ of
$H^0(Y,\mathscr L)$ and finitely many separating translates
$b_l$.  Transport an algebraic fibre frame from $y$ to the
CM point $[y,b_l]$ by the specified finite isogeny.  In the
marked analytic uniformization that isogeny has the identity
archimedean coordinate.  Hence the conversion from these
algebraic frames to the fixed automorphic frame has one
common nonzero scalar $t_y$, with any algebraic isogeny
factors put in the matrix.  The evaluation matrix is
\[
                   R_y=t_y B,
           \qquad B\in M_{L,D}(\overline\Q),\quad
                       \operatorname{rank}B=D.
 \tag{B26f}
\]
This is the same fibre-transport argument as B23, applied to
the actual line (B26b); it does not use CM of the cohomology
of $Y$.  It also retains all periods of $\mathscr C$ in $t_y$.

Include these finitely many source translates in the common
finite theta family of B19 and B25.  Assume its specified
finite Schwartz and splitting values are algebraic, exactly
as in B25.  The coefficient vector of
$K(\,\cdot\,,(1,b_l))$ in the basis $f_a$ is
$(C_x/t_x)$ times an algebraic vector.  Here $C_x$ is the
common CM theta comparison for this combined finite family,
and $t_x$ is the target coherent-fibre scalar.  Enlarging the
finite theta level is done once for this family, with its
isogeny factors retained.  Thus, for some algebraic matrix $D_0$,
\[
                    (e_a(1,b_l))_{l,a}
                         ={C_x\over t_x}D_0.
 \tag{B26g}
\]
Let $C$ be the coordinate matrix of the $e_a$ in the basis
$\eta$.  Taking an algebraic left inverse $L_B$ in (B26f)
gives the exact conclusion
\[
              C=\gamma C_0,\qquad
              \gamma={C_x\over t_x t_y},\qquad
                    C_0=L_BD_0\in M(\overline\Q).
 \tag{B26h}
\]
In particular the actual source coefficient space is defined
over $\overline\Q$ in coherent coordinates.  All its vectors
have a common scalar in the chosen coordinate list.  Neither
this statement nor Milne's CM-value criterion identifies
those coherent coordinates with Betti coordinates.

\paragraph{The exact cohomological source Gram.}
On each oriented disk use $z_\tau=x_\tau+iy_\tau$ and
$d\mu_\tau=dx_\tau dy_\tau/y_\tau^2$.  After removing the
constant line specified below, the local system is
\[
                     \mathbb V=
               \bigotimes_{\tau\in I}\mathbb V_\tau,
 \tag{B26i}
\]
where $\mathbb V_\tau$ is the standard rank-two weight-one
variation, in the complex embedding selected by the disk.
It is the corresponding factor of the first cohomology of
the source PEL abelian scheme.  The tensor can be made rational
by including its Galois conjugates; the displayed complex
summand is retained throughout the following calculation.
It has weight $s$.

Here the purity assertion has an explicit geometric
realization.  Let $\mathcal B\to Y$ be the source PEL
abelian scheme and $Z=\mathcal B\times_Y\cdots\times_Y
\mathcal B$ its $s$-fold fibre power.  Both $Y$ and $Z$
are smooth projective.  The indicated tensor local
system is a coefficient-field summand of
$R^s\pi_*\Q$, obtained by taking $R^1$ from each
factor and the prescribed $E$-embedding idempotents.
For an integer $N>1$, simultaneous multiplication
by $N$ on all fibres acts by $N^q$ on $R^q\pi_*\Q$.
All differentials in its Leray spectral sequence
therefore vanish: a differential of page $r\ge2$
would connect the distinct eigenvalues $N^q$ and
$N^{q-r+1}$.  In degree $2s$, the successive Leray
grades have distinct eigenvalues.  Hence a rational
polynomial in $[N]^*$ projects $H^{2s}(Z,\Q)$ onto
the grade $H^s(Y,R^s\pi_*\Q)$.  Indeed the product
of the distinct linear factors annihilates the
filtered vector space, so this operator is
semisimple and the usual Lagrange interpolation
polynomial is a projector.

Within that grade, multiplication by $N$ separately
on each fibre factor selects its $R^1$ summand by
the same interpolation argument, now on the
canonical fibre K\"unneth decomposition.  Finally
the $E$-endomorphism projectors select $\mathbb V$.
All these operators are algebraic graph
correspondences with the stated coefficient field.
Including the conjugate embedding projectors gives
the corresponding rational summand.  Thus the
cohomology used here is a prescribed summand of the
pure weight-$2s$ Hodge structure of $H^{2s}(Z)$.
The exact primary purity and polarization locators
are Deligne, \emph{Th\'eorie de Hodge II},
Section 2.2.1--2.2.6, printed pp.27--29.
Theorem 4.1.1(i), printed p.40, supplies the general
proper-smooth Leray degeneration statement; the
multiplication argument above proves the particular
splitting needed here directly.

There is a necessary determinant correction in this
cohomological interpretation.  The tangent bundle of the
line domain is $\operatorname{Hom}(\lambda_\tau,
\mathbb V_\tau/\lambda_\tau)$, so its canonical line is
$\lambda_\tau^2\otimes(\det\mathbb V_\tau)^{-1}$.
Consequently
\[
 \mathscr D=\mathscr C\otimes
                 \bigotimes_{\tau\in I}\det\mathbb V_\tau,
 \qquad
 \mathscr L=\Omega_Y^s\otimes F^s\mathbb V\otimes\mathscr D.
 \tag{B26i$'$}
\]
Here the second equality is the equivariant
Kodaira--Spencer identification, with its fixed
normalization.  The extra determinants are constant
CM Hodge lines.  Their selected embeddings have type
$(1,1)$, but their other conjugates need not make them
globally Tate.  Trivial determinant monodromy does not
discard these factors.  Use the product de Rham frame
of $\mathscr D$ and record its comparison with a unit
flat frame as $\Omega_{\mathscr D}$.  Thus
$\Omega_{\mathscr D}=\Omega_{\mathscr C}
\prod_{\tau\in I}\Omega_{\det\mathbb V_\tau}$ in the
chosen product frames, with every specified
Kodaira--Spencer frame factor retained.  In particular
any universal $2\pi i$ factors in these determinant
comparisons are not set equal to one.

On a component, write the actual coefficient in the form
\[
           e_a(g_z)=
              \Omega_{\mathscr D}F_a(z)
                          \prod_{\tau\in I}y_\tau^{3/2}
 \tag{B26j}
\]
in the fixed unitary automorphic frame.  This defines $F_a$;
no assertion that $F_a$ has algebraic Betti periods is being
made.  Put $v_\tau(z_\tau)=(z_\tau,1)$ and
\[
 \alpha_a=F_a(z)\,
       \bigotimes_{\tau\in I}v_\tau(z_\tau)
                   \;dz_1\wedge\cdots\wedge dz_s.
 \tag{B26k}
\]
For $\gamma_\tau\in\mathrm{SL}_2(\R)$ one has
$v_\tau(\gamma_\tau z)=
 \gamma_\tau v_\tau(z)/j(\gamma_\tau,z)$ and
$d(\gamma_\tau z)=dz/j(\gamma_\tau,z)^2$.
The weight-three transformation of $F_a$ therefore makes
(B26k) a well-defined $\mathbb V$-valued form.  It is closed
for the flat connection: its coefficients are holomorphic,
and every holomorphic derivative wedges with the already
present top product of the $dz_\tau$.
Its class has type $(2s,0)$ in the weight-$2s$ cohomology
$H^s(Y,\mathbb V)$, after the stated untwisting of
$\mathscr D$.  Restoring that line restores its actual CM
Hodge type and period; it is not a tacit Tate twist.

For the flat alternating form
$\varepsilon((a,b),(c,d))=ad-bc$ let
$B_{\mathbb V}=\bigotimes_{\tau\in I}\varepsilon_\tau$.
Define the Hermitian pairing, conjugate linear in the first
variable, on these classes by
\[
 Q(\alpha_a,\alpha_b)=
  (-1)^{s(s-1)/2}
     \int_Y B_{\mathbb V}(\alpha_b\wedge\overline{\alpha_a}).
 \tag{B26l}
\]
The convention in this display pairs coefficient tensors and
wedges the displayed base forms in that order.  Since
$\varepsilon(v,\bar v)=2iy$ and
$dz\wedge d\bar z=-2i\,dx\wedge dy$, moving the conjugate
base factors next to their partners gives exactly
\[
 Q(\alpha_a,\alpha_b)=
       4^s\int_Y\overline{F_a(z)}F_b(z)
                     \prod_{\tau\in I}y_\tau^3d\mu_\tau.
 \tag{B26m}
\]
In particular this form is positive definite on the
coefficient classes.  If a nonzero $\alpha_a$ were exact,
Stokes's theorem paired with its closed conjugate would make
the left side zero, contradicting the positive right side.
Thus no nonzero source coefficient is lost on passage to
this cohomology.

Let $\nu_S>0$ be the conversion from the specified adelic
Haar measure to the displayed disk measures and the fixed
compact-fibre and component measures.  Keep it as the
specified measure conversion, rather than renormalizing the
source quotient.  At one fixed source level, disintegration
gives the same finite-level volume and compact-fibre volume
on every adelic component.  If a deeper principal component
cover was used above, its contribution is divided by its
integer covering degree.  Thus the full disjoint-union
pairing in (B26l) includes these specified rational component
weights, and a single common positive conversion remains
in $\nu_S$.  The actual
conjugate-linear-first source Gram is
\[
 \widehat W_{ab}:=\int_{[S]}\overline{e_a(h)}e_b(h)\,dh
      ={\nu_S|\Omega_{\mathscr D}|^2\over4^s}
                                      Q(\alpha_a,\alpha_b).
 \tag{B26n}
\]
B17 uses $W_{ab}=\int e_a\overline{e_b}\,dh$ for its
coefficient matrix, so $W=\widehat W^{\mathsf t}$.
This transpose is retained when (B26n) is substituted in
the actual $TT^*$ matrix.

This gives a concrete possible CM criterion.  If the classes
of the normalized coherent source sections lie in one fixed
coefficient-embedding block of a full simple CM isotypic
summand, their algebraic de Rham classes
have a common Betti comparison period, just as in B23.
Their matrix in (B26n) is then one positive real scalar
times an algebraic Hermitian matrix.  Such a CM statement
is a statement about $H^s(Y,\mathbb V)$ with its recorded
constant CM line.  It does not follow from CM of the target
cup image.  In particular there has been no unproved
rational cohomological theta correspondence from the target
to this source group.

\paragraph{Algebraic Hecke adjoints reduce the remaining matrix.}
There is a separate reduction that needs no source CM
hypothesis.  Let $\mathcal H_{\overline\Q}$ be the finite-level
source Hecke algebra acting on $H^0(Y,\mathscr L)$, with the
actual central character imposed.  Normalize the compact
open subgroup to have finite Haar volume one in the
convolution notation.  Double-coset sums and their degrees
then have rational coefficients.  This convention does not
change the already fixed quotient measure in $\widehat W$.
The geometric Hecke correspondences and their bundle maps
are defined over $\overline\Q$, by the equivariant canonical
model above.  Thus they act algebraically in the coherent
basis.

The group is unimodular.  Change of variables in the
actual Petersson integral proves
\[
    \langle T_\phi u,v\rangle_S
        =\langle u,T_{\phi^*}v\rangle_S,
       \qquad \phi^*(h)=\overline{\phi(h^{-1})}.
 \tag{B26o}
\]
For characteristic functions this is the inverse double
coset, with the same convolution normalization.  Hence
the algebraic Hecke image is closed under the adjoint
involution of the actual positive Hermitian form.

Let $\mathcal F_{\overline\Q}$ be the Hecke saturation of
the normalized sections $\gamma^{-1}e_a$ from (B26h).
It is a finite-dimensional algebraic subspace of
$H^0(Y,\mathscr L)$.  Let $\mathcal A$ be the image of the
Hecke algebra on it.  Its complexification is a finite-dimensional
algebra of operators closed under the Hermitian adjoint.
Every invariant subspace therefore has an invariant
orthogonal complement.  Thus it is semisimple, and so is
$\mathcal A$ over $\overline\Q$.  Choose an algebraic
Wedderburn decomposition of this actual representation:
\[
 \mathcal F_{\overline\Q}
    =\bigoplus_\lambda M_\lambda\otimes V_\lambda,
 \qquad
 \mathcal A=\bigoplus_\lambda\operatorname{End}(V_\lambda),
 \qquad m_\lambda=\dim M_\lambda.
 \tag{B26p}
\]
The central block idempotents are orthogonal projections
for the actual Hermitian form.  This follows either from
the invariant-complement assertion or from the fact that
a central idempotent in a finite-dimensional adjoint-closed
operator algebra is self-adjoint.

For each $V_\lambda$ there is a positive definite algebraic
Hermitian matrix $G_\lambda$ implementing the given
involution.  To justify algebraicity, choose a finite
algebraic spanning set of the matrix algebra.  The equations
\[
        G_\lambda\rho(a)=\rho(a^*)^*G_\lambda
 \tag{B26q}
\]
are a finite real linear system with real algebraic
coefficients on the Hermitian entries.  A positive real
solution exists by restriction of the actual form to an
irreducible copy.  Algebraic points are dense in this
linear solution space, and positivity is open, so a
positive algebraic solution exists.  Algebraic changes
of basis, including an algebraic Cholesky factor, are
therefore allowed here.

Schur's lemma now gives the exact shape of the actual
Petersson Gram in (B26p):
\[
                   G_S=\bigoplus_\lambda
                                  H_\lambda\otimes G_\lambda,
 \tag{B26r}
\]
where each $H_\lambda$ is a positive definite Hermitian
matrix on the multiplicity space $M_\lambda$.  These
$H_\lambda$ are actual source period matrices; no
algebraicity is asserted for them.  For completeness,
compose the actual form on one block with the inverse
of $I\otimes G_\lambda$.  Invariance makes the resulting
operator commute with $I\otimes\operatorname{End}(V_\lambda)$,
so it has the form $H_\lambda\otimes I$.  Positivity and
Hermitian symmetry give the stated properties of
$H_\lambda$.

Let $C_1$ be the algebraic matrix of the normalized
sections $\gamma^{-1}e_a$ in (B26p).  Their actual Gram
is therefore
\[
       \widehat W=|\gamma|^2 C_1^*
          \left(\bigoplus_\lambda H_\lambda\otimes G_\lambda\right)
                                  C_1.
 \tag{B26s}
\]
This is a reduction of the Gram of the original fixed
kernel, not a replacement by norms of a different frame.

If its Hecke saturation is one irreducible module, then
there is one index $\lambda$ and $m_\lambda=1$.
Equation (B26s) becomes
\[
                 \widehat W=c\,C_1^*G_\lambda C_1,
                         \qquad c>0.
 \tag{B26t}
\]
The same conclusion holds if the coefficient vectors
occupy a single algebraic line in a single multiplicity
space.  In these cases the actual $TT^*$ inverse of B17
is one positive scalar times an algebraic matrix on
the preserved algebraic image of B25.  B23 realizes that
matrix by an algebraic CM-source correspondence.  Thus
the actual fixed-kernel correction is closed in this
stated Hecke sector, without proving the source motive
to be CM.

If several multiplicity-one sectors occur, each has its
own positive scalar; their ratios have not been shown
algebraic.  If $m_\lambda>1$, the remaining matrix has
$m_\lambda^2$ real Hermitian parameters before any further
relations.  The CM criterion after (B26n) is one way to
control these quantities, but it is an additional
source-motive assertion.

\paragraph{A finite sufficient certificate.}
The irreducible-sector condition can be certified without
assuming a global multiplicity-one theorem.  It suffices
to exhibit a finite-dimensional algebraic coherent
subspace containing all $\gamma^{-1}e_a$ and a finite
list of algebraic geometric Hecke operators and their
adjoints preserving it, whose matrices generate its full
matrix algebra.  Generation is a finite linear-algebra
test on products of these matrices; once their span has
dimension $D^2$, it is the full matrix algebra on a
$D$-dimensional space.  The proof of (B26q)--(B26t)
then applies using just that finite list.  The necessary
coherent matrices can be recovered from finite CM-value
rows with the same common-scalar argument (B26f)--(B26h),
retaining the original finite theta and isogeny data.
No such full-matrix certificate is claimed here for an
unspecified actual 032 input.

\paragraph{Source scope.}
The Fock and real representation inputs are the exact
032 Section 4.6 positive Lagrangians and B18's checked
Atobe Definition 1.6/Theorem 1.7 and Ichino Theorem 4.1
specialization.  The finite CM-value comparison is B19
and B25 with its original theta-group, tangent, splitting
and Haar normalizations.  Milne's Chapter III supplies
the algebraic coherent bundle and Hecke structures.
The cohomological factor $4^s$, its sign and the entire
finite-dimensional adjoint-algebra reduction have been
proved here.  None of these statements identifies a sum
of sectors or a changed finite kernel with one prescribed
arithmetic \texttt{MainX} tuple.


% THETA3-B27 contribution by /root/definite_bridge, 21 November 2026.
\subsection{THETA3-B27: the actual mixed source has multiplicity-free Hecke sectors}\label{B27}

For the actual rank-two source of B26, the Hermitian multiplicity
matrices in (B26r) can be reduced further: every multiplicity is one.
Consequently the Gram of the original fixed kernel depends on finitely
many positive scalar periods, one for each source Hecke sector.
This assertion does not require the source motive to be CM and does
not assert that only one sector occurs.  The argument uses only the
generic pure-inner-form case of the unitary multiplicity formula,
after excluding the nongeneric rank-two parameter directly.

\paragraph{The fixed archimedean representation.}
Retain $S=\U(P)$, the nonempty switch set $I$, the complex orientations,
the full central character, and the algebraic line $\mathscr L$ of B26.
At each $\tau\in I$ the source representation generated by a nonzero
holomorphic coefficient vector restricts to the holomorphic
discrete series of $\SU(1,1)$ of lowest compact weight three.
At every other real place its complete compact type is the fixed
one-dimensional determinant character.  In particular the entire
archimedean representation, denoted $\sigma_\infty$, is fixed, and
its specified holomorphic compact type has dimension one.

Here this conclusion includes the lowering-operator condition,
not merely a coincidence of compact weights.  The actual Fock
vacuum in B26 is killed by the lowering root.  A unitary
lowest-weight $\mathfrak{sl}_2$ module of lowest weight three is the
discrete series of that weight: successive raising vectors have
weights $3+2j$, and the commutator relations and positivity determine
the module and its invariant inner product up to a scalar.  Thus
the holomorphic sections in question cannot contribute a different
unitary archimedean representation with the same isolated compact
weight.  The fixed central character determines the extension
from $\SU(1,1)$ to $\U(1,1)$, since
$\U(1,1)=(\SU(1,1)\times\U(1))/\{\pm1\}$.
At a compact place the compact group is the entire real group,
so its specified scalar type already determines the representation.
The simultaneous contragredient convention replaces all these
holomorphic choices by their prescribed duals and leaves the
argument unchanged.

Because $P$ is definite at $\mu$, it is anisotropic over $F$.
Its unitary group and its center have compact adelic quotients.
Thus the source functions belong to the discrete automorphic
spectrum, with no continuous-spectrum issue.

\paragraph{A derived-invariance lemma at one split finite place.}
Let $D=\SU(P)$.  This is a simply connected absolutely almost
simple $F$-group of type $A_1$.  Let $v$ be a finite place split
in $E$, so $D(F_v)\simeq\mathrm{SL}_2(F_v)$ is noncompact.
Strong approximation gives
\[
       \overline{D(F)}=D(\A_F^{\{v\}}).
 \tag{B27a}
\]
The notation on the right includes every archimedean place.
This is precisely Platonov--Rapinchuk, \emph{Algebraic Groups
and Number Theory}, Theorem 7.12, printed p.427; the final part
of its proof on pp.431--432 removes the condition that the set
of omitted places contain the archimedean ones.  Their theorem
therefore applies with the singleton set $\{v\}$.

Suppose a continuous automorphic function $f$ on $S(F)\backslash
S(\A_F)$ is right invariant under $D(F_v)$.  Then it is right
invariant under all of $D(\A_F)$.  To prove this at an arbitrary
base point $g=(g^v,g_v)$, take $\gamma\in D(F)$.  Normality of
$D(F_v)$ in $S(F_v)$ and the assumed right invariance give
\[
 f(\gamma^v g^v,g_v)
   =f(\gamma^v g^v,\gamma_v g_v)=f(g).
 \tag{B27b}
\]
For any $d^v\in D(\A_F^{\{v\}})$, choose such rational
$\gamma$ converging to $d^v$ away from $v$ and apply continuity.
This proves left $D(\A_F^{\{v\}})$ invariance.  At $v$ the
assumed right invariance is also left invariance by normality.
Thus $f$ is left invariant under $D(\A_F)$, and normality once
again makes this equivalent to right invariance.  No condition
on the possibly unbounded $v$-components of the approximating
rational elements is needed.

\paragraph{Exclusion of the only nongeneric rank-two parameter.}
Use the formal global parameters and near-equivalence decomposition
of Kaletha--M\'inguez--Shin--White, \emph{Endoscopic classification
of representations: inner forms of unitary groups I},
Section 1.3 and Corollary 3.3.2, printed pp.63--70 and 151.
The latter decomposes the discrete spectrum according to
parameters using their unramified Satake classes; it is available
before the full local classification of nongeneric inner-form
packets.  In rank two the only possible nongeneric parameter is
\[
                         \psi=\chi\boxtimes S_2,
 \tag{B27c}
\]
where $\chi$ is a conjugate-self-dual automorphic character of
$\A_E^\times/E^\times$ and $S_2$ is the two-dimensional
irreducible representation of the Arthur $\mathrm{SU}(2)$.
Indeed the parameter is a sum of blocks of dimensions
$m_i n_i$; a nontrivial Arthur factor has $n_i\ge2$ and
already exhausts the rank two.  Parameters which are sums
of rank-one blocks with trivial Arthur factors are generic,
even when the sum is not discrete.

Choose a finite place $v$ split in $E$ at which the automorphic
representation, the parameter, and the fixed base-change splitting
character are unramified.  Such places exist outside any finite
set.  Here $S(F_v)\simeq\mathrm{GL}_2(F_v)$, and (B27c)
has normalized Satake eigenvalues
\[
              \alpha_v q_v^{1/2},\qquad
              \alpha_v q_v^{-1/2}
 \tag{B27d}
\]
for a nonzero scalar $\alpha_v$.  All fixed splitting-character
factors multiply both eigenvalues by the same scalar.  KMSW's
formula (1.3.2), printed p.64, and the localization in
Section 1.3.3 give exactly these exponents.  The unique spherical
irreducible representation with these eigenvalues is an
unramified character composed with the determinant.  This can
also be checked without a nonsplit packet theorem: the trivial
representation of $\mathrm{GL}_2(F_v)$ is the spherical
constituent of normalized induction of
$|\cdot|^{1/2}\otimes|\cdot|^{-1/2}$, and a determinant twist
multiplies both Satake eigenvalues by $\alpha_v$.

Consequently the local representation is trivial on $D(F_v)$.
The lemma (B27a)--(B27b) makes its automorphic realization
trivial on $D(\A_F)$.  This contradicts the lowest weight three
at any place in $I$.  Hence no actual source constituent has
parameter (B27c): every constituent is generic.

This exclusion uses a split finite place and strong approximation.
It does not infer genericity merely from an unproved compatibility
of the archimedean infinitesimal character with a nongeneric
inner-form Arthur packet.

\paragraph{The precise multiplicity formula that now applies.}
The group $S$ is a pure inner form of the quasi-split unitary
group: it is defined by a Hermitian space over $E$, rather than
by a Hermitian space over a nontrivial division algebra.
KMSW Section 0.3.3, printed p.18, proves exactly this criterion
by the relevant Galois cohomology sequence and the Brauer class.
The generic pure-inner-form case of their Theorem 5.0.5,
printed pp.205--206, is therefore applicable to every actual
source constituent.  Its two local hypotheses are expressly
proved for generic parameters and pure inner twists in the
preceding paragraph.  The theorem gives zero contribution for
a nondiscrete generic parameter and, for a discrete generic
parameter, a direct sum in which each permitted automorphic
representation occurs once.

We use this proved case, not the unrestricted starred statement
of Theorem 1.7.1.  The scope paragraph after that starred
statement, printed p.89, explicitly distinguishes the proved
generic pure-inner case from the other cases reserved for
sequels.  We also make no assertion that the local components
at every finite place are tempered; their generic global
parameter suffices here.

It follows that every irreducible automorphic representation
with the actual infinity type $\sigma_\infty$ occurs with
multiplicity one.  At a fixed finite compact open $K_S$, the
full finite-level Hecke module of holomorphic sections is
therefore multiplicity free as well:
\[
          H^0(Y,\mathscr L)
              =\bigoplus_{\sigma}
                    \sigma_f^{K_S},
 \tag{B27e}
\]
where zero summands are omitted and the equality refers to the
fixed full adelic realization and central character of B26.
Every summand on the right is a simple module for the full
$K_S$-bi-invariant Hecke algebra.  Two different representations
on the right cannot yield isomorphic nonzero Hecke modules.

For completeness, these finite-Hecke assertions are elementary
idempotent facts.  If $e_{K_S}$ is normalized compact averaging
and $V$ is an irreducible smooth representation with
$e_{K_S}V\ne0$, a nonzero invariant subspace of $e_{K_S}V$
generates $V$ under the group, and averaging that generated
space back into $e_{K_S}V$ proves its simplicity.  Conversely
a simple $e_{K_S}\mathcal H e_{K_S}$-module determines the unique
simple quotient generated by it: in the induced module, quotient
by the largest submodule annihilated by all operators
$e_{K_S}\mathcal H$; any further nonzero submodule contains the
generating $e_{K_S}$-module.  Thus the finite representation is
determined by its nonzero fixed-vector Hecke module.  Finally,
two global representations with the same finite representation
and the already fixed $\sigma_\infty$ are the same restricted
tensor product.  The selected archimedean type contributes
one dimension, not another multiplicity.

\paragraph{The original fixed-kernel Gram now has only scalar sectors.}
Let $\mathcal F_{\overline\Q}$ be the actual algebraic Hecke
saturation constructed in B26.  It is a submodule of (B27e),
so every $m_\lambda$ in (B26p) is one.  In the algebraic
decomposition and with the same positive algebraic matrices
$G_\lambda$ as in (B26q), there are actual positive real
numbers $c_\lambda$ such that
\[
 \begin{aligned}
  \mathcal F_{\overline\Q}&=\bigoplus_\lambda V_\lambda,
       &G_S&=\bigoplus_\lambda c_\lambda G_\lambda,\\
  \widehat W&=|\gamma|^2
       \sum_\lambda c_\lambda C_\lambda^*G_\lambda C_\lambda,
       &\gamma&=\frac{C_x}{t_x t_y}.
 \end{aligned}
 \tag{B27f}
\]
Here $C_\lambda$ is the algebraic coefficient matrix of
the normalized actual source coefficients in $V_\lambda$.
Every Haar, theta-group, tangent, and constant CM-line
factor in $\gamma$ and the source periods is retained from
B26.  This is the Gram of the original kernel.  Its remaining
freedom is at most one positive scalar per sector, rather
than an arbitrary positive Hermitian matrix on each
multiplicity space.

The central idempotents $p_\lambda$ of the finite algebraic
Hecke image are algebraic and are orthogonal projections for
the actual Petersson form.  Each can be represented by a
finite algebraic linear combination of finite double-coset
operators.  Apply that same operator to the source variable
of the original kernel:
\[
 K_\lambda(g,h)=p_\lambda^h K(g,h),\qquad
 K=\sum_\lambda K_\lambda.
 \tag{B27g}
\]
These are genuine theta kernels on the same $P,V$ with
the same archimedean cochain.  Indeed
$K_{\omega(s)\Phi}(g,h)=K_\Phi(g,hs)$ for a finite source
element $s$, and convolution by a finite double-coset
operator is the corresponding finite algebraic combination
of finite Weil translates.  Its coefficients are algebraic
in the finite rational Weil structure used in B25, with
the same splitting-character values retained.  No direct
sum of unrelated Hermitian spaces is introduced.

Let $T_\lambda$ be the bilinear theta integral with kernel
$K_\lambda$.  Orthogonality in the source, with exactly
the Hilbert-space convention of B17, gives
\[
 T_\lambda T_\kappa^*=0\quad(\lambda\ne\kappa),
       \qquad TT^*=\sum_\lambda T_\lambda T_\lambda^*.
 \tag{B27h}
\]
The first identity is the integral of a source coefficient
in one sector against the conjugate of a coefficient in
another.  In particular it is valid for the bilinear
theta-integral convention; no unnoticed inversion of a
source character is being used.

There is also an explicit matrix form.  Use B25's algebraic
target basis, with target Gram
$G_T=\nu_G|P_\chi|^2G_B$, and put
\[
 D_\lambda=(C_\lambda^*G_\lambda C_\lambda)^{\mathsf t},
       \qquad A_\lambda=D_\lambda G_B.
 \tag{B27i}
\]
The transpose is the one required by B17's
$W_{ab}=\int e_a\overline{e_b}$.  Thus
\[
 T_\lambda T_\lambda^*
   =\nu_G|P_\chi|^2|\gamma|^2c_\lambda A_\lambda.
 \tag{B27j}
\]
The matrix $A_\lambda$ is algebraic and is positive
semidefinite and self-adjoint for $G_B$.  Its image
$E_\lambda$ is algebraic.  On that image it is positive
definite, so its inverse is algebraic.  Equivalently,
extend that inverse by zero on the $G_B$-orthogonal
complement.  B23 realizes this algebraic matrix by a
correspondence from the finite CM source.  Therefore the
actual B17 inverse in a single nonzero sector is
\[
 (T_\lambda T_\lambda^*|_{E_\lambda})^{-1}
   =\frac{1}{\nu_G|P_\chi|^2|\gamma|^2c_\lambda}
                          (A_\lambda|_{E_\lambda})^{-1}.
 \tag{B27k}
\]
This closes its matrix-valued algebraicity issue within
each actual sector.  The one remaining scalar is its
Petersson period with the already specified normalization.

For an original nonzero lift $Tu$, the decomposition
$Tu=\sum_\lambda T_\lambda u$ has at least one nonzero
summand.  There is no assertion that the original lift
itself lies in only one $E_\lambda$, or that the target
images for different sectors are orthogonal or disjoint.
The original full $TT^*$ remains the weighted sum in
(B27h)--(B27j).  Its inverse is projectively algebraic
if the relevant $c_\lambda$ have algebraic ratios;
that comparison has not been proved here.  In particular
this result does not identify a single prescribed
\texttt{MainX} tuple or supply its arithmetic constants.

\paragraph{Primary-source and proof scope.}
The actual source weights, line, and common coherent
scalar are B26, with the exact 032 and local-theta
sources recorded there.  For the new global input the
whole cited KMSW statements and their immediate proofs
were read: Section 0.3.3, Section 1.3.2--1.3.4, Corollary
3.3.2 and its decomposition, the scope paragraph after
Theorem 1.7.1, and all of Theorem 5.0.5 and its proof.
The strong approximation source is the original authors'
Theorem 7.12 with its complete proof, including the
singleton finite-place scope.  Its application (B27b),
the rank-two exclusion, the fixed-infinity Hecke
multiplicity argument, and the displayed kernel and
matrix consequences were proved above.  No unproved
nongeneric inner-form packet theorem is needed.


\subsection{THETA3-B28: an explicit bound for the actual inverse from source CM evaluations}\label{B28}

The remaining source periods in B26 need not be determined in
order to bound the norm of the actual B17 inverse.  Its finite
CM evaluation matrix gives a lower bound for the source Gram.
For the weight-three source the local constant is explicit,
and the auxiliary source-fibre period cancels.  The statement
is a fixed-input certificate: the CM value matrix, source chart
radii and original measures remain specified inputs.

\paragraph{Data and the exact bound.}
Use B26's source $Y$ of complex dimension $s>0$ and its actual
line $\mathscr L$.  Fix the algebraic coherent basis
$f_1,\ldots,f_r$ of the actual image $E_0$ from B25, and write
\[
 K(g,h)=\sum_{a=1}^r f_a(g)e_a(h),\qquad
 G_{ab}=\langle f_a,f_b\rangle_G,\qquad
 W_{ab}=\int_{[S]}e_a(h)\overline{e_b(h)}\,dh.
 \tag{B28a}
\]
Inner products are conjugate linear in the first variable.
Thus $G$ and $W$ are positive Hermitian matrices, and the
coordinate matrix of the actual $Q_T=TT^*$ is $WG$, as in B17.
No finite restriction operator is substituted for $Q_T$.

Choose distinct source CM points $y_l=[y,b_l]$ that separate
the coefficient sections $e_a$, using the finite Hecke orbit
of the fixed CM point as in B26.  Enlarge the common finite
theta level once for these finitely many points.  Their actual
evaluation matrix has the form
\[
       E_{la}=e_a(1,b_l)={C_x\over t_x}D_{la},
       \qquad D\in M_{L,r}(\overline\Q),\quad
                     \operatorname{rank}D=r.
 \tag{B28b}
\]
The scalar and all finite isogeny factors are those of B25--B26.
In particular $C_x$ and $t_x$ have not been renormalized to one.

At each $y_l$ choose an injective product of oriented disk
charts $|w_j|<\rho_{lj}<1$ whose images in $Y$ are pairwise
disjoint.  After a deeper neat component cover if necessary,
all charts are honest manifold charts.  Its covering degree
is included in the measure below.  Let $\nu_l>0$ be the
conversion from the original adelic source measure to
$\prod_j dx_jdy_j/y_j^2$ on that chart, including its compact
fibre and component weights.  Put
\[
 \kappa_l=\nu_l(4\pi)^s
           \prod_{j=1}^s\left(\rho_{lj}^2-
                                      {\rho_{lj}^4\over2}\right),
 \qquad H=D^*\operatorname{diag}(\kappa_l)D,
 \qquad R=H^{\mathsf t}.
 \tag{B28c}
\]
Both $H$ and $R$ are positive definite.  Then the actual
source Gram and its inverse satisfy the Loewner inequalities
\[
    W\ \geq\ \left|{C_x\over t_x}\right|^2R,
    \qquad
    W^{-1}\ \leq\ \left|{t_x\over C_x}\right|^2R^{-1}.
 \tag{B28d}
\]
For $f=\sum_a q_af_a$, let
$\phi_f=T^*Q_T^{-1}f$.  This is the source vector of least
Hilbert norm among all vectors mapping to $f$.  Its exact
norm and the resulting certificate are
\[
 \|\phi_f\|_S^2=q^*W^{-1}q
       \ \leq\ \left|{t_x\over C_x}\right|^2q^*R^{-1}q.
 \tag{B28e}
\]
Likewise the operator norm on the target with Gram $G$ obeys
\[
 \|Q_T^{-1}\|_{G\to G}\ \leq\
 { |t_x/C_x|^2\over
           \lambda_{\min}(G^{1/2}RG^{1/2})}.
 \tag{B28f}
\]
The inequalities are for the original kernel and its original
measures.  They do not assert that its inverse is algebraic.

\paragraph{Proof of the finite evaluation and the local constant.}
The sections $e_a$ are linearly independent.  Otherwise a
nonzero linear functional on the basis $f_a$ would annihilate
$K(\,\cdot\,,h)$ for every $h$, contradicting the definition
of $E_0$ as the span of these vectors.  B26 proves that finite
translates of $y$ separate the whole finite-dimensional
holomorphic source space, hence a finite list separates the
$e_a$.  Repeated geometric points give proportional evaluation
functionals in a line bundle and may be removed.  The exact
common-scalar formula is B26g; its matrix therefore has rank
$r$.  A finite collection of distinct points on the smooth
quotient admits disjoint sufficiently small product charts.
All their positive radii may be chosen rational.

We compute the chart constant without any statement about
unknown Betti periods.  On one disk first move its centre to
$z=i$ by the indicated real unitary element.  Use the Cayley
coordinate
\[
 z=i{1+w\over1-w},\qquad
 y={1-|w|^2\over|1-w|^2},\qquad
 {dx\,dy\over y^2}={4\,du\,dv\over(1-|w|^2)^2}.
 \tag{B28g}
\]
For a weight-three coefficient write its holomorphic scalar
as $F(z)=(1-w)^3A(w)$.  Any fixed determinant-line comparison
factor is included in $A$; its value at zero is the actual
unitary coefficient at the centre.  The Petersson density is
therefore exactly
\[
       |F(z)|^2y^3{dx\,dy\over y^2}
               =4|A(w)|^2(1-|w|^2)\,du\,dv.
 \tag{B28h}
\]
For several disks the same formula is the product, with
$A$ holomorphic on the polydisk and $A(0)=e(y_l)$ for the
section under consideration.  The central and determinant
characters are constant along these derived-group disks;
their factors are included in this value rather than omitted.

Expand $A(w)=\sum_{\mathbf n\geq0}a_{\mathbf n}w^{\mathbf n}$.
Angular integration kills all mixed monomials.  Every
remaining term in the weighted integral is nonnegative.
The constant term gives
\[
 \int_{|w_j|<\rho_j}|A(w)|^2
                \prod_j(1-|w_j|^2)\,\prod_jdu_jdv_j
 \ \geq\ |A(0)|^2\pi^s
                  \prod_j(\rho_j^2-\rho_j^4/2).
 \tag{B28i}
\]
One may first integrate on smaller closed polydisks where
the power series converges uniformly, and then increase
their radii; positivity and monotone convergence justify
the displayed open-disk formula.  Multiplying by $4^s\nu_l$
and summing the disjoint charts proves
\[
       \|e\|_S^2\ \geq\ \sum_l\kappa_l|e(y_l)|^2.
 \tag{B28j}
\]
This is the weight-three source version of Q11's geometric
evaluation argument, with its sharper radial integral
constant computed explicitly here.

Apply (B28j) to $e=\sum_a u_ae_a$.  Its integral Gram is
$\widehat W=\overline W=W^{\mathsf t}$ in this convention,
so (B28b) gives
$\widehat W\geq|C_x/t_x|^2H$.
Transposition preserves positivity of Hermitian matrices,
and hence gives the first inequality in (B28d).  Conjugating
by $R^{-1/2}$ and inverting its positive eigenvalues gives
the second.  No diagonalization of the unknown source
period matrices was used.

The adjoint formula in B17 gives $Q_T=TT^*$ and its
invertibility on $E_0$.  Therefore $T\phi_f=f$ and
$\phi_f\in\operatorname{im}T^*=(\ker T)^\perp$.
Every other preimage differs by a vector in $\ker T$,
so orthogonality proves the least-norm assertion.  Also
\[
 \|\phi_f\|_S^2
   =\langle Q_T^{-1}f,f\rangle_G
   =q^*G(WG)^{-1}q=q^*W^{-1}q.
\]
This proves (B28e).  Under the isometry given by $G^{1/2}$,
the positive operator $Q_T$ is represented by
$G^{1/2}WG^{1/2}$.  The first inequality of (B28d)
then proves (B28f).\hfill$\square$

\paragraph{A scalar and arithmetic form of the certificate.}
Set $\kappa=\min_l\kappa_l$.  From (B28c),
$\lambda_{\min}(R)\geq\kappa\sigma_{\min}(D)^2$.
Thus the right side of (B28e) is at most
\[
       {|t_x/C_x|^2\over\kappa\sigma_{\min}(D)^2}
                          \sum_a|q_a|^2.
 \tag{B28k}
\]
If $M$ is any nonsingular $r$-row square minor of $D$,
then
\[
 \sigma_{\min}(D)\geq\sigma_{\min}(M)
              \geq { |\det M|\over\|M\|_F^{r-1}}.
 \tag{B28l}
\]
For $r=1$ the denominator is interpreted as one.  The
first assertion follows by adding the other nonnegative
row squares; the second follows by bounding each of the
other $r-1$ singular values by $\|M\|_F$ and using their
product $|\det M|$.

The algebraic determinant estimate of B7 now applies
literally.  For example let $K$ contain the entries of $M$,
$[K:\Q]=D_K$, and let the positive integer $d_0$ make
$d_0M$ integral.  If
$|\sigma(M_{ab})|\leq H_\sigma$ for every other embedding,
the norm of the nonzero algebraic integer
$d_0^r\det M$, together with Hadamard's determinant bound,
gives
\[
 |\det M|\ \geq\ d_0^{-rD_K}
       \prod_{\sigma\ne\mathrm{id}}
                      (r^{r/2}H_\sigma^r)^{-1}.
 \tag{B28m}
\]
This records the required denominator and conjugate bounds;
it does not supply uniform estimates for unspecified input
data.  Source chart radii and the original measure conversion
are likewise not assigned numerical values without an input.

In B26's coherent-source notation the coefficient scalar is
$\gamma=C_x/(t_xt_y)$ and a source CM evaluation brings back
$t_y$.  Formula (B28b) makes their cancellation exact.  The
new bound thus needs no estimate for the extra source fibre
period $t_y$ or the source multiplicity-period matrices.
The retained target theta scalar $C_x/t_x$ can be treated
using the prescribed reference normalization of B22 when
its stated data have been matched.  The algebraic matrices,
integral-frame losses and the exact archive normalization
still have to be supplied.  Nothing here proves the missing
finite-transform derivative rank of A13--A14.

\framework{This is a bound on a specified realization
functional and its genuine adjoint inverse.  The comparison
with algebraic correspondences is the framework of
Grothendieck cote~44, scans~7--8, and cote~16, scans~43--46.
B23--B27 supply the stated algebraic sectors separately;
the present analytic inequality does not itself construct
a new motivic projector.}


% THETA3-B29 contribution by /root/definite_bridge, 21 November 2026.
\subsection{THETA3-B29: actual theta sectors are orthogonal on the target}\label{B29}

Local Howe duality supplies an additional fact not used in B27:
the target images of its distinct source sectors are orthogonal.
Thus the inverse of the original fixed theta kernel is a finite
sum of positive-scalar multiples of algebraic sector inverses.
The scalar source periods remain.  This assertion concerns the
actual scalar oscillator coefficient of B17 and the actual
source sectors of B27, with the same spaces, splittings and
Haar measures throughout.

\paragraph{The local uniqueness theorem and the dual convention.}
Let $\sigma_\lambda$ denote the irreducible automorphic source
representation containing the coefficient functions in the
$\lambda$-sector of B27.  These are the representations of
$e_a(h)$, not of the functions integrated bilinearly against
them.  The latter lie in $\sigma_\lambda^\vee$.
For an irreducible target representation $\Pi$, the relevant
local occurrence is
\[
 \operatorname{Hom}_{G(F_v)\times S(F_v)}
             (\omega_v,\Pi_v\otimes\sigma_{\lambda,v})\ne0.
 \tag{B29a}
\]
In the convention where the local theta lift is the maximal
irreducible quotient of the other factor in
$\omega_v\twoheadrightarrow\sigma_{\lambda,v}\otimes
\Theta(\sigma_{\lambda,v})$, this says
$\Pi_v\simeq\theta(\sigma_{\lambda,v})$.
There is no contragredient on $\sigma_{\lambda,v}$ in
(B29a); its contragredient is the global bilinear input.

For a nonsplit finite place, Gan--Takeda,
\emph{A proof of the Howe duality conjecture}, Theorem 1.2,
printed pp.1--3, proves that the small theta lift is
irreducible or zero and that two nonzero small lifts agree
only when their inputs agree.  Its local field has
characteristic different from two and arbitrary residue
characteristic.  It therefore includes the ramified and
residue-characteristic-two completions needed here.
The spaces are Hermitian spaces over a quadratic field,
not quaternionic Hermitian spaces, so the separate
quaternionic qualification in that paper is irrelevant.

At a split finite place the groups form a type-II general
linear dual pair.  M\'inguez,
\emph{Correspondance de Howe explicite: paires duales de
type II}, Corollary 6.3(1), printed p.734, gives the
uniqueness of the other irreducible representation in
a nonzero quotient of the local oscillator representation.
Apply it also with the two members exchanged.  This is
the required injectivity for type II.  Fixed determinant
twists, the fixed inverse-transpose identification, and
the fixed nontrivial additive character of the actual
split model are those retained in B18.  They only apply
fixed bijections to the representation labels and do
not change uniqueness.

Only finite-place uniqueness will be needed below.
B27 already fixes the complete source infinity type,
so distinct source sectors differ at a finite place.
In particular the argument does not require a new
archimedean packet calculation or a separate global
irreducibility theorem for a theta lift.  It uses
irreducible quotients of the big local theta lift;
it never asserts that every big local theta lift is
itself irreducible.

\paragraph{A nonzero global constituent gives (B29a).}
The canonical cochain causes no additional representation
factor here.  B17 fixed a unit canonical frame $\kappa$
and the actual scalar oscillator vector
$\Phi=\Psi(\kappa)$, with
\[
   K(g,h)=\theta(\omega(g,h)\Phi).
 \tag{B29b}
\]
The archimedean coefficient pairing has already been
performed in this equality.  Applying a finite source
Hecke projector changes only the finite Schwartz input,
so $K_\lambda$ is still a vector in the same global
oscillator-to-automorphic theta map.  In particular it
is not merely a cohomology class compared to such a map.

Suppose a target irreducible constituent $\Pi$ occurs
in the automorphic representation generated by
$E_\lambda=\operatorname{im}T_\lambda$.
Equivalently the projection of the theta kernel to
the $\Pi$-isotypic target component and the
$\sigma_\lambda$-isotypic source component is nonzero.
Both quotients are compact, so all theta pairings used
to detect this projection converge absolutely.
Orthogonal automorphic projections commute with the
group actions.  Hence the theta map, followed by these
projections and by a choice of a nonzero irreducible
copy, gives a nonzero equivariant map from the global
oscillator representation to $\Pi\otimes\sigma_\lambda$.

Here is an explicit localization argument.  Fix a
finite place $v$.  Decompose the global Schwartz space
as the restricted tensor product with its $v$-factor.
A nonzero projected vector is detected on one finite
sum of pure tensors, hence on one choice of the
Schwartz vector away from $v$.  Pair the representation
factors away from $v$ with a suitable smooth dual
vector which detects a nonzero coefficient.  The
result is a nonzero $G(F_v)\times S(F_v)$-equivariant
map from $\omega_v$ to $\Pi_v\otimes
\sigma_{\lambda,v}$.  Admissibility ensures that a
vector fixed by an open compact subgroup lies in a
finite-dimensional fixed-vector tensor product, so
this is the usual smooth local Hom space.  This
proves (B29a) without a factorization formula for
the value of a global period.

The same convention can be checked directly by
integrating against $u\in\sigma_\lambda^\vee$ and
$v\in\Pi^\vee$:
\[
 \int_{[G]\times[S]}\theta_\Phi(g,h)u(h)v(g)\,dg\,dh.
 \tag{B29c}
\]
Changing variables after replacing $\Phi$ by
$\omega(g_0,h_0)\Phi$ moves $u,v$ by the inverse
right translations.  Thus this is an invariant
trilinear functional on
$\omega\otimes\sigma_\lambda^\vee\otimes\Pi^\vee$,
which is exactly (B29a), not its mislabeled dual.

\paragraph{Orthogonal target images and the actual full inverse.}
If one target constituent $\Pi$ occurred for both
$\lambda$ and $\eta$, local Howe uniqueness applied
to (B29a) would give
\[
             \sigma_{\lambda,v}\simeq
                   \sigma_{\eta,v}
                   \quad\hbox{for every finite }v.
 \tag{B29d}
\]
Their infinity types are already equal by B27, so the
global irreducible source representations are equal.
B27's multiplicity-one result then identifies their
finite Hecke sectors.  Consequently distinct sectors
cannot share a target irreducible automorphic
constituent.  The discrete unitary automorphic
decomposition is orthogonal, even if a target
representation has multiplicity.  It follows that
\[
 E_0=\mathop{\bigoplus}^{\perp}_{\lambda}E_\lambda,
       \qquad T_\lambda^*T_\eta=0
                       \quad(\lambda\ne\eta).
 \tag{B29e}
\]
Together with B27's source orthogonality this also
gives $T_\lambda T_\eta^*=0$.
Only nonzero sectors need be retained.  There are
at most $r=\dim E_0$ of them.

Use the matrices $A_\lambda$ of (B27i), with target
metric $G_B$.  Each is algebraic, positive
semidefinite and self-adjoint for $G_B$, and its
image is $E_\lambda$.  Equation (B29e) implies
\[
           A_\lambda A_\eta=0\quad(\lambda\ne\eta).
 \tag{B29f}
\]
Let $B_\lambda$ be the inverse of $A_\lambda$ on
$E_\lambda$, extended by zero on its $G_B$-orthogonal
complement.  This is an algebraic matrix: the image
and the complement are algebraic, and one inverts
an invertible algebraic matrix on the image.
Alternatively it is a polynomial in $A_\lambda$
which is zero at zero and is the reciprocal function
on its finitely many nonzero eigenvalues.
Put $P_\lambda=B_\lambda A_\lambda$; these are the
algebraic orthogonal projectors onto $E_\lambda$,
and their sum is the identity on $E_0$.

With every factor from B27 retained, the actual
inverse on the original fixed kernel image is
\[
 \boxed{\displaystyle
  Q_T^{-1}
    =\frac{1}{\nu_G|P_\chi|^2|\gamma|^2}
                 \sum_\lambda c_\lambda^{-1}B_\lambda,
          \qquad \gamma=\frac{C_x}{t_xt_y}.}
 \tag{B29g}
\]
Multiplication by
$Q_T=\nu_G|P_\chi|^2|\gamma|^2
\sum_\lambda c_\lambda A_\lambda$
gives $\sum_\lambda P_\lambda=1$, proving the formula.
Every $B_\lambda$ and $P_\lambda$ is realized by the
CM-source correspondence construction in B23, on
the preserved algebraic image supplied by B25.
This is an algebraic block decomposition of the
actual inverse.  It does not make the positive
numbers $c_\lambda$ algebraic or prove that their
ratios are algebraic.  B28 supplies a bound for
this inverse without requiring those ratios.

The projectors and all the scalars in (B29g) are
common to every source character or vector using
the same original kernel.  They have not been
chosen anew for each torus character.  Passing to
different complete types can change $P$ and the
kernel, and is outside this uniformity statement.

\paragraph{Exact transport along the original compact-place variation.}
There is a further uniformity statement for the
actual 032 parameter variation in B25.  Retain its
literal oscillator identity
\[
      \Phi_t=c_t\omega_S(s_t)\Phi_0,
                    \qquad c_t\ne0,
 \tag{B29h}
\]
where $s_t\in S(F\otimes\R)$ changes only the compact
target-place decompositions and all finite data and
the active canonical cochain stay fixed.  The scalar
$c_t$ includes the chosen Gaussian and constant
half-line normalizations.  It is retained exactly.
Then covariance of the scalar theta kernel and
right invariance of the original quotient Haar
measure give
\[
 \begin{aligned}
 K_t(g,h)&=c_tK_0(g,hs_t),\\
 T_t&=c_tT_0 R_S(s_t^{-1}),\\
 Q_{T_t}&=|c_t|^2Q_{T_0},
      &Q_{T_t}^{-1}&=|c_t|^{-2}Q_{T_0}^{-1}.
 \end{aligned}
 \tag{B29i}
\]
Here $(R_S(s)u)(h)=u(hs)$ is unitary on the full
source Hilbert space.  For the second line simply
substitute $h'=hs_t$ in the bilinear integral;
unitarity then proves the third line.  Thus even
the actual positive operator, not only its range,
is fixed up to the stated scalar along this
compact-place variation.  The finite source
projectors commute with $R_S(s_t)$, so (B29i)
holds in every sector separately and their target
images are unchanged.  No claim that the moved
source vector has the old compact type is needed;
the identity is on the full source Hilbert space.

Since the oscillator representation is unitary,
the modulus in (B29h) also satisfies
\[
                |c_t|^2=
                   \frac{\|\Phi_t\|_\omega^2}
                        {\|\Phi_0\|_\omega^2}.
 \tag{B29j}
\]
This does not discard a normalization: it identifies
it with the norms of the two specified Schwartz
vectors.  In a fixed real Schr\"odinger chart the
Gaussian identity
$\|\exp(\pi i x^{\mathsf t}Zx)\|_2^2
 =\det(2\operatorname{Im}Z)^{-1/2}$
computes each changing pure Gaussian factor.
Any originally included scalar half-line factors
must still be multiplied into this calculation.

\paragraph{An exact criterion for the remaining scalar gate.}
Retain only the nonzero sectors and fix one of them,
$\lambda_0$.  Write
$C=\nu_G|P_\chi|^2|\gamma|^2>0$ for the full
external factor in (B29g).  Then the following three
conditions are equivalent:
\[
 \begin{aligned}
 &Q_T=dM\quad\hbox{for some }d>0
                 \hbox{ and }M\in\operatorname{End}_{\overline\Q}(E_0);\\
 &Q_T^{-1}=d'M'\quad\hbox{for some }d'>0
                 \hbox{ and }M'\in\operatorname{End}_{\overline\Q}(E_0);\\
 &c_\lambda/c_{\lambda_0}\in
          \overline\Q\cap\R_{>0}\quad\hbox{for every retained }\lambda.
 \end{aligned}
 \tag{B29k}
\]
The first two conditions are equivalent by inversion.
If the third holds, put
$a_\lambda=c_\lambda/c_{\lambda_0}$; the exact formulas are
$Q_T=Cc_{\lambda_0}\sum_\lambda a_\lambda A_\lambda$
and
$Q_T^{-1}=(Cc_{\lambda_0})^{-1}
\sum_\lambda a_\lambda^{-1}B_\lambda$.
Their matrix factors are algebraic and their scalar
factors are positive, so no original normalization
has been removed.

Conversely, suppose $Q_T=dM$ as in the first condition.
Let $r_\lambda=\dim E_\lambda>0$.  Since
$B_\lambda A_\eta=0$ for $\lambda\ne\eta$ and
$B_\lambda A_\lambda=P_\lambda$, one has
\[
 \operatorname{Tr}(B_\lambda Q_T)=Cc_\lambda r_\lambda,
 \qquad
 \frac{c_\lambda}{c_{\lambda_0}}
   =\frac{r_{\lambda_0}}{r_\lambda}
      \frac{\operatorname{Tr}(B_\lambda M)}
           {\operatorname{Tr}(B_{\lambda_0}M)}.
 \tag{B29l}
\]
Every trace on the right is algebraic, and the
denominator is nonzero; the displayed first equality
also shows that their ratios are positive real.
This proves necessity.  In a single nonzero sector
the condition is vacuous.  With several sectors it
is precisely a period-ratio assertion, without any
remaining off-diagonal or multiplicity ambiguity.
It is not a consequence of the block decomposition
alone, and no general rationality conjecture is
being used to assert it.

\paragraph{Source scope.}
The new external inputs are the precise finite local
Howe uniqueness theorems above.  Gan--Takeda's
definitions, Theorem 1.2 and arbitrary-residue scope,
and the complete final proof in Section 6 were
read; M\'inguez's Theorem 5.1 with its complete
uniqueness proof and Corollary 6.3 were read.
The passage from the actual global kernel to local
Hom spaces, orthogonality, algebraic block inverse,
and variation identities are proved here.
No global irreducibility theorem, Galois
cohomological theta correspondence, or extra
\texttt{MainX} arithmetic family is inferred.


\section{Explicit quantitative certificates}
\label{sec:quantitative}
The estimates in this section have completely named normalizations and inputs. They do not claim a scalar-invariant positive lower bound, or a uniform degree bound for the final algebraic cycles. The latter would require arithmetic bounds for the rational source maps and extraction coefficients.

\subsection{THETA3-Q1: fixed-input perturbation of the geometric period}\label{Q1}
Fix the compact surface $X$, the total finite theta input and the determinant-line trivializations of 032 Proposition~4.6. On smooth $(2,0)$-forms use the Hermitian pairing
\[
 (\xi,\eta)=\int_X\xi\wedge\overline\eta,
 \qquad \|\xi\|^2=(\xi,\xi).
\]
For a nonzero $\xi$ put $A=\|\xi\|^2$. If
$\|\xi'-\xi\|\leq\eta\|\xi\|$, $0\leq\eta<1$, then
\begin{equation}\label{eq:perturb}
 \left|\int_X\xi\wedge\overline{\xi'}\right|
 \geq (1-\eta)A.
\end{equation}
If $\xi'=\Xi(t)$, $\xi=\Xi(t_0)$, on a convex coordinate neighborhood with $\|D\Xi\|\leq L$, it suffices that
$\|t-t_0\|\leq\eta\sqrt A/L$. For $L=0$ the assertion is immediate. A certified lower bound $A\geq A_0>0$ can replace $A$ in this sufficient radius.

On a finite etale cover $p:Y\to X$ of degree $d$, suppose
\[
 p^*\xi=\sum_r a_r\alpha_r\wedge\beta_r,
 \qquad p^*\xi'=\sum_s b_s\gamma_s\wedge\delta_s.
\]
Set $S=(\sum_r|a_r|)(\sum_s|b_s|)$. At least one of the resulting four-form periods has absolute value at least
$d(1-\eta)A/S$.
\begin{proof}
Cauchy--Schwarz gives
$|(\xi,\xi'-\xi)|\leq\|\xi\|\|\xi'-\xi\|\leq\eta A$.
The reverse triangle inequality proves \eqref{eq:perturb}. Integrating $D\Xi$ on the segment from $t_0$ to $t$ proves the sufficient radius. The pairing of $(2,0)$-forms remains positive along a smooth intermediate path; holomorphy is needed at the endpoints for the cohomological interpretation. Pullback by a degree-$d$ cover multiplies integration by $d$. Expand the two finite sums and apply the triangle inequality. Both forms are nonzero because $\eta<1$, so $S>0$.
\end{proof}
\framework{This quantifies the continuity step of 032 Corollary~4.7 on a fixed Betti realization. The cycle is still the image of a surface under the rational Hodge-source maps. The relevant polarization/correspondence framework is cote~44, scans~7--8; it does not assert that this analytic constant is a degree of a cycle.}

\subsection{THETA3-Q2: a certified Gaussian isolation bound}\label{Q2}
Let $\Lambda\subset\R^d$ be a lattice with minimum nonzero norm at least $\rho>0$. Assume $\operatorname{Im}Z\geq\lambda I$, $\lambda>0$, for the period matrix in the first theta tensor. Fix $0\ne x_0\in\Lambda$ and an integer $k\geq2\|x_0\|/\rho$. The contribution of $x_0+k\Lambda$ differs from its $x_0$ term by a vector of norm at most
\begin{equation}\label{eq:gaussian-tail}
 E(k)=2\sqrt{\frac{\pi}{e\lambda}}
 \sum_{j\geq1}(2j+3)^d e^{-b j^2},
 \qquad b=\frac{\pi\lambda k^2\rho^2}{8}.
\end{equation}
For $f_j=(2j+3)^d e^{-bj^2}$ and
\[
 q_J=\left(\frac{2J+5}{2J+3}\right)^d e^{-b(2J+1)}<1,
\]
the tail beginning at $J$ is at most $f_J/(1-q_J)$.
If projected leading gradients $u_0,v_0$ satisfy
$\|u_0\wedge v_0\|=\Delta>0$, and the two projected errors are at most $e_1,e_2$, then
\begin{equation}\label{eq:wedge-error}
 \|u\wedge v-u_0\wedge v_0\|
 \leq e_1\|v_0\|+e_2\|u_0\|+e_1e_2.
\end{equation}
An upper bound strictly smaller than $\Delta$ certifies nonvanishing.
\begin{proof}
For $m\ne0$ in $\Lambda$, $\|x_0+km\|\geq k\|m\|/2$.
The first theta tensor summand has norm at most $2\pi r e^{-\pi\lambda r^2}$, where $r=\|x_0+km\|$. The inequality
$r e^{-\pi\lambda r^2/2}\leq1/\sqrt{\pi e\lambda}$ bounds this by
$2\sqrt{\pi/(e\lambda)}e^{-\pi\lambda k^2\|m\|^2/8}$.
Balls of radius $\rho/2$ around lattice points are disjoint. Points with
$j\rho\leq\|m\|<(j+1)\rho$ number at most $(2j+3)^d$, by comparing the volume of these balls with the ball of radius $(j+3/2)\rho$. Summing proves \eqref{eq:gaussian-tail}. The successive ratios $f_{j+1}/f_j$ decrease with $j$, which gives the geometric tail bound. A projection multiplies the error by its operator norm. Expanding
$(u_0+(u-u_0))\wedge(v_0+(v-v_0))$ gives \eqref{eq:wedge-error}.
\end{proof}
\paragraph{Numerical certificate in a specified Gaussian model.}
For $d=6$, $\rho^2=2$, $\lambda=1$, $k=6$, and orthogonal projected leading terms of amplitude $c=2\pi e^{-2\pi}$, the interval-arithmetic computation supplied with the package gives
\[
 c>0.01173348,\qquad E(6)<1.765593\cdot10^{-8},
 \qquad c^2-2cE(6)-E(6)^2>\frac{137}{10^6}.
\]
The six-dimensional lattice can be taken as the Minkowski image of
$\OO_{\Q(\sqrt5)}^3$: for one coordinate its squared norm is
$2a^2+2ab+3b^2$, with minimum~2. The computation certifies the displayed Gaussian estimate. Identification with a particular global theta convention still requires its symplectic coordinates, projection and determinant-line normalization; the numerical value is not asserted to be an unnormalized global period.

\framework{This is a quantitative form of 032 Proposition~4.5, at the level of sections of the cotangent bundle. It certifies a nonzero realization before the algebraic surface correspondence is extracted; cote~44, scans~7--8, locates that later correspondence/polarization step.}

\subsection{THETA3-Q3: rational compact data with a denominator}\label{Q3}
Choose a $\Q$-basis of $E$ and coordinates on a two-dimensional $E$-space $W$. Its diagonal real embedding has dimension $D=2[E:\Q]$. Let $B$ be the invertible real embedding matrix. For any archimedean target vector $u$ and $N\geq1$, rounding $NB^{-1}u$ coordinatewise to $a\in\Z^D$ gives $w=Ba/N\in W(E)$ with
\begin{equation}\label{eq:round}
 \|w-u\|\leq\frac{\|B\|\sqrt D}{2N}.
\end{equation}
At a place let the Hermitian matrix be $H$, with $\|H\|=h$, $\|u\|\leq U$, and $|u^*Hu|\geq\eta>0$. If $\epsilon=\|w-u\|\leq1$ and
$h(2U+1)\epsilon\leq\eta/2$, the signs of $w^*Hw$ and $u^*Hu$ agree. Their Hermitian line projectors satisfy
\begin{equation}\label{eq:projector}
 \left\|\frac{w(w^*H)}{w^*Hw}-\frac{u(u^*H)}{u^*Hu}\right\|
 \leq \frac{2h(2U+1)}{\eta}
       \left(1+\frac{hU^2}{\eta}\right)\epsilon.
\end{equation}
\begin{proof}
Each rounded coordinate has error at most $1/2$, proving \eqref{eq:round}. Both the denominator difference and the numerator difference in \eqref{eq:projector} are bounded by
$h(2U\epsilon+\epsilon^2)$. Hence $|w^*Hw|\geq\eta/2$. Writing the difference of the fractions as the numerator difference divided by $w^*Hw$, plus
$u(u^*H)(u^*Hu-w^*Hw)/((w^*Hw)(u^*Hu))$, and using
$\|u(u^*H)\|\leq hU^2$, proves the formula. Apply it at every real place. The complementary projector is $1-P$, so it has the same error bound.
\end{proof}
\framework{This makes the weak-approximation endpoint in 032 Corollary~4.7 explicit relative to a field basis and prescribed analytic neighborhood. It changes multiplicity subspaces, not the signature of $V$. The resulting rational datum can be spread over a finite type base; it is not a bound for its final cycle degree.}

\subsection{THETA3-Q4 and Q5: zero-orbit selection and extraction}\label{Q4}\label{Q5}
These are the quantitative versions of A2 and A3. Let the anisotropic theta integral in A1 use probability measures, with
$\Phi=\Phi_\infty\otimes\Phi_f$, $\Phi_f(0)=1$, and
$a=\Phi_\infty(0)\ne0$. For every $0<\epsilon<1$ there is
$h_f\in N(\A_f)$ such that
\begin{equation}\label{eq:finite-selection}
 |I_{\Phi_\infty}(\omega_f(h_f)\Phi_f)|\geq(1-\epsilon)|a|,
 \qquad\|\omega_f(h_f)\Phi_f\|_2=\|\Phi_f\|_2.
\end{equation}
Suppose inverse doubling gives the finite tensor expansion
$\sum_a c_a\bigotimes_{i=1}^4\phi_{i,a}$ and the $i$th torus expansion has at most $N_{i,a}$ permitted characters. Then an individual integrated fourfold period has absolute value at least the total absolute period divided by
\begin{equation}\label{eq:loss}
 S=\sum_a |c_a|\prod_{i=1}^4N_{i,a}.
\end{equation}
Zero character coefficients may be retained; exact zero testing is unnecessary.
\begin{proof}
By A1 the continuous function $T_\Phi$ has mean $a$ on the compact quotient $[N]$. Its maximum absolute value is at least $|a|$. At a maximizing representative $(n_\infty,n_f)$ choose rational elements $\gamma_j\in N(F)$ with $\gamma_{j,\infty}\to n_\infty$. Automorphy gives
$T_\Phi(\gamma_{j,\infty},n_f)=T_\Phi(1,\gamma_{j,f}^{-1}n_f)$.
Continuity proves the inequality for some $j$. The finite Weil action is unitary. Expand the finite products and apply
$|\sum_a c_a\sum_{\mathbf j}P_{a,\mathbf j}|
\leq S\max_{a,\mathbf j}|P_{a,\mathbf j}|$.
The relevant finite torus quotient bounds the permitted character counts by its order. For actual geometric integrals the explicit measure conversion in van Hoften--Sempliner \S4.2.8 must be included, together with connected-component coefficients. More generally, if the exact extracted identity is $I=\sum_\lambda b_\lambda J_\lambda$, use $S=\sum_\lambda|b_\lambda|$.
\end{proof}
\framework{The doubled oscillator pairing is a realization-level construction. Finite rational Hodge-source extraction and pushforward of the surface supply the cycle. These are different operations, and the displayed coefficient loss records their interaction.}

\subsection{THETA3-Q6: the local positive determinant}\label{Q6}
On $\C^2$ use $d\gamma(z)=\pi^{-2}e^{-\|z\|^2}dz$. Put
$\ell_u(z)=\sum_i u_i z_i$ and
$P_{u,v}(z,w)=\ell_u(z)\ell_v(w)-\ell_v(z)\ell_u(w)$.
Then
\begin{equation}\label{eq:fock-det}
 \|P_{u,v}\|^2_{L^2(d\gamma\otimes d\gamma)}
 =2\det\begin{pmatrix}\|u\|^2&\langle u,v\rangle\\
 \langle v,u\rangle&\|v\|^2\end{pmatrix}.
\end{equation}
For orthonormal $u,v$, $P_{u,v}/\sqrt2$ has norm~1. Under a unitary Bargmann transform and the compatible doubling intertwiner,
$\delta(f\otimes\bar f)(0)=\|f\|^2$, so this normalized local doubled value is~1.
\begin{proof}
The covariance identity $\int z_i\bar z_j\,d\gamma=\delta_{ij}$ shows that the two diagonal terms of $|P|^2$ each integrate to $\|u\|^2\|v\|^2$, and the two cross terms each integrate to $|\langle u,v\rangle|^2$ with negative sign. This proves \eqref{eq:fock-det}. The remaining assertions follow from the specified unitary transforms and the doubling identity.
\end{proof}
The determinant polynomial is the shape of the $q=1,p=2$ Fock wedge cocycle in the published Bergeron--Millson--M\oe glin, Proposition~5.17 and Proposition~5.19(1), pp.~48--49. The value~1 here is a chosen convention, not an assertion about a paper's unmodified numerical constants. With the raw convention retain
$a_0=\Phi_{\infty,\mathrm{split}}(0)\ne0$ from van Hoften--Sempliner Lemma~4.2.9. Its matrix-coefficient continuity and rational isometry approximation give
$|a|\geq(1-\eta)|a_0|$, and Q4--Q5 give
$(1-\epsilon)(1-\eta)|a_0|/S$.
\framework{This is local Hermitian positivity in the automorphic realization. It is compatible with the polarization language of cote~44, scans~7--8, but is not itself an algebraic intersection form or the global Alp\"oge determinant.}

\subsection{THETA3-Q7: a finite search with a certified grid size}\label{Q7}
View the additive Siegel unipotent $N$ over $\Q$, of dimension $d$. Let $K_N\subset N(\A_f)$ be a compact open additive subgroup fixing $\Phi_f$. Set
$\Lambda=N(\Q)\cap K_N$ and
$F(b)=T_\Phi(n(b),1_f)$. Let $B$ have as columns a $\Z$-basis of $\Lambda$. If $F$ is $L$-Lipschitz on $N(\R)/\Lambda$, then among the $N_0^d$ rational grid elements
\[
 \gamma=\sum_{i=1}^d\frac{j_i}{N_0}\lambda_i,
 \qquad 0\leq j_i<N_0,
\]
one has
\begin{equation}\label{eq:grid}
 |T_\Phi(1,n_f(-\gamma))|
 \geq |a|-\frac{L\|B\|\sqrt d}{2N_0}.
\end{equation}
Thus $N_0\geq\max(1,\lceil L\|B\|\sqrt d/(2\epsilon|a|)\rceil)$ suffices for the Q4 bound. The finite Schwartz norm is unchanged; any compact subgroup of $\U(V)(\A_f)$ fixing the total doubled input still fixes its translate.
\begin{proof}
Additive strong approximation gives $N(\Q)+K_N=N(\A_f)$ and identifies
the compact quotient as
\[
 N(\Q)\backslash N(\A)/K_N \simeq N(\R)/\Lambda.
\]
Its probability mean is $a$ by A1. Choose a maximum point of $|F|$ and round its coordinates modulo $\Z^d$ to the nearest grid point. The torus distance is at most $\|B\|\sqrt d/(2N_0)$. The Lipschitz estimate and reverse triangle inequality prove the bound for $F(\gamma)$. Rational automorphy gives $F(\gamma)=T_\Phi(1,n_f(-\gamma))$. Unitarity proves the norm assertion, and commuting dual-pair actions prove the compact-level assertion. Extracting separate four-factor tensors can still require a smaller common level.
\end{proof}
Here is an explicit sufficient value of $L$. On a chosen compact fundamental set for the anisotropic quotient suppose all rational supports lie in a lattice $L_0\subset\R^D$ of minimum at least $\rho$, the integrand vanishes outside $L_0$, and
\[
 |\omega(g)\Phi(x)|\leq C(1+\|x\|)^m e^{-\lambda\|x\|^2}
\]
uniformly, with known $C,\lambda>0$ and $m\geq0$. Write the unipotent phase as $e^{2\pi i\langle b,Q(x)\rangle}$, with
$\|Q(x)\|\leq c_Q\|x\|^2$. One may take
\begin{equation}\label{eq:L-bound}
 L=2\pi c_Q C\sum_{j\geq1}(2j+3)^D((j+1)\rho)^2
       (1+(j+1)\rho)^m e^{-\lambda j^2\rho^2}.
\end{equation}
Indeed termwise differentiation in a unit $b$-direction introduces at most $2\pi c_Q\|x\|^2$; probability integration cannot enlarge the resulting uniform bound. The shell packing argument of Q2 gives \eqref{eq:L-bound}. The differentiated series converges uniformly. Its terms have polynomial growth times a Gaussian, and their successive ratios eventually decrease below~1, giving an explicit geometric tail certificate. Compactness guarantees the existence of an envelope, but a numerical application must supply its constants or a fundamental set; no effective reduction-theory bound is being silently invoked.

With the Q6 convention $a_0=1$, choose the rational isometry so $|a|\geq1/2$. Taking
$N_0\geq\max(1,\lceil2L\|B\|\sqrt d\rceil)$ gives $|I|\geq1/4$. For an exact finite extraction $I=\sum_\lambda b_\lambda J_\lambda$, some geometric mixed period therefore satisfies
\begin{equation}\label{eq:quarter}
 |J_\lambda|\geq\frac1{4\sum_\lambda|b_\lambda|}.
\end{equation}
\framework{This gives a finite search in the adelic realization with a named denominator. It preserves the total doubled level, but does not bound the complexity of every subsequent rational Hodge map. The distinction between a realization and an algebraic correspondence is the one in cote~16, scans~43--45.}

\subsection{THETA3-Q8: a pointwise lower bound for the surface norm}\label{Q8}
Let an injective holomorphic chart in $X$ contain the Euclidean ball $B_R\subset\C^2$, and write
$\xi=f(z)\,dz_1\wedge dz_2$. Suppose $|f(0)|\geq b>0$ and the real differential of $f$ has norm at most $M$ on $B_R$. For $0<r<R$ and $Mr\leq b/2$,
\begin{equation}\label{eq:local-norm}
 \int_X\xi\wedge\bar\xi\geq\frac{\pi^2}{2}b^2r^4.
\end{equation}
\begin{proof}
Integrating the differential along a segment gives $|f(z)|\geq b/2$ on $B_r$. In the complex orientation,
$dz_1\wedge dz_2\wedge d\bar z_1\wedge d\bar z_2=4\,d\vol_{\R^4}$.
Since $\vol(B_r)=\pi^2r^4/2$, integration over $B_r$ gives the bound. The integrand is nonnegative on the rest of $X$. If $M=0$, any $r<R$ works.
\end{proof}
\framework{Together with Q2 and bounds for differentiated theta series, this turns a certified nonzero coefficient into the $A_0$ needed by Q1. The injective chart is required data. This is a Hodge norm estimate, not an integral intersection multiplicity.}

\subsection{THETA3-Q9: denominators in finite Hecke extraction}\label{Q9}
Let $\Lambda=H^1(X,\Z)/\mathrm{torsion}$, and let $T_j$ be commuting integral Hodge endomorphisms induced by rational algebraic correspondences. Suppose monic $f_j,g_j\in\Z[t]$ are coprime over $\Q$ and
$f_j(T_j)g_j(T_j)=0$. Set $R_j=\Res(f_j,g_j)\ne0$. An integral Bezout identity
$u_jf_j+v_jg_j=R_j$ gives
\[
 e_j=\frac{v_j(T_j)g_j(T_j)}{R_j},\qquad e=\prod_j e_j.
\]
Then $e$ is a rational Hodge projector induced by a rational algebraic correspondence, with
$e\Lambda\subset R^{-1}\Lambda$, $R=\prod_j|R_j|$. Its image is the intersection of the $f_j$-primary factors.

Let $K$ have degree $n$ over $\Q$. For $i=1,\ldots,4$, take such a projector $e_i$ with denominator $R_i$ and $u_i\in\Lambda\otimes K$ satisfying
$d_i u_i\in\Lambda\otimes\OO_K$. Put
\[
 P=\int_X e_1u_1\cup e_2u_2\cup e_3u_3\cup e_4u_4,
 \qquad D=\prod_{i=1}^4R_i d_i.
\]
Then $DP\in\OO_K$. If $P\ne0$ and
$|\sigma(P)|\leq M_\sigma>0$ for every $\sigma\ne\sigma_0$, then
\begin{equation}\label{eq:norm-bound}
 |\sigma_0(P)|\geq D^{-n}\prod_{\sigma\ne\sigma_0}M_\sigma^{-1}.
\end{equation}
\begin{proof}
The adjugate of the Sylvester matrix gives the integral Bezout identity. The Chinese remainder theorem gives $e_j=1$ on $\ker f_j(T_j)$ and $e_j=0$ on $\ker g_j(T_j)$; equivalently $e_j^2-e_j$ is divisible by $f_jg_j$ as a rational polynomial. Commutativity gives the product projector and its image. Integral polynomial numerators give the denominators. Polynomial operations preserve Hodge endomorphisms and are realized by addition and composition of algebraic correspondences.

Cup product evaluated on the integral fundamental class is a $\Z$-valued multilinear form on $\Lambda$. Multiplication of its four arguments by $R_i d_i$ therefore gives $DP\in\OO_K$. The nonzero algebraic integer $DP$ has nonzero integral norm, whose absolute value is at least~1. Consequently
$\prod_\sigma|\sigma(P)|\geq D^{-n}$. Divide by the other bounds to get \eqref{eq:norm-bound}.
\end{proof}
In the CM application, the cited Hecke/Frobenius and weak-admissibility arguments of 032 first identify the relevant eigenspaces as pure CM Hodge sources. Their eigenvectors can be chosen over a number field in the Betti realization. Harmonic representatives have the same cup product. Arbitrarily rescaled analytic theta forms need not be those algebraic Betti vectors; their comparison scalars must be retained. The matrices $T_j$, resultants, eigenvectors and conjugate bounds are required data. The supplied papers do not give universal numerical values for them.
\framework{This is a projector in the homological realization, not an assertion of an idempotent modulo rational equivalence. It implements the algebraic-correspondence requirement of cote~16, scans~44--45, with an explicit denominator. B7 supplies a positive Gram-matrix variant when total positivity under all embeddings is established.}


% THETA3-Q10 contribution by /root/theta_032, 21 November 2026.
% GPT-6 Astra Pro; inherited reasoning-effort label not separately exposed.

\subsection{THETA3-Q10: determinant normalization and a finite certificate in every degree}\label{Q10}

\paragraph{The exact Fock determinant norm.}
Let $u_1,\ldots,u_n\in\C^n$, put
$\ell_u(z)=\sum_a u_a z_a$, and let $z_1,\ldots,z_n$ be independent
standard complex Gaussian vectors, each with probability density
$\pi^{-n}e^{-\|z\|^2}$ relative to real Lebesgue measure. Define
\[
 P_u(z_1,\ldots,z_n)=\det(\ell_{u_i}(z_j))_{1\leq i,j\leq n},
 \qquad G_u=(\langle u_i,u_j\rangle)_{i,j},
\]
with the Hermitian inner product conjugate linear in its first
variable. Then
\begin{equation}\label{Q10-Fock}
 \|P_u\|^2=n!\det G_u,\qquad
 \langle P_u,P_v\rangle
       =n!\det(\langle u_i,v_j\rangle)_{i,j}.
\end{equation}

\paragraph{Proof.}
The Gaussian covariance is
$\mathbb E[\overline{z_a}z_b]=\delta_{ab}$, so
$\mathbb E[\overline{\ell_u(z)}\ell_v(z)]
=\langle u,v\rangle$.
Expand both determinants. Independence of the $n$ columns gives
\[
 \langle P_u,P_v\rangle
 =\sum_{\sigma,\tau\in\mathfrak S_n}
   \operatorname{sgn}(\sigma)\operatorname{sgn}(\tau)
       \prod_{j=1}^n\langle u_{\sigma(j)},v_{\tau(j)}\rangle.
\]
For fixed $\sigma$, relabel $i=\sigma(j)$ and put
$\rho=\tau\circ\sigma^{-1}$. The sign product is
$\operatorname{sgn}(\rho)$ and the inner sum is the determinant
of the displayed mixed Gram matrix. There are $n!$ choices of
$\sigma$. Taking $v=u$ proves the norm identity.
\hfill$\square$

For independent $u_i$, divide by $\sqrt{n!\det G_u}$ to obtain
a unit vector. In particular the coordinate determinant divided
by $\sqrt{n!}$ has norm one. BMM, Proposition 5.19(1), p.49,
with $q=1$ and $b=n$, identifies the coefficient of the
holomorphic wedge cocycle with precisely this determinant shape.
The other active Fock coordinates and all remaining compact
Gaussians may be normalized to norm one. The unitary Bargmann
transform preserves the resulting norm. Under the standard
doubling map,
\[
 \delta(f_\infty\otimes\bar f_\infty)(0)=\|f_\infty\|^2=1.
\]
This is a chosen complete normalization, not an assertion that
an unmodified source convention already has the number one.
For the actual mixed endpoint vector retain its value
$a=\Phi_\infty(0)$; a certified continuity estimate
$|a-1|\leq\eta<1$ gives $|a|\geq1-\eta$.

\paragraph{Dimensions in the doubled model.}
For the ball of dimension $n$, $\dim_E V=n+1$. Doubling the
rank-$n$ skew-Hermitian source gives a split source of rank $2n$.
The Siegel unipotent is the additive group
$N\simeq\operatorname{Herm}_n(E)$ over $F$. A Hermitian matrix has
$n$ diagonal $F$-entries and $n(n-1)/2$ off-diagonal $E$-entries;
therefore
\[
 d_N:=\dim_\Q N=n^2[F:\Q].
\]
The Schr\"odinger position space is $V^n$ over $E$, of dimension
$n(n+1)$, and hence its real dimension at all archimedean places is
\[
 D=2n(n+1)[F:\Q].
\]
These are different dimensions: $d_N$ controls the unipotent grid,
whereas $D$ controls the position-lattice Gaussian tail.
In the grid statement below, $N$ is viewed by restriction of
scalars as a $\Q$-vector group; thus $N(\Q)$ denotes its
$F$-rational Hermitian matrices and $N(\A_f)$ their finite
$F$-adelic counterparts.

\paragraph{The explicit unipotent grid.}
Use probability measures on the compact quotients, fix
$\Phi_f(0)=1$, and fix the splitting for which the $G=\U(V)$
action in the doubled model is untwisted pullback. Let $K_N$ be
a compact open additive subgroup of $N(\A_f)$ fixing $\Phi_f$.
Set
\[
 \Lambda_N=N(\Q)\cap K_N,
 \qquad B:\R^{d_N}\longrightarrow N(\R),
\]
where the columns of $B$ are a $\Z$-basis of this full lattice.
Suppose the function
$F(b)=T_\Phi(n(b),1_f)$ is $L$-Lipschitz on
$N(\R)/\Lambda_N$. For $N_0\geq1$ define the explicit grid
\[
 \Gamma_{N_0}=
 \left\{\sum_{j=1}^{d_N}\frac{a_j}{N_0}\lambda_j:
                         0\leq a_j<N_0\right\}.
\]
Then one of its $N_0^{d_N}$ elements satisfies
\begin{equation}\label{Q10-grid}
 |T_\Phi(1,n_f(-\gamma))|
    \geq |a|-\frac{L\|B\|\sqrt{d_N}}{2N_0}.
\end{equation}

\paragraph{Proof.}
Additive strong approximation gives
$N(\A_f)=N(\Q)+K_N$ and identifies the quotient by $K_N$
with $N(\R)/\Lambda_N$. The quotient measure is its probability
Haar measure. The general rank-$r$ statement A1, applied with
$r=n$, gives its mean $a=\Phi_\infty(0)$: Siegel-unipotent
averaging annihilates every nonzero Gram matrix, and anisotropy
of $V$ leaves only the zero tuple. Hence
$\max |F|\geq|a|$.
Round the coordinates of a maximum point modulo $\Z^{d_N}$ to
the nearest grid point. Their Euclidean error is at most
$\sqrt{d_N}/(2N_0)$, so their images differ on the torus by at
most $\|B\|\sqrt{d_N}/(2N_0)$. The Lipschitz and reverse-triangle
inequalities give the bound for $|F(\gamma)|$.
Rational theta automorphy gives
$F(\gamma)=T_\Phi(1,n_f(-\gamma))$, proving
\eqref{Q10-grid}. The finite Weil action is unitary and commutes
with $G(\A_f)$, so the translate preserves both the finite
Schwartz norm and every compact target level that fixed the total
input. A factorization into individual inputs can require a
smaller common level. \hfill$\square$

If $|a|\geq1/2$, the explicit choice
\[
 N_0\geq\max\{1,\lceil2L\|B\|\sqrt{d_N}\rceil\}
\]
makes the error at most $1/4$, so the selected total integral
has absolute value at least $1/4$. Suppose the exact finite
tensor, source-character and connected-component expansion is
\[
 I=\sum_\lambda b_\lambda J_\lambda,\qquad
 S_{\rm ext}=\sum_\lambda|b_\lambda|>0,
\]
where the geometric normalization is included in $b_\lambda$.
Then at least one actual $2n$-factor integrated period satisfies
\begin{equation}\label{Q10-extract}
 |J_\lambda|\geq\frac{1}{4S_{\rm ext}}.
\end{equation}
This is the triangle inequality. If the tensor decomposition has
coefficients $c_a$ and its $2n$ torus factors have at most
$N_{i,a}$ permitted characters, a valid cost before the geometric
conversion is
$\sum_a|c_a|\prod_{i=1}^{2n}N_{i,a}$.
Zero character components may be retained in the count. The
certificate permits the displayed finite input translation; it
is not a lower bound for every preassigned unaveraged period.

\paragraph{A fully specified Lipschitz bound.}
Here is the all-rank version of the sufficient envelope in Q7.
Fix a compact representative set for the anisotropic target
quotient and a position lattice of minimum norm at least $\rho$,
containing the rational support for all its finite components.
Suppose, in the chosen Euclidean coordinates,
\[
 |\omega(g)\Phi(x)|
      \leq C(1+\|x\|)^M e^{-\lambda\|x\|^2}
\]
on that set, with $C,\lambda>0$. Write the unipotent phase as
$\exp(2\pi i\langle b,Q(x)\rangle)$ and suppose
$\|Q(x)\|\leq c_Q\|x\|^2$. Then one may take
\begin{equation}\label{Q10-L}
 L=2\pi c_QC\sum_{j\geq1}(2j+3)^D
       ((j+1)\rho)^2(1+(j+1)\rho)^M
                   e^{-\lambda j^2\rho^2}.
\end{equation}
To check this, differentiate the phase in a unit direction.
Its derivative has absolute value at most
$2\pi c_Q\|x\|^2$. The zero term contributes zero.
The shell $j\rho\leq\|x\|<(j+1)\rho$ contains at most
$(2j+3)^D$ points by the disjoint-ball argument. The displayed
bound follows by summing absolute values and integrating against
probability measure. Gaussian decay makes the differentiated
series uniformly convergent and justifies the operations.
The positive terms have an explicit polynomial factor times
$e^{-\lambda\rho^2j^2}$; after any index for which their
successive ratios are bounded by a number $q<1$, their tail is
at most its first term divided by $1-q$. This supplies a finite
certified evaluation.

For the determinant polynomial and its doubled vector one may
use $M=2n$ after the Gaussian model and its linear transformations
have been fixed: partial Fourier transform and metaplectic
changes of a polynomial Gaussian preserve the bound on its
polynomial degree, though they change the Gaussian and its
coefficients. In coordinates $Q_a(x)=x^tH_ax$ one can take
$c_Q=(\sum_a\|H_a\|^2)^{1/2}$. The lattice, envelope constants,
and geometric conversion must actually be supplied for numerical
use; no unrecorded effective reduction-theory bound is assumed.

\paragraph{A pointwise certificate for the top-form norm.}
If an injective holomorphic chart contains $B_R\subset\C^n$ and
$\xi=f(z)\,dz_1\wedge\cdots\wedge dz_n$ there, suppose
$|f(0)|\geq b>0$ and the real derivative of $f$ has norm at most
$M_1$ on the ball. For $0<r<R$ with $M_1r\leq b/2$,
\begin{equation}\label{Q10-topnorm}
 i^{n^2}\int_X\xi\wedge\bar\xi
       \geq \frac{2^{n-2}\pi^n}{n!}\,b^2r^{2n}.
\end{equation}
Indeed the segment estimate gives $|f|\geq b/2$ on $B_r$.
The positive top-form density is $2^n|f|^2$ times real
Lebesgue measure, and
$\operatorname{vol}(B_r\subset\R^{2n})=\pi^nr^{2n}/n!$.
Integration proves the formula. A determinant of first-tensor
values can be certified from individual errors $\varepsilon_i$
using the multilinear estimate
\[
 \left\|\bigwedge_i u_i-\bigwedge_i u_i^0\right\|
 \leq\prod_i(\|u_i^0\|+\varepsilon_i)-\prod_i\|u_i^0\|.
\]
Expand the wedge by subsets and apply the norm bound for each
term to prove this inequality. Together with A6's lattice tail,
it gives explicit input data for \eqref{Q10-topnorm}.

\paragraph{The definite subgroup and its local constant.}
Choose an $E$-rational line $L$ negative at the active place.
Then $H=\U(L^\perp)$ is totally definite of rank $n$. For a
canonical $n$-form, its right-regular coefficient at that line
has compact character
\[
 \kappa(\operatorname{diag}(A,d))=\det(A)d^{-n}.
\]
This follows because the tangent representation is
$\Hom(L,Q)$, whose derivative is $A/d$ and whose determinant
is $\det(A)d^{-n}$. It is trivial on the scalar center.
Suppose a unit vector $w$ in a unitary representation has that
compact type. With $\pi(A)=\det A$ on $H_\infty=\U(n)$ and
probability Haar,
\[
 \int_{\U(n)}
   \langle w,\sigma(\operatorname{diag}(A,1))w\rangle
                      \pi(A)^{-1}\,dA=1.
\]
The matrix coefficient is $\det A$ in the convention linear
in the second variable, so the integrand is one. The rank-zero
unipotent is trivial. This direct local computation is valid
for every $n$; the particular archive proposition imposing
even $k=n$ is invoked only when its additional hypotheses hold.
No global norm is identified with this local number.

\paragraph{Increasing the rank does not enlarge the canonical gap.}
For the compact sign pattern $(+,\ldots,+,-)$, the half sums are
\[
 \rho_c=\left(\frac{n-1}{2},\frac{n-3}{2},\ldots,
                         \frac{1-n}{2};0\right),\qquad
 \rho_n=\left(\frac12,\ldots,\frac12;-\frac n2\right).
\]
The right-regular canonical lowest compact weight is
$(1,\ldots,1;-n)$. Solving the archive's lowest-weight relation
$\lambda+\rho_n-\rho_c=(1,\ldots,1;-n)$ gives
\[
 \lambda=(1,\ldots,1;-n)-\rho_n+\rho_c
       =\left(\frac n2,\frac n2-1,\ldots,
                          1-\frac n2;-\frac n2\right).
\]
The compact determinant representation has tuple
$u_i=(n+3-2i)/2$, $1\leq i\leq n$.
Taking $v$ to be the first $n-1$ entries of $\lambda$ leaves
the archive's two nodes
\[
 b=\left(\frac n2-1,\frac n2\right).
\]
Their gap is exactly one in every rank.
The interlacing follows from
$u_i=\lambda_i+1/2>\lambda_i>u_{i+1}$ and the final inequality
$u_n>-\!b_0$; all cross differences are half-integral.
Thus increasing the dimension does not repair the canonical
\texttt{NodeGap} condition, whose required gap is at least
three. Invoking a general discrete-series theorem also requires
the representation identification demanded by that theorem;
the displayed direct compact integral does not require it.

For a pluricanonical power $K_X^q$, $q\geq1$, the compact weight
is $(q,\ldots,q;-nq)$. The same computation gives
\[
 \lambda_i=q+\frac n2-i\quad(1\leq i\leq n),\qquad
 \lambda_{n+1}=-nq+\frac n2,
 \qquad
 b=\left(\frac n2-q,nq-\frac n2\right).
\]
Its gap is $(n+1)q-n$. For $n=2$ this is $3q-2$, as in B9.
It enlarges the gap by changing the line bundle and its vector;
it does not identify the resulting section with the original
ordinary degree-$n$ cohomology class.


% THETA3-Q11 contribution by /root/theta_032, 21 November 2026.
% GPT-6 Astra Pro; inherited effort label not separately exposed.

\subsection{THETA3-Q11: definite periods give an explicit lower bound for the ball norm}\label{Q11}

This bound replaces the unknown upper bound for the ball Gram
matrix in A5 by local holomorphic evaluation constants.  It uses
the actual finite definite restrictions and a fixed smooth
quotient throughout.  No derivative bound for the local
holomorphic coefficient is needed.

\paragraph{Fixed geometric and period normalizations.}
Let $X$ be a compact smooth ball quotient of complex dimension
$n$, or the finite disjoint union of components of a fixed neat
Shimura level.  Let $\mathcal E\subset H^0(X,\Omega_X^n)$ be a
finite-dimensional space in its specified scalar adelic
realization.  Use the geometric Hermitian norm
\[
 \langle\xi,\zeta\rangle_X
       =i^{n^2}\int_X\zeta\wedge\overline\xi,
 \qquad \|\xi\|_X^2=i^{n^2}\int_X\xi\wedge\overline\xi.
 \tag{Q11a}
\]
It is conjugate linear in the first variable.  If an adelic
Petersson measure is used instead, its specified conversion to
(Q11a) must be included; the argument does not identify these
measures without that conversion.

Take a nonempty finite collection of restrictions $R_{a_j}$ of A5, with finite adelic
$a_j$, totally definite $H=\U(n)$, and the fixed canonical
scalar type.  The same argument applies to any such finite
collection; separation is needed only when positive definiteness
on all of $\mathcal E$ is desired.  Give each $[H]$ probability
Haar measure, choose an orthonormal basis $e_{j,\alpha}$ of its
finite restriction module, and write
\[
 B_{j,\alpha}(\xi)
     =\int_{[H]}(R_{a_j}\xi)(h)\,
                              \overline{e_{j,\alpha}(h)}\,dh .
\]
For representatives of
$H(F)\backslash H(\A_{F,f})/K_H(a_j)$, let $r_{j,t}(\xi)$
be the corresponding scalar values.  Their exact masses are
\[
 m_{j,t}=
    \frac{|\Gamma_{j,t}|^{-1}}
         {\displaystyle\sum_s|\Gamma_{j,s}|^{-1}}>0,
 \qquad
 \Gamma_{j,t}=H(F)\cap h_{j,t}K_H(a_j)h_{j,t}^{-1}.
 \tag{Q11b}
\]
Here the finite stabilizers are viewed in the compact real
group.  The formula follows by decomposing the quotient into
$\Gamma_{j,t}\backslash H(F\otimes_\Q\R)$, first giving the
real group volume one and then normalizing total quotient mass.
A stabilizer character can force a coordinate to vanish; such
coordinates may be retained with value zero.  On a neat level
the relevant finite stabilizers are trivial.  In particular the
masses are concrete rational numbers, although the form values
have not thereby been proved algebraic.

Finite-dimensional Parseval gives the exact positive quantity
\[
 \mathcal P(\xi):=\sum_{j,\alpha}|B_{j,\alpha}(\xi)|^2
          =\sum_{j,t}m_{j,t}|r_{j,t}(\xi)|^2.
 \tag{Q11c}
\]
Flatten the pairs $(j,t)$ into a single index $i$.  Each scalar
value is a linear evaluation in the canonical fibre at a point
$x_i\in X$.  Choose an injective holomorphic chart
$\phi_i:B_{R_i}\subset\C^n\longrightarrow X$ with
$\phi_i(0)=x_i$ and $0<r_i<R_i$.  Write
\[
 \phi_i^*\xi=g_i(z)\,dz_1\wedge\cdots\wedge dz_n,
 \qquad r_i^{\rm ev}(\xi)=\gamma_i g_i(0),
 \tag{Q11d}
\]
where $r_i^{\rm ev}$ denotes the scalar evaluation formerly
indexed by $(j,t)$, so it is distinct from the radius $r_i$.
The nonzero constant $\gamma_i$ is fixed by the actual
coefficient frame: explicitly it is the value of that
evaluation functional on the coordinate canonical generator
at $x_i$.  Thus it includes every prescribed frame or
rigidification factor.  Neither $\gamma_i=1$ nor its
algebraicity is assumed.

Put
\[
 A_i=\frac{2^n\pi^n r_i^{2n}}{n!},\qquad
 c_i=\frac{m_i|\gamma_i|^2}{A_i},\qquad
 C_{\rm eval}=\sum_i c_i.
 \tag{Q11e}
\]
The radii, frame constants and masses are all data for the
chosen quotient; changing its level requires recording the
corresponding changed data.

\paragraph{The explicit bound.}
For every $\xi\in\mathcal E$ one has
\[
 \boxed{\quad
       \mathcal P(\xi)\leq C_{\rm eval}\|\xi\|_X^2,
       \qquad
       \|\xi\|_X^2\geq
                \frac{\mathcal P(\xi)}{C_{\rm eval}}.
       \quad}
 \tag{Q11f}
\]
In particular any explicit lower certificate for the sum of
definite period squares gives a lower certificate for the
geometric ball norm.  This does not use an upper bound for
$G_{\rm ball}$.

\textbf{Proof, including the local constant.}
For a holomorphic function $g$ on $B_R$ and $0<r<R$,
\[
       \int_{B_r}|g(z)|^2\,dV(z)
          \geq \frac{\pi^n r^{2n}}{n!}|g(0)|^2.
 \tag{Q11g}
\]
To see the mean identity directly, fix $w\in B_r$.  The
one-variable function $u\mapsto g(uw)$ is holomorphic on a
disc containing the closed unit disc, so its circle average is
$g(0)$.  Integrate this identity in $w$.  Rotation invariance
of $B_r$ and Fubini show that the average of $g$ on $B_r$ is
$g(0)$.  Cauchy--Schwarz then gives (Q11g), using
$\operatorname{vol}(B_r)=\pi^nr^{2n}/n!$.

In coordinates $z_j=x_j+iy_j$, the order of the exterior
factors gives exactly
\[
 i^{n^2}\,dz_1\wedge\cdots\wedge dz_n\wedge
        d\bar z_1\wedge\cdots\wedge d\bar z_n
        =2^n\,dx_1\,dy_1\cdots dx_n\,dy_n.
\]
Applying (Q11g) to $g_i$ therefore proves
\[
 \int_{\phi_i(B_{r_i})}i^{n^2}\xi\wedge\bar\xi
          \geq A_i |g_i(0)|^2.
 \tag{Q11h}
\]
The chart is injective, so this is an integral over a genuine
subset of $X$, with no covering multiplicity.  Multiplication
by $m_i|\gamma_i|^2/A_i$ gives
\[
 m_i|r_i^{\rm ev}(\xi)|^2
       \leq c_i\int_{\phi_i(B_{r_i})}i^{n^2}\xi\wedge\bar\xi
       \leq c_i\|\xi\|_X^2.
\]
Summing and using (Q11c) proves (Q11f).
\hfill$\square$

The same proof retains overlap information.  If
$U_i=\phi_i(B_{r_i})$, one may replace $C_{\rm eval}$ by
\[
 C_{\rm ov}=\sup_{x\in X}\sum_{i:x\in U_i}c_i
         \leq C_{\rm eval}.
 \tag{Q11i}
\]
Indeed the sum of the local integrals is the integral against
the displayed nonnegative weight function.  In particular,
if no point lies in more than $\kappa$ chosen charts, then
$C_{\rm ov}\leq\kappa\max_i c_i$.
The single-point consequence is also sharper than a derivative
envelope: if $|g_i(0)|\geq b$, then
$\|\xi\|_X^2\geq A_i b^2$ for every recorded injective radius.
It removes the Lipschitz hypothesis from the pointwise estimate
in Q10 and its factor-four loss.

\paragraph{A mixed-period and component-extraction certificate.}
Suppose $\xi,\xi'\in H^0(X,\Omega_X^n)$ are expressed on the
same quotient and in the same canonical line bundle, and
\[
       \|\xi'-\xi\|_X\leq\eta\|\xi\|_X,
       \qquad 0\leq\eta<1.
 \tag{Q11j}
\]
They may have been constructed from different theta types;
that origin does not replace the asserted norm comparison.
Cauchy--Schwarz and (Q11f) give
\[
 \begin{split}
 \operatorname{Re}\langle\xi,\xi'\rangle_X
      &\geq(1-\eta)\|\xi\|_X^2,\\
 |\langle\xi,\xi'\rangle_X|
      &\geq\frac{1-\eta}{C_{\rm eval}}
                    \sum_{j,\alpha}|B_{j,\alpha}(\xi)|^2.
 \end{split}
 \tag{Q11k}
\]
For completeness, subtract $\|\xi\|_X^2$ from the pairing:
the difference is $\langle\xi,\xi'-\xi\rangle_X$, whose
absolute value is at most $\eta\|\xi\|_X^2$.  This proves
the first inequality and hence the second.

If the constructed comparison form has a finite expansion
$\xi'=\sum_\nu b_\nu\xi'_\nu$, set
$S=\sum_\nu|b_\nu|$.  For $\xi\ne0$, (Q11j) implies $S>0$.
Linearity in the second variable and the triangle inequality
then prove that at least one actual component satisfies
\[
 |\langle\xi,\xi'_\nu\rangle_X|
       \geq\frac{1-\eta}{S C_{\rm eval}}
                    \sum_{j,\alpha}|B_{j,\alpha}(\xi)|^2.
 \tag{Q11l}
\]
If both forms are expanded, the same proof uses the product
of their two coefficient $\ell^1$ sums.  The finite tensor and
character-extraction costs in A3 or Q7 can be inserted into
$S$ once their exact expansions are fixed.  No uniform bound
for those costs is silently supplied here.

\paragraph{A determinant certificate without the ball Gram matrix.}
Choose any basis $\xi_1,\ldots,\xi_D$ of $\mathcal E$ and form
the weighted evaluation matrix
\[
 E_{ia}=\sqrt{m_i}\,r_i^{\rm ev}(\xi_a),\qquad
 H_{\rm fin}=E^*E.
\]
For $\xi=\sum_a u_a\xi_a$ one has
$\mathcal P(\xi)=u^*H_{\rm fin}u$.  Suppose a square
$D$-row minor has determinant of modulus at least $\delta>0$
and $\operatorname{tr}H_{\rm fin}\leq T$.  If $D>1$, then
\[
 \|\xi\|_X^2
    \geq\frac{\delta^2}{C_{\rm eval}}
               \left(\frac{D-1}{T}\right)^{D-1}
                         \sum_a|u_a|^2 .
 \tag{Q11m}
\]
To prove this, Cauchy--Binet gives
$\det H_{\rm fin}\geq\delta^2$.  The product of all but
the smallest eigenvalue is at most
$(T/(D-1))^{D-1}$ by the arithmetic--geometric mean inequality.
Thus its smallest eigenvalue is at least the numerator in
(Q11m) before division by $C_{\rm eval}$.  Apply (Q11f).
For $D=1$ the corresponding bound is
$\delta^2|u_1|^2/C_{\rm eval}$.

This is a finite matrix certificate in prescribed coordinates;
an absolute positive number also requires a lower bound for
$\sum_a|u_a|^2$ for the particular input.  The exact mass and
height/denominator arguments in B5 and B7 may provide the
definite data only under their stated arithmetic normalization
hypotheses.  Rationality of the masses alone proves neither
algebraicity of the entries $r_i^{\rm ev}(\xi_a)$ nor a height
bound for their determinants.

\paragraph{Meaning for the comparison problem.}
The new constant $C_{\rm eval}$ depends on local injective
charts, the exact coefficient frames and finite definite
masses.  These data replace an unknown bound for
$G_{\rm ball}$ in this lower-bound argument.  They still must
be supplied to obtain a numerical certificate, as must any
positive definite-period lower bound and the relative error
in (Q11j).  The argument does not compute those numbers from
signature alone, nor identify a single Bessel period with the
whole ball pairing.  It gives an explicit route from the
finite definite periods to the same geometric norm that
controls the mixed theta period.

\framework{This is a comparison of positive Hermitian forms
and specified fibre evaluations.  It fits the polarization
and correspondence discussion in Grothendieck's
\emph{cote44}, scans 7--8, while retaining the distinction
between fibrewise evaluation and the cohomological operations
discussed in \emph{cote16}, scans 43--46.  The analytic estimate
does not turn the evaluation matrix or the frame correction
of A5 into a rational algebraic correspondence.}


% Non-CM and integral-model part. Uses only standard LaTeX, amsmath/amssymb,
% and the common macros authorized in the main paper.
\section{Non-CM sources and algebraic extraction}

\subsection{THETA3-C01: the precise congruence $H^1$ obstruction}\label{C01}

\paragraph{Statement.}
Let $X$ be a smooth projective congruence unitary ball quotient in the
setup of Y.~Liu, \emph{Fourier--Jacobi cycles and arithmetic relative
trace formula}, with Hermitian rank $m\geq3$. If $B$ is a simple complex
abelian variety which is not of CM type, then
\[
 \Hom_{\mathrm{Hdg}}\bigl(H^1(B,\Q),H^1(X,\Q)\bigr)=0.
\]
The same conclusion holds for any character-generated $H^1$ subspace
which is a sum of images of CM Hodge structures. Sums, tensor products,
duals, and Hodge subquotients of such structures remain of CM type.

\paragraph{Proof.}
Liu's Corollary~4.20, pp.~54--55, gives the CM isogeny decomposition of
the whole Albanese in rank $m\geq3$. Consequently $H^1(X,\Q)$ has
commutative Mumford--Tate group. A subquotient of a CM Hodge structure
is CM, because its Hodge homomorphism factors through an image of the
same torus. A nonzero Hodge morphism from $H^1(B,\Q)$ is injective:
its kernel is a rational sub-Hodge structure, and simplicity follows
from the equivalence between abelian varieties up to isogeny and
polarizable rational Hodge structures of weight one. Such an
injection into a CM structure would force $B$ to be CM, a
contradiction. A finite collection of CM structures is controlled by
a product of tori; all its tensor constructions and subquotients are
therefore controlled by a torus as well. $\square$

This is a deduction from the cited CM decomposition, not a new proof
of that decomposition. Liu's Hermitian rank is $m$, so the ball has
dimension $m-1$. In rank two the cited statement concerns the
endoscopic Albanese part, not the entire Albanese. The attached
van Hoften--Sempliner paper gives the relevant rank-three result in
Proposition~3.1.7, Theorem~3.1.11, and \S3.1.12. Increasing the ball
dimension does not remove this $H^1$ obstruction. Neither statement
applies to arbitrary noncongruence subgroups or to all higher
cohomology of the ball quotient.

\paragraph{Framework anchor.}
The relevant realization is the Albanese, not an arbitrary higher-degree
Hodge summand.  Grothendieck's letter of 12 August 1964, cote~44,
scans~7--8, explicitly places $\operatorname{Pic}^0$ and
$\operatorname{Alb}^0$ in the correspondence and polarization framework.
The CM restriction above is the additional automorphic input, not a
statement extracted from that letter.

\subsection{THETA3-C02: a trace-zero selection rule on covers}\label{C02}

\paragraph{Statement.}
Let $f:X\to S$ be a finite \'etale cover of connected smooth
projective complex surfaces of degree $d$. Then
\[
 H^1(X,\Q)=f^*H^1(S,\Q)\oplus
       \ker\bigl(f_*:H^1(X,\Q)\to H^1(S,\Q)\bigr).
\]
If $u$ belongs to the second summand and $v_1,v_2,v_3\in H^1(S,\Q)$,
then
\[
 \int_Xu\smile f^*v_1\smile f^*v_2\smile f^*v_3=0.
\]
If $H^1(S,\Q)$ is CM, every simple non-CM sub-Hodge structure of
$H^1(X,\Q)$ belongs to this trace-zero summand.

\paragraph{Proof.}
The transfer satisfies $f_*f^*=d\,\mathrm{id}$. Therefore
$d^{-1}f^*f_*$ is a projector with image $f^*H^1(S,\Q)$ and kernel
$\ker f_*$. The projection formula gives
\[
 \int_Xu\smile f^*v_1\smile f^*v_2\smile f^*v_3
 =\int_S f_*u\smile v_1\smile v_2\smile v_3=0.
\]
The restriction of $f_*$ to a simple non-CM substructure vanishes by
the argument of C01. $\square$

Thus one new non-CM factor and three old theta factors cannot give
the desired nonzero four-form period. Two new factors, for example a
holomorphic form and its conjugate, are the first possible case.

\paragraph{Framework anchor.}
The maps here are the pullback and trace of an actual finite geometric
correspondence.  Their action on the Albanese realization fits cote~44,
scans~7--8.  The vanishing calculation is a projection-formula statement;
it is not an assertion that arbitrary projectors in a realization are
represented by algebraic cycles.

\subsection{THETA3-C03: a mixed non-CM and CM period}\label{C03}

\paragraph{Hypotheses and source.}
Let $S_0=\Gamma_0\backslash D$ be a compact principal congruence
ball surface in the exact setup of 032, \S4. Suppose a compact ball
surface $X_0$ commensurable with $S_0$ has a nonconstant morphism
$a_0:X_0\to E_0$ to a non-CM elliptic curve.
Schoen's \emph{An arithmetic ball quotient surface whose Albanese
variety is not of CM type}, pp.~133--134, supplies such a
commensurability class over $\Q(\zeta_5)/\Q(\sqrt5)$, with
$j(E_0)=2^{11}$. To use it, choose 032's Hermitian space to be the
particular rank-three space in Schoen's construction. Multiplication
of the form by an element of the totally real field with suitable
real signs gives signatures $(2,1)$ and $(3,0)$. Equality of CM
fields alone would not imply commensurability. The determinant image
of an arithmetic unitary subgroup is finite: it is commensurable
with an integral subgroup, and integral norm-one units in a CM field
form a finite group. A finite-index subgroup therefore lies in
$\SU(V)$. Keep the principal congruence surface as the source of the
theta construction.

\paragraph{Statement.}
Fix an allowed CM sign function $v$ of 032 with $v(1)=1$.
There exist a sufficiently deep principal cover $S$ of $S_0$, a
common finite \'etale cover $X$ of $X_0$ and $S$, and a theta
one-form $\beta$ on $S$ in the complex span of rational source maps
from $U(v)$ such that, for a nonzero invariant form $\alpha_0$ on
$E_0$,
\[
 \xi=a^*\alpha_0\wedge f^*\beta\ne0,
 \qquad \int_X\xi\wedge\overline\xi>0,
\]
where $a:X\to E_0$ and $f:X\to S$ are the induced maps. Moreover
there are a CM abelian variety $B(v)$, a morphism $b:X\to B(v)$,
and a form $\gamma\in H^0(B(v),\Omega^1)$ with
$a^*\alpha_0\wedge b^*\gamma\ne0$. The image of $(a,b)$ has
dimension two.

\paragraph{Proof.}
The pullback $\alpha$ of $a_0^*\alpha_0$ to $D$ is nonzero. Choose
$x\in D$ with $\alpha_x\ne0$. By 032 Proposition~4.5,
pp.~21--22, finitely many actual finite Schwartz inputs give theta
one-form values spanning $T_x^*D$. At least one resulting $\beta_x$
is independent of $\alpha_x$. These finitely many inputs descend to
a sufficiently deep principal cover $S$ of $S_0$. Refine it using
032 Theorem~5.1, p.~26, and Corollary~6.7, p.~45, so that $\beta$
has its rational CM-source expansion. Only now form $X$ from the
intersection of the torsion-free commensurable lattices defining
$X_0$ and $S$. The common finite cover preserves the nonzero wedge
at a lift of $x$.

On the projective surface $X$, a nonzero holomorphic two-form has
strictly positive integral of $\xi\wedge\overline\xi$. Indeed, in
local coordinates $\xi=g\,dz_1\wedge dz_2$, this is a nonnegative
real volume form and is positive wherever $g\ne0$.

The CM extraction writes $[\beta]$ as a finite complex sum of
$h_j(u_j)$, where $h_j:U(v)\to H^1(S,\Q)$ are rational Hodge maps
and $u_j$ lie in the selected $(1,0)$ line. Holomorphic one-forms on
a compact K\"ahler manifold are the unique holomorphic
representatives of their classes, so the same expansion holds for
the forms. At least one term satisfies
$\alpha\wedge f^*h_j(u_j)\ne0$. By the Albanese realization used
in 032 Lemma~2.2, a positive integer multiple of $h_j$ is induced
by an algebraic map $S\to B(v)$. Composing with $f$ gives $b$.
If $(a,b)$ had image dimension at most one, its differential would
have rank at most one on a dense open set, forcing this wedge to
vanish. Therefore its image has dimension two, and
\[
 (a,b)_*[X]\in
 \CH^{\dim B(v)-1}(E_0\times B(v))_{\Q}
\]
is an algebraic image-surface cycle with the stated nonzero pairing.
$\square$

This construction uses two new factors, $\alpha$ and
$\overline\alpha$, and two CM theta factors. It does not realize
the new forms as adelic character theta lifts on a noncongruence
subgroup. It also does not identify a previously unknown Hodge class
as algebraic; C05 explains the limited Hodge content of this target.

\paragraph{Framework anchor.}
The surface image is an actual algebraic correspondence, and the nonzero
fourfold period is its Betti pairing.  Cote~44, scans~7--8, is the
correspondence/polarization anchor for that passage.  The extra non-CM
Albanese factor and the transverse theta form are supplied by the two
modern source constructions, not by an archival universality statement.

\subsection{THETA3-C04: the exact projector and rank gate}\label{C04}

\paragraph{Statement.}
Let $A$ be smooth projective, and let
$T\subset H^{2p}(A,\Q)(p)$ be a $d$-dimensional space of rational
Hodge classes. Suppose an algebraic correspondence acts as a
projector $e$ onto $T$. For algebraic cycles $Z_1,\ldots,Z_r$ of
codimension $p$, a basis $t_1,\ldots,t_d$ of $T$, and a
nondegenerate rational bilinear form $q_T$ on $T$, put
\[
 M_{ij}=q_T\bigl(e\,\operatorname{cl}(Z_i),t_j\bigr).
\]
If $\operatorname{rank}M=d$, every element of $T$ is algebraic.
Alternatively, if an algebra $\mathcal D$ of algebraic
correspondences preserves $T$, it suffices that the
$\mathcal D$-orbit of $e\operatorname{cl}(Z)$ spans $T$ for one
cycle $Z$. If $T$ is a simple $\mathcal D$-module, a nonzero
$e\operatorname{cl}(Z)$ suffices.

\paragraph{Proof.}
Every $e\operatorname{cl}(Z_i)$ is algebraic. Nondegeneracy
identifies $T$ with its dual, so full row rank means that these
vectors span $T$. In the orbit version, the finite-dimensional
orbit span is generated by algebraic correspondences applied to an
algebraic cycle. A nonzero submodule of a simple module is the whole
module. $\square$

A single nonzero scalar pairing does not give this rank assertion.
For example, $\Q(1,0)\subset\Q^2$ pairs nontrivially with $(1,1)$
under the ordinary dot product, although $(1,1)$ is outside that
subspace. This is a linear-algebra counterexample to an inference,
not an example of a nonalgebraic Hodge class. The pure embedding-line
divisors of 032 Proposition~3.1 supply the necessary CM extraction
in that paper. For a general non-CM $H^1$, the polarization provides
a sum of rank-one tensors; it does not supply individual
holomorphic eigenline projectors as algebraic correspondences.

\paragraph{Framework anchor.}
The distinction between an inverse or projector on cohomology and one
induced by an algebraic correspondence is explicit in cote~16,
scan~44, conditions~(iv) and~(iv bis), with the discussion continuing on
scans~45--46.  This finding isolates the same existence requirement and
adds the finite-dimensional rank test needed for the specified target.

\subsection{THETA3-C05: one non-CM elliptic factor gives a known consequence}\label{C05}

\paragraph{Statement; no novelty claim.}
Let $E_0$ be a non-CM elliptic curve, and let $B$ be a CM abelian
variety for which the rational Hodge conjecture holds. Then every
Hodge class on $E_0\times B$ is algebraic. More precisely,
\[
 \Hdg^p(E_0\times B)
 =1\otimes\Hdg^p(B)
 \oplus[\mathrm{pt}]\otimes\Hdg^{p-1}(B),
\]
where $\Hdg^p(Y)$ denotes the rational $(0,0)$ classes in
$H^{2p}(Y,\Q)(p)$.

\paragraph{Proof.}
The classical Hodge-group facts are
$\operatorname{Hg}(E_0)=\operatorname{SL}(H^1(E_0,\Q))$ and
$\operatorname{Hg}(B)=T$ for a torus $T$. For the first statement,
see Moonen, \emph{Notes on Mumford--Tate groups}, \S2,
Example~2.6, p.~11. The connected Hodge group $H$ of the product
projects onto both factors. Its derived group projects onto
$\operatorname{SL}_2$ and trivially to $T$. As a subgroup of
$\operatorname{SL}_2\times1$ with full projection to
$\operatorname{SL}_2$, it equals that group. The remaining quotient
surjects onto $T$, whence $H=\operatorname{SL}_2\times T$.
The only $\operatorname{SL}_2$ invariants in
$H^*(E_0,\Q)=\Q\oplus H^1(E_0,\Q)\oplus\Q(-1)$ occur in degrees
zero and two. Taking invariants in the K\"unneth decomposition gives
the formula. The two kinds of class are represented by a pullback of
a cycle on $B$, or by a point on $E_0$ times such a cycle. $\square$

\paragraph{Framework anchor.}
This is a consequence inside the tensor category generated by the
weight-one realizations, with the elliptic degree-two factor represented
by a divisor.  It uses the Albanese/correspondence framework of cote~44,
scans~7--8; it does not supply the additional higher-degree algebraic
correspondences discussed in cote~16, scans~43--46.



\subsection{THETA3-C12: four non-CM factors and a period of 500}\label{C12}
Use Schoen's surface $Y$ and curves from pp.~133--134 of the cited paper. The two projections to the Fermat quintic $C_5$ are
$(y_{01}:y_{03}:y_{13})$ and $(y_{02}:y_{03}:y_{23})$.
Let $C_0$ be its degree-five quotient, with equation $y^2=1-x^5$, and let $C\to C_0$ be Schoen's etale double cover, obtained from the elliptic curve
\[
 F:\quad s^2=1+x+x^2+x^3+x^4.
\]
The source proves $j(F)=2048$ and that $F$ is non-CM. There is a degree-two map $C\to F$.

There is a connected etale degree-four cover $X\to Y$ and a generically finite map $(a_1,a_2):X\to F^2$ of degree~500. If
$\int_F i\alpha\wedge\bar\alpha=1$, then
\begin{equation}\label{eq:500}
 \int_X(a_1^*\alpha\wedge a_2^*\alpha)
 \wedge\overline{(a_1^*\alpha\wedge a_2^*\alpha)}=500.
\end{equation}
All four degree-one sources are non-CM. The map
$(a_1,a_2,a_1,a_2):X\to F^4$ pushes its fundamental cycle to
$500[\Delta_{13}\cap\Delta_{24}]$, a product of diagonal divisor classes.
\begin{proof}
On $y_{03}=1$ write $u=y_{13}$ and $v=y_{23}$. The function field of $C_5^2$ has equations
$y_{01}^5=1-u^5$, $y_{02}^5=1-v^5$. Adjoining $y_{12}$ with
$y_{12}^5=u^5-v^5$ gives the function field of $Y$. The radicand has valuation~1 at the generic point of the diagonal of $C_5^2$, since $u^5-v^5=(u-v)$ times a unit there. It is not a fifth power, so the Kummer extension has degree~5. Resolution leaves function fields unchanged. Hence the product of the two curve maps gives a morphism $Y\to C_0^2$ of generic degree $5\cdot5\cdot5=125$.

The connected Galois etale cover $C^2\to C_0^2$ has degree~4 and group $(\Z/2)^2$. Its function field is linearly disjoint from the degree-125 extension, since the intersection degree divides both~4 and~125. Therefore
$X=Y\times_{C_0^2}C^2$ is connected and etale of degree~4 over $Y$. It is smooth projective and remains an arithmetic ball quotient. Its generic degree over $C^2$ is~125, and the two degree-two maps to $F$ give generic degree $125\cdot2\cdot2=500$ over $F^2$.

For $\eta=\mathrm{pr}_1^*\alpha\wedge\mathrm{pr}_2^*\alpha$ the Koszul sign gives
$\int_{F^2}\eta\wedge\bar\eta=-(-i)^2=1$.
Pushforward of the fundamental class under a proper generically finite morphism multiplies it by its degree, proving \eqref{eq:500}. Postcomposing with $(p,q)\mapsto(p,q,p,q)$ proves the cycle formula. The image is the transverse intersection of the two diagonal divisors.
\end{proof}
\framework{This realizes a four-factor non-CM period by an algebraic surface correspondence, in the sense of cote~44, scans~7--8. Its class is already generated by divisors. The new forms arise on a noncongruence cover, rather than from the original adelic character lifts. Thus neither a new exceptional Hodge-class case nor an extension of the original congruence theta packet is asserted.}


\subsection{THETA3-C13: an explicit algebraic projector onto the Weil space}\label{C13}

\paragraph{Statement.}
Let $A$ be a complex abelian variety of dimension $2n$, with an action of
an imaginary quadratic field $E$ of Weil type: $H^1(A,\Q)$ has rank
$2n$ over $E$, and the two embeddings of $E$ each occur with multiplicity
$n$ on $H^{1,0}(A)$. Put
\[
 W_E(A)=\bigwedge_E^{2n}H^1(A,\Q)\ \subset\ H^{2n}(A,\Q).
\]
Here the inclusion denotes the usual rational two-dimensional Weil
subspace, whose complexification is the sum of the two extreme
eigenspaces. Then the projector onto $W_E(A)$ is represented by an
explicit rational algebraic correspondence on $A\times A$.
Consequently a single algebraic codimension-$n$ cycle $Z$ with
$e_W\cl(Z)\ne0$ suffices to prove that every Weil class of $A$ is
algebraic. More concretely, it suffices to find a smooth projective
$n$-fold $X$, a morphism $b:X\to A$, and $w\in W_E(A)_{\C}$ such that
\begin{equation}\label{C13-period}
 \int_X b^*w\ne0.
\end{equation}
This is an extraction criterion for a supplied cycle, not an assertion
that such a cycle has been constructed for every $A$.

\paragraph{Proof and explicit denominators.}
Let $\mathcal R=E\cap\operatorname{End}(A)$, an order in $E$.
Choose $0\ne b_0\in\mathcal R$ with $\bar b_0=-b_0$, and choose an
integer $M>2n|b_0|$. Set $a=M+b_0$. After choosing one embedding of
$E$ and replacing $b_0$ by its negative if necessary, write
$a=|a|e^{i\vartheta}$ with $0<\vartheta<1/(2n)$.
The $2n+1$ numbers
\[
 \lambda_j=a^j\bar a^{\,2n-j},\qquad 0\le j\le2n,
\]
are distinct: the ratio of two of them is
$e^{2i(j-k)\vartheta}$, whose argument is nonzero and has absolute
value less than $2<2\pi$. Moreover $a^{2n}\notin\R$.
The existence of $b_0$ follows by taking $c-\bar c$ for any nonreal
$c\in\mathcal R$; since $\bar c=\Tr_{E/\Q}(c)-c$ and the trace is an
integer, this difference still belongs to $\mathcal R$.

Let $V_\tau,V_{\bar\tau}$ be the two complex eigenspaces of $H^1(A)$.
The endomorphism $T=[a]^*$ acts on
\[
 H^{2n}(A,\C)=\bigoplus_{j=0}^{2n}
       \bigwedge^jV_\tau\otimes\bigwedge^{2n-j}V_{\bar\tau}
\]
by $\lambda_j$. The two extreme summands are
$W_E(A)_\C$. Define
\[
 f(t)=(t-a^{2n})(t-\bar a^{\,2n}),\qquad
 g(t)=\prod_{j=1}^{2n-1}(t-\lambda_j).
\]
Both are monic in $\Z[t]$: their roots are algebraic integers and
their root sets are stable under the two automorphisms of $E/\Q$.
They are coprime. Put $R=\operatorname{Res}(f,g)\ne0$.
The adjugate of the integral Sylvester matrix supplies
$u,v\in\Z[t]$ with $uf+vg=R$. Thus
\[
 e_W=R^{-1}v(T)g(T)
\]
acts as one on the two extreme summands and as zero on every other
summand of $H^{2n}(A)$. For a projector on the whole cohomology, compose
with the explicit middle K\"unneth projector
\[
 \pi_{2n}=
 \prod_{\substack{0\le k\le4n\\k\ne2n}}
   \frac{[N]^*-N^k\,\mathrm{id}}{N^{2n}-N^k},\qquad N\ge2.
\]
Indeed $[N]^*$ acts by $N^k$ on $H^k(A)$. Endomorphism pullbacks,
their compositions, and the identity are induced by algebraic
correspondences. The displayed polynomial is therefore a rational
algebraic correspondence. Its idempotence is asserted modulo
homological equivalence; no Chow-idempotence is needed here.
An explicit common denominator is $R$ times the product of the
nonzero integers $N^{2n}-N^k$ in the last display.

The two extreme summands have Hodge type $(n,n)$ by the Weil-type
condition. Their rational sum has dimension two, and $T$ on it has
the irreducible quadratic minimal polynomial $f$, because
$a^{2n}\notin\Q$. If $0\ne\alpha=e_W\cl(Z)$, then
$\alpha,T\alpha$ are linearly independent over $\Q$; a dependence
would give a rational eigenvalue of $T$ on this space. Both are
classes of rational algebraic cycles. They therefore span $W_E(A)$.

Finally the middle cup pairing pairs the $j$-summand only with the
$(2n-j)$-summand. In particular $W_E(A)$ is nondegenerate for this
pairing and is orthogonal to the kernel of $e_W$. The projection
formula gives
\[
 \int_X b^*w=\int_A b_*[X]\smile w
           =\int_A e_W\cl(b_*[X])\smile w.
\]
Thus \eqref{C13-period} makes this projected class nonzero, and the
preceding argument applies. $\square$

\framework{This is a concrete projector in the homological category
of abelian motives, followed by algebraic correspondences on its
two-dimensional realization. It gives the explicit form of the
projector-and-rank step C04 in this Weil-type case. The correspondence
framework is Grothendieck, cote~16, scans~43--46, and
\emph{Standard conjectures on algebraic cycles}, pp.~193--196.
It uses no Hodge conjecture to construct the projector. The remaining
condition is a cycle with nonzero projection. In dimension four the
supplied Markman result already establishes algebraicity; this
criterion is not presented as a new proof of that theorem.}

\subsection{THETA3-C14: which curve-product constructions miss the Weil space}\label{C14}

\paragraph{Statement.}
Let $P=C_1\times\cdots\times C_r$ be a product of smooth projective
connected complex curves and let $b:P\to A$ be a morphism to an abelian
variety. The cohomology class $b_*[P]$ belongs to the subalgebra
generated by rational divisor classes on $A$ (and is zero if the
image has dimension less than $r$). If $X$ is smooth projective,
$\pi:X\to P$ is generically finite of degree $d$, and $c=b\circ\pi$, then
$c_*[X]=d\,b_*[P]$ has the same property.

Suppose now $r=n$, $A$ is as in C13, its polarization $\theta$ has
Rosati involution restricting to complex conjugation on $E$, and
$\operatorname{NS}(A)_\Q=\Q\theta$. Then
\[
 e_W\cl(c_*[X])=0.
\]
Thus a construction factoring through a product of $n$ curves cannot
satisfy C13's nonzero-period criterion for this target.

\paragraph{Proof.}
If some $C_i$ has genus zero, every morphism from that factor to an
abelian variety is constant, and the asserted pushforward is zero.
Otherwise choose base points, let $J_i=J(C_i)$, and let
$\iota_i:C_i\to J_i$ be the Abel maps. The universal property of
the Albanese variety expresses $b$ as a translation followed by a
homomorphism
\[
 q:J_1\times\cdots\times J_r\longrightarrow A
\]
composed with $\prod_i\iota_i$. Translation acts trivially on
cohomology. If $g_i=\dim J_i$ and $\theta_i$ is its principal
polarization, the Abel-curve class is
\[
 (\iota_i)_*[C_i]=\frac{\theta_i^{g_i-1}}{(g_i-1)!}.
\]
For completeness, this formula follows by pairing both sides with
$H^2(J_i)$. In a symplectic integral basis $e_j,f_j$ of $H^1(C_i)$,
the Abel map identifies $H^1(J_i)$ with $H^1(C_i)$,
$\theta_i=\sum_j e_j\wedge f_j$, and
$\int_{J_i}\prod_j e_j\wedge f_j=1$. Pairing either displayed class
with $e_j\wedge f_k$ gives $\delta_{jk}$, and pairing with
$e_j\wedge e_k$ or $f_j\wedge f_k$ gives zero. These elements span
$H^2(J_i)$, so Poincar\'e duality proves the formula.

Consequently the class of $P$ in $\prod_iJ_i$ is a product of divisor
classes. Milne, \emph{Lefschetz classes on abelian varieties} (1999),
Corollary~5.5, printed p.~25, states that direct image under a regular
map of abelian varieties preserves such classes. Applying it to $q$
proves the first assertion. The degree formula
$\pi_*[X]=d[P]$ proves the second.

For the final assertion, all codimension-$n$ divisor polynomials on
$A$ are multiples of $\theta^n$. Write
$H^1(A,\C)=V_\tau\oplus V_{\bar\tau}$. The Rosati condition gives
$[a]^*\theta=\Nm_{E/\Q}(a)\theta$ for the element $a$ chosen in C13.
The three eigenvalues on $\bigwedge^2H^1(A)$ are
$a^2,a\bar a,\bar a^2$, which are distinct. It follows that
$\theta\in V_\tau\otimes V_{\bar\tau}\subset\bigwedge^2H^1(A,\C)$.
Thus $\theta^n$ belongs to the $j=n$ summand in C13, whereas the
Weil projector retains only $j=0,2n$. Since $n\ge1$, these summands
are different and $e_W\theta^n=0$. $\square$

\framework{These are Lefschetz classes and algebraic direct images,
in the setting of Milne's Corollaries~5.3 and~5.5, pp.~25--26;
homological and numerical equivalence agree on the Lefschetz
subalgebra there. The Grothendieck anchors are cote~16,
scans~43--46, and the correspondence and positivity framework of
cote~44, scans~7--8. The finding explains the scope of C12's explicit
degree-500 non-CM construction: its diagonal-product cycle is
algebraic, but it supplies no exceptional Weil class. It does not
assert that every map from a finite cover of a curve product factors
through that product. New Albanese factors or a genuinely different
correspondence remain possible.}


\subsection{THETA3-C15: transferring curve factors to ball covers, and the countability boundary}
\label{C15}

We work over $\overline{\Q}\subset\C$ unless otherwise indicated.  Write
$C\rightsquigarrow D$ when a connected finite \'{e}tale cover of the smooth
projective curve $C$ admits a nonconstant morphism to the smooth projective
curve $D$.  This relation is directed.  It is not permissible to reverse an
arrow merely because the two curves occur in a correspondence.

\paragraph{Curve-to-ball transfer.}
Let $Y$ be a smooth compact arithmetic ball surface with a surjective morphism
$f:Y\to C_0$.  If $C_0\rightsquigarrow D$, there is a connected finite
\'{e}tale cover $X\to Y$ and a surjective morphism $X\to D$.  Consequently
$\operatorname{Alb}(X)\to J(D)$ is surjective.  Every abelian quotient of
$J(D)$ is therefore a quotient of the Albanese of an arithmetic ball cover.

\emph{Proof.}
Choose a connected finite \'{e}tale cover $C'\to C_0$ and a nonconstant
map $C'\to D$.  Any connected component $X$ of $Y\times_{C_0}C'$ is finite
\'{e}tale over $Y$.  Its image in the connected variety $Y$ is nonempty,
open, and closed, so this projection is surjective.  The image of $X$ in
$C'$ is closed by properness.  It cannot be a point, since its image in
$C_0$ is all of $C_0$; thus it is $C'$.  Composition gives the asserted
surjection $h:X\to D$.  The universal covering of $X$ is again the ball,
and its lattice has finite index in that of $Y$, so arithmeticity persists.

The pullback of a nonzero holomorphic differential on $D$ is nonzero on
$X$: this can be checked at a point where $h$ is smooth.  Thus
$h^*:H^0(D,\Omega_D^1)\to H^0(X,\Omega_X^1)$ is injective.  The dual map
on the tangent spaces of Albanese varieties is surjective.  The resulting
homomorphism $\operatorname{Alb}(X)\to J(D)$ is therefore surjective.
This also proves the assertion about further quotients. \hfill$\square$

The relation $\rightsquigarrow$ is transitive.  Indeed, if
$C'\to C$ is finite \'{e}tale, $C'\to D$ is nonconstant, and
$D'\to D$ is finite \'{e}tale with $D'\to E$ nonconstant, take a connected
component of $C'\times_DD'$.  It is finite \'{e}tale over $C'$ and maps
surjectively to $D'$, hence to $E$.  This proves transitivity, including the
connectedness qualification, without identifying any of the maps as an
isomorphism.

\paragraph{The exact Bogomolov--Tschinkel direction.}
Let $C_0$ now be Schoen's curve $y^2=1-x^5$.  Schoen constructs the morphism
$Y\to C_0$ on pp.~133--134 of his 2014 paper.  Put
\[
 D_6:\quad w^6=z(1-z),
\]
with smooth projective normalization understood.  This curve has genus two.
Bogomolov--Tschinkel, \emph{Unramified correspondences}, Proposition~1.8,
p.~2, proves
\[
 C_0\rightsquigarrow D_6.
\]
The same statement is Proposition~1.5, p.~2, of their
\emph{On curve correspondences}.  Their statement applies to every
hyperelliptic curve of genus at least two over an algebraically closed field
of characteristic different from $2,3$.  It does not assert
$D_6\rightsquigarrow C_0$ or that $D_6$ is universal in characteristic zero.
Their universality theorem over $\overline{\mathbf F}_p$, $p\geq5$, is
Theorem~1.7 of the first paper and Theorem~1.4 of the second; that theorem
must not be transported to characteristic zero.

There is nevertheless a useful positive continuation.  In
\emph{Unramified correspondences}, Example~3.1, pp.~5--6, let
$E_*\to\mathbf P^1$ be the degree-three quotient of the elliptic curve
with $j=0$, branched over $0,1,\infty$.  The paper proves
$D_6\rightsquigarrow H$ for curves $H$ admitting a map to $\mathbf P^1$
with local ramification indices at most two and branch points contained in
the image of $E_*(\overline{\Q})_{\rm tors}$.  It also states density of
these families in the indicated Hurwitz spaces.  Transitivity and the
preceding transfer lemma give actual quotients $J(H)$ of Albaneses of
covers of the fixed arithmetic ball surface $Y$.  This is a composition
of arrows in their proved directions.  In C16 we give a direct version
starting from Schoen's non-CM elliptic curve, with a full ramification
proof, an explicit new factor, and infinitely many nonisogenous factors.

\paragraph{A precise conditional universality reduction.}
If $D_6$ is universal over $\overline{\Q}$ in the sense of
Bogomolov--Tschinkel, then every abelian variety over $\overline{\Q}$ is a
quotient of the Albanese of a connected finite \'{e}tale cover of $Y$.
To see this, let $A$ be an abelian variety.  There is a smooth projective
curve $H$ with a surjection $J(H)\to A$.  For $\dim A=1$ take $H=A$;
in larger dimension one may take a smooth sufficiently ample complete
intersection curve in $A$, whose map on $H^1$ is injective by weak
Lefschetz.  Universality gives $D_6\rightsquigarrow H$, and then
$C_0\rightsquigarrow H$.  Apply the transfer lemma and compose the two
surjective homomorphisms of abelian varieties.  The hypothesis here is a
characteristic-zero universality assertion, not a consequence of the
positive-characteristic theorem.  Even under that hypothesis, the conclusion
concerns $H^1$ source factors; it does not produce the higher-degree cycle
projection or spanning property required in C04 and C13.

\paragraph{Countability.}
For a fixed compact ball surface $Y$, only countably many simple complex
abelian isogeny classes occur in the Albaneses of its finite \'{e}tale
covers.  In particular, for each $g\geq1$, a very general simple complex
abelian variety of dimension $g$ is not such a factor.

\emph{Proof.}
The compact manifold $Y$ has finitely generated fundamental group.  For
each $n$, it has only finitely many subgroups of index $n$: each such
subgroup gives a transitive permutation action on $n$ letters, and the
images of a fixed finite generating set admit only finitely many choices
in the finite group $S_n$.  Thus the connected finite covers form a
countable set.  Each Albanese has finitely many simple isogeny factors,
by complete reducibility of abelian varieties.  A countable union of
finite sets remains countable.

For a fixed complex abelian variety, there are countably many lattices
commensurable with its period lattice.  Thus its isogeny class is
countable, and allowing polarizations still gives a countable set because
integral alternating forms on a fixed lattice are countable.  The
positive-dimensional moduli space of polarized abelian varieties of
dimension $g$ has uncountably many complex points; the countable set just
obtained is contained in a countable union of proper closed points.
The same exclusion holds on its simple locus.  Finally, a nonzero rational
Hodge map from the $H^1$ of a simple abelian variety to $H^1(X)$ would,
by the equivalence for polarizable Hodge structures of weight one, make
that variety an isogeny factor of $\operatorname{Alb}(X)$.  This gives the
claimed source obstruction. \hfill$\square$

The same proof applies to any countable collection of arithmetic towers.
In particular, the number-field Hermitian data and finite adelic levels
used in the present unitary constructions give such a collection: number
fields, finite matrices over them, and finite-level choices are countable.
This countability obstruction does not conflict with the conditional
statement over $\overline{\Q}$, where the set of algebraic varieties is
already countable.  Nor does countability preclude density; the explicit
families in the next subsection are countable and dense in the relevant
complex hyperelliptic moduli spaces.

\paragraph{Framework anchor.}
The transfer lemma uses actual curve maps and their Albanese realizations,
the weight-one case of the correspondence framework in Grothendieck's
letter of 12 August 1964, cote~44, scans~7--8.  It supplies source factors;
the existence of higher-degree algebraic projectors remains the separate
issue visible in cote~16, scans~44--46.

\subsection{THETA3-C16: infinitely many non-CM elliptic factors from torsion branch points}
\label{C16}

Schoen's construction supplies a connected curve $C$ with
\[
 C\longrightarrow C_0 \quad\hbox{finite \'{e}tale of degree }2,
 \qquad C\longrightarrow F \quad\hbox{of degree }2,
\]
where $C_0:y^2=1-x^5$ and
$F:s^2=x^4+x^3+x^2+x+1$; see his pp.~133--134.  The second double cover
is ramified.  Indeed, in the function field of $F$ it is obtained by
adjoining $\sqrt{1-x}$.  The two points over $x=1$ and the two points at
infinity have odd valuation of $1-x$, so they are branch points.
Choose one branch point $q_0$ as origin on $F$.  This does not change its
$j$-invariant, which Schoen computes as $2048$.

\paragraph{The torsion-branch transfer lemma.}
Let $E$ be any elliptic curve in characteristic zero, let $C\to E$ be a
connected double cover ramified at its origin, and put
$\pi:E\to E/\{\pm1\}\simeq\mathbf P^1$.  For an odd positive integer
$N$, define
\[
 S_N=\pi(E[N]),\qquad |S_N|=(N^2+1)/2.
\]
Suppose that $H$ is a smooth connected projective curve and
$\sigma:H\to\mathbf P^1$ is a finite morphism of degree $d$, all of whose
local ramification indices are at most two, with branch locus contained
in $S_N$.  Then a connected finite \'{e}tale cover $D\to C$ of degree
at most $dN^2$ admits a surjection $D\to H$ of degree at most four.

\emph{Proof.}
Form $C_N=C\times_{E,[N]}E$.  Its projection to $C$ is finite \'{e}tale
of degree $N^2$.  It is a smooth connected curve: as a double cover of
the second copy of $E$ it is still branched, at every point of $E[N]$,
and a split degree-two cover would have no branch points.  Its composite
map $a:C_N\to E\xrightarrow{\pi}\mathbf P^1$ has degree four.
Every local ramification index of $a$ over $S_N$ is even.  Namely, above
$P\in E[N]$ the map $C_N\to E$ has ramification index two; multiplication
by $N$ is \'{e}tale and carries $P$ to the specified branch point.
The further index of $\pi$ is either one or two.  Oddness of $N$ also
shows that the only fixed point of $[-1]$ in $E[N]$ is the origin,
which proves the displayed cardinality of $S_N$.

Normalize $C_N\times_{\mathbf P^1}H$ and choose a connected component
$D$.  Its projection to $C_N$ is \'{e}tale.  Here is the complete local
check.  Where $\sigma$ is unramified, base change is unramified.  At a
ramified point of $\sigma$, choose formal parameters so that its map is
$z=u^2$, after absorbing a unit into $u$.  The other map has the form
$z=v^{2e}$ after absorbing a unit into $v$, since its ramification index
is even and the residue field is algebraically closed of characteristic
zero.  The normalization of $u^2=v^{2e}$ is the disjoint union of the
two branches $u=\pm v^e$, each unramified over the $v$-disc.  This proves
the assertion also at points where several unramified sheets of $\sigma$
lie over the same branch value.  Therefore every connected component is
finite \'{e}tale over $C_N$, with degree at most $d$.  It maps onto $H$:
both maps in the original fiber product are finite and surjective, and
each component of its normalization has dimension one.  Its degree over
$H$ is at most the total base-change degree four.  Composing with
$C_N\to C$ proves both degree bounds. \hfill$\square$

Applied to Schoen's $C$, C15 gives a connected arithmetic ball cover
$X\to Y$ of degree at most $2dN^2$ whose Albanese surjects onto $J(H)$.
For any set $B\subset S_N$ of cardinality $2g+2$, its hyperelliptic
double cover $H_B\to\mathbf P^1$ satisfies the hypotheses, and hence
\[
 \deg(X/Y)\leq4N^2,\qquad
 \operatorname{Alb}(X)\twoheadrightarrow J(H_B).
\]
As $N$ ranges through odd integers, these $H_B$ are dense in the complex
hyperelliptic moduli space of each fixed genus $g\geq1$.  Indeed, odd-order
torsion is dense in the complex torus $E(\C)$, so its image in
$\mathbf P^1(\C)$ is dense.  Any finite ordered configuration of distinct
points can be approximated by distinct points in this image.  The finitely
many odd orders occurring in such an approximation divide a common odd
$N$.  Passing continuously from configurations to the moduli of their
double covers proves density.  The curves are all defined over
$\overline{\Q}$ when $E$ is.  Density by itself does not prove that a
specified countable family avoids the CM locus.  The following arguments
give that additional conclusion explicitly.

\paragraph{An explicit new non-CM isogeny class, with cover degree at most $36$.}
The elliptic curve
\[
 E:\quad y^2=x^3-30x-25
\]
has $j(E)=2048$ and is therefore isomorphic over $\overline{\Q}$ to
Schoen's $F$ with the chosen origin.  Its nonzero three-torsion points have
four distinct $x$-coordinates, the roots of
\[
 \phi(r)=r^4-60r^2-100r-300.
\]
For completeness, the usual duplication formula yields this equation
without requiring a division-polynomial table: the condition
$x(2P)=x(P)$, with $y(P)\neq0$, gives
$3x^4+6Ax^2+12Bx-A^2=0$ on $y^2=x^3+Ax+B$; it is exactly the
nonzero three-torsion condition because $2P=P$ is impossible for such a
point.  Put $A=-30$, $B=-25$ and divide by three.

Fix a root $r$ of $\phi$, omit it from the four finite branch points, and
include infinity.  The resulting elliptic curve $H_r$ is
\[
 y^2=x^3+rx^2+(r^2-60)x+(r^3-60r-100).
\]
Its cubic is $\phi(x)/(x-r)$, so it has three distinct roots.  After
shifting $x$ by $-r/3$, its short Weierstrass coefficients are
\[
 A_r=\frac23r^2-60,\qquad
 B_r=\frac{20}{27}r^3-40r-100.
\]
Substitution in $j=1728\cdot4A_r^3/(4A_r^3+27B_r^2)$ and reduction
modulo $\phi(r)$ give
\[
 j(H_r)=-\frac{512}{405}
       (r^3+24r^2-130r-1200).
\]
Each of the four values is a root of the primitive polynomial
\begin{align*}
 P(J)={}&19683J^4-40310784J^3-5096079360J^2\\
        &+37203006717952J+6412351813189632.
\end{align*}
Both identities are polynomial identities over $\Q$ after reduction by
$\phi$; the accompanying standard-library certificate
\nolinkurl{check_C16_torsion_family.py} verifies them using exact rational
arithmetic, rather than floating-point root approximations.

Here is an explicit irreducibility certificate.  The monic reduction of
$P$ modulo $11$ is
\[
 P_0(J)=J^4-2J^3-5J^2+5J-5.
\]
In $\mathbf F_{11}[J]/(P_0)$ one has
\[
 J^{121}=5J^3-3J^2-5J+4,
 \qquad J^{14641}=J,
\]
and in $\mathbf F_{11}[J]$ one has the exact identity
\begin{align*}
 1={}&(-5J^2-5J-3)P_0(J)\\
 &+(J^3+4J^2-5J+2)(5J^3-3J^2-6J+4).
\end{align*}
The identity $J^{14641}=J$ says that all irreducible factors of $P_0$
have degree dividing four; the Bezout identity excludes factors of degree
one or two.  Thus $P_0$ is irreducible, and so is $P$ over $\Q$.
Since $P$ is primitive with leading coefficient $19683=3^9\neq1$,
its roots are not algebraic integers.  Every $H_r$ is therefore non-CM.
Moreover, it is not isogenous to $E$.  Indeed, a cyclic isogeny gives a
monic integral modular equation in the target $j$-invariant, so integrality
of the source $j$ implies integrality of the target; any isogeny factors
into cyclic isogenies.  Since $j(E)=2048$ is integral, the assertion follows.
The same argument applies place by place.  The standard source statements
are Milne, \emph{Modular Functions and Modular Forms}, version~1.31,
Theorem~6.1 and its proof, pp.~88--90, the cyclic-isogeny interpretation
in Lemma~8.5 and the following discussion, p.~99, and CM integrality in
Corollaries~12.15--12.16, p.~125.

All four branch points of $H_r$ belong to $S_3$, so the torsion-branch
lemma and C15 give the claimed ball cover of degree at most
$4\cdot3^2=36$.  This provides a non-CM elliptic isogeny class different
from Schoen's original one inside the same arithmetic covering tower.
No assertion that the four $H_r$ are mutually nonisogenous is needed.

\paragraph{Infinitely many pairwise nonisogenous non-CM factors.}
There are infinitely many pairwise nonisogenous non-CM elliptic curves
$E_1,E_2,\ldots$ over $\overline{\Q}$, each a quotient of the Albanese of
a connected finite \'{e}tale cover of Schoen's fixed ball surface $Y$.
For every finite initial set, one common connected cover of $Y$ has
$E_1\times\cdots\times E_m$ as an Albanese quotient.

\emph{Proof.}
We first produce a usable ordinary prime outside any prescribed finite
set.  The discriminant of $E:y^2=x^3-30x-25$ is
$2^4 3^6 5^3$.  If $p\equiv1\pmod{20}$, then $p\geq41$, the reduction
is good, and all its two-torsion is rational over $\mathbf F_p$:
\[
 x^3-30x-25=(x+5)(x^2-5x-5),
\]
whose remaining roots are $(5\pm3\sqrt5)/2$.  The field $\mathbf F_p$
contains the fifth roots of unity, and hence $\sqrt5$, because
$p\equiv1\pmod5$.  Consequently $4$ divides $\#E(\mathbf F_p)$ and
\[
 a_p=p+1-\#E(\mathbf F_p)\equiv2\pmod4.
\]
The Hasse bound gives $|a_p|\leq2\sqrt p<p$, so $a_p\not\equiv0\pmod p$.
This implies ordinary reduction.  One can check the last implication
directly: if $\varphi$ is Frobenius, its dual satisfies
$\widehat\varphi=[a_p]-\varphi$ and has degree $p$; its differential is
the nonzero scalar $a_p$.  It is therefore separable of degree $p$.
Since $[p]=\widehat\varphi\varphi$ and Frobenius is bijective on geometric
points, $\overline E[p](\overline{\mathbf F}_p)$ has $p$ points.
The facts used here are explicitly stated and proved in Stoll,
\emph{Elliptic Curves}, summer~2023 notes, Lemma~11.16, p.~42,
Theorem~12.1, p.~44, and Theorem~13.2, pp.~48--49.

There are infinitely many primes $p\equiv1\pmod{20}$, with an elementary
proof sufficient here.  Given a finite set of primes, choose a positive
integer $M$ divisible by $20$ and by every prime in that set.  Any prime
$p$ dividing
\[
 \Phi_{20}(M)=M^8-M^6+M^4-M^2+1>1
\]
is outside the set and is different from $2,5$.  The identity
$(M^2+1)\Phi_{20}(M)=M^{10}+1$ implies $M^{10}=-1$ modulo $p$.
The order of $M$ modulo $p$ is therefore four or twenty.  In the former
case $M^2=-1$, which gives $\Phi_{20}(M)=5$ modulo $p$, a contradiction.
Thus the order is twenty and $p\equiv1\pmod{20}$.

Fix such a prime and an embedding $\overline{\Q}\hookrightarrow
\overline{\Q}_p$.  The reduction map on $p$-torsion has a kernel of
cardinality $p$ and image $\overline E[p](\overline{\mathbf F}_p)$ of
cardinality $p$.  We spell out why no lifting assumption is hidden here.
The finite flat group scheme $E[p]$ over a good-reduction discrete
valuation ring has rank $p^2$.  After a finite unramified extension its
$p$ distinct special-fiber points are rational.  Their idempotent
decomposition lifts over this henselian ring.  Each component has rank
$p$, since translations make all special-fiber local lengths equal.
Its characteristic-zero generic fiber is \'{e}tale with $p$ geometric
points.  Every special point therefore lifts, and exactly $p$ generic
points reduce to the origin.

Choose a nonzero $P\in E[p]$ reducing to the origin, and choose
$Q,R\in E[p]$ whose reductions are nonzero and distinct modulo sign.
There are $(p-1)/2\geq2$ available nonzero sign classes.  The coordinates
satisfy
\[
 v_p(x(P))<0,\qquad
 v_p(x(Q)),v_p(x(R))\geq0,\qquad
 v_p(x(R)-x(Q))=0.
\]
For the first inequality, the integral local parameter $t=-x/y$ at the
origin gives $x=t^{-2}$ times an integral unit, and $v_p(t(P))>0$.
The others follow by reduction at finite distinct $x$-coordinates.

Let $E'_p$ be the double cover of $\mathbf P^1$ branched at
$\infty,x(P),x(Q),x(R)$.  Its branch cross ratio can be chosen as
\[
 \lambda=\frac{x(R)-x(Q)}{x(P)-x(Q)},\qquad v_p(\lambda)>0.
\]
The Legendre formula gives
\[
 j(E'_p)=256\frac{(1-\lambda+\lambda^2)^3}
                         {\lambda^2(1-\lambda)^2},
 \qquad v_p(j(E'_p))=-2v_p(\lambda)<0.
\]
This is a non-CM elliptic curve, by CM integrality.  All its branch points
lie in $S_p$, so the torsion-branch lemma gives a ball cover of degree at
most $4p^2$ whose Albanese surjects onto it.

Construct the sequence inductively.  After finitely many $E_i$ have been
chosen, the set of rational primes at which any conjugate of one of their
$j$-invariants is nonintegral is finite: each $j(E_i)$ is an algebraic
number.  Choose the next prime $p\equiv1\pmod{20}$ outside that set and
apply the construction.  All earlier $j(E_i)$ are integral at every
place above this new $p$, whereas $j(E'_p)$ is nonintegral at the chosen
place.  The local form of the monic modular-equation argument therefore
excludes an isogeny to every earlier $E_i$.  It also excludes an isogeny
to Schoen's original integral-$j$ elliptic curve.  This proves the infinite
assertion.

Finally, any finite list of the resulting connected finite \'{e}tale
ball covers has a common connected finite \'{e}tale cover, obtained by
taking a connected component of their fiber product over $Y$.  Projection
to each connected factor is surjective.  The maps to the $E_i$ persist,
and their distinct simple isogeny classes imply that the direct sum of
their $H^1$ injects into the $H^1$ of this common cover.  By complete
reducibility, the product $E_1\times\cdots\times E_m$ is an Albanese
quotient. \hfill$\square$

\paragraph{Nonzero periods with two of the new factors.}
The two independent curve projections of C12 give a generically finite
map $Y\to C_0^2$ of degree $125$.  For two curves $H_1,H_2$ furnished by
the torsion-branch lemma, choose the corresponding connected covers
$D_i\to C_0$; each $D_i\to H_i$ has degree at most four.  A connected
component $X$ of
\[
 Y\times_{C_0^2}(D_1\times D_2)
\]
is a finite \'{e}tale arithmetic ball cover of $Y$ and maps dominantly to
$H_1\times H_2$, with degree
\[
 1\leq d\leq125\cdot4\cdot4=2000.
\]
Indeed, its degree over $D_1\times D_2$ is at most the total generic
base-change degree $125$, and the remaining product map has degree at
most $16$.  If the $H_i$ are elliptic, choose differentials $\alpha_i$
normalized by $i\int_{H_i}\alpha_i\wedge\overline\alpha_i=1$.
Writing $a_i:X\to H_i$ for the two maps, the projection formula yields
\[
 \int_X a_1^*\alpha_1\wedge a_2^*\alpha_2
       \wedge\overline{a_1^*\alpha_1}
       \wedge\overline{a_2^*\alpha_2}=d>0.
\]
Thus one can obtain fourfold periods with factors from infinitely many
non-CM isogeny classes, with a positive integral degree in each instance.
The map to $H_1\times H_2\times H_1\times H_2$ repeats the two
projections; its image cycle is $d$ times the product of the two diagonals.
These are divisor-generated classes.  The construction substantially
enlarges the non-CM $H^1$ sources, but C14 still prevents these particular
product-of-curves image cycles from detecting the generic exceptional
Weil component.  The geometric factors constructed here are not supplied
by the CM character packets in the congruence $H^1$ statement of C01.

\paragraph{Framework anchor.}
The new Albanese quotients and image surfaces are actual algebraic
correspondences, as required by the realization/polarization viewpoint
of cote~44, scans~7--8.  Their explicit divisor-generated image classes
must be distinguished from the additional algebraic inverse and projector
classes discussed in cote~16, scans~43--46.  The infinite family enlarges
the available non-CM realizations without resolving that higher-degree
existence question.


\subsection{THETA3-C20: a simple non-CM abelian surface in one ball covering tower}\label{C20}

The elliptic factors in C16 can be replaced by a geometrically
simple abelian surface with no extra endomorphisms. This supplies a
non-CM simple weight-one source of dimension two, rather than a
product of the previously constructed elliptic sources.

\paragraph{Statement.}
Let
\[
 E:\ y^2=x^3-30x-25,\qquad
 \phi(x)=x^4-60x^2-100x-300.
\]
For any nonzero point $P\in E[7](\overline{\Q})$, put $a=x(P)$
and let $H_a$ be the smooth projective curve with affine equation
\[
 H_a:\quad y^2=(x-a)\phi(x).
 \tag{C20a}
\]
Then $H_a$ has genus two and
\[
 \operatorname{End}_{\overline{\Q}}J(H_a)=\Z.
 \tag{C20b}
\]
In particular $J(H_a)$ is geometrically simple and is not CM.
For Schoen's fixed arithmetic ball surface $Y$ in C15--C16, there
is a connected finite \'etale cover $X\to Y$ with
\[
 \deg(X/Y)\leq1764,\qquad
 \operatorname{Alb}(X)\twoheadrightarrow J(H_a).
 \tag{C20c}
\]
The surjection is understood over $\overline{\Q}$ and up to the
harmless choice of origins in the Albanese construction; its
associated $H^1$ inclusion is a rational Hodge morphism.

\paragraph{The division polynomial and its conjugate reductions.}
The four roots of $\phi$ are the nonzero three-torsion
$x$-coordinates, by C16. Set
\begin{align*}
 u&=x^3-30x-25,\\
 b&=x^6-150x^4-500x^3-4500x^2-3000x+22000,\\
 \psi_3&=3\phi,\qquad
 \psi_5=32u^2b-\psi_3^3,\qquad
 \psi_7=\psi_5\psi_3^3-128u^2b^3.
\end{align*}
The elementary division-polynomial recursion
$\psi_{2m+1}=\psi_{m+2}\psi_m^3-
\psi_{m-1}\psi_{m+1}^3$, with $\psi_2=2y$ and
$\psi_4=4yb$, gives these formulas. The exact source is
Silverman, \emph{The Arithmetic of Elliptic Curves},
second edition (2009), Chapter~III, Exercise~3.7,
pp.~105--106, especially parts (b), (d) and (f).
Thus $\psi_7$ has degree
24, leading coefficient seven, and its roots are exactly the
$x$-coordinates of the nonzero seven-torsion points.
One may equally obtain the same equation by the addition law
and the equality $x(4P)=x(3P)$, excluding the smaller-order
denominators; the displayed recursion fixes the normalization.

Here is an explicit finite-field irreducibility certificate.
The monic reduction of $\psi_7$ modulo 19 is
\begin{align*}
 P_{19}(x)={}&x^{24}-9x^{22}-4x^{21}-9x^{20}
       +8x^{19}+6x^{18}-4x^{16}+5x^{15}\\
 &+8x^{14}-2x^{13}-4x^{12}+7x^{11}-7x^{10}
       -9x^8-4x^7+7x^6\\
 &+4x^5+2x^4-4x^3+8x^2-7.
\end{align*}
Successive nineteenth powers modulo this polynomial give
\[
 x^{19^{24}}\equiv x,\qquad
 \gcd(P_{19},x^{19^8}-x)=
 \gcd(P_{19},x^{19^{12}}-x)=1.
 \tag{C20d}
\]
For an explicit check of the large exponents, the ascending
coefficient lists of the two intermediate remainders are
\[
\begin{split}
 r_8={}&(15,13,2,16,7,18,2,6,14,16,7,8,\\
       &\hspace{14mm}10,3,17,11,17,17,6,18,18,1,12,16),\\
 r_{12}={}&(8,18,13,9,13,1,6,13,4,16,14,12,\\
       &\hspace{14mm}14,3,11,3,12,17,5,0,12,6,0,3).
\end{split}
\]
Euclid's algorithm applied to $P_{19}$ and $r_j-x$ gives
the two unit gcds in (C20d). The first equality says that all
irreducible factors have degree dividing 24. Every proper
divisor of 24 divides 8 or 12, so the two gcds exclude all
proper factors. Hence $\psi_7$ is irreducible over $\Q$.
The independent integer-polynomial certificate
\path{documentation/check_C20_independent.py} implements this
recursion, modular powering and Euclidean division with exact
integer arithmetic. It does not use a numerical factorization.

Modulo 37, $\psi_7$ has the simple roots $7,20,35$.
In particular $K=\Q(a)$ has embeddings into $\Q_{37}$ taking
$a$ to Hensel lifts of 7 and 20. These give two places of the
same field $K$, with residue field $\F_{37}$, and two good
reductions of the same curve (C20a). No choice of unrelated
characteristic-zero curves is being compared.

\paragraph{The two point counts and absolute simplicity.}
For the two reductions, direct enumeration gives
\[
\begin{array}{c|r|r|l}
 a\bmod37 & \#H(\F_{37}) & \#H(\F_{37^2})
                    &\text{Frobenius polynomial}\\ \hline
 7 &35&1297&T^4-3T^3-32T^2-111T+1369\\
 20&37&1361&T^4-T^3-4T^2-37T+1369 .
\end{array}
 \tag{C20e}
\]
The polynomials $(x-7)\phi(x)$ and $(x-20)\phi(x)$ are
squarefree modulo 37, so these are smooth genus-two reductions.
The counts include the one point at infinity of the odd-degree
hyperelliptic equation. To specify completely the finite
calculation, use
$\F_{37^2}=\F_{37}[w]/(w^2-2)$; 2 is a nonsquare modulo 37.
For $z=r+sw$ evaluate $(z-a)\phi(z)$ by polynomial arithmetic
and count its square roots in this field. The sum of these
counts over the $37^2$ choices of $(r,s)$, plus one, gives the
second count. An independent implementation uses
$\chi_{\F_{37^2}}(v)=
\chi_{\F_{37}}(\operatorname{Norm}v)$ for $v\ne0$.
The two implementations agree and their full outputs are
included with the paper.

If $A_1=37+1-\#H(\F_{37})$, then
\[
 A_2=\frac{\#H(\F_{37^2})-37^2-1+A_1^2}{2},
 \qquad
 f(T)=T^4-A_1T^3+A_2T^2-37A_1T+37^2.
\]
This derives the two polynomials in (C20e) from the counts,
rather than assuming them. Both are irreducible modulo 5:
in each quartic quotient $T^{5^4}=T$, while
$\gcd(f,T^{5^2}-T)=1$. The respective ascending coefficient
lists of $T^{25}\bmod f$ are $(3,1,4,2)$ and $(1,2,3,2)$.
Thus both quartics are irreducible over $\Q$.
Their middle coefficients $-32$ and $-4$ are units at 37,
so the two Jacobians are ordinary and simple.

We use the exact finite-field criterion of Howe--Zhu,
\emph{On the existence of absolutely simple abelian varieties
of a given dimension over an arbitrary field},
arXiv:math/0002205v1, Theorem~6, p.~5, with proof pp.~5--6.
For $T^4+cT^3+dT^2+37cT+37^2$, its four exceptional
conditions are
$c=0$, $c^2=37+d$, $c^2=2d$ and $c^2=3d-111$.
None holds for $(c,d)=(-3,-32)$ or $(-1,-4)$.
Hence both reductions are absolutely simple.
By the same work, \S2 p.~3 and Proposition~3(2) with
Lemma~4 p.~4, their geometric endomorphism algebras are
exactly the quartic CM fields
\[
 K_1=\Q[T]/(T^4-3T^3-32T^2-111T+1369),\quad
 K_2=\Q[T]/(T^4-T^3-4T^2-37T+1369).
 \tag{C20f}
\]
Indeed absolute simplicity in the ordinary case keeps
$\Q(\pi^r)=\Q(\pi)$ over every finite residue extension.

\paragraph{Excluding every extra geometric endomorphism.}
The real subfield of a quartic ordinary Frobenius field is
generated by $s=\pi+37/\pi$. Its equation is
$s^2+c s+(d-74)=0$ for the coefficients just used.
Consequently
\[
 K_1^+=\Q(\sqrt{433}),\qquad
 K_2^+=\Q(\sqrt{313}).
 \tag{C20g}
\]
The polynomial discriminants, computed by the resultant with
the derivative, are
\[
 \operatorname{disc}(f_1)=
       2^4\,3^3\,37^2\,433^2,\qquad
 \operatorname{disc}(f_2)=
       2^4\,3^3\,11\,37^2\,313^2.
 \tag{C20h}
\]
Both are nonsquares. A quartic CM field containing an
imaginary quadratic subfield would be the compositum of
that field and its real quadratic subfield, hence a
biquadratic Galois extension. A quartic field containing
two different quadratic subfields is likewise
biquadratic. In a biquadratic extension the Galois
permutations of four embeddings are all even, so the
discriminant is a square. The discriminant of a defining
polynomial differs from the field discriminant by a
square. Thus each $K_i$ has exactly one quadratic
subfield, namely the one displayed in (C20g).
The two fields have no common nontrivial subfield up to
$\Q$-isomorphism: degree two is excluded by (C20g), and
an isomorphism in degree four would identify their unique
quadratic subfields.

Put $B=J(H_a)$. Every geometric endomorphism is defined
over some finite extension of $K$. At either of the two
places above 37 it extends to the good abelian scheme and
specializes injectively. The exact source is Conrad,
\emph{Gross--Zagier Revisited}, MSRI Publications~49
(2004), Theorem~2.1(2), printed p.~70.
After any such finite extension the geometric targets
remain the fields $K_1,K_2$, as proved above. Passing to
all geometric endomorphisms therefore embeds the finite
$\Q$-algebra
$\operatorname{End}_{\overline\Q}(B)\otimes\Q$
into each field. It is itself a field: a finite
$\Q$-subalgebra of a number field has no zero divisors,
and multiplication by a nonzero element is a bijection
of its finite-dimensional vector space. Its degree
divides four and is therefore one, two or four.
The last two degrees have just been excluded.
We obtain $\operatorname{End}_{\overline\Q}^0(B)=\Q$.
The integral endomorphism ring is an order in this
field containing $\Z$, hence is $\Z$. This proves
(C20b). A nontrivial abelian subvariety would give a
nontrivial rational idempotent by Poincar\'e
reducibility, so $B$ is geometrically simple.
A CM abelian surface would have a commutative
endomorphism algebra of dimension four, which is
impossible here.

\paragraph{Realization on the fixed ball covering tower.}
The branch points of (C20a) are infinity, the four
$x$-coordinates of $E[3]\setminus\{0\}$ modulo sign,
and $x(P)$ for a seven-torsion point. They are distinct:
an equality of the last coordinate with a three-torsion
coordinate would identify $P$ with a point of order
dividing three, up to sign. All six belong to
$S_{21}=\pi(E[21])$ of C16. The torsion-branch lemma
therefore gives a connected finite \'etale cover
$D\to C$ of degree at most $2\cdot21^2$ and a finite
map $D\to H_a$ of degree at most four. Since
$C\to C_0$ is \'etale of degree two, C15 supplies
a connected ball cover of degree at most
$4\cdot21^2=1764$ with Albanese quotient $J(H_a)$.
This proves (C20c). \hfill$\square$

The graph of the map to $H_a$ and the Abel--Jacobi
correspondence give the asserted rational $H^1$
realization. Thus this finding belongs to the
Jacobian-source and algebraic-correspondence programme
of Grothendieck's cote~44, scans~7--8; the
polarization and correspondence interpretation is
that of cote~16, scans~43--46. It supplies a new
simple source in this particular covering tower.
It does not identify a nonzero exceptional Weil
projection of an image cycle: that remains the
separate extraction problem in C13--C14.


\subsection{THETA3-C24: non-CM Weil-type sources in every even dimension}\label{C24}

The source problem in C13 can be made concrete in every even dimension.
The following construction supplies non-CM Jacobians with a balanced
$\Q(i)$-action, all as quotients of Albaneses in the single arithmetic
ball covering tower of C15--C16. It does not assert that their Weil
classes are independent of divisor classes, or that the ball surface
already supplies a nonzero projection onto their Weil spaces.

\paragraph{Statement.}
Let $n\geq2$, and let $M$ be the least power of two satisfying
$M^2\geq8n+4$. There are a smooth projective curve $H_n$ over
$\overline\Q$ of genus $2n$, a principally polarized abelian variety
$A_n=J(H_n)$, and a connected finite \'etale cover $X_n\to Y$ of
Schoen's fixed arithmetic ball surface, such that
\[
 \deg(X_n/Y)\leq196M^2<6272n+3136,\qquad
 \operatorname{Alb}(X_n)\twoheadrightarrow A_n.
 \tag{C24a}
\]
The variety $A_n$ is not of CM type. It has an action of $K=\Q(i)$
whose two embeddings each have multiplicity $n$ on
$H^0(A_n,\Omega^1)$, and whose Rosati involution is complex
conjugation. Thus
\[
 W_K(A_n)=\bigwedge_K^{2n}H^1(A_n,\Q)
\]
is an actual rational two-dimensional Hodge subspace of type $(n,n)$.
The construction gives a place above seven where the geometric
semistable toric rank of $A_n$ is exactly two. In dimension four
one may take $M=8$, so the covering degree in (C24a) is at most
$12544$.

\paragraph{Even-level torsion transfer.}
The torsion-branch lemma C16 remains valid for every positive integer
$N$, with precisely the same degree bounds. For even $N$ its branch
set has cardinality
\[
 |\pi(E[N])|=(N^2+4)/2.
 \tag{C24b}
\]
Indeed, multiplication by $N$ is \'etale in characteristic zero
regardless of parity. In the cover
$C_N=C\times_{E,[N]}E$, the map $C_N\to E$ is ramified at every
point of $E[N]$, because $C\to E$ is ramified at its origin.
Thus every ramification index of the composite to
$E/\{\pm1\}$ above $\pi(E[N])$ is even, including the indices
four at points fixed by $[-1]$. The local normalization
$u^2=v^{2e}$ used in C16 is unchanged. It gives a connected
finite \'etale cover of $C$ of degree at most $dN^2$ dominating
any degree-$d$ curve cover with ramification at most two and
branch contained in this set. The four fixed points of $[-1]$
on $E[N]$ give (C24b) by counting its orbits. Applying C15 and
Schoen's degree-two cover gives a ball cover of degree at most
$2dN^2$; for a hyperelliptic cover this is $4N^2$.

\paragraph{The torsion points and the equation.}
Use the elliptic curve from C16,
\[
 E:\quad y^2=x^3-30x-25,
 \qquad T=(-5,0),\qquad c=3\sqrt5.
\]
Fix embeddings of $\overline\Q$ in $\C$ and in
$\overline\Q_7$. The curve has good reduction at seven:
its displayed discriminant is $1458000$, a $7$-adic unit.
Choose a nonzero point $R\in E[7](\overline\Q)$ reducing
to the identity at this place. Such a point exists. There are
$49$ geometric points of $E[7]$ in characteristic zero,
whereas the geometric seven-torsion in the special elliptic
curve has at most seven points: the differential of $[7]$
vanishes, so its inseparable degree is at least seven and
its separable degree is at most seven. Properness gives
reduction of all these torsion points, and its kernel therefore
contains a nonzero point. In particular $R$ belongs to the
formal group at the chosen place.

On the Kummer line set
\[
 z(P)=\frac{x(P)+5-c}{x(P)+5+c}.
 \tag{C24c}
\]
The addition law gives
\[
 x(P+T)+5=\frac{45}{x(P)+5},\qquad
 z(P+T)=-z(P),\qquad z(-P)=z(P).
 \tag{C24d}
\]
For example the first identity follows by using slope
$y(P)/(x(P)+5)$ in the chord from $P$ to $T$ and using
$x^3-30x-25=(x+5)(x^2-5x-5)$. The two fixed points
of the involution on the Kummer line have coordinates $z=0,\infty$;
their preimages satisfy $2P=T$ and belong to $E[4]$.

The four-element group generated by $P\mapsto-P$ and
$P\mapsto P+T$ acts on $E[M]$. Exactly eight points have
nontrivial stabilizer: the four in $E[2]$ and the four satisfying
$2P=T$. These two sets are disjoint, and translation by $T$
has no fixed point. Every remaining orbit has four elements.
There are therefore
\[
 (M^2-8)/4\geq2n-1
\]
such orbits. Choose representatives $U_1,\ldots,U_{2n-1}$
of distinct ones. Put
\[
 P_1=U_1,\qquad P_2=U_1+R,\qquad
 P_j=U_{j-1}\quad(3\leq j\leq2n),\qquad a_j=z(P_j).
 \tag{C24e}
\]
These points lie in $E[7M]$. They give pairwise distinct
branch values
\[
 0,\infty, a_1,-a_1,\ldots,a_{2n},-a_{2n}.
 \tag{C24f}
\]
For the $U_j$ this is precisely the orbit condition.
The point $U_1+R$ cannot belong to the four-element orbit
of any $U_j$: such an equality would put the nonzero
seven-torsion point $R$ in $E[M]$, contrary to $(7,M)=1$.
It likewise cannot be in either excluded set. The branch
values in (C24f) all belong to $z(E[7M])$.

Define the smooth projective model of
\[
 H_n:\qquad v^2=z\prod_{j=1}^{2n}(z^2-a_j^2).
 \tag{C24g}
\]
The distinctness just proved gives a smooth hyperelliptic curve
with $4n+2$ branch points, hence genus $2n$.
Even-level transfer with $N=7M$ gives
$\deg(X_n/Y)\leq4(7M)^2=196M^2$ and the asserted
Albanese quotient. Minimality of $M$ gives
$(M/2)^2<8n+4$, proving the strict final bound in (C24a).

\paragraph{The balanced action.}
The automorphism
\[
 \alpha:(z,v)\longmapsto(-z,iv)
\]
has square equal to the hyperelliptic involution. On the
Jacobian the latter acts as $[-1]$, because for a Weierstrass
base point the divisors $P+\iota(P)$ are all linearly equivalent
to twice that base point. Thus $i\mapsto\alpha^*$ gives
a faithful $\Q(i)$-action on $A_n$. Its pullback on the
holomorphic basis is
\[
 \alpha^*\left(z^j\frac{dz}{v}\right)
   =i(-1)^j z^j\frac{dz}{v},\qquad 0\leq j<2n.
\]
There are $n$ even and $n$ odd exponents, giving the claimed
signature. A curve automorphism preserves the canonical
principal polarization of its Jacobian; hence its Rosati
adjoint is its inverse. Here $\alpha^{-1}=-\alpha$ on
$H^1$, so this involution restricts to complex conjugation
on $K$. The two $K$-eigenspaces of $H^1_\C$ each have
Hodge numbers $(n,n)$. Their top exterior powers therefore
have type $(n,n)$, as asserted for $W_K(A_n)$.

\paragraph{A two-node reduction certifying that the source is not CM.}
Pass to a finite extension of $\Q_7$ containing all the
displayed points and constants, with valuation ring $\mathcal O$.
Prime-to-seven torsion reduces injectively on the good model
of $E$. Moreover $c$ and $2c$ are units. Consequently
$a_1,a_3,\ldots,a_{2n}$ are units with pairwise distinct
reductions up to sign, none equal to zero or infinity.
The same assertions hold for $a_2$ except that
\[
 \bar a_2=\bar a_1=:a\ne0,\qquad a_2\ne\pm a_1.
\]
Write $b_j=\bar a_j$ for $j\geq3$. A proper flat double
cover of $\mathbf P^1_{\mathcal O}$ is obtained by gluing
(C24g) to its chart $u=1/z$, $w=v/z^{2n+1}$, whose equation is
\[
 w^2=u\prod_{j=1}^{2n}(1-a_j^2u^2).
\]
Each chart is free of rank two over the indicated affine
line. They give a finite cover of $\mathbf P^1$, so it is
proper and flat. The special fibre is
\[
 v^2=z(z^2-a^2)^2\prod_{j=3}^{2n}(z^2-b_j^2).
 \tag{C24h}
\]
It is geometrically irreducible: the squarefree part on the
right has the simple root $z=0$ and is not a square.
Its only singularities are the two points $z=a,-a$, each
an ordinary node. At each, the remaining factor is a unit
and $2$ is invertible, so the tangent cone is the product
of two distinct linear factors. The points at zero and
infinity and the other branch points are smooth.
The normalization is
\[
 v_0^2=z\prod_{j=3}^{2n}(z^2-b_j^2),
 \qquad v_0=v/(z^2-a^2),
\]
of genus $2n-2$. In particular the fibre is a stable
irreducible curve with two nodes.

For completeness, allowing a singular total space does not
alter the toric-rank calculation. In the completed local
ring near either node, extract a square root of the unit
factor and put $\delta=a_2-a_1$. The remaining equation
is equivalent to
\[
 XY=\delta^2/4
\]
up to a unit, by centering $z$ at $(a_1+a_2)/2$ and
factoring the difference of two squares. A finite unramified
extension suffices to split the two branches. If
$v(\delta)=e$, resolution of $XY=\pi^{2e}\cdot\mathrm{unit}$
inserts a chain between those two branches. This can also
be seen inductively in the two blow-up charts, which reduce
the exponent by one. In the dual graph it only subdivides
the corresponding loop. After any finite ramified base
extension the exponent is multiplied by its ramification
index and the same description applies. The geometric
dual graph therefore has first Betti number two over
every such extension.

Apply the N\'eron--Picard identification in SGA~7~I,
Expos\'e~IX, Theorem~12.1(d), formulas (12.1.10)--(12.1.12), printed pp.~467--468,
to a regular semistable resolution. The torus of the
special-fibre Picard variety has character lattice the
first homology of this dual graph, as in IX~12.3,
especially (12.3.6)--(12.3.7), printed pp.~471--472. Equivalently, gluing line
bundles across the two independent nodes gives two
independent $\mathbf G_m$ parameters. Hence the Jacobian
has toric rank exactly two, and cannot acquire good
reduction after any finite extension. Serre--Tate,
\emph{Good reduction of abelian varieties} (1968),
Theorem~6(a), printed p.~503, gives potential good
reduction for CM abelian varieties. For a product of
simple CM factors the same conclusion follows after a
common finite extension and isogeny, so it applies to
the general meaning of CM type used here. Thus $A_n$
is not CM. \hfill$\square$

\paragraph{What this advances, and what remains.}
For every $n\geq2$ there is now an explicit finite torsion
prescription for a non-CM source with the precise balanced
quadratic-field action required by C13, within one fixed
arithmetic ball tower. All of its weight-one realization
comes from the actual Albanese quotient (C24a). The
projector in C13 consequently applies to this specific
$W_K(A_n)$, without an additional source-existence
hypothesis. Geometric simplicity, an exact endomorphism
ring, and a non-divisor Weil class have not been inferred
from the nodal calculation. The remaining cycle problem
is still to construct an $n$-dimensional correspondence
with nonzero projection to this Weil space. Taking
products of the supplied curve map produces the
Lefschetz classes described in C14, and does not resolve
that projection problem when the Weil space is
exceptional. For $n=2$, the cited Markman theorem already
gives algebraicity; the new conclusion here is the
specified non-CM source in the ball tower.

\framework{The Albanese maps and rational Weil projector are actual
correspondences in the setting of Grothendieck's cote~44,
scans~7--8, and cote~16, scans~43--46. The new non-CM
certificate takes place directly in his theory of N\'eron
models and monodromy: SGA~7~I, IX~12.1 and IX~12.3,
with the dual-graph torus and its persistent rank. It
supplies non-CM weight-one motives and their specified
Weil submotives modulo homological equivalence; it does
not identify an unknown Hodge class with an effective
special-fibre intersection.}


\subsection{THETA3-C25: a concrete exceptional Weil source with endomorphism ring $\Z[i]$}\label{C25}

The preceding non-CM sources can be sharpened to a geometrically
simple Weil-type fourfold with no extra endomorphisms. Its Weil
space is entirely outside the divisor algebra. Thus the
projector criterion C13 now has a specified exceptional target
in the same arithmetic ball covering tower.

\paragraph{The algebraic torsion prescription.}
Keep
\[
 E:y^2=x^3-30x-25,\qquad T=(-5,0).
\]
Fix an embedding $\overline\Q\hookrightarrow\overline\Q_{109}$.
Write $s^2=2$ in $\mathbf F_{109^2}$; two is a nonsquare
modulo 109. Let $U,V,W\in E[8](\overline\Q)$ and
$R\in E[13](\overline\Q)$ be the unique prime-to-109
torsion lifts of the following points:
\[
\begin{array}{c|cc}
 &x&y\\\hline
 U&34+53s&11+96s\\
 V&2+18s&8+82s\\
 W&2+91s&8+27s\\
 R&47&20s .
\end{array}
 \tag{C25a}
\]
Here uniqueness is in the indicated finite \'etale torsion
group over the unramified quadratic extension of $\Q_{109}$.
These are algebraic points, so this prescription defines
objects over a number field, not merely over a $p$-adic field.
For $\mathcal P=\{U,U+R,V,W\}$ put
\[
 H:\quad y_H^2=F_H(x),\qquad
 F_H(x)=\bigl((x+5)^2-45\bigr)
       \prod_{P\in\mathcal P}(x-x(P))(x-x(P+T)),
 \tag{C25b}
\]
and take its smooth projective model. Then $H$ has genus four,
and its Jacobian $A$ satisfies
\[
 \operatorname{End}_{\overline\Q}(A)=\Z[i],\qquad
 \operatorname{NS}(A)_\Q=\Q\Theta.
 \tag{C25c}
\]
It is geometrically simple and is of $\Q(i)$-Weil type with
signature $(2,2)$. For a connected finite \'etale cover
$X\to Y$ of Schoen's ball surface one has
\[
 \deg(X/Y)\leq43264,\qquad
 \operatorname{Alb}(X)\twoheadrightarrow A.
 \tag{C25d}
\]
Its two-dimensional Weil space $W_{\Q(i)}(A)$ meets the
subspace generated by products of rational divisor classes
only in zero.

\paragraph{The torsion calculations and the good reduction.}
All the following calculations use the displayed addition law
of $E$ and arithmetic in $\mathbf F_{109}[s]/(s^2-2)$.
They are supplied with independent exact certificates.
The four points in (C25a) lie on $E$, and
\[
 2U=(34,0),\quad4U=0,\qquad
 4V=4W=T,\quad8V=8W=0,\qquad13R=0,
\]
with $R\ne0$. If $\varphi$ denotes 109-power Frobenius,
\[
 \varphi(U)=-U+T,\quad
 \varphi(V)=W,\quad\varphi(W)=V,\quad
 \varphi(R)=-R.
 \tag{C25e}
\]
The three four-element orbits of $U,V,W$ under negation
and translation by $T$ are distinct and avoid both
$E[2]$ and the solutions of $2P=T$. Injectivity of
prime-to-109 reduction proves the same assertions for
their lifts. The order-thirteen point cannot identify
the orbit of $U+R$ with any of these order-eight orbits.
Thus (C25b) has ten distinct branch points, all in the
image of $E[104]$ on its Kummer line. The even-level
transfer of C24 proves (C25d), since $4\cdot104^2=43264$.
The coordinate change (C24c) converts the branch divisor
to $0,\infty$ and four opposite pairs. Hence the balanced
$\Q(i)$-action and its Rosati property follow from the
explicit calculation in C24.

The Frobenius relations (C25e) hold also for the unique
torsion lifts over the unramified quadratic extension.
They preserve the entire branch divisor and its monic
equation. Thus the coefficients of (C25b) lie in
$\Q_{109}$ under the chosen embedding, and the number
field generated by those coefficients has a place with
residue field $\mathbf F_{109}$. This is stronger than
merely observing that a polynomial over a larger residue
field happens to have coefficients in the prime field.
At this place its reduction is
\[
\begin{split}
 \overline F_H(x)={}&x^{10}+62x^9+27x^8+30x^7+45x^6
       +3x^5+40x^4\\
       &+35x^3+65x^2+x+63 .
\end{split}
 \tag{C25f}
\]
Euclid's algorithm gives
$\gcd(\overline F_H,\overline F_H')=1$.
Its leading coefficient is one, so this is a smooth
genus-four reduction with two points at infinity.

\paragraph{A finite absolute-simplicity certificate.}
The exact counts on this one curve are
\[
\begin{array}{c|rrrr}
 r&1&2&3&4\\\hline
 \#H(\mathbf F_{109^r})&134&11794&1294958&141166026 .
\end{array}
 \tag{C25g}
\]
The accompanying computations enumerate all elements of
the finite fields and use quadratic characters to count
the two-sheeted cover, including its two points at infinity.
One implementation uses
$\mathbf F_{109}[u]/(u^2-2)$ and
$\mathbf F_{109}[w]/(w^4-2)$; the independent implementation
uses $u^2-6$ and $w^4-6$ instead.
For the quartic fields irreducibility follows by taking
the quadratic tower: a nonsquare $d$ has square root
$u$ whose norm is $-d$, a nonsquare because
$109\equiv1\pmod4$; thus $u$ is a nonsquare in
$\mathbf F_{109^2}$.
For each element $f$ the identity
\[
 \chi_{\mathbf F_{109^r}}(f)
  =\chi_{\mathbf F_{109}}
       \bigl(N_{\mathbf F_{109^r}/\mathbf F_{109}}f\bigr)
\]
is used, with both sides zero at zero. It follows from
the geometric-series formula for the norm exponent.
All enumerations and norm operations are exact integer
arithmetic modulo 109.

Put $s_r=109^r+1-\#H(\mathbf F_{109^r})$.
Newton's recursion
$k a_k=-\sum_{j=1}^k a_{k-j}s_j$, $a_0=1$,
followed by the reciprocal functional equation, gives
the Frobenius polynomial
\[
\begin{split}
 f(Z)={}&Z^8+24Z^7+244Z^6+1224Z^5+3510Z^4\\
       &+133416Z^3+2898964Z^2+31080696Z+141158161.
\end{split}
 \tag{C25h}
\]
The middle coefficient is a unit modulo 109, so the
reduction is ordinary. Here is an exact irreducibility
and geometric-simplicity test, rather than a numerical
root comparison.

Over $\Q(i)$ one has $f=g\overline g$, where
\[
 g(Z)=Z^4+(12+10i)Z^3+192iZ^2+(-1308+1090i)Z-11881.
 \tag{C25i}
\]
Modulo the Gaussian prime defined by $61=0$, $i=11$,
this is $Z^4-23Z^2+7Z+14$.
In its polynomial quotient over $\mathbf F_{61}$,
\[
 Z^{61^4}=Z,\qquad
 \gcd(Z^4-23Z^2+7Z+14,Z^{61^2}-Z)=1.
 \tag{C25j}
\]
These modular-power and Euclidean checks prove its
irreducibility. Hence $g$ is irreducible over $\Q(i)$.
If $\pi$ is one of its roots, its minimal polynomial over
$\Q$ is divisible by $g$ after extension to $\Q(i)$;
its rational coefficients make it divisible by
$\overline g$ as well. These two irreducible quartics are
distinct, so their product $f$ is irreducible over $\Q$.
The reduction of $A$ is consequently simple.

Write $K_\pi=\Q[Z]/(f)$ and again write $\pi$ for the
class of $Z$. We verify that $\Q(\pi^d)=K_\pi$ for
every $d>0$ by Howe--Zhu, \emph{On the existence of
absolutely simple abelian varieties of a given dimension
over an arbitrary field}, Proposition~3 and Lemma~5,
preprint pp.~4--5. Their criterion says that a first
failure would either make $f$ a polynomial in $Z^d$,
or would put a primitive $d$th root of unity in
$K_\pi$ with $\Q(\pi^d)$ a proper subfield.
The nonzero $Z^7$ coefficient rules out the first case.
In the second case $\phi(d)$ divides $[K_\pi:\Q]=8$,
so the only $d>1$ to test are
\[
 2,3,4,5,6,8,10,12,15,16,20,24,30.
 \tag{C25k}
\]
This list follows directly from
$\phi(\prod\ell^e)=\prod\ell^{e-1}(\ell-1)$:
the only possible prime factors are $2,3,5$, and their
allowed powers give precisely the list displayed.

For each $d$ in this order, express
$1,\pi^d,\ldots,\pi^{7d}$ in the basis
$1,\pi,\ldots,\pi^7$. The determinants of their
integer coefficient matrices, reduced modulo 37, are
\[
 5,11,22,30,9,18,28,14,20,33,7,17,22.
 \tag{C25l}
\]
All are nonzero. Hence each is a rational basis of
$K_\pi$, ruling out all the possible failures in
(C25k). Howe--Zhu's proposition now proves that the
reduction is absolutely simple. The independent
certificate reproduces (C25i)--(C25l), including the
power-basis determinants, by modular polynomial
arithmetic and Gaussian elimination.

The generic $A$ is geometrically simple as well.
Otherwise complete reducibility supplies a nontrivial
idempotent in its geometric rational endomorphism
algebra. After a finite extension it specializes
injectively at the good place to the rational
endomorphism algebra of the reduction; see Conrad,
\emph{Gross--Zagier revisited}, Theorem~2.1(2),
printed p.~70. This contradicts the absolute simplicity
of that reduction, whose rational endomorphism algebra
is a division algebra. Thus no such idempotent exists.

\paragraph{A rank-two torus forces the exact endomorphism ring.}
We now verify that the very same curve has the two-node
reduction used in C24 at a place above thirteen.
All nonzero points of $E[13]$ are conjugate over $\Q$.
Indeed $E(\mathbf F_7)$ has six points, so Frobenius at
seven has characteristic polynomial
$Z^2-2Z+7$ on $E[13]$.
It is irreducible over $\mathbf F_{13}$.
If $\beta$ is its root in $\mathbf F_{169}$, direct
modular powering gives
\[
 \beta^{168}=1,\quad
 \beta^{84}=-1,\quad
 \beta^{56}=9,\quad
 \beta^{24}=2+2\beta.
 \tag{C25m}
\]
The last three values are not one; since $2,3,7$ are
the prime divisors of 168, $\beta$ has order 168.
Multiplication by it acts transitively on the 168
nonzero elements of the two-dimensional $\mathbf F_{13}$
space. The global Galois action contains this Frobenius
action, proving the assertion.

The displayed model of $E$ has good reduction at thirteen.
As in C24, reduction of its 169 geometric thirteen-torsion
points has a nonzero kernel, since its special elliptic
curve has at most thirteen geometric thirteen-torsion
points. Transitivity therefore gives a place above
thirteen at which the particular $R$ in (C25a) reduces
to zero. Reduction remains injective on $E[8]$ at this
place. Apply exactly the two-node argument of C24,
with the good prime thirteen in place of seven:
the branch pairs from $U$ and $U+R$ coalesce, the
other branch values stay distinct, and the geometric
stable fibre is irreducible with two nodes. Hence
$A$ has semistable toric rank exactly two, persistent
under every further finite extension.

Let $D=\operatorname{End}_{\overline\Q}^0(A)$.
Geometric simplicity makes $D$ a division algebra.
After defining all endomorphisms over one finite
extension, the N\'eron property lets them act on the
two-dimensional geometric torus in this semistable
fibre. Its rational character space has dimension two,
so one obtains a unital algebra homomorphism
\[
                    D\longrightarrow M_2(\Q).
\]
It is injective: the kernel is a two-sided ideal of a
division algebra and the identity acts nontrivially.
Equivalently, a nonzero endomorphism of the simple
$A$ is an isogeny, and an inverse up to multiplication
by an integer induces an inverse on the torus up to
the same integer. The two-dimensional rational
character space is a nonzero vector space over the
division algebra $D$ (using the opposite algebra if
pullback conventions require it). Therefore
$\dim_\Q D$ divides two.
The existing $\Q(i)$-action forces $D=\Q(i)$.
Every integral endomorphism is an algebraic integer,
and the curve automorphism already supplies $\Z[i]$.
As this is the maximal order in $\Q(i)$, the integral
endomorphism ring is exactly $\Z[i]$.

The polarization identifies rational divisor classes
with Rosati-symmetric rational endomorphisms: Milne,
\emph{Abelian Varieties}, version~2.00 (2008),
Proposition~14.2, printed pp.~61--62. The
Rosati involution on $\Q(i)$ is conjugation, whose
fixed field is $\Q$, proving the N\'eron--Severi
assertion in (C25c). Alternatively this identification
is the standard map $L\mapsto\lambda_\Theta^{-1}\phi_L$;
its kernel is $\operatorname{Pic}^0$ and its image is
the symmetric subspace.
Every codimension-two divisor polynomial is thus a
multiple of $\Theta^2$. Under the two $K$-eigenspaces
of $H^1$, the polarization has bidegree $(1,1)$ and
$\Theta^2$ has bidegree $(2,2)$. The two Weil lines
have bidegrees $(4,0)$ and $(0,4)$. These are different
summands of $\bigwedge^4H^1(A,\C)$, so their intersection
with the divisor algebra is zero. \hfill$\square$

\paragraph{The remaining cycle test is now on a named exceptional target.}
Apply C13 to the torsion-prescribed fourfold (C25b).
One nonzero projected surface class spans all of its
Weil space under $\Q(i)$. C14 proves that surfaces
whose maps to $A$ factor through products of two
curves have zero such projection. Thus the actual
non-CM source, its simplicity, its endomorphism
algebra, and the exceptional nature of the target
space have all been decided here. The new task is
the specific projected correspondence, with these
obstructions already accounted for. The supplied
Markman theorem already proves algebraicity of
Weil classes on fourfolds; this finding is not a
new proof of that theorem. It gives an explicit
exceptional test object accessible to the ball
construction, and does not assert the Hodge
conjecture for higher-dimensional non-CM targets.

\framework{The source correspondence and Weil projector are in
the homological-motive framework of cote~44, scans~7--8,
and cote~16, scans~43--46. The persistent torus and
its endomorphism action are the N\'eron--Picard and
dual-graph mechanisms of SGA~7~I, IX~12.1(d),
printed pp.~467--468, and IX~12.3.6--7,
printed pp.~471--472. The finite-field calculations
certify a specified source and its endomorphisms;
they do not lift a vertical arithmetic intersection
to an unconstructed characteristic-zero cycle.}


\subsection{THETA3-C27: a finite cup-product test for the exceptional target and a deformation obstruction}\label{C27}

This is a reduction of the remaining construction problem,
not a construction of an exceptional cycle.  Let $A$ be the
specified fourfold in C25 and $K=\Q(i)$.  Fix a smooth
projective connected surface $X$ over $\overline\Q$ and a
base point.  Put
\[
 M_X=\Hom(\operatorname{Alb}(X),A)\otimes_\Z\Q.
\]
It is a finite-dimensional vector space over $K$, acting
by postcomposition.  Let $m=\dim_K M_X$ and choose a
$K$-basis $f_1,\ldots,f_m$.  Clear denominators so that
these are actual maps $X\to A$ through its Albanese.
Translations will not matter in any cohomological formula.

\paragraph{The exact finite test.}
For one embedding $\sigma:K\to\overline\Q$, choose an
algebraic-Betti basis $v_1,\ldots,v_4$ of
$H^1(A,\overline\Q)_\sigma$ and write
\[
 w_\sigma=v_1\wedge v_2\wedge v_3\wedge v_4.
\]
This spans one Weil line.  The individual $v_i$ are not
asserted to be of pure Hodge type over $\overline\Q$;
the exterior product has type $(2,2)$ because $A$ has
signature $(2,2)$.
Define the homogeneous quartic
\begin{equation}\label{C27-quartic}
 F_X(z_1,\ldots,z_m)
   =\int_X\bigwedge_{r=1}^4
                 \left(\sum_{j=1}^m z_j f_j^*v_r\right)
             \in\overline\Q[z_1,\ldots,z_m].
\end{equation}
Then the following are equivalent:
\begin{enumerate}
\item $F_X$ is not identically zero;
\item some algebraic map $f:X\to A$ has
      $e_W\operatorname{cl}(f_*[X])\ne0$, where $e_W$
      is the explicit algebraic Weil projector of C13.
\end{enumerate}
Thus at most $\binom{m+3}{4}$ explicitly symmetrized
cup-product coefficients decide the existence of such
a map on the given surface.  An individual unsymmetrized
four-linear coefficient is not a sufficient test.
There is also a finite constructive search: if the
quartic is nonzero, one of the at most $5^m$ actual maps
\[
                \sum_j n_jf_j,\qquad
                      n_j\in\{0,1,2,3,4\},
\]
has nonzero Weil projection.  Thus once the isotypic
basis is known, no unbounded denominator search is
needed to select the map.

\textbf{Proof.}
For $a_j\in K$, the rational homomorphism
$f_a=\sum_j a_jf_j$ satisfies
\[
 f_a^*v_r=\sum_j\sigma(a_j)f_j^*v_r,
 \qquad
       \int_X f_a^*w_\sigma
                  =F_X(\sigma(a_1),\ldots,\sigma(a_m)).
\]
If the quartic vanishes on $\{0,1,2,3,4\}^m$, view it
as a polynomial in its last variable.  Its degree in
that variable is at most four, so its values at the
five grid points force all its coefficient polynomials
to vanish on the remaining grid.  Induction on $m$
forces every coefficient to be zero.  Hence a nonzero
quartic has a nonzero grid evaluation.  Since the
chosen $f_j$ are already actual homomorphisms, the map
$f_n=\sum_jn_jf_j$ at this evaluation is actual too.
The projection formula identifies it with the pairing
of $\operatorname{cl}((f_n)_*[X])$ and $w_\sigma$.

In $\bigwedge^4H^1(A,\overline\Q)$ the class
$w_\sigma$ pairs only with the opposite pure Weil line:
wedging with any mixed eigenspace gives more than four
vectors in the four-dimensional $\sigma$ eigenspace.
The pairing with that opposite line is nondegenerate.
The cycle class is rational, so its two pure Weil
components are conjugate.  Hence this pairing is
nonzero exactly when its Weil projection is nonzero.
Conversely every map $X\to A$, after a translation,
comes from its Albanese and therefore lies in $M_X$;
the displayed evaluation gives the converse implication.
If $m=0$, both conditions are false, with the quartic
understood as zero.  The coefficient count is the
number of monomials of degree four in $m$ variables.
\hfill$\square$

When this test succeeds, it gives the actual cycle
$e_W f_*[X]$, and C13 then spans the entire rational
Weil space by algebraic endomorphisms.  It requires
no arbitrary Hodge projector.  Degree-one Hodge maps
used to determine $M_X$ are precisely homomorphisms
of abelian varieties after denominators are cleared;
the unproved datum is the particular nonzero quartic.

\paragraph{What finite covers can and cannot change.}
If $\pi:X'\to X$ is finite \'etale of degree $d$, use
the pulled-back maps $f_j\pi$ on $X'$.  For this specified
subspace of maps one has the exact identity
\begin{equation}\label{C27-cover}
                         F_{X',\pi^*M_X}=dF_X.
\end{equation}
Indeed every integrand is a pullback of the corresponding
top-degree form on $X$, and integration multiplies by
$d$.  In particular, taking another finite cover does
not repair a zero Weil pairing of existing maps.
If $M_{X'}=\pi^*M_X$, no map on $X'$ can repair it.
If the larger space contains new maps, its full quartic
has additional mixed coefficients which must actually
be tested; an increase in multiplicity alone is not
a nonvanishing proof.

For the known C25 Albanese quotient the map factors
through $H$.  The common-cover construction using two
curve projections factors through a product of two
curves.  All maps in the corresponding subspace have
zero quartic by C14.  This assertion concerns that
known subspace, not all of $M_X$ on an uncomputed
cover.  Thus a productive search in the noncongruence
tower must produce genuinely new $A$-isotypic maps
and a nonzero mixed coefficient in
\eqref{C27-quartic}, or use a codimension-two
correspondence outside this map construction.

\paragraph{A fixed ball-source map cannot follow a nontrivial first-order target deformation.}
There is a precise obstruction to one proposed
universal-family approach.  Let $f:X\to A$ generate
$A$, meaning that the induced Albanese homomorphism
is surjective.  Let $A_\epsilon$ be a first-order
deformation of $A$ over $\C[\epsilon]/(\epsilon^2)$.
If $f$ lifts to a map from the constant deformation
$X_\epsilon=X\times\operatorname{Spec}\C[\epsilon]$
to $A_\epsilon$, then
\begin{equation}\label{C27-KS}
                  \operatorname{KS}(A_\epsilon)=0
                           \quad\hbox{in }H^1(A,T_A).
\end{equation}

Here is the cotangent and obstruction calculation.
The exact tangent sequence for deformations of a
morphism is Stacks Project, Lemma~93.10.3,
Section~93.10, Tag~0E3S.  Its last term is
\[
 \operatorname{Ext}^1_X(Lf^*L_{A/\C},\mathcal O_X).
\]
Since $A$ is smooth, $L_{A/\C}=\Omega^1_A$ is locally
free in degree zero.  This term is therefore
$H^1(X,f^*T_A)$.  With the source deformation zero,
the obstruction is the pullback of the target
Kodaira--Spencer class, up to the irrelevant overall
sign in the difference convention.  This identification
can also be verified directly: choose locally trivial
deformations on affine charts of $A$; their transition
maps are $1+\epsilon v_{ab}$.  Local lifts of $f$
differ by sections of $f^*T_A$, and their coboundary
must equal the pulled-back cocycle $f^*v_{ab}$.
Thus a lift forces that cocycle class to vanish.

Invariant vector fields trivialize the tangent bundle,
so this obstruction map is exactly
\[
 \operatorname{Lie}(A)\otimes H^1(A,\mathcal O_A)
  \xrightarrow{\ 1\otimes f^*\ }
 \operatorname{Lie}(A)\otimes H^1(X,\mathcal O_X).
\]
The map $f^*$ on $H^1(\mathcal O)$ is injective:
duality identifies it with the differential of the
dual homomorphism to the surjection
$\operatorname{Alb}(X)\to A$, followed by the
Albanese identification of degree-one cohomology.
Equivalently, a polarization splits that surjection
up to isogeny, giving a left inverse in characteristic
zero.  The obstruction map is therefore injective,
which proves \eqref{C27-KS}.

This applies in particular while retaining a fixed
finite \'etale cover of Schoen's $Y$.  A finite
\'etale cover of the constant first-order deformation
of $Y$ is itself constant with its specified special
fibre, by invariance of finite \'etale covers under
nilpotent thickenings; an exact reference is Stacks
Project, Lemma~58.8.3, Tag~0BQB.
The same obstruction applies to a degree-one
polarized correspondence realizing a surjective
Albanese quotient, because it is such a rational
homomorphism after clearing denominators.

Consequently an attempted deformation from a CM
ball-derived map to the non-CM target cannot retain
that fixed source map in a nonzero target tangent
direction.  A deformation route would have to
construct a genuine relative codimension-two cycle,
with a separate flatness and class-transport argument,
rather than deform this degree-one quotient.  This
does not rule out such higher correspondences or
singular-cycle deformations, and it does not produce
one.  The exact outstanding data are now either
a nonzero finite quartic \eqref{C27-quartic} on a
new isotypic source, or that higher relative cycle.
\hfill$\square$


\subsection{THETA3-C28: the actual Schoen fibration cannot produce the exceptional cycle by base change alone}\label{C28}

For the particular target $A$ of C25, the fibration
construction has a stronger limitation than the conditional
test in C27.  Every map to $A$ from a surface obtained by
base change of the chosen Schoen curve fibration factors
through a curve.  This statement includes all maps on
those surfaces, not only the previously written Albanese
quotients.  It follows from the explicit generic fibre
and does not assume that an unspecified fixed part is CM.

\paragraph{The connected generic fibre.}
Use the Fermat projection of Schoen, printed p.~133,
\[
     f:Y\longrightarrow C_5,\qquad
             (y_{ij})\longmapsto(y_{01}:y_{03}:y_{13}),
\]
on the smooth ball quotient $Y$.  Its composition with
the degree-five quotient $C_5\to C_0$ is the map used
in C15--C16 and C25.  On the chart $y_{03}=1$, write
\[
 u=y_{13},\quad s=u^5,\quad v=y_{23},\quad
 b=y_{02},\quad c=y_{12}.
\]
The original defining equations give
\[
 y_{01}^5=1-s,\qquad
 v^5=t,\qquad b^5=1-t,\qquad c^5=s-t,
                  \qquad t=v^5.                 \tag{C28a}
\]
Consequently the geometric generic fibre is the smooth
projective normalization $C_s$ of the last three equations,
where $s\ne0,1$ is transcendental over $\C$.
The three classes $t,1-t,s-t$ are independent modulo fifth
powers in $\overline{\C(C_5)}(t)^\times$:
their valuations at $0,1,s$ show separately that each
exponent in a fifth-power relation is divisible by five.
Thus the function field has degree $5^3$ over the
$t$-line, with deck group
\[
                         G=(\Z/5\Z)^3.
\]
This also proves geometric connectedness of the generic
fibre.  Resolution of the original surface does not change
its function field or this normalized generic fibre.
Since the base curve is normal, the Stein factor of $f$
is therefore $C_5$ itself.

At $0,1,s,\infty$ the inertia order is five, and there
is no other ramification.  At infinity the three radicals
have the same simple pole up to units, so its inertia is
the diagonal cyclic subgroup, rather than a group of
order $125$.  Riemann--Hurwitz gives
\[
 2g(C_s)-2=125(-2)+4\cdot125(1-1/5)=150,
                     \qquad g(C_s)=76.          \tag{C28b}
\]

\paragraph{All rational character factors have dimension at most four.}
A nontrivial deck character is specified by
\[
 (a_0,a_s,a_1,a_\infty)\in\mathbf F_5^4,
       \qquad a_0+a_s+a_1+a_\infty=0.
\]
The last entry records its value on infinity inertia.
Let $r$ be the number of nonzero entries.  The cyclic
quotient $C_\chi=C_s/\ker\chi$ has degree five over
$\mathbf P^1$ and is branched at exactly these $r$
points, each with inertia five.  Hence
\[
                      g(C_\chi)=2(r-2).          \tag{C28c}
\]
Here $r\in\{2,3,4\}$; a single nonzero entry is
incompatible with the displayed sum condition.

The rational character orbits consist of
$\chi,\chi^2,\chi^3,\chi^4$.  Pullback from $C_\chi$
identifies its first cohomology with the sum
of those four eigenspaces in the first cohomology of $C_s$.
Here one may use \'etale cohomology with coefficients
$\Q_\ell(\zeta_5)$, for a fixed $\ell\ne5$, over
the algebraically closed generic field.  The orbit
idempotents themselves belong to $\Q[G]$ and act by
algebraic correspondences.  Thus no embedding of an
extension of $\C(C_5)$ into $\C$ fixing the constants
is being assumed.
Indeed pullback and trace multiply by $\deg(C_s/C_\chi)$,
and averaging over $\ker\chi$ is the projector onto
its invariants.  The trivial-character part is zero,
since $C_s/G=\mathbf P^1$.
Distinct rational orbits give disjoint summands, and
together they exhaust first cohomology.  The corresponding
quotient maps and pullbacks therefore give an isogeny
\[
             J(C_s)\ \sim\ \prod_{[\chi]} J(C_\chi),
                                                        \tag{C28d}
\]
omitting genus-zero factors.  One can see the isogeny
directly on first cohomology: the product map has an
isomorphism there, and therefore has finite kernel and
full-dimensional image.

For completeness the exact character counts are
\[
\begin{array}{c|ccc}
 r&2&3&4\\\hline
 \text{characters}&24&48&52\\
 \text{rational orbits}&6&12&13\\
 \dim J(C_\chi)&0&2&4 .
\end{array}                                             \tag{C28e}
\]
For $r=2$, choose the two positions and one nonzero
entry, giving $6\cdot4=24$.  For $r=3$, choose the
zero position and two nonzero entries whose sum is
nonzero, giving $4(4^2-4)=48$.  Subtracting these and
the trivial character from $5^3$ leaves $52$.
Division by four gives the orbit counts.  The dimensions
sum to $12\cdot2+13\cdot4=76$, as required.

Every positive-dimensional factor in (C28d) has a
unital action by $E_5=\Q(\zeta_5)$.  Its order-five
deck automorphism has no invariant first cohomology,
so its minimal polynomial on that cohomology is
$1+T+T^2+T^3+T^4$; this produces the stated field action.
In particular the twelve dimension-two factors are
CM abelian surfaces.  No assertion about finite
monodromy of the thirteen dimension-four factors is
needed for the argument below.

\paragraph{The C25 target cannot occur in a fibre Jacobian.}
The target $A$ is geometrically simple of dimension four,
with $\operatorname{End}^0(A)=\Q(i)$ by C25.
These properties persist over algebraically closed
extensions of the constant field.  A precise source
for the invariance of abelian subvarieties and
homomorphisms under regular field extension is Milne,
\emph{Abelian Varieties}, version~2.00, Corollary~16.4(a)--(b),
printed p.~68; passage to an arbitrary algebraically
closed extension follows by finite presentation.

If there were a nonzero homomorphism $J(C_s)\to A$
over the algebraic closure of the base function field,
(C28d) would give a nonzero homomorphism from one of
its factors to $A$.  Simplicity of $A$ makes that
homomorphism surjective.  A dimension-two factor
cannot surject.  A dimension-four factor would be
isogenous to $A$, and its $E_5$ action would induce
a unital embedding
\[
                  \Q(\zeta_5)\hookrightarrow\Q(i),
\]
which is impossible by the field degrees.  We have proved
\[
                         \Hom(J(C_s),A)=0.       \tag{C28f}
\]
This argument does not presume that a dimension-four
factor is simple; a surjection between equal-dimensional
abelian varieties is already an isogeny.

\paragraph{Descent of every map on every such base change.}
Let $D\to C_5$ be any dominant morphism of smooth
projective connected curves, and let $X_D$ be a smooth
projective model of the main component of
$Y\times_{C_5}D$, with its induced map to $D$.
Geometric connectedness of the generic fibre gives
a unique main component.  By (C28f), every morphism
$h:X_D\to A$ is constant on the geometric generic
fibre: a map from a smooth projective curve to an
abelian variety, after choosing a point and translating,
is induced by its Jacobian homomorphism.

The resulting constant point is rational over
$\C(D)$.  To check this, extend scalars to an algebraic
closure of that field; each field automorphism fixes
the original map and hence fixes its unique constant
value.  Thus it descends to a rational map $D\dashrightarrow A$.
Properness of $A$ extends this map across every point
of the smooth curve $D$ by the valuative criterion.
The composite $X_D\to D\to A$ equals $h$ on the
generic fibre and hence on a dense open subset;
separatedness of $A$ makes them equal everywhere.
Therefore every $h$ factors through $D$.

For the original base $C_0$, form the normalization
of $C_5\times_{C_0}D_0$ and apply the preceding
argument to each connected component.  When
$D_0\to C_0$ is connected finite \'etale, the
degree-five Galois cover $C_5\to C_0$ pulls back
either to a connected degree-five cover or to five
copies of $D_0$, since its monodromy subgroup has
order either five or one.  These are exactly the
possible main components of the surface base change.
In every case the conclusion is factorization through
the corresponding Stein base curve.

It follows that the full quartic $F_{X_D}$ of C27
is zero for every surface in this single-fibration
base-change tower: the image of each map has dimension
at most one, so the pullback of every degree-four
class integrates to zero.  This includes the C25
cover and every further cover obtained by changing
only this base curve.  A map-generated exceptional
cycle on the specified $A$ must therefore use a
surface cover that genuinely changes the generic
fibre, or leave this map construction in favour of
a higher correspondence.  The proof does not rule
out either route, and does not claim the construction
of the exceptional cycle. \hfill$\square$


\subsection{THETA3-C30: simultaneous Fermat-base changes add only CM Albanese factors}\label{C30}

C28 rules out changing the base of one Schoen fibration.
There is also an exact restriction on the natural attempt
to change two bases simultaneously and use a new Prym
factor.  In this construction every additional Albanese
factor is a quotient of a power of the Fermat-quintic
Jacobian.  Thus no additional copy of the C25 target
appears, and its full C27 quartic is still zero.

\paragraph{Statement for the actual two-fibration cover.}
Let $C=C_5$ be the Fermat quintic
\[
                 U^5+V^5=W^5,
\]
and use the two morphisms
\[
 Y\longrightarrow C\times C,\qquad
 (y_{ij})\longmapsto
 \bigl((y_{13}:y_{01}:y_{03}),
       (y_{23}:y_{02}:y_{03})\bigr)
                                                        \tag{C30a}
\]
on the smooth Schoen ball quotient.  These are the
Fermat projections in Schoen, printed p.~133; C12
computed the generic degree of their product to be five.
Let $f_i:D_i\to C$ be connected finite \'etale covers
of degrees $d_i$, and set
\[
               X=Y\times_{C\times C}(D_1\times D_2).
                                                        \tag{C30b}
\]
Then $X$ is connected, smooth and projective, and is a
finite \'etale arithmetic ball cover of $Y$.  There is
an isogeny
\[
 \operatorname{Alb}(X)\ \sim\ J(D_1)\times J(D_2)\times B,
                                                        \tag{C30c}
\]
where $B$ is a quotient of
$J(C)^{75d_1d_2}$.  In particular $B$ is CM.
The isogeny is compatible, up to isogeny, with the
Albanese map induced by $X\to D_1\times D_2$.

For the simple fourfold $A$ of C25, every morphism
$h:X\to A$ therefore satisfies
\[
 [N]h=b\circ\pi+a_0
                                                        \tag{C30d}
\]
for some positive integer $N$, some morphism
$b:D_1\times D_2\to A$, and a constant $a_0\in A$;
here $\pi$ is the displayed map from $X$ to the product.
Consequently
\[
                    F_X=0
                                                        \tag{C30e}
\]
for the entire quartic in C27.  This includes all
Albanese maps on these covers, not only sums of the two
initially supplied curve maps.

\paragraph{The exact cyclic model and its root line.}
Let $H=\mathcal O_C(1)$, and denote the pulled-back
sections of $f_i^*H$ by $U_i,V_i,W_i$.  On
$P=D_1\times D_2$ put
\[
 H_i=f_i^*H,\qquad
 L=\operatorname{pr}_1^*H_1\otimes
                  \operatorname{pr}_2^*H_2.
\]
The degrees of the $H_i$ are $5d_i>0$.  Define
\[
 s=U_1^5W_2^5-W_1^5U_2^5\in H^0(P,L^5),
 \qquad
 Z=\operatorname{Spec}_P
       \left(\bigoplus_{j=0}^4L^{-j}\right),
                                                        \tag{C30f}
\]
where multiplication uses $s$ to identify the fifth
power of the root with $s$.  Thus $Z\to P$ is finite
flat of degree five, and its direct image of the
structure sheaf is precisely the displayed sum.

On $W_1W_2\ne0$, write $t_i=(U_i/W_i)^5$.
The function field is obtained by adjoining $w$ with
\[
                         w^5=t_1-t_2.
                                                        \tag{C30g}
\]
This is exactly the function field of (C30b): in the
original equations it is the remaining coordinate
$w=y_{12}/y_{03}$.  The branch divisor has smooth
points at which $t_1-t_2$ has a simple zero.  Its
pullback to $P$ still has a simple zero, because both
base changes are \'etale.  Hence this function is not
a fifth power.  Since five is prime and the ground
field contains its roots of unity, the cyclic field
extension has degree five.  In particular $Z$ is
integral.  The same calculation proves connectedness
of (C30b): every component of its finite \'etale map
to connected $Y$ would dominate $Y$, while its generic
algebra over $P$ is the single degree-five field just
computed.  Smoothness, projectivity and the arithmetic
ball property follow from its finite \'etale map to $Y$.

\paragraph{Every singularity is an ordinary Fermat cone.}
The map
\[
          t:C\longrightarrow\mathbf P^1,
                    t=U^5/W^5
\]
has degree $25$, with five points of ramification
index five over each of $0,1,\infty$, and no other
ramification.  Away from these three values, $t$ is
an \'etale coordinate.  Thus the divisor $s=0$ is
smooth there.  At $t=0$ use the local parameter
$x=U/W$, so that $t=x^5$.  At $t=1$ use
$x=V/W$, so that $t=1-x^5$.  At infinity use
$x=W/U$, so that $t=x^{-5}$.  These are actual local
parameters and remain so after the \'etale maps $f_i$.

At a pair of points above the same one of
$0,1,\infty$, the section in (C30f), in the indicated
local trivializations, is $x_1^5-x_2^5$ or its
negative.  At infinity this follows by trivializing
the two factors with $U_1,U_2$, which gives
$s/(U_1^5U_2^5)=x_2^5-x_1^5$.
Consequently the local cyclic surface is
\[
                       z^5=x_1^5-x_2^5,
                                                        \tag{C30h}
\]
after changing one coordinate by a fifth root of $-1$
when needed.  These are all its singular points.
There are exactly
\[
             3(5d_1)(5d_2)=75d_1d_2
                                                        \tag{C30i}
\]
of them.  At a point with just one $W_i=0$, the other
nonzero, the section in the relevant trivialization
is a unit, so the cyclic cover is \'etale and smooth.
This checks the points omitted by the affine equation
(C30g).

The surface is a hypersurface over the smooth surface
$P$, hence Cohen--Macaulay, and it is nonsingular in
codimension one.  It is therefore normal by Serre's
criterion, EGA~IV$_2$, Theorem~5.8.6, printed p.~108
(also Stacks, Lemma~10.157.4, tag \texttt{031S}).
Blowing up the vertex of (C30h) resolves it in one
step.  Indeed its exceptional divisor is the smooth
plane curve $z^5=x_1^5-x_2^5$, and in each standard
blow-up chart the strict transform is the corresponding
affine chart of that smooth curve times the radial
coordinate.  Thus the exceptional divisor is a copy
of $C$ with normal bundle $\mathcal O_C(-1)$.
Let
\[
                       r:\widetilde Z\longrightarrow Z
\]
be the resulting smooth projective resolution.  Its
exceptional divisor is a disjoint union of the
$75d_1d_2$ Fermat curves in (C30i).

The surfaces $X$ and $\widetilde Z$ are smooth
projective models of the same function field.  Their
Albanese varieties, maps to abelian varieties and
degree-four pullback integrals are consequently
identified birationally.  One may use Milne,
\emph{Abelian Varieties}, Theorem~3.2, printed p.~17:
a rational map from a nonsingular variety to an
abelian variety extends everywhere.  Applying it in
both directions gives the Albanese identification.
For the integrals, use a common smooth resolution;
each birational map has degree one.

\paragraph{No new coherent first cohomology on the finite model.}
For $1\leq j\leq4$, each $H_i^{-j}$ has negative
degree, so $H^0(D_i,H_i^{-j})=0$.  The coherent
K\"unneth formula therefore gives
\[
 \begin{aligned}
 H^1(P,L^{-j})={}&
 H^1(D_1,H_1^{-j})\otimes H^0(D_2,H_2^{-j})\\
 &\quad\oplus
 H^0(D_1,H_1^{-j})\otimes H^1(D_2,H_2^{-j})=0.
 \end{aligned}
                                                        \tag{C30j}
\]
In this case the formula follows directly by taking
affine \v Cech covers of the two curves: the product
cover computes the tensor product of the two
complexes, and the K\"unneth decomposition is ordinary
linear algebra over the ground field.
Finite pushforward and (C30f) now imply
\[
 H^1(Z,\mathcal O_Z)=H^1(P,\mathcal O_P),\qquad
 h^1(Z,\mathcal O_Z)=g(D_1)+g(D_2).
                                                        \tag{C30k}
\]
The isomorphism is the pullback from $P$, not merely
an equality of dimensions.

\paragraph{Resolving the cones adds only Fermat-Jacobian factors.}
Let $V=\operatorname{Alb}(\widetilde Z)$, and let
$B\subset V$ be the sum of the images of the
Jacobians of all its exceptional curves.  Thus $B$
is a quotient of $J(C)^{75d_1d_2}$.  Put $Q=V/B$.
The composite $\widetilde Z\to Q$ is constant on
each exceptional curve, because its induced Jacobian
homomorphism has been killed.  It is therefore constant
on every fibre of $r$.

This map descends to a morphism $Z\to Q$.  Here is
a direct descent argument that does not presume
rational singularities.  Take the reduced image of
its graph in $Z\times Q$.  Projection to $Z$ is
proper and birational, and its geometric fibres
consist of one point.  It is finite by the proper
finite-fibre criterion (Stacks, Lemma~37.44.1,
tag \texttt{02LS}).  A finite birational morphism
to normal $Z$ is an isomorphism, since its local
coordinate algebras are integral subalgebras of
the same function field.  This proves the descent.

The image of $\widetilde Z\to Q$ generates $Q$,
so pullback gives an injection
\[
         H^1(Q,\mathcal O_Q)
                         \hookrightarrow H^1(\widetilde Z,
                                                      \mathcal O_{\widetilde Z}).
\]
For example this follows from the surjection
$\operatorname{Alb}(\widetilde Z)\to Q$ and the
usual first-cohomology identification for the smooth
Albanese map.  It can also be seen by pulling back
holomorphic one-forms: a finite sum of differences
of the image generates a dominant map to $Q$, so
a form vanishing on the image must be zero; complex
conjugation gives the assertion for $H^1(\mathcal O)$.
As the map factors through $Z$, its pullback factors
through $H^1(Z,\mathcal O_Z)$.  Thus (C30k) yields
\[
                      \dim Q\leq g(D_1)+g(D_2).
                                                        \tag{C30l}
\]

The morphism $\widetilde Z\to P$ is dominant, and
the product Abel map generates
$J(D_1)\times J(D_2)$.  Hence it induces a surjection
\[
                    V\longrightarrow J(D_1)\times J(D_2).
\]
Every exceptional curve maps to a point of $P$, so
this map kills $B$ and factors through a surjection
from $Q$.  Its target has dimension equal to the
right side of (C30l).  It follows that
\[
                       Q\longrightarrow J(D_1)\times J(D_2)
\]
is an isogeny.  Poincar\'e complete reducibility now
gives (C30c), compatibly with this quotient; a full
proof of the required complement construction is
Milne, \emph{Abelian Varieties}, Proposition~10.1,
printed p.~42.  In particular no vanishing of
$R^1r_*\mathcal O_{\widetilde Z}$, and no assertion
that the cone singularities are rational, has entered
the proof.

For completeness, $J(C)$ is a product up to isogeny
of three CM abelian surfaces.  Its deck group over
the $t$-line is $(\Z/5\Z)^2$.  A character is given
by three entries in $\mathbf F_5$ of sum zero,
describing the inertia at $0,1,\infty$.  The only
characters with nonzero first cohomology have all
three entries nonzero; there are $4^2-4=12$ of them,
in three rational orbits of size four.  The associated
cyclic quotients have genus two by
Riemann--Hurwitz.  Pullback and trace give their
Jacobian product isogenous to $J(C)$, exactly as in
C28.  Each genus-two factor has a unital
$\Q(\zeta_5)$ action, of degree four, and is CM.
Every quotient of a power of this product is CM by
complete reducibility.  This also recovers the CM
Fermat input in Schoen's Remark~(ii), printed p.~134.

\paragraph{The C25 quotient and the entire quartic.}
A nonzero homomorphism from one of these genus-two
factors to the simple fourfold $A$ would have to
surject, which dimensions prohibit.  Consequently
\[
                         \Hom(B,A)=0.
\]
Every homomorphism $V\to A$ therefore factors through
$Q$, and the displayed isogeny $Q\to J(D_1)\times
J(D_2)$ makes it factor through that product after
multiplication by one positive integer.  Translating
the image of the chosen base point gives exactly
(C30d).  All isogeny denominators are retained in $N$.

Let $w$ lie in the two-dimensional Weil space of $A$.
Translations act trivially on cohomology, and
$[N]^*w=N^4w$.  Since $\pi$ has generic degree five,
the projection formula gives
\[
 N^4\int_X h^*w
       =5\int_{D_1\times D_2}b^*w=0.
                                                        \tag{C30m}
\]
The final equality is C14: the image of a product
of two curves is a Lefschetz class, and
$\operatorname{NS}(A)_\Q=\Q\theta$ for the C25
target.  Its Weil projection is zero.  Thus every
actual map has zero pairing, which proves that the
full quartic in C27 vanishes, including on its finite
search grid.

If the two original covers are instead finite
\'etale covers of $C_0$, first normalize their fibre
products with $C_5\to C_0$.  Each connected component
is a finite \'etale cover of $C_5$.  Applying the
argument to each pair of components proves the same
conclusion for every component of the simultaneous
$C_0$-base-change construction.  The statement does
not cover arbitrary finite \'etale covers of the
resulting ball surfaces.  Such a cover need not
arise from two base curves, and its new first
cohomology is not controlled by (C30f).

\framework{This is a restriction obtained from actual
normal models, a specified resolution, and the
Albanese universal property.  The original equations
and exceptional CM curves are Schoen, pp.~133--134.
The exact normality criterion is EGA~IV$_2$~5.8.6,
p.~108; graph descent uses the finite-morphism
criterion above.  The Albanese and realization
viewpoint is Grothendieck, cote~44, scans~7--8.
Together with C14 it excludes the new-Prym attempt
in this particular two-base tower.  It does not
construct an exceptional cycle or exclude covers
whose Albanese factors lie outside this tower.}


\subsection{THETA3-C32: an explicit exceptional cycle for the accessible cyclic-four Jacobian}\label{C32}

The C25 target admits a constructive higher-correspondence
route, although its C27 map quartic vanishes on the covers
in C28 and C30.  The construction below uses the actual
curve equation, a plane in a symmetric power of its quotient
line, and three blowups.  It does not insert a previously
supplied Weil cycle or use Markman's theorem.

\paragraph{Statement and relation to the published construction.}
Let
\[
 H:\quad y^2=z\prod_{j=1}^4(z^2-a_j^2),
 \qquad a_j\ne0,\quad a_j^2\ne a_k^2\ (j\ne k),
                                                        \tag{C32a}
\]
be defined over $\overline\Q$, and let
$A=J(H)$ with its action by
$\rho(z,y)=(-z,iy)$.  This includes the specific
simple fourfold of C25.  There is an explicitly
constructed integral algebraic cycle
$C_A\in\CH^2(A)$, allowed to have signed components,
such that
\[
 0\ne\operatorname{cl}(C_A)\in W_{\Q(i)}(A),\qquad
                   \int_A\operatorname{cl}(C_A)^2=73728.
                                                        \tag{C32b}
\]
Its pullback by the endomorphism $[2+i]$ supplies an
independent rational direction of the Weil space:
the two eigenvalues there are $(2+i)^4=-7+24i$
and its conjugate.  The scalar automorphism $[i]$
itself acts trivially on this fourth exterior power
and is not used for that step.

The underlying method is Schoen,
\emph{Hodge classes on self-products of a variety
with an automorphism}, Compositio~65 (1988),
Theorem~2.0, printed p.~11, and Corollary~3.1,
p.~24.  Its paired branch hypothesis applies here.
We give a complete specialized construction and
nonvanishing proof.  In particular the exceptional
cohomology argument below treats inertia of order
two inside the cyclic group of order four directly.
It does not assume that such inertia is detected
by a prime-degree quotient.

If $f:X\to H$ is the actual surjection from the
compact arithmetic ball cover in C25, the same
equations produce a surface cycle on $X^4$ whose
image in $A$ is a positive integer multiple of
$C_A$.  The precise multiplier is given below.
This is a higher correspondence on a product of
ball surfaces; the cycle's surface components are
not asserted to be ball quotients.

\paragraph{The cyclic cover and its two-dimensional tensor.}
Put $t=z^2$ and $b_j=a_j^2$.  Equation (C32a)
becomes the connected cyclic-four cover
\[
       y^4=t\prod_{j=1}^4(t-b_j)^2
                                                        \tag{C32c}
\]
of $\mathbf P^1$.  Its six branch points have
weights
\[
                 (1,3,2,2,2,2)\pmod4
                                                        \tag{C32d}
\]
at $0,\infty,b_1,b_2,b_3,b_4$, respectively.
Pair them as $(0,\infty),(b_1,b_2),(b_3,b_4)$.
Each pair has total weight four.  The involution
$\rho^2$ is hyperelliptic, hence acts as $-1$ on
$H^1(H,\Q)$.  Thus the only eigenvalues of $\rho$
on first cohomology are $i,-i$, each with multiplicity
four.  On holomorphic differentials the basis
$z^jdz/y$, $0\leq j\leq3$, gives multiplicity two
for each eigenvalue.  The determinant lines are
therefore of type $(2,2)$.

Let
\[
 N=\{(e_1,e_2,e_3,e_4)\in(\Z/4\Z)^4:
                                      \sum e_i=0\},
 \qquad G=N\rtimes\mathfrak S_4,
 \qquad W_0=H^4/G.
                                                        \tag{C32e}
\]
Write $q:H^4\to W_0$ for the quotient.  Its generic
degree is $|G|=4^3\cdot4!=1536$.  The residual
cyclic group is generated by $\delta$, induced by
applying $\rho$ on the first factor.  The map
\[
             \beta_0:W_0\longrightarrow
                  \operatorname{Sym}^4(\mathbf P^1)=\mathbf P^4
\]
has degree four.  On the open set of four distinct
unramified points its root coordinate is the product
of their four $y$-coordinates; its fourth power is
the product of the four right sides in (C32c).
These function fields determine the normal cover
everywhere.

Rational cohomology of a finite quotient is the
invariant cohomology, as follows by equivariant
triangulation and averaging over the finite group.
The part of $H^4(W_0,\Q)$ on which $\delta^2=-1$
is rationally two-dimensional.  To verify this
without a dimension assumption, use K\"unneth on
$H^4$.  Invariance under $N$ forces the four
cyclic characters in a tensor to agree.  If a
factor has degree zero or two its character is
trivial; such a tensor cannot have primitive
residual character.  The primitive part is thus
in $H^1(H)^{\otimes4}$.  Permutations of the four
odd-degree factors include their Koszul signs,
so their invariant tensors are
\[
     \bigwedge^4 H^1(H,\C)_i
          \ \oplus\ 
     \bigwedge^4 H^1(H,\C)_{-i}.
                                                        \tag{C32f}
\]
Each summand has dimension one.  Denote this
rational two-dimensional subspace by $U_0$.

\paragraph{The explicit plane.}
For each branch point $c$ let $S_c\subset\mathbf P^4$
be the hyperplane of effective degree-four divisors
containing $c$, with defining evaluation form $\ell_c$.
Any five of the six evaluation forms are independent:
a nonzero binary quartic cannot vanish at five
distinct points.  Consequently their unique linear
relation has all coefficients nonzero.  Rescale
them, over $\overline\Q$, to write this relation as
\[
       u_1+v_1+u_2+v_2+u_3+v_3=0,
                                                        \tag{C32g}
\]
where $(u_i,v_i)$ correspond to the three paired
branch points.  This realizes $\mathbf P^4$ as
the indicated hyperplane in $\mathbf P^5$.

Let $P_0\simeq\mathbf P^2$ be the plane
\[
                         u_i+v_i=0\quad(1\leq i\leq3).
                                                        \tag{C32h}
\]
Equivalently it is spanned by the three quartics
whose zero divisors consist of all branch points
except one chosen pair.  On $P_0$, the two
hyperplanes of a pair restrict to the same line;
the three lines are its coordinate lines.
The fourth-root equation restricts to a fourth
power up to a nonzero algebraic scalar, because
the exponents in each pair sum to four.

More explicitly, if $F(T)$ represents a quartic
with leading coefficient $\ell_\infty$, the norm
of the right side of (C32c) is
\[
 \frac{\ell_0\prod_{j=1}^4\ell_{b_j}^2}
                         {\ell_\infty^9}
 =\frac{\ell_0\ell_\infty^3
                 \prod_{j=1}^4\ell_{b_j}^2}
                         {\ell_\infty^{12}}.
                                                        \tag{C32i}
\]
Thus $W_0$ is the normalization of the cyclic
cover with root line $\mathcal O_{\mathbf P^4}(3)$
and branch weights (C32d).  On $P_0$ the numerator
in the last expression is a scalar times
$(u_1u_2u_3)^4$.  A fourth root of that scalar
is the only additional constant needed to label
the four sheets over the plane.

\paragraph{The three blowups and the normal bundle.}
Blow up, in order, the strict transforms of
\[
                 K_i=\{u_i=v_i=0\},\qquad i=1,2,3,
\]
and call the successive smooth varieties $M_i$.
The six hyperplanes and their strict transforms
are normal-crossing divisors; the centers are
intersections of two of them.  Let $P_i$ be the
strict transform of $P_0$.  At the $i$th center
its intersection with $P_{i-1}$ is a Cartier
line, so $P_i\simeq P_0$.  The two strict branch
hyperplanes of that pair no longer meet $P_i$.
After all three steps $P=P_3$ misses all the
strict branch hyperplanes.

Here is also the full local normal-bundle
calculation.  After the first two blowups,
$M_2$ is the incidence variety resolving
\[
        ([u_1:v_1],[u_2:v_2]):\mathbf P^4
                                         \dashrightarrow\mathbf P^1\times\mathbf P^1.
\]
It is a $\mathbf P^2$ bundle: after writing
$(u_i,v_i)=(a_i t_i,b_i t_i)$ for $i=1,2$,
the single relation is
$(a_1+b_1)t_1+(a_2+b_2)t_2+u_3+v_3=0$.
The coefficient of $v_3$ is one, so the kernel
is a rank-three vector bundle.  These incidence
equations are the two blowup charts: pulling
back $(u_1,v_1)$ and then the strict ideal
$(u_2,v_2)$ makes each principal.  The fibre at
$[a_i:b_i]=[1:-1]$ is exactly $P_2$.

Near that fibre use base coordinates $r_1,r_2$
and fibre coordinates $[t_1:t_2:u_3]$, so that
$v_3=-u_3-r_1t_1-r_2t_2$.  The final center has
ideal $(u_3,r_1t_1+r_2t_2)$.  The blowup chart
containing the whole strict plane has equation
\[
                    r_1t_1+r_2t_2=u_3w,
\]
and the plane is $r_1=r_2=w=0$.  Its normal
bundle is therefore the kernel in
\[
 0\longrightarrow N_{P/M_3}\longrightarrow
 \mathcal O_P^3
 \xrightarrow{(t_1,t_2,-u_3)}\mathcal O_P(1)
 \longrightarrow0.
                                                        \tag{C32j}
\]
This is the Euler sequence, so
$N_{P/M_3}\simeq\Omega^1_{\mathbf P^2}(1)$ and
\[
 c(N_{P/M_3})=(1+H)^{-1}=1-H+H^2,
                    \int_P c_2(N_{P/M_3})=1.
                                                        \tag{C32k}
\]

Let $W_i$ be the normalization of the main cyclic
cover after each blowup.  The exceptional divisor
has total branch multiplicity four, so division
of the root by its local exceptional equation
removes this multiplicity.  The new branch
divisor is precisely the union of the strict
transforms, with the original weights.  Hence
$W_3\to M_3$ is \'etale over a neighborhood of
$P$.  Since $\mathbf P^2$ has no nontrivial
connected finite \'etale covers, its preimage
is a disjoint union
\[
                      Q_0,Q_1,Q_2,Q_3\simeq\mathbf P^2,
\]
labelled so that $\delta Q_j=Q_{j+1}$.
The normal bundle of each is (C32j).

\paragraph{The inertia argument, including weights two modulo four.}
We verify that the primitive part of cohomology is
unchanged by every map $\varphi_i:W_i\to W_{i-1}$.
Let $E_i\to K_i$ be the exceptional $\mathbf P^1$
bundle on the base, let $F_i$ be its reduced
inverse image in $W_i$, and let $K_i'$ be its
image in $W_{i-1}$.  Write the paired weights
as $a,4-a$, and put $d=\gcd(4,a)$.  Thus
$a=1$ or $2$ in this construction.

The bundle $E_i$ is the projectivization of the
sum of the two normal line bundles of the
intersecting branch divisors.  Its two sections
come from the strict branch divisors.  All other
branch components restrict to vertical divisors
pulled back from $K_i$, because the original
divisors are normal crossing.  In a ratio
coordinate $t$ on this bundle, the normal cyclic
cover $F_i$ has equation
\[
                       v^4=t^a u,
                                                        \tag{C32l}
\]
where $u$ comes from the base, followed by
normalization.  This description includes points
on the vertical branch divisors.  It also identifies
this normalization with the actual reduced inverse
image of $E_i$: locally analytically one extracts
a fourth root of the unit in the branch equation,
after which the radial exceptional coordinate is
absent from the normalized monomial equation.
Setting that coordinate to zero gives exactly
the displayed normal cyclic cover.

There is a global connected $\mathbb G_m$ action
on $F_i$, fixing $K_i$, with local formula
\[
              t\longmapsto\lambda^{4/d}t,
              \qquad v\longmapsto\lambda^{a/d}v.
                                                        \tag{C32m}
\]
To check gluing, use the diagonal scaling action
on the sum of the two normal line bundles.
The root line on $E_i$ is $\mathcal O_{E_i}(1)$
times a line from $K_i$: before restriction it
is the old root line tensored with
$\mathcal O(-E_i)$.  Give this root line the
linearization of weight $a/d$ in the displayed
chart.  On the other chart $t'=t^{-1}$ the
root is $v'=v/t$, with weight $-(4-a)/d$,
so the formulas agree on the overlap.  Base
transition functions have weight zero.
Normalization is functorial for these
automorphisms and preserves the action.

For $\lambda\in\mu_{4/d}$, the base coordinate
$t$ is unchanged, and the resulting automorphisms
are exactly the deck subgroup
$I=\mu_{4/d}\subset\mu_4$.  Here
$\gcd(a/d,4/d)=1$.  Consequently $I$ acts
trivially on rational cohomology of $F_i$:
each element belongs to a connected group
action and is homotopic to the identity.
It also fixes $K_i'$ pointwise.  Indeed a
$\mathbb G_m$ orbit in $F_i$ maps into one
finite fibre of $K_i'\to K_i$; its image is
connected and is therefore a point.  For
elements of $I$ this is the given equivariant
deck action.  A primitive order-four character
is nontrivial on $I$, both when $I=\mu_4$
and when $I=\mu_2$.  Hence neither
$H^*(F_i,\Q)$ nor $H^*(K_i',\Q)$ has a
primitive component.

The map is an isomorphism off these exceptional
sets, and the quotient pairs give a homeomorphism
$W_i/F_i\simeq W_{i-1}/K_i'$.  This follows
directly from properness and the compact
Hausdorff topology after collapsing each entire
closed exceptional set to one point.  Their
relative cohomologies therefore agree.  Apply
the long exact sequences of the two pairs and
the primitive idempotent.  The preceding
vanishing proves
\[
 \varphi_i^*:H^*(W_{i-1},\Q)_{\rm prim}
                    \xrightarrow{\ \sim\ }
                    H^*(W_i,\Q)_{\rm prim}.
                                                        \tag{C32n}
\]
This supplies the exceptional-cohomology step
without a prime-quotient assumption.  In
particular a weight-two branch of the
cyclic-four cover is not mistakenly declared
visible in its degree-two quotient.

\paragraph{Nonzero cycle and exact intersection number.}
The normal covers $W_i$ have finite quotient
singularities.  Locally, after extracting a
root of the analytic unit, their equations
are fourth roots of products of normal-crossing
coordinates.  Substituting fourth roots for
those coordinates realizes each local branch
as a finite abelian quotient of a polydisc.
They are therefore oriented rational homology
manifolds: rational cohomology of a local
quotient sphere is the invariant cohomology
of the sphere, and a complex linear action
preserves its orientation.  In particular
rational Poincar\'e duality, the degree formula
and the projection formula apply to these
proper maps.  The planes $Q_j$ lie in their
smooth loci, so their ordinary fundamental
classes have their usual self-intersection
interpretation.

Set
\[
              C_3=Q_0-Q_2,\qquad c_3=\operatorname{cl}(C_3).
\]
Its class is primitive, since $\delta^2$ sends
it to its negative.  The two planes are
disjoint, and (C32k) gives
\[
                         \int_{W_3}c_3^2=2.
                                                        \tag{C32o}
\]
Write $\varphi:W_3\to W_0$, put
$C_0=\varphi_*C_3$, and let $c_0$ be its
cohomology class.  Proper cycle classes and
Gysin agree.  The degree-one identity
$\varphi_*\varphi^*=1$, together with (C32n),
implies
\[
       \varphi^*c_0=c_3,\qquad
       c_0\in U_0,\qquad
                         \int_{W_0}c_0^2=2.
                                                        \tag{C32p}
\]
Thus no nonzero primitive class is lost in
blowing the construction down.

To obtain an actual cycle on the smooth
$H^4$, take a smooth projective resolution
$T$ of the main component of
$H^4\times_{W_0}W_3$.  Denote its maps by
$p:T\to H^4$ and $q':T\to W_3$.
The first is birational and
$q\circ p=\varphi\circ q'$.
Each $Q_j\hookrightarrow W_3$ is a regular
codimension-two embedding into a smooth
neighborhood.  Its refined pullback along
$q'$ is consequently an integral cycle on
$T$, including any excess-intersection
terms.  Define
\[
 C_H=p_*\bigl(q'^![Q_0]-q'^![Q_2]\bigr)\in\CH^2(H^4).
                                                        \tag{C32q}
\]
The notation means the refined intersection
with $[T]$, pushed from the fibre product
to $T$.  Naturality of the Thom class for
a regular embedding proves the cycle-class
identity for this operation.  Since
$p_*p^*=1$, the commuting diagram and
(C32p) give
\[
 \operatorname{cl}(C_H)=p_*q'^*c_3
                      =p_*p^*q^*c_0=q^*c_0,
 \qquad
              \int_{H^4}\operatorname{cl}(C_H)^2=3072.
                                                        \tag{C32r}
\]
The last integer is $1536\cdot2$.
This construction does not presume that the
finite quotient map $q$ is flat at its
singular points.

\paragraph{The actual exceptional cycle on the Jacobian.}
Choose a base point on $H$ and let
$a:H^4\to A$ be the sum of the four Abel
maps.  It has generic degree $4!=24$:
$\operatorname{Sym}^4H\to\operatorname{Pic}^4H$
is birational, because a general divisor of
degree $g=4$ has one section by Riemann--Roch.
One can check the latter by choosing four
points which impose independent conditions
on the four-dimensional canonical space.

Pullback by $a$ identifies $W_{\Q(i)}(A)$
with $q^*U_0$.  Indeed it sends a first
cohomology class to the sum of its copies
in the four factors.  The cup product of
two classes with the same eigenvalue $i$
or $-i$ vanishes on one curve, since its
eigenvalue would be $-1$ in the trivial
representation $H^2(H)$.  Expanding a
fourfold wedge therefore leaves precisely
the alternating tensors in (C32f), with
nonzero coefficient.  This is an isomorphism
of the rational two-dimensional spaces.

Define $C_A=a_*C_H$.  If
$\operatorname{cl}(C_H)=a^*w_0$, then
$\operatorname{cl}(C_A)=24w_0$ is a nonzero
pure Weil class.  Finally
\[
 \int_A\operatorname{cl}(C_A)^2
   =24\int_{H^4}\operatorname{cl}(C_H)^2
   =24\cdot3072=73728.
                                                        \tag{C32s}
\]
All maps, planes, blowups, normalizations
and cycles are defined over a finite
extension of the original number field.
The only sheet choice above the plane
is a fourth root of the algebraic scalar
in (C32i); a projective resolution is
available in characteristic zero.

\paragraph{A cycle constructed on the known ball power.}
Let $f:X\to H$ be the actual map from
C25.  It is flat: at a point of the
smooth base curve a uniformizer pulls
back to a nonzerodivisor on the smooth
integral surface, so its local module
is torsion free, hence flat, over the
discrete valuation ring.  Choose a very
ample line bundle $L_X$ and four general
smooth divisors $D_i\in|L_X|$ which
dominate $H$.  Each $D_i\to H$ is finite
flat of the same positive degree
\[
                       d=\deg(L_X|_{X_\eta}).
\]
Pull the explicit cycle (C32q) back to
$D_1\times\cdots\times D_4$, and include
it into $X^4$; denote the resulting
surface cycle by $\Gamma$.  Equivalently
one pulls its equations through the four
ball fibrations and takes one relative
hyperplane section in every ball factor.
The finite degree formula and projection
formula give
\[
       (a\circ f^4)_*\Gamma=d^4 C_A,
 \qquad
       \int_A\operatorname{cl}((a\circ f^4)_*\Gamma)^2
                                      =d^8\cdot73728>0.
                                                        \tag{C32t}
\]
This proves nonvanishing of the actual
higher-correspondence cycle.  If
$L_X=K_X^{\otimes m}$ is chosen very
ample and $e$ is the degree of the
Stein base curve over $H$, C28's genus
$76$ fibre gives the additional explicit
value $d=150me$ by adjunction.

The source of this nonzero class is the
signed split-plane cycle and its refined
pullbacks, not the fundamental class of
a curve product or of a ball surface.
It therefore does not contradict C14,
C28 or C30.  It supplies a concrete
exceptional cycle built from the
accessible curve and ball fibrations,
while the literal single-ball theta
packet comparison remains a separate
problem.

\paragraph{The full C24 family in every even dimension.}
The same calculation gives explicit Weil cycles for
every member of C24, in every dimension $2n\geq4$.
Precisely, let
\[
 H_n:\quad y^2=z\prod_{j=1}^{2n}(z^2-a_j^2),
 \qquad A_n=J(H_n),
                                                        \tag{C32u}
\]
with the same nonzero and distinct-square conditions.
Then an integral codimension-$n$ cycle $C_{A_n}$
is constructed by the preceding operations, and
\[
 \begin{split}
 0\ne\operatorname{cl}(C_{A_n})&\in W_{\Q(i)}(A_n),\\
 \int_{A_n}\operatorname{cl}(C_{A_n})^2
       &=2(-1)^n4^{2n-1}\bigl((2n)!\bigr)^2.
 \end{split}
                                                        \tag{C32v}
\]
The sign is the usual middle-degree primitive
polarization sign; the polarized square is positive.
For the particular non-CM members obtained in C24,
these are actual cycles on the supplied targets.
Their being outside the divisor algebra follows
when that target has the additional N\'eron--Severi
property used in C14.  That property is proved in
C25 for the chosen fourfold; it has not been
proved here for every higher-dimensional C24
member.  Algebraicity and the nonzero pure Weil
projection in (C32v) do not require that extra
assertion.

Here are all changes needed in the proof.  The
cyclic-four equation is
\[
            y^4=t\prod_{j=1}^{2n}(t-a_j^2)^2.
\]
Its branch weights are $1,3$ and $2n$ copies of
$2$, paired into $n+1$ pairs.  Its primitive
first-cohomology rank over $\Q(i)$ is $2n$.
The $2n$ holomorphic differentials $z^jdz/y$
give multiplicities $(n,n)$, so the determinant
lines have type $(n,n)$.

In $\operatorname{Sym}^{2n}\mathbf P^1
=\mathbf P^{2n}$, the $2n+2$ evaluation
hyperplanes have one relation, with every
coefficient nonzero.  After rescaling write
it as $\sum_{i=1}^{n+1}(u_i+v_i)=0$.
The plane $P_0\simeq\mathbf P^n$ is
$u_i+v_i=0$ for all $i$; it is spanned by
the $n+1$ divisors omitting one pair.  The
cyclic root line is $\mathcal O(n+1)$:
the total branch weight is $4(n+1)$.
The same norm calculation as (C32i) shows
that its restriction to $P_0$ is a fourth
power up to an algebraic scalar.

Blow up the $n+1$ paired hyperplane
intersections.  The first $n$ blowups give
the incidence $\mathbf P^n$ bundle over
$(\mathbf P^1)^n$ resolving the first
$n$ pair ratios.  The proof is exactly
the displayed incidence calculation:
\[
       \sum_{i=1}^{n}(a_i+b_i)t_i+u_{n+1}+v_{n+1}=0
\]
is a surjective relation on a rank-$(n+2)$
bundle.  The desired $\mathbf P^n$ is
the fibre where every ratio is $[1:-1]$.
For the final blowup its local equation
near the strict plane is
\[
                  \sum_{i=1}^{n}r_it_i=u_{n+1}w.
\]
Hence its normal bundle is the kernel of
\[
       \mathcal O_{\mathbf P^n}^{n+1}
        \xrightarrow{(t_1,\ldots,t_n,-u_{n+1})}
                          \mathcal O_{\mathbf P^n}(1),
\]
namely $\Omega^1_{\mathbf P^n}(1)$.  Its
top Chern number is $(-1)^n$.
The normalized cyclic cover is \'etale near
the strict plane and splits into four
copies.  Their signed difference $Q_0-Q_2$
has square $2(-1)^n$.

Every exceptional step has weights $1,3$
or $2,2$.  The global action (C32m), the
vanishing of primitive exceptional
cohomology, and the relative-pair proof
apply without any dimension change.
The regular embedding now has codimension
$n$, so the graph-resolution refined
pullback (C32q) works in that codimension.
The finite quotient degree is
$4^{2n-1}(2n)!$, and the addition map
$H_n^{2n}\to J(H_n)$ has degree $(2n)!$.
K\"unneth with the same character and
Koszul-sign argument identifies its
primitive subspace with the Weil space.
Multiplying these two degrees by
$2(-1)^n$ proves (C32v).

The second rational direction can again
be obtained using $[2+i]^*$.  Indeed
$(2+i)^{2n}$ is never real: otherwise
$((2+i)/(2-i))^{2n}=1$, whereas
$(2+i)/(2-i)=(3+4i)/5$ is not a root
of unity in $\Q(i)$, whose roots of
unity are $1,-1,i,-i$.  Thus the two
Weil eigenvalues are distinct and
nonrational, and the two rational
cycle classes are independent.

Finally C24 supplies the actual ball
cover $X_n\to H_n$.  Cutting its
$2n$-fold power by one dominating
very ample curve in each factor and
pulling back the same explicit cycle
gives an $n$-dimensional cycle whose
image in $A_n$ is $d^{2n}C_{A_n}$.
Its square is $d^{4n}$ times (C32v).
Thus the all-dimensional conclusion
is a specialized Schoen construction
on this cyclic-four family, together
with the already constructed ball
access and the exact degree and
intersection normalizations proved
here.  It is not a proof for arbitrary
Weil-type abelian $2n$-folds.

\paragraph{Actual integral specialization at a good prime.}
Let $K$ be a number field over which the chosen curve, its
$\Q(i)$-action and the constructed cycle are defined, and
let $R$ be a discrete valuation ring above a finite prime
of $K$. Suppose $A_n$ extends to an abelian scheme
$\mathcal A_n/R$, with the order action extended. This
holds, for example, at any odd prime where all $a_j$ are
units and all $a_j^2-a_k^2$ are units, after adjoining the
chosen sheet label. Take the schematic closure in
$\mathcal A_n$ of each integral component of $C_{A_n}$,
with its signed multiplicity. Each such closure is flat
over $R$: its local rings are torsion free over the
discrete valuation ring. It has relative dimension $n$.
Its special fibre therefore defines an actual integral
cycle $C_{A_n,s}$ of codimension $n$.

For every $\ell$ different from the residue characteristic,
smooth proper base change and compatibility of cycle
classes with specialization give
\[
 \operatorname{sp}\bigl(\operatorname{cl}_{\ell}(C_{A_n})\bigr)
       =\operatorname{cl}_{\ell}(C_{A_n,s}),\qquad
 \deg\bigl(C_{A_n,s}\cdot C_{A_n,s}\bigr)
       =2(-1)^n4^{2n-1}\bigl((2n)!\bigr)^2.
                                                        \tag{C32w}
\]
Here the intersection is the Chow intersection on the
smooth special-fibre abelian variety; its degree equals
the cup-product trace. These are precisely the
specialization operations and exact SGA locators recorded
in C17. The specialization is equivariant for the extended
order action and the exterior-power decomposition of
$H^1$. Thus the class remains in the rank-two determinant
Weil summand, and (C32w) proves that it is nonzero. The
integer in (C32w) remains a nonzero element of $\Q_\ell$
even when the residue characteristic divides it. No
claim of ordinary proper intersection of the signed
components is needed.

This argument deliberately uses closure of the already
constructed integral cycle; it does not require a new
resolution of a mixed-characteristic quotient. If the
ball surface, its map to the curve and the chosen
dominating divisors also extend smoothly, with the
divisors finite flat of degree $d$, the proper pushforward
and flat pullback identities in (C32t) specialize as
well. Such simultaneous models exist after excluding
finitely many primes, since all the stated maps and
equations have been fixed over a number field. This
gives the same actual ball-derived special-fibre cycle
and the factor $d^{2n}$. Its identification with an
interior Kudla--Rapoport moduli cycle is not a consequence
of specialization alone.

\framework{Schoen's original construction is
Theorem~2.0 and Lemmas~2.3--2.8, printed
pp.~11--19; the passage to Jacobians is
Corollary~3.1, p.~24.  The proof above
specializes the plane and normal bundle
and provides the full order-four inertia
argument.  Its regular-embedding pullback,
proper pushforward and graph resolution
are the correspondence operations in
C17 and C23, with the Grothendieck
anchors cote~44, scans~7--8, and
cote~16, scans~43--46.  No new theorem
about all non-CM Hodge classes is inferred
from this particular cyclic-cover family.}


\subsection{THETA3-C36: accessible exceptional Weil targets and cycles in every even dimension}\label{C36}

The concrete fourfold of C25 and its cycle from C32 give an
additional family in every even dimension. These targets are
products. The conclusion below does not assert simplicity or
N\'eron--Severi rank one for the particular higher-genus
Jacobians in C24.

\paragraph{The targets and their access from one ball cover.}
Let $J=J(H)$ be the simple fourfold of C25, with
$\operatorname{End}^0(J)=K=\Q(i)$ and
$\operatorname{NS}(J)_\Q=\Q\Theta$. Let
$C_J\in\CH^2(J)$ be the integral cycle constructed in
C32, so that its class lies in the $K$-Weil space and
\[
                 \int_J\operatorname{cl}(C_J)^2=73728.
 \tag{C36a}
\]
For $n\ge2$ put $m=n-2$ and choose $m$ pairwise
nonisogenous non-CM elliptic curves $F_1,\ldots,F_m$
from C16. Define
\[
                  A_n=J\times\prod_{j=1}^m F_j^2.
 \tag{C36b}
\]
On each $F_j^2$ let $i$ act by
$\rho_j(x,y)=(-y,x)$, and use the given action on $J$.
The product polarization has Rosati involution
$i\mapsto-i$. On the two-dimensional holomorphic tangent
space of $F_j^2$, the eigenvalues are $i,-i$, each once.
Thus $A_n$ has dimension $2n$ and $K$-signature $(n,n)$.
It is non-CM because it has the non-CM factor $J$.

There is a connected finite \'etale arithmetic ball cover
$X_n\to Y$ of the same Schoen surface with a surjection
\[
                    \operatorname{Alb}(X_n)
                               \longrightarrow A_n.
 \tag{C36c}
\]
Here the two copies of every $F_j$ are obtained from the
two different curve projections, rather than counted
twice from one Albanese map.

\textbf{Proof of (C36c).}
C24's even-level transfer, applied in C25, supplies a
connected finite \'etale cover $D_0\to C_0$ with a
nonconstant map $D_0\to H$. The C16 construction gives
connected finite \'etale covers $D_j\to C_0$ with
nonconstant maps $D_j\to F_j$. Let $D$ be any connected
component of
$D_0\times_{C_0}\cdots\times_{C_0}D_m$.
It is finite \'etale over $C_0$. Its projection onto every
connected $D_j$ is finite \'etale, has nonempty image,
and that image is open and closed. Hence each projection
is surjective. In particular $D$ maps nonconstantly to
$H$ and to every $F_j$.

The induced homomorphism
\[
                  J(D)\longrightarrow
                         J\times\prod_{j=1}^mF_j
 \tag{C36d}
\]
is surjective. Each component is surjective by pullback
of holomorphic differentials along the corresponding
nonconstant curve map. Its target factors are simple
and pairwise nonisogenous: $J$ has dimension four and
each $F_j$ has dimension one, and the latter were chosen
pairwise nonisogenous. Complete reducibility therefore
shows that their simple rational first-Hodge-structure
images in $H^1(D,\Q)$ have direct sum. Equivalently,
the dual of (C36d) is injective up to isogeny, and the
image of (C36d) has full dimension. An image abelian
subvariety of full dimension is the whole target.

C12 supplies a generically finite map
$Y\to C_0^2$ of degree $125$. Take a connected component
\[
                   X_n\subset Y\times_{C_0^2}(D\times D).
 \tag{C36e}
\]
It is finite \'etale over $Y$, and surjects onto $Y$
by the same open-and-closed argument. Its proper image
in $D^2$ maps surjectively onto $C_0^2$. Since
$D^2\to C_0^2$ is finite, that image has dimension two;
as $D^2$ is irreducible, the image is all of $D^2$.
Thus $X_n\to D^2$ is dominant and generically finite.
Pullback on $H^1$ is injective: for a generically finite
proper map of smooth projective surfaces, its push-pull
on rational cohomology is multiplication by the nonzero
generic degree. The corresponding map
$\operatorname{Alb}(X_n)\to J(D)^2$ is consequently
surjective. Compose its first $J(D)$ factor with all
components of (C36d), and its second factor with the
$F_j$ components. The result is (C36c).

If $e_j=\deg(D_j/C_0)$ and $e=\deg(D/C_0)$, this
construction retains the finite bounds
\[
 e\le\prod_{j=0}^m e_j,\qquad
 \deg(X_n/Y)\le e^2,\qquad
 \deg(X_n/D^2)\le125.
 \tag{C36f}
\]
For the specific sources above one may take
$e_0\le43264$ and $e_j\le4p_j^2$, where $p_j$ is the
odd torsion prime used to construct $F_j$ in C16.
The cover is compact arithmetic because it is finite
\'etale over the compact arithmetic ball surface $Y$.
All maps and varieties are defined over $\overline\Q$;
a finite extension defines the finitely many choices.
\hfill$\square$

\paragraph{Two explicit Weil divisors on an elliptic square.}
Fix $F=F_j$ and write $\omega_1,\omega_2$ for the
pullbacks of $[0]\in H^2(F,\Q)$. Choose a symplectic
integral basis $u,v\in H^1(F,\Q)$ with
$\int_Fu\wedge v=1$, and use subscripts for the two
factors. For the indicated $i$-action,
$u_1+i u_2$ and $v_1+i v_2$ have eigenvalue $i$.
Put
\[
 \begin{split}
 \beta_+&=(u_1+i u_2)\wedge(v_1+i v_2),\\
 \beta_-&=(u_1-i u_2)\wedge(v_1-i v_2).
 \end{split}
 \tag{C36g}
\]
They span the two embedding lines of
$\bigwedge_K^2H^1(F^2,\Q)$. Both have Hodge type
$(1,1)$. More explicitly,
\[
 \beta_\pm=D\pm iE,\qquad
 D=\omega_1-\omega_2,\qquad
 E=\omega_1+\omega_2-[\Delta_F].
 \tag{C36h}
\]
Indeed the difference map $(x,y)\mapsto x-y$ pulls
$u\wedge v$ back to
$(u_1-u_2)\wedge(v_1-v_2)$ and its zero divisor is
the diagonal. This gives exactly the cross term in
(C36h). Thus $D$ and $E$ are integral divisor classes.
Direct exterior multiplication gives
\[
 \beta_+^2=\beta_-^2=0,\qquad
 \int_{F^2}\beta_+\beta_-=-4,\qquad
 D=\tfrac12(\beta_++\beta_-),\qquad
 \int_{F^2}\beta_\pm D=-2.
 \tag{C36i}
\]
These signs will fix both of the quantitative identities
below. Write $D_j$ also for the actual divisor
$\operatorname{pr}_1^*[0]-\operatorname{pr}_2^*[0]$
on the $j$th elliptic square.

\paragraph{An algebraic projector and the resulting rational cycle.}
Use the C13 projector onto the two $K$-Weil lines in
$H^{2n}(A_n,\Q)$. It can be specified without any
choice of an abstract Hodge projector. Put $b=2+i$,
\[
 f_n(t)=(t-b^{2n})(t-\bar b^{2n}),\qquad
 g_n(t)=\prod_{r=1}^{2n-1}
                       (t-b^r\bar b^{2n-r}).
 \tag{C36j}
\]
Both polynomials lie in $\Z[t]$. Their roots are
distinct because $b/\bar b=(3+4i)/5$ is not a root
of unity in $\Q(i)$. Choose rational polynomials
$r_n,s_n$ by the Sylvester-matrix Bezout identity, so
$r_ng_n+s_nf_n=1$ and their denominators divide
$|\operatorname{Res}(f_n,g_n)|$, and put
$e_n=r_ng_n$. Then $e_n([b]^*)$ is the required
cohomological projector, realized by a rational
combination of algebraic graph correspondences.
Define
\[
 Z_n=e_n([b]^*)\left(C_J\times
                              \prod_{j=1}^mD_j\right)
                  \in\CH^n(A_n)_\Q.
 \tag{C36k}
\]
No Chow-idempotence assertion is needed: the polynomial
is idempotent on the specified cohomology, which is
what is used below.

Let $u_+,u_-$ be the two $K$-embedding components of
$\operatorname{cl}(C_J)$, so
$\operatorname{cl}(C_J)=u_++u_-$.
The determinant of the direct sum of the $K$-modules
in (C36b) is the tensor product of their determinants.
The all-$i$ and all-$-i$ terms are therefore exactly
the terms selected by $e_n$. Equations (C36h)--(C36i)
give the explicit class
\[
 \operatorname{cl}(Z_n)=2^{-m}
 \left(u_+\otimes\bigotimes_{j=1}^m\beta_{j,+}
       +u_-\otimes\bigotimes_{j=1}^m\beta_{j,-}\right).
 \tag{C36l}
\]
It is a nonzero pure $K$-Weil class. If
$\pi:A_n\to J$ is projection, then
\[
 \operatorname{cl}\left(
   \pi_*\left(Z_n\cdot\prod_{j=1}^m
                         \operatorname{pr}_j^*D_j\right)\right)
          =(-1)^m\operatorname{cl}(C_J),
 \qquad
 \int_{A_n}\operatorname{cl}(Z_n)^2
          =(-1)^m73728.
 \tag{C36m}
\]
The first identity uses $2^{-m}(-2)^m=(-1)^m$.
For the second, the squares of each single embedding
line vanish, and the two cross terms in (C36l) give
$2^{-2m}(-4)^m\int_J\operatorname{cl}(C_J)^2$.
Every factor has even cohomological degree, so no
additional K\"unneth sign occurs.

The coefficients of $Z_n$ are rational as stated.
One explicit denominator bound is
\[
                 d_n=|\operatorname{Res}(f_n,g_n)|>0.
\]
The Sylvester-matrix adjugate gives the Bezout identity
with coefficients in $d_n^{-1}\Z$, so $d_ne_n$ has
integral coefficients. Thus $d_nZ_n$ is represented
by an integral signed cycle. Its self-intersection
is $d_n^2(-1)^m73728$, and the first identity in
(C36m) acquires the factor $d_n$. Alternatively any
smaller actual denominator of the chosen projector
may be retained explicitly.

\paragraph{Exceptionality of the entire Weil space.}
Let $B=\prod_jF_j^2$. There is no nonzero homomorphism
between $J$ and $B$: simplicity of $J$ and the distinct
simple isogeny types exclude it. The normalized
Poincar\'e description of line bundles on a product
therefore gives
\[
             \operatorname{NS}(J\times B)_\Q
              =\operatorname{NS}(J)_\Q
                         \oplus\operatorname{NS}(B)_\Q.
 \tag{C36n}
\]
Consequently the subalgebra generated by divisor classes
is the tensor product of the two divisor subalgebras.
Multiplication by $\prod_jD_j$ and proper pushforward
along $\pi$ sends it into the divisor algebra of $J$:
expand a polynomial in the two pullback families and
integrate its $B$-part.

On the two-dimensional $K$-Weil space of $A_n$, that
same contraction is an isomorphism onto the Weil space
of $J$. In the two embedding lines its multipliers
are both $(-2)^m$, by (C36i). They are nonzero; the
formula also respects rational conjugation.
C25 proves that the Weil space of $J$ meets its divisor
algebra only in zero. Applying the contraction shows
\[
 W_K(A_n)\ \cap\
       \Q[\hbox{rational divisor classes on }A_n]=0.
 \tag{C36o}
\]
In particular (C36k) is an actual exceptional algebraic
cycle in every dimension $2n$, $n\ge2$.
Applying $[2+i]^*$ gives a second independent rational
Weil class, since $(2+i)^{2n}$ is nonreal. Thus these
two algebraic cycles span the full Weil space.

The construction is an explicit product consequence of
the checked fourfold cycle and the accessible elliptic
factors. It does not prove that the original genus-$2n$
Jacobians of C24 have N\'eron--Severi rank one, and it
uses no general assertion about Weil classes on arbitrary
non-CM abelian varieties. The common ball cover proves
the source realization (C36c); the exceptional cycle
itself is the algebraic correspondence (C36k), rather
than the image of a surface factoring through a product
of two curves.


\section{Integral models, theta frames and specialization}

\subsection{The precise unramified local normalization}

Let $k_0/\Q_p$ be finite with $p$ odd, residue cardinality $q$, and
uniformizer $\pi$. Let $k/k_0$ be unramified quadratic and
$\mathcal O=\mathcal O_k$. We use Li--Zhang's principally polarized
unitary Rapoport--Zink space $\mathcal N_n$ of signature $(1,n-1)$,
and its nonsplit rank-$n$ Hermitian space $\mathbb V_n$ of special
quasi-homomorphisms. For a full-rank lattice
$L=\mathcal O x_1+\cdots+\mathcal O x_n\subset\mathbb V_n$, define
\[
 \operatorname{Int}(L)=\chi\left(\mathcal N_n,
  \mathcal O_{\mathcal Z(x_1)}\otimes^{\mathbf L}\cdots
                    \otimes^{\mathbf L}\mathcal O_{\mathcal Z(x_n)}\right).
\]
This is a derived Euler characteristic, not a definition by the
ordinary scheme-theoretic intersection length. The source proves
basis independence. Li--Zhang, \emph{Kudla--Rapoport cycles and
derivatives of local densities}, Theorem~1.2.1, equivalently
Theorem~3.4.1, identifies it with
\[
 \operatorname{Int}(L)=-\left.\frac{d}{dX}
                           \operatorname{Den}(X,L)\right|_{X=1}.
\]
Their Theorem~3.5.1 gives the exact polynomial
\begin{align*}
 \operatorname{Den}(X,L)
 &=\sum_{L\subset N\subset N^\vee}
      X^{2\operatorname{length}_{\mathcal O}(N/L)}m(t(N);X),\\
 m(t;X)&=\prod_{i=0}^{t-1}\bigl(1-(-q)^iX\bigr).
\end{align*}
Here $t(N)$ is the number of positive elementary divisors. The
nonsplit space in this convention has odd determinant valuation.
The interpolation normalization is
\[
 \operatorname{Den}((-q)^{-k},L)
 =\frac{\operatorname{Den}(\langle1\rangle^{n+k},L)}
        {\operatorname{Den}(\langle1\rangle^{n+k},\langle1\rangle^n)}.
\]
The identity $\operatorname{Int}=-\operatorname{Den}'(1)$ has no
omitted logarithm. Conversion to a Whittaker derivative introduces
$\log(q^2)$ and explicit unramified factors as in the source's
Remark~3.4.2. The following calculations use this precise
normalization. Author-PDF locators are \S1.2, pp.~3--4;
Theorem~3.4.1, p.~17; Theorem~3.5.1, p.~18; and
Corollary~3.5.3, p.~19.

\subsection{THETA3-C06: an explicit negative derived intersection}\label{C06}

\paragraph{Statement.}
For $n=3$, let $L$ have Hermitian Gram matrix $\pi I_3$.
Then in the specified nonsplit local model,
\[
 \operatorname{Den}(X,L)=(1-X)\bigl(1+(q-q^2)X+X^2\bigr),
 \qquad \operatorname{Int}(L)=q+2-q^2<0.
\]
In particular the value at $q=3$ is $-4$.

\paragraph{Proof.}
The determinant valuation is three, hence odd, so the ambient
Hermitian space is the required nonsplit $\mathbb V_3$. Every
integral overlattice $N\supset L$ has
\[
 \operatorname{val}(N)=3-2\operatorname{length}_{\mathcal O}(N/L)
 \geq0.
\]
Thus its index length is zero or one. The length-zero term is $L$,
of type three. Since $L^\vee=\pi^{-1}L$, a length-one overlattice
has the form $N=L+\mathcal O\pi^{-1}v$, where $v\bmod\pi$ spans
a line in $L/\pi L$. Its pairings with $L$ are integral. Its
remaining integrality condition is
$h(v,v)\in\pi^2\mathcal O$, equivalently that the residual line
is isotropic for the nondegenerate Hermitian form $I_3$ over
$\mathbb F_{q^2}/\mathbb F_q$. Distinct residual lines give
distinct overlattices. Such an $N$ has determinant valuation one,
so its elementary divisors are $(0,0,1)$ and $t(N)=1$.

Here is a direct count of these lines. Let $a_r$ count vectors in
$\mathbb F_{q^2}^r$ whose coordinate norms sum to zero. The norm
map has one preimage of zero and $q+1$ preimages of each nonzero
element, so
\[
 a_1=1,\qquad
 a_r=(q+1)q^{2r-2}-q a_{r-1}.
\]
Consequently $a_2=q^3+q^2-q$ and $a_3=q^5-q^3+q^2$. Dividing the
nonzero isotropic-vector count by $q^2-1$ gives
\[
 \frac{q^5-q^3+q^2-1}{q^2-1}=q^3+1
\]
isotropic lines. The density formula therefore becomes
\begin{align*}
 \operatorname{Den}(X,L)
 &=(1-X)(1+qX)(1-q^2X)+(q^3+1)X^2(1-X)\\
 &=(1-X)\bigl(1+(q-q^2)X+X^2\bigr).
\end{align*}
Differentiating and invoking Li--Zhang gives
$\operatorname{Int}(L)=q+2-q^2$. Since
$q^2-q-2=(q-2)(q+1)>0$ for odd $q\geq3$, it is negative.
$\square$

This is a specialization of the source theorem, not a claim of a
new local-density theorem. It is consistent with the independent
horizontal and vertical calculation in Li--Zhang's proof of
Lemma~6.2.1, pp.~35--36: at $a_3=1$ the two terms are
$(q+1)+(1-q^2)$. A separate finite-field script in this run counted
$28,126,344,1332$ lines for $q=3,5,7,11$, respectively, and obtained
intersection values $-4,-18,-40,-108$. These finite checks do not
replace the general proof.

Thus a positive definite Hermitian norm cannot be identified with
an arbitrary local arithmetic intersection without controlling the
derived signs and the chosen cycle combination. Positivity of
representation densities at their interpolation arguments also
does not imply positivity of the derivative at one.

\paragraph{Framework anchor.}
The precise intersection-theoretic analogue is SGA~4$\frac12$,
\emph{Cycle}, Remark~2.3.9, p.~150, where the multiplicity is an
alternating Tor sum.  The negative example prevents identifying this
derived local number with the positive polarized pairing of cote~44,
scans~7--8, without an additional comparison theorem.

\subsection{THETA3-C07: a positive local sector with an exact bound}\label{C07}

\paragraph{Statement.}
Let $M$ be a self-dual Hermitian $\mathcal O$-lattice of rank $n-1$,
and let $x\perp M$ satisfy $h(x,x)=\pi^{2a+1}$ with $a\geq0$.
For $L=M\perp\mathcal O x$,
\[
 \operatorname{Den}(X,L)=(1-X)\sum_{j=0}^aX^{2j},
 \qquad \operatorname{Int}(L)=a+1\geq1.
\]

\paragraph{Proof.}
An integral overlattice $N\supset L$ splits as
$M\perp(N\cap M_k^\perp)$. Indeed, for $y\in N$, its orthogonal
projection to $M_k$ belongs to $M$ because $h(y,M)\subset\mathcal O$
and the pairing on $M$ is perfect. Subtract this projection.
The complementary lattice is $\pi^{-j}\mathcal O x$. Containment
and integrality give $j\geq0$ and $2a+1-2j\geq0$, hence
$0\leq j\leq a$. Each overlattice has index length $j$ and type
one, since its remaining positive elementary divisor is the odd
integer $2a+1-2j$. The density formula is the displayed sum.
Differentiation at one and the local theorem give $a+1$.
$\square$

The rank-one deformation interpretation after cancellation of the
self-dual part appears in Li--Zhang \S2.11 and Examples~2.4.3 and
3.5.2. The preceding proof does not require that interpretation.
It gives a normalized positive local certificate; comparison with
a different archimedean mixed period still requires a matching
theorem for its test data and factors.

\paragraph{Framework anchor.}
Here the exact local model permits a positive calculation, while the
cycle-class and intersection comparison still has the hypotheses of
SGA~4$\frac12$, \emph{Cycle}, Theorem~2.3.8 and Remark~2.3.9,
pp.~148--150.  Its connection to the correspondence/polarization problem
of cote~44, scans~7--8, requires the geometric identification in C17.

\subsection{THETA3-C08: vertical positivity does not detect a generic class}\label{C08}

\paragraph{Statement and counterexample.}
A positive ordinary local intersection on an integral model need
not detect a nonzero generic-fibre cycle. Let $R$ be a discrete
valuation ring with uniformizer $\pi$, and consider $\mathbb P^1_R$.
Its special fibre $F$ and a section $H$ meet with local intersection
ring $R/(\pi)$ and length one. Nevertheless $F$ has empty generic
fibre, so its restriction to the generic Chow group and cohomology
is zero. $\square$

An actual horizontal flat cycle of the required relative
codimension has an algebraic generic fibre by construction. The
remaining issue is identifying its class with the desired Hodge
tensor. The specialization criterion below makes the available
functoriality precise; it does not construct the missing cycle.

The basic-locus uniformization retains a geometric factor as well.
Vollaard--Wedhorn, \emph{The supersingular locus of the Shimura
variety of $\operatorname{GU}(1,n-1)$ II}, equation~(5.4.2), p.~34,
is
\[
 I(\Q)\backslash
 \bigl(\mathcal N_{\mathrm{red}}\times G(\mathbb A_f^p)/K^p\bigr)
 \simeq \mathcal M^{\mathrm{ss}}.
\]
It does not identify the whole supersingular locus with a finite
Shimura set. In rank three the Rapoport--Zink factor contributes
Fermat curves, as recalled in the source's introduction, p.~4.
Definite mass formulas can describe finite indexing data, while
local intersections depend on the retained geometry.

Li--Zhang's Theorem~4.2.1, p.~21, decomposes the horizontal part of
a rank-$(n-1)$ special cycle into primitive quasi-canonical lifting
cycles indexed by integral overlattices of type at most one.
Theorem~6.1.3, p.~34, compares the horizontal intersection function
with the corresponding horizontal density derivative. These are
more refined existing inputs than an unqualified positivity claim
for the total intersection.

\paragraph{Framework anchor.}
Cote~16, scans~44--46, explicitly discusses specialization in the
standard-conjecture framework.  The example above explains why a
vertical local intersection is not, by itself, a generic realization
class.  C17 states the actual relative-correspondence hypotheses under
which the cited SGA specialization comparison applies.

\subsection{THETA3-C09: finite lifting valuations and the infinite-product obstruction}\label{C09}

\paragraph{Statement.}
For $e_s=q^{s-1}(q+1)$,
\[
 \sum_{s\geq1}\frac1{e_s}=\frac{q}{q^2-1}.
\]
If $t_s\in\C_p$ have normalized valuations $v(t_s)=1/e_s$, the
infinite product $\prod_{s\geq1}t_s$ does not converge in $\C_p$.

\paragraph{Proof.}
The sum is geometric. The valuations of its partial products tend
to the displayed finite value, so the absolute values tend to a
nonzero number. A convergent sequence of these products would have
a nonzero limit, forcing its successive ratios $t_s$ to tend to
one. But $v(t_s)>0$ implies $|t_s|<1$ and $|1-t_s|=1$ for every
$s$. This is impossible. $\square$

The relevant primary source is Gross, \emph{On canonical and
quasi-canonical liftings}, Proposition~5.3(2),(3), manuscript p.~9,
proof pp.~9--11, in the original MPI/SFB~85-30 scan appended to
Gross--Zagier. These are PDF pages~77--78 of the retrieved combined
scan. In the unramified quadratic case it gives the degree
$[W_s:W]=e_s$ and valuation $v_{W_s}(t_s)=1$ of the formal modulus.
The original scan was visually checked. Li--Zhang \S4.1, p.~20,
uses the same extension degree $q^s(1+q^{-1})$.

Normalize the base valuation by $v(\pi)=1$. Then
$v(t_s)=1/e_s$. The canonical modulus $t_0\in\pi W$ satisfies
$v_{W_s}(t_s-t_0)=1$. In $W[[t]]$ the intersection with the canonical
section is therefore
\[
 W_s/(t_s-t_0),
\]
of ordinary length one over $W_s$ and over $W$, since the residue
fields agree. The branch degree $e_s$, its local intersection
length one, and the base-normalized coordinate valuation $1/e_s$
are different invariants. For finite $S$,
\[
 v\left(\prod_{s=1}^St_s\right)
 =\frac{q(1-q^{-S})}{q^2-1},\qquad
 v\left(\prod_{s=1}^S
             \operatorname{Norm}_{W_s/W}(t_s)\right)=S.
\]
The norm uses complete conjugate orbits and has valuation one at
each level. None of these finite identities is itself an identity
for a theta section. A comparison by divisors must specify the
section, its integral frame, and which branches or norms occur.
For the comparison with C11 over $\Q_p$, set this $q=p$; the
Lubin--Tate residue cardinality in C11 is instead $Q=p^2$.

\paragraph{Framework anchor.}
Gross's deformation divisors supply a concrete local intersection problem.
The exact normalization matters before applying the intersection and
cycle-class formalism of SGA~4$\frac12$, \emph{Cycle}, pp.~148--150.
The correspondence viewpoint in cote~44, scans~7--8, does not identify
these local divisors with a chosen analytic theta section.

\subsection{THETA3-C10: non-CM arithmetic theta and its realization}\label{C10}

\paragraph{Elementary observation.}
If an algebraic cycle $Z$ is homologically trivial, every ordinary
cohomological pairing with $Z$ is zero. A nonzero arithmetic height
may detect information about rational equivalence but does not
supply a nonzero ordinary Hodge class.

\paragraph{Proof.}
Homological triviality is precisely
$\operatorname{cl}_{\mathrm B}(Z)=0$. Its cup product with every
Betti class is zero, and the same statement holds in the compatible
\'etale realization. A height does not change this vanishing.
$\square$

There is an existing non-character extension relevant to this
boundary. Li--Liu, \emph{Chow groups and $L$-derivatives of
automorphic motives for unitary groups, II}, Example~1.10,
pp.~6--7, includes the representation attached to
$\operatorname{Sym}^{2r-1}$, $r\geq2$, of a non-CM elliptic curve,
under its modularity and local hypotheses. Good reduction at all
nonsplit finite places gives $S_\pi=\varnothing$; when
$r[F:\Q]$ is odd the stated case applies. Theorem~1.4 produces a
nonzero localized $\CH^r_{\mathrm{hom}}$ from a simple central zero,
with the technical Hypothesis~6.6 removable in this empty-$S_\pi$
case by Remark~1.9. Theorem~1.5's height formula also requires the
stated generating-series modularity hypothesis. These cited source
results are carried as background, not reproved here.

Thus non-CM input is available in arithmetic theta theory, but this
output belongs to homologically invisible cycles and
Abel--Jacobi or motivic extension data. A non-CM Hodge-class problem
still requires an algebraic correspondence with the relevant
nonzero realization and the rank condition C04.

\paragraph{Framework anchor.}
The output in this source is homologically trivial, so its nonvanishing
height does not give a nonzero class in the pure cohomological realization.
This is the precise distinction from the algebraic cohomology classes
and correspondence inverses discussed in cote~16, scans~43--46.

\subsection{THETA3-C11: the exact local derivative slope, with sharpness}\label{C11}

\paragraph{Statement.}
Let $p$ be odd, put $Q=p^2$, and choose $\pi\in\C_p$ with
$v_p(\pi)=1$. Set
\[
 \ell(T)=\sum_{r\geq0}\frac{T^{Q^r}}{\pi^r},\qquad
 E(Z)=\ell^{-1}(Z),\qquad D=\ell'(T)^{-1}\frac{d}{dT}.
\]
For every $f\in\mathcal O_{\C_p}[[T]]$ and every integer $a\geq1$,
\begin{align*}
 v_p(D^af(0))
 &\geq v_p(a!)-\frac{a-1}{p^2-1}\\
 &=\frac{ap}{p^2-1}+\frac1{p^2-1}
                     -\frac{s_p(a)}{p-1},
\end{align*}
where $s_p(a)$ is the base-$p$ digit sum and $v_p(0)=+\infty$.
Moreover $D^p=0$ after reduction modulo $p$.
The bound is attained infinitely often: for $f(T)=T$ and $a=Q^r$,
\[
 v_p(D^{Q^r}T(0))=\frac{(Q^r-1)p}{p^2-1}.
\]

\paragraph{Proof of the bound.}
Choose $\beta$ with $\beta^{Q-1}=\pi$, so
$v_p(\beta)=1/(Q-1)$. The scaled logarithm is
\[
 L(T)=\frac{\ell(\beta T)}\beta
 =\sum_{r\geq0}
     \pi^{(Q^r-1)/(Q-1)-r}T^{Q^r}.
\]
It is integral and has linear coefficient one, since for $r\geq1$
\[
 \frac{Q^r-1}{Q-1}-r=1+Q+\cdots+Q^{r-1}-r\geq0.
\]
Recursive formal inversion over any commutative ring therefore gives
\[
 U(T)=\frac{E(\beta T)}\beta
       \in T\mathcal O_{\C_p}[[T]],\qquad U'(0)=1.
\]
Integrality of $f$ implies
\[
 f(E(\beta T))-f(0)=f(\beta U(T))-f(0)
                         \in\beta T\mathcal O_{\C_p}[[T]].
\]
The chain rule in the logarithmic coordinate gives
\[
 \frac{d}{dZ}f(E(Z))=(Df)(E(Z)),\qquad
 f(E(Z))=\sum_{a\geq0}\frac{D^af(0)}{a!}Z^a.
\]
Its coefficient after substituting $Z=\beta T$ is
$\beta^aD^af(0)/a!$, of valuation at least $v_p(\beta)$ for
$a\geq1$. This proves the bound. Legendre's formula
$v_p(a!)=(a-s_p(a))/(p-1)$ gives the second expression. Since
$s_p(a)\leq(p-1)(1+\lfloor\log_pa\rfloor)$, the error from the
linear term is explicitly logarithmic.

\paragraph{Proof of nilpotence.}
We have
\[
 \ell'(T)=\sum_{r\geq0}(Q/\pi)^rT^{Q^r-1}\equiv1\pmod p.
\]
Thus $D$ reduces to $d/dT$ over $\mathcal O_{\C_p}/p$.
The $p$-th iterate of $d/dT$ is zero in characteristic $p$, because
its coefficient on every monomial contains a product of $p$
consecutive integers. This also holds for this nonreduced
coefficient ring.

\paragraph{Proof of sharpness.}
In $L(T)$ the exponent of $\pi$ is zero for $r=0,1$ and positive
for every $r\geq2$. Modulo the maximal ideal,
$\overline L(T)=T+T^Q$. In characteristic $p$ its inverse is
\[
 V(T)=\sum_{j\geq0}(-1)^jT^{Q^j},\qquad V(T)+V(T)^Q=T.
\]
Formal inversion commutes with reduction. Hence the
$T^{Q^r}$ coefficient of $U$ is a unit. Since
$E(Z)=\beta U(Z/\beta)$, the coefficient of $Z^{Q^r}$ in $E$
has valuation exactly $-(Q^r-1)/(Q-1)$. For $f=T$ and $a=Q^r$,
$D^aT(0)=a![Z^a]E(Z)$ and $s_p(a)=1$, giving
\[
 v_p(D^aT(0))=
 \frac{a-1}{p-1}-\frac{a-1}{p^2-1}
 =\frac{(a-1)p}{p^2-1}.
\]
This proves all assertions. $\square$

\paragraph{Precise relation to the CM local source.}
The coefficient is
\[
 \frac{p}{p^2-1}=\frac1p+\frac1{p^3}+\frac1{p^5}+\cdots,
\]
exactly the linear slope in Alp\"oge's step~2, p.~3.
For this arithmetic application, now require $\pi$ to be a
uniformizer of the unramified quadratic extension of $\Q_p$.
In that local-field setting the displayed logarithm is a permitted
coordinate for the height-two Lubin--Tate formal group.
The purely formal theorem above allows arbitrary $\pi\in\C_p$
of valuation one and does not assert an integral formal group for
that wider class of parameters.
Bannai--Kobayashi, \emph{Integral structures on $p$-adic Fourier
theory}, Proposition~3.5 and proof, pp.~531--532, first use
$\sum T^{q^r}/\pi^r$ and then treat general integral parameters.
An integral change of parameter with unit linear coefficient
preserves integral coefficient functions. Under a formal-group
isomorphism the normalized invariant derivations differ by a
constant unit, so the valuation statement is unchanged.
In their CM elliptic application, \S5, pp.~543--544, $p$ is odd
inert and the curve has good supersingular reduction; its formal
group is the height-two Lubin--Tate group for
$\pi=-\epsilon p$. The preceding proof allows this unit multiple.

The source provides actual integral functions to which the estimate
applies. On p.~546 it sets, for $b$ prime to $p$,
\[
 \wp_b(z)=\wp(z)-b^2\wp(bz),\qquad
 \phi(t)=\wp_b(\lambda(t))\in\mathcal O_K[[t]],
\]
and constructs an integral modified zeta function as well.
Proposition~5.5, pp.~546--547, identifies the invariant derivatives
and Fourier moments with modified Bernoulli--Hurwitz values.
Theorem~5.1, p.~544, states Katz's stronger floor-divisibility
result for the specified modified values, re-proved in that section.
These are existing CM elliptic applications of the same slope; no
new proof of Katz's full theorem is claimed here.

The slope also equals the valuation of a Lubin--Tate period.
Schneider--Teitelbaum, \emph{$p$-adic Fourier theory},
Lemma~3.4(b), gives
\[
 v_p(\Omega_{\mathrm{LT}})=\frac1{p-1}-\frac1{e(q-1)}.
\]
Putting $e=1$, $q=p^2$ gives $p/(p^2-1)$; the source's appendix
gives the period calculation. The $p$-adic period
$\Omega_{\mathrm{LT}}$ is distinct from the complex period
$\Omega_{\C}$ in the analytic normalization of Alp\"oge's theta
values.

\paragraph{Divided powers give a stronger, testable local identity.}
Katz, \emph{Divisibilities, congruences, and Cartier duality},
Identity~2.4 and Congruences~2.5, p.~670, work in the integral
algebra of invariant differential operators of the coordinate model
$[p](T)=pT+\alpha T^{Q}$, with $\alpha$ a unit and $Q=p^2$.
Define $D(n)$ by
\[
 f(X\mathbin\oplus Y)=\sum_{n\geq0}(D(n)f)(X)Y^n,
 \qquad D=D(1).
\]
Then $D(n)=D^n/n!$ for $n<Q$, and
\[
 D(n)^p\equiv
 \begin{cases}
 0,&p\nmid n,\\
 \alpha^{n/p}D(n/p),&p\mid n
 \end{cases}
 \pmod{p\operatorname{Diff}(G)}.
\]
In particular $D^p\equiv0$ and $D(p)^p\equiv\alpha D$ modulo $p$.
For completeness, the operator-valued translation series
$\Phi(U)=\sum D(n)U^n$ satisfies
$\Phi([p](U))=\Phi(U)^p$ by iterated translation.
Reduction modulo $p$ and comparison of the $U^{pn}$ coefficients
give the congruence. Over characteristic zero,
$\Phi(U)=\exp(D\lambda(U))$; since
$\lambda(U)\equiv U\pmod{U^Q}$, comparison below degree $Q$
gives the divided-power identity. This reconstructs the source's
elementary operator argument and exhibits more integral information
than nilpotence alone.

Katz's Main Theorem~3.11, p.~672, gives the stronger bound
\[
 v_p(D^nf(0))\geq\left\lfloor\frac{np}{p^2-1}\right\rfloor
\]
for integral $f$ satisfying the explicit torsion-trace equation
\[
 \sum_{u\in G[p]}f(T\mathbin\oplus u)=0.
\]
This theorem is cited as background, not reproved here.
In \S~IV, equations~(4.8)--(4.10), p.~676, Katz constructs such an
integral function on the supersingular CM elliptic formal group:
$f=\wp_b-[p]^*\wp_b$. Equations~(4.13)--(4.16), p.~677, identify
its derivatives with the modified Bernoulli--Hurwitz values.
Thus the stronger divisibility has a concrete functional equation
and actual integral elliptic examples. Matching a proposed theta
section to this function, or proving its own torsion-trace identity
in a compatible integral frame, remains a specific additional
comparison.

\paragraph{An integral trace-zero operator and the optimal uniform error one.}
In Katz's exact model $[p](T)=pT+\alpha T^Q$, over the ring $R$
of integers of a finite extension of $\Q_p$, the cited theorem
implies the stronger uniform estimate
\[
 v_p(D^af(0))\geq
       \left\lfloor\frac{ap}{p^2-1}\right\rfloor-1
       \qquad(f\in R[[T]],\ a\geq1).
\]
The constant one is optimal uniformly in $f$ and $a$.
Here is a complete reduction to Katz's theorem. Put
$A=R[[S]]$, $B=R[[T]]$, $S=[p](T)$. Weierstrass division for
\[
 T^Q+(p/\alpha)T-S/\alpha=0
\]
gives a free $A$-basis $1,T,\ldots,T^{Q-1}$ of $B$.
Newton identities yield
\[
 \Tr(1)=Q,\qquad \Tr(T^j)=0\ (1\leq j<Q-1),\qquad
 \Tr(T^{Q-1})=-(Q-1)p/\alpha.
\]
Thus $\Tr_{B/A}(B)=pA$. Define the integral functions
\[
 h=\Tr_{B/A}(f)/p,\qquad
 g=pf-[p]^*h.
\]
Their trace satisfies $\Tr(g)=p\Tr(f)-Qh=0$.
Over the split generic kernel this trace is the sum over translations
by $G[p]$, so $g$ satisfies exactly Katz's torsion-trace equation.
The equality holds integrally as well, since the formal rings are
$p$-torsion free. Katz's theorem therefore applies to $g$.
The invariant derivative satisfies $D[p]^*=p[p]^*D$, and $D$
preserves the integral formal functions. Hence
\[
 pD^af(0)=D^ag(0)+p^aD^ah(0)
\]
has valuation at least $\lfloor ap/(p^2-1)\rfloor$; the second
summand has valuation at least $a$, which is no smaller.
Dividing by $p$ proves the estimate. For $a=Q-1$ and
$f=T^{Q-1}$, one has $D^af(0)=a!$ and
\[
 v_p(a!)=p-1,
 \qquad\left\lfloor\frac{ap}{p^2-1}\right\rfloor=p.
\]
A smaller uniform loss is therefore impossible. $\square$

The argument initially has the exact DVR and coordinate hypotheses
of Katz. It includes the CM elliptic groups described in his \S~IV.
After integral base change from this model, it also holds for
$\mathcal O_{\C_p}$ coefficients: a fixed derivative uses only a
finite jet; approximate its integral coefficients by algebraic
integral coefficients in a finite extension and pass to the closed
valuation ball. Integral changes of formal parameter with unit
linear coefficient preserve the valuation of the normalized
invariant derivative. This is a deduction from Katz's theorem,
without a priority claim. It is not an automatic assertion about
an arbitrary analytic theta normalization.

\paragraph{The remaining theta-frame gate.}
A theta function is a section of a line bundle. Extension as an
invertible sheaf over the formal integral neighborhood gives some
integral frame, but differentiating its coefficient depends on the
frame or connection. The comparison still required is an equality
between the normalized analytic derivatives
$D^a\theta(0)/\Omega_{\C}^a$ and $D_G^af(0)$ for an integral
coefficient function in a specified compatible frame. Existence of
an arbitrary integral trivialization does not prove that equality;
a horizontal or cubical analytic frame need not be integral.

This is an actual obstruction to inferring the derivatives from
the divisor. On the formal characteristic-zero germ
$\C_p[[T]]$, take the trivial line bundle, section one, and frame
$e$. Replace the frame by
$e'=(1+p^{-M}T^a)e$. The two frames agree at the origin and the
section has the same empty divisor, but its coefficient becomes
$f'=(1+p^{-M}T^a)^{-1}$ and
\[
 D^af'(0)=-a!p^{-M}.
\]
Its valuation is arbitrarily negative as $M$ grows. The new frame
is a generic formal frame, not an integral one; no assertion about
invertibility on the whole rigid open unit disc is involved.
Even the exact integral Cartier divisor does not determine higher
jets: for $2\leq a<Q$, all
$f_N(T)=T(1+p^NT^{a-1})$ define the divisor $(T)$ and have first
derivative one, whereas
\[
 D^af_N(0)=a!p^N,
\]
because $E(Z)=Z+O(Z^Q)$. The uniform lower bound remains valid,
but the precise derivatives cannot be recovered from that divisor.

\paragraph{A concrete reduced-theta failure of integrality.}
The distinction occurs for a usual normalized reduced theta.
Take $E:y^2=x^3-x$, $\omega=dx/(2y)$, and $t=-x/y$.
At $p=7$ this is a good supersingular CM elliptic curve.
After replacing $y$ by $2y$, it is precisely the example in
Bannai--Kobayashi, arXiv:math/0610163v4, \S4, pp.~49--50.
Their Example~1.9, pp.~10--11, defines
$\theta(z)=\exp(-e_2^*z^2/2)\sigma(z)$ with $\theta'(0)=1$.
For the square lattice $e_2^*=0$, by its $i$-symmetry, also recorded
in Figure~1, p.~49. Thus $\theta=\sigma$ in this normalization.

The equation $t^2=x/(x^2-1)$ gives
\[
 x=\frac{1+\sqrt{1+4t^4}}{2t^2},\qquad y=-x/t,
 \qquad \omega=\frac{dt}{\sqrt{1+4t^4}}.
\]
Consequently
\[
 \lambda(t)=t-\frac25t^5+\frac23t^9+O(t^{13}).
\]
The Weierstrass equation $\wp''=6\wp^2-2$ gives
$\wp(z)=z^{-2}+z^2/5+z^6/75+O(z^{10})$.
Using $(\log\sigma)''=-\wp$ and the sigma normalization yields
\[
 \log(\sigma(z)/z)=-\frac{z^4}{60}-\frac{z^8}{4200}+O(z^{12}),
 \qquad
 \sigma(z)=z-\frac{z^5}{60}-\frac{z^9}{10080}+O(z^{13}).
\]
Exact composition therefore gives
\[
 \theta(\lambda(t))
 =t-\frac5{12}t^5+\frac{1411}{2016}t^9+O(t^{13}).
\]
Since $2016=7\cdot288$ and $7\nmid1411$, the $t^9$ coefficient
has valuation $-1$ at $7$. This actual normalized reduced theta is
not an integral coefficient function in the stated formal
parameter. The finite calculation makes no assertion about an
asymptotic radius of convergence and does not refute Alp\"oge's
accepted asymptotic theorem. It proves that analytic reducedness
alone does not provide the integral-frame hypothesis needed above.
A separate rational symbolic calculation reproduced these
coefficients; it is not a formal proof certificate.

\paragraph{Framework anchor.}
This finding concerns an integral function, invariant derivation, and
specified local frame.  It reaches the correspondence/polarization
problem of cote~44, scans~7--8, only after a section and integral-model
comparison such as C17 is supplied.  The explicit reduced-theta example
shows why that normalization datum cannot be suppressed.

\subsection{THETA3-C17: the cycle-moduli comparison and its specialization criterion}
\label{C17}

The specialization statement below is a standard consequence of the cited
EGA/SGA functorialities.  The contribution here is the explicit comparison
criterion and identification of the geometric data that an application
to the proposed bridge would have to supply.

\paragraph{The two cycles have different definitions and targets.}
The 032 construction starts with maps from a ball surface to fixed
CM abelian varieties, induced through the Albanese. For
$f:S\to A=\prod_iB_i$, it produces the image-surface cycle
$f_*[S]\in\CH_2(A)$. It is a cycle on the fixed abelian variety
$A$.

By contrast, Li--Zhang \S2.3, p.~11, defines a local special cycle
$\mathcal Z(L)$ by an extension condition on special
quasi-homomorphisms: after applying the framing quasi-isogenies,
each $x\in L$ must extend from the special fibre to an actual
$\mathcal O_k$-linear homomorphism of the indicated
$p$-divisible groups. Globally, \S\S13.2--13.3, p.~82, starts
with the unitary moduli object $(A_0,A)$, its actions,
polarizations, and level, and considers tuples of special
homomorphisms in the allowed prime-to-$p$ isogeny category with
prescribed Hermitian Gram matrix $T$ and prescribed finite test
function. This is a moduli condition on a unitary Shimura model.
It is not supplied by an arbitrary map from its surface to fixed
abelian varieties. In particular the two definitions neither have
the same ambient variety nor automatically the same codimension.

A literal comparison therefore requires a specified algebraic
correspondence from the relevant Shimura model to the fixed target
$A$, a theorem identifying its action on the selected special cycle
with the projected image-surface cycle, and integral extensions
compatible with the desired local model. This is an identification
gate, not an assertion that such a correspondence is impossible.
C04 still supplies the separate rank requirement on the target
Hodge summand.

\paragraph{Exact primary EGA and SGA locators.}
EGA~IV$_3$, Theorem~8.8.2, pp.~28--29, and Scholie~8.8.3,
pp.~33--34, descend finite-presentation schemes, morphisms, and
finite collections of relations through filtered limits.
EGA~IV$_2$, Theorem~6.9.1, p.~153, is generic flatness. These give
spreading of finite specified geometric data and flatness after
shrinking a suitable integral base. They do not produce a desired
cycle, and shrinking does not ensure good behavior at an arbitrarily
specified bad prime.

For specialization use the original SGA~4, tome~III,
Expos\'e~XVI, Theorem~2.1, p.~213, and Corollary~2.3, p.~215,
with locally constant prime-to-residue-characteristic coefficients.
The locally constant hypothesis is explicit in those statements;
one need not rely on the misprinted formulation of Corollary~2.2.
For cycle classes use SGA~4$\frac12$, \emph{La classe de cohomologie
associ\'ee \`a un cycle}: Definition~2.3.2, p.~145;
the perfect-weight construction~2.3.3, p.~146;
Lemma~2.3.7 and Theorem~2.3.8, p.~148; and
Remark~2.3.9, p.~150. Lemma~2.3.7 gives base-change compatibility.
Theorem~2.3.8 gives pullback and proper-intersection compatibility
under its fibrewise codimension and perfect-weight hypotheses.
Remark~2.3.9 explicitly uses alternating sums of Tor for proper
intersection multiplicities. These original scans were inspected;
none of these source theorems is claimed to have been reproved here.

\paragraph{A precise cohomological criterion.}
Let $R$ be a strictly henselian discrete valuation ring,
$\ell$ invertible in $R$, and $X,Y$ smooth proper over $R$.
Let $\Gamma\subset X\times_RY$ be a relative flat cycle of pure
codimension $c$ with the preceding cycle-class hypotheses; for
example its structure sheaf is a perfect weight on the regular
ambient product. Fix geometric generic and special fibres and the
specialization isomorphisms. Then the action of $\Gamma$ commutes
with specialization:
\[
 \operatorname{sp}_Y\circ(\Gamma_{\bar\eta})_*
       =(\Gamma_{\bar s})_*\circ\operatorname{sp}_X.
\]
The correspondence action has the usual degree shift
$2c-2\dim(X/R)$ and Tate twist $c-\dim(X/R)$.

\paragraph{Proof from the cited functorialities.}
Write the action as pullback to the product, cup product with
$\operatorname{cl}(\Gamma)$, and the proper smooth projection's
trace. The cited specialization theorem identifies the fibre
cohomologies. Pullback and cup product commute with base change;
Lemma~2.3.7 identifies the two restrictions of the cycle class;
and the trace functoriality in the construction~2.3.1--2.3.4
commutes with that base change. Combining the three operations
proves the identity. Flatness makes the fibre restriction of
$\mathcal O_\Gamma$ agree with its derived base change, so it is
the stipulated fibre cycle class. $\square$

Consequently, once actual correspondences and their relative cycle
classes are provided, their cohomological action and any certified
rank calculation can be transported through this comparison. The
criterion neither lifts arbitrary special-fibre cycles nor proves
that an arbitrary Hodge projector is algebraic. The same
cycle-class/intersection compatibility shows that homological
triviality implies numerical triviality: pair a zero cycle class
with each complementary algebraic class and apply the trace.
The converse is not supplied by EGA/SGA or by the local
intersection formula. Passing among a geometric correspondence,
its homological realization, and a desired Hodge tensor therefore
requires the explicit maps and rank statements described above.

\paragraph{Archival anchor.}
Cote~16, scan~43, states the complementary-dimension issue for algebraic
classes; scans~44--45 formulate cohomological inverse maps induced by
algebraic correspondences; scan~46 continues the discussion of
specialization and hyperplane reductions.  C17 concerns the functoriality
of already specified relative correspondences.  It does not establish the
existence of the conjectural correspondences in that letter.  Cote~44,
scans~7--8, supplies the related Albanese and polarization interpretation.


\subsection{THETA3-C18: tame theta ratios have an integral common frame}\label{C18}

Let $p\geq5$, let $R$ be the ring of integers of a finite extension of
$\Q_p$, and let $\mathcal E/R$ be the smooth elliptic curve
$y^2=4x^3-g_2x-g_3$ with $g_2,g_3\in R$ and unit discriminant.
Use its integral parameter $t=-2x/y$ and its formal logarithm
$\lambda(t)$ normalized by the invariant differential $dx/y$.
The letter $\mathcal E$ here denotes a curve, not a CM field.
Enlarge $R$ by a finite extension whenever the torsion points below
must be made rational.

For a compatible system of division points $\widetilde T=(T_n)$,
suppose $T_1$ is a nonzero torsion point of order $N$ prime to $p$.
Write $s_T(t)$ for its algebraic sigma series in
Katagiri--Kobayashi, \emph{p-adic properties of division polynomials
and algebraic sigma functions}, arXiv:2609.07024v1,
Definition~3.1, p.~9. Their Remark~3.2 on that page identifies it over
$\C$ with
\[
 s_T(t)=\sigma_{z_T}(\lambda(t)),\qquad
 \sigma_{z_T}(z)=
 e^{-\eta(z_T)(z+z_T/2)}\sigma(z+z_T).
\]
The torsion lift, differential, and normalization are part of this
notation. Theorem~2.7, pp.~7--8, supplies the algebraic value and
Taylor-coefficient statement. For an arbitrary compatible system,
one can choose a rational lattice lift realizing the finite datum
$T_{4N}$ on which these constructions depend.
The following is a precise integral consequence of their
construction, rather than an integrality assertion for the raw sigma
series.

\paragraph{Common-factor statement.}
There are units $u_T(t)\in R[[t]]^\times$ such that
\begin{equation}\label{C18-factor}
 s_T(t)=q(t)u_T(t),\qquad
 q(t)=\frac{\sigma(\lambda(t))}{t},\qquad
 u_T(0)=\sigma_{\rm alg}(\widetilde T)\in R^\times.
\end{equation}
The same $q(t)\in\operatorname{Frac}(R)[[t]]^\times$ works for every
such tame torsion point. Consequently, for two lists of the same
length,
\begin{equation}\label{C18-ratio}
 R_{\mathbf T,\mathbf U}(t):=
 \frac{\prod_{i=1}^r s_{T_i}(t)}
      {\prod_{i=1}^r s_{U_i}(t)}
       =\frac{\prod_i u_{T_i}(t)}{\prod_i u_{U_i}(t)}
       \in R[[t]]^\times.
\end{equation}
After pullback by a common prime-to-$p$ multiplication isogeny,
this is the completion of a rational function on
$\mathcal E$ that is a unit along the identity residue disc.
On that cover it is the ratio of sections of equal tensor powers
of the same line bundle; before pullback a torsion flat twist can
remain. This integral-unit conclusion is specific to equal tensor
degrees and prime-to-$p$ torsion. A weaker derivative estimate for
arbitrary integer tensor degree is proved below.

\paragraph{Proof, including the integral model.}
The rational function $g_T$ used in Definition~3.1 has
\[
 \operatorname{div}(g_T)=[N]^*(\tau_{T_1}^*(O)-(O)),\qquad
 (t g_T)(O)=N^{-1}\sigma_{\rm alg}(\widetilde T).
\]
By their Lemma~3.3(1), pp.~9--10, the latter sigma value is a
unit. Here is the resulting local integral calculation.
Because $N$ is prime to $p$, all the torsion sections in this
divisor are disjoint over $R$ after the finite extension. The
only one through the special-fiber origin is the simple pole
along $O$. Any vertical contribution to the divisor of $g_T$
is an integer multiple of the smooth irreducible special fiber.
In the regular local ring at its origin, cancellation of the
horizontal pole gives $t g_T=\varpi^a v$ with $v$ a unit and
$a\in\Z$. Its restriction to $O$ is a unit by the displayed
normalization. Hence $a=0$, and
\[
 h_T(t):=t g_T(t)\in R[[t]]^\times.
\]
This argument can equivalently be made directly with the
horizontal rational function in the proof of Lemma~3.3.

Multiplication by $N$ on the integral formal group has expansion
$[N](t)=Nt+\cdots$ with $N\in R^\times$. Its compositional inverse
$a_N(t)=[N^{-1}](t)$ belongs to $tR[[t]]$, with linear coefficient
$N^{-1}$. Therefore
\[
 u_T(t):=t g_T(a_N(t))
       =\frac{t}{a_N(t)}h_T(a_N(t))\in R[[t]]^\times,
 \qquad u_T(0)=N\,h_T(0)=\sigma_{\rm alg}(\widetilde T).
\]
Definition~3.1 now gives \eqref{C18-factor}, and cancellation
gives \eqref{C18-ratio}.

To verify the rational-function assertion, take a common multiple
$M$ of all the orders $N_i,N_i'$ in the two lists, still prime
to $p$. The formal group identity
$a_{N_i}([M](t))=[M/N_i](t)$ gives directly
\[
 R_{\mathbf T,\mathbf U}([M](t))=
 \frac{\prod_i[M/N_i]^*g_{T_i}(t)}
      {\prod_i[M/N_i']^*g_{U_i}(t)}.
\]
The common sigma factors cancel before this substitution.
The right side is an explicit rational function on $\mathcal E$
over the same finite torsion extension. Its completion is an
integral unit by \eqref{C18-ratio} and the integral multiplication
series. This proves the asserted rationality on the finite \'etale
isogeny cover without asserting descent to the original curve.
\hfill$\square$

\paragraph{Derivative consequence.}
Suppose the integral formal parameter has been identified by a
unit-linear integral change of coordinates with the exact height-two
model of C11, $[p](T)=pT+\alpha T^{p^2}$ with $\alpha\in R^\times$.
For its normalized invariant derivation $D$, C11 applies to
\eqref{C18-ratio} and gives, for every $a\geq1$,
\[
 v_p\bigl(D^a R_{\mathbf T,\mathbf U}(0)\bigr)
       \geq\left\lfloor\frac{ap}{p^2-1}\right\rfloor-1.
\]
The assertion survives multiplication by an integral unit and
unit-linear integral changes of parameter, with the invariant
differential normalization retained. This is an actual theta
quotient to which the integral-function estimate applies.

For a single translated theta section the same calculation gives
the exact connection conversion
\[
 q^{-1}D(qf)=Df+\frac{Dq}{q}f.
\]
It separates the integral function $u_T$ from the common connection
term $Dq/q$. The coefficient example in C11 shows why $q$ cannot
simply be declared integral. For a tensor quotient of total degree
zero this common factor cancels as in \eqref{C18-ratio}; for a
tensor of degree $d$ it remains as $q^d$. Thus the remaining
comparison with Alp\"oge's harmonic covariant theta derivatives
is the specified connection and its degree, rather than an
unspecified raw-theta integrality hypothesis. This statement does
not identify those archive derivatives with the degree-zero
quotients without their actual connection data.

\paragraph{The derivative of the common connection is integral and rational.}
Put $A=Dq/q$, with $D$ normalized by $dx/y$.  The classical identity
$(\log\sigma)''=-\wp$ and $x=\wp(z)$ give the exact identity
\begin{equation}\label{C18-connection-rational}
 B:=DA=D^2\log q=-x-D^2\log t
   =\frac{(Dt)^2-tD^2t-t^2x}{t^2}.
\end{equation}
This is a rational function on the elliptic curve, and
$B\in t^2R[[t]]$.  Here is a direct integral verification, which does
not assume that $q$ or $A$ is integral.  Put $u=t^2x$.  The Weierstrass
equation and $t=-2x/y$ give
\[
 4u^2(u-1)=g_2u t^4+g_3t^6.
\]
The implicit-function derivative at $(u,t)=(1,0)$ is $4\in R^\times$.
Thus $u\in1+t^4R[[t]]$, so $x=t^{-2}+t^2R[[t]]$.
Since $Dx=y$ and $Dy=6x^2-g_2/2$, direct differentiation gives
\[
 Dt=\frac{4x^3+g_2x+2g_3}{4x^3-g_2x-g_3}
    =1+\frac{g_2t^4}{2u}+\frac{3g_3t^6}{4u^2}
       \in1+t^4R[[t]].
\]
Consequently $D=(Dt)\,d/dt$ preserves $R[[t]]$ and
$D^2t\in t^3R[[t]]$.  The numerator in
\eqref{C18-connection-rational} belongs to $t^4R[[t]]$, proving the
claim.  An explicit expression, useful independently of the parameter,
is
\begin{align*}
 B&=9x-\frac{g_2}{2x}-\frac{g_3}{x^2}
       -\frac{(6x^2-g_2/2)^2}{y^2}\\
  &=\frac{-5g_2x^4-13g_3x^3+(g_2^2/4)x^2
                +(3g_2g_3/2)x+g_3^2}{x^2y^2}.
\end{align*}
Since $q\in1+t^4\operatorname{Frac}(R)[[t]]$, both
$\log q(0)$ and $D\log q(0)$ are zero.  Under the exact height-two
coordinate hypothesis already stated, C11 now gives
\[
 v_p\bigl(D^a\log q(0)\bigr)
   \geq\left\lfloor\frac{(a-2)p}{p^2-1}\right\rfloor-1
       \qquad(a\geq3),
\]
while $D^2\log q(0)=B(0)=0$.
The integral rational derivative need not have an integral primitive.
For example, C11's curve with $g_2=4,g_3=0$ gives
\[
 A=-\frac53t^3+\frac{11}{7}t^7+O(t^{11}),\qquad
 B=-5t^2+t^6+O(t^{10}).
\]
Thus at $p=7$ the connection term itself is not integral in this frame.

\paragraph{A weaker positive slope for every tame tensor degree.}
Retain the exact height-two coordinate hypothesis, so that the
operator congruence from C11 holds:
\[
 D^p\bigl(R[[t]]\bigr)\subset pR[[t]].
\]
For every integer $d$, every $v(t)\in R[[t]]$, and every $n\geq1$,
the series $f=q^dv$ satisfies
\begin{equation}\label{C18-weaker-slope}
 v_p\bigl(D^nf(0)\bigr)
 \geq\max\left\{0,\ v_p(n!)-\frac{np}{p^2-1}\right\}
 =\max\left\{0,\frac{n}{p^2-1}
                         -\frac{s_p(n)}{p-1}\right\}.
\end{equation}
Here $s_p(n)$ is the sum of the base-$p$ digits of $n$.
This applies in particular to every finite product or quotient of the
tame translated sigma series: its integer tensor degree is $d$, and
the remaining factor $v$ is an integral unit.  Degree zero still has
the stronger C11 estimate with slope $p/(p^2-1)$.
The slope in \eqref{C18-weaker-slope} is $1/(p^2-1)$; it does not assert
the stronger slope for a single theta section.

\emph{Proof.}
The operator congruence implies
$D^k(R[[t]])\subset p^{\lfloor k/p\rfloor}R[[t]]$ by iteration.
Let $l_j=D^j\log q(0)$.  In the additive coordinate $z=\lambda(t)$,
both $t(z)$ and $\sigma(z)$ are odd.  Hence $q(t(z))$ and its formal
logarithm are even.  As the logarithm starts in degree four,
$l_j=0$ for odd $j$ and for $j<4$.  For even $j\geq4$,
\[
 l_j=D^{j-2}B(0),\qquad
 v_p(l_j)\geq\left\lfloor\frac{j-2}{p}\right\rfloor.
\]
For every such $j$ one has the elementary inequality
\begin{equation}\label{C18-Bell-weight}
 v_p(j!)-\left\lfloor\frac{j-2}{p}\right\rfloor
          +\frac1{p-1}\leq\frac{jp}{p^2-1}.
\end{equation}
For $j\leq p-1$ this is immediate; for $j=p+1$ it is an equality.
The remaining even integers satisfy $j\geq p+3$.  Using
$v_p(j!)\leq(j-1)/(p-1)$ and
$\lfloor(j-2)/p\rfloor\geq(j-p-1)/p$, the left side is at most
$j/(p(p-1))+(p+1)/p$.  Its difference from the right side is
nonpositive because already at $j=p+3$ the right side minus this
upper bound is
$(p^2-3p-2)/(p(p^2-1))>0$; the difference then increases with $j$.

Bell's formula for the exponential gives
\[
 D^n(q^d)(0)=
 n!\!\sum_{\substack{m_j\geq0,\ \sum jm_j=n\\j\geq4\ {\rm even}}}
       \prod_j\frac{(d l_j)^{m_j}}{m_j!(j!)^{m_j}}.
\]
It is also a polynomial with integer coefficients in the integral
quantities $d l_j$, so these derivatives are integral.  In each term,
use $v_p(m_j!)\leq m_j/(p-1)$ and
\eqref{C18-Bell-weight}.  The resulting valuation is at least
$v_p(n!)-np/(p^2-1)$.  This proves the stated bound for $q^d$,
including negative $d$; $d=0$ is immediate.

Finally write the additive-coordinate series for $q^d$ and $v$ in
exponential Taylor form.  The coefficient of $z^k$ in $q^d$ has
valuation at least $-kp/(p^2-1)$.  The coefficient of $z^a$ in $v$
has valuation at least
\[
 \left\lfloor\frac ap\right\rfloor-v_p(a!)
  =-\sum_{r\geq2}\left\lfloor\frac a{p^r}\right\rfloor
  \geq-\frac{a}{p(p-1)}
  \geq-\frac{ap}{p^2-1}.
\]
Convolution gives the same coefficient bound for $q^dv$, and
multiplication by $n!$ gives the factorial lower bound in
\eqref{C18-weaker-slope}.  Integrality of all its derivatives follows
separately from Leibniz's formula, because the derivatives of $q^d$
and $v$ are integral.  Legendre's factorial formula gives the final
digit-sum expression. \hfill$\square$

\framework{This is an integral trivialization and specialization
calculation for a line-bundle ratio on a smooth proper model.
Its geometric setting is EGA~IV$_2$, Theorem~6.9.1, p.~153
(flat spreading after localization), and SGA~4$\frac12$,
\emph{Cycle}, Definition~2.3.2, p.~145, and Theorem~2.3.8,
p.~148 (cycle classes and pullback). The connection normalization
is additional arithmetic data, not a change to Grothendieck's
cycle-class or polarization statements.}


\subsection{THETA3-C19: intrinsic CM derivative bounds and the reduced theta frame}\label{C19}

This subsection removes the exact Lubin--Tate coordinate hypothesis
from the derivative conclusions of C18 for a good-reduction CM
elliptic curve at an inert prime.  The argument uses integral
translation operators and the CM action, rather than an asserted
integral isomorphism with a specified polynomial formal group.

\paragraph{Hypotheses and normalization.}
Let $p\geq5$ be inert in an imaginary quadratic field $K$, put
$q=p^2$ and $N=q-1$, and fix the embedding $K\hookrightarrow\C_p$
through which the CM action acts on invariant differentials.
Let $E$ be an elliptic curve over a number field with CM by an order
$\mathcal O\subset K$ whose conductor is prime to $p$.
After a finite extension, suppose that $E$ has good reduction over
$R=\mathcal O_L$, where $L/\Q_p$ is finite and contains $K_p$,
that all its CM endomorphisms are defined, and that its chosen
invariant differential $\omega$ generates the integral invariant
differentials of the smooth model $\mathcal E/R$.
CM endomorphisms extend to this abelian scheme by the N\'eron
mapping property: see Conrad, \emph{Minimal models for elliptic
curves}, author version dated 21 November 2015,
Definition~5.1 and Example~5.3, p.~13.
Uniqueness of the extension preserves the group-homomorphism
identities, since they already hold on the generic fiber.  Thus
$\mathcal O\otimes\Z_p=\mathcal O_{K_p}$, with residue field
$\mathbb F_q$, acts on its formal completion.
The completion of this action is proved below, so no continuity
assumption about an unspecified action is needed.

Choose an integral formal parameter $t$ with
$\omega=(1+O(t))\,dt$, and let $D$ be the invariant derivation dual
to $\omega$.  Thus $Dt(0)=1$.  All valuations below satisfy
$v_p(p)=1$.  The ring $R$ may be ramified over $\Z_p$; in
particular, $R/pR$ is not assumed reduced.

\paragraph{Statements.}
For every $f\in R[[t]]$ and every $n\geq1$,
\begin{equation}\label{C19-strong}
 v_p\bigl(D^nf(0)\bigr)
 \geq\max\left\{0,
       \left\lfloor\frac{np}{p^2-1}\right\rfloor-1\right\}.
\end{equation}
The loss of one cannot be replaced uniformly by zero: at
$n=p^2-1$, the function $f=t^{p^2-1}$ has
$v_p(D^nf(0))=p-1$, whereas
$\lfloor np/(p^2-1)\rfloor=p$.
If the torsion-translation trace of $f$ is zero, the stronger bound
\begin{equation}\label{C19-trace-zero}
 v_p\bigl(D^nf(0)\bigr)
       \geq\left\lfloor\frac{np}{p^2-1}\right\rfloor
\end{equation}
holds.  Independently, the full integral operator satisfies
\begin{equation}\label{C19-Dp}
 D^p(R[[t]])\subset pR[[t]].
\end{equation}

For the theta assertion use a short Weierstrass equation and
differential defined over a number field and integral at the
chosen place.  Write the equation as
$y^2=4x^3-g_2x-g_3$ with unit discriminant,
$\omega=dx/y$, and $t=-2x/y$.  Let $\Lambda$ be the complex period
lattice for this algebraic differential, and let $e_2^*(\Lambda)$
be the coefficient in the classical sigma quasiperiod formula.
Then this coefficient is algebraic and
\begin{equation}\label{C19-e2-integral}
 e_2^*(\Lambda)\in R.
\end{equation}
For the tame algebraic sigma translates in C18, put
\[
 s_T^{\rm red}(t)=
 \exp\bigl(-e_2^*(\Lambda)\lambda(t)^2/2\bigr)s_T(t).
\]
These are the reduced, or Fock-translated, theta normalizations
with the same torsion lifts as C18.  Every finite product and
quotient of these series of total integer tensor degree $d$ has
the form $Q(t)^d v(t)$ with $v\in R[[t]]^\times$ and
\[
 Q(t)=\exp\bigl(-e_2^*(\Lambda)\lambda(t)^2/2\bigr)
            \frac{\sigma(\lambda(t))}{t}.
\]
For every $v\in R[[t]]$, every integer $d$, and $n\geq1$,
\begin{equation}\label{C19-theta-weak}
 v_p\bigl(D^n(Q^dv)(0)\bigr)
 \geq\max\left\{0,
       \frac{n}{p^2-1}-\frac{s_p(n)}{p-1}\right\}.
\end{equation}
Here $s_p(n)$ is the base-$p$ digit sum.  In degree zero the common
factor cancels, and \eqref{C19-strong} applies directly.
These are respectively slopes $1/(p^2-1)$ and $p/(p^2-1)$;
no equality of the two estimates is asserted.

\paragraph{Proof of the operator congruence over $R/pR$.}
Let $S=R/pR$, and reduce the formal group and $D$ to $S$.
A continuous $S$-linear derivation of $S[[t]]$ is of the form
$b(t)\,d/dt$.  Translation invariance identifies such a derivation
with its tangent value: if $D=w(t)\,d/dt$ and $w(0)=1$, the
invariance identity for the group law gives
$b(t)=b(0)w(t)$.  Thus invariant derivations form the free
rank-one $S$-module generated by $D$.
In characteristic $p$, the iterated Leibniz formula shows that
$D^p$ is a derivation, since its middle binomial coefficients
vanish.  It is still translation invariant.  Consequently
$D^p=H D$ for some $H\in S$.

Choose $a\in\mathcal O$ whose image in $\mathbb F_q$ lies outside
$\mathbb F_p$.  The formal endomorphism $[a]$ has unit tangent
coefficient $a$, and so is a formal automorphism.  The invariant
differential identity $[a]^*\omega=a\omega$ gives
$D[a]^*=a[a]^*D$.  Comparing its $p$-fold iterate with
$D^p=HD$ gives $(a^p-a)H=0$ in $S$.
The element $a^p-a$ has nonzero residue and is a unit in $S$,
including when $S$ is nonreduced.  Hence $H=0$, proving
\eqref{C19-Dp}.  Iteration also gives
\begin{equation}\label{C19-iterated-Dp}
 D^k(R[[t]])\subset p^{\lfloor k/p\rfloor}R[[t]]
       \qquad(k\geq0).
\end{equation}

\paragraph{Completion and linearization of the CM action.}
For completeness, let $I=(p,t)\subset R[[t]]$.
For every integral one-dimensional formal group,
$[p](t)=pt+O(t^2)\in I^2$ and
$[p^r](t)\in I^{r+1}$, by induction: substitution of
$u\in I^m$ into $[p](u)$ gives an element of $I^{m+1}$.
If $a-b\in p^r\mathcal O$, then the formal group identity
$[a]=F([b],[a-b])$ shows that
$[a](t)-[b](t)\in I^{r+1}$.
Approximating an element of $\mathcal O_{K_p}$ by elements of
$\mathcal O$ therefore gives a convergent sequence of formal
endomorphisms.  The ring $R[[t]]$ is $I$-adically complete;
for each fixed coefficient this convergence is $p$-adic.
The limit respects the group law, addition, composition, and
tangent coefficients.  This constructs the asserted continuous
$\mathcal O_{K_p}$ action.

Let $\mu_N\subset\mathcal O_{K_p}^\times$ be the Teichm\"uller
roots of unity.  The integral parameter
\[
 X(t)=\frac1N\sum_{\zeta\in\mu_N}\zeta^{-1}[\zeta](t)
\]
has linear coefficient one and satisfies
$X([\eta](t))=\eta X(t)$ for every $\eta\in\mu_N$.
It is thus an invertible integral change of parameter, in which
$[\zeta](X)=\zeta X$.
Let $L(X)$ be the formal logarithm normalized by $\omega$.
Its identity $L(\zeta X)=\zeta L(X)$ implies that only powers
congruent to one modulo $N$ occur.  In particular,
\begin{equation}\label{C19-log-gap}
 L(X)=X+O(X^q).
\end{equation}

Define integral translation operators $D^{[j]}$ by
\[
 f(F(X,Y))=\sum_{j\geq0}(D^{[j]}f)(X)Y^j.
\]
They preserve $R[[X]]$, are $R$-linear and $p$-adically continuous,
and commute with one another.  The last assertion follows by
comparing the coefficients in
$f(F(F(X,Y),Z))$, which is symmetric in $Y,Z$.
Over $L$, the invariant differential Taylor formula is
\[
 f(F(X,Y))=\sum_{k\geq0}\frac{D^kf(X)}{k!}L(Y)^k.
\]
Therefore \eqref{C19-log-gap} gives
\begin{equation}\label{C19-small-Hasse}
 D^{[j]}=D^j/j!\qquad(0\leq j<q).
\end{equation}
No such identification is used for the higher-index operators.

\paragraph{The torsion power sums.}
The generic $p$-torsion is a one-dimensional vector space over
$\mathcal O/p\mathcal O=\mathbb F_q$.  The special fiber's
geometric $p$-torsion is also an $\mathbb F_q$-module, but it has
at most $p$ points: multiplication by $p$ has degree $p^2$ and
inseparable degree at least $p$, since its differential is zero.
It must consequently have only the identity point.
Thus all the generic $p$-torsion reduces to the identity, and
the formal multiplication map has degree $q$.
For any nonzero torsion point with coordinate $\xi$, all its
nonzero torsion coordinates are $\zeta\xi$, $\zeta\in\mu_N$.

Weierstrass preparation, followed by this description of the
roots over characteristic zero, gives
\[
 [p](X)=U(X)(X^q-\beta X),\qquad
 U\in R[[X]]^\times,\quad
 \beta=\xi^N=-p/\alpha,\quad \alpha=U(0)\in R^\times.
\]
Indeed, the prepared polynomial is monic of degree $q$, and its
simple roots are exactly the $q$ torsion coordinates; comparison
of the derivative at zero gives $-U(0)\beta=p$.
Consequently the torsion power sums are exactly
\begin{equation}\label{C19-power-sums}
 \sum_{u\in E[p]}X(u)^j=
 \begin{cases}
 q,&j=0,\\
 0,&j>0\text{ and }N\nmid j,\\
 N(-p/\alpha)^r,&j=rN, r\geq1.
 \end{cases}
\end{equation}
Only the torsion polynomial has this special form; no assertion
that $[p](X)$ itself equals $pX+\alpha X^q$ has been made.

\paragraph{The trace-zero estimate without a polynomial coordinate.}
Define the torsion-translation trace after a finite extension
containing the torsion by
\[
 \mathcal T f(X)=\sum_{u\in E[p]}f(F(X,X(u))).
\]
All $X(u)$ lie in the maximal ideal.  Taylor substitution and
summation thus converge coefficientwise $p$-adically for every
integral $f$.  If $\mathcal T f=0$, equations
\eqref{C19-power-sums} and \eqref{C19-small-Hasse} give
\[
 0=p^2f+N\sum_{r\geq1}(-p/\alpha)^rD^{[rN]}f,
 \qquad
 D^{[N]}f=\frac{\alpha p}{N}(1+H)f,
\]
where
\begin{equation}\label{C19-H}
 H=N\sum_{r\geq2}(-1)^r p^{r-2}\alpha^{-r}D^{[rN]}.
\end{equation}
This series is an integral continuous operator on $R[[X]]$:
each summand preserves the integral ring and its scalar has
valuation tending to infinity.  The ring is complete for this
$p$-adic topology.  Every translation operator commutes with
$D$ and with $\mathcal T$; hence $H$ does also, and all these
operators preserve the trace-zero submodule.

Since $v_p(N!)=p-1$, multiplication by $N!$ gives on that
submodule
\[
 D^N=p^p u(1+H),\qquad
 u=\frac{\alpha N!}{Np^{p-1}}\in R^\times.
\]
Write $n=mN+r$ with $0\leq r<N$.  The commutation and
trace-zero stability just proved allow iteration, giving
\[
 D^nf=p^{pm}u^m r!D^{[r]}(1+H)^m f.
\]
The last factor is integral, so its value at zero has valuation
at least zero.  Furthermore
\[
 v_p(r!)=\left\lfloor\frac rp\right\rfloor
        =\left\lfloor\frac{rp}{p^2-1}\right\rfloor
          \qquad(0\leq r<p^2-1).
\]
For the second equality, write $r=ap+b$, $0\leq b<p$;
then $rp=a(p^2-1)+(bp+a)$, and
$0\leq bp+a<p^2-1$ in the stated range.
This proves \eqref{C19-trace-zero}.

This is the divisibility argument of Katz,
\emph{Divisibilities, congruences, and Cartier duality},
\S III, equations (3.1)--(3.8), pp.~671--672, and Part I of
the proof of Theorem~3.11, pp.~672--673, applied using
\eqref{C19-power-sums} rather than an exact polynomial
coordinate.  The stronger Cartier-operator congruences used
for the other conclusions of that theorem are not claimed here.

\paragraph{The trace ideal and removal of the trace-zero condition.}
Put $A=R[[S]]$, $B=R[[X]]$, with $S=[p](X)$.
This is a finite free extension of degree $q$.
Weierstrass preparation gives
\[
 [p](X)-S=U_S(X)P_S(X),
\]
where $P_S$ is monic of degree $q$ and $U_S$ is a unit.
Write $\omega=a(X)dX$ with $a(X)$ a unit.
The invariant differential identity implies
\[
 [p]'(X)=p\,\frac{a(X)}{a([p](X))}.
\]
In $B=A[X]/(P_S)$ it follows that
$P_S'(X)=p\,v(X)$ with $v\in B^\times$.

Here is an elementary proof of the exact trace ideal, avoiding
an uncomputed different.  If $\ell_{
\rm top}:B\to A$ takes the coefficient of $X^{q-1}$ in the
unique representative of degree less than $q$, then
\[
 \operatorname{Tr}_{B/A}\bigl(b/P_S'(X)\bigr)
      =\ell_{\rm top}(b).
\]
This identity follows over the fraction field by Lagrange
interpolation at the distinct roots of $P_S$; it then holds
in the fraction algebra for every $b\in B$.
Substituting $vb$ gives
\[
 \operatorname{Tr}_{B/A}(b)/p=\ell_{\rm top}(vb)\in A.
\]
Taking $b=v^{-1}X^{q-1}$ makes the right side one.
We have therefore proved the exact equality
\begin{equation}\label{C19-trace-ideal}
 \operatorname{Tr}_{B/A}(B)=pA.
\end{equation}

For any integral $f$, set
$h=\operatorname{Tr}_{B/A}(f)/p\in A$ and
$g=pf-[p]^*h$.  Since the rank is $q=p^2$,
$\operatorname{Tr}(g)=p\operatorname{Tr}(f)-qh=0$.
The algebra trace, pulled back to $B$, equals $\mathcal T$:
after adjoining the torsion, the roots of
$[p](Y)=[p](X)$ are precisely $F(X,X(u))$.
The equality can first be checked in the separable fraction
algebra and then in the integral power-series ring.
Thus \eqref{C19-trace-zero} applies to $g$.
On the other hand, invariant differentiation gives
\[
 D^n[p]^*h(0)=p^nD^nh(0),
\]
where on the right $D$ is the same invariant derivation in the
target parameter.  Since $D$ preserves the integral ring and
$\lfloor np/(p^2-1)\rfloor\leq n$, the identity
$pf=g+[p]^*h$ proves \eqref{C19-strong}.
The additional lower bound zero follows from integrality of
$D^nf$.  Finally, $D=w(t)d/dt$ with $w(0)=1$ gives
$D^n(t^n)(0)=n!$: terms with fewer than $n$ ordinary derivatives
still vanish at zero.  For $n=p^2-1$, Legendre's formula gives
$v_p(n!)=p-1$, proving the claimed sharpness.
The estimates also extend to $f\in\mathcal O_{\C_p}[[t]]$
on the base-changed fixed formal group.  For fixed $n$, the
value $D^nf(0)$ depends only on the first $n+1$ coefficients
of $f$.  Approximate these by integral algebraic coefficients,
apply the finite-extension assertion, and use continuity and
closedness of each stated valuation inequality.

\paragraph{Integral de Rham CM splitting and $e_2^*$.}
Use the short Weierstrass model specified above and set
$\eta=x\omega$.  C18's integral expansion gives
\[
 x=t^{-2}+O(t^2),\qquad
 \omega=(1+O(t^4))dt,\qquad
 \eta=t^{-2}dt+O(t^2)dt.
\]
On $\mathcal E\setminus O$ the differential $\eta$ is regular.
Near $O$ the differential $\eta-d(-1/t)$ is regular.  These
two differentials and the transition function $-1/t$ define
an integral \v Cech--de Rham class.  The image of this class
in $H^1(\mathcal E,\mathcal O_{\mathcal E})$ pairs with
$\omega$ by the residue $\pm\operatorname{Res}_O(\omega/t)
=\pm1$.  Relative Serre duality and the Hodge exact sequence
therefore show that $(\omega,\eta)$ is an integral basis of
$H^1_{\rm dR}(\mathcal E/R)$.  Equivalently, its de Rham cup
pairing is a unit; the sign of this pairing is immaterial here.
For a curve, the Hodge exact sequence used here follows also
directly from the two-term \v Cech--de Rham complex: the map
$H^1(\mathcal O)\to H^1(\Omega^1)$ induced by $d$ is zero,
as the residue of an exact differential is zero.

Choose the global CM element $a$ above.  Both $a$ and
$a-\bar a$ are units in $R$, since reduction of complex
conjugation on $K_p$ is the $p$-power automorphism of
$\mathbb F_q$.  The integral endomorphism acts as
\[
 [a]^*\omega=a\omega,
 \qquad [a]^*\eta=b\eta+c\omega,
 \qquad b,c\in R.
\]
The cup product is multiplied by $\deg[a]=a\bar a$;
comparison in the integral basis gives $ab=a\bar a$ and
hence $b=\bar a$.

Here is the comparison with the specified analytic coefficient.
Under complex uniformization, $x=\wp(z)$ and $\omega=dz$.
Writing $\zeta=\sigma'/\sigma$, one has
$\zeta'=-\wp$.  Thus the periods of $\eta$ are the negatives
of the sigma quasiperiods.  The classical formula is
\[
 \eta_{\rm per}(\gamma)
    =e_2^*(\Lambda)\gamma+C\bar\gamma
       \qquad(\gamma\in\Lambda),
\]
for the nonzero area constant $C$.
Consequently the periods of
$\eta+e_2^*(\Lambda)\omega$ are $-C\bar\gamma$;
this class is a $\bar a$-eigenvector for $[a]^*$.
Comparison with the integral matrix above gives exactly
\[
 c=(\bar a-a)e_2^*(\Lambda).
\]
It proves algebraicity and \eqref{C19-e2-integral}, since
$\bar a-a$ is a unit.  The source for the quasiperiod formula
and its normalization is Bannai--Kobayashi,
\emph{Algebraic theta functions and the $p$-adic interpolation
of Eisenstein--Kronecker numbers}, arXiv:math/0610163v4,
Example~1.9, pp.~10--11.  That example explicitly gives
$\theta_{\rm red}(z)=e^{-e_2^*z^2/2}\sigma(z)$ and
$\theta_{\rm red}'(0)=1$.

\paragraph{Completion of the reduced theta estimate.}
Let $q_0(t)=\sigma(\lambda(t))/t$; this is the common factor
denoted $q(t)$ in C18 and is distinct from the number $q=p^2$.
C18 proves $s_T=q_0u_T$, with $u_T$ an integral unit, and
\[
 B=D^2\log q_0=-x-D^2\log t\in t^2R[[t]].
\]
The sigma quasiperiod identity also gives, by direct substitution
in the translation formula of C18,
\[
 s_T^{\rm red}(t)=Q(t)u_T(t),\qquad
 Q(t)=e^{-e_2^*\lambda(t)^2/2}q_0(t).
\]
In particular every stated product or quotient has the announced
factorization.  No integrality of the raw series $Q$ is assumed.

Put $l_j=D^j\log Q(0)$.  In the additive coordinate
$z=\lambda(t)$, the series $Q$ and its logarithm are even,
and $Q(0)=1$.  Hence $l_0=l_1=0$, all odd $l_j$ vanish,
$l_2=-e_2^*\in R$, and for even $j\geq4$,
\[
 l_j=D^{j-2}B(0),\qquad
 v_p(l_j)\geq\left\lfloor\frac{j-2}{p}\right\rfloor
\]
by \eqref{C19-iterated-Dp}.
The Bell-polynomial proof in C18 now applies verbatim, allowing
the additional even index $j=2$.  Its required numerical
inequality for that index is
\[
 v_p(2!)-0+\frac1{p-1}
       \leq\frac{2p}{p^2-1},
\]
which holds for $p\geq5$.  For clarity, in each exponential
Bell term the inequality
\[
 v_p(j!)-\left\lfloor\frac{j-2}{p}\right\rfloor
       +\frac1{p-1}\leq\frac{jp}{p^2-1}
       \qquad(j\geq2\text{ even})
\]
and $v_p(m!)\leq m/(p-1)$ bound the total denominator loss
by $np/(p^2-1)$.  All $D^n(Q^d)(0)$ are also integral,
since exponential Bell polynomials have integer coefficients
in the integral variables $d l_j$, including for $d<0$.
Thus
\[
 v_p(D^n(Q^d)(0))
       \geq\max\{0,v_p(n!)-np/(p^2-1)\}.
\]
For $v\in R[[t]]$, \eqref{C19-iterated-Dp} bounds the
coefficient of $z^a$ in $v(t(z))$ below by
$\lfloor a/p\rfloor-v_p(a!)\geq-a/(p(p-1))
\geq-ap/(p^2-1)$.
Convolution gives the same factorial bound for $Q^dv$,
and Leibniz's rule gives integrality of its derivatives.
Legendre's formula yields \eqref{C19-theta-weak}.
The degree-zero assertion follows from cancellation of $Q$
and the already proved \eqref{C19-strong}.
\hfill$\square$

\paragraph{Formal-module generalization and several independent factors.}
The formal-group assertions have a slightly wider scope than
the elliptic application.  Let $R$ be any complete discrete
valuation ring of mixed characteristic $(0,p)$ containing
the integers of the unramified quadratic extension of $\Q_p$.
Let $G/R$ be a one-dimensional smooth commutative formal group
with a strict action of those integers, meaning that the action
of $a$ on the tangent space is multiplication by $a$.
Assume that $[p]:G\to G$ is finite flat of degree $p^2$ and
that its reduction has its entire kernel supported at the
identity.  Equivalently in this setting one may assume that
the closed-fiber formal group has height two.
Then all of \eqref{C19-Dp}, \eqref{C19-trace-zero},
\eqref{C19-trace-ideal}, and \eqref{C19-strong} hold.
For these formal assertions alone $p\geq2$ suffices.

Indeed, the integral action is now supplied as a hypothesis,
and averaging over its Teichm\"uller units gives the same
coordinate $X$.  The prepared $p$-torsion polynomial has
$p^2$ simple roots in characteristic zero, all in the maximal
ideal.  Its roots form a one-dimensional vector space over
$\mathbb F_{p^2}$, so the same orbit and power-sum calculation
applies.  Every subsequent translation-operator, convergence,
invariant-differential and trace argument was purely formal.
The elementary factorial identities remain valid at $p=2,3$.
No elliptic sigma or de Rham assertion at those two primes
is included in this generalization.

The proof in fact gives operator divisibility, not only a
bound at the origin.  With
\[
 b(n)=\max\{0,\lfloor np/(p^2-1)\rfloor-1\},
\]
one has $D^nR[[t]]\subset p^{b(n)}R[[t]]$.
For a product $G_1\times\cdots\times G_r$ of such formal
modules and its independently normalized derivations $D_i$,
it follows for every integral series $f$ that
\[
 v_p\bigl(D_1^{n_1}\cdots D_r^{n_r}f(0)\bigr)
       \geq\sum_{i=1}^r b(n_i).
\]
To justify coefficient rings in this last assertion, apply the
one-variable proof coefficientwise in the other variables:
all translation operators preserve their integral coefficients,
the same $p$-adically convergent sums still converge, and the
finite-free trace identity is preserved by this power-series
base change.  The $D_i$ commute, so their integral operator
divisibilities multiply.

Thus these bounds apply whenever the formal completion of a
good-reduction CM abelian scheme has been decomposed into the
specified one-dimensional height-two factors by its integral
local CM idempotents.  The precise decomposition, its integral
differentials and the theta-section connection are still part
of the input when using this multi-factor consequence.

\paragraph{Scope of the arithmetic bridge.}
The only formal-group hypothesis used here beyond good reduction
is the maximal local unramified quadratic CM action; a CM order
with conductor divisible by $p$ is not covered by this proof.
The arguments apply after every finite extension carrying the
stated good model and normalized differential.  In particular,
they never replace $R/pR$ by its reduced quotient.
The result supplies intrinsic derivative estimates for the actual
tame theta functions with the named differential and torsion
normalization.  A comparison with a differently normalized
harmonic or covariant derivative still uses its explicit analytic
normalization; the bounds do not make that identification by
the existence of an arbitrary line-bundle trivialization.

The de Rham part is a concrete realization of the integral
correspondence and polarization setting, and its algebraic
splitting is computed rather than assumed.  It is compatible
with the Albanese and polarization discussion in Grothendieck's
cote~44, scans~7--8, and with C17's precisely stated cycle-class
specialization formalism.  Those archive passages are not cited
as proofs of the local valuation estimates established here.


\subsection{THETA3-C21: split CM formal factors and higher-dimensional theta ratios}\label{C21}

\paragraph{Statement.}
Let $A$ be a $g$-dimensional abelian variety over a number field,
with CM by an order $\mathcal O$ in a CM field $E$ of degree
$2g$.  Write $F$ for the maximal totally real subfield of $E$.
Let $p\geq5$ split completely in $F$, and suppose every prime
$\nu$ of $F$ above $p$ is unramified and inert in $E$.
Assume that $\mathcal O$ is maximal at $p$.
After a finite extension choose a good abelian scheme
$\mathcal A/R$, with $R$ the integers of a finite extension
of $\Q_p$, and define all CM endomorphisms over it.
Enlarge the field further to contain the CM embeddings whenever
needed.  Then
\begin{equation}\label{C21-formal-product}
 \widehat{\mathcal A}\simeq
       \prod_{\nu\mid p}G_\nu,
\end{equation}
where every $G_\nu$ is an integral one-dimensional formal
module for the unramified quadratic extension of $\Q_p$,
with height-two special fiber.

Choose integral invariant differentials on the factors and
write $D_\nu$ for their dual invariant derivations, normalized
by integral parameters with $D_\nu t_\nu(0)=1$.
For $f\in R[[t_\nu:\nu\mid p]]$ and a multi-index
$\mathbf n=(n_\nu)$, put
\[
 b_p(n)=\max\left\{0,
            \left\lfloor\frac{np}{p^2-1}\right\rfloor-1\right\}.
\]
Then the intrinsic mixed-derivative estimate is
\begin{equation}\label{C21-mixed-bound}
 v_p\left(\left(\prod_{\nu\mid p}D_\nu^{n_\nu}\right)f(0)\right)
       \geq\sum_{\nu\mid p}b_p(n_\nu).
\end{equation}
In particular the right side is at least
$p|\mathbf n|/(p^2-1)-2r$, where $r$ is the number of
nonzero entries in $\mathbf n$.

There is a direct theta application with a concrete integral
hypothesis.  Suppose $s_1,\ldots,s_m$ and
$u_1,\ldots,u_m$ are integral sections of the same power
$\mathscr L^{\otimes k}$ of a line bundle on $\mathcal A$,
and every $u_i$ is a unit at the identity in an integral local
frame.  Their equal-degree ratio
\begin{equation}\label{C21-section-ratio}
 f=\frac{\prod_i s_i}{\prod_i u_i}
\end{equation}
belongs to the integral formal ring and satisfies
\eqref{C21-mixed-bound}.  The same conclusion holds after
an explicitly chosen common prime-to-$p$ isogeny pullback
has trivialized the flat torsion twists of translated theta
sections, provided the resulting sections have the stated
integral normalization and the denominator is a unit at the
identity.  No assertion that all higher-dimensional tame
theta values are automatically units is used.

\paragraph{Integral CM idempotents and their heights.}
The local splitting hypotheses give the exact algebra
\[
 \mathcal O\otimes\Z_p
       =\mathcal O_E\otimes\Z_p
       \simeq\prod_{\nu\mid p}\mathcal O_{E_\nu},
 \qquad E_\nu/\Q_p\text{ unramified quadratic}.
\]
Let $e_\nu$ be its mutually orthogonal idempotents.
These act on $\mathcal A[p^n]$ for every $n$: the action
factors through $\mathcal O/p^n\mathcal O$, and the
idempotents are compatible as $n$ varies.
An idempotent on a finite flat commutative group splits it
into its image and its complementary image, with inverse
maps given by the idempotent and its complement.
Thus
\[
 \mathcal A[p^\infty]=\prod_{\nu\mid p}\mathcal G_\nu,
       \qquad \mathcal G_\nu=e_\nu\mathcal A[p^\infty].
\]
The factors are finite flat at every level; equivalently,
their coordinate algebras are direct retracts of the finite
flat coordinate algebra before splitting.

The rational Betti homology of $A$ is a rank-one $E$-module:
it is a faithful $E$-vector space of rational dimension $2g$.
Using complex uniformization, its $p$-adic completion identifies
the rational Tate module as a rank-one
$E\otimes\Q_p$-module.  Hence every $\mathcal G_\nu$
has height two.
The CM type selects one embedding out of every conjugate
pair of embeddings of $E$.  Because $p$ splits completely
in $F$, the two embeddings extending a fixed $\nu$ form
exactly such a pair.  Therefore
$e_\nu\operatorname{Lie}(\mathcal A)$ is an integral direct
summand of rank one.  Each $\mathcal G_\nu$ has dimension
one, and its $\mathcal O_{E_\nu}$ action on the tangent
line is through one specified embedding into $R$.

Here is why no \'etale part remains in its special fiber.
The Tate module of that \'etale part is a free $\Z_p$-module
of some rank $h_{\rm et}\leq2$, stable under
$\mathcal O_{E_\nu}$.  Its rationalization is an
$E_\nu$-vector space, so $h_{\rm et}$ is even.
It cannot be two: the remaining connected part would have
height zero, making the whole factor \'etale of dimension
zero, whereas its dimension is one.  Thus $h_{\rm et}=0$.
The special fiber is connected.  Equivalently its geometric
$p$-torsion, an $\mathbb F_{p^2}$-module, has only the
identity point.  All the finite kernels are therefore supported
at the identity in the special fiber.

\paragraph{A direct construction of the formal product.}
The completed local CM action exists on
$\widehat{\mathcal A}$ without an algebraization claim for
the idempotents.  In arbitrary integral coordinates
$t_1,\ldots,t_g$, put $I=(p,t_1,\ldots,t_g)$.
Each component of $[p](\mathbf t)$ equals $pt_i$ plus
terms of total degree at least two.  The same induction as
in C19 gives $[p^n](\mathbf t)\in I^{n+1}$ componentwise.
Approximating local CM elements by global ones consequently
converges $I$-adically to formal endomorphisms, respecting
addition, composition and tangent maps.

Choose the initial coordinates so that their tangent
covectors respect the rank-one direct-sum decomposition
defined by the $e_\nu$.  Set
\[
 X_\nu=t_\nu\circ e_\nu.
\]
The linear part of the tuple $(X_\nu)$ is the identity.
The integral formal inverse-function theorem therefore makes
it a system of integral coordinates.  Orthogonality gives
$X_\mu(e_\nu P)=0$ for $\mu\ne\nu$, and
$X_\nu(e_\nu P)=X_\nu(P)$.
Thus the image of $e_\nu$ is precisely the coordinate axis
on which all other $X_\mu$ vanish.  Because the $e_\nu$ are
group homomorphisms and their sum is the identity, the group
law is the product of its restrictions to these axes.
This proves \eqref{C21-formal-product} with each factor
smooth of dimension one.

The connectedness proved above identifies the $p^n$-kernels
of these formal factors with the entire finite flat groups
$e_\nu\mathcal A[p^n]$.  In particular $[p]$ on each
$G_\nu$ is finite flat of degree $p^2$, with height-two
special fiber.  Identifying $\mathcal O_{E_\nu}$ through
its tangent embedding makes its action strict, exactly as
required by the formal-module generalization of C19.
The exact basic references are Tate, \emph{$p$-Divisible
Groups}, \S1.4, p.~160 (connected--\'etale sequence),
\S2.2, Proposition~1, p.~162 with proof pp.~162--164
(connected groups and divisible formal groups), the
dimension definition following (4), p.~164, and
\S2.3, Proposition~3 and Example~(b), p.~166
(invariance of height and dimension under reduction and
the formal completion of an abelian scheme).
The idempotent splitting and coordinate construction have
been given explicitly above.

\paragraph{Mixed derivatives.}
C19 gives the full one-variable operator divisibility
\[
 D_\nu^nR[[X_\nu]]\subset
              p^{b_p(n)}R[[X_\nu]].
\]
Apply it coefficientwise in the other formal variables.
The factors are independent, so their derivations commute,
and successive application multiplies these integral powers
of $p$.  This proves \eqref{C21-mixed-bound}.
For $n>0$, $b_p(n)\geq np/(p^2-1)-2$; summing over the
nonzero entries gives the displayed total-degree consequence.
As in C19, finite-jet approximation extends the assertion to
integral $\C_p$ coefficients on the fixed base-changed group.
Rescaling any differential by an integral unit preserves the
valuation bound.  A nonunit analytic period rescaling must
instead be retained explicitly in the derivative formula.

\paragraph{Integral sections, common frames, and torsion twists.}
Trivialize $\mathscr L$ on the formal neighborhood of the
identity by an integral frame.  The numerator sections in
\eqref{C21-section-ratio} then have integral coefficients,
and every denominator series has a unit constant term.
The quotient is consequently an integral formal function.
Changing the common frame multiplies each section by the
same factor to the same tensor degree, so that factor cancels
exactly.  In the canonical complex theta trivialization this
includes the common quadratic exponential potential.
Thus the canonical theta ratio represents the same formal
function in the exponential coordinates normalized by the
chosen invariant differentials.  The analytic matching is
the degree-zero case of A8.  If the supplied theta sections
are rescaled by constants, their ratio is correspondingly
rescaled by their product quotient $c$; the resulting bound
is \eqref{C21-mixed-bound} plus $v_p(c)$.

For translated sections, $\tau_T^*\mathscr L$ differs from
$\mathscr L$ by the flat line bundle attached to
$\phi_{\mathscr L}(T)\in\mathcal A^\vee$.
If all the translations have prime-to-$p$ torsion order,
choose a common prime-to-$p$ integer $M$ killing these
dual torsion points.  Pullback by $[M]$ trivializes their
flat twists: on rigidified degree-zero line bundles it
corresponds to multiplication by $M$ on the dual variety.
Choose these integral identifications and retain the
integrality and unit-denominator conditions in the statement.
The pulled-back equal-degree ratio is then a rational function
on the isogeny cover whose completion is integral.
This assertion does not require it to descend to the
unpulled-back abelian variety.

On the formal neighborhood $[M]$ is an integral isomorphism,
with inverse having integral coefficients because its tangent
matrix is $M$ times the identity and $M\in R^\times$.
Also
\[
 D_\nu[M]^*=M[M]^*D_\nu.
\]
For a mixed derivative at the identity, passage to the
pullback therefore multiplies its value by
$M^{|\mathbf n|}$, a unit.  The same valuation estimate
holds in the original local translated-theta coordinates
with the specified identifications. \hfill$\square$

\paragraph{Remaining scope.}
This supplies the strong local estimate on an actual
higher-dimensional CM formal group at the stated split-in-$F$,
inert-in-$E$ primes, and on canonical theta ratios meeting
the explicit integral-section hypothesis.  Positive tensor
degree retains the canonical connection; C22 shows why
that issue cannot be removed by declaring every theta
section to have the ratio estimate.  Neither the formal
idempotents nor this analytic estimate construct the
special-cycle correspondence required in C17.


\subsection{THETA3-C22: the weaker slope is sharp for an actual tame theta section}\label{C22}

The stronger integral-function estimate in C19 cannot be extended
with a bounded loss to every individual tame reduced theta section.
The following explicit example determines the optimal possible
linear slope for that proposed extension.

\paragraph{Statement with the actual theta normalization.}
Take
\[
 E:\quad y^2=4x^3-4x,\qquad p=7,\qquad
 \omega=dx/y,\qquad t=-2x/y,
 \qquad T=(0,0).
\]
This curve has good reduction at 7 and CM by $\Z[i]$, and 7 is
inert in $\Q(i)$.  The point $T$ is a nonzero point of order
two.  Let $z_T$ be a half-period representing $T$, with a
compatible torsion lift as in C18, and use the actual Fock
translation
\[
 \sigma_{z_T}(z)=
 e^{-\eta_{\rm per}(z_T)(z+z_T/2)}\sigma(z+z_T).
\]
Here the real-linear quasiperiod map is extended to the
half-period in the indicated formula.  The period lattice is
$i$-stable, so $e_2^*=0$ and this is already the reduced theta
translation.  Normalize its nonzero constant by putting
\[
 F(z)=\frac{\sigma_{z_T}(z)}{\sigma_{z_T}(0)},\qquad
 f_T(t)=F(\lambda(t)),\qquad F(0)=1.
\]
The algebraic value $\sigma_{z_T}(0)$ is a 7-adic unit by
Katagiri--Kobayashi, arXiv:2609.07024v1,
Theorem~2.7, pp.~7--8, and Lemma~3.3(1), pp.~9--10.
Thus this normalization does not change any derivative valuation.
If $D$ is the invariant derivation dual to $\omega$, then for
every integer $r\geq0$, with $n_r=8\cdot7^r$,
\begin{equation}\label{C22-exact}
 v_7\bigl(D^{n_r}f_T(0)\bigr)
      =\frac{7^r-1}{6}
      =\frac{n_r}{48}-\frac16.
\end{equation}
Consequently no bound
$v_7(D^nf_T(0))\geq c n-C$ with $C$ independent of $n$
can hold for any $c>1/48$.
In particular the coefficient $7/48$ from the degree-zero
integral-function theorem cannot hold for this degree-one
theta section, even with an arbitrarily large fixed loss.

\paragraph{Evenness and the logarithmic derivatives.}
The sigma normalization and quasiperiod identities are those of
Bannai--Kobayashi, arXiv:math/0610163v4,
Example~1.9, pp.~10--11.
The automorphism $(x,y)\mapsto(-x,iy)$ pulls $\omega$ back to
$i\omega$, so the lattice is $i$-stable.  The weight-two
transformation of $e_2^*$, or comparison in
$\eta_{\rm per}(\gamma)=e_2^*\gamma+C\bar\gamma$, gives
$e_2^*=-e_2^*$ and hence $e_2^*=0$.

Since $2z_T$ is a lattice element not divisible by two in the
lattice, the sigma multiplier at $2z_T$ is $-1$.
The translation formula and oddness of sigma therefore give
\[
 \sigma(z+z_T)
     =-e^{2\eta_{\rm per}(z_T)z}\sigma(z-z_T),
 \qquad
 \sigma(z_T-z)
     =e^{-2\eta_{\rm per}(z_T)z}\sigma(z+z_T).
\]
After multiplication by the defining exponential factor these
identities say $F(-z)=F(z)$.  Thus $\log F$ is even, with
constant and linear terms zero.
The second derivative has the particularly simple algebraic form
\begin{equation}\label{C22-B}
 \frac{d^2}{dz^2}\log F(z)=-\wp(z+z_T).
\end{equation}
The right side is $-x(T+P(t))$ in the formal parameter at the
identity.  It belongs to $R[[t]]$ after a finite extension
carrying the torsion data: translation by $T$ is integral,
and its residue disc is disjoint from the sole pole of $x$,
because $T$ is nonzero prime-to-7 torsion.

Put $l_j=(d/dz)^j\log F(0)$.  Equation \eqref{C22-B}
and the intrinsic congruence $D^7R[[t]]\subset7R[[t]]$ of
C19 give
\begin{equation}\label{C22-log-bound}
 l_j=0\quad(j\text{ odd}),\qquad
 v_7(l_j)\geq\left\lfloor\frac{j-2}{7}\right\rfloor
       \quad(j\geq2\text{ even}).
\end{equation}
The decisive logarithmic derivative is a unit.  Indeed,
$Dx=y$ and $Dy=6x^2-2$ imply by direct differentiation
\[
 D^2x=6x^2-2,\qquad
 D^4x=120x^3-72x,\qquad
 D^6x=5040x^4-4032x^2+144.
\]
Evaluating at $T$ gives
\begin{equation}\label{C22-unit}
 l_8=-D^6x(T)=-144\in R^\times.
\end{equation}
The same computation gives $l_2=0$ and $l_4=2$, although
only \eqref{C22-unit} is needed for sharpness.

\paragraph{The unique lowest-valuation Bell term.}
Set $s=7/48$.  For a nonzero even logarithmic derivative
define the surplus
\[
 \epsilon_j=v_7(l_j)-v_7(j!)+js-\frac16.
\]
We have $\epsilon_8=0$, and
\begin{equation}\label{C22-strict}
 \epsilon_j>0\qquad(j\ne8\text{ even},\ l_j\ne0).
\end{equation}
Here is the strict inequality in full, rather than an appeal
only to the non-strict Bell bound of C18.
For $j=2,4,6$, \eqref{C22-log-bound} gives
$\epsilon_j\geq(7j-8)/48>0$.
For every even $j\geq10$, use
\[
 v_7(j!)\leq\frac{j-1}{6},\qquad
 \left\lfloor\frac{j-2}{7}\right\rfloor
       \geq\frac{j-8}{7}.
\]
It follows that
\[
 \epsilon_j\geq
   \frac{j-8}{7}-\frac{j-1}{6}
          +\frac{7j}{48}-\frac16
      =\frac{41j-384}{336}>0.
\]
At $j=8$, $v_7(8!)=1$ and \eqref{C22-unit} gives the
claimed equality.  Vanishing logarithmic derivatives simply
contribute no Bell monomials.

Write $F(z)=\sum_n a_nz^n$.  Its coefficient formula is
\[
 a_n=
 \sum_{\substack{m_j\geq0,\ \sum jm_j=n\\j\geq2\ {\rm even}}}
       \prod_j\frac{l_j^{m_j}}{m_j!(j!)^{m_j}}.
\]
For any nonzero summand, add $ns$ to its valuation and use
Legendre's formula.  The result is exactly
\begin{equation}\label{C22-partition}
 \sum_j m_j\epsilon_j
       +\frac16\sum_j s_7(m_j).
\end{equation}
Take $n=8\cdot7^r$.
The partition with $m_8=7^r$ and all other multiplicities
zero has value $1/6$ in \eqref{C22-partition}.
It is the only such partition: any other partition uses an
index $j\ne8$, so the first sum is strictly positive by
\eqref{C22-strict}; the second sum is at least $1/6$,
since some multiplicity is positive.
Thus the coefficient sum has a unique term of least valuation,
and the nonarchimedean triangle rule gives
\[
 v_7(a_n)=-\frac{7n}{48}+\frac16.
\]
Since the base-7 digits of $8\cdot7^r$ sum to two,
$v_7(n!)=(n-2)/6$.  Multiplying the coefficient by $n!$
therefore gives
\[
 v_7(F^{(n)}(0))=
    \frac{n-2}{6}-\frac{7n}{48}+\frac16
    =\frac{n}{48}-\frac16.
\]
The identity $D=d/dz$ in the additive coordinate proves
\eqref{C22-exact}.  Finally, if $c>1/48$, subtracting the
right side of \eqref{C22-exact} from $c n_r$ tends to
infinity.  This proves the asserted failure of every bounded-loss
estimate with the larger slope. \hfill$\square$

\paragraph{A second instance over $\Q(\sqrt{-3})$.}
The same phenomenon occurs on
$E:y^2=4x^3-4$ at $p=5$, with the nonzero point of
order two $T=(1,0)$ and $\omega=dx/y$.
This curve has good reduction at 5, CM by
$\Z[\zeta_3]$, and inert reduction of the CM prime.
The $\zeta_3$-symmetry forces $e_2^*=0$.
The same half-period argument makes the unit-normalized
actual Fock translate $F$ even, and again
$D^2\log F=-x(T+P(t))$ is integral.
Here $Dx=y$, $Dy=6x^2$ and
$D^4x=120x^3-48$, so
$l_6=-D^4x(T)=-72$ is a 5-adic unit.

For completeness the strict surplus calculation is also
explicit in this case.  Put
$\epsilon_j=v_5(l_j)-v_5(j!)+5j/24-1/4$.
The bound $v_5(l_j)\geq\lfloor(j-2)/5\rfloor$
gives $\epsilon_j>0$ for $j=2,4$;
$\epsilon_6=0$; and for all even $j\geq8$ it gives
\[
 \epsilon_j\geq
    \frac{j-6}{5}-\frac{j-1}{4}
          +\frac{5j}{24}-\frac14
       =\frac{19j-144}{120}>0.
\]
The identical Bell-partition argument, now with the
unique multiplicity $m_6=5^r$, yields
\[
 v_5\bigl(D^{6\cdot5^r}f_T(0)\bigr)
       =\frac{5^r-1}{4}
       =\frac{6\cdot5^r}{24}-\frac14.
\]
Thus the obstruction is present for this second CM field
and its stated actual theta normalization as well.

\paragraph{Positive tensor powers.}
In either example the same exact valuation formula holds
with $F$ replaced by $F^d$ for every positive integer $d$
prime to the displayed prime.  Indeed, the normalized
tensor section is still even and has constant term one,
and every logarithmic derivative $l_j$ is replaced by
$d l_j$.  Its valuation, and hence every surplus
$\epsilon_j$, is unchanged.  The same unique least
Bell monomial proves the same equality at every stated
index.  The obstruction therefore persists for these
actual positive tensor powers, with no adjustable
nonintegral scalar normalization.

\paragraph{What this decides.}
C19's weaker estimate for actual reduced theta tensors has
linear coefficient $1/(p^2-1)$, and this example shows that
coefficient is optimal for the class under consideration.
Equal-degree ratios still have the stronger intrinsic bound
because their common canonical frame cancels.
The example rules out transferring that stronger bound to
all individual sections by an unspecified torsion trace or
frame choice.  It makes no assertion that a differently
normalized archive theorem has these same degree, period,
or jet hypotheses; those data must be matched explicitly.


\subsection{THETA3-C23: the actual A6 cycle has an integral model, with a precise KR comparison test}\label{C23}

\paragraph{Statement.}
Let $R$ be a discrete valuation ring with fraction field $L$
of characteristic zero and residue characteristic $p$.
Suppose the A6 data are defined over $L$, the smooth projective
$n$-fold $X$ has a smooth proper model $\mathcal X/R$, and
all the auxiliary abelian varieties $B_i,B_i'$ have abelian
scheme models $\mathcal B_i,\mathcal B_i'$ over $R$.
Write $\mathcal A'=\prod_{i=1}^{2n}\mathcal B_i'$ and
$g=\dim(\mathcal A'/R)$.
After the finite field extensions needed to define the finitely
many maps and divisor classes, the maps $f_i:X\to B_i$ and
all divisor kernels in A6 extend over $R$.

More precisely, write each selected pure cohomological divisor class as
\[
 D_i=\sum_j c_{ij}c_1(L_{ij}),\qquad c_{ij}\in M,
\]
for a number field $M$ and actual line bundles $L_{ij}$ on
$B_i\times B_i'$.
There are line bundles $\mathcal L_{ij}$ extending them and
maps $\mathfrak f_i:\mathcal X\to\mathcal B_i$ extending
$f_i$.  If
\[
 \pi:\mathcal X\times_R\mathcal A'\longrightarrow\mathcal A',
 \qquad
 g_i:\mathcal X\times_R\mathcal A'\longrightarrow
                   \mathcal B_i\times_R\mathcal B_i'
\]
are the projection and the map using $\mathfrak f_i$ and
the $i$th copied factor, put
\begin{equation}\label{C23-integral-cycle}
 \mathcal C=
  \pi_*\left(
    \prod_{i=1}^{2n}
        \left(\sum_jc_{ij}c_1(g_i^*\mathcal L_{ij})\right)
        \cap[\mathcal X\times_R\mathcal A']\right)
       \in\CH^n(\mathcal A')\otimes_\Z M.
\end{equation}
Its generic fibre is exactly the A6 output cycle for the chosen
divisor representatives.  Its refined special fibre is the same
explicit formula with every object reduced modulo $p$.
In particular, the A6 nonzero cohomology class remains nonzero
in prime-to-$p$ \'etale cohomology after specialization.

Here ``integral model'' refers to extension over $R$.
It does not remove the algebraic coefficients $c_{ij}$ or their
denominators.  One may clear them in a fixed coefficient field;
any resulting denominator at $p$ must be retained in a numerical
valuation comparison.

\paragraph{Why the maps extend at the specified good prime.}
An abelian scheme over a discrete valuation ring is the N\'eron
model of its generic fibre.  Its defining mapping property
therefore extends a generic-fibre map from the smooth
$R$-scheme $\mathcal X$ uniquely to $\mathcal B_i$.
It also extends every homomorphism between the auxiliary
abelian varieties uniquely.  These are applications at the
given prime, not spreading assertions obtained by deleting it.
An exact source is Conrad, \emph{Minimal models for elliptic
curves}, author version 21 November 2015, Definition~5.1 and
Example~5.3, p.~13; the latter records the same N\'eron property
for any abelian scheme.

The hypothesis that $\mathcal X$ is smooth remains substantive.
For the corresponding hyperspecial unitary moduli problem,
Li--Zhang, \S11.2, pp.~75--76, gives a smooth model under its
explicit hypotheses (H1)--(H6) and conditions (M1)--(M9).
One must match the selected component, prime and level to those
hypotheses; an arbitrary principal level at $p$ is not
hyperspecial.  Properness of the chosen model is also retained
in the statement above.

\paragraph{Why the divisor kernels extend, and their normalized form.}
For any smooth proper geometrically connected $R$-scheme
$Y$, the restriction map
\begin{equation}\label{C23-Pic}
                         \operatorname{Pic}(Y)
                           \longrightarrow\operatorname{Pic}(Y_L)
\end{equation}
is an isomorphism.  Here is the elementary divisor proof in
the situation used here.  The scheme $Y$ is regular and
integral, and its special fibre is smooth and geometrically
connected, hence geometrically integral.  A line bundle on
$Y_L$ has a nonzero rational section.  Close its finitely many
prime divisors in $Y$; their closures have codimension one.
Regular local rings are factorial, so this closed Weil divisor
is Cartier and defines an extension.  If an extended line
bundle is generically trivial, a generic trivialization gives
a Cartier divisor supported on the special fibre.  As that
fibre is integral, this divisor is an integer multiple of the
whole fibre, which is the principal divisor of a uniformizer.
Thus the line bundle was trivial.  This proves both assertions
in \eqref{C23-Pic}.  The factoriality input is Stacks Project,
Tag~0AG0; the rest of the argument has been given explicitly.
Apply it to $Y=\mathcal B_i\times_R\mathcal B_i'$.

The pure $H^1\otimes H^1$ kernel can, moreover, be represented
by normalized Poincar\'e correspondences.
For an individual line bundle $L$ on $B\times B'$, replace it by
\[
 L^{\rm norm}=L\otimes
       \operatorname{pr}_1^*(L|_{B\times0})^{-1}
       \otimes\operatorname{pr}_2^*(L|_{0\times B'})^{-1},
\]
with the constant fibre at $(0,0)$ rigidified.
Its first Chern class is the mixed K\"unneth component of
$c_1(L)$.  Since each $D_i$ is already purely mixed, this
operation leaves its cohomological class unchanged.
We may choose these normalized representatives before forming
\eqref{C23-integral-cycle}; its exact generic-fibre identity
then refers to that chosen cycle representative.  No equality
of unnormalized Chow classes is inferred from equality of
their cohomology classes.
The fibres of $L^{\rm norm}$ on $\{b\}\times B'$ are
algebraically trivial and the fibre at $b=0$ is trivial.
The universal property of the dual abelian variety yields
a map $h:B\to(B')^\vee$ with $h(0)=0$, hence a homomorphism,
and $L^{\rm norm}$ is the corresponding pullback of the
Poincar\'e bundle.  The exact universal property is Milne,
\emph{Abelian Varieties}, version~2.00 (2008), \S8,
p.~37, Remark~8.7(c); the zero-preserving homomorphism
property is \S16, Corollary~16.2(a), p.~67.

The homomorphism extends by the N\'eron property, and the
Poincar\'e kernel extends by \eqref{C23-Pic}; their compatibility
follows by uniqueness of the line-bundle extension.
Thus this construction does not require lifting an arbitrary
special-fibre divisor or inventing an integral Hodge projector.
It starts from the exact finite list of A6 correspondences.

\paragraph{Verification of the cycle and of specialization.}
The relative dimension of $\mathcal X\times_R\mathcal A'$
is $n+g$.  The product in \eqref{C23-integral-cycle} has
$2n$ first Chern operations, leaving relative dimension
$g-n$; proper pushforward therefore has codimension $n$
on $\mathcal A'$.
On the generic fibre, apply the projection formula to
$f\times1:X\times A'\to A\times A'$.
It identifies \eqref{C23-integral-cycle} with
\[
 p'_*(p^*(f_*[X])\cdot\prod_iq_i^*D_i),
\]
which is the actual A6 construction.  In the selected
cohomological component it is $(-1)^n k\,t_\beta$, with
the same nonzero $k$ and the same Koszul sign as A6.

The special-fibre immersion is the Cartier divisor defined
by a uniformizer.  Its refined Gysin map takes the fundamental
class of the smooth product to the special-fibre fundamental
class, commutes with first Chern operations, and commutes
with proper pushforward.  Consequently it takes
\eqref{C23-integral-cycle} to exactly the stated fibre formula.
Explicit algebraic references are Stacks Project,
Section~42.29, Tag~02T7, especially Lemma~42.29.8,
and Lemma~42.30.5, Tag~0B73, together with
Section~42.26, Tag~0AYA.  The original prime-to-$p$
cycle-class and specialization locators in C17 apply to this
formula.  Smooth proper base change is an isomorphism, so
the nonzero generic class cannot become zero.
CM endomorphism actions also extend, hence the specified
eigencomponent is preserved.  This does not say that the
subalgebra generated by special-fibre divisors has the same
dimension as in characteristic zero.

No proper intersection of chosen effective divisors is needed
to define \eqref{C23-integral-cycle}: first Chern operations
are defined for line bundles.  When a local intersection is
computed using equations, its appropriate derived Koszul
complex must still be used unless regularity of that sequence
has been established.

\paragraph{The exact common category: balanced Poincar\'e biextensions.}
There is a precise moduli interpretation of the normalized
kernels which sharpens the comparison problem.  Let $A_0,A$
be abelian schemes over a base $S$, with actions $\iota_0,\iota$
of the same commutative order $\mathcal O$.  Write $P_A$
for the Poincar\'e biextension on $A\times_S A^\vee$.
For a homomorphism $x:A_0\to A$, set
\[
                   L_x=(x\times1)^*P_A.
\]
This is a line bundle together with its rigidifications on
both zero axes and its compatible addition laws in both
variables.  Then $x\mapsto L_x$ is a natural equivalence
between homomorphisms $A_0\to A$ and isomorphism classes of
such biextensions on $A_0\times_S A^\vee$.
The exact primary source is SGA~7~I, Expos\'e~VIII,
\S\S3.1--3.2, printed pp.~254--255, in particular
\[
 \operatorname{Biext}^1(P,Q;\mathbb G_m)
            \simeq\Hom(P,Q^\vee)
                         \tag{VIII.3.2.2}
\]
and the symmetric formula~(VIII.3.2.4).
Apply this with $P=A_0,Q=A^\vee$ and use abelian biduality.
Section~3.1 identifies the same objects with divisorial
correspondences.  The construction by pullback of the
Poincar\'e bundle is compatible with every base change.
An automorphism of a biextension induces a bilinear map
$A_0\times A^\vee\to\mathbb G_m$, hence is trivial because
$\Hom(A^\vee,\mathbb G_m)=0$; the latter is exactly
formula~(VIII.3.2.1).  Thus the relevant isomorphisms, when
they exist, are unique.

Under this equivalence, $\mathcal O$-linearity of $x$ is
equivalent to the balanced identities
\begin{equation}\label{C23-balanced}
 (\iota_0(a)\times1)^*L_x
       \simeq(1\times\iota(a)^\vee)^*L_x
                         \qquad(a\in\mathcal O)
\end{equation}
as biextensions.  Indeed the left member corresponds to
$x\iota_0(a)$ and the right member corresponds to
$\iota(a)x$, by functoriality of the Poincar\'e bundle.
Here $\iota(a)^\vee$ is the actual dual endomorphism;
there is no additional conjugation in
\eqref{C23-balanced}.  Rosati conjugation enters when the
polarizations are used to form a Hermitian pairing.

Consequently, for a nilpotent immersion $S_0\subset S$
and a fixed $\mathcal O$-linear $x_0$ over $S_0$, the
following lifting problems are equivalent: lift $x_0$ as an
$\mathcal O$-linear homomorphism between the prescribed
abelian schemes, or lift $L_{x_0}$ as an $\mathcal O$-balanced
biextension, with both rigidifications and both addition
laws.  To prove this assertion, associate $L_x$ to a lift
$x$; conversely apply the displayed equivalence over $S$.
Its compatibility with restriction to $S_0$ and its
injectivity identify the resulting homomorphism with $x_0$
after restriction.  Equation~\eqref{C23-balanced} gives
the required linearity in both directions.  The structures
in this equivalence are exact; the refinement below explains
when they can be recovered from an actual line-bundle lift.
An equality of cohomology classes is not such a lift.

With polarizations $\lambda_0,\lambda$, the Gram entries
of a tuple are recovered from these kernels by recovering
their homomorphisms and taking
\[
           T_{ij}=\lambda_0^{-1}x_j^\vee\lambda x_i.
\]
Thus the condition of prescribed Gram matrix, and any
specified framing identifications, can be imposed on
either side without loss of information.  If the
semi-global convention uses $\Hom\otimes\Z_{(p)}$, clear
denominators by a specified integer prime to $p$ and retain
the corresponding Gram rescaling.  On the $p$-divisible
groups such multiplication is invertible, so it does not
alter the local lifting condition.  Multiplication by $p$
is not covered by this argument; a special
quasi-homomorphism with a $p$ denominator still has to meet
the original genuine-homomorphism extension condition.

This identifies an exact category shared by the normalized
A6 kernels and the global special-homomorphism moduli
problem.  To identify their cycles one must still realize
the auxiliary pair as the specified tautological PEL pair,
with the required dimensions, signatures, polarizations,
level and framing, and match its particular kernels.
The equivalence does not provide those identifications.
For an arbitrary normalized divisor representative,
membership in the balanced category for a specified pair of
$\mathcal O$-actions remains a condition to check.

For the rank-one CM factors of A6 this condition can be
arranged explicitly.  Let $B$ have CM by a field $E$,
with $H^1(B,\Q)$ of rank one over $E$, and choose an
$E$-compatible polarization $\lambda:B\to B^\vee$;
its Rosati involution on $E$ is complex conjugation.
Give $A=B^\vee$ the conjugate-dual action
\[
             \iota_A(a)=\iota_B(\bar a)^\vee.
\]
For every $a$ in the acting order, the homomorphism
$h_a=\lambda\iota_B(a):B\to A$ is $\mathcal O$-linear.
Indeed, for $b\in\mathcal O$, Rosati compatibility gives
\[
 h_a\iota_B(b)=\lambda\iota_B(ab)
      =\iota_B(\bar b)^\vee\lambda\iota_B(a)
      =\iota_A(b)h_a.
\]
Let $L_{h_a}$ be its normalized Poincar\'e kernel on
$B\times A^\vee=B\times B$.
It equals $(\iota_B(a)\times1)^*L_\lambda$.
The mixed polarization class $c_1(L_\lambda)$ pairs each
$E$-eigenline indexed by $\tau$ with the one indexed by
$\bar\tau$, and every such component is nonzero by
nondegeneracy of the polarization and the rank-one
hypothesis.

To isolate the first eigenline $\sigma$, write its
idempotent in $E\otimes_\Q\overline\Q$ as
$e_\sigma=\sum_j c_j a_j$, with $a_j$ in the acting
order and $c_j$ algebraic.  Then
\[
                     \sum_jc_jc_1(L_{h_{a_j}})
\]
is exactly the nonzero $\sigma\otimes\bar\sigma$
component of the polarization kernel.  One nonzero
algebraic rescaling makes it the selected pure divisor
class of A6.  Thus one can choose the actual finite
list of divisor representatives from balanced kernels,
with all coefficients and denominators explicit.
After defining the polarization and the finitely many
endomorphisms, each extends over the given good base.
The target action just constructed must still be matched
to the tautological PEL action and its signature;
the construction does not identify them automatically.

\paragraph{On the fixed infinitesimal PEL pair, a genuine line-bundle lift suffices.}
There is a further reduction for the actual deformation
problem.  Let $S$ be Artinian local, let $S_0\subset S$
be a nilpotent closed subscheme, and fix abelian schemes
$A_0,A$ over $S$ with their order actions.  Suppose
$x_0:(A_0)_{S_0}\to A_{S_0}$ is $\mathcal O$-linear.
Then $x_0$ lifts if and only if the actual line bundle
$L_{x_0}$ lifts to $A_0\times_S A^\vee$.
No chosen addition laws on the proposed line-bundle
lift are needed as extra input.

Here is the full proof.  A homomorphism lift gives a
line-bundle lift by Poincar\'e pullback.  Conversely,
normalize any proposed line-bundle lift $L$ by setting
\[
 N=L\otimes\operatorname{pr}_1^*(L|_{A_0\times0})^{-1}
       \otimes\operatorname{pr}_2^*(L|_{0\times A^\vee})^{-1}
       \otimes\pi^*(L|_{(0,0)}),
\]
where $\pi:A_0\times_S A^\vee\to S$.
This formula supplies compatible rigidifications on
both zero axes; it still reduces to the normalized
$L_{x_0}$.  SGA~7~I, Expos\'e~VII, Example~2.9.5,
printed p.~175 (following Corollary~2.9.4, p.~174),
states precisely that, for abelian schemes over any
base, the category of such birigidified line bundles
is equivalent to the category of biextensions.
Its argument uses the theorem of the square for the
two maps to the relative Picard functors in~(2.9.5.1).
Thus $N$ has the required unique biextension structure.
The classification in Expos\'e~VIII now produces a
homomorphism $x:A_0\to A$ reducing to $x_0$.

It remains to prove that $x$ is $\mathcal O$-linear;
this follows from elementary nilpotent rigidity.
If a morphism $h:B\to C$ of abelian schemes over
Artinian local $S$ reduces to zero over $S_0$,
its topological image lies at the zero point of the
special fibre.  Hence it factors through an affine
open neighbourhood of that point in $C$.
Since $H^0(B,\mathcal O_B)=H^0(S,\mathcal O_S)$,
the resulting map to this affine open factors through
$S$.  If $h$ preserves zero, it is therefore the zero
map.  Apply this to
\[
                h=x\iota_0(a)-\iota(a)x
                              \qquad(a\in\mathcal O).
\]
Its reduction is zero because $x_0$ is linear, and it
preserves zero.  Thus $x$ is linear for every $a$.
The same argument proves uniqueness of a homomorphism
lift.  This proves the asserted equivalence.

Consequently the KR obstruction below is also the
obstruction to lifting this specified normalized
Poincar\'e line bundle on the prescribed abelian
pair.  This statement concerns an actual line bundle
and a fixed infinitesimal moduli problem; neither its
first Chern class nor an unspecified auxiliary pair
can replace those data.

\paragraph{The actual KR lifting equation.}
The integral construction above removes an extension problem;
it does not yet impose the special-homomorphism moduli condition.
That condition has an exact local test.
For definiteness take the unramified quadratic local situation
over $\Q_p$ in Kudla--Rapoport I.
Let $A\to A_0=A/I$ be a square-zero thickening in the
nilpotent-$p$ deformation category, with compatible divided
powers on $I$.  Fix an ambient unitary deformation over $A$
whose reduction over $A_0$ belongs to $\mathcal Z(x)$.
Write $\mathbb D(\overline E)$ and $\mathbb D(X)$ for the
covariant Dieudonn\'e crystals evaluated on the thickening,
with Hodge filtrations $F_{\overline E}$ and $F_X$.
The crystalline map induced by the special homomorphism is
$\alpha_x:\mathbb D(\overline E)\to\mathbb D(X)$.
Its obstruction to lifting is
\begin{equation}\label{C23-obstruction}
 F_{\overline E,0}\xrightarrow{\alpha_x}
       \mathbb D(X)_0\longrightarrow(\operatorname{Lie}X)_0.
\end{equation}
Here the subscript $0$ labels the indicated
$\mathcal O_{K_p}$ eigenspace; it is not reduction to $A_0$.
The map reduces to zero over $A_0$.
Both its source and target have rank one, so after choosing
generators it is one element of $I$, and the lift exists
exactly when that element vanishes.
This is the explicit proof of Kudla--Rapoport,
\emph{Special cycles on unitary Shimura varieties I},
Proposition~3.5 and equations~(3.2)--(3.3), pp.~15--16
of the author preprint.  The proposition also proves that a
nonempty $\mathcal Z(x)$ is a relative Cartier divisor.
Li--Zhang, \S\S2.1--2.4, pp.~10--12, fixes the corresponding
polarization, signature, framing and derived-cycle conventions.

Thus a candidate comparison must supply a PEL moduli map
and actual special quasi-homomorphisms whose obstruction
equations match the intended divisors.  A homomorphism between
the fixed auxiliary good abelian schemes already extends on
the entire base.  Its lifting obstruction in that constant
family is identically zero.  Such a fixed-target homomorphism,
by itself, cannot define a nonempty proper KR divisor in a
varying unitary deformation space.

\paragraph{A sufficient equation-level comparison, with derived correction.}
There is a useful reduction of the required verification.
If $D,D'$ are closed subschemes of an $R$-flat scheme and
both are $R$-flat, then equality of their generic-fibre
closed subschemes implies $D=D'$ whenever they are compared
in the same ambient scheme.  Indeed their ideals are saturated
with respect to the uniformizer:
\[
 I_D=I_D[1/\varpi_R]\cap\mathcal O,
       \qquad
 I_{D'}=I_{D'}[1/\varpi_R]\cap\mathcal O,
\]
because the two quotient structure sheaves have no
$\varpi_R$-torsion.  The displayed equality follows immediately.

In particular, suppose a map from a regular comparison space
$T$ to the KR deformation space has been specified, and its
pullbacks of the selected KR divisors are $R$-flat Cartier
divisors.  If chosen flat Cartier representatives of the A6
kernel terms agree with these pullbacks on the generic fibre,
they agree integrally.  Locally their equations differ by units.
The corresponding two-term complexes
$[\mathcal O_T\xrightarrow{s_i}\mathcal O_T]$ are then
isomorphic, as are their derived tensor products.
This is an equality of supported perfect complexes, so it
retains every higher Tor term.
Any subsequently specified proper correspondence and finite
Euler characteristic respects this equality.
If the equations form a regular sequence, the Koszul
complex has only its ordinary quotient as homology; otherwise
the alternating higher-homology terms cannot be discarded.

The flatness condition is exact and necessary for this
generic-to-integral reduction.  In $R[[t]]$, the divisors
$(t)$ and $(\varpi_Rt)$ have the same generic divisor and
the same associated trivial line bundle.
Intersecting with the horizontal section $t=\varpi_R$
gives respectively the quotients $R/(\varpi_R)$ and
$R/(\varpi_R^2)$, of lengths one and two.
Both intersections are regular-sequence intersections, so
the discrepancy is already present before higher Tor occurs.
The second divisor has a vertical component and is not
$R$-flat.  Hence a generic identity or an identity of first
Chern classes does not determine the required arithmetic length.

\paragraph{A full-rank local KR cycle cannot be the horizontal A6 cycle.}
Let the local unitary space have rank $N$ and signature
$(1,N-1)$ in Li--Zhang's convention.
For a full-rank lattice $L$ of special quasi-homomorphisms,
$\mathcal Z(L)$ has empty horizontal part.
Li--Zhang, Lemma~2.10.1, p.~15, proves this directly:
a characteristic-zero lift would give an isogeny
$\overline E^{\,N}\to X$, forcing signature $(0,N)$,
contrary to $(1,N-1)$.
The derived full-rank cycle is consequently supported on
the special fibre.

Pullback, tensor product and proper pushforward over $R$
preserve the property of being generically zero.
Thus no correspondence made from those operations can carry
this full-rank vertical derived cycle to the A6 cycle with
its nonzero generic cohomology class.
A literal cycle comparison must instead use a horizontal
special-cycle stage (necessarily of rank less than $N$),
with a separately identified arithmetic intersection supplying
the final vertical contribution.
The full-rank density derivative computes that derived
arithmetic intersection; it is not itself an integral
representative of the horizontal A6 class.
\hfill$\square$

\paragraph{What remains to be constructed.}
At a prime satisfying the good-model hypotheses, the actual
A6 map and finite divisor kernel now have the explicit
integral realization \eqref{C23-integral-cycle}; this part of
the comparison is available.  The unresolved KR step is a
specific PEL moduli map and a generic divisor identification
whose pullbacks are flat, or a direct identity of the
supported perfect complexes defined by
\eqref{C23-obstruction}.  Those data would make the
derived correction and the normalization testable.
They are not consequences of the automorphic packet label,
the equality of Hodge classes, or smooth proper base change.


\subsection{THETA3-C26: an integral unit on a two-copy isotropic CM direction}\label{C26}

\paragraph{Statement and normalization.}
Use the elliptic CM hypotheses of C19: $p\geq5$ is inert in the
imaginary quadratic CM field $K$, the CM order $\mathcal O$ is
maximal at $p$, and $E$ has a good model over a complete discrete
valuation ring $R$ finite over $\Z_p$.  All CM endomorphisms are
defined, and their action on the chosen integral invariant
differential is the fixed embedding $K\hookrightarrow\operatorname{Frac}(R)$.
Let $t$ be an integral parameter, let $\lambda(t)=t+\cdots$
be its logarithm for this differential, and let $D\lambda=1$.
Use the actual reduced theta normalized by $\theta'(0)=1$ and put
\[
 q(t)=\frac{\theta(\lambda(t))}{t}\in1+t\C_p[[t]].
\]
No integrality of $q$ is assumed.

For every $b\in\mathcal O_{K_p}^{\times}$ the completed CM
action of C19 satisfies
\begin{equation}\label{C26-unit-isogeny}
 U_b(t):=\frac{q([b](t))}{q(t)^{b\bar b}}
           \in1+tR[[t]].
\end{equation}
The exponent on the right is a $p$-adic integer and is interpreted
coefficientwise by the binomial series at the constant term one.
In particular, if $b\bar b=-1$, then
\begin{equation}\label{C26-two-copy-unit}
                  q(t)q([b](t))\in1+tR[[t]].
\end{equation}
Such a $b$ exists: the Teichm\"uller group
$\mu_{p^2-1}\subset\mathcal O_{K_p}^{\times}$ contains an element
with $b^{p+1}=-1$, because $2(p+1)$ divides $p^2-1$.

For actual nonzero tame translated reduced theta sections, write the
factorization already proved in C18--C19 as
\[
 s_T(t)=q(t)u_T(t),\qquad u_T(t)\in R[[t]]^{\times},
\]
after a finite extension carrying the prime-to-$p$ torsion data.
Then the positive-degree two-section pullback
\begin{equation}\label{C26-theta-pullback}
                 f(t)=s_{T_1}(t)s_{T_2}([b](t))
\end{equation}
is an integral unit whenever $b\bar b=-1$.
Consequently
\begin{equation}\label{C26-strong}
 v_p(D^nf(0))\geq
     \max\left\{0,\left\lfloor\frac{np}{p^2-1}\right\rfloor-1\right\}.
\end{equation}
The same statement holds for positive tensor powers of this
pullback.  It concerns a specified isotropic formal direction;
C22's one-copy derivative has no such cancellation.

\paragraph{The global CM isogeny quotient is exactly rational.}
First take a nonzero $a\in\mathcal O$ such that
$a\equiv1\pmod{2\mathcal O}$ and $p\nmid a\bar a$.
Write $d=a\bar a\in\Z_{>0}$; it is odd.
Let $\Lambda$ be the complex period lattice of the algebraic
differential.  Bannai--Kobayashi, arXiv:math/0610163v4,
Theorem~1.7, Definition~1.8 and Example~1.9, pp.~10--11,
give the exact factors of automorphy
\[
 \theta(z+\gamma)
 =\alpha(\gamma)
   \exp\!\left(\pi H(z,\gamma)+\frac\pi2H(\gamma,\gamma)\right)
   \theta(z),\qquad
 H(z,w)=\frac{z\bar w}{\pi A},
\]
where $\alpha(\gamma)=-1$ off $2\Lambda$ and equals one on
$2\Lambda$.  Since $a\Lambda\subset\Lambda$,
$a\equiv1\pmod2$ gives $a\gamma\equiv\gamma\pmod{2\Lambda}$.
Thus $\alpha(a\gamma)=\alpha(\gamma)=\alpha(\gamma)^d$.
Also $H(az,a\gamma)=dH(z,\gamma)$.
All the factors of automorphy therefore cancel in
\[
                 R_a(z)=\frac{\theta(az)}{\theta(z)^d}.
\]
This is a meromorphic function on the elliptic curve, with
\begin{equation}\label{C26-rational-divisor}
       \operatorname{div}(R_a)=[a]^*(O)-d(O),
       \qquad R_a(z)=a z^{1-d}(1+O(z)).
\end{equation}
The function is algebraic over the field of definition of the
CM endomorphism and the differential.  Indeed the displayed
divisor is defined there.  The vector space of rational sections
with this prescribed principal divisor is one-dimensional after
extension to $\C$; ordinary field-extension compatibility of
global sections descends this one-dimensional space.  Fixing the
nonzero Laurent leading coefficient to be $a$ selects a unique
element over the original field.  This also shows that no
uncontrolled transcendental scalar enters $R_a$.

\paragraph{The primitive local quotient is an integral unit.}
The isogeny $[a]$ on the good model has degree $d$ prime to $p$.
Its kernel is finite \'etale, and its nonidentity geometric
points do not specialize to the identity.  Equivalently,
$[a](t)/(a t)\in R[[t]]^{\times}$ and the only kernel point
in the identity formal disc is $O$.
It follows from \eqref{C26-rational-divisor} that the rational
function
\[
              V_a(t)=\frac{R_a(\lambda(t))t^{d-1}}a
\]
has no horizontal zero or pole on that formal disc, and has
constant term one.  Here the notation $R_a(\lambda(t))$ means
the expansion of the algebraic function $R_a$ in the parameter
$t$.

For precision, the completed local ring is the regular factorial
ring $R[[t]]$.  After cancelling the order $1-d$ at the identity,
the divisor of $V_a$ in this local ring is supported, if anywhere,
on its sole vertical height-one prime $(\varpi_R)$.
Hence $V_a=\varpi_R^m v$ for a unit
$v\in R[[t]]^{\times}$ and some integer $m$.
Evaluation at $t=0$ gives $1=\varpi_R^m v(0)$, forcing $m=0$.
Thus $V_a\in1+tR[[t]]$.
The factoriality used here is the regular-local-ring theorem;
an exact reference is Stacks Project, Tag~0AG0.

The logarithm intertwines CM actions:
$\lambda([a](t))=a\lambda(t)$.  Therefore
\begin{align*}
 \frac{q([a](t))}{q(t)^d}
 &=R_a(\lambda(t))\frac{t^d}{[a](t)}\\
 &=V_a(t)\frac{a t}{[a](t)}\in1+tR[[t]].
\end{align*}
This proves \eqref{C26-unit-isogeny} for the specified global
CM elements.  The argument used the actual canonical reduced
theta and its algebraic Laurent normalization, not just its
divisor class.

\paragraph{Passage to the completed CM action.}
For a fixed $b\in\mathcal O_{K_p}^{\times}$, choose
$a_j\in\mathcal O$ with
\[
 a_j\equiv1\pmod{2\mathcal O},\qquad
 a_j\longrightarrow b\quad\hbox{in }\mathcal O_{K_p}.
\]
The Chinese remainder theorem and maximality at $p$ supply
such a sequence, and all its sufficiently late terms have
$p\nmid a_j\bar a_j$.  C19's completed action gives
coefficientwise convergence $[a_j](t)\to[b](t)$ in $R[[t]]$;
also $a_j\bar a_j\to b\bar b$ in $\Z_p$.

For every fixed power of $t$, the coefficient of $q(t)^u$
is a finite sum of binomial polynomials $\binom{u}{m}$
multiplied by fixed coefficients of $(q-1)^m$.
It is therefore continuous in $u\in\Z_p$, even though the
coefficients of $q$ need not be integral.
Composition with $[a_j](t)$ is likewise coefficientwise
continuous, because these series have zero constant term.
It follows that
\[
 \frac{q([a_j](t))}{q(t)^{a_j\bar a_j}}
          \longrightarrow
 \frac{q([b](t))}{q(t)^{b\bar b}}
\]
coefficientwise.  Every series on the left has integral
coefficients and constant term one.  The valuation ring $R$
is closed in its completion, so the limit lies in
$1+tR[[t]]$.  This proves \eqref{C26-unit-isogeny} and
then \eqref{C26-two-copy-unit}.

\paragraph{Actual theta tensors and their differential contraction.}
The integral-unit factorization of the nonzero tame translates comes
from Katagiri--Kobayashi, arXiv:2609.07024v1,
Definition~3.1 and Remark~3.2, p.~9, and Lemma~3.3(1),
pp.~9--10, together with C19's identification of the reduced
Gaussian factor.  Their constants are units in this normalization.
Using \eqref{C26-two-copy-unit}, the right side of
\eqref{C26-theta-pullback} becomes
\[
       \bigl(q(t)q([b](t))\bigr)
          u_{T_1}(t)u_{T_2}([b](t)),
\]
a product of integral units.  C19 now proves
\eqref{C26-strong}.  Positive integral powers remain integral
units, so the same estimate applies to them.

More generally, take units $b_1,\ldots,b_r$ in
$\mathcal O_{K_p}$ and nonnegative integers $d_1,\ldots,d_r$
satisfying
\[
                  \sum_i d_i b_i\bar b_i=0.
\]
Multiplying \eqref{C26-unit-isogeny} with these powers cancels
the entire remaining exponent of $q$.  Hence
$\prod_i s_{T_i}([b_i](t))^{d_i}$ is an integral unit and
satisfies the same strong derivative estimate.
This is a rank-one isotropic direction with the stated
unit coordinates; no assertion for arbitrary nonunit columns
or higher-rank blocks is used.

The map $t\mapsto([b_1](t),\ldots,[b_r](t))$ is an integral
formal group homomorphism.  Its invariant tangent vector is
$(b_1,\ldots,b_r)$, so its $n$th derivative is precisely the
contraction by
$(\sum_i b_iD_i)^n$ of the product theta section at the origin.
This identifies the differential being estimated.  For
$b_i\notin\mathcal O$, the map is formal and need not
algebraize to a map of complex elliptic curves.
The global CM approximants above justify the finite-jet limit.
\hfill$\square$

\paragraph{Relation to the remaining bridge.}
This supplies an explicit positive-degree canonical-frame
cancellation on an actual CM elliptic formal group, including
its primitive integral normalization.  The isotropic condition
is additional geometric data, so it does not conflict with
C22.  Matching a full archive determinant block, its unit minor
and its normalizing factors still requires the exact dictionary
in A8 and the archive comparison.  Nor does this local
theta identity identify the integral cycle with a
Kudla--Rapoport moduli condition.


\subsection{THETA3-C29: actual horizontal boundary special divisors realize the elliptic Poincar\'e kernels}\label{C29}

There is a concrete integral special-cycle realization of the
normalized kernels in C23 in the CM elliptic-power sector.
It uses labelled horizontal boundary special divisors on a
larger unitary Shimura variety. The resulting identity is an
identity of line bundles and actual rational sections over
the integral base. It is not an identification with a
full-rank vertical local intersection.

\paragraph{The precise boundary datum.}
Let $k$ be an imaginary quadratic field of odd discriminant,
and let $\mathcal E$ be a CM elliptic curve with its
$\mathcal O_k$ action and principal polarization over a
base $S\to\mathcal M_{(1,0)}$. One may take
$k=\Q(\sqrt{-7})$. Fix a positive definite self-dual
Hermitian $\mathcal O_k$-lattice $L_0$ of rank $r\geq2$,
and put
\[
                  \mathcal B=\mathcal E\otimes_{\mathcal O_k}L_0.
\]
For a nonzero vector $x\in L_0$, write
$j_x:\mathcal B\to\mathcal E$ for the homomorphism
given by its Hermitian linear functional, with the convention
of the chosen polarization. The additivity relation is
$j_{x+y}=j_x+j_y$. Put $Z(x)=j_x^{-1}(0)$.

These are actual unitary boundary data. Choose the unitary
space with lattice $L_0\perp\mathbb H$, where the
hyperbolic plane has Gram matrix
$\left(\begin{smallmatrix}0&1\\1&0\end{smallmatrix}\right)$.
It has signature $(r+1,1)$ and a rational cusp. In the
notation of Bruinier--Howard--Kudla--Rapoport--Yang,
\emph{Modularity of generating series of divisors on unitary
Shimura varieties}, corrected 2017 author version,
Proposition~3.4.3, pp.~46--48, its boundary abelian scheme
is $B_\Phi=\mathcal E\otimes L_0$ after the stated
CM-base identification. Proposition~3.6.1, pp.~52--53,
identifies each labelled $Z_\Phi(x)$ with $Z(x)$ and
proves that it is an effective Cartier divisor flat over
the CM base. Theorem~3.7.1(4),(8), pp.~53--55,
identifies their sum at a fixed norm with the formal
boundary description of the closure of the corresponding
interior Kudla--Rapoport divisor.

We work with these labelled boundary components. A global
unlabelled fixed-norm divisor sums several such components;
it is not silently replaced by a chosen summand. The formal
boundary chart, and the finite stabilizer quotient when
present, remain part of this comparison. All statements
below may be pulled to a neat level and to the chosen
base $S$.

\paragraph{The normalized line bundle and its section.}
Let $\omega=\operatorname{Lie}(\mathcal E)^{-1}$ and let
$\pi:\mathcal E\to S$. With the principal self-duality,
let $\mathcal P$ be the normalized Poincar\'e bundle on
$\mathcal E\times_S\mathcal E$, and put
\[
       \mathcal Q=\mathcal O_{\mathcal E}(0)\otimes\pi^*\omega.
\]
This is rigidified along zero. The theorem of the square,
in its rigidified form, gives the canonical identity
\[
 \mathcal P\simeq m^*\mathcal Q\otimes
          \operatorname{pr}_1^*\mathcal Q^{-1}\otimes
          \operatorname{pr}_2^*\mathcal Q^{-1}.
 \tag{C29a}
\]
This is the same biextension normalization as C23;
SGA~7~I, VII~2.9.5, p.~175, and VIII~3.1--3.2,
pp.~254--255, identify the rigidified correspondence.

There is a specified integral section $s$ of
$\mathcal Q^{24}$ whose divisor is exactly $24(0)$.
The precise input is the above author version,
Proposition~5.1.4, pp.~69--70: its section is
$\Theta^{24}$ of $J_{0,12}=\mathcal Q^{24}$,
where $\Theta=i\vartheta_1/\eta$ in the stated analytic
convention. That proposition also provides the nowhere
vanishing section of $\omega^{12}$ fixing the associated
Hodge-line normalization. In particular there is no
unrecorded vertical divisor in $s$.

Let $x,y,x+y$ all be nonzero, and set
\[
        \mathcal P_{x,y}=(j_x,j_y)^*\mathcal P.
\]
Then the following rational section is defined on the
dense complement of the indicated divisors:
\[
  s_{x,y}=\frac{j_{x+y}^*s}{j_x^*s\,j_y^*s}
                 \quad\hbox{of }\mathcal P_{x,y}^{24}.
 \tag{C29b}
\]
The exact integral divisor identity is
\[
 \operatorname{div}(s_{x,y})=
       24\bigl(Z(x+y)-Z(x)-Z(y)\bigr),
 \qquad
 \mathcal P_{x,y}^{24}\simeq
       \mathcal O\bigl(24Z(x+y)-24Z(x)-24Z(y)\bigr).
 \tag{C29c}
\]
Consequently
\[
       c_1(\mathcal P_{x,y})
                  =[Z(x+y)]-[Z(x)]-[Z(y)]
                         \quad\hbox{in }\CH^1(\mathcal B)_\Q.
 \tag{C29d}
\]

\textbf{Proof.}
Pull (C29a) back by $(j_x,j_y)$ and take its twenty-fourth
power. It identifies the line of (C29b) with
$\mathcal P_{x,y}^{24}$, with no unspecified scalar.
Each $j_v^*s$ is nonzero as a rational section because
$j_v$ is flat and surjective for $v\ne0$, as used in
the cited proof of Proposition~3.6.1. Its divisor is
$24Z(v)$. Divisors add under tensor products and subtract
under division of nonzero rational sections. This proves
both equalities in (C29c), over the integral base itself.
Taking first Chern classes and dividing by $24$ proves
(C29d). No proper intersection or reducedness of the
individual kernels was required.\hfill$\square$

The power $24$ is deliberately retained: it removes the
theta multiplier and the torsion Hodge-line ambiguity
using the specified integral Jacobi section. At residue
characteristics two or three, dividing a cycle class by
$24$ has a valuation cost. The integral line-bundle and
section statement (C29c) itself has not inverted that
integer. If $x+y=0$, one may replace $y$ by $2y$:
then $x,y, x+2y$ are nonzero, and biextension bilinearity
gives $\mathcal P_{x,2y}=\mathcal P_{x,y}^2$.
Thus the same construction realizes $48c_1(\mathcal P_{x,y})$
as an actual integral combination of horizontal divisors.
If either $x$ or $y$ is zero, the normalized kernel is
trivial and no divisor formula is needed.

\paragraph{An actual C23 kernel, with every label visible.}
Take $L_0=\mathcal O_k^2$ with its standard positive
Hermitian form, so $\mathcal B=\mathcal E^2$. For
$a\in\mathcal O_k\setminus\{0\}$ choose $x,y$ so that
\[
       j_x(P,Q)=[a]P,\qquad j_y(P,Q)=Q.
\]
If the chosen Hermitian convention is linear in the first
entry, $x$ uses the coordinate $\bar a$; this is just the
identity $j_x(P,Q)=\langle(P,Q),x\rangle$.
Their norms are $N_{k/\Q}(a)$, $1$, and
$N_{k/\Q}(a)+1$ for $x,y,x+y$, respectively.
The line $\mathcal P_{x,y}$ is exactly
\[
                     ([a]\operatorname{pr}_1,
                           \operatorname{pr}_2)^*\mathcal P.
 \tag{C29e}
\]
Under the principal polarization this is the normalized
kernel of $h_a=\lambda[a]$ in C23. Its divisor expression
uses precisely the three labelled horizontal boundary
special divisors of these norms. The ambient unitary
variety has signature $(3,1)$; its boundary is a surface.
This is a specified realization, not just an equality of
names for two kinds of divisor.

For powers of $\mathcal E$, an
$\mathcal O_k$-linear homomorphism is a matrix of
endomorphisms $[a]$. Biextension additivity expresses
its normalized Poincar\'e kernel as the tensor product
of the coordinate kernels (C29e), with inverses for
negative coefficients. Apply (C29c) to each nonzero
entry. The selected embedding projectors in C23 are
finite algebraic linear combinations of these classes,
so the same formulas realize those pure kernels with
their exact coefficients. If the auxiliary CM factors
are merely isogenous to such powers, one must specify
the isogenies and retain their pullback and denominator
factors; no integral inverse isogeny is presumed.

\paragraph{Substitution into the actual integral A6 cycle.}
Suppose the auxiliary kernels in C23 belong to this
specified elliptic-power sector, with the required
isogeny factors already fixed. For each term write the
identity (C29c), or its $48$-multiple in the exceptional
zero-sum case, on its appropriate boundary abelian
scheme. Pull the resulting line-bundle identities back
by the actual maps $g_i$ of C23. Substitute their first
Chern identities into its displayed formula for
$\mathcal C$ and expand the finite product.
This expresses the actual cycle $\mathcal C$ as a
finite algebraic-coefficient combination of proper
pushforwards of intersections of these pulled-back
horizontal boundary special-divisor classes.

This assertion is an exact identity in the same Chow
group as C23, because first Chern operators are additive
under tensor product and commute with the indicated
pullbacks and proper pushforward. It is not a new
nonvanishing argument: the nonzero realization is the
one already proved in A6/C23. It now has a specified
special-divisor expression. The factors $24$ and $48$,
the original coefficient field, and every isogeny
denominator remain in the expression.

There is a corresponding statement at the level of
perfect complexes for each effective term. An effective
Cartier divisor $Z(v)$ has the two-term resolution
\[
          [\mathcal O(-Z(v))\xrightarrow{s_v}\mathcal O].
\]
Under an arbitrary comparison map, its derived pullback
is the pulled-back two-term complex. Tensoring several
such complexes retains all their higher Tor groups.
If the pulled-back equations form a regular sequence,
this complex reduces to the usual intersection quotient.
If not, replacing it by the ordinary scheme-theoretic
intersection would lose information. The equality of
first Chern operations above does not assert an equality
between one effective divisor and a signed divisor, or
between their individual structure sheaves.

\paragraph{Why a signed expression is necessary for a pure kernel.}
Let $B_1,B_2$ be positive-dimensional complex abelian
varieties, and let a nonzero class $\gamma$ lie in the
mixed summand $H^1(B_1)\otimes H^1(B_2)$ of $H^2$.
Choose ample classes $\theta_i$ on the factors and put
$g_i=\dim B_i$, $\theta=\operatorname{pr}_1^*\theta_1+
\operatorname{pr}_2^*\theta_2$. Then
\[
       \int_{B_1\times B_2}\gamma\,\theta^{g_1+g_2-1}=0.
 \tag{C29f}
\]
Every summand has odd total cohomological degree on
each factor, whereas a nonzero integral over a factor
requires degree $2g_i$. Thus each term vanishes.
A nonzero effective divisor has positive degree for
an ample polarization: a sufficiently high very ample
power computes that degree by a positive projective
intersection number. Consequently a nonzero pure
mixed kernel class cannot be represented by one
effective divisor. The difference in (C29d) is
therefore essential; no single effective KR divisor
can replace a nonzero pure kernel on this product.

\paragraph{The exact advance and the remaining interior comparison.}
This proves an actual horizontal integral special-cycle
realization for the elliptic-power kernels, on the
specified boundary of a larger unitary integral model.
The fixed-norm sums are the boundary limits of
interior special divisors by the cited toroidal theorem.
The original compact one-active-place ball has no such
rational boundary. Nor does a general CM factor of
dimension greater than one become an elliptic power
merely by calling its polarization a Poincar\'e bundle.
Thus the remaining general comparison is a
correspondence from the chosen compact 032 model to
this formal boundary or to a suitable higher-CM
boundary model, or a direct interior PEL identification.
The construction here supplies the actual horizontal
kernels for its stated sector and preserves the
separate vertical correction in the arithmetic theory.

\framework{The line identities are rigidified divisorial
correspondences in SGA~7~I, VII~2.9.5, p.~175, and
VIII~3.1--3.2, pp.~254--255. Their substitution into
an image cycle is the correspondence construction of
cote~44, scans~7--8. The proper pushforward and
specialization operations use C17's exact SGA/EGA
locators. The signed formula and the derived
pullback distinguish algebraic cycle classes from
local arithmetic lengths.}


\subsection{THETA3-C31: higher-CM kernels from horizontal homomorphism-lifting cycles on an actual boundary}\label{C31}

The elliptic-power restriction in C29 can be removed at a specified
good prime if one allows higher-codimension boundary cycles and a
polarization contraction. The cycles below are actual loci where
a specified homomorphism lifts to the tautological semiabelian
extension. Their realization as closures of interior special cycles
is a further comparison, not part of the assertion.

\paragraph{Data and assertion.}
Let $B/R$ be an abelian scheme of relative dimension $d$, where
$R$ is a characteristic-zero discrete valuation ring with residue
characteristic $p$. Suppose its generic fibre has rank-one CM by
the maximal order $\mathcal O_E$, and fix an $E$-compatible
polarization $\lambda$ of degree prime to $p$. Assume $p$ is
unramified in $E$. Fix a rigidified line bundle $Q$ inducing
$\lambda$. Give $B^\vee$ the conjugate-dual action
$a\mapsto\iota_B(\bar a)^\vee$, so that $\lambda$ is
$E$-linear; character-lattice actions below use the same
involution convention as Lan. These hypotheses hold for the chosen CM data outside
an explicitly retained finite set of primes. Neither a prime
dividing the polarization degree nor a ramified PEL prime is
silently included.

Let $P_\lambda$ be the normalized Poincar\'e kernel on $B\times B$
and, for $a\in\mathcal O_E$, put
\[
 P_a=([a]\operatorname{pr}_1,\operatorname{pr}_2)^*P_\lambda.
 \tag{C31a}
\]
On $B^2\times B$, with coordinates $(P,Q',T)$, define
\[
 \begin{aligned}
 j_+(P,Q')&=[a]P+Q',\\
 j_1(P,Q')&=[a]P,\qquad j_2(P,Q')=Q',\\
 \Gamma_j&=\{T=j(P,Q')\}.
 \end{aligned}
 \tag{C31b}
\]
Each $\Gamma_j$ is smooth over $R$, is a regular immersion of
codimension $d$, and is a labelled component of an actual
horizontal homomorphism-lifting locus on a prime-to-$p$ isogeny
pullback of a toroidal boundary stratum of a larger unitary PEL
model. If $\pi_{12}:B^3\to B^2$ is projection, then
\[
 c_1(P_a)=\pi_{12*}\left(
   ([\Gamma_{j_+}]-[\Gamma_{j_1}]-[\Gamma_{j_2}])
                  \cdot\operatorname{pr}_3^*c_1(Q)\right)
                  \quad\hbox{in }\CH^1(B^2).
 \tag{C31c}
\]
This is an integral Chow identity on the indicated pullback.
Its graph construction has no factor $24$. It does not provide
the separate rational theta section of C29. Descent from the
isogeny pullback retains its actual degree; that degree is a
unit at the specified prime.

\paragraph{The extension obstruction, as a scheme-theoretic locus.}
For an abelian scheme $A/S$, the normalized Poincar\'e bundle
on $A\times A^\vee$, regarded as a $\mathbb G_m$-torsor,
is the universal extension of $A$ by $\mathbb G_m$ over
$A^\vee$. The addition law in its first variable gives the
group law. The relative Picard universal property and the
rigidification identify its class with the parameter itself.
This is the relative form of Milne, \emph{Abelian Varieties},
Propositions 9.2--9.3, printed p.~41; its biextension and
base-change formulation is SGA~7~I, VII~2.9.5, p.~175,
and VIII~3.1--3.2, pp.~254--255.

More generally an extension $G$ of $A$ by a split torus with
character lattice $X$ is classified by
$c:X\to A^\vee$: push out by each character and apply the
preceding construction. For $h:B\to A$, a lift
$\widetilde h:B\to G$ exists exactly when the pulled-back
extension is split, equivalently
\[
                        h^\vee c=0.                 \tag{C31d}
\]
The assertion is functorial under every base change, so it
defines the actual zero fibre of the morphism in (C31d),
including its scheme structure. A splitting is unique because
$\Hom(B,\mathbb G_m)=0$, hence $\Hom(B,T)=0$ for a split
torus. One can verify the first equality directly: a map
from the proper geometrically connected abelian scheme to
the affine group factors through the base, and a group map
must preserve zero. If $G$ and $h$ carry compatible
$\mathcal O_E$-actions, this uniqueness also makes the lift
$\mathcal O_E$-linear: the two composites with any element
of the order lift the same homomorphism and therefore agree.

\paragraph{An actual unitary boundary carrying these extensions.}
Use the zero-dimensional unitary PEL datum of $A=B^3$ with
its product polarization and its full CM type. Take the
rank-one $\mathcal O_E$-lattice $\mathcal Q=\mathcal O_E$
in Lan, \emph{Toroidal compactifications of PEL-type Kuga
families}, corrected author version, Section~3A,
pp.~27--30. Its construction adds the two dual isotropic
lattices to $H_1(A)$ and gives a larger PEL datum with a
rational boundary. A top-dimensional cone gives an actual
boundary abelian stratum, prime-to-$p$ isogenous to
$\Hom_{\mathcal O_E}(\mathcal Q,A)=A$. The neat level
conditions in that section make it an abelian scheme,
rather than an unspecified torsor. The tautological
semiabelian extension on this stratum is described in
Section~3B, p.~31, and in the paragraph before Lemma~4.9
and its proof, p.~46. Its extension parameter is
$c:X\to A^\vee$, with $X$ of rank one over $\mathcal O_E$.
The order action and the polarization are the given ones.

This use includes the author's errata items (4), (8), (13),
and (23): the finite \'etale cover of the smaller
moduli space, the actual torus-torsor chart, both level
conditions and the reduction to principal level are retained.
No directedness assertion about the set of compactifications
is needed. The unramified and prime-to-$p$ polarization
hypotheses above place the construction in Lan's good-prime
base. After acquiring a sufficiently fine prime-to-$p$ level,
pull the boundary stratum to the specified CM base $R$.

Here is why the isogeny ambiguity does not change the claimed
labelled graphs. Write $C_\Phi$ for the boundary parameter
abelian scheme. Forgetting the auxiliary parameter and level
defines
\[
 F:C_\Phi\longrightarrow
       \Hom_{\mathcal O_E}(X,A^\vee)^\circ.
\]
This is an actual prime-to-$p$ isogeny. Indeed the polarization
relation is $c\phi_{\rm cusp}=\lambda_Ac^\vee$. The chosen
cusp lattice map and $\lambda_A$ are isogenies prime to $p$;
solving this equation for $c^\vee$, and acquiring the specified
prime-to-$p$ level lifts, are finite \'etale isogenies.
Taking the identity component preserves surjectivity onto the
connected parameter variety. The component and pairing
conditions are precisely those in Lan's construction and its
corrected principal-level description in Lemma~4.9.

A rank-one projective lattice $X$ is free locally at all
primes over $p$. Choose a generator at those finitely many
primes and clear denominators only away from $p$. This gives
an $\mathcal O_E$-linear injection $i:X\hookrightarrow
\mathcal O_E$ of index prime to $p$. Define explicitly
\[
        \varphi(P)(x)=i(x)\cdot\lambda(P).
\]
The dual action is the conjugate-dual one fixed above.
This formula and the prime-to-$p$ degree of $\lambda$ give
a prime-to-$p$ isogeny
\[
       \varphi:B\longrightarrow
               \Hom_{\mathcal O_E}(X,B^\vee)^\circ.
\]
For $A=B^3$ take its third power. Choose an integer $m$
prime to $p$ annihilating $\ker F$.
Then $[m]F^{-1}$ is an actual homomorphism. Composing with
$\varphi^3$ supplies an actual isogeny pullback on which
the extension parameter is $[m]\varphi^3(P,Q',T)$.
All finite indices and isogenies in this construction are
recorded choices, not integral inverse isogenies.

For any of the three maps $j$ in (C31b), let
$h_j:B\to B^3$ be the polarization adjoint of
\[
                 (P,Q',T)\longmapsto j(P,Q')-T.
\]
It is an actual $\mathcal O_E$-linear homomorphism.
For example the first is the column $(\bar a,1,-1)$;
the other two are $(\bar a,0,-1)$ and $(0,1,-1)$.
The product polarization and Rosati conjugation verify
these formulas directly. Equation (C31d) on the isogeny
pullback is now
\[
               [m]\varphi\bigl(j(P,Q')-T\bigr)=0.
 \tag{C31e}
\]
The kernel of $[m]\varphi$ is finite \'etale over
$R$. After a finite \'etale extension supplying its
labels, (C31e) is a disjoint union of translates of the
graph $T=j(P,Q')$. Its zero-labelled component is exactly
$\Gamma_j$. In particular it is flat and smooth, with
no vertical component or hidden intersection multiplicity.
This proves the boundary and homomorphism-lifting assertions.
The zero label is essential: the entire unlabelled
lifting locus is not silently replaced by this component.

\paragraph{The polarization contraction.}
The rigidified theorem of the square gives
\[
 P_a\simeq j_+^*Q\otimes(j_1^*Q)^{-1}
                         \otimes(j_2^*Q)^{-1}.
 \tag{C31f}
\]
If $i_j:B^2\hookrightarrow B^3$ is the graph immersion,
then $\pi_{12}i_j=\mathrm{id}$ and
$\operatorname{pr}_3i_j=j$. The projection formula for
the first Chern operator therefore gives
\[
 \pi_{12*}\bigl([\Gamma_j]\cdot
                  \operatorname{pr}_3^*c_1(Q)\bigr)
     =c_1(j^*Q)\cap[B^2].
 \tag{C31g}
\]
Taking the signed sum and using (C31f) proves (C31c).
Every graph is a section of the smooth morphism
$B^3\to B^2$, so it is a regular immersion of codimension
$d$ and is isomorphic to $B^2$. This proves its remaining
geometric properties without a transversality assumption.
The codimension count is $(d+1)-d=1$ after pushforward.

Apply (C31c) to the finite list of $a$'s in the CM
embedding projector of C23, retaining its algebraic
coefficients. It expresses every selected pure kernel of
A6 as a finite sum of these specified boundary
homomorphism-lifting cycles contracted with a polarization.
Substitution in C23's formula for $\mathcal C$ now works
for arbitrary rank-one CM factors satisfying the good-prime
hypotheses, not only elliptic powers. The first Chern
operations and proper pushforwards are exactly those in
that formula. For an arbitrary further pullback, use the
derived pullback of the regular graph's Koszul complex;
otherwise excess Tor terms could be lost. The Chow identity
itself requires no ordinary proper-intersection assertion.

\paragraph{Remaining comparison.}
These are actual labelled horizontal cycles on a specified
unitary boundary isogeny chart, with a precise moduli
meaning. They are of codimension $d$ before the
polarization contraction. When $d>1$, one cannot replace
their kernels by divisors by analogy with an elliptic
curve. The argument has not identified each boundary
lifting locus with the full formal completion of a
particular interior Kudla--Rapoport cycle. That requires
compatibility with the remaining degeneration data and
its generic-fibre special-homomorphism condition. Nor
has it supplied an interior correspondence from the
original compact 032 ball to the larger model. Those
are the remaining comparisons; the higher-CM boundary
host and its horizontal kernel contraction are now
specified.

\framework{The universal extension and its splitting
obstruction are Poincar\'e/biextension constructions in
SGA~7~I, VII~2.9.5, p.~175, and VIII~3.1--3.2,
pp.~254--255. Lan's Section~3B uses Raynaud's
semiabelian quotient and extension results in SGA~7~I,
IX~1.4 and IX~2.4. Formula (C31g) is an explicit
divisorial correspondence and pushforward of the kind
recorded in cote~44, scans~7--8; it uses C23's exact
first-Chern and specialization framework.}


\subsection{THETA3-C33: the higher-CM graph is a labelled horizontal interior special-cycle closure}\label{C33}

The remaining formal comparison in C31 can be made on its
specified good-prime toroidal chart. The result concerns the
labelled horizontal closure. It does not identify additional
vertical components of an integral special-homomorphism
moduli space, nor a correspondence from the original compact
ball to this larger PEL model.

\paragraph{Statement and the precise closure.}
Keep all the data and prime-to-$p$ choices of C31. Write
$A=B^3$ for the abelian part at the cusp, and let
$G^{\natural}$ be its Raynaud extension by the cusp torus.
The associated degeneration datum also contains a lattice
$Y$, a period map $\iota:Y_\eta\to G^{\natural}_\eta$,
and the polarization and level structures. Work on the
finite \'etale formal pullback of a smooth toroidal chart
which labels the kernel of $[m]\varphi$ in C31. Finite
\'etale covers of its closed stratum extend uniquely to
its formal completion, so this is an actual chart of the
same formal PEL model.
This is the finite \'etale equivalence for henselian
pairs, Stacks Project, Lemma~15.13.2; at each
nilpotent stage it is Lemma~15.11.2, Tag~0ALI.

For one of C31's columns
\[
 h_+=(\bar a,1,-1),\qquad
 h_1=(\bar a,0,-1),\qquad h_2=(0,1,-1):B\longrightarrow A,
                                                        \tag{C33a}
\]
select the zero-labelled locus where $h^\vee c=0$.
It is precisely the corresponding graph $\Gamma_j$ in
C31, pulled back in the toroidal directions. Then this
formal graph is the scheme-theoretic closure of the
labelled characteristic-zero interior special-homomorphism
locus. Its Hermitian norm is, respectively,
\[
              a\bar a+2,\qquad a\bar a+1,\qquad 2.
                                                        \tag{C33b}
\]
In particular C31's graph contraction expresses the actual
Poincar\'e class using these labelled horizontal interior
closures, not just abstract boundary loci.

To specify ``horizontal'' without ambiguity, choose local
smooth cone coordinates $q_1,\ldots,q_r$, and put
$q=q_1\cdots q_r$. The abelian interior is obtained by
inverting $q$. In a mixed-characteristic chart we take
the closure from the locus where both $p$ and $q$ are
invertible. Thus the closure ideal is saturated by $pq$.
If a chosen integral interior component is already its
horizontal closure, one may equivalently take its
$q$-saturated closure. We do not assert that arbitrary
$p$-vertical components disappear under $q$-saturation.

\paragraph{The exact fully faithful theorem.}
The source is Lan's corrected author revision of
\emph{Arithmetic compactifications of PEL-type Shimura
varieties}, Theorem~4.4.16, printed pp.~254--255, and
especially Corollary~4.5.5.5 with its complete proof,
p.~303. The latter explicitly treats an arbitrary
homomorphism of generic abelian varieties: it gives
uniquely a morphism of degeneration data
\[
 (f_Y,f^{\natural}):(Y_1,G_1^{\natural})
                     \longrightarrow(Y_2,G_2^{\natural}),
 \qquad f^{\natural}\iota_1=\iota_2f_Y.                 \tag{C33c}
\]
It is not restricted to polarization-preserving
isomorphisms. Its hypotheses, stated in Section~4.1,
printed pp.~207--208, are a noetherian normal domain
complete for an ideal, with the indicated semiabelian
and isotrivial-torus hypotheses. The normal completed
charts below satisfy them. Corollary~4.5.5.5 identifies
its assertion with Faltings--Chai, Chapter~III,
Theorem~7.1. We use the corrected author's statement
and proof, not an inference from an equivalence of
polarized groupoids alone.

The compatibility with duality and polarization is
explicit in Lan's Construction~4.5.4.27 and
Lemma~4.5.4.28, printed pp.~299--300. The actual PEL
formal family is the Mumford family under Lan's
\emph{Toroidal compactifications of PEL-type Kuga
families}, Theorem~1.41(5), pp.~16--17. Hence applying
(C33c) compares to the actual interior abelian family
on that chart.

\paragraph{There is no extra source-lattice condition.}
Let $S=\operatorname{Spec}(R')$ be one of the normal
complete chart bases in the preceding theorem. The
good-reduction source $B/S$ has degeneration datum
\[
                         [0\longrightarrow B].
\]
A morphism from it to $[Y\longrightarrow G^{\natural}]$
therefore consists exactly of a group homomorphism
$\widetilde h:B\to G^{\natural}$: the lattice map
$0\to Y$ is zero, and the compatibility in (C33c)
is vacuous. Such a lift of $h:B\to A$ exists if
and only if the pulled-back torus extension splits.
By C31 this is precisely
\[
                              h^\vee c=0.              \tag{C33d}
\]
The splitting is unique because $\Hom(B,T)=0$.
It is automatically $\mathcal O_E$-linear by uniqueness.
Thus the fully faithful theorem yields a unique actual
homomorphism $f:B\to G$ of degenerating group schemes,
and in particular its interior abelian-fibre
homomorphism, exactly on (C33d).

For clarity the dual period condition gives nothing
additional either. The dual datum has the form
$[X\to G^{\vee,\natural}]$, with abelian quotient
$A^\vee$ and projected period map $c:X\to A^\vee$.
The dual map to $[0\to B^\vee]$ has semiabelian
extension part $h^\vee\circ\pi^\vee$ and zero lattice
map. Its period condition is exactly $h^\vee c=0$,
again (C33d). Consequently the target's period
trivialization and toroidal $q$ parameters do not
supply a further equation for this source. The
polarization relation between $c$ and $c^\vee$ is
retained in the given target degeneration datum;
it is not discarded or separately postulated.

\paragraph{The norm is exactly the abelian-part norm.}
Functoriality, duality and the polarization construction
identify the abelian part of $f^\vee\lambda_G f$ with
$h^\vee\lambda_Ah$. Its source and target both have
zero torus and lattice parts, so this determines the
whole homomorphism of good abelian schemes. Thus
\[
       \lambda_B^{-1}f^\vee\lambda_Gf
          =\lambda_B^{-1}h^\vee\lambda_Ah.              \tag{C33e}
\]
Here the inverse polarization is an equality of
quasi-homomorphisms; its right side in the three
cases is the actual integral endomorphism in
(C33b). Indeed $\lambda_A$ is the product
polarization and its Rosati involution sends
$a$ to $\bar a$, so the three matrix products
are $a\bar a+1+1$, $a\bar a+1$, and $1+1$.
All are totally positive. No factor from the
period lattice, the torus, or the prime-to-$p$
boundary isogeny changes these norms. The
isogeny changes the extension parameter as
recorded in C31, and its zero-labelled graph
has already accounted for that change.

\paragraph{Scheme structure of the horizontal closure.}
We give the argument for one graph; it applies to
all three. Let $\mathfrak S$ be the selected formal
chart. \'Etale locally its toroidal direction is a
smooth formal torus embedding over the boundary
abelian parameter scheme. The graph $\mathfrak Z$
is the inverse image of $\Gamma_j$. Its equation
is independent of those directions. Therefore
$\mathfrak Z$ is smooth over the original
mixed-characteristic base, is normal and flat,
and multiplication by each $q_i$ is injective
on its local rings. In particular its ideal
$I_{\mathfrak Z}$ is saturated for both $p$
and $q$.

On each normal complete local ring of
$\mathfrak Z$, (C33d) splits the extension,
so (C33c) constructs the actual interior
homomorphism with norm (C33b). Uniqueness
makes these homomorphisms agree on overlaps.
This gives the inclusion of $\mathfrak Z$
in the closure of the selected interior
locus: its open part with $pq$ invertible
is schematically dense.

We verify the converse with the actual
scheme structures. In characteristic zero
the labelled special-homomorphism moduli
locus is smooth after neat level. Here is
an infinitesimal justification, which is
stronger than a description of its set
of complex points. Write $t=f^\dagger f$
for the totally positive element (C33b),
where $f^\dagger=\lambda_B^{-1}f^\vee\lambda_G$.
It is invertible in $E$, and
\[
                  e=f t^{-1}f^\dagger
\]
is a self-adjoint idempotent in
$\operatorname{End}^0_{\mathcal O_E}(G)$.
It splits the target up to isogeny into
the fixed factor $B$ and its orthogonal
abelian complement $C$. This assertion
works over arbitrary local Artinian
characteristic-zero test rings: clear
one integer denominator, and take the
images of the two complementary integral
endomorphisms. Equivalently take their
connected kernels. The resulting product
map is an isogeny, and its finite kernel
is \'etale because every nonzero integer
is invertible on the test ring.
The image and connected-kernel construction
is valid over these Artinian bases: the
formal logarithm in characteristic zero
identifies the completed group law with
its Lie algebra, where the idempotent is
an actual projection. Its kernel is
formally smooth of the constant required
rank; the connected proper kernel is
therefore an abelian scheme.

After one finite \'etale torsion level,
that kernel and all its action, isotropy
and level constraints have unique lifts
across every nilpotent thickening. The
polarization on the first factor is the
fixed polarization $\lambda_B t$; its
mixed term with $C$ is zero because
$e$ is self-adjoint. Consequently the
deformation functor of the target
together with the labelled map $f$
is exactly the deformation functor of
the complementary PEL abelian variety,
with those discrete isogeny data fixed.
Conversely, deforming the complement
and taking the quotient by the uniquely
lifted finite kernel reconstructs the
target, its polarization and $f$.
These two constructions are inverse
on Artinian test rings and on their
isomorphisms.

The complementary PEL functor is smooth
in characteristic zero. In the orientation
selected by the source CM type, its
signatures are $(3,1)$ instead of
$(4,1)$ at each real place; the conjugate
orientation reverses both pairs. This
is the ordinary smooth unitary PEL
moduli problem, also supplied by Lan's
Theorem~1.8 in the cited Kuga-family
paper. Thus the actual labelled locus
is smooth and reduced, of codimension
$d=[E^+:\Q]$. The unramifiedness of the
Hom functor by rigidity is used only
to choose its local labelled closed
immersion; it is not used to deduce
reducedness. This last localization is
the \'etale local structure of unramified
morphisms, Stacks Project, Lemma~41.17.1,
Tag~04HH.

Its scheme-theoretic horizontal closure
is reduced: locally its ideal is the
contraction of a radical ideal after
inverting $pq$, which is radical.
The rings in question are excellent;
completion along the boundary preserves
reducedness, and normalization is finite.
For exact algebra references, completion
of a G-ring is a regular map by Stacks
Project, Lemma~15.51.14, Tag~0AH2;
reducedness and normality ascend under
such a map by Lemmas~10.163.6 and
10.163.8, Tags~0C21 and~0C22.
Finite normalization follows from
Lemma~10.162.10(4), Tag~032Y, or from
the Nagata property of excellent rings.
Normalize a component of this closure
which meets the selected formal chart,
and complete at a boundary point. After
choosing a local domain component this
is a normal complete base to which
(C33c) applies. Its generic interior
homomorphism extends uniquely, and its
Raynaud abelian-part homomorphism is the
specified label $h$. Indeed equality
with that boundary label holds on the
closed fibre, and rigidity gives equality
on every Artinian thickening; separatedness
and completeness give equality over the
complete base. This is the same
normalized homomorphism-rigidity argument
used in C23. Consequently (C33d)
holds identically on that normal component.
It therefore maps into $\mathfrak Z$.
All components of the reduced closure
satisfy this, so the closure ideal contains
$I_{\mathfrak Z}$. Together with the first
inclusion this proves equality of the
reduced closed formal subschemes. Equivalently,
their ideals agree after inverting $pq$,
and both are $pq$-saturated.

This proof retains the label: other finite
kernel labels give their corresponding
translated graphs, and other special
homomorphisms are not identified with the
zero label. It also retains the horizontal
qualification. In particular it gives no
formula for a derived vertical intersection
length, and does not change the sign or
coherence constraints in C06 and the
arithmetic Siegel--Weil comparisons.

\paragraph{Consequences for the integral kernel formula.}
Let $i$ denote the inclusion of the
closed boundary stratum. The $q_i$
form a regular sequence on every
closure just constructed, so its
refined boundary pullback is exactly
$[\Gamma_j]$, with no excess term.
Substitute these refined restrictions
for $\Gamma_{j_+},\Gamma_{j_1},\Gamma_{j_2}$
in (C31c). The actual integral Chow
identity and polarization contraction
therefore remain an identity on the
boundary abelian parameter scheme,
now expressed as restrictions of the
specified interior horizontal closures.
This does not assume a proper projection
forgetting $T$ on the entire toroidal
chart. The C23/A6 finite CM
projector therefore uses specified
labelled interior horizontal special
cycles in this larger unitary model,
with the exact prime-to-$p$ descent
indices already retained. This closes
the degeneration-data condition left
open in C31. An interior correspondence
from the original compact 032 ball,
and a comparison of vertical derived
intersection numbers, remain additional
geometric requirements.

\framework{The lifting obstruction is the
Poincar\'e extension class from SGA~7~I,
VII~2.9.5 and VIII~3.1--3.2, as in C31.
The new input here is the fully faithful
Mumford degeneration theorem for arbitrary
homomorphisms, Lan Corollary~4.5.5.5,
p.~303. BHKRY Remark~3.6.2, printed
p.~53, describes the elliptic instance
as homomorphisms into the universal
one-motive with its Hermitian norm;
Theorem~3.7.1(8), p.~54, identifies its
formal closures. The proof above keeps
the higher-CM source and proves the
labelled horizontal comparison in the
specified chart.}


\subsection{THETA3-C34: a relative PEL boundary correspondence recovers the chosen ball-derived cycle}\label{C34}

The original compact ball can be included in the boundary
comparison of C31--C33. The construction below yields an
exact identity for the chosen A6/C23 cycle, using refined
boundary restrictions of actual labelled horizontal
special-cycle closures on the specified formal charts. Every proper pushforward used in
the identity occurs on the stated boundary abelian schemes
or on the smooth proper ball model. It does not assert an
unprovided proper projection of an entire toroidal chart.

\paragraph{The retained integral hypotheses.}
Let $E$ be the CM field of 032, $F=E^+$ and
$d=[F:\Q]$. Let $\mathcal X/R$ be the chosen smooth
proper integral model of the compact ball quotient,
with its universal PEL abelian scheme
$\mathcal A_{\rm ref}$ of relative dimension $3d$.
For this rank-three 032 ball, write
$n=\dim(X)=2$, so $s=2n=4$ in the formulas below.
Assume that the chosen prime and level permit this
model, as required in C23. Let the rank-one CM
factors $B_i$ of A6, $1\leq i\leq s=2n$, have
abelian scheme models $\mathcal B_i/R$, full
$\mathcal O_E$ action and compatible polarizations
$\lambda_i$ of degree prime to $p$. The prime is
unramified in $E$, and the reference PEL lattice is
self-dual at $p$. All finitely many maps
$f_i:X\to B_i$, polarizations, endomorphisms and
levels are defined after a fixed finite extension.
We retain the coefficients and denominators of
C23's finite CM embedding projectors.

One may pass to a common finite \'etale principal
level cover $u:\mathcal X'\to\mathcal X$ and a
finite base extension to make all objects and cusp
labels honest schemes. Write its degree as
$\delta$. The degree is recorded and is not
assumed prime to $p$: a level prime to $p$ gives
an \'etale cover, but its degree may still be
divisible by $p$. Below we work on this cover;
if none is needed, $\delta=1$.

The maps extend at the specified prime by the
N\'eron property, as proved in C23:
\[
        \mathfrak f_i:\mathcal X'\longrightarrow\mathcal B_i.
                                                        \tag{C34a}
\]
Thus none of the constructions below replaces
an actual Albanese map by an unspecified
cohomological map.

\paragraph{The actual PEL hosts and the original ball.}
For one factor $B_i$, let $P_i$ be the rank-one
Hermitian $E$-space underlying its polarized
CM first homology. Use the good integral
lattice determined by $\mathcal B_i$ and its
polarization. Add a hyperbolic $E$-plane
$H_i$ and form
\[
                          W_i=H_i\perp P_i^3.          \tag{C34b}
\]
Lan's explicit construction in Section~3A of
\emph{Toroidal compactifications of PEL-type
Kuga families}, pp.~27--30, is the integral
version of this operation: its lattice is
$\mathcal Q_{-2}\oplus L\oplus\mathcal Q_0$,
with the perfect trace pairing between the
outer summands. Take $\mathcal Q=\mathcal O_E$.
The different and polarization degrees are
units at the specified good prime, so the
new lattice has the required good-prime
pairing. Compatible neat prime-to-$p$ levels
satisfying both conditions on p.~29 give
an actual toroidal boundary stratum as in
C31, not an arbitrary extension parameter
space asserted to be a stratum.

At a real place where $P_i$ is positive,
$W_i$ has signature $(4,1)$; where it is
negative, the signature is $(1,4)$. In
either case the quotient by the chosen
hyperbolic cusp is $P_i^3$, which is
compact at every real place. Its CM
boundary point therefore has abelian
scheme $\mathcal B_i^3$.
The boundary abelian parameter scheme
is prime-to-$p$ isogenous to
$\mathcal B_i^3$ as in C31. This treats
every CM type separately and does not
silently replace switched types by the
reference type.

More explicitly, use C31's actual
isogeny $F_i:C_{\Phi_i}\to
\Hom_{\mathcal O_E}(X_i,(\mathcal B_i^3)^\vee)^\circ$,
choose an injection $X_i\hookrightarrow\mathcal O_E$
of index prime to $p$, and let
$\varphi_i:\mathcal B_i\to
\Hom_{\mathcal O_E}(X_i,\mathcal B_i^\vee)^\circ$
be the map obtained from $\lambda_i$.
Choose $m_i$ prime to $p$ annihilating
$\ker(F_i)$. The actual homomorphism
$[m_i]F_i^{-1}\varphi_i^3$ supplies
the chart on $\mathcal B_i^3$; its
extension parameter is
$[m_i]\varphi_i^3(P,Q,T)$.
Write $\Delta_i$ for its prime-to-$p$
isogeny degree. The selected zero
kernel label gives exactly the graph
in (C34h), with multiplicity one;
the unlabelled inverse image has
the other finite \'etale translates
and is not substituted for it.

Let $\mathscr M_i$ be this actual unitary
PEL model, and choose the toroidal
compactification supplied by Lan's
Theorem~1.41, pp.~16--17. The product
\[
             \mathscr P=\mathcal X'\times_R
                             \prod_{i=1}^{s}\mathscr M_i
                                                        \tag{C34c}
\]
is the product of the actual PEL
models, and is the host used here.
With compatible markings, a common
similitude factor and the chosen
components, its direct-sum moduli
realization uses the product
endomorphism algebra. The required
algebraic map is also given directly
on tuples by their direct-sum
polarizations. Its universal
abelian scheme is the direct sum of
$\mathcal A_{\rm ref}$ and the small
unitary families. The indicated boundary
stratum, after the recorded prime-to-$p$
isogeny charts, has parameter space
\[
             \mathcal X'\times_R
                             \prod_{i=1}^{s}\mathcal B_i^3.
                                                        \tag{C34d}
\]
In particular the original ball and its
universal family are present in the
actual moduli construction, not merely
in a matching of Hodge structures.

If desired the interior product moduli
map is the direct-sum PEL map into the
single Hermitian host
\[
             V\perp\bigoplus_{i=1}^{s}(H_i\perp P_i^3),
                                                        \tag{C34e}
\]
of $E$-rank $3+5s$. For one $B_i$ this
is the rank-eight host proposed above.
At the active ball place, where the
selected type has $v_i(\mu)=1$, its
signature is $(6,2)$ and its cusp
quotient has signature $(5,1)$. At an
inactive place the two possibilities
are respectively
\[
 \begin{array}{c|c|c}
 P_i^3&V\perp H_i\perp P_i^3&V\perp P_i^3\\ \hline
 (3,0)&(7,1)&(6,0)\\
 (0,3)&(4,4)&(3,3).
 \end{array}                                           \tag{C34f}
\]
Thus the larger boundary Shimura base
can be positive dimensional at switch
places. The direct-sum map fixing the
CM block still exists. We do not need
to identify an ambient refined pullback
with a smaller cycle across that
additional codimension; the intrinsic
product host (C34c) keeps all its
special-cycle conditions exact.

\paragraph{Why the relative degeneration actually splits.}
There is also a direct description inside
the larger boundary datum. Write its
abelian part as
$\mathcal A_{\rm ref}\times\prod_i\mathcal B_i^3$.
Set the reference-family components of
both $c$ and $c^\vee$ equal to zero,
and use the block-diagonal extension
and period data of the small hosts.
The polarization relation
$c\phi=\lambda_Ac^\vee$ is satisfied.
If one first sets only $c_{\rm ref}=0$,
then $c^\vee_{\rm ref}$ lies in a
prime-to-$p$ polarization kernel;
selecting its zero label is necessary.
It is supplied by the same finite
\'etale label refinement. It is not
inferred by inverting a polarization
integrally.

The Raynaud extension now splits
canonically as the product of
$\mathcal A_{\rm ref}$ and the small
Raynaud extensions. Indeed its extension
class is zero on the reference factor,
and a splitting there is unique because
$\Hom(\mathcal A_{\rm ref},T)=0$.
The projected periods in that factor
are zero by the imposed
$c^\vee_{\rm ref}=0$. Cross-period
trivializations are the rigidified
trivial ones in the chosen direct-sum
datum. Hence the full degeneration
datum is
\[
 [0\longrightarrow\mathcal A_{\rm ref}]
       \ \oplus\ \bigoplus_i[Y_i\longrightarrow G_i^{\natural}].
                                                        \tag{C34g}
\]
Mumford's construction is compatible
with products and arbitrary
homomorphisms, so (C34g) gives the
actual product family in (C34c).
This compatibility is explicit in
Lan's corrected revision,
Lemma~4.5.4.13, Construction~4.5.4.14,
Theorem~4.4.16 and
Corollary~4.5.5.5. There is consequently
no unproved lattice-period condition
in the relative construction over the
ball.

\paragraph{The labelled cycles on the relative chart.}
For a fixed $i$ and $a\in\mathcal O_E$,
C31--C33 produce three labelled
horizontal interior closures whose
refined boundary restrictions, on the
$\mathcal B_i^3$ isogeny chart, are
\[
 \begin{aligned}
 \Gamma_{i,+}(a)&=\{T=[a]P+Q\},\\
 \Gamma_{i,1}(a)&=\{T=[a]P\},\\
 \Gamma_{i,2}&=\{T=Q\}.
 \end{aligned}                                         \tag{C34h}
\]
Their norms are $a\bar a+2$, $a\bar a+1$
and $2$. On (C34c) take their product
with $\mathcal X'$ and the other
factors. The homomorphisms have zero
reference-family component, so their
lifting equations and norms are
unchanged. The $q$ parameters remain
free, the closures are smooth over
this relative boundary base, and
refined boundary restriction has no
excess term. All finite isogeny
indices in C31 are retained.

For one $i$, let $Q_i$ be the fixed
rigidified polarization bundle on
$\mathcal B_i$, and let
$\pi_{i,12}:\mathcal B_i^3\to\mathcal B_i^2$
be the proper projection. If
$\mathcal P_{i,a}$ is the exact
normalized Poincar\'e kernel of C23,
then the already proved integral
identity is
\[
 c_1(\mathcal P_{i,a})=
 \pi_{i,12*}\left(
  ([\Gamma_{i,+}(a)]-[\Gamma_{i,1}(a)]-[\Gamma_{i,2}])
                         c_1(\operatorname{pr}_T^*Q_i)\right).
                                                        \tag{C34i}
\]
Each graph on the right is now the
specified refined boundary restriction
of an actual labelled interior
horizontal closure. Taking product
with $\mathcal X'$ commutes with this
formula. Notice that $\pi_{i,12}$ is
proper on the displayed abelian
parameter space; no formal toric
projection is used.

\paragraph{The exact pullback along the ball maps.}
We use the normalized representatives
from C23, so write
\[
 D_i=\sum_j c_{ij}c_1(\mathcal P_{i,a_{ij}}),
                                                        \tag{C34j}
\]
including the stated algebraic
rescaling for the selected pure
eigencomponent. More precisely these
are the $\lambda_i[a]$ embedding-projector
combinations fixed in C23 and C31.
If a normalized representative was
originally specified through its
cohomology class, the same identity
holds in $\operatorname{Pic}\otimes M$
and hence in $\CH^1\otimes M$:
normalization on both axes kills
$\operatorname{Pic}^0(B_i\times B_i)$,
and the cycle-class map on normalized
biextensions is injective. Equivalently,
such a biextension is a homomorphism
$B_i\to B_i^\vee$, and a homomorphism
acting trivially on $H^1$ is zero.
For the latter assertion, analytic
uniformization identifies the
homomorphism with a complex-linear
map taking the period lattice to the
target lattice; zero on the lattice
forces the linear map to be zero.
Tensoring this injection with the
coefficient field preserves it.
Uniqueness of line-bundle extension
in C23 then gives the same exact
identity over the chosen DVR.
Let
$\mathcal A'=\prod_i\mathcal B_i$,
with the copy actions as fixed in
C23. On
$\mathcal X'\times_R\mathcal A'$
define
\[
 b_i(x,(Q_k)_k)=(\mathfrak f_i(x),Q_i)
                         \in\mathcal B_i^2.
                                                        \tag{C34k}
\]
Pull the \emph{line bundle} in (C34i)
back by $b_i$. This is an actual
line-bundle pullback with its
rigidifications, not just a selected
cohomological component. Its first
Chern operation is consequently
exact on every cycle. If the graph
presentation is used instead, take
its refined pullback; since every
$\Gamma$ is the graph of a map to
$\mathcal B_i$, this pullback is the
corresponding graph over
$\mathcal X'\times\mathcal A'$.
It is a regular immersion, without
an assumed transversality of
$\mathfrak f_i$. The proper graph
projection and the first Chern
operation give the same pulled-back
line-bundle class.

With $\pi:\mathcal X'\times_R\mathcal A'
\to\mathcal A'$ proper, the expression
\[
 \mathcal C'=\pi_*\left(
       \prod_{i=1}^{2n}
         \sum_jc_{ij}c_1(b_i^*\mathcal P_{i,a_{ij}})
       \cap[\mathcal X'\times_R\mathcal A']\right)
                                                        \tag{C34l}
\]
is exactly C23's chosen integral
A6 cycle on the level cover. Formula
(C34i), with the specified boundary
restrictions, expresses every one of
its factors through the actual
horizontal special-cycle closures.
There is no arbitrary Hodge projector
in this step: the finite embedding
projector and its coefficients are
precisely those already chosen in
(C34j).

Finally the finite \'etale degree
formula gives
\[
                          \mathcal C'=\delta\mathcal C,
                                                        \tag{C34m}
\]
where $\mathcal C$ is C23's output
on the original model. Indeed all
kernels are pulled back along
$u\times1$, and
$(u\times1)_*[\mathcal X'\times\mathcal A']
=\delta[\mathcal X\times\mathcal A']$;
the equality is the projection
formula. Thus division by the
recorded $\delta$ recovers exactly
the chosen original cycle. Any
$p$-valuation of $\delta$, and of
its CM projector coefficients,
remains in a quantitative integral
comparison. The same identity holds
after refined specialization, by
C23.

\paragraph{What this comparison establishes.}
For the specified smooth proper
model and normalized representatives,
the full diagram from the original
ball maps to the boundary parameter
spaces, the labelled interior
horizontal closures, their refined
boundary restrictions, and the
proper Chern-class contraction is
now given. It recovers the actual
chosen A6 cycle, not merely a cycle
in the same Hodge summand. The
generic and special-fibre
cohomological nonvanishing therefore
have the exact normalization already
proved in A6 and C23.

This is a boundary correspondence
using interior horizontal closures.
It is not an assertion that a
compact interior image meets a
boundary, or a formula for vertical
local Rapoport--Zink intersections.
Nor is it an interior Chow
pushforward along an unspecified
proper map of toroidal charts.
Those distinctions are compatible
with, rather than restrictions on,
the exact cycle identity (C34m).

\framework{Lan's Section~3A supplies
actual Kuga boundary strata and
integral lattice choices; its
Theorem~1.41 supplies the actual
Mumford formal family. The
Poincar\'e universal property,
line-bundle lifting and refined
intersection operations are those
proved and sourced in C23 and C31.
C33 supplies the exact horizontal
interior closure comparison. Their
composition here uses the given
Albanese maps, which is the
realization/correspondence operation
anchored by cote~44, scans~7--8,
and the explicit cycle/intersection
framework of cote~16,
scans~43--46.}


\subsection{THETA3-C35: exact intersections and local normalization for the chosen horizontal boundary cycles}\label{C35}

The three labelled closures used in C33--C34 have an exact
intersection calculation. It concerns those same cycles on their
specified good-prime formal charts. Their refined boundary
intersection is a finite flat CM torsion scheme, so its ordinary
length and derived intersection agree. Its Hermitian Gram matrix
also determines an explicit coherent local-density polynomial.
This supplies a quantitative datum for the chosen-cycle comparison;
it is separate from a full-rank vertical Rapoport--Zink intersection.

\paragraph{The actual labels and their Gram matrix.}
Keep C31's rank-one CM abelian scheme $B/R$, of relative dimension
$d$, its $E$-compatible polarization $\lambda$, and every
prime-to-$p$ isogeny chart and zero label used in C33. Choose
$0\ne a\in\mathcal O_E$, and put $A=a\bar a\in F$.
The three special homomorphisms have abelian-part columns
\[
 h_+=(\bar a,1,-1),\qquad
 h_1=(\bar a,0,-1),\qquad
 h_2=(0,1,-1):B\longrightarrow B^3.
 \tag{C35a}
\]
Their Hermitian pairing is $\langle h,k\rangle=
\lambda^{-1}h^\vee\lambda_{B^3}k$, with the product
polarization and Rosati conjugation fixed in C31.
In the displayed order its matrix is
\[
 T(a)=
 \begin{pmatrix}
 A+2&A+1&2\\
 A+1&A+1&1\\
 2&1&2
 \end{pmatrix},\qquad \det T(a)=A.
 \tag{C35b}
\]
Thus the individual norms are exactly the norms in C33,
and the determinant has no omitted polarization factor.
The integral change of generators
\[
 k_1=h_+-h_2=(\bar a,0,0),\quad
 k_2=h_+-h_1=(0,1,0),\quad
 k_3=h_1-h_++h_2=(0,0,-1)
 \tag{C35c}
\]
has determinant $\pm1$. Consequently $T(a)$ is integrally
Hermitian-congruent to $\operatorname{diag}(A,1,1)$.
In particular it is positive definite at every real place.
These assertions follow by the displayed multiplication and
change of basis, without changing the chosen cycle labels.

\paragraph{The exact derived intersection on the boundary.}
Write $(P,Q,T)$ for the coordinates of $B^3$ and put
\[
 \Gamma_+=\{T=[a]P+Q\},\quad
 \Gamma_1=\{T=[a]P\},\quad
 \Gamma_2=\{T=Q\}.
\]
These are precisely the refined boundary restrictions of the
three labelled horizontal interior closures in C33. Then
\[
 \mathcal O_{\Gamma_+}
       \mathbin{\otimes}^{\mathbf L}_{\mathcal O_{B^3}}
 \mathcal O_{\Gamma_1}
       \mathbin{\otimes}^{\mathbf L}_{\mathcal O_{B^3}}
 \mathcal O_{\Gamma_2}
 \simeq \mathcal O_{\{(P,0,0):[a]P=0\}},
 \tag{C35d}
\]
where the right side is in degree zero. This intersection is
finite flat over $R$, of rank
\[
       \deg[a]=N_{E/\Q}(a)=N_{F/\Q}(A).
 \tag{C35e}
\]
All the norms in (C35e) are positive. In particular every
geometric fibre has this total length.

\textbf{Proof.}
Consider the homomorphism of abelian schemes
\[
 \Phi:B^3\longrightarrow B^3,\qquad
 (P,Q,T)\longmapsto
       (T-[a]P-Q,\ T-[a]P,\ T-Q).
 \tag{C35f}
\]
If its three output coordinates are $u,v,w$, then
$Q=v-u$, $T=w+v-u$, and $[a]P=w-u$.
These identities show, by integral row and column operations,
that $\Phi$ is equivalent to $[a]\times1\times1$.
Its scheme-theoretic zero fibre is therefore the scheme on
the right of (C35d).

Here $[a]$ is finite flat, including when $p$ divides $a$.
For completeness choose $b\in\mathcal O_E$ and an integer
$N>0$ with $ab=N$. Its fibrewise kernel is contained in
$B[N]$, so the proper morphism $[a]$ is quasi-finite and
hence finite. It is surjective on every geometric fibre,
because its composite with $[b]$ is the surjective
multiplication $[N]$. Both source and target are smooth
over the discrete valuation ring. Their local rings are
regular, and the dimension equality for this finite
surjective morphism holds. Miracle flatness therefore
makes $[a]$ flat. The exact commutative-algebra statement
is Stacks Project, Lemma~10.128.1, Tag~00R4, with its
regular-target, Cohen--Macaulay-source and dimension
hypotheses. Multiplication by a nonzero integer on a
geometric abelian fibre is finite locally free by
Stacks Project, Lemma~39.9.8, Tag~0BFG.

The identity section of the smooth scheme $B^3/R$ is a
regular immersion of codimension $3d$. Its local equations
are a regular sequence. Pullback by the finite flat map
$\Phi$ preserves that regular sequence. Grouping its
$3d$ equations into the three blocks of $d$ equations
for $\Gamma_+,\Gamma_1,\Gamma_2$ gives the tensor product
of their Koszul resolutions. This tensor product has no
positive Koszul homology and resolves the ordinary zero
fibre. It proves (C35d), including the absence of all
higher Tor terms.

Finally the generic CM homology has rank one over $E$,
so the determinant of $[a]$ on rational first homology is
$N_{E/\Q}(a)$. The integral lattice index is the degree
of this isogeny. This is also Milne, \emph{Abelian
Varieties}, Lemma~10.8 and its following degree formula,
printed p.~46, or Proposition~10.23, pp.~52--53.
Finite flatness keeps this rank unchanged on the special
fibre. Since $E$ is CM, $N_{E/\Q}(a)=N_{F/\Q}(a\bar a)>0$.
This proves (C35e).\hfill$\square$

\paragraph{Refined intersection of the actual interior closures.}
Let $\mathcal Z_+,\mathcal Z_1,\mathcal Z_2$ denote
C33's labelled horizontal closures on the specified
formal chart. Their ideals are exactly the pullbacks of
the graph ideals above, with the toroidal parameters
$q_1,\ldots,q_r$ free. Locally a smooth cone chart is a
formal power-series extension of its boundary base; a
possibly twisted torus bundle does not affect this local
assertion. Flat base change of (C35d) proves that the
derived intersection of these very closures is concentrated
in degree zero. Its boundary restriction is $B[a]$.
The $q_i$ remain a regular sequence on the intersection,
because its boundary coordinate algebra is finite flat
over $R$. Thus the refined boundary Gysin pullback has
no additional Tor or excess factor:
\[
 i^!\bigl([\mathcal Z_+]\cdot[\mathcal Z_1]
                         \cdot[\mathcal Z_2]\bigr)
       =[B[a]]\quad\hbox{on the indicated boundary chart}.
 \tag{C35g}
\]
This is a local formal-chart intersection identity; no
unprovided proper morphism from the entire toroidal
chart is used.

The same calculation holds in C34 after taking the product
with $\mathcal X'$ and every other spectator factor and
imposing its zero reference-extension and zero projected-period
conditions. Hence it computes the refined intersection of
the actual closures used before C34's proper polarization
contraction. Every chart isogeny degree and the source-cover
degree $\delta$ remain as recorded in C34. There is no
extra degree in (C35g): it concerns the specified isogeny
pullback and its zero-labelled component. Forgetting those
labels includes translated components, which cannot be
replaced by the single scheme $B[a]$.

\paragraph{Exact local length at the inert CM primes of C21.}
Assume now that $p$ splits completely in $F$, is inert in
$E$ at each prime of $F$ above it, and the full order action
and good reduction satisfy C21. After the finite extension
used there, the formal group is the product of $d$ strict
one-dimensional height-two CM formal modules. Write
$r_\nu=\operatorname{ord}_p(a_\nu)$ in
$\mathcal O_E\otimes\Z_p=\prod_{\nu=1}^d
\mathcal O_{\Q_{p^2}}$.
Then the completed intersection at the identity in the
geometric special fibre has length
\[
                 p^{\,2\sum_{\nu=1}^d r_\nu}.
 \tag{C35h}
\]
Indeed a unit in each local order acts as a formal
coordinate automorphism, and $[p^{r_\nu}]$ on its
height-two factor has finite flat degree $p^{2r_\nu}$.
That factor's special fibre is connected and is supported
at the origin, as proved in C21. Tensoring their finite
local coordinate algebras multiplies their lengths.
The part of $B[a]$ of order prime to $p$ is finite
\'etale and contributes one reduced identity point.
For $a=p^r$, the whole special fibre is therefore local,
and its length is exactly $p^{2dr}$.

\paragraph{The matching coherent density and the precise remaining distinction.}
At one of these inert primes put $q=p$ and let $L_r$ be
the rank-three Hermitian lattice of (C35b) for $a=p^r$.
By (C35c) it has Gram matrix
$\operatorname{diag}(p^{2r},1,1)$, and hence
fundamental invariants $(0,0,2r)$. In Li--Zhang's
normalization, retained in C06--C07, its local Siegel
series is exactly
\[
 \operatorname{Den}(X,L_r)
   =(1-X)\sum_{j=0}^{r-1}X^{2j}+X^{2r}
   =\sum_{j=0}^{2r}(-X)^j.
 \tag{C35i}
\]
For $r=0$ the first sum is empty and the polynomial is $1$.
To prove the formula, split off the self-dual rank-two
summand. Every integral overlattice splits off this same
summand by its perfect pairing. The remaining rank-one
overlattices are $p^{-j}\mathcal O x$, $0\le j\le r$.
For $j<r$ their type is one and their weight factor is
$1-X$; for $j=r$ the lattice is self-dual and its weight
factor is $1$. Their index exponents are $2j$.
Li--Zhang, Theorem~3.5.1, printed p.~18, gives (C35i).
Its definition in Section~3.2, p.~16, applies to this
coherent, even-valuation lattice as well as to odd valuation.

Consequently
\[
 \operatorname{Den}(1,L_r)=1,\qquad
 -\operatorname{Den}'(1,L_r)=-r,\qquad
 \operatorname{Den}(-q,L_r)=\sum_{j=0}^{2r}q^j.
 \tag{C35j}
\]
For $r\ge1$ the exact one-factor length of our horizontal
intersection can thus also be written
\[
 q^{2r}=
 \frac{q}{q+1}
 \left(\operatorname{Den}(-q,L_r)
               -\operatorname{Den}(-q,L_{r-1})\right).
 \tag{C35k}
\]
This is an equality for the polynomial belonging to the
actual Gram matrix, with its factor $q/(q+1)$ specified.
Here $X=-q$ is a formal value of the polynomial; it is
not one of the standard interpolation points
$X=(-q)^{-k}$ with $k\geq0$ and nonnegative additional
representation rank. Equation (C35k) is an algebraic
identity between the displayed polynomial and the proved
length. It is not an arithmetic Siegel--Weil theorem or
an application of a derivative intersection formula.

Indeed $L_r$ has even determinant valuation and lies in
the split local Hermitian space. Li--Zhang's
Theorem~1.2.1, equivalently Theorem~3.4.1, concerns a
full-rank lattice in the \emph{nonsplit} special
quasi-homomorphism space, which has odd determinant
valuation; see Section~1.7, printed p.~8.
The source defines its derived local density in the
odd-valuation case in Definition~3.2.1, p.~16.
Thus the formal number $-\operatorname{Den}'(1,L_r)$
in (C35j) is not the derived Rapoport--Zink intersection
of the cycles considered here. The positive length
(C35h), the coherent density (C35i), the specified
Hermitian norms, and the refined horizontal comparison
(C35g) are now all fixed. An identification with a
vertical Rapoport--Zink arithmetic degree still requires
its own change of local model and special-cycle data;
C35 does not assert that identification.


\part{Unresolved comparisons and current gates}
\section{Statements not asserted as proved}

There is no unchecked theorem silently used in the preceding part.
The implications with hypotheses, including B8 and C15's universality
reduction, are proved implications; their hypotheses are not supplied
by this run. The following are the remaining identifications.

\paragraph{G1: the exact iterated character-independent functional.}
The common oscillator kernel, A1's normalized zero orbit, B11's
restriction seesaw, A5's finite-frame identity and A7's actual single
theta lift have been established. B17 now makes the actual wedge
a literal double lift $V\to P\to V$ by the genuine adjoint
kernel and a common positive inverse $Q^{-1}$.
It specializes the analytic last-step Bessel seesaw without a
new occurrence hypothesis. The operator is common to every
character for fixed spaces, kernel and levels.
The remaining \texttt{MainX} question is its stronger arithmetic
uniform relation-preservation statement. To apply it, one must
construct its actual common geometry,
tuples, archimedean maps, definite class kernels and arithmetic links,
and identify its normalized functional on the selected 032 vectors.
The concrete all-place tower fails the canonical node-gap and
signature requirements. The abstract two-step case is more flexible:
with global Witt index zero, one can take $m=2$ and the Bessel
unipotent rank zero. Its middle space must be totally definite.
A7's natural multiplicity space is totally definite precisely when
the selected first-list CM types agree; that direct realization then
gives a divisor-generated balanced line. B16 supplies a conditional
pole criterion for a different coherent definite source.
The arithmetic kernel/covolume and terminal-mass equations,
and an arithmetic interpretation of the analytic correction,
are the remaining matching problems. B17's upper space is
$A\perp\mathbb H$, of signature $(3,1)$ at every real place;
it is not a hidden definite inner form.

B19 narrows the arithmetic comparison to a finite CM-jet
functional on the actual polarized tori at a rational negative
line. Finite Heisenberg equivariance gives one theta-null
comparison scalar for every fixed-type finite input family;
the tangent periods, levels, stabilizers and Haar volumes are
explicit. The exact archive class-kernel identity can now be
tested on a finite basis after its polarized datum and frames
are matched. B20 supplies the literal isotropic doubled columns
for the exponent-one value using even integral reference sections.
It does not identify the resulting global automorphic functions
or provide high-order integral constant jets. Across different
balanced CM types, the theta-null determinant-period factors
still require a proof; balance of tangent lines alone is
insufficient to identify the full Hodge determinant.
A10 now constructs the common stabilized product theta module,
including the exact reference multipliers and a correspondence
proof of the balanced tangent-period relation. B21 identifies
the raw scalar obstruction as the specific ratio
$D_1D_2/(D_3D_4)$ of full Hodge period determinants. It proves
by an explicit complete-type calculation that this determinant
tensor need not balance. No transcendence conclusion follows;
the raw period relation, or its retained correction, and the
global weight/signature comparison remain explicit.

A11 settles the actual generic geometric host problem. Its
degree-at-most-four field extension is finite-unramified over
the real subfield, and its minimal CM-type padding, compatible
principalization and finite theta interpolation produce the
literal \path{Geo.Tor.Host/Datum/Obj}. The transported root
in \texttt{Obj.geom} accounts for every coefficient embedding.
The dimension cost is at most $64$ times the initial dimension.
This geometric host is different from the stronger global
\nolinkurl{Formalism.BasicNotions.Host} used by \texttt{D7};
its additional Galois frame, arithmetic parameters and the
prescribed \texttt{StdPair} and tuple still need construction.
A16 proves that field-preserving enlargement cannot construct this
stronger host for arbitrary CM fields: its field is necessarily
abelian Galois, whereas the displayed nonnormal quartic CM field
and its nonabelian Galois normal closure cannot embed into it.
Thus the universal literal-host enlargement route is refuted;
the general comparison must use the field-independent proved
identities with their own arithmetic hypotheses.
The interpolation preserves algebraic first jets and records
the analytic comparison scalars, rather than identifying the
entire global theta family.

B23 gives an algebraic interpretation of A5's correction on
each full CM Hodge block. The non-Tate cup image of the actual
finite CM sources is an algebraic retract, constructed using
CM Hodge only on the source product. In a full multiplicity
and embedding block, the inverse is an algebraic correspondence
times $\nu_G|P_\chi|^2/|t_x|^2$. The common canonical-fibre
period $t_x$ is the specified Kodaira--Spencer determinant
divided by $(2\pi i)^2$. This removes the matrix-algebraicity
gate in that setting. The scalar and the raw theta-coordinate
comparisons remain explicit; the inverse of one fixed B17
integral kernel has not been identified with this operator.

A12 goes beyond abstract interpolation on the standard side.
For a specified compatible \texttt{StdPair}, finite
Heisenberg congruence projectors have actual
\path{IsThetaAt/IsRatDatum} witnesses. An explicit Gaussian
tail certificate makes their first-derivative matrix invertible;
the algebraic inverse retains the vacuum scalar, theta scalar,
common tangent period and isogeny matrices. When the prescribed
global \texttt{xiAt} has the stated support and period bounds,
its equality with the newly constructed array is exactly an
equality on one finite quotient. This does not supply the
stronger host, change the original tuple by definition, or
identify the global theta functions. The full jet value at
the origin is not prescribed by the derivative spanning result.

B24 identifies the actual B17 integral-kernel inverse in the
natural totally definite same-type sector. The source infinity
type is a character of the full compact group; its integral
reduces to finitely many rational masses. A common B19 CM
evaluation factor makes the kernel matrix projectively
algebraic and its image algebraically defined in the full CM
block. The actual inverse is consequently an explicit positive
scalar times an algebraic correspondence. A7 shows that this
natural sector is divisor-generated. For the general
exceptional balanced type the remaining operator condition is
projective algebraicity of the specified noncompact-source
Gram matrix, or a different coherent definite realization.

B25 removes the image-preservation problem in the finite-period
identity even for the actual noncompact middle source. Cayley
weak approximation and rational theta automorphy show that
finite adelic source translates span the same finite theta
image. Their common CM evaluation scalar makes that image
algebraic in coherent coordinates. The algebraic finite-frame
inverse can therefore preserve it. For every pair in this
image, including 032's actual fixed-input compact-place
variation, one has the exact identity
\[
 \langle f,g\rangle_G
   =s_\chi\sum_{j,\alpha}
       \overline{\mathcal P_{j,\alpha}(K_0f)}
                     \mathcal P_{j,\alpha}(g),
\]
with one algebraic correspondence $K_0$ and one named positive
scalar. Each corrected vector remains in the original B17
theta image, so the literal last-step seesaw applies there.
This gives the fully proved finite definite-period form of
the proposed comparison. The actual integral inverse $TT^*$
still has its own Gram-period matrix; the displayed finite
restriction inverse does not replace it, and the finite sum
is not asserted to be one prescribed arithmetic MainX tuple.

A13 now realizes any stipulated admissible finite tame character
family in a literal valid \texttt{TC} on an existing full
\texttt{Pre}. The coefficient enlargement preserves the entire
\texttt{Base}; the source constructor retains the chosen ambient
datum $J$, and the proved universal laws supply \texttt{Laws}
and \texttt{Law0}. With the unchanged \texttt{HasLetters},
\texttt{Later} and a same-$J$ arithmetic-jet field follow.
A14 supplies that specialized $p$-primary \texttt{TwDat} shape
for any prescribed finite local primary family at the stated
allowed split places. Starting from an existing full witness,
it reconstructs \texttt{ColW} while preserving the literal
\texttt{Pins}, \texttt{KsExact}, \texttt{FieldOK} and
\texttt{ValLL0}. The required coefficient-prime orbits are
created by a disjoint real multiquadratic extension split at $p$.
The actual arrays vanish at zero and have full prescribed rank
on their support. Their local incidence and determinant-character
spans, and the exact root-of-unity obstruction to global tame
characters, are proved. The remaining finite gate is the actual
Heisenberg/Fourier transform of these arrays followed by the
period-normalized first-derivative matrix. Raw array-window
selection alone is not the Gaussian-window construction of A12.
An arbitrary 032 normalization has not been asserted to satisfy
the literal \texttt{AmbIn.Ok} conditions.

B26 further reduces the actual noncompact-source integral Gram.
Its coefficient sections lie in the algebraic weight-three
source line, with the central CM determinant factor retained,
and their coherent coordinates have the common scalar
$\gamma=C_x/(t_x t_y)$. The Gram is an algebraic compression
of $\bigoplus_\lambda H_\lambda\otimes G_\lambda$, where
$G_\lambda$ is algebraic and $H_\lambda$ is the actual positive
multiplicity-period matrix. A single irreducible Hecke sector
therefore closes the actual $TT^*$ inverse up to one positive
scalar. Distinct sector norm ratios and higher multiplicity
matrices remain. Its cohomological source is
$H^s(Y,\bigotimes_\tau\mathbb V_\tau)$ with the recorded
constant line $\mathscr D=\mathscr C\otimes\bigotimes_\tau
\det\mathbb V_\tau$; neither the determinant comparison periods
nor the exact $4^{-s}$ measure factor is discarded.

B27 proves that the actual rank-two mixed source has
multiplicity-free finite Hecke sectors. The only nongeneric
rank-two parameter would be a determinant character at a
split finite place; strong approximation would then force
global derived-group invariance, contradicting the actual
weight-three discrete series. The proved generic pure-inner
multiplicity formula applies. Thus every $H_\lambda$ above
is a positive scalar $c_\lambda$. Each source-sector kernel
has a projectively algebraic inverse on its actual target
image, while the full kernel still has an algebraic
compression of these different scalar periods. B29 now proves
that their actual target images are orthogonal, by finite local
Howe uniqueness and the already fixed source infinity type.
The full inverse is therefore the exact sum
$(\nu_G|P_\chi|^2|\gamma|^2)^{-1}
\sum_\lambda c_\lambda^{-1}B_\lambda$ of algebraic sector
inverses with retained reciprocal scalar weights. Their exact
ratios remain. The operator also changes by exactly
$|c_t|^{-2}$ along the original compact-place variation.
B28's bound controls its norm without evaluating the source
periods.

C37 supplies the simultaneous analytic rank theorem for several
ordered rank-one differences at the same prime. It retains arbitrary
lattice representatives and all their moving character phases. A15 proves the exact ordered modulation identity and finite-block
injectivity from selfdual lattice eigencharacters and Fourier
inversion. A nonzero permitted array therefore supplies a nonzero
baseline in that block. The remaining application check is the
actual fixed-data and genuine-period-lattice transport as the
allowed archive prime grows; an arbitrary changing datum has
not been identified with a fixed positive section.

\paragraph{G2: turning the explicit certificates into a uniform
arithmetic lower bound for the prescribed iterated family.}
Q1--Q11 and B5--B7 give actual formulas under specified inputs.
The finite unipotent search yields $|J|\geq1/(4S_{\rm ext})$ with
an explicit grid size after its Gaussian envelope and lattice basis
are supplied. An arithmetic Gram bound needs a coefficient field,
denominator and bounds at all embeddings. Q9 provides resultant
denominators for specified integral Hecke matrices. Q11 transfers
the definite period-square certificate directly to the geometric
ball norm using explicit injective-chart constants; no unknown
ball Gram upper bound is needed. The remaining
step is the normalization and overlap data of the particular
iterated family in G1. Positivity alone cannot supply it, as B4
and B6 show. No universal cycle-degree estimate is asserted.
B19 now supplies the actual algebraic finite coefficient to
which the norm argument applies: its nonzero value has a lower
bound in terms of its denominator and its other conjugates,
multiplied by the specified common CM period and measure factor.
The arithmetic input bounds have not been supplied uniformly.
B22 removes the unknown absolute CM-period size by a specified
algebraic normalization: an even reference and an odd first jet
have explicit Gaussian-tail bounds, and the first-jet frame
makes the combined theta and tangent scalar exactly that
bounded analytic derivative. The resulting algebraic
denominators and conjugate bounds, and any loss of integral
saturation, must still be recorded. A9 also proves an explicit
constant-loss determinant derivative estimate under an integral
product-formal splitting, valid in every theta degree after
isotropic frame cancellation.

B28 supplies a further fixed-input certificate for the genuine
integral-kernel correction. On disjoint source polydisks,
weight-three integration gives
$\kappa_l=\nu_l(4\pi)^s\prod_j(\rho_{lj}^2-\rho_{lj}^4/2)$.
The actual source CM evaluation matrix is $(C_x/t_x)D$,
so $W\geq|C_x/t_x|^2(D^*\operatorname{diag}(\kappa_l)D)^{\mathsf t}$.
This bounds the least-norm source preimage and the actual
$TT^*$ inverse. The extra source-fibre period cancels exactly.
The remaining numerical data are the specified finite
algebraic minor and its denominators/conjugates, chart radii,
original measures and the retained target theta factor.
No uniform estimate for unspecified arithmetic data is supplied.

\paragraph{G3: an exceptional non-CM Hodge summand.}
C16 supplies infinitely many non-CM elliptic source factors, and
C20 adds a geometrically simple genus-two Jacobian with
endomorphism ring $\Z$, on a ball cover of degree at most1764.
C24 gives non-CM balanced $\Q(i)$-Weil sources in every
even dimension. C25 specifies a geometrically simple fourfold
with endomorphism ring $\Z[i]$, N\'eron--Severi rank one,
and entirely exceptional Weil space; its ball cover has degree
at most43264. Thus the source and the exceptional target
are both concrete.
C13 gives an explicit algebraic projector onto a Weil space.
C32 now supplies the nonzero projected cycle for the named
C25 target. Its integral class is pure Weil and has exact
self-intersection $73728$. Pullback by $[2+i]$ supplies the
second rational direction, so the exceptional Weil-space
extraction is complete for this target and the displayed
cyclic-four family.
The known curve-product maps in C12, C16 and C20 are subject to the
same C14 vanishing whenever the target has the generic Weil-type
divisor algebra. C32 uses a signed split-plane cycle, refined
pullback and relative hyperplane sections on the fourth power
of the actual ball cover; its image has self-intersection
$d^8\cdot73728$. The remaining task is to extend this
higher-correspondence construction to other non-CM Hodge
summands, or to identify a literal single-ball theta realization.
C15's characteristic-zero universality implication is a source
reduction, not a proof of that conjecture or of the extraction gate.
C27 makes the map-based version a finite quartic cup-product
test on the $A$-isotypic Albanese homomorphism space. Pulling
back existing zero pairings to another finite cover cannot
repair them; genuinely new maps and mixed cup coefficients
are required. A fixed generating source map also has an
injective target Kodaira--Spencer obstruction, so a deformation
argument must use a higher relative cycle instead of retaining
that degree-one quotient.
The countability result restricts fixed arithmetic towers over
$\C$; it does not preclude the dense algebraic families constructed
here or a further deformation argument.

C28 now excludes the entire single-fibration base-change
tower for the specified C25 target, including all previously
unknown maps on those covers. The actual geometric generic
fibre has genus $76$; its Jacobian splits into factors of
dimensions two and four carrying $\Q(\zeta_5)$. It has no
homomorphism to the simple fourfold with endomorphism algebra
$\Q(i)$. Consequently every map on every such base change
factors through its Stein base curve and the full C27 quartic
vanishes. A map-based solution must therefore change the
generic fibre, or be replaced by a higher correspondence.

C30 also excludes the natural simultaneous two-Fermat-base
construction. Its finite cyclic model has no additional
coherent first cohomology, and resolving its explicit Fermat
cones adds only CM Jacobian quotients. The Albanese is
isogenous to the product of the two base Jacobians and these
CM factors. Thus every map to C25 factors, up to multiplication,
through the product of the two curves; C14 again makes the
entire quartic zero. An arbitrary finite cover need not arise
this way, and remains outside this calculation.

C32 also constructs integral pure Weil cycles on every C24
cyclic-four Jacobian of dimension $2n$, with exact square
$2(-1)^n4^{2n-1}((2n)!)^2$. Its $[2+i]$ pullback spans
the other rational direction, and the actual ball power has
image multiplier $d^{2n}$. The remaining scope is other
non-CM Hodge summands and the additional N\'eron--Severi
condition needed to call every higher-dimensional member
exceptional outside the divisor algebra.

The C28 and C30 exclusions therefore remain precise statements
about maps in those towers, not barriers to the constructed
higher correspondence. C32 uses Schoen's paired-branch geometry
and checks the order-four exceptional inertia directly. It
does not import Markman's cycle. The supplied fourfold theorem
already asserts algebraicity; the new result here is the
specified cycle, its normalization and its actual ball-power
realization.

\paragraph{G4: the chosen integral correspondence and connection.}
C17 identifies the exact horizontal correspondence and local-model
data needed to compare the 032 image cycle with a Kudla--Rapoport
special cycle. C06--C08 prevent replacing those data by an
undifferentiated positivity claim. C11 proves the sharp integral
height-two derivative estimate. C19 removes its artificial
polynomial-coordinate hypothesis and proves integrality of the
CM de Rham normalization coefficient $e_2^*$.
A8 identifies the actual completed 032 tensor with the matched
canonical harmonic jet, including its embedding labels and
factorials. On a matched CM elliptic factor, A8/C19 give the
actual reduced-theta derivative bounds and their prime-to-$p$
level normalization. C21 proves the actual product decomposition
at primes split in $F$ and inert in $E$, and the strong
mixed-derivative bound for canonically matched equal-degree
ratios with integral sections and unit denominator.
For arbitrary integer tensor degree the elliptic bound has
the weaker positive slope, which C22 proves optimal through
exact infinite subsequences on two actual tame theta sections.
Thus transferring the ratio slope to all individual sections
is decided negatively in this normalization.
A9 resolves the positive-degree frame problem for the precise
archive isotropic columns: all higher canonical gauge terms
cancel to an integral unit, and the ordinary formal depth
equals the actual local theta depth. C26 proves an independent
actual elliptic two-column instance. B20 completes the first
determinant value to these columns under explicit integral
evaluation and saturation hypotheses. These conclusions do
not extend the raw one-copy bound refuted in C22.

C23 constructs the actual A6 cycle over a specified good model
and specializes its nonzero class. Its pure CM kernels are
finite combinations of normalized Poincar\'e biextensions,
the exact category of linear special homomorphisms in SGA~7~I,
VIII~3.1--3.2. C29 and C31 realize these kernels by signed
labelled homomorphism-lifting graphs on actual PEL boundary
charts, for elliptic and higher CM factors respectively.
C29 retains the exact integral multiplier $24$ or $48$;
C31's higher-CM construction instead uses the specified
prime-to-$p$ isogeny chart and zero kernel labels.

C33 proves that these graphs are exactly the labelled horizontal
interior special-cycle closures on the stated formal charts.
The good source has zero period lattice, so Lan's fully faithful
theorem for arbitrary homomorphisms identifies the extension
condition with $h^\vee c=0$. The three positive Hermitian norms,
characteristic-zero deformation functor and saturated scheme
structure are explicit. Refined boundary restriction gives the
graph before the proper polarization contraction on its
abelian parameter space.

C34 supplies the relative correspondence to the original compact
ball on the literal product of its good model and the actual
small PEL hosts. Both reference extension and dual-period
coordinates are zero, and the degeneration splits as stated.
Pulling back the exact normalized line bundles along the actual
ball maps and applying the proper Chern-class contraction gives
$\mathcal C'=\delta\mathcal C$ for the chosen original cycle.
The finite cover degree, including any $p$-divisibility, and all
CM projector coefficients remain explicit. Thus the chosen
horizontal-cycle comparison is proved at these good models.
The remaining quantitative comparison concerns the exact theta
normalizations and integral denominators. Vertical local
Rapoport--Zink cycles and their derived intersection numbers
require their own equations; they are not inferred from this
horizontal identity. No proper projection of the full toroidal
formal chart is used.

C35 fixes an exact local quantitative datum for those same
horizontal closures. Their Hermitian Gram is integrally
congruent to $\operatorname{diag}(a\bar a,1,1)$; their derived
intersection has no higher Tor, and its refined boundary
restriction is $B[a]$ with degree $N_{E/\Q}(a)$.
At the specified inert factors its identity length is
$p^{2\sum r_\nu}$. Its coherent density polynomial is also
computed with the same Gram. Its formal derivative is not
identified with a vertical arithmetic intersection.

These gates are recorded with dated changes in \texttt{LEDGER.md}.
The mandated three-hour progress review has not been replaced by
this mathematical inventory. Its timing is controlled by the actual
run clock and recorded separately.


\appendix
\section{Realization and archival crosswalk for every finding}

The following anchors locate the mathematical framework of a
finding, not its theta formula in Grothendieck's handwriting.
Write G44 for cote~44, scans~7--8 (12 August 1964), and G16 for
cote~16, scans~43--46 (4 January 1967). The original scans are
\url{https://grothendieck.umontpellier.fr/44.pdf} and
\url{https://grothendieck.umontpellier.fr/16.pdf}.
The published \emph{Standard conjectures on algebraic cycles},
pp.~196--198, is an additional polarization/realization anchor.

\begin{longtable}{@{}p{.18\linewidth}p{.48\linewidth}p{.27\linewidth}@{}}
\toprule Finding IDs & Geometric or motivic meaning & Exact framework anchor\\\midrule\endhead
A1--A3 & Oscillator period functionals on a Betti realization; extraction before an algebraic image cycle. & G44, scans 7--8.\\
A4 & Restriction to ample complete intersections; compatibility of the cup pairing. & G16, scans 43--45.\\
A5 & Positive analytic recovery operator on canonical forms; algebraicity is a separate question. & G44, scans 7--8; G16, scans 44--45.\\
A6 & Actual divisor contractions and image-cycle pushforwards, hence algebraic correspondences modulo homological equivalence. & G44, scans 7--8; G16, scans 43--46.\\
A7 & Exact analytic theta realization; its source vector is not asserted to be a motivic correspondence. & G44, scans 7--8.\\
A8 & Rigidified theta line bundle and canonical harmonic jet; actual elliptic integral normalization. & G44, scans 7--8; G16, scans 43--46.\\
A9 & Canonical biextension-frame cancellation on integral isotropic columns; local derivative-depth equality and a split-formal determinant bound. & SGA7 I, VIII 3.1--3.2, pp.254--255; G44, scans 7--8, for the further correspondence problem.\\
A10 & A common finite theta module on a product and a balanced tangent-period relation from CM algebraic correspondences. & G44, scans 7--8; G16, scans 43--46.\\
B0--B3 & Hermitian realization data and the local representation needed for a period comparison. & G44, scans 7--8, as the eventual correspondence setting.\\
B4--B8 & Positivity, determinant and overlap estimates after fixing realization norms. & G44, scans 7--8; published Standard conjectures, pp.196--198.\\
B9--B12 & Local and global definite-subgroup coefficients of the canonical line; evaluation is not pullback to degree-two cohomology of a point. & G44, scans 7--8.\\
B13 & Uniform functional on an iterated theta family, with arithmetic links retained. & G44, scans 7--8; G16, scans 44--45.\\
B15 & Explicit local theta representations, global splitting characters and exact normalized two-step kernels. & G44, scans 7--8; G16, scans 44--45.\\
B16--B18 & Conditional coherent descent and an unconditional analytic adjoint seesaw on the actual theta image. & G44, scans 7--8; G16, scans 44--45.\\
B19--B20 & Algebraic theta modules at CM points, common period lines and exact first-jet isotropic completion. & G44, scans 7--8; SGA7 I, VIII 3.1--3.2, pp.254--255, for rigidified correspondences.\\
B21 & Full Hodge-determinant period correction and the distinction between a balanced tangent tensor and a top determinant tensor. & G44, scans 7--8, for algebraic correspondence realizations; exact modular-form sources in B21.\\
B22 & Explicit reference-jet normalization of the analytic comparison factor; arithmetic denominator and integral-frame costs retained. & G44, scans 7--8; G16, scans 43--46, for the separate rational-correspondence comparison.\\
B14 & Algebraic diagonal contraction of specified tensor realizations, contrasted with a different coherent line-bundle power. & G16, scans 43--46.\\
Q1--Q3 & Quantitative selection of forms and rational multiplicity data preceding cycle construction. & G44, scans 7--8.\\
Q4--Q8 & Explicit normalized period and polarization certificates; Q4,Q5 are aliases of A2,A3. & G44, scans 7--8; published Standard conjectures, pp.196--198.\\
Q9 & Algebraic endomorphism polynomials; idempotence modulo homological equivalence, with denominators. & G16, scans 43--45.\\
Q11 & Definite finite periods control the geometric ball norm through explicit holomorphic chart constants. & G44, scans 7--8; G16, scans 43--46.\\
Q10 & All-degree determinant and period certificates for the A6 cycle construction. & G44, scans 7--8; G16, scans 43--46.\\
C01--C03 & Albanese realization, trace projector and genuine non-CM factors on noncongruence covers. & G44, scans 7--8.\\
C04--C05 & Correspondence-module extraction and its Hodge-group scope. & G44, scans 7--8; G16, scans 43--45.\\
C06--C08 & Derived local intersections and the horizontal/special-fiber distinction. & SGA4$\frac12$, Cycle, Rem.2.3.9 p.150; C17's exact specialization locators.\\
C09--C11 & Local deformation and derivative valuations; an integral frame is additional input to a cycle comparison. & EGA IV$_2$, Thm.6.9.1 p.153; G44 scans 7--8 as the eventual realization setting.\\
C12 & Image cycle with degree500; product of diagonal divisors. & G44, scans 7--8; G16, scans 43--45.\\
C13--C14 & Explicit Weil projector and its action on the Lefschetz algebra. & G16, scans 43--46; G44, scans 7--8.\\
C15--C16 & Albanese quotients from finite \'etale covers and exact image-cycle degrees. & G44, scans 7--8; G16, scans 43--46.\\
C17 & Relative algebraic correspondences and specialization. & EGA IV$_3$, 8.8.2--8.8.3 pp.28--34; SGA4 III, XVI 2.1 p.213, 2.3 p.215; SGA4$\frac12$, Cycle 2.3.2--2.3.9 pp.145--150.\\
C18 & Integral line-bundle ratios on an explicit tame isogeny cover; connection data retained. & EGA IV$_2$, 6.9.1 p.153; SGA4$\frac12$, Cycle 2.3.2 p.145, 2.3.8 p.148.\\
C19 & Integral de Rham CM eigensplitting and intrinsic height-two derivative estimates, with the actual reduced theta frame. & G44, scans 7--8; C17's exact specialization framework.\\
C21--C22 & Local CM product factors and integral canonical theta ratios; an exact sharpness example for individual sections. & G44, scans 7--8; C17 specialization framework; Tate, p-Divisible Groups, connected formal factors.\\
C20 & A simple non-CM abelian-surface quotient of a ball-cover Albanese; injective specialization of endomorphisms. & G44, scans 7--8; G16, scans 43--46; EGA IV$_3$, 11.10 as cited by Conrad, Thm.2.1(2) p.70.\\
C23 & Integral first-Chern correspondence, balanced Poincar\'e kernels and the exact horizontal special-homomorphism lifting gate. & SGA7 I, VIII 3.1--3.2, pp.254--255; C17's exact specialization locators.\\
C24--C25 & Non-CM balanced Weil sources and a simple exceptional fourfold; persistent toric rank controls endomorphisms. & SGA7 I, IX 12.1(d), pp.467--468, and 12.3.6--7, pp.471--472; G44, scans 7--8.\\
C26 & Positive-degree actual theta-frame cancellation along a completed CM-isotropic direction. & SGA7 I, VIII 3.1--3.2, pp.254--255, for biextensions; G44, scans 7--8, for the separate algebraic-correspondence step.\\
C27 & Finite cup-product test for a Weil projection and the obstruction to deforming a fixed generating Albanese quotient. & G44, scans 7--8; Stacks Lemma 93.10.3, Tag 0E3S, and Lemma 58.8.3, Tag 0BQB.\\
A11--A12 & Actual finite polarized theta hosts and rational Heisenberg inputs; derivative interpolation with all realization factors. & G44, scans 7--8; G16, scans 43--46, for the subsequent correspondence comparison.\\
A13--A15 & Literal arithmetic character families, full coefficient-field transport, normalized jets, exact finite theta transform and faithful Heisenberg action. & G44, scans 7--8; G16, scans 43--46, for the separate algebraic inverse.\\
A16 & Endomorphism-field preservation and an explicit obstruction to universal literal-host enlargement. & G44, scans 7--8; G16, scans 43--46, for the separate correspondence question.\\
B23--B25 & Algebraic retract of the CM cup image and image-preserving positive inverse correspondences; distinct integral and finite-frame operators. & G44, scans 7--8; G16, scans 43--46.\\
B26--B28 & Kuga realization, determinant periods, actual multiplicity-free Hecke sectors and an explicit source-evaluation bound for the genuine inverse. & G44, scans 7--8; G16, scans 43--46; Deligne Hodge II 2.2.1--2.2.6 pp.27--29 and 4.1.1(i) p.40.\\
B29 & Orthogonal actual target sectors and the scalar-weighted algebraic-correspondence decomposition of the original inverse. & G44, scans 7--8; G16, scans 43--46.\\
C28, C30 & Exact Albanese decompositions and curve-factorization exclusions for the named exceptional target. & G44, scans 7--8; G16, scans 43--46.\\
C29, C31 & Actual elliptic and higher-CM horizontal boundary special cycles and integral Poincar\'e contraction. & SGA7 I, VII 2.9.5 p.175, VIII 3.1--3.2 pp.254--255; G44, scans 7--8.\\
C32 & Explicit integral pure Weil cycles and their ball-power correspondences, including good-prime specialization. & G44, scans 7--8; G16, scans 43--46; C17 specialization locators.\\
C33--C34 & Scheme-theoretic horizontal interior closures and exact recovery of the chosen compact-ball cycle by a relative PEL boundary correspondence. & SGA7 I, VII 2.9.5 p.175, VIII 3.1--3.2 pp.254--255; G44, scans 7--8.\\
C35 & Exact Tor-free horizontal special-cycle intersection, finite-flat CM torsion degree and the matched coherent density polynomial. & SGA7 I, VII 2.9.5 p.175 and VIII 3.1--3.2 pp.254--255; C17 specialization; G44, scans 7--8.\\
C36 & Explicit exceptional Weil cycles and common ball-cover Albanese access for product targets in every even dimension. & G44, scans 7--8; G16, scans 43--46, for projector correspondences and contraction modulo homological equivalence.\\
C37 & Simultaneous first-jet spanning for ordered rank-one theta differences, retaining arbitrary lattice lifts and their character phases. & G44, scans 7--8, for the later realization comparison; this analytic rank result alone is not a motivic projector.\\
\bottomrule
\end{longtable}

\section{Source register and exact version conventions}

All cited supplied works remain in the package under their original
filenames. In particular, 032 means
\emph{The rational Hodge conjecture for CM abelian varieties}; its
printed theorem/page locators are used throughout. The companion
van Hoften--Sempliner paper is the supplied Hexagon 2610.00259v2.
Milne's supplied Hexagon 2611.00244v1, Theorems~2.1, 3.1 and~3.5,
is carried as the stated CM-Hodge, finite-field Tate and Hodge-standard
background. It is distinct from his 1999 article and modular-forms notes.

The supplied Markman source is \emph{Cycles on abelian $2n$-folds of
Weil type from secant sheaves on abelian $n$-folds},
\path{LIT_Markman_Hodge_abelian_fourfolds_2502.03415/250531.tex}.
Its theorem labelled \texttt{thm-algebraicity} in Section~1.5 and
corollary in Section~1.6 cover the stated Weil-sixfold and
abelian-fourfold results. These are background when determining
whether a proposed non-CM target is already known; this run does
not claim to improve those general results.

The external primary-source URLs and exact reading locators are
listed in \path{documentation/external_sources_C.tsv} and the
dated source-check records. Important additional versions are:
\begin{itemize}[leftmargin=*]
\item BMM: Bergeron--Millson--M\oe glin, published 2016 version,
125 PDF pages; Lemmas~5.5, 5.18, Propositions~5.17, 5.19,
Theorem~9.1. The 97-page preprint has different page numbers.
\item Liu: \emph{Fourier--Jacobi cycles and arithmetic relative
trace formula}, \url{https://yhzhumath.github.io/FJ.pdf};
Definitions~4.3, 4.5, 4.12; Theorems~4.15, 4.18;
Corollary~4.20; Appendix~D.
\item Haan: arXiv:1403.5061v7, Theorem~3.4, p.8,
\url{https://arxiv.org/pdf/1403.5061v7}.
\item Gan--Qiu--Takeda: arXiv:1207.4709v3,
Sections~11.3--11.7; its stated dimension range is retained.
\item Gan--Takeda: \emph{A proof of the Howe duality conjecture},
\url{https://arxiv.org/pdf/1407.1995}, Theorem~1.2,
printed pp.1--3; finite local field of characteristic different
from two, arbitrary residue characteristic, Hermitian dual pairs.
The separate quaternionic qualification is not used.
\item M\'inguez: \emph{Correspondance de Howe explicite:
paires duales de type II}, Annales scientifiques de l'ENS
41 (2008), 717--741; Theorem~5.1 and Corollary~6.3(1),
printed p.734 for the latter;
\url{https://www.numdam.org/item/10.24033/asens.2080.pdf}.
\item Gross: \emph{Incoherent definite spaces and Shimura varieties},
25 February 2018, Theorem~3.1, pp.6--7.
\item Milne: \emph{Lefschetz classes on abelian varieties},
\url{https://jmilne.org/math/articles/1999aP.pdf},
Proposition~5.4 and Corollary~5.5, p.25, used in C14.
\item Katagiri--Kobayashi: arXiv:2609.07024v1,
Definition~3.1, Remark~3.2 and Lemma~3.3, pp.9--11,
\url{https://arxiv.org/pdf/2609.07024}.
\item Howe--Zhu: arXiv:math/0002205v1,
\S2 p.3, Proposition~3 and Lemma~4 p.4,
Lemma~5 pp.4--5 and Theorem~6 pp.5--6;
\url{https://arxiv.org/pdf/math/0002205}.
\item SGA~7~I: original scan
\url{https://wstein.org/sga/sga/pdf/sga7-1.pdf};
Expos\'e VIII 3.1--3.2, printed pp.254--255;
Expos\'e IX 12.1(d), printed pp.467--468, and
12.3.6--7, printed pp.471--472.
\item Serre--Tate: \emph{Good reduction of abelian varieties},
Theorem~6(a), printed p.503,
\url{https://www.wstein.org/papers/bib/Serre-Tate-Good_Reduction_of_Abelian_Varieties.pdf}.
\item Milne: \emph{Abelian Varieties}, version~2.00 (2008),
Remark~8.7(c), p.37; Proposition~14.2, pp.61--62;
Corollary~16.2(a), p.67; Corollary~16.4(a)--(b), p.68;
Propositions~9.2--9.3, p.41, and Proposition~10.1, p.42.
\item Bruinier--Howard--Kudla--Rapoport--Yang,
\emph{Modularity of generating series of divisors on unitary
Shimura varieties}, corrected 2017 author version,
\url{https://www.math.uni-bonn.de/people/rapoport/myalggeom/preprints/Modularity_arXiv_v1_corrected.pdf};
Proposition~3.4.3, pp.46--48; Proposition~3.6.1, pp.52--53;
Theorem~3.7.1, pp.53--55; Definition~5.1.1, p.67;
Proposition~5.1.4, pp.69--70. These are author-version pages.
\item Lan, \emph{Toroidal compactifications of PEL-type Kuga
families}, corrected 74-page author version,
\url{https://www.kwlan.org/articles/cpt-fam.pdf};
Section~3A, pp.27--30; Section~3B, p.31;
Lemma~4.9 and proof, p.46; author errata items (4), (8),
(13), (23), \url{https://www.kwlan.org/articles/cpt-fam-err.pdf}.
The published work is Algebra Number Theory 6 (2012),
885--966; its page numbering differs.
\item Schoen: \emph{Hodge classes on self-products of a variety
with an automorphism}, Compositio Mathematica 65 (1988),
3--32. Theorem~2.0, p.~11; Lemmas~2.2--2.9,
pp.~13--19; Corollary~3.1, pp.~24--25.
\url{https://www.numdam.org/item/CM_1988__65_1_3_0.pdf}.
C32 supplies a specialized complete order-four construction,
including the weight-two inertia argument.
\item Lan: \emph{Arithmetic compactifications of PEL-type
Shimura varieties}, corrected author revision,
\url{https://www.kwlan.org/articles/cpt-PEL-type-thesis-revision.pdf}.
Section~4.1, pp.~207--208; Theorem~4.4.16,
pp.~254--255; Construction~4.5.4.27 and
Lemma~4.5.4.28, pp.~299--300; Corollary~4.5.5.5
and its complete arbitrary-homomorphism proof, p.~303.
These locators refer to the corrected revision's printed
pages, not PDF page indices.
\item Milne: \emph{Algebraic Number Theory}, version 3.08,
19 July 2020, Theorem~8.31 and Corollary~8.32,
printed p.~147, \url{https://www.jmilne.org/math/CourseNotes/ANT.pdf}.
\item Milne: \emph{Canonical models of mixed Shimura varieties
and automorphic vector bundles}, corrected 2018 version of the
1990 work, \url{https://jmilne.org/math/xnotes/AA.pdf};
Chapter III, Remark~2.3 and Example~2.4 pp.55--56,
Theorem~5.1 and Remark~5.2 p.61, Sections~6--7 pp.61--63.
\item Deligne: \emph{Th\'eorie de Hodge II}, Publications
Math\'ematiques de l'IH\'ES 40 (1971), 5--57;
2.2.1--2.2.6 pp.27--29 and Theorem~4.1.1(i) p.40.
\item Kaletha--M\'inguez--Shin--White:
\emph{Endoscopic classification of representations: inner forms
of unitary groups I}, arXiv:1409.3731,
\url{https://arxiv.org/pdf/1409.3731};
Section~0.3.3 p.18, formula~(1.3.2) p.64,
Corollary~3.3.2 p.151, the scope paragraph after
Theorem~1.7.1 p.89, and Theorem~5.0.5 with proof
pp.205--206. Only the proved generic pure-inner case
is used for multiplicities.
\item Platonov--Rapinchuk: \emph{Algebraic Groups and Number
Theory} (1994), Theorem~7.12 p.427 and proof pp.428--432;
\url{https://uva.theopenscholar.com/files/andrei-rapinchuk/files/agnt_english.pdf}.
The final pp.431--432 permit omission of one finite
place without omitting the archimedean places.
\item Conrad: \emph{Gross--Zagier Revisited}, MSRI
Publications~49 (2004), Theorem~2.1(2), p.70,
\url{https://library.slmath.org/books/Book49/files/05conrad.pdf}.
\item Silverman: \emph{The Arithmetic of Elliptic Curves},
second edition (2009), Chapter~III, Exercise~3.7,
pp.105--106.
\end{itemize}

Alp\"oge locators use the pinned mirror commit given at the start
of the paper. The full relative file path, declaration, and where
available line numbers are recorded with the relevant finding.
\path{documentation/alpoge_mirror/pinned_manifest.json} and
\path{documentation/step2_archive_manifest.json} preserve recovery
URLs and hashes. A failed retrieval is explicitly recorded rather
than represented as a source reading. No source-module text was
executed, and no Lean build certificate is claimed.

The mathematical deliverable also contains:
\begin{itemize}[leftmargin=*]
\item \path{CHECK_EVENTS.tsv} and \path{QUESTIONS.tsv};
\item \path{LEDGER.md} and \path{IMPLICATIONS.md};
\item \path{TRANSCRIPTION_NOTES.tsv}.
\end{itemize}
An empty transcription register means that no
supplied text-of-record reading was available for a justified
Grothendieck correction; it is not a claim that all possible
transcriptions are correct.


\end{document}
